% Wspolna preambula: skrypt Logika matematyczna i teoria mnogosci
\documentclass[11pt,a4paper,oneside,openany]{book}
\usepackage[T1]{fontenc}
\usepackage[utf8]{inputenc}
\usepackage[polish]{babel}
\usepackage{lmodern}
\usepackage{microtype}
\usepackage{amsmath,amssymb,amsthm,mathtools}
\usepackage{stmaryrd} % wartosci boolowskie: \llbracket i \rrbracket
\usepackage{enumitem}
\usepackage{xcolor}
\usepackage{geometry}
\usepackage{tikz}
\usepackage{titlesec}
\usetikzlibrary{arrows.meta,positioning,calc}
\usepackage{xr-hyper}
\usepackage[unicode,colorlinks=true,linkcolor=blue!45!black,urlcolor=blue!45!black]{hyperref}
\externaldocument{model_theory}
\geometry{margin=27mm,headheight=14pt}
\setlength{\parindent}{0pt}
\setlength{\parskip}{2.5pt}
\setlength{\emergencystretch}{2em}
\setlist[itemize]{topsep=3pt,itemsep=2pt}
\setlist[enumerate]{topsep=3pt,itemsep=2pt}
\numberwithin{equation}{chapter}
\theoremstyle{plain}
\newtheorem{twierdzenie}{Twierdzenie}[chapter]
\newtheorem{lemat}[twierdzenie]{Lemat}
\newtheorem{wniosek}[twierdzenie]{Wniosek}
\newtheorem{stwierdzenie}[twierdzenie]{Stwierdzenie}
\theoremstyle{definition}
\newtheorem{definicja}[twierdzenie]{Definicja}
\newtheorem{przyklad}[twierdzenie]{Przykład}
\theoremstyle{remark}
\newtheorem{uwaga}[twierdzenie]{Uwaga}
\renewcommand{\proofname}{Dowód}
\newcommand{\N}{\mathbb N}
\newcommand{\Z}{\mathbb Z}
\newcommand{\Q}{\mathbb Q}
\newcommand{\R}{\mathbb R}
\newcommand{\calP}{\mathcal P}
\newcommand{\id}{\operatorname{id}}
\newcommand{\dom}{\operatorname{dom}}
\newcommand{\ran}{\operatorname{ran}}
\newcommand{\fld}{\operatorname{fld}}
\newcommand{\diag}{\operatorname{diag}}
\newcommand{\symdiff}{\mathbin{\triangle}}
\newcommand{\dotminus}{\mathbin{\dot-}}
\newcommand{\set}[1]{\left\{#1\right\}}

% Zwarty, jednostronny sklad skryptu: bez pustych stron dla parzystosci.
\raggedbottom
\titleformat{\chapter}[display]{\normalfont\bfseries\LARGE\raggedright}{\chaptername\ \thechapter}{5pt}{\LARGE}
\titlespacing*{\chapter}{0pt}{0pt}{16pt}
\titlespacing*{\section}{0pt}{10pt plus 2pt minus 1pt}{5pt}

% BEGIN GLOBAL CODE LISTINGS (generated from latex/pl/code-listings.tex)
\usepackage{xcolor}
\usepackage{listings}
\definecolor{codebg}{RGB}{248,249,251}
\definecolor{codeborder}{RGB}{180,191,203}
\definecolor{codecomment}{RGB}{61,113,79}
\definecolor{codekeyword}{RGB}{41,78,132}
\lstdefinelanguage{Julia}{
  morekeywords={abstract,baremodule,begin,break,catch,const,continue,do,else,elseif,
    end,export,false,finally,for,function,global,if,import,in,isa,let,local,macro,
    module,mutable,primitive,quote,return,struct,true,try,type,using,where,while,
    missing,nothing,println,print},
  sensitive=true,morecomment=[l]{\#},morecomment=[n]{\#=}{=\#},
  morestring=[b]",
}
\lstdefinestyle{skrypt}{
  backgroundcolor=\color{codebg},frame=single,rulecolor=\color{codeborder},
  framesep=5pt,frameround=tttt,framerule=.5pt,
  basicstyle=\ttfamily\footnotesize,keywordstyle=\color{codekeyword}\bfseries,
  commentstyle=\color{codecomment}\itshape,stringstyle=\color{purple!65!black},
  numbers=left,numberstyle=\tiny\color{gray!75!black},numbersep=8pt,
  stepnumber=1,breaklines=true,breakatwhitespace=false,
  columns=fullflexible,keepspaces=true,showstringspaces=false,tabsize=4,
  linewidth=.90\linewidth,captionpos=b,aboveskip=11pt,belowskip=12pt,
  xleftmargin=12pt,xrightmargin=2pt,
}
\lstdefinestyle{pythonpde}{style=skrypt,language=Python}
\lstdefinestyle{python}{style=skrypt,language=Python}
\lstdefinestyle{julia}{style=skrypt,language=Julia}
\lstset{style=skrypt}
\renewcommand{\lstlistingname}{Kod}
\renewcommand{\lstlistlistingname}{Spis kodów}
% END GLOBAL CODE LISTINGS
\begin{document}
\frontmatter
\hypersetup{pageanchor=false} % bez zdublowanej kotwicy strony tytulowej
\title{Logika matematyczna i teoria mnogości}
\author{}
\date{}
\maketitle
\hypersetup{pageanchor=true}
\tableofcontents
\mainmatter
\chapter{Zdania i spójniki logiczne}\label{ch:propositional}

\section{Po co matematyce logika?}\label{sec:logic-motivation}
Wyobraźmy sobie rozumowanie: „Jeżeli liczba jest podzielna przez $4$,
to jest parzysta. Liczba $12$ jest podzielna przez $4$, a zatem jest
parzysta”. Nie trzeba jeszcze znać żadnej zaawansowanej teorii, aby
zauważyć, że wniosek wynika z podanych przesłanek. Logika bada
\emph{formę} takich wnioskowań: co sprawia, że przejście od przesłanek
do wniosku jest poprawne, niezależnie od tego, czy rozmawiamy o
liczbach, figurach czy programach komputerowych.

Od początku oddzielamy dwa zagadnienia. \emph{Prawdziwość zdania}
dotyczy jego wartości logicznej w ustalonej interpretacji. \emph{Poprawność wnioskowania} dotyczy tego, czy
możliwa jest sytuacja, w której wszystkie przesłanki są prawdziwe,
ale wniosek fałszywy. Zobaczymy, że nie są to te same pytania.

\begin{przyklad}[Prawdziwa implikacja z fałszywym poprzednikiem]
„Jeżeli $2$ jest liczbą nieparzystą, to $2+2=4$”. Cała wypowiedź
jest prawdziwa w logice klasycznej, mimo że jej poprzednik jest
fałszywy. Znaczenie tego pozornie dziwnego faktu wyjaśnimy przy
implikacji, zamiast przyjmować go bez komentarza.
\end{przyklad}

\section{Zdania, pytania i wypowiedzi ze zmiennymi}\label{sec:statements}
\begin{definicja}[Zdanie w sensie logiki klasycznej]\label{def:statement}
\emph{Zdaniem} nazywamy wypowiedź oznajmującą, której w ustalonej
interpretacji można przypisać dokładnie jedną z dwóch wartości:
\emph{prawdę} (oznaczaną $1$) albo \emph{fałsz} (oznaczany $0$).
W tym rozdziale posługujemy się klasyczną logiką dwuwartościową.
\end{definicja}

\begin{uwaga}[Granice wstępnej definicji zdania]\label{rem:statement-scope-xli}
Samo oznajmujące brzmienie wypowiedzi nie gwarantuje, że jest
ona zdaniem o jednoznacznej wartości logicznej. W szczególności
samoodnosząca się wypowiedź „To zdanie jest fałszywe” prowadzi
przy naiwnym zastosowaniu reguły dwuwartościowości do sprzeczności:
jeżeli jest prawdziwe, to --- zgodnie ze swą treścią --- jest
fałszywe, a jeżeli jest fałszywe, to jego treść jest prawdziwa.
Dlatego w ścisłym rachunku logicznym ustalamy \emph{język}
i jego \emph{interpretację}; nie przyznajemy automatycznie wartości
logicznych wszystkim napisom języka naturalnego. Składnię
rachunku zdań zdefiniujemy niżej, a składnię logiki pierwszego
rzędu --- w rozdziale~\ref{ch:syntax-fol}.
\end{uwaga}

\begin{przyklad}[Zdania i wypowiedzi niebędące zdaniami]
Wypowiedzi „$2+3=5$” oraz „$7<4$” są zdaniami: pierwsze jest
prawdziwe, drugie fałszywe. Pytanie „Ile wynosi $2+3$?” i polecenie
„Oblicz tę sumę” nie są zdaniami w powyższym sensie. Napis „$x>4$”
nie jest jeszcze zdaniem o ustalonej wartości logicznej: jego
prawdziwość zależy od wartości $x$. Po podstawieniu $x=7$ otrzymujemy
zdanie prawdziwe, a po podstawieniu $x=2$ --- fałszywe.
\end{przyklad}

\begin{uwaga}[Zdanie a umiejętność rozstrzygnięcia prawdy]
Definicja nie wymaga, abyśmy potrafili rozpoznać, czy dane zdanie
jest prawdziwe. Matematyczne twierdzenie może pozostawać nierozstrzygnięte
przez wiele lat, a mimo to w ustalonej interpretacji klasycznej
jest zdaniem. Nie wolno utożsamiać „nie znam odpowiedzi” z
„wypowiedź nie ma wartości logicznej”.
\end{uwaga}

\begin{definicja}[Zdania proste i złożone]
Litery zdaniowe $p,q,r,\ldots$ są symbolami, które po ustaleniu
interpretacji mogą reprezentować całe zdania (na przykład „$6$ jest
parzyste”), a nie liczby. Za ich pomocą i przy użyciu spójników
budujemy \emph{formuły zdaniowe}. W rachunku zdań litery traktujemy
jako symbole formalne: w różnych wartościowaniach ta sama litera
może otrzymać różną wartość logiczną. Po zastąpieniu liter
konkretnymi zdaniami o ustalonych wartościach logicznych
otrzymujemy zdanie; jest ono złożone, jeżeli użyliśmy
spójników. Nawiasy określają budowę formuły.
\end{definicja}

\begin{definicja}[Skończona budowa formuł rachunku zdań]
\label{def:propositional-grammar-xli}
Wybierzmy litery zdaniowe $p_1,p_2,\ldots$.
\emph{Formuły rachunku zdań} powstają według następujących reguł:
\begin{enumerate}[label=(\roman*)]
\item każda litera zdaniowa jest formułą;
\item jeżeli $\alpha$ jest formułą, to $\neg\alpha$ jest formułą;
\item jeżeli $\alpha,\beta$ są formułami, to
      $(\alpha\land\beta)$, $(\alpha\lor\beta)$,
      $(\alpha\Rightarrow\beta)$ i
      $(\alpha\Leftrightarrow\beta)$ są formułami;
\item formułami są tylko napisy otrzymane przez skończoną liczbę
      zastosowań reguł (i)--(iii).
\end{enumerate}
Jest to \emph{definicja indukcyjna składni}: mówimy, z jakich
prostych napisów buduje się poprawne napisy, jeszcze bez pytania
o ich prawdziwość. Umowy dotyczące opuszczania nawiasów
podamy w sekcji~\ref{sec:biconditional}.
\end{definicja}

\begin{przyklad}[Jednoznaczny rozbiór nawiasów]\label{ex:parse-tree-xli}
Napis $p\Rightarrow(\neg q\lor r)$ ma w pełnym zapisie postać
$(p\Rightarrow(\neg q\lor r))$. Najpierw powstają
formuły $p,q,r$; następnie $\neg q$ i
$(\neg q\lor r)$, a na końcu cała implikacja.
Inny napis $((p\Rightarrow\neg q)\land r)$ ma inną budowę.
Gdy $p=0$, $q=1$, $r=0$, pierwszy ma wartość $1$
(fałszywy poprzednik implikacji), a drugi $0$
(jeden z członów koniunkcji jest fałszywy). Nie wolno więc
usuwać nawiasów w sposób zmieniający budowę formuły.
\end{przyklad}

\begin{center}
\begin{tikzpicture}[
  level distance=11mm,
  level 1/.style={sibling distance=36mm},
  level 2/.style={sibling distance=24mm},
  every node/.style={draw,rounded corners=2pt,inner sep=3pt,font=\small}]
\node {$\Rightarrow$}
 child {node {$p$}}
 child {node {$\lor$}
   child {node {$\neg$} child {node {$q$}}}
   child {node {$r$}}};
\end{tikzpicture}
\par\smallskip
\small Schemat budowy formuły
$p\Rightarrow(\neg q\lor r)$: liście są literami, węzły spójnikami.
\end{center}

\section{Negacja, koniunkcja i alternatywa}\label{sec:connectives-first}
\begin{definicja}[Trzy podstawowe spójniki]\label{def:not-and-or}
Niech $p,q$ oznaczają zdania.
\begin{itemize}
\item \emph{Negacja} $\neg p$ (czytamy: „nie $p$”) jest prawdziwa
      dokładnie wtedy, gdy $p$ jest fałszywe.
\item \emph{Koniunkcja} $p\land q$ („$p$ i $q$”) jest prawdziwa
      dokładnie wtedy, gdy \emph{oba} zdania są prawdziwe.
\item \emph{Alternatywa} $p\lor q$ („$p$ lub $q$”) jest prawdziwa
      dokładnie wtedy, gdy \emph{co najmniej jedno} z dwóch zdań
      jest prawdziwe. W szczególności oba mogą być prawdziwe.
\end{itemize}
Reguły te wyraża następująca tablica prawdy:
\[
\begin{array}{c c|c c c}
p&q&\neg p&p\land q&p\lor q\\\hline
1&1&0&1&1\\
1&0&0&0&1\\
0&1&1&0&1\\
0&0&1&0&0
\end{array}
\]
\end{definicja}

\begin{przyklad}[„Lub” w języku naturalnym]
Niech $p$ oznacza „$6$ jest parzyste”, a $q$ „$6$ jest
podzielne przez $3$”. Oba zdania są prawdziwe, a więc także
$p\lor q$ jest prawdziwe. Jest to \emph{lub niewykluczające}.
W języku codziennym „albo herbata, albo kawa” może wykluczać
jednoczesne wystąpienie obu zdarzeń. W logice matematycznej
\emph{wyłączność} trzeba wyrazić osobno: $(p\lor q)\land\neg(p\land q)$.
\end{przyklad}

\begin{przyklad}[Negacja zdania złożonego]
Ustalmy liczbę całkowitą $n$. Niech $p$ oznacza zdanie „$n$ jest
liczbą dodatnią”, a $q$ zdanie „$n$ jest liczbą parzystą”. Negacja
zdania „$n$ jest liczbą dodatnią i parzystą” nie brzmi „$n$ jest
liczbą niedodatnią i nieparzystą”. Wystarczy, że zawiedzie
\emph{jeden} z dwóch warunków: $n$ nie jest liczbą dodatnią
\emph{lub} nie jest liczbą parzystą. Za chwilę uzasadnimy to formalnie.
\end{przyklad}


\begin{uwaga}[Jak czytać nawiasy w zapisie logicznym?]
Jeżeli $p$ oznacza „pada deszcz”, a $q$ „wieje wiatr”, to
$\neg(p\land q)$ oznacza zaprzeczenie całego zdania „pada i wieje”.
Natomiast $(\neg p\land q)$ oznacza „nie pada, ale wieje”.
To nie są te same wypowiedzi: gdy pada deszcz, lecz nie wieje,
pierwsza jest prawdziwa, a druga fałszywa. Dlatego w tym skrypcie
nawiasy nie są ozdobą: wskazują dokładnie, do jakiej części
wypowiedzi należy dany spójnik.
\end{uwaga}

\section{Implikacja: kiedy „jeżeli..., to...” jest prawdziwe?}
\label{sec:implication}
\begin{definicja}[Implikacja materialna]\label{def:implication}
Dla zdań $p,q$ \emph{implikacją} $p\Rightarrow q$ nazywamy
zdanie fałszywe dokładnie wtedy, gdy $p$ jest prawdziwe, a $q$
fałszywe. Zdanie $p$ to \emph{poprzednik} (założenie), a $q$ to
\emph{następnik} (konkluzja) implikacji.
\[
\begin{array}{c c|c}
p&q&p\Rightarrow q\\\hline
1&1&1\\
1&0&0\\
0&1&1\\
0&0&1
\end{array}
\]
\end{definicja}

\paragraph{Dlaczego dokładnie taka tablica?}
Warunek „jeżeli $p$, to $q$” ma wykluczać sytuację, gdy $p$
zachodzi, a $q$ nie zachodzi. Właśnie taka sytuacja jest
\emph{kontrprzykładem} do implikacji. Gdy $p$ nie zachodzi,
obietnica nie została naruszona, niezależnie od wartości $q$.
Tak rozumiana implikacja jest regułą o wartości logicznej, a nie
stwierdzeniem związku przyczynowego między dwiema wypowiedziami.

\begin{przyklad}[Fałszywy poprzednik nie daje kontrprzykładu]
Zdanie „jeżeli $3$ jest parzyste, to $5$ jest większe od $2$”
jest prawdziwe. Prawdziwe jest także zdanie „jeżeli $3$ jest
parzyste, to $5$ jest mniejsze od $2$”. W obu przypadkach
poprzednik jest fałszywy, a zatem nie występuje układ
„prawdziwy poprzednik --- fałszywy następnik”.
\end{przyklad}

\begin{definicja}[Warunek konieczny i wystarczający]
Niech $p,q$ będą formułami zdaniowymi i niech zostanie ustalony
zbiór dopuszczalnych wartościowań ich liter. Jeżeli
$p\Rightarrow q$ jest prawdziwa przy każdym z tych wartościowań,
to $p$ nazywamy \emph{warunkiem wystarczającym} dla $q$,
a $q$ --- \emph{warunkiem koniecznym} dla $p$.
Przy braku dodatkowych ograniczeń rozważamy wszystkie
wartościowania. Nie wystarcza sprawdzenie prawdziwości
implikacji tylko przy jednym wartościowaniu, na przykład takim,
w którym $p$ jest fałszywe. Jeżeli przy każdym dopuszczalnym
wartościowaniu prawdziwe są obie implikacje $p\Rightarrow q$
i $q\Rightarrow p$, to $p$ jest dla $q$ warunkiem zarówno
koniecznym, jak i wystarczającym. Są to sposoby opisu
zależności logicznych, a nie osobne twierdzenia.
\end{definicja}

\begin{przyklad}[Konieczny nie znaczy wystarczający]
Rozpatrujemy liczby całkowite $n$ oraz zdania „$4\mid n$” i „$2\mid n$”.
Dopuszczalne są tutaj tylko te pary wartości logicznych, które
pochodzą od tej samej liczby $n$, a nie wszystkie niezależne
przypisania wartości dwóm literom zdaniowym. Jeżeli $4\mid n$,
to $n=4k=2(2k)$ dla pewnego $k\in\mathbb Z$, zatem $2\mid n$.
Podzielność przez $4$ jest więc warunkiem wystarczającym
parzystości, a parzystość --- koniecznym warunkiem podzielności
przez $4$. Nie jest jednak warunkiem wystarczającym: liczba $6$
jest parzysta, ale nie jest podzielna przez $4$. Nie oznacza to,
że sama formuła $p\Rightarrow q$ jest tautologią, gdy $p$ i $q$
traktujemy jako niezależne litery zdaniowe.
\end{przyklad}

\begin{definicja}[Implikacja odwrotna, przeciwna i przeciwstawna]
Dla $p\Rightarrow q$ rozważamy:
\begin{itemize}
\item \emph{implikację odwrotną} $q\Rightarrow p$;
\item \emph{implikację przeciwną} $\neg p\Rightarrow\neg q$;
\item \emph{implikację przeciwstawną} $\neg q\Rightarrow\neg p$.
\end{itemize}
W dalszym ciągu udowodnimy, że implikacja jest logicznie
równoważna swojej implikacji przeciwstawnej. Implikacje odwrotna
i przeciwna na ogół nie są z nią logicznie równoważne.
\end{definicja}

\begin{przyklad}[Dlaczego nie wystarczy odwrócić strzałki?]
„Jeżeli liczba jest podzielna przez $4$, to jest parzysta”
jest prawdziwe. Zdanie odwrotne „jeżeli liczba jest parzysta,
to jest podzielna przez $4$” jest fałszywe: kontrprzykładem
jest $6$. Zdanie przeciwstawne „jeżeli liczba nie jest
parzysta, to nie jest podzielna przez $4$” jest prawdziwe.
\end{przyklad}

\section{Równoważność i inne użyteczne skróty}\label{sec:biconditional}
\begin{definicja}[Równoważność]\label{def:biconditional}
Dla zdań $p,q$ formuła $p\Leftrightarrow q$
(czytamy: „$p$ wtedy i tylko wtedy, gdy $q$”)
jest prawdziwa dokładnie wtedy, gdy $p$ i $q$ mają tę samą
wartość logiczną. Odpowiada koniunkcji obu implikacji:
$(p\Rightarrow q)\land(q\Rightarrow p)$.
\[
\begin{array}{c c|c}
p&q&p\Leftrightarrow q\\\hline
1&1&1\\
1&0&0\\
0&1&0\\
0&0&1
\end{array}
\]
\end{definicja}

\begin{przyklad}[Dowód równoważności wymaga dwóch kierunków]
Zdanie „liczba całkowita $n$ jest parzysta wtedy i tylko wtedy,
gdy $n$ da się zapisać jako $2k$ dla pewnej liczby całkowitej $k$”
to w istocie definicja parzystości. Aby udowodnić \emph{inną}
równoważność, często trzeba odrębnie przeprowadzić dowody
$p\Rightarrow q$ oraz $q\Rightarrow p$.
\end{przyklad}

\begin{uwaga}[Umowa dotycząca nawiasów]
Najwyższy priorytet ma negacja $\neg$, potem koniunkcja
$\land$, następnie alternatywa $\lor$, implikacja $\Rightarrow$
i równoważność $\Leftrightarrow$. Jeżeli jakiekolwiek wątpliwości
są możliwe, zapisujemy nawiasy. Dla przykładu
$\neg p\land q$ oznacza $(\neg p\land q)$, a nie
$\neg(p\land q)$. Ciągi spójników dwuargumentowych tego samego
priorytetu również zapisujemy z nawiasami: nie posługujemy się
niejednoznacznym zapisem $p\Rightarrow q\Rightarrow r$.
Skrót $\varphi_1\land\cdots\land\varphi_n$ ($n\ge1$)
oznacza natomiast koniunkcję grupowaną od lewej: jest to
$C_n$, gdzie $C_1=\varphi_1$ oraz
$C_{k+1}=(C_k\land\varphi_{k+1})$ dla $1\le k<n$.
To jedynie konwencja zapisu, nie nowy spójnik.
\end{uwaga}

\section{Jak badać formuły? Wartościowania i tablice prawdy}
\label{sec:valuations}
\begin{definicja}[Wartościowanie]\label{def:valuation}
\emph{Wartościowaniem} dla rozważanej formuły nazywamy
przypisanie każdej występującej w niej literze zdaniowej
wartości $0$ lub $1$. Wartość całej formuły obliczamy od
najbardziej wewnętrznych jej części, korzystając z tablic
spójników logicznych. Ponieważ każde wystąpienie spójnika ma
argumenty będące krótszymi formułami, takie obliczanie kończy
się po skończonej liczbie kroków. Dwa wystąpienia tej samej litery
mają w jednym wartościowaniu tę samą wartość. Rozszerzoną
w ten sposób wartość formuły $\varphi$ oznaczamy $v(\varphi)$.
\end{definicja}

\begin{twierdzenie}[Jednoznaczność wartości formuły]
\label{thm:valuation-well-defined-xli}
Niech $\varphi$ będzie formułą zbudowaną zgodnie
z definicją~\ref{def:propositional-grammar-xli} i niech
$v$ przypisuje każdej literze występującej w $\varphi$
wartość $0$ albo $1$. Wówczas z tablic spójników można
wyznaczyć \emph{dokładnie jedną} wartość $v(\varphi)$.
\end{twierdzenie}
\begin{proof}
Rozpatrujemy kolejne etapy budowy formuły. W etapie
początkowym formuła jest pojedynczą literą; ma wartość
przyznaną jej przez $v$. Załóżmy, że wartości wszystkich
formuł powstałych w pierwszych $m$ etapach są jednoznaczne.
Jeśli w kolejnym etapie utworzono negację $\neg\alpha$,
to wartość $\alpha$ jest już znana, a tabela negacji
wskazuje dokładnie jedną wartość nowej formuły.
Jeżeli utworzono $(\alpha\circ\beta)$, gdzie $\circ$
jest jednym z czterech spójników dwuargumentowych,
to obie wartości $v(\alpha),v(\beta)$ są znane;
odpowiednia tabela wskazuje dokładnie jedną wartość wyniku.
Budowa $\varphi$ jest skończona, więc po skończonej
liczbie takich kroków otrzymamy jej wartość.
W pełnym zapisie pojedyncza litera jest rozpoznawana
jednoznacznie. Każda inna formuła zaczyna się od $\neg$
albo od nawiasu otwierającego. W pierwszym przypadku reszta
napisu stanowi argument negacji. W drugim spójnik dwuargumentowy
znajduje się wewnątrz zewnętrznej pary nawiasów, poza wszystkimi
wewnętrznymi parami; położenie tego spójnika wyznacza obie
podformuły. Stąd rozbiór formuły jest jednoznaczny.
Powtarzając to rozumowanie dla podformuł, widzimy, że żaden
krok obliczania wartości nie pozostawia dowolności.
\end{proof}

\begin{przyklad}[Obliczenie od wewnątrz]
Rozważmy $\varphi=(p\Rightarrow q)\land\neg q$.
Gdy $p=0$ i $q=0$, otrzymujemy kolejno:
$p\Rightarrow q=1$, $\neg q=1$, a zatem $\varphi=1$.
Gdy $p=1$ i $q=0$, mamy $p\Rightarrow q=0$ i
$\neg q=1$, więc $\varphi=0$. Zmiana jednego zdania
prostego może zmienić wartość całej formuły.
\end{przyklad}

\begin{definicja}[Tautologia, sprzeczność, spełnialność]\label{def:tautology}
Formułę nazywamy \emph{tautologią}, gdy przy \emph{każdym}
wartościowaniu jest prawdziwa. Jest \emph{sprzeczna}, gdy
przy każdym wartościowaniu jest fałszywa.
Jest \emph{spełnialna}, gdy istnieje \emph{przynajmniej jedno}
wartościowanie, w którym jest prawdziwa. Formuły
$\varphi$ i $\psi$ są \emph{logicznie równoważne},
gdy przy każdym wspólnym wartościowaniu mają taką samą
wartość logiczną; zapisujemy $\varphi\equiv\psi$.
Symbol $\equiv$ opisuje relację między formułami;
nie jest kolejnym spójnikiem tworzącym zdanie złożone.
Przy porównywaniu dwóch formuł wartościowanie obejmuje
\emph{wszystkie} litery występujące w którejkolwiek z nich,
także jeśli część liter występuje tylko w jednej formule.
\end{definicja}

\begin{przyklad}[Trzy różne sytuacje]
$p\lor\neg p$ jest tautologią, gdyż niezależnie od wartości
$p$ co najmniej jedno ze zdań $p,\neg p$ jest prawdziwe.
$p\land\neg p$ jest sprzecznością, gdyż nigdy nie mogą
być prawdziwe równocześnie. Formuła $p\land q$ jest
spełnialna (weźmy $p=q=1$), ale nie jest tautologią
(weźmy $p=1,q=0$).
\end{przyklad}

\begin{twierdzenie}[Test równoważności]\label{thm:equiv-test}
Dla dowolnych formuł zdaniowych $\varphi,\psi$ zachodzi:
$\varphi\equiv\psi$ wtedy i tylko wtedy, gdy
$\varphi\Leftrightarrow\psi$ jest tautologią.
\end{twierdzenie}
\begin{proof}
Ustalmy dowolne wartościowanie występujących w formułach
liter zdaniowych. Z definicji równoważności spójnikowej
$\varphi\Leftrightarrow\psi$ jest w tym wartościowaniu
prawdziwe dokładnie wtedy, gdy obie formuły mają tę samą
wartość. Zatem żądanie, aby działo się tak dla wszystkich
wartościowań, jest jednocześnie definicją tautologii
$\varphi\Leftrightarrow\psi$ i definicją relacji
$\varphi\equiv\psi$. Dowiedliśmy obu kierunków.
\end{proof}

\begin{twierdzenie}[Skończony test tablicowy]\label{thm:truth-table-decision-xli}
Jeżeli formuła $\varphi$ zawiera dokładnie $m$ różnych
liter zdaniowych, gdzie $m\ge1$, to ma dokładnie $2^m$
wartościowań. Zbadanie wszystkich tych wartościowań
pozwala rozstrzygnąć, czy $\varphi$ jest tautologią,
sprzecznością albo formułą spełnialną.
\end{twierdzenie}
\begin{proof}
Pierwszej literze możemy przypisać dwie wartości.
Po ustaleniu wartości pierwszych $k$ liter dla
następnej są znów dwa niezależne wybory. Zatem
liczba pełnych przypisań wynosi $2\cdot\ldots\cdot2=2^m$.
Każde przypisanie określa wartość $\varphi$ na mocy
twierdzenia~\ref{thm:valuation-well-defined-xli}.
Jeżeli każda z $2^m$ wartości jest równa $1$, definicja
daje tautologię. Jeśli każda jest równa $0$, otrzymujemy
sprzeczność. W pozostałym przypadku pojawia się co
najmniej jedna wartość $1$ i co najmniej jedna wartość $0$:
formuła jest spełnialna, ale nie jest tautologią.
\end{proof}

\begin{przyklad}[Trzy litery, lecz tylko osiem przypadków]
Rozważmy $(p\land q)\Rightarrow r$. Litery $p,q,r$
są trzy, więc wartościowań jest $2^3=8$. Implikacja
ma wartość $0$ tylko przy $p=q=1$, $r=0$;
w każdym z siedmiu pozostałych wierszy ma wartość $1$.
Jest zatem spełnialna, ale nie jest tautologią.
Powtórzone wystąpienia tej samej litery nie są
nowymi niezależnymi wyborami, np. $p\land p$
ma tylko dwa wartościowania, a nie cztery.
\end{przyklad}

\section{Podstawowe prawa rachunku zdań z dowodami}\label{sec:propositional-laws}
\begin{twierdzenie}[Prawa de Morgana]\label{thm:de-morgan-logic}
Dla dowolnych formuł $p,q$ zachodzą równoważności logiczne
\[
\neg(p\land q)\equiv(\neg p\lor\neg q),\qquad
\neg(p\lor q)\equiv(\neg p\land\neg q).
\]
\end{twierdzenie}
\begin{proof}
Pokażemy wszystkie cztery możliwe pary wartości formuł
$p$ i $q$. Jeżeli formuły mają wspólne litery, nie każda
para musi wystąpić przy pewnym wartościowaniu. Wystarcza jednak
sprawdzenie wszystkich czterech możliwości, bo każde
wartościowanie daje jedną z nich.
\[
\begin{array}{cc|cc|cc}
p&q&\neg(p\land q)&\neg p\lor\neg q&
\neg(p\lor q)&\neg p\land\neg q\\\hline
1&1&0&0&0&0\\
1&0&1&1&0&0\\
0&1&1&1&0&0\\
0&0&1&1&1&1
\end{array}
\]
Dwie kolumny opisujące pierwszą parę formuł są identyczne
w każdym wierszu; to samo dotyczy drugiej pary.
Z definicji równoważności logicznej otrzymujemy obie tezy.
\end{proof}

\begin{twierdzenie}[Rozwinięcie i przeciwstawienie implikacji]
\label{thm:implication-laws}
Dla dowolnych formuł $p,q$ zachodzą
\[
(p\Rightarrow q)\equiv(\neg p\lor q),\qquad
(p\Rightarrow q)\equiv(\neg q\Rightarrow\neg p).
\]
W szczególności $(p\Rightarrow q)\land p\Rightarrow q$
jest tautologią (\emph{modus ponens}).
\end{twierdzenie}
\begin{proof}
Rozpatrzmy cztery możliwe pary wartości formuł $p,q$;
nie wymagamy, aby każda z nich była realizowana przy pewnym
wspólnym wartościowaniu:
\[
\begin{array}{cc|ccc}
p&q&p\Rightarrow q&\neg p\lor q&\neg q\Rightarrow\neg p\\\hline
1&1&1&1&1\\
1&0&0&0&0\\
0&1&1&1&1\\
0&0&1&1&1
\end{array}
\]
W każdym wierszu trzy ostatnie wartości są równe, co
dowodzi obu równoważności. Aby sprawdzić ostatnią tezę,
rozpatrzmy dowolne wartościowanie. Jeśli koniunkcja
$(p\Rightarrow q)\land p$ jest fałszywa, całe zdanie
warunkowe jest prawdziwe z definicji implikacji.
Jeśli koniunkcja jest prawdziwa, to $p$ jest prawdziwe,
a ponieważ $p\Rightarrow q$ też jest prawdziwe, $q$
nie może być fałszywe: w przeciwnym razie implikacja
byłaby fałszywa. Zatem $q$ jest prawdziwe, więc
cała rozpatrywana implikacja jest prawdziwa również
i w tym przypadku.
\end{proof}

\begin{twierdzenie}[Negacja implikacji]\label{thm:negated-implication-xli}
Dla dowolnych formuł $p,q$ zachodzi
$\neg(p\Rightarrow q)\equiv p\land\neg q$.
\end{twierdzenie}
\begin{proof}
Ustalmy wartościowanie. Na mocy definicji implikacji
$p\Rightarrow q$ jest fałszywe wtedy i tylko wtedy,
gdy jednocześnie $p=1$ i $q=0$. Jego negacja jest
więc prawdziwa dokładnie przy $p=1$ i $q=0$.
Z drugiej strony $p\land\neg q$ jest prawdziwe
dokładnie wtedy, gdy $p=1$ oraz $\neg q=1$,
czyli $p=1$ i $q=0$. W obu przypadkach zbiór
wartościowań dających wartość $1$ jest taki sam,
a pozostałe wartościowania dają $0$. Otrzymaliśmy zatem
żądaną równoważność logiczną.
\end{proof}

\begin{przyklad}[Obliczenie kontrprzykładu do implikacji]
Niech $p$ oznacza „$12$ jest podzielne przez $4$”,
a $q$ --- „$12$ jest liczbą nieparzystą”. Wtedy
$p=1$, $q=0$, więc $p\Rightarrow q=0$ oraz
$p\land\neg q=1$. Zmienienie poprzednika na
zdanie fałszywe usunęłoby kontrprzykład do implikacji.
\end{przyklad}

\begin{twierdzenie}[Podwójna negacja i rozdzielność]
\label{thm:logic-distributivity}
Dla dowolnych formuł $p,q,r$ zachodzą
\begin{align*}
\neg\neg p&\equiv p,\\
p\land(q\lor r)&\equiv(p\land q)\lor(p\land r),\\
p\lor(q\land r)&\equiv(p\lor q)\land(p\lor r).
\end{align*}
\end{twierdzenie}
\begin{proof}
Dla podwójnej negacji wystarczy zauważyć, że przy
$p=1$ mamy $\neg p=0$ i $\neg\neg p=1$, zaś przy
$p=0$ mamy $\neg p=1$ i $\neg\neg p=0$.

Wykażmy pierwszą rozdzielność, analizując dwa przypadki.
Jeśli $p=0$, to lewa strona jest fałszywa, a prawa
również: oba składniki jej alternatywy zawierają
fałszywy czynnik $p$. Jeśli $p=1$, to lewa strona
przyjmuje dokładnie wartość $q\lor r$, a prawa strona
jest $(1\land q)\lor(1\land r)$, która także
przyjmuje wartość $q\lor r$. Oba przypadki wyczerpują
wszystkie wartościowania, ponieważ $p$ musi mieć wartość
$0$ lub $1$, niezależnie od wartości $q,r$.

Podobnie przy $p=1$ obie strony drugiej rozdzielności
są prawdziwe, a przy $p=0$ obie sprowadzają się do
$q\land r$. Stąd także druga równoważność.
\end{proof}

\begin{uwaga}[Dlaczego przykład liczbowy nie zastępuje dowodu?]
Sprawdzenie prawa logicznego na jednym przykładzie potwierdza
jedynie jedno wartościowanie. Aby wykazać tautologię,
trzeba zbadać \emph{wszystkie} wartościowania albo podać
ogólny argument obejmujący każdą możliwość. Z kolei do
\emph{obalenia} zdania głoszącego równoważność wystarczy
\emph{jedno} wartościowanie, przy którym formuły są różne.
\end{uwaga}

\section{Postacie normalne i zupełność funkcyjna}
\label{sec:lst-propositional-normal-forms}
Tablica prawdy pozwala odczytać wartości danej formuły.
Można również rozwiązać zadanie odwrotne: mając ustaloną
tablicę wartości, zbudować formułę, która ją realizuje.
W ten sposób otrzymamy systematyczną metodę zapisywania
i upraszczania formuł.

\begin{definicja}[Literały i postacie normalne]
\label{def:lst-dnf-cnf}
Literałem nazywamy literę zdaniową $p$ albo jej negację $\neg p$.
Formuła jest w alternatywnej postaci normalnej
(ang. \emph{disjunctive normal form}, DNF), gdy jest
alternatywą koniunkcji literałów. Jest w koniunkcyjnej
postaci normalnej (ang. \emph{conjunctive normal form}, CNF),
gdy jest koniunkcją alternatyw literałów.
Alternatywę literałów w CNF nazywamy klauzulą.
\end{definicja}
Pojedynczy literał jest zarówno koniunkcją, jak i alternatywą
jednoelementową. Dla $n\geq1$ puste alternatywy i koniunkcje
możemy zastępować odpowiednio przez $p_1\land\neg p_1$
oraz $p_1\lor\neg p_1$. Symbole $\bot,\top$ traktujemy
w tej części jako skróty tych formuł, a nie nowe symbole
w definicji składni z
Definicji~\ref{def:propositional-grammar-xli}.

\begin{twierdzenie}[Konstrukcja postaci normalnych z tablicy prawdy]
\label{thm:lst-canonical-normal-forms}
Niech $n\geq1$ i $f:\{0,1\}^n\to\{0,1\}$.
Dla $a=(a_1,\ldots,a_n)\in\{0,1\}^n$ określamy
\[
\begin{split}
 \ell_i^a&=\begin{cases}p_i,&a_i=1,\\\neg p_i,&a_i=0,\end{cases}\\
 m_a&=\bigwedge_{i=1}^n\ell_i^a,\qquad
 c_a=\bigvee_{i=1}^n\neg\ell_i^a.
\end{split}
\]
W wyrażeniu $c_a$ podwójną negację zastępujemy
odpowiednim literałem dodatnim.
Wówczas formuły
\[
 D_f=\bigvee_{a:\,f(a)=1}m_a,\qquad
 C_f=\bigwedge_{a:\,f(a)=0}c_a
\]
mają tablicę prawdy określoną przez $f$. W szczególności
każda formuła jest równoważna pewnej formule w DNF
oraz pewnej formule w CNF.
\end{twierdzenie}
\begin{proof}
Ustalmy wartościowanie $p_i=b_i$ dla
$b=(b_1,\ldots,b_n)\in\{0,1\}^n$.
Literał $\ell_i^a$ ma wartość $1$ dokładnie wtedy,
gdy $b_i=a_i$. Koniunkcja $m_a$ ma zatem wartość $1$
dokładnie dla $b=a$. Alternatywa wszystkich $m_a$,
dla których $f(a)=1$, ma wartość $1$ wtedy i tylko wtedy,
gdy $f(b)=1$. Otrzymujemy tablicę $D_f$.
Z kolei $c_a$ ma wartość $0$ dokładnie wtedy,
gdy wszystkie $\neg\ell_i^a$ mają wartość $0$,
czyli gdy $b=a$. Koniunkcja klauzul odpowiadających
wierszom, w których $f$ przyjmuje $0$, ma wartość $0$
dokładnie w tych wierszach. Otrzymujemy tablicę $C_f$.
W przypadku funkcji stałej obowiązuje podana wyżej
konwencja dotycząca pustych działań. Dla tablicy
dowolnej formuły podstawiamy jej funkcję wartości
za $f$ i stosujemy definicję równoważności logicznej.
\end{proof}
Wszystkie litery $p_1,\ldots,p_n$ występujące
w poszczególnych członach dają postacie pełne (kanoniczne).
Dla ustalonej kolejności liter i wierszy wybór członów
jest jednoznaczny, ponieważ każdy odpowiada dokładnie
jednemu wierszowi tablicy prawdy. Nie oznacza to
jednoznaczności dowolnej DNF lub CNF: równoważna
formuła może zawierać mniej literałów.

\begin{przyklad}[Dokładnie jedna z dwóch przesłanek]
\label{ex:lst-xor-normal-forms}
Niech $f(p,q)$ ma wartość $1$, gdy dokładnie jedna
z liter $p,q$ ma wartość $1$.
\[
\begin{array}{c|c|c|l|l}
 p&q&f&\text{człon DNF}&\text{klauzula CNF}\\ \hline
 0&0&0&-&p\lor q\\
 0&1&1&\neg p\land q&-\\
 1&0&1&p\land\neg q&-\\
 1&1&0&-&\neg p\lor\neg q
\end{array}
\]
Z wierszy prawdziwych odczytujemy
$(\neg p\land q)\lor(p\land\neg q)$,
a z fałszywych $(p\lor q)\land(\neg p\lor\neg q)$.
Obie formuły wyrażają tę samą funkcję alternatywy
rozłącznej. W drugim zapisie pierwszy nawias wymaga
co najmniej jednej prawdy, a drugi wyklucza dwie prawdy.
\end{przyklad}

\begin{przyklad}[Upraszczanie postaci pełnej]
\label{ex:lst-dnf-simplification}
Z dwóch sąsiednich wierszy tablicy, różniących się
tylko wartością $q$, możemy otrzymać krótszy człon:
\[
 (p\land q)\lor(p\land\neg q)
 \equiv p\land(q\lor\neg q)\equiv p.
\]
Pierwsza równoważność wynika z rozdzielności
(Twierdzenie~\ref{thm:logic-distributivity}),
a druga z faktu, że $q\lor\neg q$ przyjmuje zawsze
wartość $1$. Konstrukcja z tablicy daje więc
poprawną formułę, lecz nie gwarantuje najmniejszej
liczby spójników.
\end{przyklad}

\begin{definicja}[Zupełność funkcyjna spójników]
\label{def:lst-functional-completeness}
Rodzinę spójników nazywamy funkcyjnie zupełną, jeżeli
dla każdego $n\geq1$ i każdej funkcji
$f:\{0,1\}^n\to\{0,1\}$ istnieje formuła w literach
$p_1,\ldots,p_n$, używająca wyłącznie spójników
tej rodziny, której funkcją wartości jest $f$.
\end{definicja}

\begin{twierdzenie}\label{thm:lst-complete-connectives}
Funkcyjnie zupełne są rodziny
$\{\neg,\land,\lor\}$, $\{\neg,\land\}$,
$\{\neg,\lor\}$ i $\{\neg,\Rightarrow\}$.
\end{twierdzenie}
\begin{proof}
Pierwsza teza wynika z konstrukcji DNF
w Twierdzeniu~\ref{thm:lst-canonical-normal-forms}.
Alternatywę można zastąpić przez
$\neg(\neg p\land\neg q)$, a koniunkcję przez
$\neg(\neg p\lor\neg q)$, korzystając z praw
de Morgana (Twierdzenie~\ref{thm:de-morgan-logic}).
Dla rodziny $\{\neg,\Rightarrow\}$ mamy
$p\lor q\equiv(\neg p\Rightarrow q)$:
obie formuły są fałszywe dokładnie dla $p=q=0$.
Stąd możemy zbudować alternatywę, negację i,
przez prawo de Morgana, koniunkcję.
Sukcesywne zastępowanie spójników w dowolnej DNF
zachowuje wartość całej formuły, co wynika
z rekurencyjnej definicji wartościowania.
\end{proof}

\begin{twierdzenie}[Pojedynczy spójnik wystarcza]
\label{thm:lst-nand-nor}
Każda z jednoelementowych rodzin zawierających
spójnik NAND (kreska Sheffera) lub NOR (strzałka Peirce'a)
jest funkcyjnie zupełna.
\end{twierdzenie}
\begin{proof}
Piszemy $p\uparrow q=\neg(p\land q)$ oraz
$p\downarrow q=\neg(p\lor q)$. Dla NAND zachodzi
\[
 \neg p\equiv p\uparrow p,\qquad
 p\land q\equiv(p\uparrow q)\uparrow(p\uparrow q).
\]
Dla NOR zachodzi
\[
 \neg p\equiv p\downarrow p,\qquad
 p\lor q\equiv(p\downarrow q)\downarrow(p\downarrow q).
\]
Wzory wynikają bezpośrednio z idempotentności
koniunkcji i alternatywy oraz prawa podwójnej negacji;
można też sprawdzić cztery wiersze tablicy.
Twierdzenie~\ref{thm:lst-complete-connectives}
daje zupełność obu rodzin.
\end{proof}

\begin{stwierdzenie}[Rodziny niezupełne funkcyjnie]
\label{prop:lst-incomplete-connectives}
Rodzina $\{\land,\lor\}$ nie jest funkcyjnie zupełna,
nawet po dołączeniu obu stałych logicznych.
Sama implikacja bez stałej fałszu również nie jest
funkcyjnie zupełna.
\end{stwierdzenie}
\begin{proof}
Na $\{0,1\}^n$ przyjmujemy porządek współrzędny,
$0<1$. Rzuty i funkcje stałe są niemalejące,
a minimum oraz maksimum dwóch funkcji niemalejących
są niemalejące. Indukcja po budowie formuły
pokazuje, że każda formuła z koniunkcji,
alternatywy i stałych wyznacza funkcję niemalejącą.
Negacja nie jest niemalejąca, ponieważ jej wartości
w $0$ i $1$ to odpowiednio $1$ i $0$.
Dla samej implikacji indukcja pokazuje, że przy
wartościowaniu wszystkich liter przez $1$ każda
formuła ma wartość $1$: w kroku indukcyjnym
korzystamy z $1\Rightarrow1=1$. Nie można więc
zbudować negacji ani funkcji stale równej $0$.
\end{proof}

\begin{wniosek}[Ile istnieje funkcji prawdziwościowych?]
\label{cor:lst-boolean-function-count}
Dla ustalonej listy $n\geq1$ liter zdaniowych istnieje
dokładnie $2^{2^n}$ klas równoważności logicznej formuł
w tych literach.
\end{wniosek}
\begin{proof}
Tablica ma $2^n$ wierszy, w każdym można niezależnie
wybrać $0$ albo $1$. Jest zatem $2^{2^n}$ tablic.
Każdą realizuje formuła z
Twierdzenia~\ref{thm:lst-canonical-normal-forms}.
Dwie formuły są równoważne dokładnie wtedy, gdy
mają tę samą tablicę prawdy. Stąd podana liczba klas.
\end{proof}
Przykładowo dla dwóch liter istnieje $16$ funkcji
prawdziwościowych, chociaż samych formuł jest
nieskończenie wiele. Zupełność funkcyjna dotyczy
reprezentowania funkcji prawdziwościowych; nie jest
tym samym co pełność rachunku dowodów, która dotyczy
przejścia od konsekwencji semantycznej do dowodliwości.

\section{Przesłanki, wnioski i kontrprzykłady}\label{sec:logical-consequence}
\begin{definicja}[Poprawność wnioskowania zdaniowego]
\label{def:consequence}
Mówimy, że formuła $\psi$ \emph{wynika logicznie}
z przesłanek $\varphi_1,\ldots,\varphi_n$, gdy nie istnieje
\emph{wspólne} wartościowanie wszystkich liter występujących
w przesłankach lub we wniosku, przy którym wszystkie przesłanki
są prawdziwe, a $\psi$ jest fałszywa. Zapisujemy
$\varphi_1,\ldots,\varphi_n\models\psi$.
Przy braku przesłanek $\models\psi$ oznacza, że
$\psi$ jest tautologią.
\end{definicja}

\begin{twierdzenie}[Kryterium poprawności wnioskowania]
\label{thm:inference-test}
Niech $\varphi_1,\ldots,\varphi_n,\psi$ będą formułami
zdaniowymi, gdzie $n\ge1$. Następujące warunki są
równoważne:
\begin{enumerate}[label=(\roman*)]
\item $\varphi_1,\ldots,\varphi_n\models\psi$;
\item $((\varphi_1\land\cdots\land\varphi_n)\Rightarrow\psi)$
      jest tautologią.
\end{enumerate}
\end{twierdzenie}
\begin{proof}
Ustalmy dowolne wartościowanie. Implikacja z punktu
(ii) jest fałszywa dokładnie wtedy, gdy jej
poprzednik jest prawdziwy, a następnik fałszywy.
Poprzednik jest koniunkcją wszystkich przesłanek,
więc jest prawdziwy dokładnie wtedy, gdy wszystkie
przesłanki są prawdziwe. Zatem istnienie
wartościowania obalającego tautologiczność w (ii)
jest równoważne istnieniu wartościowania obalającego
poprawność w (i). Zaprzeczenie obu tych stwierdzeń
daje żądaną równoważność.
\end{proof}

\begin{twierdzenie}[Jednolite kryterium wynikania dla skończonej listy przesłanek]
\label{thm:empty-premises-consequence-lxiii}
Niech $n\in\{0,1,2,\ldots\}$, niech $\varphi_1,\ldots,\varphi_n$ będą formułami
rachunku zdań, a $\psi$ dowolną formułą. Dla $n=0$ przyjmijmy, że
koniunkcja pustej listy ma stałą wartość logiczną $1$ (oznaczamy ją
metajęzykowo przez $\top$). Wówczas
\[
\varphi_1,\ldots,\varphi_n\models\psi
\quad\Longleftrightarrow\quad
\models\bigl((\varphi_1\land\cdots\land\varphi_n)\Rightarrow\psi\bigr),
\]
gdzie przy $n=0$ prawą stronę rozumiemy jako $\models(\top\Rightarrow\psi)$.
Symbol $\top$ jest tutaj skrótem semantycznym, a nie dodatkową
literą w gramatyce z definicji~\ref{def:propositional-grammar-xli}.
\end{twierdzenie}
\begin{proof}
Dla $n\ge1$ jest to twierdzenie~\ref{thm:inference-test}.
Dla $n=0$ definicja~\ref{def:consequence} mówi, że
$\models\psi$ wtedy i tylko wtedy, gdy każde wartościowanie
nadaje $\psi$ wartość $1$. Ponieważ $\top$ ma zawsze wartość $1$,
implikacja $\top\Rightarrow\psi$ ma przy każdym wartościowaniu
dokładnie tę samą wartość co $\psi$. Stąd równoważność również
dla pustej listy. Nie trzeba wybierać żadnego wartościowania
spełniającego przesłanki: przesłanek po prostu nie ma.
\end{proof}
\begin{przyklad}[Pusta lista nie jest listą sprzeczną]
Dla $\psi=p\lor\neg p$ zarówno $\psi$, jak i
$\top\Rightarrow\psi$ mają wartość $1$ przy $p=0$
i przy $p=1$. Dlatego oba stwierdzenia w kryterium są
prawdziwe. Dla $\psi=p$ oba stwierdzenia o wszystkich
wartościowaniach są fałszywe: $p=0$ daje kontrprzykład.
Natomiast dla listy $(p,\neg p)$
wynikanie dowolnej formuły jest prawdziwe z innego powodu:
żadne wartościowanie nie spełnia wszystkich przesłanek.
\end{przyklad}

\begin{uwaga}[Przesłanki sprzeczne i brak przesłanek]
\label{rem:explosion-semantic-xli}
Nie należy mylić prawdziwości przesłanek
z poprawnością reguły wnioskowania. Dla dowolnych formuł
$p,q$ zachodzi $p,\neg p\models q$. Rzeczywiście,
żadne wartościowanie nie czyni obu przesłanek prawdziwymi:
jeżeli $p=1$, to $\neg p=0$, a jeżeli $p=0$, to pierwsza
przesłanka jest fałszywa. Nie istnieje więc wartościowanie
z prawdziwymi wszystkimi przesłankami i fałszywym wnioskiem.
Jest to semantyczna postać \emph{zasady eksplozji}; nie znaczy,
że dowolne zdanie $q$ jest prawdziwe. Podobnie
$\models p\lor\neg p$, lecz nie zachodzi
$\models p$ dla samej litery $p$, bo wartościowanie
$p=0$ jest kontrprzykładem.
\end{uwaga}

\begin{przyklad}[Niepoprawne wnioskowanie: potwierdzanie następnika]
Rozumowanie „$p\Rightarrow q$; $q$; zatem $p$” jest
niepoprawne. Wystarczy przyjąć $p=0$, $q=1$:
wtedy obie przesłanki są prawdziwe, ale wniosek
$p$ jest fałszywy. Konkretny przykład: „jeżeli
liczba jest podzielna przez $4$, to jest parzysta;
liczba $6$ jest parzysta; zatem jest podzielna przez $4$”.
\end{przyklad}

\begin{przyklad}[Niepoprawne wnioskowanie: zaprzeczanie poprzednikowi]
Z $p\Rightarrow q$ oraz $\neg p$ nie wynika $\neg q$.
Przyjmując $p=0$, $q=1$, otrzymujemy prawdziwe
przesłanki i fałszywy wniosek. Jest to ten sam
błąd co twierdzenie, że liczba niepodzielna przez $4$
nie może być parzysta.
\end{przyklad}

\begin{twierdzenie}[Trzy poprawne reguły wnioskowania]
\label{thm:valid-arguments}
Dla dowolnych formuł $p,q,r$ zachodzą
\begin{align*}
p\Rightarrow q,\quad\neg q&\models\neg p&&\text{(modus tollens)},\\
p\Rightarrow q,\quad q\Rightarrow r&\models p\Rightarrow r
   &&\text{(przechodniość implikacji)},\\
p\lor q,\quad\neg p&\models q&&\text{(sylogizm rozłączny)}.
\end{align*}
\end{twierdzenie}
\begin{proof}
W pierwszym przypadku przyjmijmy, że obie przesłanki
są prawdziwe. Ponieważ $\neg q$ jest prawdziwe,
$q$ jest fałszywe. Gdyby $p$ było prawdziwe,
$p\Rightarrow q$ byłoby fałszywe, sprzecznie z
założeniem. Zatem $p$ jest fałszywe, czyli
$\neg p$ jest prawdziwe.

W drugim przypadku załóżmy prawdziwość obu implikacji.
Chcemy uzasadnić $p\Rightarrow r$. Gdy $p$ jest
fałszywe, implikacja ta jest prawdziwa z definicji.
Gdy $p$ jest prawdziwe, pierwsza przesłanka wymusza
prawdziwość $q$, a druga prawdziwość $r$.
Wtedy również $p\Rightarrow r$ jest prawdziwa.

W trzecim przypadku $p\lor q$ jest prawdziwe, więc
co najmniej jedna z wartości $p,q$ wynosi $1$.
Jednocześnie $\neg p$ jest prawdziwe, a więc $p=0$.
Pozostaje $q=1$, co jest dokładnie tezą.
\end{proof}

\begin{przyklad}[Jak odczytywać wzory, aby nie mylić spójnika z wnioskowaniem]
Zapis $(p\Rightarrow q)$ jest \emph{formułą}: można obliczyć
jego wartość z tablicy prawdy. Natomiast zapis
$p\models q$ \emph{nie jest formułą rachunku zdań}, lecz
metamatematycznym stwierdzeniem, że nie istnieje wartościowanie
z $p=1$ i $q=0$. W przypadku jednej przesłanki zachodzi
$p\models q$ wtedy i tylko wtedy, gdy $p\Rightarrow q$ jest
tautologią. Podobnie $p\equiv q$ opisuje zgodność wartości
formuł przy wszystkich wartościowaniach, a
$p\Leftrightarrow q$ jest formułą, której wartość badamy.
\end{przyklad}

\section{Zbiory przesłanek i konsekwencja logiczna}
\label{sec:lst-propositional-consequence}
Dotychczas rozważaliśmy skończone listy przesłanek.
Możemy jednak ustalić także zbiór nieskończony, na przykład
nieskończoną listę warunków, które mają zachodzić jednocześnie.
Nadal każda pojedyncza formuła ma skończoną budowę.

\begin{definicja}[Spełnialność zbioru formuł i wynikanie]
\label{def:lst-propositional-theory}
Niech $\Gamma$ będzie dowolnym zbiorem formuł zdaniowych.
Wartościowanie $v$ spełnia $\Gamma$, gdy $v(\gamma)=1$
dla każdego $\gamma\in\Gamma$. Zbiór jest spełnialny,
gdy istnieje takie wartościowanie.
Piszemy $\Gamma\models\varphi$, gdy każde wartościowanie
spełniające $\Gamma$ nadaje formule $\varphi$ wartość $1$.
\end{definicja}
Definicja rozszerza Definicję~\ref{def:consequence}.
Zbiór pusty jest spełnialny: spełnia go każde wartościowanie.
Nie wystarcza natomiast, by każda formuła z $\Gamma$
była z osobna spełnialna. Na przykład $p$ oraz $\neg p$
są z osobna spełnialne, ale nie mają wspólnego wartościowania
spełniającego obie formuły.

\begin{twierdzenie}[Kryterium przez niespełnialność]
\label{thm:lst-propositional-refutation}
Dla dowolnego zbioru przesłanek $\Gamma$ i formuł
$\varphi,\psi$ zachodzą równoważności
\[
\begin{split}
 \Gamma\models\varphi
 &\iff \Gamma\cup\{\neg\varphi\}
          \text{ jest niespełnialny},\\
 \Gamma\cup\{\varphi\}\models\psi
 &\iff\Gamma\models(\varphi\Rightarrow\psi).
\end{split}
\]
\end{twierdzenie}
\begin{proof}
W pierwszym przypadku zaprzeczenie lewej strony oznacza
istnienie wartościowania spełniającego $\Gamma$ i nadającego
$\varphi$ wartość $0$, czyli nadającego $\neg\varphi$ wartość $1$.
Jest to dokładnie spełnialność zbioru po prawej stronie.
W drugim przypadku, przy wartościowaniu spełniającym $\Gamma$,
implikacja $\varphi\Rightarrow\psi$ jest fałszywa dokładnie
wtedy, gdy $\varphi$ jest prawdziwa, a $\psi$ fałszywa.
Takie wartościowanie obala obie strony równoważności.
\end{proof}
Drugą równoważność nazywamy semantyczną postacią twierdzenia
o dedukcji. Nie zakładamy tu żadnego formalnego rachunku
dowodów. Syntaktyczne twierdzenie o dedukcji dla dedukcji
naturalnej omawiamy w teorii modeli.

\begin{stwierdzenie}[Własności konsekwencji]
\label{prop:lst-propositional-consequence-properties}
Dla zbiorów formuł $\Gamma,\Delta$ i formuły $\varphi$
zachodzą następujące własności:
\begin{enumerate}
 \item jeśli $\varphi\in\Gamma$, to $\Gamma\models\varphi$;
 \item jeśli $\Gamma\subseteq\Delta$ i $\Gamma\models\varphi$,
       to $\Delta\models\varphi$;
 \item jeśli $\Gamma\models\delta$ dla każdego $\delta\in\Delta$
       i $\Delta\models\varphi$, to $\Gamma\models\varphi$.
\end{enumerate}
\end{stwierdzenie}
\begin{proof}
Każde wartościowanie spełniające $\Gamma$ spełnia w szczególności
każdą należącą do niego formułę, co daje pierwszy punkt.
Jeśli $\Gamma\subseteq\Delta$, to każde wartościowanie spełniające
$\Delta$ spełnia także $\Gamma$. Daje to drugi punkt.
W trzecim, dowolne wartościowanie spełniające $\Gamma$
spełnia wszystkie formuły z $\Delta$, a więc także $\varphi$.
Argumenty pozostają ważne dla zbiorów niespełnialnych:
wówczas nie istnieje wartościowanie obalające konsekwencję.
\end{proof}
Druga własność nazywa się monotonicznością: dodanie przesłanek
nie unieważnia żadnego dotychczasowego wniosku klasycznego.
Trzecia pozwala zastąpić cały łańcuch pośrednich wnioskowań
jednym wynikaniem od przesłanek początkowych do końcowego wniosku.

\begin{twierdzenie}[Jednoczesne podstawianie formuł]
\label{thm:lst-propositional-substitution}
Niech $\sigma$ przyporządkowuje każdej literze zdaniowej $p$
formułę $\sigma(p)$. Rozszerzamy $\sigma$ na wszystkie formuły,
zachowując ich spójniki i podstawiając jednocześnie
przypisane formuły w miejsca liter. Jeżeli
$\Gamma\models\varphi$, to $\sigma(\Gamma)\models\sigma(\varphi)$,
gdzie $\sigma(\Gamma)=\{\sigma(\gamma):\gamma\in\Gamma\}$.
\end{twierdzenie}
\begin{proof}
Ustalmy dowolne wartościowanie $v$ i określmy wartościowanie
liter $w(p)=v(\sigma(p))$. Indukcja po budowie formuły $\theta$
pokazuje, że $v(\sigma(\theta))=w(\theta)$.
Dla litery wynika to z definicji $w$, a dla spójnika
obie strony powstają przez zastosowanie tej samej funkcji
prawdziwościowej do równych wartości argumentów.
Jeżeli $v$ spełnia $\sigma(\Gamma)$, to $w$ spełnia $\Gamma$.
Z założenia $w(\varphi)=1$, zatem $v(\sigma(\varphi))=1$.
\end{proof}
W szczególności podstawienie dowolnych formuł w miejsce liter
tautologii daje tautologię. Na przykład do tautologii
$p\lor\neg p$ możemy podstawić formułę $q\land r$ za $p$
i otrzymać $(q\land r)\lor\neg(q\land r)$.
Nie podstawiamy następnie rekurencyjnie za litery pojawiające się
we wstawionej formule: podstawienie jest jednoczesne.

\section{Skończoność przesłanek}
\label{sec:lst-propositional-compactness}
Wynikanie z nieskończonej liczby warunków może wyglądać na
istotnie silniejsze od wynikania ze skończonej listy.
W klasycznym rachunku zdań zachodzi jednak następujące
twierdzenie. Jego dowód podamy dla ustalonego wcześniej
przeliczalnego języka liter zdaniowych.

\begin{twierdzenie}[Zwartość rachunku zdań]
\label{thm:lst-propositional-compactness}
Zbiór formuł $\Gamma$ w języku o przeliczalnym zbiorze liter
zdaniowych jest spełnialny wtedy i tylko wtedy,
gdy każdy jego skończony podzbiór jest spełnialny.
\end{twierdzenie}
\begin{proof}
Jeżeli $v$ spełnia $\Gamma$, to spełnia każdy jego podzbiór.
Dowodzimy implikacji odwrotnej.
W przypadku skończonego zbioru liter dołączamy w razie potrzeby
przeliczalnie wiele nowych liter niewystępujących w $\Gamma$,
co nie zmienia spełnialności.
Oznaczmy litery przez $p_1,p_2,\ldots$.
Dla skończonego przypisania $a_1,\ldots,a_n\in\{0,1\}$
rozpatrujemy własność: każdy skończony podzbiór $\Gamma$
ma wartościowanie spełniające go i przyjmujące
wartości $a_i$ na literach $p_i$ dla $i\leq n$.

Przypisanie puste ma tę własność z założenia.
Jeśli ma ją przypisanie długości $n$, to co najmniej jedno
z jego dwóch rozszerzeń przez $a_{n+1}=0$ albo $a_{n+1}=1$
także ma tę własność.
Gdyby nie miało jej żadne, istniałyby dwa skończone podzbiory
$\Gamma_0,\Gamma_1\subseteq\Gamma$: pierwszy nie miałby
wartościowania zgodnego z rozszerzeniem przez $0$,
a drugi zgodnego z rozszerzeniem przez $1$.
Skończony zbiór $\Gamma_0\cup\Gamma_1$ nie miałby wtedy
wartościowania zgodnego z wyjściowym przypisaniem,
ponieważ dowolne wartościowanie musi nadać $p_{n+1}$
jedną z tych dwóch wartości. Otrzymujemy sprzeczność.

Budujemy więc ciąg $a_1,a_2,\ldots$, na każdym kroku
wybierając rozszerzenie przez $0$, jeśli zachowuje własność,
a w przeciwnym razie przez $1$. Definiujemy $v(p_i)=a_i$.
Dla dowolnej $\gamma\in\Gamma$ istnieje $n$ takie, że
wszystkie jej litery należą do $p_1,\ldots,p_n$,
gdyż formuła jest skończona.
Własność z etapu $n$ zastosowana do zbioru $\{\gamma\}$
daje wartościowanie $w$ spełniające $\gamma$ i zgodne
z $v$ na tych literach. Z rekurencyjnej definicji
wartościowania wynika $v(\gamma)=w(\gamma)=1$.
Tak jest dla każdej $\gamma$, zatem $v$ spełnia $\Gamma$.
\end{proof}

\begin{wniosek}[Skończony zbiór istotnych przesłanek]
\label{cor:lst-propositional-finite-support}
W przeliczalnym języku zdaniowym $\Gamma\models\varphi$
wtedy i tylko wtedy, gdy istnieje skończony podzbiór
$\Gamma_0\subseteq\Gamma$ taki, że $\Gamma_0\models\varphi$.
\end{wniosek}
\begin{proof}
Jeżeli $\Gamma\models\varphi$, to
$\Gamma\cup\{\neg\varphi\}$ jest niespełnialny na mocy
Twierdzenia~\ref{thm:lst-propositional-refutation}.
Ze zwartości wynika, że pewien jego skończony podzbiór
$\Sigma$ jest niespełnialny. Niech $\Gamma_0=\Sigma\cap\Gamma$.
Mamy $\Sigma\subseteq\Gamma_0\cup\{\neg\varphi\}$,
więc ten ostatni zbiór jest również niespełnialny.
Ponownie stosujemy kryterium przez niespełnialność.
Kierunek odwrotny wynika z monotoniczności
(Stwierdzenie~\ref{prop:lst-propositional-consequence-properties}).
\end{proof}

\begin{przyklad}[Nieskończony łańcuch implikacji]
\label{ex:lst-propositional-chain}
Rozważmy $\Gamma=\{p_1\}\cup\{p_i\Rightarrow p_{i+1}:i\geq1\}$.
Dla każdego ustalonego $n$ mamy $\Gamma\models p_n$:
od $p_1$ przechodzimy kolejno do $p_2,\ldots,p_n$.
Wystarcza skończony podzbiór złożony z $p_1$ i pierwszych
$n-1$ implikacji. Nie istnieje jednak jeden skończony
podzbiór $\Gamma_0$ wystarczający do wyprowadzenia wszystkich $p_n$.
Wybierzmy $N\geq1$ tak, by wszystkie litery występujące
w $\Gamma_0$ miały indeks co najwyżej $N$.
Wartościowanie nadające $p_1,\ldots,p_N$ wartość $1$
i pozostałym literom wartość $0$ spełnia $\Gamma_0$,
lecz nie spełnia $p_{N+1}$.
\end{przyklad}
Zwartość nie daje algorytmu rozstrzygającego spełnialność
dowolnie zadanego nieskończonego zbioru formuł.
Mówi o istnieniu skończonej przeszkody w razie niespełnialności,
nie o sposobie jej znalezienia przy dowolnej reprezentacji danych.
Zwartość logiki pierwszego rzędu ma osobny dowód
w teorii modeli.

\section{Czego jeszcze nie potrafi rachunek zdań?}\label{sec:why-quantifiers}
Niech $P(x)$ oznacza wypowiedź „liczba $x$ jest parzysta”.
Możemy sprawdzić $P(2)$ albo $P(3)$, ale sama logika zdań
nie ma jeszcze symbolu pozwalającego zbudować wypowiedź
„każda liczba ma pewną własność” lub „istnieje liczba
o tej własności”. Nie odróżnia też jednoznacznie wypowiedzi
„każdy ma swojego przyjaciela” od „istnieje ktoś,
kto jest przyjacielem wszystkich”. Do tego potrzebujemy
\emph{kwantyfikatorów}; wprowadzimy je w następnym rozdziale.

\begin{uwaga}[Pierwsze trzy rozróżnienia, które należy zapamiętać]
\emph{Zdanie prawdziwe} nie musi być \emph{tautologią}.
\emph{Implikacja} na ogół nie jest logicznie równoważna
swojej implikacji odwrotnej.
\emph{Poprawne wnioskowanie} nie zakłada samo przez się,
że jego przesłanki są faktycznie prawdziwe.
\end{uwaga}

\chapter{Kwantyfikatory i podstawowe metody dowodzenia}\label{ch:quantifiers}

\section{Dlaczego same spójniki nie wystarczają?}\label{sec:quantifier-motivation}
Zdania „$2$ jest parzyste” i „$7$ jest parzyste” mają
określone wartości logiczne. Matematyka posługuje się jednak
przede wszystkim wypowiedziami dotyczącymi całych rodzin
obiektów: „każda liczba parzysta ma pewną własność” albo
„istnieje liczba spełniająca równanie”. Samymi spójnikami
$\neg,\land,\lor,\Rightarrow,\Leftrightarrow$ nie umiemy
wygodnie zapisać, do \emph{ilu} obiektów odnosi się taki warunek.

Potrzebne są dwa nowe symbole: $\forall$ („dla każdego”)
oraz $\exists$ („istnieje”). W przeciwieństwie do zwykłych
spójników związane są one ze zmienną i ze wskazanym
\emph{zakresem jej wartości}.

\section{Dziedzina zmiennej i predykat}\label{sec:predicates}
\begin{definicja}[Dziedzina i predykat --- ujęcie wstępne]
\label{def:predicate}
Ustalamy pewną \emph{dziedzinę rozważań} $D$, czyli kolekcję
obiektów, o których chcemy mówić (może być także
pusta). \emph{Predykat} $P(x)$
to warunek o zmiennej $x$, który po wskazaniu obiektu
$x$ z $D$ przyjmuje wartość prawdy albo fałszu.
Samo $P(x)$ jest na ogół \emph{wypowiedzią otwartą},
a nie zdaniem o ustalonej wartości logicznej.
W tym rozdziale pojęcie dziedziny traktujemy intuicyjnie;
jego teoriomnogościową definicję podamy później.
\end{definicja}

\begin{uwaga}[Parametry i predykaty wieloargumentowe]
\label{rem:predicates-arity-xli}
Można również rozważać predykat $R(x,y)$ dwóch zmiennych
o dziedzinach odpowiednio $D$ i $E$: po wskazaniu
$a\in D$ i $b\in E$ otrzymujemy zdanie $R(a,b)$.
Przykładowo dla $D=E=\N$ predykat $R(x,y)$ może
oznaczać $x<y$. Napis $R(x,3)$ ma nadal wolną
zmienną $x$, natomiast $R(2,3)$ jest prawdziwym
zdaniem. Aby rozstrzygać wartość logiczną wyrażenia
z kwantyfikatorem i dodatkowymi zmiennymi wolnymi,
trzeba wcześniej ustalić ich wartości.
Ścisłą semantykę przypisań do zmiennych podamy w
rozdziale~\ref{ch:semantics-fol}.
\end{uwaga}

\begin{przyklad}[Ta sama wypowiedź, różna dziedzina]
Rozważmy $P(x)$: „$x^2=1$”. Jeśli dopuszczamy tylko
nieujemne liczby całkowite, jedynym rozwiązaniem jest $1$.
Jeśli dopuszczamy wszystkie liczby całkowite, rozwiązaniami
są $1$ oraz $-1$. Jeśli zaś dopuszczamy wszystkie liczby
rzeczywiste, nadal otrzymujemy te dwa rozwiązania.
Bez wskazania dziedziny wypowiedź „istnieje dokładnie jedno
rozwiązanie” nie jest dostatecznie określona.
\end{przyklad}

\begin{uwaga}[Podstawowa notacja liczbowa]
$\N=\{0,1,2,\ldots\}$ będzie oznaczać liczby naturalne,
$\Z$ liczby całkowite, a $\R$ liczby rzeczywiste.
Zapisy te wykorzystujemy na razie jako znane nazwy
kolekcji liczb. Dopiero w części teoriomnogościowej
uzasadnimy formalnie istnienie i konstrukcję takich zbiorów.
\end{uwaga}

\section{Kwantyfikator ogólny}\label{sec:forall}
\begin{definicja}[„Dla każdego”]\label{def:forall}
Zdanie $\forall x\in D\;P(x)$ jest prawdziwe wtedy i tylko
wtedy, gdy $P(a)$ jest prawdziwe dla \emph{każdego} obiektu
$a$ z dziedziny $D$. Jest fałszywe dokładnie wtedy, gdy
\emph{istnieje kontrprzykład}: pewien obiekt $a$ z $D$,
dla którego $P(a)$ jest fałszywe.
\end{definicja}

\begin{przyklad}[Wystarczy jeden kontrprzykład]
Zdanie $\forall n\in\N\;(n<10)$ jest fałszywe:
liczba $10$ jest naturalna, lecz nierówność $10<10$
jest fałszywa. Nie trzeba w tym celu analizować
pozostałych liczb naturalnych.
\end{przyklad}

\begin{przyklad}[Dlaczego trzeba ustalić dziedzinę?]
Wypowiedź „każde $x$ spełnia $x^2\ge0$” jest prawdziwa
dla liczb rzeczywistych. Przy przejściu do liczb zespolonych
nie staje się po prostu fałszywa: bez dodatkowego
określenia znaczenia symbolu $\ge$ nie ma ustalonej
wartości logicznej, ponieważ na liczbach zespolonych
nie przyjmujemy standardowo takiego porządku.
Nie wolno zmieniać dziedziny bez sprawdzenia
sensu użytych symboli.
\end{przyklad}

\section{Kwantyfikator egzystencjalny}\label{sec:exists}
\begin{definicja}[„Istnieje”]\label{def:exists}
Zdanie $\exists x\in D\;P(x)$ jest prawdziwe dokładnie
wtedy, gdy \emph{istnieje co najmniej jeden}
obiekt $a$ z dziedziny $D$, dla którego $P(a)$ jest
prawdziwe. Taki obiekt $a$ nazywamy \emph{świadkiem istnienia}.
Uwaga: prawdziwość zdania egzystencjalnego nie wymaga,
abyśmy potrafili obliczyć lub jawnie nazwać jego świadka.
W dowodach konstrukcyjnych go podajemy, ale w logice
klasycznej możliwe są także dowody niekonstrukcyjne.
Aby obalić zdanie egzystencjalne, trzeba wykazać,
że \emph{żaden} obiekt z $D$ nie spełnia warunku.
\end{definicja}

\begin{przyklad}[Wskazanie świadka]
Zdanie $\exists n\in\N\;(n^2=25)$ jest prawdziwe,
gdyż $n=5$ jest liczbą naturalną i $5^2=25$.
Zdanie $\exists n\in\N\;(n^2=-1)$ jest fałszywe,
ponieważ dla każdej liczby naturalnej $n$ zachodzi
$n^2\ge0$. W drugim przypadku pojedyncze
przykłady $n=0,1,2$ nie wystarczą do dowodu;
argument musi objąć wszystkie liczby naturalne.
\end{przyklad}

\begin{uwaga}[Dowód istnienia nie zawsze daje algorytm]
Nie należy utożsamiać warunku $\exists x\in D\;P(x)$
z istnieniem procedury, która dla danego $P$ wyszuka
świadka. Istnienie można czasem uzasadnić pośrednio,
wykazując sprzeczność założenia $\forall x\in D\;\neg P(x)$.
Jest to dowód klasyczny i nie musi wyznaczać konkretnego
$x$. W rozdziałach o obliczalności przekonamy się, dlaczego
takie rozróżnienie jest fundamentalne.
\end{uwaga}

\begin{przyklad}[Prawdziwość nie oznacza jednoznaczności]
Zdanie $\exists x\in\Z\;(x^2=1)$ jest prawdziwe,
ale nie mówi, że świadek jest jeden: spełniają je
$x=1$ i $x=-1$. Osobno wprowadzimy zapis
„istnieje dokładnie jeden”.
\end{przyklad}

\begin{uwaga}[Pusta dziedzina]
Jeżeli $D$ jest puste, to $\forall x\in D\;P(x)$ jest
prawdziwe: nie istnieje w $D$ obiekt obalający $P$.
Natomiast $\exists x\in D\;P(x)$ jest fałszywe:
nie można wskazać świadka w pustej dziedzinie.
Nie zakładamy przez to istnienia jakiegoś „pustego obiektu”.
To szczególne przypadki definicji obu kwantyfikatorów.
\end{uwaga}

\begin{uwaga}[Niepusta dziedzina struktur pierwszego rzędu]
\label{rem:empty-domain-semantics-xli}
W tym rozdziale dopuszczamy także pustą dziedzinę $D$
przy \emph{kwantyfikatorze ograniczonym}.
Nie należy z tego wyciągać wniosku, że dowolna
struktura klasycznej logiki pierwszego rzędu może
mieć puste uniwersum: w standardowo przyjętej
semantyce tej logiki uniwersum struktury jest niepuste.
Nie przeszkadza to interpretować $\forall x\in D$
i $\exists x\in D$ dla pustego \emph{podzbioru}
$D$ niepustego uniwersum, korzystając z rozwinięć
podanych w sekcji~\ref{sec:quantifier-scope}.
\end{uwaga}

\section{Zapisy skrócone i zakres kwantyfikatora}\label{sec:quantifier-scope}
\begin{definicja}[Zakres działania kwantyfikatora]
W zapisie $\forall x\in D\;(P(x)\Rightarrow Q(x))$
\emph{zakresem} kwantyfikatora jest cała formuła
w nawiasie. Wystąpienia zmiennej $x$ w tym zakresie są związane
przez dany kwantyfikator, o ile nie zostały związane
przez wewnętrzny kwantyfikator tej samej zmiennej.
Wystąpienie zmiennej, którego nie wiąże żaden
kwantyfikator, nazywamy \emph{wolnym}.
\end{definicja}

\begin{przyklad}[Zmienna wolna i związana]
Niech $y$ przebiega $\R$. W formule
$(\forall x\in\N\;x^2\ge0)\land(y>0)$
zmienna $x$ jest związana, natomiast $y$ pozostaje wolna:
wartość całego wyrażenia zależy jeszcze od wyboru $y$.
Jeśli dopiszemy $\exists y\in\R$ przed całą formułą,
uzyskamy zdanie bez zmiennych wolnych.
\end{przyklad}

\begin{uwaga}[Niebezpieczne nawiasy]
Formuły $\forall x\;(P(x)\Rightarrow Q(x))$
i $(\forall x\;P(x))\Rightarrow Q(x)$ nie wyrażają
tego samego: w drugiej formule wystąpienie $x$ w
$Q(x)$ jest wolne. Przy pierwszych obliczeniach
zawsze wpisujemy zakres kwantyfikatora w nawiasach.
\end{uwaga}

\begin{twierdzenie}[Zmiana nazwy zmiennej związanej]
\label{thm:alpha-renaming-xli}
Niech $D$ będzie ustaloną dziedziną, a $P(x,z)$
predykatem, w którym $z$ oznacza ewentualne dodatkowe
parametry. Niech $y$ będzie nową nazwą zmiennej,
która nie występuje w $P(x,z)$. Wtedy, po ustaleniu
wartości parametrów $z$, zdania
$\forall x\in D\;P(x,z)$ i
$\forall y\in D\;P(y,z)$ są równoważne; podobnie
$\exists x\in D\;P(x,z)$ i
$\exists y\in D\;P(y,z)$ są równoważne.
W drugich formułach $P(y,z)$ oznacza wynik zastąpienia
\emph{wolnych} wystąpień $x$ przez $y$.
\end{twierdzenie}
\begin{proof}
Ustalmy wartości wszystkich parametrów $z$.
Dla każdego obiektu $a\in D$ zdania otrzymane
przez podstawienie $x=a$ w $P(x,z)$ i przez
podstawienie $y=a$ w $P(y,z)$ mają tę samą treść:
zmieniliśmy nazwę miejsca podstawienia, a nie
sam warunek na $a$. Zdanie ogólne po obu stronach
wymaga zatem prawdziwości tej samej rodziny zdań
$P(a,z)$ dla wszystkich $a\in D$.
Zdanie egzystencjalne po obu stronach wymaga
istnienia świadka w tej samej rodzinie.
W obu przypadkach warunki prawdziwości są identyczne.
\end{proof}

\begin{przyklad}[Niebezpieczeństwo przechwycenia zmiennej]
\label{ex:variable-capture-xli}
Przemianowanie $x$ na nazwę występującą pod innym
kwantyfikatorem wymaga sprawdzenia, czy nie dojdzie
do przechwycenia zmiennej. Na przykład
$\forall x\in\N\;\exists y\in\N\;(x<y)$ jest prawdziwe:
dla każdego $x$ wybieramy $y=x+1$.
Błędne przemianowanie zewnętrznej zmiennej $x$ na $y$
(łącznie z jej wystąpieniem w nierówności) prowadzi do
$\forall y\in\N\;\exists y\in\N\;(y<y)$;
wewnętrzny kwantyfikator wiąże teraz oba wystąpienia
$y$ w nierówności, która jest zawsze fałszywa.
Poprawna zmiana obu nazw z zachowaniem ich odrębności
daje np.
$\forall u\in\N\;\exists v\in\N\;(u<v)$.
\end{przyklad}

\begin{przyklad}[Wszystkie elementy spełniające warunek]
Wypowiedź „każda liczba naturalna podzielna przez $4$
jest parzysta” zapisujemy
\[
\forall n\in\N\;\bigl(4\mid n\Rightarrow2\mid n\bigr),
\]
gdzie $a\mid b$ czytamy „$a$ dzieli $b$”.
Zapis $4\mid n$ znaczy, że $n=4k$ dla pewnej
liczby całkowitej $k$. Gdy $4$ nie dzieli $n$,
nie naruszamy prawdziwości implikacji.
\end{przyklad}

\begin{uwaga}[Skrót formalny]
Załóżmy, że $D$ jest podzbiorem ustalonego niepustego
uniwersum $U$, a warunek $P(x)$ ma określoną wartość
logiczną dla każdego $x\in U$. Wtedy zapis
$\forall x\in D\;P(x)$ jest skrótem formuły
$\forall x\;(x\in D\Rightarrow P(x))$;
podobnie $\exists x\in D\;P(x)$ oznacza
$\exists x\;(x\in D\land P(x))$.
Bez takiego wspólnego uniwersum i interpretacji $P$
poza $D$ nie musimy mieć prawa używać tych rozwinięć:
poprzednie definicje kwantyfikatorów ograniczonych
wymagały wartości $P(x)$ jedynie dla $x\in D$.
Symbol $\in$ będzie zdefiniowany dokładnie w następnym
rozdziale; na razie oznacza po prostu „jest elementem”.
\end{uwaga}

\begin{uwaga}[Kiedy rozwinięcia skrótów są prawidłowe?]
W zapisie $\forall x\,(x\in D\Rightarrow P(x))$ zmienna
$x$ przebiega całe ustalone uniwersum, a $x\in D$ jest
dodatkowym warunkiem. Jeśli $D$ jest puste, poprzednik
implikacji jest zawsze fałszywy, zatem kwantyfikator
ogólny jest prawdziwy. W zapisie
$\exists x\,(x\in D\land P(x))$ koniunkcja nie może
być prawdziwa dla żadnego $x$.
Właśnie dlatego w ograniczeniu kwantyfikatora ogólnego
występuje \emph{implikacja}, a w ograniczeniu kwantyfikatora
egzystencjalnego --- \emph{koniunkcja}; ich zamiana
zmieniłaby sens sformułowania.
\end{uwaga}

\begin{przyklad}[Jak odróżnić „nie wszyscy” od „nikt”?]
Weźmy wypowiedź „wszyscy uczniowie zdali egzamin”.
Jej negacją jest „istnieje uczeń, który nie zdał”, a nie
„żaden uczeń nie zdał”. Pierwsze zdanie kwestionuje choć jeden
przypadek; drugie kwestionuje wszystkie. Analogicznie negacja
„istnieje uczeń, który zdał” brzmi „żaden uczeń nie zdał”.
To przygotowuje nas do ścisłych wzorów z kwantyfikatorami.
\end{przyklad}

\section{Negowanie kwantyfikatorów}\label{sec:quantifier-negation}
\begin{twierdzenie}[Prawa de Morgana dla kwantyfikatorów]
\label{thm:negation-quantifiers}
Niech $D$ będzie dowolną dziedziną (także pustą), a
$P(x)$ predykatem o zmiennej $x$ z dziedziny $D$.
W klasycznej logice zachodzą równoważności
\[
\neg\bigl(\forall x\in D\;P(x)\bigr)
\ \Leftrightarrow\ \exists x\in D\;\neg P(x),
\]
\[
\neg\bigl(\exists x\in D\;P(x)\bigr)
\ \Leftrightarrow\ \forall x\in D\;\neg P(x).
\]
\end{twierdzenie}
\begin{proof}
Pierwsza równoważność: zdanie $\forall x\in D\;P(x)$
jest fałszywe dokładnie wtedy, gdy istnieje w $D$
co najmniej jeden kontrprzykład dla $P$, czyli
obiekt $a$ taki, że $P(a)$ jest fałszywe.
Na mocy definicji negacji jest to dokładnie warunek,
że $\neg P(a)$ jest prawdziwe, a więc że
$\exists x\in D\;\neg P(x)$ zachodzi.

Druga równoważność: $\exists x\in D\;P(x)$
jest fałszywe dokładnie wtedy, gdy nie ma w $D$
żadnego świadka dla $P$. To znaczy, że dla
każdego obiektu $a$ z $D$ formuła $P(a)$
jest fałszywa, czyli $\neg P(a)$ jest prawdziwa.
Definicja kwantyfikatora ogólnego daje
$\forall x\in D\;\neg P(x)$.
W obu dowodach przechodziliśmy przez równoważne
warunki, więc uzyskaliśmy obie strony implikacji.
\end{proof}

\begin{przyklad}[Negowanie zdania z codzienną treścią]
„Każdy student zdał egzamin” negujemy jako
„Istnieje student, który nie zdał egzaminu”.
Nie jako „Żaden student nie zdał egzaminu”:
ta druga wypowiedź jest znacznie silniejsza.
Podobnie zaprzeczenie „Istnieje rozwiązanie równania”
brzmi „Żaden dopuszczalny obiekt nie jest rozwiązaniem”,
a nie „Istnieje obiekt, który nie jest rozwiązaniem”.
\end{przyklad}

\begin{przyklad}[Negacja implikacji pod kwantyfikatorem]
Rozważmy zdanie
$\forall n\in\N\;(n>2\Rightarrow n^2>4)$.
Jego zaprzeczenie ma postać
$\exists n\in\N\;(n>2\land n^2\le4)$.
Dlaczego koniunkcja? Ponieważ
$\neg(p\Rightarrow q)\equiv p\land\neg q$
(twierdzenie~\ref{thm:negated-implication-xli}):
aby obalić implikację, musimy mieć zarazem
prawdziwy poprzednik i fałszywy następnik.
\end{przyklad}

\section{Kolejność kwantyfikatorów ma znaczenie}
\label{sec:quantifier-order}
\begin{przyklad}[„Każdemu można dobrać” a „istnieje jeden dla wszystkich”]
Rozpatrzmy na liczbach naturalnych warunek $x<y$.
Zdanie
\[
\forall x\in\N\;\exists y\in\N\;(x<y)
\]
jest prawdziwe: dla dowolnie wskazanego $x$
możemy wybrać $y=x+1$. Wybrany $y$ \emph{zależy od $x$}.
Natomiast zdanie
\[
\exists y\in\N\;\forall x\in\N\;(x<y)
\]
jest fałszywe. Gdyby istniało takie $y$,
można byłoby wstawić $x=y$, otrzymując $y<y$,
co jest fałszem. W szczególności nie istnieje
jedna liczba naturalna większa od wszystkich
liczb naturalnych.
\end{przyklad}

\begin{twierdzenie}[Podwójne kwantyfikatory tego samego typu]
\label{thm:same-quantifier-order}
Niech $D,E$ będą dziedzinami, a $P(x,y)$ predykatem
dla $x\in D$ i $y\in E$. Wtedy
\[
\forall x\in D\;\forall y\in E\;P(x,y)
\ \Leftrightarrow\
\forall y\in E\;\forall x\in D\;P(x,y),
\]
\[
\exists x\in D\;\exists y\in E\;P(x,y)
\ \Leftrightarrow\
\exists y\in E\;\exists x\in D\;P(x,y).
\]
\end{twierdzenie}
\begin{proof}
Pierwsze zdanie po lewej jest prawdziwe dokładnie
wtedy, gdy dla \emph{każdej} pary dopuszczalnych
obiektów $a\in D$, $b\in E$ warunek $P(a,b)$
jest prawdziwy. Kolejność, w jakiej wymienimy te
dwa obiekty, nie zmienia wartości $P(a,b)$.
To samo kryterium charakteryzuje zdanie po prawej.

Druga para zdań ma po obu stronach dokładnie
ten sam warunek: można wskazać \emph{co najmniej
jedną} parę $a\in D$, $b\in E$, dla której
$P(a,b)$ jest prawdziwe. Wskazanie najpierw $a$,
a potem $b$ albo odwrotnie nie zmienia
istnienia takiej pary.
\end{proof}

\begin{twierdzenie}[Jedna poprawna zamiana kwantyfikatorów różnych typów]
\label{thm:swap-mixed}
Niech $D,E$ będą dziedzinami, a $P(x,y)$ predykatem
określonym dla $x\in D$ i $y\in E$. Wtedy
\[
\Bigl(\exists y\in E\;\forall x\in D\;P(x,y)\Bigr)
\Rightarrow
\Bigl(\forall x\in D\;\exists y\in E\;P(x,y)\Bigr).
\]
Implikacja odwrotna nie zachodzi ogólnie.
\end{twierdzenie}
\begin{proof}
Załóżmy, że po lewej istnieje obiekt $b\in E$,
dla którego $P(a,b)$ zachodzi dla każdego $a\in D$.
Weźmy dowolny $a\in D$. Ten sam, już wybrany
obiekt $b\in E$ spełnia $P(a,b)$, a więc
istnieje odpowiedni $y\in E$. Ponieważ wybór
$a\in D$ był dowolny, otrzymujemy zdanie po prawej.
Jeżeli dziedzina $D$ jest pusta, wniosek po prawej jest
prawdziwy niezależnie od poprzednika, więc
implikacja również pozostaje prawdziwa.

Aby pokazać, że odwrotność jest fałszywa,
weźmy $D=E=\N$ i $P(x,y)$ określony przez $x<y$.
Jak wykazaliśmy w poprzednim przykładzie,
$\forall x\in\N\;\exists y\in\N\;(x<y)$
jest prawdziwe, a
$\exists y\in\N\;\forall x\in\N\;(x<y)$
jest fałszywe. Jest to kontrprzykład do odwrotności.
\end{proof}

\begin{uwaga}[Puste dziedziny w zamianie kwantyfikatorów]
\label{rem:mixed-quantifiers-empty-domains-lxiii}
W twierdzeniu~\ref{thm:swap-mixed} poprawność implikacji
nie wymaga niepustości $D$ ani $E$, o ile dopuszczamy
semantykę kwantyfikatorów po pustych zbiorach.
Jeżeli $D=\varnothing$, to zdanie
$\forall x\in D\;\exists y\in E\;P(x,y)$ jest prawdziwe,
niezależnie od $E$; zdanie
$\exists y\in E\;\forall x\in D\;P(x,y)$
jest natomiast prawdziwe dokładnie wtedy, gdy $E\ne\varnothing$.
Jeżeli $D\ne\varnothing$ i $E=\varnothing$, obie strony są
fałszywe, ponieważ nie ma świadka $y$.
Wszystkie te przypadki zachowują prawdziwość implikacji,
ale pokazują, że nie można bez dodatkowych założeń
zamienić jej w równoważność.
\end{uwaga}
\begin{przyklad}[Kontrprzykład o pustym zbiorze świadków]
\label{ex:mixed-quantifiers-empty-lxiii}
Weźmy $D=\varnothing$ i $E=\varnothing$.
Wtedy $\forall x\in D\;\exists y\in E\;P(x,y)$ jest
prawdziwe na mocy pustego zakresu kwantyfikatora ogólnego,
a $\exists y\in E\;\forall x\in D\;P(x,y)$ jest fałszywe,
ponieważ kwantyfikator egzystencjalny nie ma świadka.
Jest to kontrprzykład do implikacji odwrotnej niezależny
od arytmetyki liczb naturalnych. W standardowej semantyce
pierwszego rzędu nośnik całej struktury jest niepusty;
tutaj $D$ i $E$ są zakresami kwantyfikatorów ograniczonych,
które mogą być puste.
\end{przyklad}

\begin{przyklad}[Negowanie dwóch kwantyfikatorów]
Zaprzeczeniem
$\forall x\in\R\;\exists y\in\R\;(x+y=0)$
jest
$\exists x\in\R\;\forall y\in\R\;(x+y\ne0)$.
Przenosimy negację kolejno przez oba kwantyfikatory,
każdorazowo zmieniając ich typ, aż trafi do
wewnętrznego warunku. Oryginalne zdanie jest
prawdziwe: wystarczy dla danego $x$ wybrać $y=-x$.
Jego negacja jest zatem fałszywa.
\end{przyklad}

\section{Istnienie dokładnie jednego obiektu}\label{sec:unique-existence}
\begin{definicja}[Kwantyfikator jednoznacznego istnienia]
\label{def:unique}
Zapis $\exists!x\in D\;P(x)$ czytamy:
„istnieje \emph{dokładnie jeden} obiekt $x$ z $D$,
dla którego $P(x)$ jest prawdziwe”. Oznacza on
łącznie dwa warunki:
\begin{enumerate}[label=(\roman*)]
\item \emph{istnienie}: pewien $a\in D$ spełnia $P(a)$;
\item \emph{jednoznaczność}: jeśli $a,b\in D$ spełniają
      $P$, to $a=b$.
\end{enumerate}
Nie jest to nowy podstawowy kwantyfikator: jest to
wygodny skrót łączący $\exists$ i $\forall$.
\end{definicja}

\begin{twierdzenie}[Rozwinięcie symbolu $\exists!$]
\label{thm:unique-expansion-xli}
Dla dowolnej dziedziny $D$ (także pustej) i predykatu
$P(x)$ następujące zdania są równoważne:
\begin{align*}
&\exists!x\in D\;P(x),\\
&\exists x\in D\;\Bigl(P(x)\land
  \forall y\in D\,(P(y)\Rightarrow y=x)\Bigr).
\end{align*}
\end{twierdzenie}
\begin{proof}
Załóżmy, że pierwsze zdanie jest prawdziwe.
Z definicji istnieje świadek $a\in D$ spełniający $P(a)$.
Jeżeli dowolny $b\in D$ spełnia $P(b)$, jednoznaczność
wymusza $b=a$. Zatem dla $x=a$ prawdziwa jest
koniunkcja w zakresie kwantyfikatora egzystencjalnego
w drugim zdaniu, więc to zdanie zachodzi.

Odwrotnie, niech drugie zdanie będzie prawdziwe
i niech $a\in D$ będzie jego świadkiem.
Pierwszy człon koniunkcji daje $P(a)$, czyli istnienie.
Aby wykazać jednoznaczność, weźmy dowolne $b\in D$
spełniające $P(b)$. Drugi człon koniunkcji zastosowany
do $y=b$ daje $P(b)\Rightarrow b=a$, zatem $b=a$.
To jest dokładnie definicja jednoznacznego istnienia.
Jeżeli $D$ jest puste, oba zdania są fałszywe,
ponieważ w obu musi wystąpić świadek $x\in D$.
\end{proof}

\begin{przyklad}[Jednoznaczne rozwiązanie równania]
Zdanie $\exists!x\in\R\;(2x+3=7)$ jest prawdziwe.
Istnienie: $x=2$ spełnia równość, bo $2\cdot2+3=7$.
Jednoznaczność: jeśli $2x+3=7$, odejmujemy $3$ od
obu stron i dostajemy $2x=4$; dzieląc przez $2$,
uzyskujemy $x=2$. Każde rozwiązanie musi więc
być tym samym obiektem.
\end{przyklad}

\begin{przyklad}[Istnienie bez jednoznaczności]
$\exists x\in\Z\;(x^2=4)$ jest prawdziwe,
lecz $\exists!x\in\Z\;(x^2=4)$ jest fałszywe,
ponieważ $2^2=4$ i $(-2)^2=4$, a $2\ne-2$.
\end{przyklad}

\section{Jak budować matematyczny dowód?}\label{sec:proof-methods}
\paragraph{Co oznacza udowodnić twierdzenie?}
Twierdzenie ma \emph{tezę} i może mieć \emph{założenia}.
Dowód to skończony ciąg uzasadnionych przejść
od przyjętych założeń i wcześniej ustalonych faktów
do tezy. W dalszych partiach książki zdefiniujemy
formalny rachunek dowodowy, ale już teraz możemy
nauczyć się jego najważniejszych poprawnych schematów.

\subsection{Dowód bezpośredni}\label{sec:direct-proof}
Aby wykazać $p\Rightarrow q$, zakładamy $p$
i na podstawie definicji oraz znanych twierdzeń
wyprowadzamy $q$. Nie wystarczy sprawdzić
kilku przykładów spełniających $p$.

\begin{przyklad}[Suma dwóch liczb parzystych]
\label{ex:even-sum}
\emph{Teza.} Dla dowolnych liczb całkowitych $a,b$,
jeżeli obie są parzyste, to ich suma jest parzysta.

\emph{Dowód.} Z definicji parzystości $a=2k$
i $b=2l$ dla pewnych liczb całkowitych $k,l$.
Dodając równości, otrzymujemy
$a+b=2k+2l=2(k+l)$. Suma $k+l$ jest liczbą
całkowitą, więc $a+b$ ma postać dwukrotności
liczby całkowitej. Z definicji jest zatem parzysta.
\end{przyklad}

\begin{uwaga}[Parzyste i nieparzyste liczby całkowite]
Korzystamy z elementarnego faktu arytmetycznego, że każda
liczba całkowita $n$ ma dokładnie jedną z postaci
$n=2k$ oraz $n=2k+1$, gdzie $k\in\Z$.
Jest to szczególny przypadek dzielenia z resztą przez $2$:
reszta jest równa $0$ albo $1$. Pierwsza postać definiuje
parzystość, druga nieparzystość. Nie zachodzą jednocześnie,
ponieważ równość $2k=2l+1$ prowadziłaby do
$2(k-l)=1$, co jest niemożliwe dla liczb całkowitych.
Formalny dowód istnienia reszty z dzielenia podamy przy
budowaniu arytmetyki liczb naturalnych; nie wyprowadzamy
go z samych praw logiki.
\end{uwaga}

\subsection{Dowód przez przeciwstawienie}\label{sec:contrapositive-proof}
Z rozdziału~\ref{ch:propositional} wiemy, że
$p\Rightarrow q$ jest równoważne
$\neg q\Rightarrow\neg p$. Czasem znacznie
łatwiej udowodnić drugi z tych warunków.

\begin{przyklad}[Jeśli kwadrat jest parzysty, liczba jest parzysta]
\label{ex:even-square}
\emph{Teza.} Dla każdego $n\in\Z$, jeśli $n^2$
jest parzyste, to $n$ jest parzyste.

\emph{Dowód.} Udowodnimy implikację przeciwstawną.
Załóżmy, że $n$ jest \emph{nieparzyste}. Wtedy
$n=2k+1$ dla pewnego $k\in\Z$. Podnosimy
do kwadratu:
\[
n^2=(2k+1)^2=4k^2+4k+1=2(2k^2+2k)+1.
\]
Ponieważ $2k^2+2k$ jest liczbą całkowitą,
$n^2$ ma postać $2m+1$, a więc jest nieparzyste.
Pokazaliśmy, że z nieparzystości $n$ wynika
nieparzystość $n^2$. Na mocy równoważności
implikacji z przeciwstawną wynika stąd teza.
\end{przyklad}

\subsection{Dowód nie wprost}\label{sec:contradiction-proof}
Aby udowodnić zdanie $q$, możemy założyć jego
negację $\neg q$ i wyprowadzić sprzeczność:
jednoczesną prawdziwość pewnego zdania $r$
i jego negacji. W klasycznej logice oznacza to,
że $\neg q$ nie może być prawdziwe, więc $q$
musi być prawdziwe.

\begin{przyklad}[Niewymierność pierwiastka z dwóch]
\label{ex:sqrt-two}
Przypomnijmy znaczenie słów: liczba rzeczywista $t$
jest \emph{wymierna}, jeżeli $t=a/b$ dla pewnych
liczb całkowitych $a,b$ z $b\ne0$; w przeciwnym
razie jest \emph{niewymierna}. Przez $\sqrt2$
rozumiemy dodatnią liczbę rzeczywistą, której
kwadrat jest równy $2$. Wykorzystamy elementarne
twierdzenie arytmetyki, że ułamek można skrócić
do postaci, w której licznik i mianownik nie mają
wspólnego dodatniego dzielnika większego od $1$.
Formalne uzasadnienie tego faktu zostanie podane
w części poświęconej konstrukcji liczb naturalnych.

\emph{Teza.} Liczba $\sqrt2$ jest niewymierna.

\emph{Dowód.} Przypuśćmy przeciwnie, że $\sqrt2$
jest wymierna. Istnieją więc liczby całkowite
$a,b$ z $b\ne0$ i $\sqrt2=a/b$. Po skróceniu
ułamka możemy przyjąć, że $a$ i $b$ nie mają
wspólnego dzielnika większego od $1$; to znaczy,
że ułamek jest nieskracalny. Podnosząc równość
do kwadratu i mnożąc przez $b^2$, otrzymujemy
$a^2=2b^2$. Zatem $a^2$ jest parzyste, więc
$a$ jest parzyste na mocy Przykładu~\ref{ex:even-square}.
Możemy więc zapisać $a=2k$ dla pewnego $k\in\Z$.
Podstawiając do poprzedniej równości, mamy
$4k^2=2b^2$, a po podzieleniu przez $2$:
$b^2=2k^2$. Także $b^2$ jest parzyste, więc
$b$ jest parzyste. Wynika stąd, że zarówno $a$,
jak i $b$ są podzielne przez $2$, co przeczy
nieskracalności ułamka $a/b$. Otrzymana
sprzeczność obala przypuszczenie o wymierności
$\sqrt2$; dowodzi to tezy.
\end{przyklad}

\begin{uwaga}[Gdzie użyliśmy logiki klasycznej?]
Przejście „założenie $\neg q$ prowadzi do sprzeczności,
zatem $q$” korzysta z zasady usuwania podwójnej
negacji. W logice intuicjonistycznej trzeba staranniej
rozróżniać dowód $\neg\neg q$ od dowodu $q$.
W tej książce na razie pracujemy w logice klasycznej.
\end{uwaga}

\subsection{Dowód przez przypadki}\label{sec:proof-by-cases}
Jeżeli wiadomo, że zachodzi $p\lor q$, można
osobno wykazać tezę $r$ przy założeniu $p$,
a następnie przy założeniu $q$. Oznacza to,
że bez względu na to, która z dopuszczalnych
możliwości zachodzi, otrzymujemy $r$.

\begin{przyklad}[Iloczyn dwóch kolejnych liczb jest parzysty]
\emph{Teza.} Dla każdej liczby całkowitej $n$
iloczyn $n(n+1)$ jest parzysty.

\emph{Dowód.} Każda liczba całkowita jest parzysta
albo nieparzysta. W pierwszym przypadku $n=2k$
dla pewnego $k\in\Z$, więc
$n(n+1)=2[k(n+1)]$ jest parzyste. W drugim
przypadku $n=2k+1$, zatem $n+1=2(k+1)$
i $n(n+1)=2[n(k+1)]$ jest parzyste.
Oba przypadki wyczerpują możliwości, więc
teza zachodzi dla każdego $n\in\Z$.
\end{przyklad}

\subsection{Jak dowodzić równoważności i jednoznaczności?}
\label{sec:proof-biconditional}
Aby uzasadnić $p\Leftrightarrow q$, udowadniamy
osobno $p\Rightarrow q$ i $q\Rightarrow p$.
Aby uzasadnić $\exists!x\in D\;P(x)$ dla ustalonej
dziedziny $D$, wykazujemy istnienie obiektu z $D$
spełniającego $P$ (na przykład wskazując świadka),
a następnie porównujemy dowolne dwa obiekty z $D$
spełniające $P$ i wykazujemy, że muszą być równe.
Dowód samego istnienia nie uzasadnia jednoznaczności.

\begin{przyklad}[Równoważność dla parzystości]
\emph{Teza.} Dla każdej liczby całkowitej $n$
liczba $n$ jest parzysta wtedy i tylko wtedy,
gdy liczba $n+2$ jest parzysta.

\emph{Dowód.} Jeśli liczba $n$ jest parzysta, to $n=2k$
dla pewnego $k\in\Z$, a więc $n+2=2(k+1)$,
czyli liczba $n+2$ jest parzysta. Odwrotnie,
jeżeli liczba $n+2$ jest parzysta, to $n+2=2l$
dla pewnego $l\in\Z$. Po odjęciu $2$
uzyskujemy $n=2(l-1)$, więc również liczba $n$
jest parzysta. Wykazaliśmy oba kierunki.
\end{przyklad}

\section{Indukcja matematyczna --- pierwsze zastosowanie}
\label{sec:induction}
Liczb naturalnych jest nieskończenie wiele.
Nie da się zakończyć skończonego dowodu przez oddzielne
sprawdzanie wszystkich warunków $P(0),P(1),P(2),\ldots$.
Indukcja jest metodą,
która łączy sprawdzenie jednego przypadku z
regułą przejścia do następnego. Jej formalne
podstawy, wraz z konstrukcją liczb naturalnych,
rozwiną późniejsze rozdziały.

\begin{twierdzenie}[Zasada indukcji --- postać elementarna]
\label{thm:induction}
Niech $P(n)$ będzie predykatem dla $n\in\N$.
Jeśli zachodzą dwa warunki:
\begin{enumerate}[label=(\roman*)]
\item $P(0)$ jest prawdziwe (\emph{baza indukcji});
\item dla każdego $n\in\N$ z prawdziwości $P(n)$
      wynika prawdziwość $P(n+1)$ (\emph{krok indukcyjny}),
\end{enumerate}
to $P(n)$ jest prawdziwe dla każdego $n\in\N$.
\end{twierdzenie}

\begin{uwaga}[Status zasady indukcji]
Na tym etapie traktujemy zasadę indukcji jako
podstawową własność liczb naturalnych, a nie jako
wniosek z samych tablic prawdy. Później pokażemy,
jak zasada ta wiąże się z budową $\N$ i aksjomatami
arytmetyki. Nie wolno uznać jej za „oczywisty”
wniosek z faktu, że sprawdziliśmy kilka liczb.
\end{uwaga}

\begin{przyklad}[Wzór na sumę pierwszych liczb naturalnych]
\label{ex:gauss-induction}
\emph{Teza.} Dla każdego $n\in\N$ zachodzi
\[
0+1+\cdots+n=\frac{n(n+1)}2.
\]
\emph{Dowód.} Dla $n=0$ lewa strona wynosi $0$,
a prawa $0\cdot1/2=0$; baza zachodzi.

Ustalmy teraz dowolne $n\in\N$ i załóżmy,
że wzór jest prawdziwy dla tej liczby:
$0+1+\cdots+n=n(n+1)/2$. Chcemy wykazać
wzór dla $n+1$. Dodajemy $n+1$ do obu stron
równości przyjętej w założeniu indukcyjnym:
\begin{align*}
0+1+\cdots+n+(n+1)
&=\frac{n(n+1)}2+(n+1)\\
&=\frac{n(n+1)+2(n+1)}2\\
&=\frac{(n+1)(n+2)}2.
\end{align*}
Ostatni wyraz jest dokładnie prawą stroną
wzoru po zastąpieniu $n$ przez $n+1$.
Udowodniliśmy krok indukcyjny, więc
na mocy zasady indukcji wzór obowiązuje
dla wszystkich liczb naturalnych.
\end{przyklad}

\begin{przyklad}[Suma kolejnych liczb nieparzystych]
\label{ex:odd-sum-induction-xli}
Dla każdego $n\in\N$ zachodzi
$1+3+\cdots+(2n+1)=(n+1)^2$; po lewej jest
$n+1$ kolejnych liczb nieparzystych.
Dla $n=0$ lewa strona wynosi $1$, a prawa $(0+1)^2=1$.
Przyjmijmy teraz dla ustalonego $n\in\N$, że
$1+3+\cdots+(2n+1)=(n+1)^2$. Aby otrzymać wzór
dla $n+1$, trzeba dodać następny składnik, czyli
$2(n+1)+1=2n+3$. Z założenia indukcyjnego:
\[
1+3+\cdots+(2n+1)+(2n+3)
 =(n+1)^2+2n+3=(n+2)^2.
\]
Jest to dokładnie prawa strona wzoru dla $n+1$.
Warunki zasady indukcji są spełnione, więc
równość zachodzi dla każdego $n\in\N$.
\end{przyklad}

\begin{przyklad}[Dlaczego dowód dla jednego $n$ nie stanowi kroku indukcyjnego]
Załóżmy, że chcemy udowodnić $P(n)$ dla wszystkich
$n\in\N$. Sprawdzenie $P(0)$ i $P(1)$ nie wystarcza;
potrzebujemy zdania $\forall n\in\N\,(P(n)\Rightarrow P(n+1))$.
W kroku wybieramy \emph{dowolne} $n$, zakładamy dla niego
$P(n)$, a następnie wyprowadzamy $P(n+1)$.
Nie można zakładać $P(n+1)$ --- to właśnie teza kroku.
\end{przyklad}

\begin{przyklad}[Fałszywy krok indukcyjny ukryty w przykładach]
\label{ex:induction-counterexample-xli}
Rozważmy $P(n)$: „$n<3$” dla $n\in\N$.
Zdania $P(0)$, $P(1)$ i $P(2)$ są prawdziwe.
Co więcej, implikacje $P(0)\Rightarrow P(1)$
i $P(1)\Rightarrow P(2)$ są prawdziwe.
Jednak $P(3)$ jest fałszywe. Dokładnie dla $n=2$
implikacja $P(2)\Rightarrow P(3)$ jest fałszywa:
jej poprzednik ma wartość $1$, a następnik $0$.
Zatem z poprawnej bazy i kilku sprawdzonych
przejść nie wynika jeszcze warunek
$\forall n\in\N\,(P(n)\Rightarrow P(n+1))$.
\end{przyklad}

\section{Typowe pomyłki przy dowodzeniu}\label{sec:proof-errors}
\begin{przyklad}[Wstawienie liczby nie dowodzi zdania ogólnego]
Dla $n=2$ zachodzi $n^2\ge n$, ale nie dowodzi
to jeszcze twierdzenia
$\forall n\in\N\;(n^2\ge n)$.
Aby je wykazać, trzeba wyprowadzić nierówność
bez ograniczania się do jednego $n$.
Na przykład dla $n=0$ obie strony są równe,
a dla $n\ge1$ mamy $n^2-n=n(n-1)\ge0$.
Ta argumentacja rzeczywiście obejmuje każdy
możliwy naturalny $n$.
\end{przyklad}

\begin{przyklad}[Wybór świadka zależnego od zmiennej]
W prawdziwym zdaniu
$\forall x\in\R\;\exists y\in\R\;(x+y=0)$
wybrany świadek $y=-x$ zależy od rozważanego $x$.
Nie uprawnia to do zapisania
$\exists y\in\R\;\forall x\in\R\;(x+y=0)$.
Takie zdanie wymagałoby jednego ustalonego $y$,
działającego jednocześnie dla wszystkich $x$.
Dla $x=0$ wymuszałoby $y=0$, a dla $x=1$
wymuszałoby $y=-1$, co jest niemożliwe.
\end{przyklad}

\begin{przyklad}[Założenie tezy nie jest dowodem]
Chcąc wykazać $p\Rightarrow q$, wolno przyjąć
$p$, ale nie wolno bez uzasadnienia przyjąć $q$.
Z kolei w dowodzie nie wprost można przyjąć
$\neg q$, ale wtedy trzeba rzeczywiście
wyprowadzić sprzeczność z przyjętymi założeniami
i znanymi faktami.
\end{przyklad}

\section{Co wiemy przed wprowadzeniem zbiorów?}
\label{sec:logic-to-sets}
Potrafimy już czytać formuły z negacją, koniunkcją,
alternatywą, implikacją i równoważnością. Potrafimy
odróżnić $\forall$ od $\exists$, negować oba
kwantyfikatory, rozpoznawać zakresy zmiennych,
a także budować dowody bezpośrednie, przez
przeciwstawienie, nie wprost, przez przypadki
i przez indukcję. Te narzędzia są niezbędne,
aby precyzyjnie zdefiniować \emph{zawieranie zbiorów},
\emph{równość zbiorów}, \emph{relacje} i
\emph{funkcje}. Od tego rozpoczniemy następną część.

\chapter{Zbiory i operacje na zbiorach}\label{ch:sets}

\section{Dlaczego po logice przechodzimy do zbiorów?}
W rozdziałach~\ref{ch:propositional}--\ref{ch:quantifiers}
nauczyliśmy się odczytywać implikacje, równoważności
i kwantyfikatory. Teraz zastosujemy te narzędzia do opisu
kolekcji obiektów. W matematyce często interesuje nas
najpierw pewna kolekcja obiektów,
a dopiero potem struktury i działania na tej kolekcji. Zbiór punktów
spełniających równanie, zbiór rozwiązań nierówności, rodzina otwartych
podzbiorów przestrzeni czy zbiór wszystkich funkcji o określonych
własnościach to przejawy tej samej idei. Nawet do wypowiedzenia
stwierdzenia „każda liczba z $A$ należy do $B$” potrzebujemy pojęcia
przynależności i zawierania.

\begin{uwaga}[Dwa poziomy ścisłości]\label{rem:sets-metatheory}
W rozdziałach początkowych pracujemy z ustalonymi zbiorami i operacjami,
których istnienie przyjmujemy zgodnie ze zwykłą praktyką matematyczną.
Nie przyjmujemy zasady „każdy warunek $P(x)$ wyznacza zbiór
$\{x:P(x)\}$” bez określenia, skąd pochodzą elementy $x$. Podstawy
aksjomatyczne wrócą wraz z teorią ZFC. Wszystkie poniższe definicje
zbiorów przez warunek należy rozumieć jako wycinanie podzbioru z
wcześniej ustalonego zbioru albo jako skrót dla konstrukcji, której
istnienie jest już uzasadnione.
\end{uwaga}

\begin{definicja}[Przynależność i podzbiór]\label{def:subset}
Niech $A,B$ będą zbiorami. Piszemy $x\in A$, gdy $x$ jest elementem
$A$, oraz $x\notin A$ w przeciwnym przypadku. Mówimy, że $A$ jest
\emph{podzbiorem} $B$, i piszemy $A\subseteq B$, gdy
\[
\forall x\;(x\in A\Longrightarrow x\in B).
\]
Piszemy $A\subsetneq B$, gdy $A\subseteq B$ i $A\ne B$.
\end{definicja}

\begin{przyklad}\label{ex:subset-basic}
Dla $A=\{1,2\}$ i $B=\{1,2,3\}$ mamy $A\subsetneq B$,
ale $\{1,2\}\notin B$: elementy $B$ są liczbami, a nie zbiorem $\{1,2\}$.
Z kolei dla $C=\{\{1,2\},3\}$ zachodzi $\{1,2\}\in C$,
natomiast $1\notin C$. Rozróżnienie $\in$ i $\subseteq$ będzie
niezbędne przy definiowaniu relacji i funkcji.
\end{przyklad}

\begin{definicja}[Zbiór pusty, równość i ekstensjonalność]
Zbiór bez elementów oznaczamy $\varnothing$. Zbiory $A,B$ są równe
wtedy i tylko wtedy, gdy mają dokładnie te same elementy:
\[
A=B\quad\Longleftrightarrow\quad
\forall x\;(x\in A\Longleftrightarrow x\in B).
\]
Ostatnią zasadę nazywamy \emph{ekstensjonalnością zbiorów}.
\end{definicja}

\begin{uwaga}[Co stwierdza ekstensjonalność?]\label{rem:equality-extensionality-xlii}
Znak $=$ jest symbolem równości: oznacza identyczność obiektów.
Kierunek $A=B\Rightarrow\forall x\,(x\in A\Leftrightarrow x\in B)$
wynika z możliwości zastępowania równych obiektów w każdym zdaniu.
Kierunek odwrotny jest zasadą \emph{ekstensjonalności zbiorów}:
zbiory mające te same elementy są równe. Nie należy traktować
całej tej równoważności jako definicji znaku równości w logice;
w aksjomatycznej teorii ZF kierunek odwrotny będzie osobnym
aksjomatem. W niniejszym rozdziale przyjmujemy go jawnie.
\end{uwaga}

\begin{twierdzenie}[Kryterium dwóch zawierań]\label{thm:two-inclusions}
Dla dowolnych zbiorów $A,B$ zachodzi
$A=B$ wtedy i tylko wtedy, gdy $A\subseteq B$ oraz $B\subseteq A$.
Ponadto $\varnothing\subseteq A$ dla każdego zbioru $A$.
\end{twierdzenie}
\begin{proof}
Załóżmy najpierw $A=B$. Każdy element $A$ jest wtedy elementem
$B$, więc $A\subseteq B$; zamieniając role zbiorów otrzymujemy
$B\subseteq A$. Odwrotnie, załóżmy oba zawierania i ustalmy
dowolny obiekt $x$. Jeżeli $x\in A$, pierwsze zawieranie daje
$x\in B$. Jeżeli $x\in B$, drugie daje $x\in A$. Zatem
$x\in A\Longleftrightarrow x\in B$ dla każdego $x$, skąd $A=B$
na mocy ekstensjonalności. Wreszcie warunek
$x\in\varnothing\Longrightarrow x\in A$ jest spełniony dla
każdego $x$: nie istnieje $x\in\varnothing$, który mógłby go
naruszyć. To dowodzi ostatniej tezy.
\end{proof}

\begin{uwaga}[O kwantyfikatorze i pustym zbiorze]
W zdaniu „każdy element zbioru pustego ma własność $P$” nie jest
ukryte twierdzenie o własnościach jakiegoś szczególnego elementu.
Zdanie jest prawdziwe, ponieważ nie istnieje kontrprzykład.
Natomiast „istnieje element zbioru pustego o własności $P$” jest
fałszywe. To dokładnie zachowanie kwantyfikatorów nad pustą dziedziną
omówione w rozdziale~\ref{ch:quantifiers}.
\end{uwaga}


\begin{przyklad}[Czytanie zapisu zbioru przez warunek]
Zapis $A=\{n\in\N:n<3\}$ czytamy: „zbiór tych liczb
naturalnych $n$, dla których $n<3$”. Dla przyjętej konwencji
$\N=\{0,1,2,\ldots\}$ mamy $A=\{0,1,2\}$.
Dwukropek oznacza „takich, że”, a nie dzielenie.
Aby wykazać $n\in A$, musimy sprawdzić dwa fakty:
$n\in\N$ oraz $n<3$. Aby obalić $n\in A$, wystarczy
pokazać, że przynajmniej jeden z tych dwóch warunków zawodzi.
\end{przyklad}

\section{Operacje skończone}\label{sec:set-basic-ops}
\begin{definicja}[Suma, przecięcie, różnica]
Dla zbiorów $A,B$ określamy
\begin{align*}
A\cup B&=\{x:x\in A\text{ lub }x\in B\},\\
A\cap B&=\{x:x\in A\text{ i }x\in B\},\\
A\setminus B&=\{x\in A:x\notin B\}.
\end{align*}
Zbiory $A$ i $B$ nazywamy \emph{rozłącznymi}, gdy
$A\cap B=\varnothing$.
\end{definicja}

\begin{uwaga}[Istnienie zbiorów a zapis warunkowy]\label{rem:finite-set-existence-xlii}
W zapisie $\{x:x\in A\lor x\in B\}$ zmienna $x$ nie
przebiega przez „zbiór wszystkich obiektów”. Jest to skrót
\emph{charakteryzacji} sumy $A\cup B$ przez jej elementy.
W ZF istnienie $A\cup B$ otrzymujemy z aksjomatu pary
i aksjomatu sumy, stosując sumę elementów $\{A,B\}$.
Przecięcie $A\cap B$ i różnicę $A\setminus B$ otrzymujemy
przez wycinanie odpowiednio z $A$ zbiorów
$\{x\in A:x\in B\}$ oraz $\{x\in A:x\notin B\}$.
Nie jest tu potrzebna nieograniczona zasada tworzenia zbioru
z dowolnego warunku. Aksjomaty te omówimy dokładnie
w rozdziale~\ref{ch:zf-axioms}.
\end{uwaga}

\begin{przyklad}[Od zbioru rozwiązań do działań na zbiorach]
Niech $A=\{x\in\R:x^2<4\}=(-2,2)$ oraz
$B=\{x\in\R:x\geq 0\}=[0,\infty)$.
Wtedy $A\cap B=[0,2)$, $A\cup B=(-2,\infty)$,
$A\setminus B=(-2,0)$ oraz $B\setminus A=[2,\infty)$.
Punkt $2$ nie należy do $A$, więc należy do $B\setminus A$.
\end{przyklad}

\begin{center}
\begin{tikzpicture}[scale=.92]
  \draw[thick] (-.85,0) circle (1.25);
  \draw[thick] (.85,0) circle (1.25);
  \begin{scope}
    \clip (-.85,0) circle (1.25);
    \fill[blue!22] (.85,0) circle (1.25);
  \end{scope}
  \draw[thick] (-.85,0) circle (1.25);
  \draw[thick] (.85,0) circle (1.25);
  \node at (-1.47,0) {$A$};
  \node at (1.47,0) {$B$};
  \node[font=\small] at (0,-1.6) {Zacieniowana część to $A\cap B$.};
\end{tikzpicture}
\end{center}
\begin{twierdzenie}[Prawa przemienności, łączności i rozdzielności]
\label{thm:set-lattice}
Dla dowolnych zbiorów $A,B,C$ zachodzą
\begin{align*}
A\cup B&=B\cup A,& A\cap B&=B\cap A,\\
(A\cup B)\cup C&=A\cup(B\cup C),&
(A\cap B)\cap C&=A\cap(B\cap C),\\
A\cap(B\cup C)&=(A\cap B)\cup(A\cap C),&
A\cup(B\cap C)&=(A\cup B)\cap(A\cup C).
\end{align*}
\end{twierdzenie}
\begin{proof}
We wszystkich przypadkach używamy kryterium dwóch zawierań
(Twierdzenie~\ref{thm:two-inclusions}) lub, równoważnie, sprawdzamy
przynależność dowolnego elementu. Dla przemienności sumy:
$x\in A\cup B$ oznacza, że $x\in A$ lub $x\in B$; kolejność
tych dwóch możliwości nie ma znaczenia, więc jest to równoważne
$x\in B\cup A$. Identyczny argument ze spójnikiem „i” daje
przemienność przecięcia.

Dla łączności sumy warunek $x\in (A\cup B)\cup C$
oznacza $(x\in A\lor x\in B)\lor x\in C$, czyli
$x\in A\lor(x\in B\lor x\in C)$, czyli
$x\in A\cup(B\cup C)$. Analogiczna zamiana nawiasów dla
spójnika „i” dowodzi łączności przecięcia.

Wykażmy pierwszą rozdzielność w obu kierunkach. Jeżeli
$x\in A\cap(B\cup C)$, to $x\in A$ oraz
$(x\in B\lor x\in C)$. Jeżeli zachodzi $x\in B$,
to $x\in A\cap B$, w przeciwnym rozpatrywanym przypadku
$x\in C$, więc $x\in A\cap C$. W obu przypadkach
$x\in(A\cap B)\cup(A\cap C)$. Odwrotnie, jeśli $x$ należy
do tej ostatniej sumy, to należy do $A\cap B$ lub do $A\cap C$;
w obu przypadkach $x\in A$ i co najmniej jeden z warunków
$x\in B$, $x\in C$ jest spełniony. Stąd
$x\in A\cap(B\cup C)$. Druga rozdzielność wynika analogicznie:
$x\in A\cup(B\cap C)$ zachodzi wtedy i tylko wtedy, gdy
$(x\in A\lor x\in B)$ oraz $(x\in A\lor x\in C)$,
czyli gdy $x\in(A\cup B)\cap(A\cup C)$.
\end{proof}

\begin{definicja}[Dopełnienie względem zbioru odniesienia]
Jeżeli $A\subseteq U$ i $U$ jest ustalonym zbiorem odniesienia,
to \emph{dopełnieniem $A$ w $U$} nazywamy
$A^{c,U}:=U\setminus A$. Jeśli $U$ jest jasny z kontekstu,
piszemy $A^c$. Dopełnienie zawsze zależy od wyboru $U$.
\end{definicja}

\begin{twierdzenie}[Prawa de Morgana]\label{thm:de-morgan-finite}
Jeżeli $A,B\subseteq U$, to
\[
(A\cup B)^{c,U}=A^{c,U}\cap B^{c,U},\qquad
(A\cap B)^{c,U}=A^{c,U}\cup B^{c,U}.
\]
\end{twierdzenie}
\begin{proof}
Weźmy dowolne $x\in U$. Przynależność
$x\in (A\cup B)^{c,U}$ oznacza, że $x\notin A\cup B$.
Definicja sumy mówi, że nie zachodzi ani $x\in A$, ani
$x\in B$. Ponieważ $x\in U$, są to równoważnie warunki
$x\in A^{c,U}$ i $x\in B^{c,U}$. Otrzymujemy
$x\in A^{c,U}\cap B^{c,U}$. Wszystkie kroki są
równoważnościami, więc udowodniona została równość zbiorów.
W drugim przypadku $x\notin A\cap B$ oznacza,
że $x\notin A$ lub $x\notin B$, co przy $x\in U$
jest równoważne przynależności do sumy dopełnień.
Żaden z porównywanych zbiorów nie ma elementów poza $U$,
więc dowody dla $x\in U$ są wystarczające.
\end{proof}

\begin{twierdzenie}[Różnica zbiorów a dwa spójniki]\label{thm:difference-laws-xlii}
Dla dowolnych zbiorów $A,B,C$ zachodzą tożsamości
\[
(A\setminus B)\setminus C=A\setminus(B\cup C),
\qquad
A\setminus(B\setminus C)=(A\setminus B)\cup(A\cap C).
\]
W szczególności różnica zbiorów nie jest na ogół łączna.
\end{twierdzenie}
\begin{proof}
Ustalmy obiekt $x$. Na mocy definicji różnicy:
\begin{align*}
x\in(A\setminus B)\setminus C
&\iff x\in A\setminus B\text{ i }x\notin C\\
&\iff x\in A\text{ i }x\notin B\text{ i }x\notin C\\
&\iff x\in A\text{ i }x\notin B\cup C\\
&\iff x\in A\setminus(B\cup C).
\end{align*}
Każde przejście jest równoważnością, więc pierwszy wzór
wynika z ekstensjonalności. W drugim wzorze kluczowe
jest zanegowanie koniunkcji określającej $B\setminus C$:
\begin{align*}
x\in A\setminus(B\setminus C)
&\iff (x\in A)\land\neg((x\in B)\land(x\notin C))\\
&\iff (x\in A)\land((x\notin B)\lor(x\in C))\\
&\iff ((x\in A)\land(x\notin B))
  \lor((x\in A)\land(x\in C))\\
&\iff x\in(A\setminus B)\cup(A\cap C).
\end{align*}
W trzecim przejściu użyliśmy rozdzielności koniunkcji
względem alternatywy. Ostatni warunek oznacza dokładnie
przynależność do zbioru po prawej stronie.
\end{proof}
\begin{przyklad}[Kiedy zmiana nawiasów daje inny wynik?]\label{ex:difference-parentheses-xlii}
Dla $A=B=C=\{1\}$ mamy $B\setminus C=\varnothing$,
więc $A\setminus(B\setminus C)=\{1\}$.
Tymczasem $A\setminus B=\varnothing$, a zatem
$(A\setminus B)\setminus C=\varnothing$.
Pierwsza tożsamość twierdzenia oblicza drugie z tych wyrażeń
jako $A\setminus(B\cup C)$; druga tożsamość oblicza
pierwsze jako $(A\setminus B)\cup(A\cap C)$.
\end{przyklad}


\begin{twierdzenie}[Zawieranie a działania na zbiorach]\label{thm:subset-test}
Dla dowolnych zbiorów $A,B$ następujące warunki są równoważne:
\begin{enumerate}[label=(\roman*)]
\item $A\subseteq B$;
\item $A\cap B=A$;
\item $A\cup B=B$.
\end{enumerate}
\end{twierdzenie}
\begin{proof}
Załóżmy $A\subseteq B$. Jeśli $x\in A\cap B$, to
$x\in A$, stąd $A\cap B\subseteq A$. Z drugiej strony
jeśli $x\in A$, to również $x\in B$, więc
$x\in A\cap B$. Otrzymujemy (ii).

Niech teraz zachodzi (ii). Dla dowolnego $x\in A$ mamy
$x\in A\cap B$ z równości zbiorów, więc $x\in B$.
To dowodzi (i).

Następnie z (i) wynika (iii): jeśli $x\in A\cup B$,
to $x\in A$ albo $x\in B$; w pierwszym przypadku
$x\in B$ przez zawieranie, a w drugim bezpośrednio.
Zatem $A\cup B\subseteq B$. Odwrotnie każdy $x\in B$
należy do $A\cup B$, więc mamy równość.
Wreszcie jeśli zachodzi (iii) i $x\in A$, to
$x\in A\cup B=B$, a więc $x\in B$, co daje (i).
\end{proof}

\begin{przyklad}[Szybkie testowanie zawierania]
Dla $A=\{1,3\}$ i $B=\{1,2,3\}$ mamy
$A\cap B=A$ oraz $A\cup B=B$. Jeśli natomiast
$A=\{1,4\}$, to $A\cup B\ne B$, gdyż $4$ występuje
po lewej, lecz nie po prawej.
\end{przyklad}

\begin{definicja}[Różnica symetryczna]
Różnicą symetryczną $A$ i $B$ nazywamy
$A\symdiff B:=(A\setminus B)\cup(B\setminus A)$.
\end{definicja}
\begin{przyklad}
Dla $A=\{1,2,3\}$ i $B=\{2,3,4\}$ mamy
$A\symdiff B=\{1,4\}$. Różnica symetryczna zatrzymuje
dokładnie te elementy, które należą do \emph{jednego} ze zbiorów,
ale nie do obydwu.
\end{przyklad}

\begin{twierdzenie}[Własności różnicy symetrycznej]\label{thm:symdiff}
Dla dowolnych zbiorów $A,B,C$ zachodzi
$A\symdiff B=B\symdiff A$,
$(A\symdiff B)\symdiff C=A\symdiff(B\symdiff C)$,
$A\symdiff A=\varnothing$ i $A\symdiff\varnothing=A$.
\end{twierdzenie}
\begin{proof}
Definicja bezpośrednio pokazuje przemienność. Ponadto żaden
$x$ nie może jednocześnie należeć i nie należeć do $A$,
więc $A\symdiff A=\varnothing$; przy $B=\varnothing$
definicja daje $A\symdiff\varnothing=A$.

Pozostaje łączność. Ustalmy dowolny obiekt $x$.
Niech $p,q,r$ oznaczają odpowiednio zdania
$x\in A$, $x\in B$, $x\in C$. Z definicji różnicy
symetrycznej wynika, że $x\in A\symdiff B$ wtedy
i tylko wtedy, gdy dokładnie jedno spośród $p,q$
jest prawdziwe. Wszystkie osiem możliwych ustawień
wartości logicznych przedstawia tablica:
\[
\begin{array}{ccc|cc}
p&q&r&(p\mathbin{\triangle}q)\mathbin{\triangle}r
  &p\mathbin{\triangle}(q\mathbin{\triangle}r)\\\hline
0&0&0&0&0\\
0&0&1&1&1\\
0&1&0&1&1\\
0&1&1&0&0\\
1&0&0&1&1\\
1&0&1&0&0\\
1&1&0&0&0\\
1&1&1&1&1
\end{array}
\]
W dwóch ostatnich kolumnach symbol $\triangle$ oznacza
operację „dokładnie jedna wartość jest równa $1$”
wykonywaną na wartościach $0,1$, a nie na zbiorach.
Dla przykładu w wierszu $(1,0,1)$ mamy najpierw
$1\triangle0=1$, potem $1\triangle1=0$;
przy przeciwnym nawiasowaniu $0\triangle1=1$,
a następnie $1\triangle1=0$. W każdym wierszu
obie kolumny są zgodne. Na mocy kryterium równości
zbiorów oba nawiasowania mają więc dokładnie
te same elementy.
\end{proof}

\section{Zbiór potęgowy i rodziny indeksowane}\label{sec:families}
\begin{definicja}[Zbiór potęgowy]\label{def:powerset}
Dla zbioru $X$ oznaczamy przez
$\calP(X):=\{A:A\subseteq X\}$ zbiór wszystkich jego podzbiorów.
\end{definicja}
\begin{przyklad}
Dla $X=\{a,b\}$ mamy
$\calP(X)=\{\varnothing,\{a\},\{b\},\{a,b\}\}$.
W szczególności $\varnothing\in\calP(X)$ oraz
$X\in\calP(X)$, także gdy $X=\varnothing$.
Zauważmy, że $\calP(\varnothing)=\{\varnothing\}$ nie jest
zbiorem pustym.
\end{przyklad}

\begin{uwaga}[Element zbioru potęgowego nie musi być elementem $X$]
\label{rem:powerset-membership-xlii}
Niech $X=\{\varnothing\}$. Wówczas $X$ ma jeden element,
a jego zbiór potęgowy to
$\calP(X)=\{\varnothing,\{\varnothing\}\}$.
Zachodzi $\varnothing\in X$ oraz $X\in\calP(X)$,
lecz $X\notin X$: jedynym elementem $X$ jest
$\varnothing$, a $X\ne\varnothing$. Ponadto
$\calP(\varnothing)=\{\varnothing\}$, podczas gdy
$\calP(X)$ ma dwa elementy. Dzięki temu widać różnicę
między relacjami $\in$ i $\subseteq$ także w przypadku
zbiorów, których elementy same są zbiorami.
\end{uwaga}


\begin{twierdzenie}[Liczba podzbiorów zbioru skończonego]
\label{thm:powerset-finite}
Jeżeli zbiór $X$ ma dokładnie $n$ elementów, gdzie $n\in\N$,
to $\calP(X)$ ma dokładnie $2^n$ elementów.
\end{twierdzenie}
\begin{proof}
Przeprowadzimy indukcję po $n$. Jeżeli $n=0$, to $X=\varnothing$,
a $\calP(X)=\{\varnothing\}$ jest jednoelementowy;
$2^0=1$.
Załóżmy, że dla każdego zbioru mającego $n$ elementów
jego zbiór potęgowy ma $2^n$ elementów. Niech $X$ ma $n+1$
elementów. Wybierzmy jeden z nich, oznaczmy go $a$,
i połóżmy $Y=X\setminus\{a\}$. Wtedy $Y$ ma $n$ elementów.
Każdy podzbiór $C\subseteq X$ należy do dokładnie jednej
z dwóch grup: albo $a\notin C$, albo $a\in C$.
W pierwszej grupie znajdują się dokładnie podzbiory $Y$;
na mocy założenia indukcyjnego jest ich $2^n$.
W drugiej grupie każdy zbiór ma jednoznaczną postać
$D\cup\{a\}$, gdzie $D\subseteq Y$. Każdy wybór $D$
daje dokładnie jeden zbiór z drugiej grupy,
a każdy zbiór z tej grupy wyznacza dokładnie jeden
$D$ przez usunięcie $a$. Dlatego i w tej grupie jest
$2^n$ podzbiorów.
Grupy są rozłączne i wyczerpują wszystkie podzbiory $X$.
Łącznie jest ich $2^n+2^n=2^{n+1}$, co zamyka indukcję.
\end{proof}

\begin{twierdzenie}[Zbiór potęgowy i przecięcia]
\label{thm:powerset-intersection}
Dla dowolnych zbiorów $A,B$ zachodzi
$\calP(A\cap B)=\calP(A)\cap\calP(B)$.
Ponadto $\calP(A)\cup\calP(B)\subseteq\calP(A\cup B)$,
lecz drugie zawieranie nie musi być równością.
\end{twierdzenie}
\begin{proof}
Dla dowolnego zbioru $C$ zachodzi
$C\in\calP(A\cap B)$ wtedy i tylko wtedy, gdy
$C\subseteq A\cap B$. To oznacza, że każdy element $C$
należy jednocześnie do $A$ i $B$, czyli
$C\subseteq A$ oraz $C\subseteq B$.
Z definicji zbioru potęgowego jest to równoważne
$C\in\calP(A)\cap\calP(B)$; pierwsza równość jest dowiedziona.
Jeśli zaś $C\in\calP(A)\cup\calP(B)$, to
$C\subseteq A$ lub $C\subseteq B$, a więc zawsze
$C\subseteq A\cup B$. Stąd drugie zawieranie.
Że może ono być ostre, widzimy dla $A=\{1\}$,
$B=\{2\}$: zbiór $\{1,2\}$ należy do
$\calP(A\cup B)$, ale nie do
$\calP(A)\cup\calP(B)$.
\end{proof}

\begin{twierdzenie}[Monotoniczność zbioru potęgowego]
\label{thm:power-monotonic}
Dla dowolnych zbiorów $A,B$ zachodzi
$A\subseteq B$ wtedy i tylko wtedy, gdy
$\calP(A)\subseteq\calP(B)$.
\end{twierdzenie}
\begin{proof}
Jeżeli $A\subseteq B$ i $C\in\calP(A)$,
to $C\subseteq A$, a z przechodniości zawierania
$C\subseteq B$. Stąd $C\in\calP(B)$, a więc
$\calP(A)\subseteq\calP(B)$.
Odwrotnie, niech $\calP(A)\subseteq\calP(B)$.
Ponieważ $A\subseteq A$, mamy $A\in\calP(A)$.
Zatem $A\in\calP(B)$, co z definicji oznacza
$A\subseteq B$.
\end{proof}

\begin{definicja}[Rodzina indeksowana]
Niech $I$ będzie zbiorem i niech dla każdego $i\in I$
ustalono zbiór $A_i$. Kolekcję $(A_i)_{i\in I}$ nazywamy
\emph{rodziną indeksowaną zbiorów}. Ustalamy jeden
zbiór odniesienia $U$ i zakładamy $A_i\subseteq U$
dla wszystkich $i\in I$. Dla dowolnego $I$ definiujemy
\[
\bigcup_{i\in I}A_i
=\{x\in U:\exists i\in I\;x\in A_i\},\qquad
\bigcap_{i\in I}A_i
=\{x\in U:\forall i\in I\;x\in A_i\}.
\]
Przy $I=\varnothing$ powyższe definicje dają
$\bigcup_{i\in\varnothing}A_i=\varnothing$ oraz
$\bigcap_{i\in\varnothing}A_i=U$.
\end{definicja}

\begin{przyklad}[Powtórzenia w rodzinie indeksowanej]
\label{ex:indexed-repetitions-xlii}
Dla $I=\{1,2,3\}$ i $U=\{a,b\}$ określmy
$A_1=\{a\}$, $A_2=\{a\}$, $A_3=\{b\}$.
Indeksy $1$ i $2$ są różne, lecz wyznaczają ten sam zbiór.
Mimo powtórzenia $\bigcup_{i\in I}A_i=\{a,b\}$ oraz
$\bigcap_{i\in I}A_i=\varnothing$. Rodzina indeksowana
nie jest tym samym co zbiór jej różnych wartości:
informacja, który indeks odpowiada któremu zbiorowi,
może mieć znaczenie. Nie potrzebujemy do tego żadnego
wybierania elementów z poszczególnych $A_i$.
\end{przyklad}


\begin{uwaga}[Dlaczego puste przecięcie wymaga otoczenia?]
Dla pustego zbioru indeksów warunek „$x\in A_i$ dla każdego $i$”
jest spełniony niezależnie od $x$. Bez ustalenia, spośród jakich
obiektów wybieramy $x$, nie ma tu jednego zbioru wszystkich takich
obiektów. Dlatego przyjmujemy puste przecięcie względem
konkretnego $U$, a nie „zbioru wszystkich zbiorów”.
\end{uwaga}

\begin{przyklad}[Nie tylko skończone sumy i przecięcia]
Niech $A_n=(-1/n,1/n)\subseteq\R$ dla $n\ge1$.
Wówczas $A_{n+1}\subseteq A_n$ i
$\bigcap_{n\ge1}A_n=\{0\}$: liczba $0$ należy do każdego
$A_n$, a dla $x\ne0$ wybieramy $n>1/|x|$ i otrzymujemy
$x\notin A_n$. Natomiast
$\bigcup_{n\ge1}A_n=(-1,1)$, bo największy zbiór to $A_1$.
\end{przyklad}

\begin{twierdzenie}[Dowolne prawa de Morgana]\label{thm:arbitrary-de-morgan}
Niech $U$ będzie zbiorem, $I$ dowolnym zbiorem indeksów,
a $A_i\subseteq U$ dla każdego $i\in I$. Wtedy
\[
U\setminus\bigcup_{i\in I}A_i
=\bigcap_{i\in I}(U\setminus A_i),\qquad
U\setminus\bigcap_{i\in I}A_i
=\bigcup_{i\in I}(U\setminus A_i),
\]
gdzie przy pustym $I$ stosujemy umowę o przecięciu względem $U$.
\end{twierdzenie}
\begin{proof}
Ustalmy dowolne $x\in U$. Wtedy
\begin{align*}
x\in U\setminus\bigcup_{i\in I}A_i
&\iff \neg\bigl(\exists i\in I:\;x\in A_i\bigr)\\
&\iff \forall i\in I:\;x\notin A_i\\
&\iff x\in\bigcap_{i\in I}(U\setminus A_i).
\end{align*}
W drugim wzorze analogicznie:
\begin{align*}
x\in U\setminus\bigcap_{i\in I}A_i
&\iff \neg\bigl(\forall i\in I:\;x\in A_i\bigr)\\
&\iff \exists i\in I:\;x\notin A_i\\
&\iff x\in\bigcup_{i\in I}(U\setminus A_i).
\end{align*}
W drugim kroku użyliśmy praw negowania kwantyfikatorów
udowodnionych w Twierdzeniu~\ref{thm:negation-quantifiers}. Wszystkie
porównywane zbiory są podzbiorami $U$, więc stosujemy ekstensjonalność.
Dla $I=\varnothing$ pierwsza tożsamość daje
$U\setminus\varnothing=U$, druga
$U\setminus U=\varnothing$, zgodnie z przyjętą umową.
\end{proof}

\begin{twierdzenie}[Rozdzielność względem dowolnych rodzin]
\label{thm:arbitrary-distributivity}
Niech $I$ będzie dowolnym zbiorem indeksów, a $B$ oraz
$A_i$ ($i\in I$) niech będą podzbiorami jednego zbioru $U$.
Przyjmujemy, że przecięcie pustej rodziny podzbiorów $U$ jest równe $U$. Wtedy
\[
B\cap\bigcup_{i\in I}A_i
=\bigcup_{i\in I}(B\cap A_i),\qquad
B\cup\bigcap_{i\in I}A_i
=\bigcap_{i\in I}(B\cup A_i).
\]
\end{twierdzenie}
\begin{proof}
Niech $x\in U$. Aby należeć do pierwszego zbioru, $x$
musi należeć do $B$ i do co najmniej jednego $A_i$.
Oznacza to, że dla pewnego $i\in I$ jednocześnie
$x\in B$ oraz $x\in A_i$; jest to dokładnie warunek
przynależności do sumy po prawej stronie.

Dla drugiego wzoru załóżmy najpierw, że
$x\in B\cup\bigcap_i A_i$. Jeśli $x\in B$, to dla każdego
$i$ należy do $B\cup A_i$. Jeśli $x\notin B$, to z założenia
$x$ należy do wszystkich $A_i$, a zatem także do wszystkich
$B\cup A_i$. W obu przypadkach $x$ należy do przecięcia
po prawej stronie.
Odwrotnie, jeżeli $x$ należy do wszystkich $B\cup A_i$
i $x\in B$, to należy do sumy po lewej. Jeśli zaś $x\notin B$,
to z $x\in B\cup A_i$ dla każdego $i$ wynika $x\in A_i$
dla każdego $i$. Stąd $x\in\bigcap_i A_i$, a więc znowu
$x$ należy do sumy po lewej. Oba zawierania dają równość.
\end{proof}

\begin{przyklad}[Pusta rodzina jako sprawdzian hipotez]
Jeżeli $I=\varnothing$, to w pierwszej równości twierdzenia obie
strony są puste: po lewej $B\cap\varnothing$, po prawej suma
pustej rodziny. W drugiej równości po lewej otrzymujemy
$B\cup U=U$, a po prawej przecięcie pustej rodziny, czyli $U$.
Zatem obie równości pozostają prawdziwe bez hipotezy $I\ne\varnothing$.
\end{przyklad}

\begin{twierdzenie}[Różnica i dowolne sumy oraz przecięcia]
\label{thm:relative-complement-families}
Niech $U$ będzie zbiorem, $A\subseteq U$,
$I$ zbiorem indeksów oraz $B_i\subseteq U$ dla każdego $i\in I$.
Wówczas, przy konwencji przecięcia pustej rodziny względem $U$,
\[
 A\setminus\bigcup_{i\in I}B_i=\bigcap_{i\in I}(A\setminus B_i)
 \quad\text{dla }I\ne\varnothing,
 \qquad
 A\setminus\bigcap_{i\in I}B_i=\bigcup_{i\in I}(A\setminus B_i).
\]
Gdy $I=\varnothing$, pierwsza równość zachodzi po zastąpieniu
prawego pustego przecięcia przecięciem \emph{względem $A$},
a druga zachodzi bez zmian.
\end{twierdzenie}
\begin{proof}
Gdy $I\ne\varnothing$, dla dowolnego $x\in U$ mamy
\begin{align*}
 x\in A\setminus\bigcup_i B_i
 &\iff (x\in A)\land\neg\exists i\in I\ (x\in B_i)\\
 &\iff (x\in A)\land\forall i\in I\ (x\notin B_i)\\
 &\iff \forall i\in I\ ((x\in A)\land(x\notin B_i))\\
 &\iff x\in\bigcap_i(A\setminus B_i).
\end{align*}
Przejście z drugiego do trzeciego warunku wymaga $I\ne\varnothing$:
stały warunek $x\in A$ można wówczas przenieść pod kwantyfikator
$\forall i\in I$, a z otrzymanego warunku uniwersalnego
odzyskać $x\in A$. Przejście z trzeciego do czwartego warunku
wynika już z definicji przecięcia względem $U$.
Jeśli $I$ jest puste, lewa strona pierwszej równości
to $A$, a przecięcie puste względem $U$ to $U$; dlatego bez
wskazanej modyfikacji równość może być fałszywa.
Dla drugiej równości nie trzeba wymagać niepustości $I$:
\begin{align*}
 x\in A\setminus\bigcap_i B_i
 &\iff (x\in A)\land\neg\forall i\in I\ (x\in B_i)\\
 &\iff (x\in A)\land\exists i\in I\ (x\notin B_i)\\
 &\iff \exists i\in I\ ((x\in A)\land(x\notin B_i))\\
 &\iff x\in\bigcup_i(A\setminus B_i).
\end{align*}
W przypadku pustego $I$ zdania egzystencjalne są fałszywe,
więc obie strony drugiej równości są puste.
\end{proof}

\begin{uwaga}[Poprawna konwencja przy różnicach i pustej rodzinie]
\label{rem:empty-relative-intersection-xlii}
Dla dowolnego $I$ można zapisać pierwszy wzór
Twierdzenia~\ref{thm:relative-complement-families}
bez wyjątku, jeśli \emph{wszystkie} zbiory $A\setminus B_i$
rozpatrujemy jako podzbiory $A$, a puste przecięcie
w tym wzorze liczymy względem $A$:
\[
A\setminus\bigcup_{i\in I}B_i
=\{x\in A:\forall i\in I\;x\notin B_i\}.
\]
Dla $I=\varnothing$ prawa strona jest całym $A$.
Jeśli natomiast puste przecięcie po prawej stronie
zdefiniowano względem większego $U$, wynosi ono $U$;
wtedy równanie nie jest prawdziwe, gdy $A\ne U$.
Przeniesienie stałego warunku $x\in A$ pod kwantyfikator
$\forall i\in I$ jest równoważnym przekształceniem warunków
dla $x\in U$ tylko przy $I\ne\varnothing$.
Dla pustego $I$ można zachować równość \emph{zbiorów},
licząc przecięcie względem $A$; nie oznacza to jednak
równoważności obu warunków dla dowolnego $x\in U$.
\end{uwaga}


\section{Ciągi zbiorów i granice mnogościowe}\label{sec:set-limits}
W analizie często obserwujemy nie pojedynczy zbiór, ale zbiór
zmieniający się wraz z numerem etapu. Dwa naturalne pytania to:
„które punkty pozostają w zbiorach od pewnego momentu?” i
„które punkty wracają do zbiorów nieskończenie wiele razy?”.

\begin{definicja}[Granica dolna i górna zbiorów]\label{def:set-liminf-limsup}
Niech $(A_n)_{n\ge1}$ będzie ciągiem podzbiorów zbioru $U$.
Definiujemy
\[
\liminf_{n\to\infty}A_n
:=\bigcup_{m\ge1}\bigcap_{n\ge m}A_n,\qquad
\limsup_{n\to\infty}A_n
:=\bigcap_{m\ge1}\bigcup_{n\ge m}A_n.
\]
Pierwszy zbiór zawiera elementy należące do wszystkich $A_n$
od pewnego miejsca; drugi --- elementy należące do nieskończenie
wielu $A_n$.
\end{definicja}

\begin{twierdzenie}[Znaczenie obu granic]\label{thm:set-lims}
Niech $(A_n)_{n\ge1}$ będzie ciągiem podzbiorów $U$. Wtedy:
\begin{enumerate}[label=(\roman*)]
\item $x\in\liminf_n A_n$ wtedy i tylko wtedy, gdy istnieje
$m\ge1$ takie, że $x\in A_n$ dla wszystkich $n\ge m$;
\item $x\in\limsup_n A_n$ wtedy i tylko wtedy, gdy dla każdego
$m\ge1$ istnieje $n\ge m$ z $x\in A_n$;
\item $\liminf_n A_n\subseteq\limsup_n A_n$;
\item $U\setminus\limsup_n A_n
=\liminf_n (U\setminus A_n)$ i
$U\setminus\liminf_n A_n
=\limsup_n(U\setminus A_n)$.
\end{enumerate}
\end{twierdzenie}
\begin{proof}
Dla (i) rozwijamy definicję sumy:
$x\in\bigcup_m\bigcap_{n\ge m}A_n$ oznacza istnienie $m$
takiego, że $x$ należy do całego przecięcia po $n\ge m$;
to z definicji przecięcia znaczy $x\in A_n$ dla wszystkich
$n\ge m$. W (ii) kolejność kwantyfikatorów jest odwrotna:
$x$ musi należeć do $\bigcup_{n\ge m}A_n$ dla każdego
$m$, zatem dla każdego $m$ można wskazać odpowiednie $n\ge m$.

Jeżeli zachodzi (i), wybierzmy odpowiadające mu $m_0$.
Dla dowolnego $m$ weźmy $n\ge\max\{m,m_0\}$.
Wówczas $x\in A_n$, więc spełniony jest warunek (ii).
To pokazuje (iii). Wreszcie (iv) otrzymujemy dwukrotnie
stosując dowolne prawa de Morgana
(Twierdzenie~\ref{thm:arbitrary-de-morgan}); dla przykładu
\begin{align*}
U\setminus\limsup_n A_n
&=U\setminus\bigcap_m\bigcup_{n\ge m}A_n\\
&=\bigcup_m\left(U\setminus\bigcup_{n\ge m}A_n\right)\\
&=\bigcup_m\bigcap_{n\ge m}(U\setminus A_n).
\end{align*}
Ostatni zbiór to $\liminf_n(U\setminus A_n)$.
Dla drugiej tożsamości zapisujemy osobno wszystkie kroki:
\begin{align*}
U\setminus\liminf_n A_n
&=U\setminus\bigcup_m\bigcap_{n\ge m}A_n\\
&=\bigcap_m\left(U\setminus\bigcap_{n\ge m}A_n\right)\\
&=\bigcap_m\bigcup_{n\ge m}(U\setminus A_n)\\
&=\limsup_n(U\setminus A_n).
\end{align*}
W obu dowodach zastosowaliśmy prawo de Morgana do rodziny
zewnętrznej i ponownie do rodziny wewnętrznej.
\end{proof}

\begin{przyklad}[Zmienność bez granicy]
Niech $A_n=\{0\}$ dla parzystych $n$ oraz $A_n=\{1\}$
dla nieparzystych $n$. Ani $0$, ani $1$ nie występują w
$A_n$ od pewnego miejsca zawsze, więc $\liminf_n A_n=\varnothing$.
Obie liczby należą jednak do nieskończenie wielu zbiorów,
więc $\limsup_n A_n=\{0,1\}$.
\end{przyklad}

\begin{przyklad}[Przedziały skupiające się wokół zera]
\label{ex:set-limits-boundary-xlii}
Rozważmy w $U=\R$ dwa ciągi:
$A_n=[-1/n,1/n]$ oraz $B_n=(0,1/n)$ dla $n\ge1$.
W pierwszym przypadku $0\in A_n$ dla każdego $n$.
Jeśli $x\ne0$, wybierzmy liczbę naturalną $N>1/|x|$.
Dla każdego $n\ge N$ zachodzi $1/n<|x|$,
więc $x\notin A_n$. Z definicji granic wynika
$\liminf A_n=\limsup A_n=\{0\}$.
W drugim przypadku $0\notin B_n$ dla żadnego $n$;
liczby niedodatnie również do nich nie należą.
Jeżeli $x>0$, wybieramy $N>1/x$. Wtedy dla wszystkich
$n\ge N$ mamy $1/n<x$, czyli $x\notin B_n$.
Każda liczba rzeczywista należy więc do co najwyżej
skończenie wielu zbiorów $B_n$, stąd
$\liminf B_n=\limsup B_n=\varnothing$.
Przedziały $B_n$ zbliżają się geometrycznie do zera,
lecz granica mnogościowa nie zawiera zera: brak
przynależności w każdym kroku jest tu decydujący.
\end{przyklad}


\begin{twierdzenie}[Monotoniczne ciągi zbiorów]\label{thm:monotone-sets}
Jeżeli $A_n\subseteq A_{n+1}$ dla wszystkich $n\ge1$,
to $\liminf_n A_n=\limsup_n A_n=\bigcup_{n\ge1}A_n$.
Jeżeli natomiast $A_{n+1}\subseteq A_n$ dla wszystkich
$n\ge1$, to $\liminf_n A_n=\limsup_n A_n=\bigcap_{n\ge1}A_n$.
\end{twierdzenie}
\begin{proof}
Załóżmy, że ciąg jest rosnący. Jeśli $x$ należy do
$\bigcup_n A_n$, to istnieje $m$ takie, że $x\in A_m$.
Przez monotoniczność $x\in A_n$ dla wszystkich $n\ge m$,
a zatem $x\in\liminf_n A_n$ na mocy
Twierdzenia~\ref{thm:set-lims}(i). Odwrotnie,
$x\in\liminf_n A_n$ implikuje przynależność do przynajmniej
jednego $A_m$, więc do ich sumy. Mamy równość
$\liminf_n A_n=\bigcup_n A_n$. Ponieważ
$\liminf_n A_n\subseteq\limsup_n A_n$ i każdy element
$\limsup_n A_n$ należy do przynajmniej jednego $A_n$
(na przykład do któregoś $A_n$ z $n\ge1$), otrzymujemy też
równość z granicą górną.

Teraz niech ciąg będzie malejący. Jeśli $x\in\bigcap_n A_n$,
to $x$ należy do każdego $A_n$, więc do obu granic.
Jeżeli $x\notin\bigcap_n A_n$, istnieje $m$ takie, że
$x\notin A_m$. Z malejącości wynika $x\notin A_n$
dla wszystkich $n\ge m$. Wtedy $x\notin\limsup_n A_n$
na mocy Twierdzenia~\ref{thm:set-lims}(ii).
Ponieważ granica dolna zawiera się w górnej, dostajemy
obydwie żądane równości.
\end{proof}

\begin{przyklad}[Przecięcie pustej rodziny a logika]
Niech $U=\{a,b\}$ i $I=\varnothing$. Nie ma ani
jednego zbioru $A_i$ do sprawdzenia. Zdanie
„dla każdego $i\in I$ zachodzi $x\in A_i$”
jest prawdziwe dla $x=a$ i $x=b$, dlatego
$\bigcap_{i\in I}A_i=U$. Natomiast zdanie
„istnieje $i\in I$ z $x\in A_i$” jest dla obu
$x$ fałszywe, więc $\bigcup_{i\in I}A_i=\varnothing$.
To zastosowanie rozdziału~\ref{ch:quantifiers}, a nie
wyjątek od logiki.
\end{przyklad}

\begin{twierdzenie}[Niezmienniczość granic mnogościowych przy skończonej zmianie]\label{thm:lxiv-finite-set-modifications}
Niech $U$ będzie zbiorem, a $(A_n)_{n\ge1}$ i $(B_n)_{n\ge1}$
ciągami jego podzbiorów. Jeżeli istnieje $N\ge1$ takie, że
$A_n=B_n$ dla każdego $n\ge N$, to
$\liminf_n A_n=\liminf_n B_n$ oraz
$\limsup_n A_n=\limsup_n B_n$.
\end{twierdzenie}
\begin{proof}
Ustalmy $x\in U$. Na mocy Twierdzenia~\ref{thm:set-lims}(i)
warunek $x\in\liminf_n A_n$ oznacza, że istnieje $m\ge1$,
dla którego $x\in A_n$ przy każdym $n\ge m$.
Wówczas dla każdego $n\ge\max\{m,N\}$ mamy $x\in B_n$,
więc $x\in\liminf_n B_n$. Zamieniając role obu ciągów,
otrzymujemy także zawieranie odwrotne.

Podobnie, jeżeli $x\in\limsup_n A_n$, to dla dowolnego
$m\ge1$ istnieje $n\ge\max\{m,N\}$ z $x\in A_n$:
wynika to z Twierdzenia~\ref{thm:set-lims}(ii) zastosowanego
do indeksu $\max\{m,N\}$. Ponieważ $n\ge N$, zachodzi
$x\in B_n$. Stąd $x\in\limsup_n B_n$; zamiana ról
ciągów daje drugie zawieranie. Nie korzystaliśmy z aksjomatu wyboru.
\end{proof}
\begin{przyklad}[Zmiana pierwszego wyrazu a granice]\label{ex:lxiv-finite-set-modifications}
Niech $U=\{0,1\}$, $A_1=\{0\}$, $B_1=\{1\}$,
a dla każdego $n\ge2$ niech $A_n=B_n=\{0\}$.
Wtedy $\bigcup_n A_n=\{0\}$, natomiast
$\bigcup_n B_n=\{0,1\}$, więc zwykłe sumy są różne.
Jednak obie granice dolne i obie granice górne są równe $\{0\}$:
element $0$ należy do wszystkich zbiorów od drugiego wyrazu,
a element $1$ występuje tylko raz, w $B_1$.
\end{przyklad}

\section{Dwa ostrzeżenia na zakończenie}
\begin{przyklad}[Pozorne podobieństwo dwóch operacji]
Nie należy traktować odejmowania zbiorów tak jak odejmowania
liczb: znak $\setminus$ nie jest łączny. Pełny kontrprzykład,
a także poprawne wzory z oboma nawiasowaniami,
zapisano w Przykładzie~\ref{ex:difference-parentheses-xlii}
i Twierdzeniu~\ref{thm:difference-laws-xlii}.
\end{przyklad}

\begin{uwaga}[Paradoks Russella]\label{rem:russell}
Przypuśćmy na chwilę, że wolno tworzyć zbiór
$R=\{x:x\notin x\}$ obejmujący wszystkie obiekty o takiej
własności. Pytanie $R\in R$ prowadzi wówczas do równoważności
$R\in R\iff R\notin R$, czyli do sprzeczności.
Nie oznacza to, że działania na konkretnych zbiorach są
sprzeczne. Oznacza, że zakres dopuszczalnych konstrukcji zbiorów
musi zostać określony aksjomatycznie. Później zobaczymy,
że oddzielanie \emph{wewnątrz ustalonego zbioru} jest bezpieczną
operacją dopuszczoną przez ZF.
\end{uwaga}

\chapter{Pary uporządkowane i relacje}\label{ch:relations}

\section{Od nieuporządkowanych zbiorów do uporządkowanych par}
Korzystamy z pojęć i praw działań na zbiorach
z rozdziału~\ref{ch:sets}. Zbiór $\{a,b\}$ nie rozróżnia kolejności: zachodzi
$\{a,b\}=\{b,a\}$. Tymczasem punkt płaszczyzny z pierwszą
współrzędną $2$ i drugą $5$ powinien być inny niż punkt o
współrzędnych $5$ i $2$. Podobnie zapis „$a$ pozostaje w relacji z
$b$” może być prawdziwy bez prawdziwości zdania
„$b$ pozostaje w relacji z $a$”. Potrzebujemy zatem pojęcia pary,
w której pozycja elementu ma znaczenie.

\begin{definicja}[Para uporządkowana --- konstrukcja Kuratowskiego]
\label{def:kuratowski-pair}
Dla dowolnych zbiorów $a,b$ definiujemy
\[
(a,b):=\bigl\{\{a\},\{a,b\}\bigr\}.
\]
Jest to \emph{para uporządkowana} o pierwszej współrzędnej $a$
i drugiej współrzędnej $b$. Konstrukcję tę wybieramy jako
jedno z możliwych kodowań uporządkowanej pary za pomocą zbiorów.
\end{definicja}

\begin{twierdzenie}[Jednoznaczność współrzędnych pary]
\label{thm:kuratowski}
Dla dowolnych zbiorów $a,b,c,d$ zachodzi
\[
(a,b)=(c,d)\quad\Longleftrightarrow\quad a=c\text{ i }b=d.
\]
\end{twierdzenie}
\begin{proof}
Jeżeli $a=c$ i $b=d$, podstawienie do definicji
bezpośrednio daje równość par.

Załóżmy odwrotnie, że $(a,b)=(c,d)$. Zauważmy, że
\[
\bigcap\{\{a\},\{a,b\}\}=\{a\},\qquad
\bigcup\{\{a\},\{a,b\}\}=\{a,b\}.
\]
Są to przecięcie i suma elementów \emph{niepustej} rodziny
kodującej parę. Analogiczne równości zachodzą dla $c,d$.
Równość par oznacza równość kodujących je rodzin zbiorów,
a więc równość ich przecięć i sum. Otrzymujemy
$\{a\}=\{c\}$ oraz $\{a,b\}=\{c,d\}$.
Z pierwszej równości wynika $a=c$.

Aby wywnioskować $b=d$ z $\{a,b\}=\{a,d\}$,
rozważmy dwa przypadki. Jeśli $b=a$, to lewa strona
jest singletonem $\{a\}$. Równość wymusza
$d\in\{a\}$, zatem $d=a=b$. Jeżeli $b\ne a$,
to $b\in\{a,d\}$ i $b\ne a$, więc $b=d$.
W obu przypadkach druga współrzędna jest jednoznaczna.
\end{proof}

\begin{przyklad}[Dwa przypadki kodowania pary]\label{ex:kuratowski-cases-xlii}
Jeżeli $a=b$, to $\{a,b\}=\{a\}$, zatem
$(a,a)=\{\{a\}\}$, czyli kod jest zbiorem
jednoelementowym, którego jedynym elementem jest $\{a\}$.
Jeżeli $a\ne b$, to $\{a\}\ne\{a,b\}$ i kod
$(a,b)=\{\{a\},\{a,b\}\}$ ma dwa różne elementy.
Nie należy mylić liczby elementów \emph{samego kodu}
z liczbą współrzędnych: uporządkowana para zawsze
ma dwie współrzędne, nawet jeśli są one równe.
\end{przyklad}


\begin{uwaga}[Dlaczego definicja jest umowna?]
W wielu działach matematyki nie korzystamy z wnętrza kodowania
$(a,b)$; potrzebujemy wyłącznie własności
z Twierdzenia~\ref{thm:kuratowski}. Można wybrać inne
kodowanie spełniające to samo kryterium. Takie rozróżnienie
między obiektem a jego reprezentacją będzie ważne również
przy budowie liczb naturalnych w teorii ZF.
\end{uwaga}

\begin{definicja}[Iloczyn kartezjański]\label{def:cartesian}
Dla zbiorów $A,B$ definiujemy
$A\times B:=\{(a,b):a\in A,\ b\in B\}$.
Jeśli $A=\varnothing$ lub $B=\varnothing$, to
$A\times B=\varnothing$.
Dla $n\ge2$ możemy analogicznie rozważać uporządkowane
$n$-tki: definiujemy je rekurencyjnie, np.
$(a_1,a_2,a_3):=((a_1,a_2),a_3)$, a dla większego
$n$ dodajemy kolejne współrzędne. Iloczyn
$A_1\times\cdots\times A_n$ składa się z takich
$n$-tek, dla których $a_i\in A_i$ dla każdego $i$.
Warto zapamiętać, że zapis trójki jest tutaj kodowaniem,
a nie utożsamieniem dwóch różnych nawiasowań.
\end{definicja}

\begin{twierdzenie}[Dlaczego iloczyn kartezjański jest zbiorem?]
\label{thm:cartesian-existence-xlii}
Dla dowolnych zbiorów $A,B$ oraz $U:=A\cup B$ zachodzi
\[
A\times B\subseteq\calP\bigl(\calP(U)\bigr).
\]
W konsekwencji, gdy dostępne są operacje tworzenia pary,
sumy, zbioru potęgowego i wycinania podzbioru według
warunku, zbiór $A\times B$ rzeczywiście istnieje.
\end{twierdzenie}
\begin{proof}
Weźmy dowolną parę $(a,b)\in A\times B$. Z definicji
$a\in A$ i $b\in B$, a więc oba należą do $U$.
Wobec tego $\{a\}\subseteq U$ i $\{a,b\}\subseteq U$.
Definicja zbioru potęgowego daje
$\{a\}\in\calP(U)$ i $\{a,b\}\in\calP(U)$.
Konstrukcja Kuratowskiego mówi, że
$(a,b)=\{\{a\},\{a,b\}\}$; oba elementy tego zbioru
należą do $\calP(U)$, stąd $(a,b)\subseteq\calP(U)$.
Ponownie korzystając z definicji zbioru potęgowego,
otrzymujemy $(a,b)\in\calP(\calP(U))$.
Dowolność pary daje żądane zawieranie.

Mamy zatem \emph{ustalony zbiór kandydatów}
$\calP(\calP(U))$. Wycinając z niego te elementy $z$,
dla których istnieją $a\in A$ i $b\in B$ spełniające
$z=\{\{a\},\{a,b\}\}$, otrzymujemy dokładnie
$A\times B$. Zasada wycinania nie wymaga
nieograniczonego tworzenia zbiorów. W ZF jest to
zastosowanie odpowiednich aksjomatów, które
omówimy w rozdziale~\ref{ch:zf-axioms}.
\end{proof}


\begin{przyklad}
Jeśli $A=\{1,2\}$ i $B=\{x,y\}$, to
$A\times B=\{(1,x),(1,y),(2,x),(2,y)\}$.
Zwykle $A\times B\ne B\times A$, nawet gdy
$|A|=|B|$. Na przykład $(1,x)\in A\times B$,
ale $(1,x)\notin B\times A$ przy założeniu, że
$1,2,x,y$ są parami różne.
\end{przyklad}

\begin{twierdzenie}[Iloczyn zachowuje zawieranie]
\label{thm:cart-monotonic}
Niech $A,A',B,B'$ będą zbiorami. Jeżeli
$A\subseteq A'$ i $B\subseteq B'$, to
$A\times B\subseteq A'\times B'$.
Ponadto dla dowolnych zbiorów $A,B,C$ zachodzi
\[
A\times(B\cup C)=(A\times B)\cup(A\times C),\qquad
A\times(B\cap C)=(A\times B)\cap(A\times C).
\]
\end{twierdzenie}
\begin{proof}
Weźmy $(a,b)\in A\times B$. Wtedy $a\in A$, $b\in B$,
a więc z założeń $a\in A'$ oraz $b\in B'$.
Zatem $(a,b)\in A'\times B'$, co dowodzi zawierania.

Dla sumy ustalmy dowolną parę $(a,b)$.
Należy ona do $A\times(B\cup C)$ wtedy i tylko wtedy,
gdy $a\in A$ oraz $b\in B\cup C$, co znaczy,
że $a\in A$ i ($b\in B$ lub $b\in C$).
To dokładnie warunek $(a,b)\in A\times B$ lub
$(a,b)\in A\times C$. Dla przecięcia
zamieniamy w tym samym rozumowaniu „lub” na „i”.
Kryterium jednoznaczności par
(Twierdzenie~\ref{thm:kuratowski}) gwarantuje, że każda para
w iloczynie jest rozpoznawana po swoich współrzędnych.
\end{proof}

\begin{uwaga}[Wielokrotne iloczyny]
Istnieje kanoniczna bijekcja między $(A\times B)\times C$
a $A\times(B\times C)$, wysyłająca $((a,b),c)$ na
$(a,(b,c))$, ale przy przyjętym kodowaniu uporządkowanych par
nie muszą to być dosłownie równe zbiory.
Zwyczajowe zapisy trójek ukrywają tę różnicę.
\end{uwaga}

\begin{przyklad}[Ile jest par w skończonym iloczynie?]
Jeśli $A=\{a_1,a_2\}$ oraz $B=\{b_1,b_2,b_3\}$,
to dla każdego z dwóch pierwszych elementów mamy trzy
możliwe drugie współrzędne; otrzymujemy sześć par.
Ogólnie, jeżeli $A$ ma $m$ elementów, a $B$ ma
$n$ elementów, to $A\times B$ ma $mn$ elementów.
Równość nie wymaga, aby $A$ i $B$ były rozłączne,
gdyż pozycja współrzędnej rozstrzyga o parze.
\end{przyklad}

\section{Definicja relacji}\label{sec:relation-definition}
\begin{definicja}[Relacja binarna]\label{def:binary-relation}
Niech $A,B$ będą zbiorami. \emph{Relacją binarną z $A$ do $B$}
nazywamy dowolny podzbiór $R\subseteq A\times B$.
Zapis $aRb$ jest skrótem warunku $(a,b)\in R$.
Jeśli $A=B$, mówimy również o relacji \emph{na $A$}.
Dla ustalonej dodatniej liczby naturalnej $n$ relacją
$n$-argumentową na zbiorach $A_1,\ldots,A_n$ nazywamy
dowolny podzbiór $A_1\times\cdots\times A_n$;
na potrzeby tego zapisu jednoelementową krotkę
rozumiemy po prostu jako jej jedyną współrzędną.
\end{definicja}

\begin{przyklad}[Liczba możliwych relacji]\label{ex:number-relations-xlii}
Niech $A=\{1,2\}$ i $B=\{u,v\}$, gdzie wszystkie
wskazane elementy w odpowiednich zbiorach są różne.
Iloczyn $A\times B$ składa się z czterech par:
$(1,u),(1,v),(2,u),(2,v)$. Relacja z $A$ do $B$
jest \emph{dowolnym podzbiorem} tego iloczynu.
Dla każdej z czterech par niezależnie podejmujemy
decyzję „należy” albo „nie należy”, co daje $2^4=16$
różnych relacji. Wśród nich są relacja pusta, pełna
i na przykład relacja zawierająca tylko $(1,u)$.
Ogólnie dla skończonych $A,B$ o odpowiednio $m,n$
elementach istnieje dokładnie $2^{mn}$ relacji
z $A$ do $B$: wynika to z Twierdzenia~\ref{thm:powerset-finite}
i liczby elementów iloczynu kartezjańskiego.
Nie każda z tych relacji jest funkcją.
\end{przyklad}



\begin{przyklad}[Mniejsze niż]
Zbiór $R_<:=\{(m,n)\in\Z\times\Z:m<n\}$
jest relacją na $\Z$. Zapis $mR_<n$ jest po prostu
innym zapisem $m<n$. W szczególności $2R_<3$,
ale $3R_<2$ nie zachodzi.
\end{przyklad}

\begin{przyklad}[Podzielność i wybór zbioru bazowego]
Na $\N_{>0}:=\{1,2,3,\ldots\}$ relacja
$m\mid n\iff \exists k\in\N_{>0}\colon n=mk$
ma inne własności niż podzielność na $\Z$.
Na przykład $2\mid 6$ i $3\mid 6$, ale $2\nmid3$.
Na $\Z$ mamy $2\mid(-2)$ i $(-2)\mid2$ przy $2\ne-2$.
Będzie to ważny kontrprzykład do antysymetrii.
\end{przyklad}

\begin{definicja}[Dziedzina i zbiór wartości relacji]
\label{def:relation-dom-ran}
Niech $A,B$ będą zbiorami i $R\subseteq A\times B$. Określamy
\[
\dom(R)=\{a\in A:\exists b\in B\ (a,b)\in R\},\qquad
\ran(R)=\{b\in B:\exists a\in A\ (a,b)\in R\}.
\]
Zbiór $\dom(R)$ zawiera argumenty występujące na pierwszej,
a $\ran(R)$ na drugiej współrzędnej jakiejś pary z $R$.
\end{definicja}

\begin{przyklad}[Relacja z luką w dziedzinie]
Niech $A=\{1,2,3\}$, $B=\{u,v\}$ i
$R=\{(1,u),(1,v),(3,v)\}$.
Wtedy $\dom(R)=\{1,3\}$ oraz $\ran(R)=\{u,v\}$.
Element $2$ należy do ustalonego $A$, ale nie do dziedziny
relacji $R$. Co więcej, jednemu elementowi $1$ przypisano
\emph{dwa} elementy $u,v$. Relacja nie musi więc być funkcją;
precyzyjną definicję funkcji podamy w rozdziale o odwzorowaniach.
\end{przyklad}

\begin{definicja}[Obraz i przeciwobraz względem relacji]
\label{def:relation-image}
Niech $A,B$ będą zbiorami, $R\subseteq A\times B$,
$X\subseteq A$ i $Y\subseteq B$. Definiujemy
\[
R[X]:=\{b\in B:\exists a\in X\ (a,b)\in R\},\qquad
R^{-1}[Y]:=\{a\in A:\exists b\in Y\ (a,b)\in R\}.
\]
Nazywamy je odpowiednio \emph{obrazem $X$} i
\emph{przeciwobrazem $Y$ względem relacji}.
Dla ogólnej relacji przeciwobraz definiujemy więc przez
\emph{istnienie} powiązanego elementu. Gdy relacja będzie wykresem
funkcji, odzyskamy zwykłe przeciwobrazy funkcji. W szczególności
nie wolno automatycznie przenosić wszystkich wzorów znanych
z przeciwobrazów funkcji na dowolne relacje.
\end{definicja}

\begin{twierdzenie}[Obrazy relacji i działania na zbiorach]
\label{thm:relation-image-unions}
Niech $A,B,I$ będą zbiorami, $R\subseteq A\times B$,
a $(X_i)_{i\in I}$ rodziną podzbiorów $A$. Wtedy
\[
R\left[\bigcup_{i\in I}X_i\right]
=\bigcup_{i\in I}R[X_i].
\]
Dla $X,Z\subseteq A$ zachodzi ponadto
$R[X\cap Z]\subseteq R[X]\cap R[Z]$.
\end{twierdzenie}
\begin{proof}
Dla $b\in B$ przynależność do lewej strony pierwszej
równości oznacza istnienie $a\in\bigcup_i X_i$
takiego, że $(a,b)\in R$. Z definicji sumy istnieje
zatem indeks $i$ z $a\in X_i$. Otrzymujemy $b\in R[X_i]$,
a więc $b$ należy do prawej strony. Odwrotnie,
jeżeli $b\in R[X_i]$ dla pewnego $i$, istnieje
$a\in X_i\subseteq\bigcup_jX_j$ z $(a,b)\in R$;
stąd $b$ należy do lewej strony.

Jeżeli $b\in R[X\cap Z]$, wybierzmy $a\in X\cap Z$
z $(a,b)\in R$. Ten sam element $a$ świadczy jednocześnie
o $b\in R[X]$ i $b\in R[Z]$. Stąd drugie zawieranie.
\end{proof}

\begin{definicja}[Jednoznaczność po obu stronach relacji]
\label{def:relation-unique-fibers-xlii}
Niech $A,B$ będą zbiorami i niech $R\subseteq A\times B$. Relację $R$ nazwiemy
\emph{prawostronnie jednoznaczną}, gdy dla wszystkich
$a\in A$ i $b_1,b_2\in B$ warunki $aRb_1$ i $aRb_2$
implikują $b_1=b_2$. Relację $R$ nazwiemy
\emph{lewostronnie jednoznaczną}, gdy dla wszystkich
$b\in B$ i $a_1,a_2\in A$ warunki $a_1Rb$ i $a_2Rb$
implikują $a_1=a_2$. Nie żądamy tu istnienia
żadnej pary dla każdego elementu $A$ lub $B$.
\end{definicja}

\begin{twierdzenie}[Kiedy relacyjne obrazy zachowują przecięcia?]
\label{thm:relation-images-intersections-xlii}
Niech $A,B$ będą dowolnymi zbiorami, a
$R\subseteq A\times B$. Następujące warunki są
równoważne:
\begin{enumerate}[label=(\roman*)]
\item $R$ jest lewostronnie jednoznaczna;
\item dla wszystkich $X,Z\subseteq A$ zachodzi
$R[X\cap Z]=R[X]\cap R[Z]$.
\end{enumerate}
Ponadto następujące warunki są równoważne:
\begin{enumerate}[label=(\roman*)]
\item $R$ jest prawostronnie jednoznaczna;
\item dla wszystkich $Y,W\subseteq B$ zachodzi
$R^{-1}[Y\cap W]=R^{-1}[Y]\cap R^{-1}[W]$.
\end{enumerate}
\end{twierdzenie}
\begin{proof}
Załóżmy najpierw lewostronną jednoznaczność. Znamy już
zawieranie $R[X\cap Z]\subseteq R[X]\cap R[Z]$
(Twierdzenie~\ref{thm:relation-image-unions}).
Dla zawierania odwrotnego weźmy $b\in R[X]\cap R[Z]$.
Istnieje $a_1\in X$ takie, że $a_1Rb$, i istnieje
$a_2\in Z$ takie, że $a_2Rb$. Lewostronna jednoznaczność
daje $a_1=a_2$. Ten wspólny element należy więc do
$X\cap Z$ i jest związany z $b$, stąd $b\in R[X\cap Z]$.
To dowodzi równości.

Załóżmy teraz równość obrazów przecięć dla \emph{wszystkich}
$X,Z\subseteq A$. Przypuśćmy, że $a_1Rb$ i $a_2Rb$.
Gdyby $a_1\ne a_2$, to dla singletonów
$X=\{a_1\}$ i $Z=\{a_2\}$ mielibyśmy
$X\cap Z=\varnothing$, a zatem $R[X\cap Z]=\varnothing$.
Z drugiej strony $b$ należy do obu zbiorów $R[X]$
i $R[Z]$, więc jest elementem ich przecięcia,
sprzecznie z założoną równością. Koniecznie $a_1=a_2$.

Dla drugiej części dowód prowadzimy z pełną kontrolą
orientacji par. Przy prawostronnej jednoznaczności
weźmy $a\in R^{-1}[Y]\cap R^{-1}[W]$.
Istnieje $b_1\in Y$ z $aRb_1$ i istnieje $b_2\in W$
z $aRb_2$. Jednoznaczność daje $b_1=b_2$,
więc ten element leży w $Y\cap W$, skąd
$a\in R^{-1}[Y\cap W]$. Przeciwne zawieranie
wynika bez założeń z definicji obrazu relacji odwrotnej.

Wreszcie jeśli druga równość obowiązuje dla wszystkich
$Y,W$, to z $aRb_1$ i $aRb_2$ oraz $b_1\ne b_2$
dostajemy sprzeczność, biorąc $Y=\{b_1\}$,
$W=\{b_2\}$: lewy przeciwobraz ich pustego przecięcia
jest pusty, natomiast $a$ należy do prawej strony.
Zatem prawostronna jednoznaczność jest konieczna.
\end{proof}


\begin{przyklad}[Dlaczego przecięcie daje tylko zawieranie?]
Niech $A=\{1,2\}$, $B=\{u\}$ i
$R=\{(1,u),(2,u)\}$. Dla $X=\{1\}$ i $Z=\{2\}$
mamy $X\cap Z=\varnothing$, a więc
$R[X\cap Z]=\varnothing$. Jednak
$R[X]\cap R[Z]=\{u\}$.
Obrazy różnych argumentów mogą się przecinać, nawet gdy
same zbiory argumentów są rozłączne.
\end{przyklad}

\begin{przyklad}[Dlaczego relacyjny przeciwobraz nie zachowuje przecięć]
Niech $A=\{a\}$, $B=\{u,v\}$ oraz
$R=\{(a,u),(a,v)\}\subseteq A\times B$.
Dla $Y=\{u\}$ i $Z=\{v\}$ mamy
$R^{-1}[Y\cap Z]=R^{-1}[\varnothing]=\varnothing$, ale
$R^{-1}[Y]\cap R^{-1}[Z]=\{a\}$.
Powodem jest to, że istnieją \emph{dwa różne} elementy
związane z $a$. Odróżnijmy to od przyszłego pojęcia funkcji,
gdzie każdemu argumentowi przypisany jest tylko jeden wynik.
\end{przyklad}


\begin{przyklad}[Relacja jako lista par]
Niech $A=\{1,2\}$, $B=\{a,b\}$ i
$R=\{(1,a),(2,a)\}\subseteq A\times B$.
Mamy $1Ra$ i $2Ra$, ale $1Rb$ nie zachodzi.
Relacja mówi wyłącznie o tym, które pary należą do $R$;
nie zobowiązuje nas do przypisania każdemu elementowi $A$
jednej i tylko jednej wartości. Ta dodatkowa własność
odróżni funkcje w rozdziale~\ref{ch:functions}.
\end{przyklad}

\section{Relacja odwrotna, identycznościowa i składanie}
\label{sec:relation-algebra}
\begin{definicja}[Relacja odwrotna i identycznościowa]
\label{def:inverse-identity-rel}
Niech $A,B$ będą zbiorami. Dla $R\subseteq A\times B$
relacją odwrotną jest
\[
R^{-1}:=\{(b,a)\in B\times A:(a,b)\in R\}.
\]
Dla dowolnego zbioru $A$ relacja identycznościowa ma postać
$\Delta_A:=\{(a,a):a\in A\}$.
Relacją pełną na $A$ jest $A\times A$, zaś relacją pustą
na $A$ --- $\varnothing\subseteq A\times A$.
\end{definicja}

\begin{twierdzenie}[Obraz względem relacji odwrotnej]
\label{thm:inverse-rel-image}
Niech $A,B$ będą zbiorami, $R\subseteq A\times B$,
$X\subseteq A$ i $Y\subseteq B$. Wówczas
\[
R^{-1}[Y]=\{a\in A:\exists b\in Y\ (a,b)\in R\},
\qquad R[X]=\{b\in B:\exists a\in X\ (b,a)\in R^{-1}\}.
\]
Ponadto $\dom(R)=R^{-1}[B]$ oraz $\ran(R)=R[A]$.
\end{twierdzenie}
\begin{proof}
Weźmy $a\in A$. Z definicji obrazu dla relacji odwrotnej,
$a\in R^{-1}[Y]$ oznacza istnienie $b\in Y$ takiego,
że $(b,a)\in R^{-1}$. Z definicji relacji odwrotnej
zachodzi to wtedy i tylko wtedy, gdy $(a,b)\in R$.
Otrzymaliśmy pierwszą równość. Podobnie $b\in R[X]$
wtedy i tylko wtedy, gdy pewien $a\in X$ spełnia
$(a,b)\in R$, co po odwróceniu pary jest równoważne
$(b,a)\in R^{-1}$. To dowodzi drugiego wzoru.
Jeżeli w pierwszym wzorze podstawimy $Y=B$, otrzymamy
dokładnie definicję $\dom(R)$. Podstawiając $X=A$ w
drugim wzorze, otrzymamy definicję $\ran(R)$.
\end{proof}

\begin{uwaga}[Dwa zastosowania $R^{-1}$]
Symbol $R^{-1}$ oznacza relację powstałą przez odwrócenie
współrzędnych. Zapis $R^{-1}[Y]$ w
Definicji~\ref{def:relation-image} to obraz $Y$ względem
tej właśnie relacji odwrotnej; są to zgodne znaczenia.
\end{uwaga}

\begin{definicja}[Składanie relacji]\label{def:relation-composition}
Niech $A,B,C$ będą zbiorami, $R\subseteq A\times B$
i $S\subseteq B\times C$. Relację
$S\circ R\subseteq A\times C$ definiujemy przez
\[
(a,c)\in S\circ R\quad\Longleftrightarrow\quad
\exists b\in B:\;(a,b)\in R\ \text{oraz}\ (b,c)\in S.
\]
Najpierw przechodzimy relacją $R$ z $A$ do $B$,
a następnie relacją $S$ z $B$ do $C$.
\end{definicja}

\begin{przyklad}[Wyszukiwanie dwustopniowe]
Niech $A=\{a_1,a_2\}$, $B=\{b_1,b_2\}$,
$C=\{c_1,c_2\}$ oraz
$R=\{(a_1,b_1),(a_1,b_2),(a_2,b_2)\}$,
$S=\{(b_1,c_1),(b_2,c_2)\}$.
Wtedy
$S\circ R=\{(a_1,c_1),(a_1,c_2),(a_2,c_2)\}$.
Na przykład $(a_1,c_2)$ należy do składania, gdyż
możemy wskazać pośredni element $b_2$.
\end{przyklad}

\begin{center}
\begin{tikzpicture}[>=Stealth, every node/.style={font=\small}]
\node (a1) at (0,1.2) {$a_1$}; \node (a2) at (0,0) {$a_2$};
\node (b1) at (3,1.2) {$b_1$}; \node (b2) at (3,0) {$b_2$};
\node (c1) at (6,1.2) {$c_1$}; \node (c2) at (6,0) {$c_2$};
\draw[->,thick] (a1)--(b1); \draw[->,thick] (a1)--(b2);
\draw[->,thick] (a2)--(b2); \draw[->,thick] (b1)--(c1);
\draw[->,thick] (b2)--(c2);
\node at (1.5,1.7) {$R$}; \node at (4.5,1.7) {$S$};
\end{tikzpicture}

{\small Przejścia przez punkt pośredni opisują relację $S\circ R$.}
\end{center}

\begin{lemat}[Monotoniczność składania]
\label{lem:composition-monotone}
Niech $A,B,C$ będą zbiorami, $R_1,R_2\subseteq A\times B$
oraz $S_1,S_2\subseteq B\times C$. Jeśli
$R_1\subseteq R_2$ i $S_1\subseteq S_2$,
to $S_1\circ R_1\subseteq S_2\circ R_2$.
\end{lemat}
\begin{proof}
Weźmy $(a,c)\in S_1\circ R_1$. Istnieje wtedy $b\in B$
takie, że $(a,b)\in R_1$ i $(b,c)\in S_1$.
Z założonych zawierań wynika $(a,b)\in R_2$
oraz $(b,c)\in S_2$. Ten sam element pośredni $b$
świadczy zatem o $(a,c)\in S_2\circ R_2$.
\end{proof}

\begin{twierdzenie}[Łączność składania relacji]\label{thm:relation-assoc}
Niech $A,B,C,D$ będą zbiorami, $R\subseteq A\times B$,
$S\subseteq B\times C$ i $T\subseteq C\times D$. Wtedy
\[
T\circ(S\circ R)=(T\circ S)\circ R.
\]
\end{twierdzenie}
\begin{proof}
Ustalmy $(a,d)\in A\times D$.
Z definicji składania $(a,d)\in T\circ(S\circ R)$
wtedy i tylko wtedy, gdy istnieje $c\in C$ taki, że
$(a,c)\in S\circ R$ oraz $(c,d)\in T$.
Pierwszy z tych warunków oznacza istnienie $b\in B$
z $(a,b)\in R$ i $(b,c)\in S$. Łącznie otrzymujemy:
istnieją $b\in B$ i $c\in C$ takie, że
\[
(a,b)\in R,\qquad (b,c)\in S,\qquad (c,d)\in T.
\]
To samo zdanie możemy odczytać inaczej: istnieje $b\in B$
taki, że $(a,b)\in R$ oraz istnieje $c\in C$,
dla którego $(b,c)\in S$ i $(c,d)\in T$.
Druga część jest równoważna $(b,d)\in T\circ S$.
Stąd całość jest równoważna
$(a,d)\in(T\circ S)\circ R$.
Ponieważ każdy krok można odwrócić, obie relacje zawierają
te same pary.
\end{proof}

\begin{twierdzenie}[Relacje identycznościowe i odwracanie składania]
\label{thm:relation-identities}
Niech $A,B$ będą zbiorami i $R\subseteq A\times B$.
Wówczas $R\circ\Delta_A=R=\Delta_B\circ R$ oraz
$(R^{-1})^{-1}=R$.
Jeżeli dodatkowo $S\subseteq B\times C$, to
\[
(S\circ R)^{-1}=R^{-1}\circ S^{-1}.
\]
\end{twierdzenie}
\begin{proof}
Dla $(a,b)\in A\times B$ warunek
$(a,b)\in R\circ\Delta_A$ oznacza, że istnieje
$a'\in A$ z $(a,a')\in\Delta_A$ i $(a',b)\in R$.
Pierwszy warunek wymusza $a'=a$, a drugi redukuje się do
$(a,b)\in R$. Dowodzi to pierwszej równości.
Analogicznie przy $\Delta_B\circ R$ pośredni element
$b'$ musi być równy $b$. Dwukrotne odwrócenie pary
$(a,b)$ przywraca tę samą parę, więc $(R^{-1})^{-1}=R$.

Wreszcie $(c,a)\in(S\circ R)^{-1}$
wtedy i tylko wtedy, gdy $(a,c)\in S\circ R$,
czyli gdy dla pewnego $b\in B$ zachodzą
$(a,b)\in R$ i $(b,c)\in S$.
Po odwróceniu współrzędnych są to warunki
$(c,b)\in S^{-1}$ i $(b,a)\in R^{-1}$,
a więc $(c,a)\in R^{-1}\circ S^{-1}$.
Równoważność działa w obu kierunkach.
\end{proof}

\section{Własności relacji na zbiorze}\label{sec:relation-properties}
\begin{definicja}[Sześć podstawowych własności]\label{def:relation-properties}
Niech $A$ będzie zbiorem, a $R\subseteq A\times A$.
Relację $R$ nazywamy:
\begin{enumerate}[label=(\roman*)]
\item \emph{zwrotną}, gdy $\forall a\in A\;aRa$;
\item \emph{przeciwzwrotną}, gdy $\forall a\in A\;\neg(aRa)$;
\item \emph{symetryczną}, gdy $\forall a,b\in A\;(aRb\Rightarrow bRa)$;
\item \emph{antysymetryczną}, gdy
$\forall a,b\in A\;((aRb\land bRa)\Rightarrow a=b)$;
\item \emph{asymetryczną}, gdy
$\forall a,b\in A\;(aRb\Rightarrow\neg(bRa))$;
\item \emph{przechodnią}, gdy
$\forall a,b,c\in A\;((aRb\land bRc)\Rightarrow aRc)$.
\end{enumerate}
\end{definicja}

\begin{uwaga}[Antysymetria nie znaczy braku symetrii]
Relacja $\le$ na $\R$ jest antysymetryczna, ponieważ
$x\le y$ i $y\le x$ implikują $x=y$. Jest też zwrotna:
dla każdej liczby $x$ zachodzi $x\le x$. W szczególności,
gdy $x=y$, oba warunki z definicji antysymetrii są spełnione.
Antysymetria dopuszcza zatem pary $(x,x)$, natomiast
asymetria wyklucza je całkowicie.
\end{uwaga}

\begin{twierdzenie}[Asymetria i podstawowe własności]
\label{thm:asymmetric}
Niech $A$ będzie zbiorem, a $R\subseteq A\times A$.
Wówczas $R$ jest asymetryczna wtedy i tylko wtedy, gdy
jest jednocześnie przeciwzwrotna i antysymetryczna.
Ponadto $R$ jest symetryczna wtedy i tylko wtedy,
gdy $R=R^{-1}$, a przechodnia wtedy i tylko wtedy,
gdy $R\circ R\subseteq R$.
\end{twierdzenie}
\begin{proof}
Załóżmy asymetrię. Gdyby istniał $a\in A$ z $aRa$,
to podstawiając $b=a$ do definicji asymetrii otrzymalibyśmy
$\neg(aRa)$, sprzeczność. Stąd przeciwzwrotność.
Jeśli $aRb$ i $bRa$, to z pierwszego warunku i asymetrii
wynika $\neg(bRa)$, także sprzeczność; sytuacja ta
nigdy więc nie występuje. W szczególności implikacja
$(aRb\land bRa)\Rightarrow a=b$ jest zawsze prawdziwa,
czyli $R$ jest antysymetryczna.

Odwrotnie, załóżmy przeciwzwrotność i antysymetrię.
Weźmy $a,b\in A$ z $aRb$. Jeśli również $bRa$,
antysymetria dałaby $a=b$, a wówczas $aRa$,
co przeczy przeciwzwrotności. Zatem $\neg(bRa)$,
czyli $R$ jest asymetryczna.

Dla symetrii zauważmy, że $R\subseteq R^{-1}$ oznacza,
że z $(a,b)\in R$ wynika $(a,b)\in R^{-1}$,
czyli $(b,a)\in R$. To dokładnie definicja symetrii.
Jeżeli $R\subseteq R^{-1}$, odwrócenie obu stron
daje $R^{-1}\subseteq R$; otrzymujemy więc $R=R^{-1}$.
Odwrotnie, jeżeli $R=R^{-1}$, to w szczególności
$R\subseteq R^{-1}$, a wyżej pokazaliśmy, że takie zawieranie
jest równoważne symetrii relacji $R$.

Wreszcie para $(a,c)$ należy do $R\circ R$ dokładnie
wtedy, gdy istnieje $b\in A$ z $aRb$ i $bRc$.
Warunek $R\circ R\subseteq R$ mówi zatem, że dla
każdej takiej pary zachodzi $aRc$, co jest definicją
przechodniości.
\end{proof}

\begin{przyklad}[Jedna relacja, wiele własności]
Na $\Z$ relacja $<$ jest przeciwzwrotna, asymetryczna
i przechodnia; nie jest zwrotna ani symetryczna.
Relacja $\le$ jest zwrotna, antysymetryczna i przechodnia.
Relacja „mieć tę samą parzystość” na $\Z$ jest zwrotna,
symetryczna i przechodnia, ale nie jest antysymetryczna,
gdyż $2$ i $4$ są różne, a mają tę samą parzystość.
Na $\Z$ relacja podzielności nie jest antysymetryczna,
gdyż $2\mid(-2)$ i $(-2)\mid2$, ale $2\ne-2$.
\end{przyklad}

\begin{przyklad}[Niezależność przechodniości od innych własności]
Niech $A=\{1,2,3\}$, $R=\{(1,2),(2,3)\}$.
Relacja ta jest asymetryczna, a więc przeciwzwrotna
i antysymetryczna, lecz nie jest przechodnia:
$1R2$ i $2R3$, podczas gdy $1R3$ nie zachodzi.
Relacja pusta jest z kolei symetryczna, asymetryczna
i przechodnia na dowolnym $A$; gdy $A\ne\varnothing$,
nie jest zwrotna. Na $A=\varnothing$ staje się również
zwrotna, ponieważ nie istnieje element, który naruszałby
warunek zwrotności.
\end{przyklad}

\section{Najmniejsze relacje o zadanych własnościach}
\label{sec:relation-closures}
Często początkowo znamy tylko bezpośrednie połączenia między
obiektami. Chcemy jednak rozszerzyć je tak, aby wolno było
traktować połączenia jako osiągalność po dowolnej liczbie kroków.
To podstawowa motywacja domknięcia przechodniego.

\begin{definicja}[Potęgi relacji i ścieżki]
\label{def:relation-powers}
Niech $A$ będzie zbiorem i $R\subseteq A\times A$.
Definiujemy rekurencyjnie
$R^0:=\Delta_A$ i $R^{n+1}:=R\circ R^n$ dla
$n\in\N:=\{0,1,2,\ldots\}$. Para $(a,b)\in R^n$
oznacza istnienie ścieżki o długości $n$ od $a$ do $b$,
tj. elementów $a=a_0,a_1,\ldots,a_n=b$ takich,
że $a_{j-1}Ra_j$ dla $j=1,\ldots,n$.
Dla $n=0$ ścieżka składa się z jednego punktu $a=b$.
\end{definicja}

\begin{lemat}[Ścieżki a potęgi relacji]\label{lem:relation-paths-xlii}
Niech $A$ będzie zbiorem, $R\subseteq A\times A$,
$a,b\in A$ oraz $n\in\N$. Wtedy $(a,b)\in R^n$
wtedy i tylko wtedy, gdy istnieją
$a_0,\ldots,a_n\in A$ takie, że
$a_0=a$, $a_n=b$ oraz $a_{j-1}Ra_j$ dla
każdego $j\in\{1,\ldots,n\}$.
Dla $n=0$ ostatnia lista warunków jest pusta,
ale nadal wymagamy $a_0=a=b$.
\end{lemat}
\begin{proof}
Dowodzimy równoważności przez indukcję po $n$.
Dla $n=0$ z definicji $R^0=\Delta_A$ mamy
$(a,b)\in R^0$ wtedy i tylko wtedy, gdy $a=b$.
To dokładnie warunek istnienia ścieżki długości zero,
składającej się z jednego wierzchołka $a_0=a=b$.

Załóżmy twierdzenie dla ustalonej liczby $n$.
Z definicji $R^{n+1}=R\circ R^n$ para $(a,b)$
należy do $R^{n+1}$ wtedy i tylko wtedy, gdy
istnieje punkt $c\in A$ taki, że $(a,c)\in R^n$
i $cRb$. Z założenia indukcyjnego pierwszy warunek
jest równoważny istnieniu ciągu
$a=a_0,\ldots,a_n=c$ z kolejnymi krawędziami w $R$.
Dołączenie $a_{n+1}=b$ daje ścieżkę długości $n+1$
od $a$ do $b$. Odwrotnie, ze ścieżki
$a=a_0,\ldots,a_n,a_{n+1}=b$ wybieramy $c=a_n$;
jej pierwsze $n$ krawędzi, na mocy założenia
indukcyjnego, świadczy o $(a,c)\in R^n$,
a ostatnia o $cRb$. Definicja składania daje
$(a,b)\in R^{n+1}$. To zamyka obie implikacje.
\end{proof}


\begin{lemat}[Potęgi relacji i łączenie ścieżek]
\label{lem:relation-powers}
Niech $A$ będzie zbiorem, $R\subseteq A\times A$
i $m,n\in\N$. Wówczas zachodzi $R^{m+n}=R^m\circ R^n$.
\end{lemat}
\begin{proof}
Prowadzimy indukcję po $m$. Dla $m=0$ prawa strona
jest równa $\Delta_A\circ R^n=R^n$ na mocy
Twierdzenia~\ref{thm:relation-identities}; lewa strona
to $R^{0+n}=R^n$. Załóżmy równość dla ustalonego $m$
i każdego $n$. Wówczas, korzystając z definicji potęgi
i łączności składania
(Twierdzenie~\ref{thm:relation-assoc}), mamy
\[
R^{(m+1)+n}=R\circ R^{m+n}
=R\circ(R^m\circ R^n)
=(R\circ R^m)\circ R^n=R^{m+1}\circ R^n.
\]
Krok indukcyjny kończy dowód.
\end{proof}

\begin{definicja}[Domknięcia zwrotne, symetryczne, przechodnie]
\label{def:relation-closure}
Niech $A$ będzie zbiorem i $R\subseteq A\times A$.
Definiujemy
\[
R^{\mathrm{ref}}:=R\cup\Delta_A,\qquad
R^{\mathrm{sym}}:=R\cup R^{-1},\qquad
R^+:=\bigcup_{n\ge1}R^n,\qquad
R^*:=\bigcup_{n\ge0}R^n.
\]
Relacje $R^+$ oraz $R^*$ nazywamy odpowiednio
\emph{domknięciem przechodnim} i
\emph{domknięciem zwrotno-przechodnim} relacji $R$.
\end{definicja}

\begin{twierdzenie}[Minimalność domknięć]\label{thm:relation-closure}
Niech $A$ będzie zbiorem i $R\subseteq A\times A$.
Wtedy $R^{\mathrm{ref}}$ jest najmniejszą (względem zawierania)
relacją zwrotną na $A$, która zawiera $R$,
a $R^{\mathrm{sym}}$ jest najmniejszą relacją symetryczną
na $A$, która zawiera $R$. Ponadto $R^+$ jest najmniejszą
relacją przechodnią na $A$ zawierającą $R$,
a $R^*$ jest najmniejszą relacją zwrotną i przechodnią
na $A$ zawierającą $R$.
\end{twierdzenie}
\begin{proof}
Relacja $R\cup\Delta_A$ zawiera $R$ i wszystkie pary
$(a,a)$, więc jest zwrotna. Jeżeli zwrotna relacja
$S\subseteq A\times A$ zawiera $R$, musi zawierać
każdą parę przekątnej $\Delta_A$, a więc również
$R\cup\Delta_A$. To dowodzi minimalności
$R^{\mathrm{ref}}$.

Jeśli $(a,b)\in R\cup R^{-1}$, to
$(a,b)\in R$ lub $(b,a)\in R$. W pierwszym przypadku
$(b,a)\in R^{-1}$, a w drugim $(b,a)\in R$.
Zatem $R\cup R^{-1}$ jest symetryczna. Jeżeli symetryczna
relacja $S$ zawiera $R$, to zawiera także $R^{-1}$,
więc zawiera sumę $R\cup R^{-1}$.

Przejdźmy do $R^+$. Ponieważ $R^1=R$, mamy
$R\subseteq R^+$. Jeśli $(a,b)\in R^+$ i $(b,c)\in R^+$,
wybierzmy dodatnie liczby naturalne $m,n$ takie, że
$(a,b)\in R^m$ i $(b,c)\in R^n$.
Z definicji składania $(a,c)\in R^n\circ R^m$.
Na mocy Lematu~\ref{lem:relation-powers} prawa strona
jest równa $R^{m+n}$. Ponieważ $m+n\ge1$,
mamy $(a,c)\in R^+$. Zatem $R^+$ jest przechodnia.
Niech teraz $S$ będzie dowolną przechodnią relacją
zawierającą $R$. Pokażemy indukcyjnie, że
$R^n\subseteq S$ dla wszystkich $n\ge1$.
Dla $n=1$ jest to założenie $R\subseteq S$.
Jeśli $R^n\subseteq S$, to
$R^{n+1}=R\circ R^n\subseteq S\circ S\subseteq S$:
pierwsze zawieranie wynika z Lematu~\ref{lem:composition-monotone},
a drugie z przechodniości $S$ (zob.
Twierdzenie~\ref{thm:asymmetric}).
Stąd każdy składnik sumy $R^+$ zawiera się w $S$,
więc $R^+\subseteq S$.

Wreszcie $R^*=\Delta_A\cup R^+$.
Jest zwrotna, gdyż zawiera $\Delta_A$.
Jest też przechodnia: przy złożeniu dwóch par z
$R^*$ możemy dobrać $m,n\ge0$ i użyć
$R^n\circ R^m=R^{m+n}\subseteq R^*$.
Jeśli $S$ jest zwrotna i przechodnia oraz zawiera $R$,
to zawiera $R^+$ z poprzedniego argumentu i
$\Delta_A$ ze zwrotności, a zatem
$R^*=R^+\cup\Delta_A\subseteq S$.
\end{proof}

\begin{przyklad}[Zależność w grafie]
Niech $A=\{a,b,c,d\}$, $R=\{(a,b),(b,c),(c,d)\}$.
W $R^+$ pojawią się między innymi pary $(a,c)$,
$(b,d)$ oraz $(a,d)$, których nie ma w $R$.
Oznacza to, że możemy przejść od $a$ do $d$ po trzech
krawędziach. W $R^*$ dodatkowo obecne będą
$(a,a),(b,b),(c,c),(d,d)$, odpowiadające ścieżkom
o długości zero.
\end{przyklad}

\begin{uwaga}[Zapowiedź relacji równoważności i porządków]
Relacja jest \emph{relacją równoważności}, gdy jest
zwrotna, symetryczna i przechodnia; jej klasy i ilorazy
zbadamy w rozdziale 5. Relacja jest \emph{porządkiem
częściowym}, gdy jest zwrotna, antysymetryczna i przechodnia;
porządki, kresy i diagramy Hassego będą tematem rozdziału 6.
Niniejszy rozdział dostarcza narzędzi potrzebnych do obu teorii.
\end{uwaga}

\begin{twierdzenie}[Najmniejsza relacja równoważności generowana przez relację]
\label{thm:equiv-generated-prelim}
Niech $A$ będzie zbiorem, a $R\subseteq A\times A$ relacją.
Wtedy relacja $E:=(R\cup R^{-1})^*$ jest zwrotna, symetryczna
i przechodnia. Jest też najmniejszą relacją o tych trzech
własnościach zawierającą $R$.
\end{twierdzenie}
\begin{proof}
Oznaczmy $S=R\cup R^{-1}$. Z definicji $S$ jest symetryczna:
jeśli $(a,b)\in S$, to $(b,a)\in S$. Relacja $S^*$ zawiera $S$
i jest zwrotna i przechodnia z
Twierdzenia~\ref{thm:relation-closure}. Aby wykazać symetrię $S^*$,
weźmy $(a,b)\in S^*$. Istnieje wtedy ciąg
$a=a_0,a_1,\ldots,a_n=b$, w którym każda kolejna para
$(a_{j-1},a_j)$ należy do $S$ (dla $n=0$ mamy $a=b$).
Dzięki symetrii $S$ wszystkie pary
$(a_j,a_{j-1})$ też należą do $S$. Odwracając kolejność
punktów, otrzymujemy ścieżkę od $b$ do $a$, zatem
$(b,a)\in S^*$. To dowodzi symetrii.

Niech teraz $T$ będzie dowolną relacją zwrotną, symetryczną
i przechodnią zawierającą $R$. Ze zwrotności i symetrii wynika
$R^{-1}\subseteq T$, zatem $S\subseteq T$.
Z minimalności domknięcia zwrotno-przechodniego
(\ref{thm:relation-closure}) dostajemy $S^*\subseteq T$.
Dokładnie to znaczy, że $E$ jest najmniejszą z żądanych relacji.
\end{proof}

\begin{przyklad}[Spójne składowe grafu]
Dla $A=\{1,2,3,4,5\}$ i $R=\{(1,2),(2,3),(4,5)\}$
relacja $(R\cup R^{-1})^*$ łączy każdą parę elementów z
$\{1,2,3\}$ oraz każdą parę z $\{4,5\}$, ale nie łączy
elementów należących do różnych tych grup.
W rozdziale o relacjach równoważności pokażemy,
że te dwie grupy są klasami równoważności.
\end{przyklad}

\begin{center}
\begin{tikzpicture}[>=Stealth,every node/.style={font=\small}]
  \node[circle,draw,minimum size=6.5mm] (v1) at (0,0) {$1$};
  \node[circle,draw,minimum size=6.5mm] (v2) at (1.5,.65) {$2$};
  \node[circle,draw,minimum size=6.5mm] (v3) at (3,0) {$3$};
  \node[circle,draw,minimum size=6.5mm] (v4) at (5.0,.65) {$4$};
  \node[circle,draw,minimum size=6.5mm] (v5) at (6.5,0) {$5$};
  \draw[very thick] (v1)--(v2)--(v3);
  \draw[very thick] (v4)--(v5);
  \node[below=8mm of v2] {$\{1,2,3\}$};
  \node at (5.75,-.9) {$\{4,5\}$};
  \node[font=\small,anchor=north] at (3.3,-1.6)
    {Dwie składowe relacji $\bigl(R\cup R^{-1}\bigr)^*$.};
\end{tikzpicture}
\end{center}


\section{Relacje a macierze zero-jedynkowe}
\label{sec:relation-matrices}
\begin{przyklad}[Konkretne obliczenie składania]
Dla skończonych zbiorów $A=\{a_1,\ldots,a_m\}$ i
$B=\{b_1,\ldots,b_n\}$ relację $R\subseteq A\times B$
można zakodować macierzą $M_R\in\{0,1\}^{m\times n}$:
$(M_R)_{ij}=1$ wtedy i tylko wtedy, gdy $a_iRb_j$.
Jeśli $S\subseteq B\times C$, to macierz składania
ma współczynniki
\[
(M_{S\circ R})_{ik}
=\bigvee_{j=1}^{n}\bigl((M_R)_{ij}\land(M_S)_{jk}\bigr),
\]
gdzie $\lor$ i $\land$ oznaczają operacje logiczne.
Przy $B=\varnothing$ mamy $n=0$ i przyjmujemy,
że alternatywa pustej rodziny ma wartość $0$;
odpowiada to brakowi elementu pośredniego i daje
relację pustą $S\circ R=\varnothing$.
Aby współczynnik był równy $1$, musi istnieć przynajmniej
jedno $j$ takie, że oba czynniki mają wartość $1$;
jest to dokładnie warunek istnienia punktu pośredniego
$b_j$ w definicji składania.

Na przykład dla $A=B=C=\{1,2\}$ i
$R=\{(1,1),(1,2)\}$ oraz $S=\{(1,2),(2,2)\}$ mamy
\[
M_R=\begin{pmatrix}1&1\\0&0\end{pmatrix},\qquad
M_S=\begin{pmatrix}0&1\\0&1\end{pmatrix},\qquad
M_{S\circ R}=\begin{pmatrix}0&1\\0&0\end{pmatrix}.
\]
Takie macierze i iloczyn boolowski stanowią podstawę
algorytmów obliczania osiągalności w skończonych grafach.
\end{przyklad}

\begin{twierdzenie}[Macierz relacji odwrotnej]\label{thm:lxiv-relation-matrix-transpose}
Niech $A=\{a_1,\ldots,a_m\}$ i $B=\{b_1,\ldots,b_n\}$
będą skończonymi zbiorami z ustalonymi wyliczeniami,
a $R\subseteq A\times B$ relacją. Wówczas macierz relacji
odwrotnej $R^{-1}\subseteq B\times A$, przy tych samych
wyliczeniach, jest transpozycją macierzy $R$:
$M_{R^{-1}}=M_R^{\mathsf T}$. Równość obowiązuje także
wtedy, gdy jeden z tych zbiorów jest pusty.
\end{twierdzenie}
\begin{proof}
Dla dowolnych $1\le j\le n$ i $1\le i\le m$,
z definicji macierzy relacji mamy
$(M_{R^{-1}})_{ji}=1$ wtedy i tylko wtedy, gdy
$b_jR^{-1}a_i$. Z definicji relacji odwrotnej jest to
równoważne $a_iRb_j$, a więc $(M_R)_{ij}=1$.
Ponieważ oba współczynniki należą do $\{0,1\}$,
są równe. Jeżeli $m=0$ lub $n=0$, nie istnieje
żadna para indeksów do sprawdzenia, a obie macierze
mają odpowiednio rozmiary $n\times m$ i $m\times n$;
transpozycja nadal ma wymagany rozmiar i wpisy.
\end{proof}
\begin{przyklad}[Transpozycja a kierunek strzałek]\label{ex:lxiv-relation-matrix-transpose}
Niech $A=\{1,2\}$, $B=\{u,v,w\}$ i
$R=\{(1,u),(1,w),(2,v)\}$. Wówczas
\[
M_R=\begin{pmatrix}1&0&1\\0&1&0\end{pmatrix},\qquad
M_{R^{-1}}=\begin{pmatrix}1&0\\0&1\\1&0\end{pmatrix}.
\]
Na przykład wpis $(3,1)$ drugiej macierzy wynosi $1$,
ponieważ $wR^{-1}1$, czyli $(1,w)\in R$.
\end{przyklad}

\begin{uwaga}[Podsumowanie metody dowodzenia]
Najczęstsza technika w tych rozdziałach to wzięcie
dowolnego elementu $x$ lub pary $(a,b)$ i zastąpienie
przynależności do zbioru rozwinięciem jego definicji.
W szczególności prawa składania relacji są prawami
manipulowania kwantyfikatorami egzystencjalnymi:
„istnieje element pośredni”. W następnych rozdziałach
będziemy coraz częściej korzystać z tej metody.
\end{uwaga}

\chapter{Relacje równoważności i zbiory ilorazowe}\label{ch:equivalence}

\section{Dlaczego sama równość nie wystarcza?}\label{sec:eq-motivation}
Równość $a=b$ jest najbardziej rygorystycznym sposobem uznania dwóch
obiektów za nierozróżnialne. W matematyce często chcemy jednak pomijać
pewien rodzaj różnicy. Na przykład liczby $2$ i $7$ nie są równe,
ale przy dzieleniu przez $5$ dają tę samą resztę. Dwa napisy mogą się
różnić kolejnością liter, lecz mieć tę samą długość. Dwa ułamki
$1/2$ i $2/4$ są różnymi parami liczb, chociaż opisują tę samą
liczbę wymierną. Chcemy opisać wspólny mechanizm takich utożsamień.

W rozdziale~\ref{ch:relations} poznaliśmy zwrotność, symetrię i
przechodniość. Łącznie są one dokładnie własnościami pozwalającymi
rozsądnie grupować obiekty, które w wybranym znaczeniu uważamy za
równoważne. Najpierw pokażemy to na przykładach, a potem sformalizujemy.

\begin{definicja}[Relacja równoważności]\label{def:equivalence}
Niech $X$ będzie zbiorem, a $E\subseteq X\times X$ relacją.
Nazywamy $E$ \emph{relacją równoważności na $X$}, jeśli
\begin{enumerate}[label=(\roman*)]
\item dla każdego $x\in X$ zachodzi $xEx$ (zwrotność);
\item dla wszystkich $x,y\in X$, z $xEy$ wynika $yEx$ (symetria);
\item dla wszystkich $x,y,z\in X$, z $xEy$ i $yEz$ wynika $xEz$
(przechodniość).
\end{enumerate}
Najczęściej zapisujemy $x\sim y$ zamiast $xEy$.
\end{definicja}

\begin{przyklad}[Równość i relacja całkowita]
Na dowolnym zbiorze $X$ relacja $x\sim y\iff x=y$ jest
równoważnością: każdy element jest równy samemu sobie, równość
można odwrócić i składać. Relacja $x\sim y$ dla \emph{każdych}
$x,y\in X$ także jest równoważnością. W pierwszym przypadku
niczego nie utożsamiamy, w drugim utożsamiamy wszystkie elementy.
Obie relacje są równoważnościami również dla $X=\varnothing$.
\end{przyklad}

\begin{przyklad}[Te same reszty z dzielenia]\label{ex:eq-congruence}
Ustalmy liczbę naturalną $m\ge1$. Na $\Z$ określmy
\[
 a\sim_m b\quad\Longleftrightarrow\quad
 m\mid(a-b),
\]
gdzie $m\mid t$ znaczy, że istnieje $k\in\Z$ z $t=mk$.
Zwrotność wynika z $a-a=m\cdot0$. Jeśli $a-b=mk$, to
$b-a=m(-k)$, zatem zachodzi symetria. Jeśli
$a-b=mk$ oraz $b-c=m\ell$, to dodając równości dostajemy
$a-c=m(k+\ell)$; to dowodzi przechodniości.
Dla $m=5$ liczby $2,7,-3,12$ są równoważne. Dla $m=1$
wszystkie liczby całkowite są równoważne.
\end{przyklad}

\begin{przyklad}[Bliskość nie jest równoważnością]
Na $\R$ ustalmy $xRy\iff |x-y|<1$. Relacja jest zwrotna
i symetryczna, ale nieprzechodnia: $0R\frac34$ i
$\frac34R\frac32$, natomiast $0\not R\frac32$.
Sama intuicja podobieństwa nie wystarcza, jeśli nie można
przenosić jej przez pośrednie obiekty.
\end{przyklad}

\section{Klasy równoważności i ich podstawowe własności}
\label{sec:eq-classes}
\begin{definicja}[Klasa równoważności]\label{def:eq-class}
Niech $X$ będzie zbiorem, a $\sim$ relacją równoważności na $X$.
Dla $x\in X$ \emph{klasą równoważności elementu $x$} nazywamy
\[
 [x]_{\sim}:=\{y\in X:y\sim x\}.
\]
Gdy relacja jest jasna, piszemy krótko $[x]$. Zbiór
wszystkich klas równoważności oznaczamy
\[
 X/{\sim}:=\{[x]_{\sim}:x\in X\}
\]
i nazywamy \emph{zbiorem ilorazowym} zbioru $X$ przez $\sim$.
Każdy element $X/{\sim}$ jest sam w sobie podzbiorem $X$.
\end{definicja}

\begin{uwaga}[Istnienie zbioru ilorazowego, także dla pustego $X$]
Każda klasa $[x]_{\sim}$ jest podzbiorem $X$, więc
$X/{\sim}\subseteq\mathcal P(X)$. Aby otrzymać iloraz jako
zbiór w ZF, można wyciąć z już istniejącego zbioru
$\mathcal P(X)$ te i tylko te podzbiory $C$, dla których
istnieje $x\in X$ spełniający
$\forall y\in X\,(y\in C\iff y\sim x)$.
To jest zastosowanie schematu wyróżniania, a nie
nieograniczonego tworzenia zbioru z warunku.
Jeśli $X=\varnothing$, jedyna relacja na $X$ jest pusta,
nie istnieje żaden $x\in X$, a więc $X/{\sim}=\varnothing$.
W przypadku niepustego $X$ każda klasa jest niepusta,
ale nie oznacza to, że sam zbiór ilorazowy ma jeden element:
liczba klas zależy od przyjętej relacji.
\end{uwaga}



\begin{uwaga}[Element, singleton i klasa]
Dla $x\in X$ mamy trzy potencjalnie różne obiekty: $x$,
$\{x\}$ oraz $[x]$. Zwrotność zapewnia $x\in[x]$,
a więc $[x]\ne\varnothing$, lecz nie musi zachodzić
$[x]=\{x\}$. Na przykład w relacji kongruencji modulo $5$
klasa $[2]$ zawiera nieskończenie wiele liczb całkowitych.
\end{uwaga}

\begin{twierdzenie}[Kiedy dwie klasy są równe?]\label{thm:class-equality}
Niech $X$ będzie zbiorem, a $\sim$ relacją równoważności na $X$.
Dla dowolnych $x,y\in X$ następujące warunki są równoważne:
\begin{enumerate}[label=(\roman*)]
\item $x\sim y$;
\item $[x]=[y]$;
\item $[x]\cap[y]\ne\varnothing$.
\end{enumerate}
\end{twierdzenie}
\begin{proof}
Najpierw załóżmy $x\sim y$. Aby wykazać $[x]\subseteq[y]$,
wybierzmy dowolny $z\in[x]$. Definicja klasy daje $z\sim x$.
Z przechodniości $z\sim x$ i $x\sim y$ wynika $z\sim y$,
a więc $z\in[y]$. W drugą stronę, jeśli $z\in[y]$, to
$z\sim y$; z symetrii $x\sim y$ mamy $y\sim x$.
Przechodniość daje $z\sim x$, zatem $z\in[x]$.
Pokazaliśmy oba zawierania, więc $[x]=[y]$.

Jeżeli $[x]=[y]$, to dzięki zwrotności $x\in[x]$,
a zatem $x\in[x]\cap[y]$. Przecięcie jest niepuste.

Wreszcie załóżmy $[x]\cap[y]\ne\varnothing$.
Istnieje wtedy $z\in X$ należący do obu klas.
Oznacza to $z\sim x$ i $z\sim y$.
Z symetrii pierwszej relacji wynika $x\sim z$.
Łącząc $x\sim z$ z $z\sim y$, otrzymujemy
$x\sim y$ przez przechodniość. Zamknęliśmy ciąg implikacji.
\end{proof}

\begin{wniosek}[Klasy nie przecinają się częściowo]\label{cor:eq-disjoint}
Niech $X$ będzie zbiorem, a $\sim$ relacją równoważności na $X$.
Dla dowolnych $x,y\in X$ klasy $[x]$ i $[y]$
są albo rozłączne, albo równe. Ponadto
\[
 X=\bigcup_{x\in X}[x].
\]
\end{wniosek}
\begin{proof}
Jeśli przecięcie klas jest niepuste, to na mocy
Twierdzenia~\ref{thm:class-equality} klasy są równe.
Daje to pierwszą tezę. Każda klasa jest podzbiorem $X$,
więc ich suma zawiera się w $X$. Odwrotnie, dowolny
$x\in X$ należy do $[x]$ przez zwrotność. Zatem
$x$ należy do sumy wszystkich klas. Stąd równość.
Gdy $X=\varnothing$, prawa strona jest pustą sumą i nadal
otrzymujemy równość.
\end{proof}

\begin{przyklad}[Kongruencje modulo $3$]
Na $\Z$ przy relacji $\sim_3$ mamy dokładnie trzy różne klasy:
\[
 [0]=\{3k:k\in\Z\},\quad
 [1]=\{3k+1:k\in\Z\},\quad
 [2]=\{3k+2:k\in\Z\}.
\]
Każda liczba całkowita należy do dokładnie jednej z nich:
dzieląc ją przez $3$ otrzymujemy dokładnie jedną z reszt
$0,1,2$. Klasy $[3]$ i $[0]$ nie są dwiema różnymi klasami;
ponieważ $3\sim_3 0$, są tym samym podzbiorem $\Z$.
\end{przyklad}

\begin{przyklad}[Przykład geometryczny]
Na płaszczyźnie $\R^2$ ustalmy
$(x,y)\sim(u,v)\iff y=v$. Zwrotność, symetria i
przechodniość wynikają z tych własności równości drugiej
współrzędnej. Klasa punktu $(x_0,y_0)$ ma postać
$\{(x,y_0):x\in\R\}$, czyli jest prostą poziomą.
Zbiór ilorazowy składa się z \emph{prostych poziomych},
a nie z pojedynczych punktów płaszczyzny.
\end{przyklad}

\begin{center}
\begin{tikzpicture}[x=0.82cm,y=0.82cm]
  \begin{scope}[shift={(0,0)}]
    \draw[->] (-.2,0)--(4.4,0) node[right] {$x$};
    \draw[->] (0,-.2)--(0,3.3) node[above] {$y$};
    \foreach \h in {.65,1.45,2.25} {
      \draw[blue!75!black,thick] (.18,\h)--(4.0,\h);
    }
    \filldraw (1.1,1.45) circle (1.7pt) node[above left] {$p$};
    \filldraw (3.1,1.45) circle (1.7pt) node[above] {$q$};
    \node[font=\small,anchor=north] at (2.0,-.5) {$y=v$: klasy poziome};
  \end{scope}
  \begin{scope}[shift={(6.3,0)}]
    \draw[->] (-.2,0)--(4.4,0) node[right] {$x$};
    \draw[->] (0,-.2)--(0,3.3) node[above] {$y$};
    \begin{scope}
      \clip (.18,.1) rectangle (4,3);
      \foreach \c in {-2,-1,0,1,2} {
        \draw[red!75!black,thick] (-.8,{-.8+\c})--(5,{5+\c});
      }
    \end{scope}
    \node[font=\small,anchor=north] at (2,-.5) {$x-y=u-v$: klasy ukośne};
  \end{scope}
\end{tikzpicture}

{\small Po lewej klasy relacji $(x,y)\sim(u,v)\iff y=v$;
po prawej klasy \emph{innej} relacji
$(x,y)\approx(u,v)\iff x-y=u-v$.
Na każdym panelu różne klasy nie mają punktów wspólnych.}
\end{center}


\begin{uwaga}[Dlaczego czerwone i niebieskie proste opisują różne relacje?]
Na rysunku powyżej dwie poziome proste nie przecinają się,
ponieważ stanowią dwie różne klasy \emph{tej samej} relacji
„mieć jednakową drugą współrzędną”. Czerwone proste
są klasami innej relacji, „mieć jednakową różnicę współrzędnych”.
Czerwona i niebieska prosta mogą się przeciąć i nie narusza to
Twierdzenia~\ref{thm:class-equality}: twierdzenie o rozłączności
różnych klas dotyczy klas \emph{jednej ustalonej} relacji.
\end{uwaga}

\section{Podziały zbioru i równoważności}\label{sec:eq-partition}
\begin{definicja}[Podział zbioru]\label{def:partition}
Niech $X$ będzie zbiorem. \emph{Podziałem $X$} nazywamy
rodzinę $\mathcal D$ podzbiorów $X$ spełniającą:
\begin{enumerate}[label=(\roman*)]
\item każdy $D\in\mathcal D$ jest niepusty;
\item jeśli $D,E\in\mathcal D$ i $D\ne E$, to
$D\cap E=\varnothing$;
\item $\bigcup_{D\in\mathcal D}D=X$.
\end{enumerate}
Dla $X=\varnothing$ rodzinę pustą uznajemy za podział;
co więcej, jest ona jedynym podziałem pustego zbioru: każdy
blok musiałby być niepustym podzbiorem $\varnothing$, co jest niemożliwe.
Elementy $\mathcal D$ nazywamy \emph{blokami podziału}.
\end{definicja}

\begin{twierdzenie}[Równoważności i podziały]\label{thm:eq-partition}
Niech $X$ będzie dowolnym zbiorem.
\begin{enumerate}[label=(\roman*)]
\item Każda relacja równoważności $\sim$ na $X$ wyznacza
podział $\mathcal D_{\sim}:=X/{\sim}$.
\item Każdy podział $\mathcal D$ zbioru $X$ wyznacza relację
$\sim_{\mathcal D}$, określoną przez
\[
 x\sim_{\mathcal D}y\iff
 \text{istnieje }D\in\mathcal D\text{ takie, że }x,y\in D.
\]
Dla wszystkich $x,y\in X$ formuła ta określa relację
równoważności, której klasami są dokładnie bloki $\mathcal D$.
\item Przejścia opisane w (i) i (ii) są wzajemnie odwrotne:
wyznaczają jednoznaczną odpowiedniość między wszystkimi
relacjami równoważności na $X$ a wszystkimi podziałami $X$.
\end{enumerate}
\end{twierdzenie}
\begin{proof}
(i) Z definicji klasy $[x]\subseteq X$, a zwrotność
gwarantuje $x\in[x]$, czyli klasa jest niepusta.
Wniosek~\ref{cor:eq-disjoint} mówi, że różne klasy są rozłączne
i że ich suma jest całym $X$. To są dokładnie trzy warunki
w definicji podziału.

(ii) Ponieważ $\mathcal D$ pokrywa $X$, dla każdego
$x\in X$ istnieje blok $D$ zawierający $x$.
Wtedy $x,x\in D$, więc $x\sim_{\mathcal D}x$.
Jeżeli $x,y\in D$, to także $y,x\in D$, więc zachodzi
symetria. Dla przechodniości załóżmy
$x\sim_{\mathcal D}y$ oraz $y\sim_{\mathcal D}z$.
Istnieją bloki $D,E\in\mathcal D$ takie, że
$x,y\in D$ i $y,z\in E$. W szczególności $y\in D\cap E$.
Przecięcie nie jest puste, więc z warunku rozłączności
różnych bloków musi zachodzić $D=E$. Wtedy $x,z\in D$,
zatem $x\sim_{\mathcal D}z$.

Niech teraz $D\in\mathcal D$ i $x\in D$ (taki $x$
istnieje, bo bloki są niepuste). Jeśli $y\in D$, to
$x\sim_{\mathcal D}y$, a więc $y\in[x]$.
Odwrotnie, jeżeli $y\in[x]$, to $y$ i $x$ należą do
pewnego wspólnego bloku $E$. Ponieważ $x\in D\cap E$,
wynika $E=D$, a zatem $y\in D$. Stąd $[x]=D$.
Każdy blok jest więc klasą. Odwrotnie, dla dowolnego
$x\in X$ istnieje blok $D\in\mathcal D$ zawierający $x$
(z warunku pokrycia), a pokazany argument daje $[x]=D$.
Zatem każda klasa jest blokiem, także gdy $X=\varnothing$
(wówczas oba zbiory klas i bloków są puste).

(iii) Jeśli zaczęliśmy od relacji $\sim$, to dla
$x,y\in X$ warunek przynależności do wspólnej klasy
jest równoważny $x\sim y$ na mocy
Twierdzenia~\ref{thm:class-equality}.
Odzyskujemy więc tę samą relację. Jeśli zaczęliśmy
od podziału $\mathcal D$, to (ii) pokazało, że klasy
wyznaczonej relacji są dokładnie blokami $\mathcal D$.
Odzyskujemy więc dokładnie ten sam podział.
\end{proof}

\begin{przyklad}[Wszystkie podziały trójelementowego zbioru]
\label{ex:eq-partitions-three-xliii}
Niech $X=\{a,b,c\}$. Każdy podział ma dokładnie jeden,
dwa albo trzy bloki. Przy jednym bloku istnieje tylko
$\{\{a,b,c\}\}$. Przy trzech blokach istnieje tylko
$\{\{a\},\{b\},\{c\}\}$. Przy dwóch blokach jeden ma
jeden element, a drugi dwa; otrzymujemy dokładnie
\begin{gather*}\{\{a\},\{b,c\}\},\quad \{\{b\},\{a,c\}\},\quad\{\{c\},\{a,b\}\}.\end{gather*}
Zatem wszystkich podziałów jest pięć, a na mocy
Twierdzenia~\ref{thm:eq-partition} także relacji
równoważności na $X$ jest pięć. Nie wolno tu liczyć
uporządkowanych list bloków: zamiana kolejności bloków
nie tworzy nowego podziału.
\end{przyklad}



\begin{przyklad}[Podziały skończonego zbioru]
Dla $X=\{a,b,c\}$ jednymi z możliwych podziałów są
$\{\{a\},\{b\},\{c\}\}$,
$\{\{a,b\},\{c\}\}$ oraz $\{\{a,b,c\}\}$.
Pierwszy odpowiada równości, trzeci relacji pełnej,
a drugi utożsamieniu wyłącznie elementów $a,b$.
Zwróćmy uwagę na różnicę pomiędzy zbiorem $X$ a
zbiorem jego bloków: w drugim przykładzie mamy dwa
bloki, choć zbiór $X$ ma trzy elementy.
\end{przyklad}

\section{Uogólnianie i uszczegóławianie identyfikacji}\label{sec:eq-refinement}
\begin{definicja}[Relacja drobniejsza i grubsza]
Niech $X$ będzie zbiorem. Jeżeli $E,F$ są relacjami
równoważności na $X$, to $E$ nazywamy \emph{drobniejszą} od $F$,
gdy $E\subseteq F$ jako podzbiory $X\times X$.
Oznacza to, że $E$ utożsamia nie więcej par niż $F$.
W tej samej sytuacji relację $F$ nazywamy \emph{grubszą}
od $E$.
\end{definicja}

\begin{twierdzenie}[Drobniejsza relacja ma mniejsze klasy]
\label{thm:eq-refinement}
Niech $X$ będzie zbiorem, a $E,F$ relacjami
równoważności na $X$. Wtedy $E\subseteq F$ wtedy i tylko wtedy,
gdy $[x]_E\subseteq[x]_F$ dla każdego $x\in X$.
\end{twierdzenie}
\begin{proof}
Załóżmy $E\subseteq F$. Jeśli $y\in[x]_E$, to $yEx$,
a z zawierania relacji wynika $yFx$. Zatem
$y\in[x]_F$. Udowodniliśmy zawieranie klas.

Odwrotnie, załóżmy zawieranie klas dla wszystkich $x$.
Jeśli $(x,y)\in E$, to $xEy$ i stąd
$x\in[y]_E\subseteq[y]_F$. Definicja klasy daje $xFy$,
czyli $(x,y)\in F$. To dowodzi $E\subseteq F$.
\end{proof}

\begin{przyklad}[Reszty modulo $4$ i modulo $2$]
Na $\Z$ kongruencja modulo $4$ jest drobniejsza niż
kongruencja modulo $2$: podzielność przez $4$ implikuje
podzielność przez $2$. Klasa $[1]_{4}$ składa się z liczb
$\ldots,-7,-3,1,5,9,\ldots$, a klasa $[1]_{2}$
z wszystkich liczb nieparzystych. Zatem
$[1]_4\subsetneq[1]_2$.
\end{przyklad}

\begin{twierdzenie}[Przecięcie rodzin równoważności]
\label{thm:eq-intersection}
Niech $X$ będzie zbiorem, $I$ niepustym zbiorem indeksów,
a dla każdego $i\in I$ niech $E_i$ będzie relacją
równoważności na $X$.
Wtedy $E:=\bigcap_{i\in I}E_i$ jest relacją
równoważności na $X$.
\end{twierdzenie}
\begin{proof}
Dla każdego $x\in X$ i każdego $i\in I$ zachodzi
$(x,x)\in E_i$, więc $(x,x)$ należy do przecięcia.
Jeżeli $(x,y)\in E$, to $(x,y)\in E_i$ dla wszystkich $i$.
Symetria każdej $E_i$ daje $(y,x)\in E_i$ dla wszystkich
$i$, a więc $(y,x)\in E$. Jeżeli $(x,y),(y,z)\in E$,
to obie pary należą do każdej $E_i$; przechodniość
wszystkich $E_i$ daje $(x,z)\in E_i$ dla każdego $i$,
czyli $(x,z)\in E$. Wszystkie trzy własności zostały
wykazane. Dla pustego $I$ można przyjąć za przecięcie
całe $X\times X$, które także jest równoważnością.
\end{proof}

\begin{twierdzenie}[Równoważność generowana przez relację]
\label{thm:eq-generated}
Niech $X$ będzie zbiorem, a $R\subseteq X\times X$
dowolną relacją. Istnieje dokładnie jedna relacja równoważności
na $X$ zawierająca $R$ i \emph{najmniejsza względem zawierania}.
Jest nią
\[
 \langle R\rangle_{\mathrm{eq}}=(R\cup R^{-1})^*.
\]
W szczególności dla wszystkich $x,y\in X$ zachodzi
$x\langle R\rangle_{\mathrm{eq}}y$ wtedy i tylko wtedy,
gdy dla pewnego $n\in\{0,1,2,\ldots\}$ istnieje łańcuch
$x=x_0,\ldots,x_n=y$, w którym dla każdego $j=1,\ldots,n$
zachodzi $x_{j-1}Rx_j$ lub $x_jRx_{j-1}$.
\end{twierdzenie}
\begin{proof}
W Twierdzeniu~\ref{thm:equiv-generated-prelim} już wykazaliśmy,
że $(R\cup R^{-1})^*$ jest zwrotna, symetryczna,
przechodnia i najmniejsza spośród wszystkich relacji
o tych własnościach zawierających $R$.
Dowód tamtego twierdzenia podawał również równoważne
kryterium ścieżkowe: każda para w $S^*$, gdzie
$S=R\cup R^{-1}$, jest połączona ścieżką z krawędziami
z $S$, a krawędź z $S$ jest krawędzią z $R$ w jednym
z dwóch kierunków. W drugą stronę dowolna taka ścieżka
wyznacza parę należącą do pewnego $S^n\subseteq S^*$.
Jedyność najmniejszej relacji wynika z kryterium dwóch zawierań:
jeśli $E$ i $F$ byłyby obie najmniejsze, to $E\subseteq F$
i $F\subseteq E$, więc $E=F$.
\end{proof}

\begin{twierdzenie}[Przecięcie i najmniejsze wspólne rozszerzenie]
\label{thm:eq-meet-join-xliii}
Niech $E,F$ będą relacjami równoważności na zbiorze $X$.
Relacja $E\cap F$ jest największą względem zawierania
relacją równoważności zawartą jednocześnie w $E$ i $F$.
Relacja $\langle E\cup F\rangle_{\mathrm{eq}}$
jest najmniejszą relacją równoważności zawierającą jednocześnie
$E$ i $F$. Sama suma $E\cup F$ nie musi być przechodnia.
\end{twierdzenie}
\begin{proof}
Z Twierdzenia~\ref{thm:eq-intersection} wynika, że $E\cap F$
jest równoważnością. Z definicji przecięcia zawiera się w
$E$ oraz w $F$. Jeżeli równoważność $G$ spełnia
$G\subseteq E$ i $G\subseteq F$, to każda para z $G$
należy do obu tych relacji, czyli do $E\cap F$.
To uzasadnia największość, nie tylko samo istnienie.

Niech $H=\langle E\cup F\rangle_{\mathrm{eq}}$.
Na mocy Twierdzenia~\ref{thm:eq-generated} relacja $H$
jest równoważnością i zawiera $E\cup F$, a więc zarówno
$E$, jak i $F$. Jeśli inna równoważność $K$ zawiera obie
relacje, to $E\cup F\subseteq K$. Minimalność
równoważności generowanej daje $H\subseteq K$.
Ostatnią część tezy, że suma $E\cup F$ nie musi być przechodnia,
uzasadnia kontrprzykład~\ref{ex:eq-union-not-equivalence-xliii}.
\end{proof}

\begin{przyklad}[Suma dwóch równoważności może nie być równoważnością]
\label{ex:eq-union-not-equivalence-xliii}
Niech $X=\{a,b,c\}$. Niech $E$ odpowiada podziałowi
$\{\{a,b\},\{c\}\}$, zaś $F$ podziałowi
$\{\{a\},\{b,c\}\}$. Obie relacje są równoważnościami,
ponieważ powstały z podziałów. W $E\cup F$ zachodzi
$a(E\cup F)b$ (z relacji $E$) i $b(E\cup F)c$
(z relacji $F$), lecz $(a,c)\notin E\cup F$:
$a,c$ leżą w różnych blokach obu podziałów.
Zatem suma nie jest przechodnia. Jej równoważność
generowana utożsamia wszystkie trzy elementy, a przecięcie
$E\cap F$ jest relacją równości $\Delta_X$.
\end{przyklad}



\begin{przyklad}[Utożsamienia generowane przez krawędzie]
Dla $X=\{a,b,c,d,e\}$ i
$R=\{(a,b),(b,c),(d,e)\}$ relacja generowana ma dwie klasy:
$\{a,b,c\}$ i $\{d,e\}$. Nawet jeśli w $R$ nie było pary
$(a,c)$, pojawia się ona przez przechodniość. Nie pojawia
się natomiast $(a,d)$, bo nie ma łańcucha łączącego oba zbiory.
\end{przyklad}

\section{Podzbiory nasycone i konstrukcje na ilorazach}
\label{sec:eq-saturated}
\begin{definicja}[Zbiór nasycony]
Niech $X$ będzie zbiorem, a $\sim$ relacją równoważności na $X$.
Podzbiór $A\subseteq X$ nazywamy \emph{nasyconym}, jeśli
\[
 x\in A,\ x\sim y\quad\Longrightarrow\quad y\in A
 \qquad(x,y\in X).
\]
Inaczej mówiąc, wraz z każdym elementem $A$ zawiera całą jego klasę.
\end{definicja}

\begin{twierdzenie}[Charakteryzacja nasycenia]
\label{thm:eq-saturated}
Niech $X$ będzie zbiorem, $\sim$ relacją równoważności
na $X$, a $A\subseteq X$. Następujące warunki są równoważne:
\begin{enumerate}[label=(\roman*)]
\item $A$ jest nasycony;
\item $A=\bigcup_{x\in A}[x]$;
\item $A$ jest sumą pewnej rodziny klas równoważności.
\end{enumerate}
Ponadto suma, przecięcie i dopełnienie względem $X$
dowolnych zbiorów nasyconych są nasycone.
\end{twierdzenie}
\begin{proof}
(i)$\Rightarrow$(ii). Dla dowolnego $x\in A$ z nasycenia
mamy $[x]\subseteq A$, zatem suma wszystkich tych klas
zawiera się w $A$. Odwrotnie, każdy $x\in A$ należy do
$[x]$ przez zwrotność, więc należy do sumy.

(ii)$\Rightarrow$(iii). Jest to szczególny przypadek sumy
rodziny klas indeksowanej przez elementy $A$.

(iii)$\Rightarrow$(i). Załóżmy, że $A$ jest sumą klas,
$x\in A$ i $x\sim y$. Istnieje klasa $[z]$ zawarta w
przyjętej sumie, do której należy $x$. Na mocy
Twierdzenia~\ref{thm:class-equality} z $x\sim y$ wynika
$[x]=[y]$, a z $x\in[z]$ wynika $[x]=[z]$.
W konsekwencji $y\in[y]=[z]\subseteq A$.

Jeżeli $A_i$ są nasycone, a $x\in\bigcup_i A_i$ i $x\sim y$,
wybieramy $i$ z $x\in A_i$; wówczas $y\in A_i$,
więc $y$ należy do sumy. Jeżeli $x\in\bigcap_i A_i$
i $x\sim y$, to $x\in A_i$ dla każdego $i$, więc
$y\in A_i$ dla każdego $i$ i należy do przecięcia.
W przypadku pustego przecięcia otrzymujemy $X$, który
jest nasycony. Wreszcie dla dopełnienia załóżmy
$x\notin A$ oraz $x\sim y$. Gdyby $y\in A$, to przez
symetrię $y\sim x$, a nasycenie $A$ dawałoby $x\in A$,
sprzeczność. Zatem $y\notin A$.
\end{proof}

\begin{twierdzenie}[Nasycenie jako najmniejsze domknięcie]
\label{thm:eq-saturation-closure-xliii}
Niech $E$ będzie relacją równoważności na zbiorze $X$.
Dla każdego $A\subseteq X$ określmy
\[
 \operatorname{Sat}_E(A)
 :=\{y\in X:\exists x\in A\ (xEy)\}.
\]
Wtedy $\operatorname{Sat}_E(A)$ jest najmniejszym ze względu
na zawieranie podzbiorem nasyconym, który zawiera $A$.
Dla wszystkich $A,B\subseteq X$ zachodzą
\[
 A\subseteq\operatorname{Sat}_E(A),\quad
 A\subseteq B\Rightarrow
 \operatorname{Sat}_E(A)\subseteq\operatorname{Sat}_E(B),\quad
 \operatorname{Sat}_E(\operatorname{Sat}_E(A))=
 \operatorname{Sat}_E(A).
\]
Ponadto nasycenie zachowuje dowolne sumy rodzin podzbiorów
$X$, lecz nie musi zachowywać przecięć.
\end{twierdzenie}
\begin{proof}
Jeżeli $x\in A$, to $xEx$ przez zwrotność, zatem
$x\in\operatorname{Sat}_E(A)$. To daje pierwsze zawieranie.
Jeżeli $y\in\operatorname{Sat}_E(A)$ oraz $yEz$, to
istnieje $x\in A$ z $xEy$. Przechodniość daje $xEz$,
a więc $z\in\operatorname{Sat}_E(A)$. Jest to warunek
nasycenia. Jeśli nasycony zbiór $B$ zawiera $A$, a
$y\in\operatorname{Sat}_E(A)$, istnieje $x\in A\subseteq B$
z $xEy$, stąd z nasycenia $y\in B$. Otrzymaliśmy
minimalność.

Jeśli $A\subseteq B$, każdy świadek $x\in A$ w definicji
$\operatorname{Sat}_E(A)$ należy też do $B$, co daje
monotoniczność. Zastosujmy teraz minimalność do nasyconego
zbioru $\operatorname{Sat}_E(A)$: jego nasycenie jest
w nim zawarte, a pierwsze zawieranie daje kierunek odwrotny.
To jest idempotentność.

Dla rodziny $(A_i)_{i\in I}$ i $y\in X$ mamy
$y\in\operatorname{Sat}_E(\bigcup_i A_i)$ wtedy i tylko
wtedy, gdy istnieją $i\in I$ oraz $x\in A_i$ z $xEy$.
To z kolei zachodzi dokładnie wtedy, gdy
$y\in\bigcup_i\operatorname{Sat}_E(A_i)$.
Dowód obejmuje $I=\varnothing$: obie strony są wtedy puste.
Niezachowywanie przecięć pokazuje
Przykład~\ref{ex:eq-sat-intersection-xliii}.
\end{proof}
\begin{przyklad}[Dlaczego nasycenie nie zachowuje przecięć]
\label{ex:eq-sat-intersection-xliii}
Na $X=\Z$ przyjmijmy kongruencję modulo $2$ oraz
$A=\{0\}$, $B=\{2\}$. Mamy $A\cap B=\varnothing$,
więc $\operatorname{Sat}_E(A\cap B)=\varnothing$.
Jednak $\operatorname{Sat}_E(A)=\operatorname{Sat}_E(B)$
jest zbiorem wszystkich liczb parzystych, toteż przecięcie
obu nasyceń jest niepuste. Zawsze zachodzi tylko zawieranie
$\operatorname{Sat}_E(A\cap B)\subseteq
\operatorname{Sat}_E(A)\cap\operatorname{Sat}_E(B)$,
wynikające z monotoniczności.
\end{przyklad}



\begin{przyklad}[Podzbiór nie jest nasycony]
Niech $X=\Z$, a $\sim$ będzie kongruencją modulo $2$.
Zbiór liczb parzystych jest nasycony (to jedna klasa),
ale $A=\{0,2\}$ nie jest: $0\in A$ i $0\sim 4$,
lecz $4\notin A$. Suma \emph{wszystkich} klas dotykających
$A$ to cały zbiór liczb parzystych.
\end{przyklad}

\begin{definicja}[Dobrze określone działanie na klasach]
Niech $X$ będzie zbiorem z dwuargumentowym działaniem
$\star$, tzn. regułą, która każdej parze $(x,y)\in X\times X$
przypisuje jednoznaczny element $x\star y\in X$.
Niech $\sim$ będzie relacją równoważności na $X$.
Mówimy, że $\sim$ jest \emph{zgodna z działaniem $\star$}, jeśli
\[
 x\sim x',\ y\sim y'\quad\Longrightarrow\quad
 x\star y\sim x'\star y'
\]
dla wszystkich $x,x',y,y'\in X$.
Wówczas na klasach chcemy przyjąć regułę
$[x]\widehat\star[y]:=[x\star y]$.
\end{definicja}

\begin{twierdzenie}[Kryterium niezależności od reprezentantów]
\label{thm:eq-operation}
Niech $X$ będzie zbiorem, $\star$ działaniem na $X$,
a $\sim$ relacją równoważności na $X$. Wzór
$[x]\widehat\star[y]=[x\star y]$ dla $x,y\in X$
wyznacza jednoznaczne działanie na $X/{\sim}$ niezależne
od wyboru reprezentantów wtedy i tylko wtedy, gdy $\sim$ jest
zgodna z $\star$.
\end{twierdzenie}
\begin{proof}
Załóżmy zgodność. Ta sama klasa może mieć dwa zapisy
$[x]=[x']$ i $[y]=[y']$. Na mocy
Twierdzenia~\ref{thm:class-equality} oznacza to
$x\sim x'$ i $y\sim y'$. Zgodność daje
$x\star y\sim x'\star y'$, więc znów na mocy
Twierdzenia~\ref{thm:class-equality}
$[x\star y]=[x'\star y']$. Wynik nie zależy od zapisów klas.

Odwrotnie, załóżmy niezależność od reprezentantów.
Jeżeli $x\sim x'$ i $y\sim y'$, to $[x]=[x']$ i
$[y]=[y']$. Obliczając wynik działania na tych samych
klasach na dwa sposoby, dostajemy
$[x\star y]=[x'\star y']$. Ponowne użycie
Twierdzenia~\ref{thm:class-equality} prowadzi do
$x\star y\sim x'\star y'$. Jest to żądana zgodność.
\end{proof}

\begin{przyklad}[Arytmetyka klas reszt]\label{ex:eq-mod-arithmetic}
Na $\Z$ rozważmy kongruencję modulo $m\ge1$.
Jeżeli $a\sim_m a'$ i $b\sim_m b'$, to istnieją
$k,\ell\in\Z$ z $a-a'=mk$ i $b-b'=m\ell$.
Dla dodawania mamy
$(a+b)-(a'+b')=m(k+\ell)$, więc
$a+b\sim_m a'+b'$.
Dla mnożenia obliczamy
\[
 ab-a'b'=(a-a')b+a'(b-b')=m(kb+a'\ell),
\]
zatem $ab\sim_m a'b'$. Oba działania są zgodne z relacją.
Możemy więc bez dwuznaczności napisać
$[a]+[b]:=[a+b]$ i $[a]\cdot[b]:=[ab]$.
Przykładowo modulo $5$ zachodzi
$[3]+[4]=[7]=[2]$ oraz $[3]\cdot[4]=[12]=[2]$.
Ta konstrukcja prowadzi do pierścienia $\Z/m\Z$;
pojęcie pierścienia będzie rozwijane w skrypcie algebry.
\end{przyklad}

\begin{przyklad}[Działanie, które \emph{nie} schodzi na iloraz]
Niech na $\Z$ dana będzie kongruencja modulo $2$.
Rozważmy regułę $a\star b:=\max\{a,b\}$.
Mamy $0\sim_2 2$ i $1\sim_2 1$, lecz
$\max\{0,1\}=1$ oraz $\max\{2,1\}=2$, a liczby
$1,2$ nie są kongruentne modulo $2$. Dlatego nie istnieje
jednoznaczne działanie $[a]\widehat\star[b]:=[\max\{a,b\}]$.
Ten przykład pokazuje, dlaczego nie wolno pominąć sprawdzania
niezależności od reprezentantów.
\end{przyklad}


\begin{przyklad}[Działanie, które nie przechodzi na iloraz]
Niech $X=\Z$, a $a\sim b$ oznacza zgodność modulo $2$.
Zdefiniujmy działanie $a\star b:=a^2+b$. Ta operacja jest
zgodna z relacją, gdyż z $a\equiv a'\pmod 2$ wynika
$a^2\equiv (a')^2\pmod 2$, a zgodność $b\equiv b'\pmod 2$
pozwala dodać te kongruencje. Natomiast reguła
$a\diamond b:=\lfloor a/2\rfloor+b$ nie jest zgodna:
$0\sim 2$ i $0\sim 0$, lecz
$0\diamond0=0$ oraz $2\diamond0=1$ nie są równoważne.
Dlatego zapis $[a]\widehat\diamond[b]=[a\diamond b]$
nie definiuje działania na klasach: dla tej samej pary
klas $[0]=[2]$ i $[0]$ otrzymalibyśmy różne wyniki.
\end{przyklad}

\begin{uwaga}[Klasy a wybór reprezentantów]
Aby zdefiniować działanie na ilorazie, nie trzeba wybierać
po jednym reprezentancie z każdej klasy. Wystarczy dowieść,
że wynik nie zależy od reprezentantów. Rozróżnienie to stanie
się ważne w teorii aksjomatu wyboru: wybór reprezentanta
\emph{dla jednej wskazanej klasy} jest czymś innym niż
jednoczesny wybór reprezentantów dla całej rodziny klas.
\end{uwaga}

\begin{twierdzenie}[Własność uniwersalna odwzorowania ilorazowego]
\label{thm:lxv-quotient-factorization}
Niech $X,Y$ będą zbiorami, $E$ relacją równoważności na $X$,
a $q:X\to X/{E}$ odwzorowaniem $q(x)=[x]_E$.
Dla funkcji $f:X\to Y$ następujące warunki są równoważne:
\begin{enumerate}[label=(\roman*)]
\item dla wszystkich $x,x'\in X$, z $xEx'$ wynika $f(x)=f(x')$;
\item istnieje dokładnie jedna funkcja $\bar f:X/{E}\to Y$
taka, że $f=\bar f\circ q$.
\end{enumerate}
Twierdzenie pozostaje prawdziwe dla $X=\varnothing$ i dowolnego $Y$.
\end{twierdzenie}
\begin{proof}
Najpierw zauważmy, że $q$ jest funkcją: dla każdego $x\in X$
klasa $[x]_E$ istnieje i należy do $X/{E}$. Jest też surjekcją,
ponieważ każdy element ilorazu jest z definicji klasą pewnego $x\in X$.
Załóżmy (i). Dla klasy $C\in X/{E}$ określmy $\bar f(C)$ jako
$f(x)$ dla dowolnego $x\in X$ spełniającego $C=[x]_E$.
Taki $x$ istnieje z definicji ilorazu. Jeśli $C=[x]_E=[x']_E$,
to z Twierdzenia~\ref{thm:class-equality} mamy $xEx'$, więc
$f(x)=f(x')$. Wartość jest zatem niezależna od reprezentanta.
Nie wybieramy jednocześnie reprezentantów klas: graf funkcji
$\bar f$ jest jednoznacznie określoną relacją
$\{(C,y)\in(X/{E})\times Y:\exists x\in C\ (f(x)=y)\}$.
Dla $x\in X$ otrzymujemy $\bar f(q(x))=\bar f([x]_E)=f(x)$.
Jeśli $g:X/{E}\to Y$ również spełnia $f=g\circ q$,
to dla dowolnej klasy $C=[x]_E$ zachodzi
$g(C)=g(q(x))=f(x)=\bar f(C)$; stąd $g=\bar f$.

Odwrotnie, niech $f=\bar f\circ q$. Jeśli $xEx'$, to
$[x]_E=[x']_E$, a więc
$f(x)=\bar f([x]_E)=\bar f([x']_E)=f(x')$.
Dla $X=\varnothing$ także $X/{E}=\varnothing$;
jedyna funkcja z pustego zbioru do $Y$ spełnia oba warunki,
więc argument obejmuje również ten przypadek.
\end{proof}
\begin{przyklad}[Faktoryzacja kongruencji i jej brak]
\label{ex:lxv-quotient-factorization}
Na $\Z$ niech $aEb$ oznacza $a\equiv b\pmod 3$.
Funkcja $f(a)$ przypisująca resztę z dzielenia $a$ przez $3$
w zbiorze $\{0,1,2\}$ jest stała na każdej klasie:
$\bar f([0])=0$, $\bar f([1])=1$, $\bar f([2])=2$.
W szczególności $f(4)=\bar f([4])=\bar f([1])=1$.
Natomiast funkcja $h(a)=a$ nie jest stała na klasach,
ponieważ $0E3$, ale $h(0)=0\ne3=h(3)$.
Nie może więc istnieć funkcja $\bar h$ z $h=\bar h\circ q$.
\end{przyklad}

\section{Ułamki jako klasy par i podsumowanie}\label{sec:eq-fractions}
\begin{przyklad}[Konstrukcja liczb wymiernych z par liczb całkowitych]
Niech $D:=\{(p,q)\in\Z\times\Z:q\ne0\}$. Na $D$
rozważmy relację
\[
 (p,q)\sim(r,s)\quad\Longleftrightarrow\quad ps=rq.
\]
Udowodnimy, że jest to równoważność.
Zwrotność: $pq=pq$. Symetria: z $ps=rq$ wynika
$rq=ps$. Przechodniość wymaga większej ostrożności.
Załóżmy $(p,q)\sim(r,s)$ i $(r,s)\sim(t,u)$.
Mamy $ps=rq$ i $ru=ts$, przy czym $q,s,u\ne0$.
Mnożąc pierwszą równość przez $u$, a drugą przez $q$,
dostajemy $psu=rqu=tsq$. A zatem $(pu-tq)s=0$.
Ponieważ $s\ne0$ i w $\Z$ iloczyn dwóch liczb jest
zerem tylko wtedy, gdy przynajmniej jeden czynnik jest zerem,
wynika $pu=tq$. To znaczy $(p,q)\sim(t,u)$.

Zbiór ilorazowy $D/{\sim}$ jest standardowym modelem
zbioru liczb wymiernych: klasę $[(p,q)]$ zapisujemy $p/q$.
Teraz równanie $1/2=2/4$ nie jest gołosłowną umową,
lecz skrótem sprawdzonej równości klas:
$1\cdot4=2\cdot2$. Dodawanie i mnożenie na klasach
wymaga osobnego sprawdzenia zgodności, tak jak w
Twierdzeniu~\ref{thm:eq-operation}.
\end{przyklad}

\begin{twierdzenie}[Dodawanie i mnożenie ułamków nie zależą od zapisu]
\label{thm:eq-fraction-operations-xliii}
Niech $D=\{(p,q)\in\Z\times\Z:q\ne0\}$,
a $\sim$ będzie relacją równoważności na $D$ określoną przez
$(p,q)\sim(r,s)\iff ps=rq$ dla $(p,q),(r,s)\in D$.
Dla wszystkich $(p,q),(r,s)\in D$ wzory
\[
 [(p,q)]+[(r,s)]:=[(ps+rq,qs)],\qquad
 [(p,q)]\cdot[(r,s)]:=[(pr,qs)]
\]
wyznaczają dobrze określone działania na $D/{\sim}$.
W każdym z tych wzorów obliczenie nie zależy od
reprezentantów klas, a powstający mianownik jest niezerowy.
\end{twierdzenie}
\begin{proof}
Skoro $q,s\ne0$, to $qs\ne0$ w $\Z$, zatem obie
pary po prawej stronie należą do $D$. Pozostaje
niezależność od reprezentantów.
Załóżmy $(p,q)\sim(p',q')$ i $(r,s)\sim(r',s')$.
Z definicji oznacza to
$ p q'=p'q$ oraz $r s'=r's$, przy czym
$q,q',s,s'$ są niezerowe.

Dla mnożenia sprawdzamy warunek równoważności par:
\[
 (pr)(q's')=(pq')(rs')=(p'q)(r's)=(p'r')(qs).
\]
Wykorzystaliśmy dwie założone równości, a więc
$(pr,qs)\sim(p'r',q's')$.

Dla dodawania porównujemy liczniki na krzyż:
\begin{align*}
 (ps+rq)(q's')
 &= (pq')ss'+(rs')qq'\\
 &= (p'q)ss'+(r's)qq'\\
 &= (p's'+r'q')(qs).
\end{align*}
Pierwsza i ostatnia równość są zwykłym przestawianiem
czynników, a środkowa korzysta z $pq'=p'q$ i
$rs'=r's$. W ten sposób otrzymujemy
$(ps+rq,qs)\sim(p's'+r'q',q's')$.
Twierdzenie~\ref{thm:eq-operation} daje zatem oba
jednoznaczne działania na zbiorze klas.
\end{proof}
\begin{przyklad}[Dwa zapisy tego samego dodawania]
\label{ex:eq-fraction-arithmetic-xliii}
Dla reprezentantów $(1,2)$ i $(1,3)$ obliczamy
$[(1,2)]+[(1,3)]=[(5,6)]$.
Wybierzmy teraz równoważne pary $(2,4)$ i $(3,9)$.
Ich suma ma postać $[(2\cdot9+3\cdot4,4\cdot9)]
=[(30,36)]$. Sprawdzamy bez odwołania do intuicji
ułamków, że $(5,6)\sim(30,36)$, ponieważ
$5\cdot36=6\cdot30=180$.
Przy mnożeniu
$[(1,2)]\cdot[(1,3)]=[(1,6)]$ i
$[(2,4)]\cdot[(3,9)]=[(6,36)]$; obie klasy są równe,
bo $1\cdot36=6\cdot6$.
\end{przyklad}



\begin{uwaga}[Co w kolejnym rozdziale?]
Relacja równoważności grupuje obiekty w klasy, ale nie musi
porządkować tych klas od mniejszej do większej.
W rozdziale~\ref{ch:orders} zbadamy inną kombinację własności:
zwrotność, \emph{antysymetrię} i przechodniość. Dzięki temu
powstaną relacje porządku i pojęcia takie jak kres czy element
maksymalny. W rozdziale poświęconym funkcjom bardziej systematycznie
omówimy odwzorowanie kanoniczne $x\mapsto[x]$ i jego zastosowania;
własność uniwersalną tego odwzorowania udowodniliśmy już
w Twierdzeniu~\ref{thm:lxv-quotient-factorization}.
\end{uwaga}

\chapter{Relacje porządku: od porównywania do kresów}\label{ch:orders}

\section{Dlaczego porządek to nie zawsze ustawienie w jednym szeregu?}
W zbiorze liczb rzeczywistych możemy porównać każdą parę liczb:
jedna jest mniejsza lub równa drugiej. Na zbiorze podzbiorów
ustalonego zbioru sytuacja jest inna. Dla $A=\{1\}$ i $B=\{2\}$
nie zachodzi ani $A\subseteq B$, ani $B\subseteq A$.
Mimo to zawieranie jest sensownym sposobem porządkowania podzbiorów.
Podobnie dwie różne liczby dodatnie nie zawsze dzielą się nawzajem,
a mimo to podzielność wyraża relację „bycia prostszym czynnikiem”.

Relacja równoważności z rozdziału~\ref{ch:equivalence} mówiła,
które obiekty uznajemy za nierozróżnialne. Relacja porządku
mówi, które obiekty można porównać i co z takiego porównania wynika.
W pierwszej kolejności rozróżnimy \emph{preporządek}, \emph{porządek
częściowy} i \emph{porządek liniowy}.

\begin{definicja}[Preporządek i porządek częściowy]\label{def:partial-order}
Niech $P$ będzie zbiorem, a $\preceq\subseteq P\times P$
relacją. Nazywamy ją \emph{preporządkiem}, jeśli jest zwrotna
i przechodnia. Nazywamy ją \emph{porządkiem częściowym}, jeśli
jest zwrotna, antysymetryczna i przechodnia.
Parę $(P,\preceq)$, w której $\preceq$ jest porządkiem
częściowym, nazywamy \emph{zbiorem częściowo uporządkowanym}.
Dwa elementy $a,b\in P$ są \emph{porównywalne}, jeśli
$a\preceq b$ lub $b\preceq a$. W przeciwnym razie są
\emph{nieporównywalne}.
\end{definicja}

\begin{przyklad}[Zawieranie podzbiorów]\label{ex:order-powerset}
Niech $X$ będzie zbiorem. Na $\mathcal P(X)$ relacja
$A\preceq B\iff A\subseteq B$
jest porządkiem częściowym: zwrotność zawierania wynika
z implikacji $x\in A\Rightarrow x\in A$, przechodniość
z dwóch kolejnych zawierań $A\subseteq B\subseteq C$,
a antysymetria z kryterium dwóch zawierań
(Twierdzenie~\ref{thm:two-inclusions}).
Dla $X=\{1,2\}$ podzbiory $\{1\}$ i $\{2\}$ są
nieporównywalne, ale obydwa leżą między $\varnothing$ i $X$.
\end{przyklad}

\begin{przyklad}[Podzielność wymaga wyboru zbioru podstawowego]
Na $\N_{>0}$ relacja $a\preceq b\iff a\mid b$ jest
zwrotna, gdyż $a=a\cdot1$. Jest przechodnia: z
$b=ak$ i $c=b\ell$ wynika $c=a(k\ell)$.
Jest też antysymetryczna: jeśli $b=ak$ i $a=b\ell$,
gdzie $a,b,k,\ell$ są dodatnimi liczbami całkowitymi, to
$a=a(k\ell)$, skąd $k\ell=1$, więc $k=\ell=1$ i $a=b$.
Na $\Z\setminus\{0\}$ podzielność \emph{nie} jest
antysymetryczna: $2\mid -2$ oraz $-2\mid2$, lecz $2\ne-2$.
\end{przyklad}

\begin{przyklad}[Preporządek, który nie jest porządkiem]
Na $\Z$ rozważmy $a\preceq b\iff |a|\le |b|$.
Zwrotność i przechodniość wynikają z własności porządku
liczb nieujemnych. Jednak $2\preceq -2$ i
$-2\preceq2$, a $2\ne-2$. Brakuje antysymetrii.
To pokazuje, dlaczego sama zwrotność i przechodniość
nie wystarczają do uzyskania porządku częściowego.
\end{przyklad}

\section{Porządki ostre, liniowe i przechodzenie od jednego do drugiego}
\label{sec:order-strict}
\begin{definicja}[Porządek ostry i porządek liniowy]\label{def:strict-total}
Relację $\prec$ na zbiorze $P$ nazywamy \emph{ostrym porządkiem
częściowym}, jeśli jest przeciwzwrotna i przechodnia.
Porządek częściowy $\preceq$ jest \emph{liniowy} (całkowity),
jeśli dla każdych $a,b\in P$ zachodzi $a\preceq b$ lub
$b\preceq a$. Ostry porządek jest \emph{ostrym porządkiem
liniowym}, jeżeli dla każdych różnych $a,b\in P$ zachodzi
$a\prec b$ lub $b\prec a$.
\end{definicja}

\begin{twierdzenie}[Przejście między porządkiem ostrym i nieostrym]
\label{thm:order-strict-correspondence}
Niech $P$ będzie zbiorem.
\begin{enumerate}[label=(\roman*)]
\item Jeżeli $\preceq$ jest porządkiem częściowym na $P$,
to relacja $\prec$ na $P$ określona dla $a,b\in P$ wzorem
$a\prec b\iff(a\preceq b\text{ i }a\ne b)$ jest ostrym
porządkiem częściowym.
\item Jeżeli $\prec$ jest ostrym porządkiem częściowym na $P$,
to relacja $\preceq$ na $P$ określona dla $a,b\in P$ wzorem
$a\preceq b\iff(a\prec b\text{ lub }a=b)$
jest porządkiem częściowym.
\item Obie konstrukcje są wzajemnie odwrotne i zachowują
liniowość.
\end{enumerate}
\end{twierdzenie}
\begin{proof}
(i) Ponieważ $a\ne a$ jest fałszem, żadne $a$ nie
spełnia $a\prec a$. Dla przechodniości załóżmy
$a\prec b$ i $b\prec c$. Z definicji mamy
$a\preceq b$ i $b\preceq c$, więc z przechodniości
$a\preceq c$. Gdyby $a=c$, to równocześnie zachodziłyby
$a\preceq b$ i $b\preceq a$. Antysymetria dawałaby
$a=b$, sprzecznie z $a\prec b$. Zatem $a\ne c$, a więc
$a\prec c$.

(ii) Nowa relacja jest zwrotna, bo $a=a$. Jeśli
$a\preceq b$ i $b\preceq a$, to albo $a=b$, albo
$a\prec b$ i $b\prec a$ (w przypadkach mieszanych
równość też dałaby $a=b$). Druga możliwość jest niemożliwa:
z przechodniości $a\prec a$, wbrew przeciwzwrotności.
Mamy zatem antysymetrię. Jeśli $a\preceq b$ i
$b\preceq c$, rozpatrujemy, czy każda z tych dwóch
relacji pochodzi z równości czy z $\prec$. Jeśli którakolwiek
jest równością, wynik $a\preceq c$ otrzymujemy bezpośrednio
z drugiej. Jeśli obie są ostre, daje go przechodniość $\prec$.

(iii) W pierwszej konstrukcji usuwamy dokładnie pary
$(a,a)$, a w drugiej dodajemy je z powrotem. Porządek ostry
nie zawiera par przekątnych, a porządek częściowy zawiera
wszystkie, więc odzyskujemy odpowiednio pierwotną relację.
Porównywalność różnych elementów jest niezmieniona;
dlatego liniowość obu wersji jest równoważna.
\end{proof}

\begin{przyklad}[Uwaga o słowie „ostry”]
Relacja $<$ na $\R$ jest ostrym porządkiem liniowym.
Dla dowolnego zbioru $X$ relacja właściwego zawierania
$\subsetneq$ na $\mathcal P(X)$ jest ostrym porządkiem częściowym,
ale gdy $X$ ma co najmniej dwa różne elementy,
nie jest liniowa: singletony tych elementów są
nieporównywalne. Ostrość i liniowość są niezależnymi cechami.
\end{przyklad}

\section{Łańcuchy, antyłańcuchy i diagramy Hassego}
\label{sec:order-hasse}
\begin{definicja}[Łańcuch i antyłańcuch]\label{def:chain-antichain}
Niech $(P,\preceq)$ będzie zbiorem częściowo uporządkowanym.
Podzbiór $C\subseteq P$ nazywamy \emph{łańcuchem}, jeśli
dowolne dwa elementy $C$ są porównywalne.
Podzbiór $D\subseteq P$ nazywamy \emph{antyłańcuchem}, jeśli
żadne dwa \emph{różne} elementy $D$ nie są porównywalne.
\end{definicja}

\begin{przyklad}[Podzbiory zbioru trójelementowego]
Na $\mathcal P(\{1,2,3\})$ rozważmy rodzinę
$\mathcal C=\{\varnothing,\{1\},\{1,2\},\{1,2,3\}\}$.
Jest ona łańcuchem, gdyż
$\varnothing\subsetneq\{1\}\subsetneq\{1,2\}\subsetneq\{1,2,3\}$.
Z kolei
$\{\{1\},\{2\},\{3\}\}$ jest antyłańcuchem:
żaden z trzech jednoelementowych zbiorów nie jest
podzbiorem innego. Te same elementy mogą brać udział
w różnych łańcuchach; łańcuch nie musi być maksymalny.
\end{przyklad}

\begin{definicja}[Relacja pokrywania]
W zbiorze częściowo uporządkowanym $(P,\preceq)$ oznaczmy
przez $a\prec b$ warunek $a\preceq b$ i $a\ne b$.
Mówimy, że $b$ \emph{pokrywa} $a$, i piszemy $a\lessdot b$, jeśli
$a\prec b$ oraz nie istnieje $c\in P$ takie, że
$a\prec c\prec b$.
Dla skończonego porządku \emph{diagram Hassego} ma po
jednym wierzchołku na element zbioru i krawędź prowadzącą
w górę od $a$ do $b$ dokładnie wtedy, gdy $a\lessdot b$.
Pomijamy pętle i krawędzie wynikające z przechodniości.
\end{definicja}

\begin{center}
\begin{tikzpicture}[scale=.85, every node/.style={font=\small}]
\node (e) at (0,0) {$\varnothing$};
\node (a) at (-1.4,1.5) {$\{1\}$};
\node (b) at (1.4,1.5) {$\{2\}$};
\node (t) at (0,3.0) {$\{1,2\}$};
\draw[thick] (e)--(a)--(t)--(b)--(e);
\node[font=\small,align=left,anchor=west] at (2.3,1.5)
{Brak krawędzi między\\$\{1\}$ i $\{2\}$: są\\nieporównywalne.};
\end{tikzpicture}

{\small Diagram Hassego porządku $(\mathcal P(\{1,2\}),\subseteq)$.}
\end{center}

\begin{uwaga}[Dlaczego na rysunku nie ma strzałek?]
Na diagramie Hassego umownie czytamy krawędzie od dołu do góry.
Na przykład $\varnothing\subseteq\{1,2\}$ odczytujemy
z dwukrawędziowej drogi przechodzącej przez $\{1\}$,
chociaż nie narysowano osobnej krawędzi od dolnego do górnego
wierzchołka. Dla skończonego porządku relację wyznaczają
osiągalność po krawędziach i równość elementów.
\end{uwaga}


\begin{przyklad}[Dwa elementy minimalne i brak najmniejszego]
Niech zbiór $P=\{a,b,c\}$ ma tylko porównania
$a\le c$, $b\le c$ oraz porównania wynikające ze zwrotności.
Elementy $a$ i $b$ są minimalne: nie mają elementów ściśle mniejszych.
Nie ma jednak elementu najmniejszego, bo $a\not\le b$ i $b\not\le a$.
W porządku liniowym taka sytuacja nie jest możliwa, ponieważ
dowolne dwa elementy można porównać.
\end{przyklad}

\section{Elementy najmniejsze, minimalne i ich odpowiedniki}
\label{sec:order-minmax}
\begin{definicja}[Minimum, maksimum, elementy minimalne i maksymalne]
\label{def:order-extrema}
Niech $(P,\preceq)$ będzie zbiorem częściowo uporządkowanym,
a $A\subseteq P$ podzbiorem. W poniższych warunkach
$a\prec b$ oznacza $a\preceq b$ i $a\ne b$.
\begin{enumerate}[label=(\roman*)]
\item $m\in A$ jest \emph{elementem najmniejszym} $A$,
jeśli $m\preceq a$ dla każdego $a\in A$.
\item $M\in A$ jest \emph{elementem największym} $A$,
jeśli $a\preceq M$ dla każdego $a\in A$.
\item $m\in A$ jest \emph{elementem minimalnym} $A$,
jeśli nie istnieje $a\in A$ z $a\prec m$.
\item $M\in A$ jest \emph{elementem maksymalnym} $A$,
jeśli nie istnieje $a\in A$ z $M\prec a$.
\end{enumerate}
Porządki w definicjach są ograniczeniami porządku z $P$ do $A$.
\end{definicja}

\begin{twierdzenie}[Jedyność elementów najmniejszych i największych]
\label{thm:order-least-unique}
W zbiorze częściowo uporządkowanym $(P,\preceq)$ każdy
podzbiór $A\subseteq P$ ma co najwyżej jeden element
najmniejszy i co najwyżej jeden element największy.
Każdy element najmniejszy jest minimalny, a każdy
największy jest maksymalny.
\end{twierdzenie}
\begin{proof}
Niech $m_1,m_2$ będą elementami najmniejszymi $A$.
Z definicji $m_1\preceq m_2$ i $m_2\preceq m_1$.
Antysymetria daje $m_1=m_2$. Podobnie, jeśli
$M_1,M_2$ są największe, to $M_1\preceq M_2$ i
$M_2\preceq M_1$, więc $M_1=M_2$.

Jeżeli $m$ jest najmniejszy, to dla każdego $a\in A$
zachodzi $m\preceq a$. Gdyby istniał $a\in A$ z
$a\prec m$, mielibyśmy także $a\preceq m$ i $a\ne m$,
co wraz z $m\preceq a$ przeczyłoby antysymetrii.
Zatem taki $a$ nie istnieje i $m$ jest minimalny.
Jeżeli zaś $M$ jest największy, to każdy $a\in A$ spełnia
$a\preceq M$. Gdyby istniał $b\in A$ z $M\prec b$,
mielibyśmy $M\preceq b$ i $M\ne b$. Ponieważ $M$ jest największy,
mamy także $b\preceq M$. Antysymetria wymusiłaby $b=M$,
wbrew nierówności. Zatem $M$ jest maksymalny.
\end{proof}

\begin{przyklad}[Dwa różne elementy minimalne]
Niech $P=\mathcal P(\{1,2\})$ z porządkiem zawierania,
a $A=\{\{1\},\{2\}\}$. Oba elementy $A$ są minimalne:
w $A$ nie ma od nich mniejszych elementów. Żaden nie jest
jednak najmniejszy, gdyż $\{1\}\nsubseteq\{2\}$
i $\{2\}\nsubseteq\{1\}$. Analogicznie oba są
maksymalne, ale nie ma elementu największego.
Jest to podstawowy kontrprzykład do mylenia minimum z
minimalnością.
\end{przyklad}

\begin{twierdzenie}[Istnienie ekstremów w skończonym porządku]
\label{thm:finite-poset-extrema}
Niech $(P,\preceq)$ będzie zbiorem częściowo
uporządkowanym, a $A\subseteq P$ niepustym skończonym
podzbiorem. Wtedy $A$ ma element minimalny i element
maksymalny. Niekoniecznie ma element najmniejszy lub
największy.
\end{twierdzenie}
\begin{proof}
Wybierzmy $a_0\in A$; można to zrobić, ponieważ $A$
jest niepusty. Jeśli $a_0$ jest minimalny, mamy tezę.
W przeciwnym razie istnieje $a_1\in A$ z
$a_1\prec a_0$. Jeżeli $a_1$ nie jest minimalny,
znajdujemy $a_2\in A$ z $a_2\prec a_1$ itd.
W każdym kroku otrzymujemy element różny od wszystkich
wcześniejszych: gdyby $a_j=a_i$ dla $j>i$, to z
przechodniości wynikłoby $a_i\prec a_i$, sprzecznie
z przeciwzwrotnością ostrego porządku
(Twierdzenie~\ref{thm:order-strict-correspondence}).
Ponieważ $A$ jest skończony, nie możemy wytworzyć
dowolnie długiego ciągu parami różnych elementów $A$.
Proces musi zakończyć się na elemencie minimalnym.

Aby uzyskać element maksymalny, powtarzamy argument,
zastępując każdy krok $a_{j+1}\prec a_j$ krokiem
$a_j\prec a_{j+1}$. Ponownie skończoność wymusza
zakończenie. Przykład z dwoma singletonami powyżej
pokazuje, że minimum i maksimum nie muszą istnieć.
\end{proof}


\begin{twierdzenie}[Odczytywanie porządku z diagramu Hassego]
\label{thm:hasse-paths-xliii}
Niech $(P,\preceq)$ będzie \emph{skończonym} zbiorem
częściowo uporządkowanym. Dla $x,y\in P$ piszemy $x\prec y$,
gdy $x\preceq y$ i $x\ne y$, oraz $x\lessdot y$,
gdy $x\prec y$ i nie istnieje $z\in P$ z $x\prec z\prec y$.
Dla różnych $a,b\in P$ zachodzi $a\prec b$
wtedy i tylko wtedy, gdy istnieją
$a=x_0,x_1,\ldots,x_k=b$ z $k\ge1$ i
$x_{j-1}\lessdot x_j$ dla każdego $j=1,\ldots,k$.
Zatem porządek można odtworzyć z diagramu Hassego,
dodając równość elementów i przechodniość wzdłuż
rosnących dróg.
\end{twierdzenie}
\begin{proof}
Jeżeli istnieje taka droga, każda krawędź pokrywania
spełnia $x_{j-1}\prec x_j$. Przechodniość daje
$a\prec b$ (dla $k=1$ jest to pierwsza krawędź).

W przeciwnym kierunku załóżmy $a\prec b$.
Rozważmy skończony, niepusty zbiór
$S=\{x\in P:a\prec x\preceq b\}$; należy do niego $b$.
Na mocy Twierdzenia~\ref{thm:finite-poset-extrema}
$S$ ma element minimalny $x_1$. Mamy $a\prec x_1$.
Gdyby istniało $z$ z $a\prec z\prec x_1$, to
$z\preceq x_1\preceq b$ i $z\in S$,
co przeczyłoby minimalności $x_1$. Zatem $a\lessdot x_1$.
Jeśli $x_1=b$, kończymy. W przeciwnym razie stosujemy
identyczną konstrukcję do pary $x_1\prec b$.
Każdy krok daje większy element $x_j\preceq b$;
żaden nie może powtórzyć wcześniejszego, bo wówczas
z przechodniości wynikałoby $x_j\prec x_j$.
Skończoność $P$ wymusza zakończenie i osiągnięcie $b$.
\end{proof}
\begin{przyklad}[Dlaczego skończoność jest istotna]
\label{ex:hasse-infinite-caution-xliii}
Na $(\Q,\le)$ nie istnieje żadna para $a\lessdot b$:
jeśli $a<b$, liczba $c=(a+b)/2$ jest wymierna i
spełnia $a<c<b$. Diagram z samymi krawędziami
pokrywania nie miałby więc żadnej krawędzi, choć
wiele liczb jest porównywalnych. Odczytywanie porządku
z dróg diagramu Hassego bez założenia skończoności
jest w ogólności błędne.
\end{przyklad}


\begin{twierdzenie}[Jedyny element minimalny w skończonym porządku]
\label{thm:unique-minimal-finite-xliii}
Niech $(P,\preceq)$ będzie zbiorem częściowo uporządkowanym,
a $A\subseteq P$ niepustym, skończonym podzbiorem.
Wtedy $A$ ma dokładnie jeden element minimalny wtedy i tylko
wtedy, gdy ma element najmniejszy.
Analogicznie $A$ ma dokładnie jeden element maksymalny
wtedy i tylko wtedy, gdy ma element największy.
\end{twierdzenie}
\begin{proof}
Jeżeli $m$ jest elementem najmniejszym $A$, to jest
minimalny na mocy Twierdzenia~\ref{thm:order-least-unique}.
Gdyby $u\ne m$ był innym elementem minimalnym, to
$m\preceq u$ i $m\ne u$, czyli $m\prec u$,
sprzecznie z minimalnością $u$. Zatem $m$ jest jedynym
minimalnym elementem.

Odwrotnie, niech $m$ będzie jedynym minimalnym elementem $A$.
Wybierzmy dowolny $a\in A$. Zbiór
$D_a=\{x\in A:x\preceq a\}$ jest niepusty (zawiera $a$)
i skończony. Na mocy
Twierdzenia~\ref{thm:finite-poset-extrema} ma element
minimalny $u$. Gdyby istniał $v\in A$ z $v\prec u$,
to $v\preceq u\preceq a$, a więc $v\in D_a$;
byłaby to sprzeczność z minimalnością $u$ w $D_a$.
Zatem $u$ jest minimalny w całym $A$, czyli $u=m$.
Ponieważ $u\in D_a$, mamy $m\preceq a$. Element $a$
był dowolny, więc $m$ jest najmniejszy w $A$.
Dla elementów maksymalnych rozumowanie działa z odwróconym
kierunkiem porządku: w miejscu $D_a$ stosujemy
$\{x\in A:a\preceq x\}$.
\end{proof}

\begin{przyklad}[Jedyny element minimalny nie musi być najmniejszy]
\label{ex:unique-minimal-infinite-xliii}
Skończoności nie można usunąć z poprzedniego twierdzenia.
Niech $P=\{m\}\cup\{b_n:n\in\N\}$, przy czym
wszystkie oznaczenia są różnymi elementami.
Ustalmy zwrotne porównania oraz
$b_j\preceq b_i$ dokładnie wtedy, gdy $j\ge i$;
nie porównujemy $m$ z żadnym $b_i$.
Wtedy $P$ jest zbiorem częściowo uporządkowanym.
Element $m$ jest minimalny, bo nic nie leży ściśle
poniżej niego. Żaden $b_i$ nie jest minimalny, gdyż
$b_{i+1}\prec b_i$. Zatem $m$ jest \emph{jedynym}
elementem minimalnym, lecz nie jest elementem najmniejszym:
nie zachodzi $m\preceq b_0$.
W podzbiorze $\{b_n:n\in\N\}$ nie ma żadnego
elementu minimalnego. To właśnie umożliwia kontrprzykład.

\begin{center}
\begin{tikzpicture}[x=.92cm,y=.78cm, every node/.style={font=\small}]
 \node (m) at (-1.9,1.45) {$m$};
 \node (b0) at (1.25,3.2) {$b_0$};
 \node (b1) at (1.25,2.35) {$b_1$};
 \node (b2) at (1.25,1.50) {$b_2$};
 \node (b3) at (1.25,.65) {$b_3$};
 \draw[thick] (b0)--(b1)--(b2)--(b3);
 \node at (1.25,-.06) {$\vdots$};
 \node[anchor=north,align=center] at (-1.9,.92) {jedyny element\\minimalny};
 \node[anchor=west,align=left] at (2.05,1.6) {łańcuch bez\\elementu minimalnego};
\end{tikzpicture}

{\small Schemat nieskończonego porządku: $m$ jest nieporównywalny
z każdym $b_n$, a łańcuch po prawej biegnie w dół bez końca.}
\end{center}
\end{przyklad}

\section{Ograniczenia, kres dolny i kres górny}\label{sec:order-bounds}
\begin{definicja}[Ograniczenia i kresy]\label{def:bounds-sup-inf}
Niech $(P,\preceq)$ będzie zbiorem częściowo uporządkowanym,
a $A\subseteq P$.
Element $u\in P$ jest \emph{ograniczeniem górnym} $A$,
jeżeli $a\preceq u$ dla wszystkich $a\in A$.
Element $\ell\in P$ jest \emph{ograniczeniem dolnym} $A$,
jeżeli $\ell\preceq a$ dla wszystkich $a\in A$.

Jeśli zbiór wszystkich ograniczeń górnych $A$ ma element
najmniejszy, nazywamy go \emph{kresem górnym} $A$ i piszemy
$\sup_P A$. Jeśli zbiór wszystkich ograniczeń dolnych
ma element największy, nazywamy go \emph{kresem dolnym} $A$
i piszemy $\inf_P A$. Symbolu kresu używamy tylko, gdy
taki element istnieje w $P$.
\end{definicja}

\begin{uwaga}[Kres nie musi należeć do rozpatrywanego zbioru]
W $P=\R$ z naturalnym porządkiem zbiór
$A=(0,1)$ ma $\sup_{\R}A=1$ i $\inf_{\R}A=0$,
ale żaden z tych kresów nie należy do $A$.
W zbiorze liczb wymiernych $\Q$ zbiór
$B=\{q\in\Q:q\ge0,\ q^2<2\}$ nie ma kresu górnego
\emph{w $\Q$}: jedyny rzeczywisty kres to $\sqrt2$,
który nie należy do $\Q$. Stąd w definicji zawsze
wskazujemy, w jakim porządku szukamy kresu.
\end{uwaga}

\begin{twierdzenie}[Kryterium kresu górnego]\label{thm:sup-characterization}
Niech $(P,\preceq)$ będzie zbiorem częściowo uporządkowanym,
a $A\subseteq P$. Element $s\in P$ jest kresem górnym $A$
wtedy i tylko wtedy, gdy
\begin{enumerate}[label=(\roman*)]
\item dla każdego $a\in A$ zachodzi $a\preceq s$;
\item dla każdego $u\in P$, jeżeli $a\preceq u$ dla każdego
$a\in A$, to $s\preceq u$.
\end{enumerate}
Analogicznie $t\in P$ jest kresem dolnym $A$ wtedy
i tylko wtedy, gdy $t\preceq a$ dla wszystkich $a\in A$
oraz $\ell\preceq t$ dla każdego ograniczenia dolnego
$\ell$ zbioru $A$.
\end{twierdzenie}
\begin{proof}
Z definicji $s=\sup_P A$ znaczy, że jest
\emph{najmniejszym} elementem zbioru wszystkich ograniczeń
górnych. To dokładnie oznacza, że sam jest ograniczeniem
górnym (warunek (i)) oraz że leży poniżej każdego innego
ograniczenia górnego (warunek (ii)). Każdy z tych dwóch
warunków jest konieczny, a ich koniunkcja wystarcza.
Dla kresu dolnego odwracamy kierunek relacji i powtarzamy
ten sam argument. Jedyność kresów wynika z
Twierdzenia~\ref{thm:order-least-unique} zastosowanego
do zbiorów ograniczeń.
\end{proof}

\begin{przyklad}[Suma i przecięcie jako kresy]\label{ex:orders-sets-lattice}
Niech $X$ będzie zbiorem, a $A,B\subseteq X$. W porządku
$(\mathcal P(X),\subseteq)$ zachodzą równości
\[
 \sup_{\mathcal P(X)}\{A,B\}=A\cup B,\qquad
 \inf_{\mathcal P(X)}\{A,B\}=A\cap B.
\]
Istotnie, $A\subseteq A\cup B$ i $B\subseteq A\cup B$,
więc suma jest ograniczeniem górnym pary $\{A,B\}$.
Jeśli $U$ jest dowolnym ograniczeniem górnym, to
$A\subseteq U$ i $B\subseteq U$, z czego dla każdego
$x\in A\cup B$ wynika $x\in U$. A zatem
$A\cup B\subseteq U$, więc suma jest najmniejszym
ograniczeniem górnym. Dla przecięcia analogicznie:
$A\cap B$ zawiera się w obu zbiorach, a każdy
podzbiór zawarty jednocześnie w $A$ i $B$ zawiera się
w $A\cap B$.
\end{przyklad}

\begin{przyklad}[Największy wspólny dzielnik i najmniejsza wspólna wielokrotność]
Niech $a,b\in\N_{>0}$. W porządku podzielności na $\N_{>0}$
ograniczeniami
dolnymi pary $\{a,b\}$ są dokładnie wspólne dzielniki
$a,b$, a górnymi --- ich wspólne wielokrotności.
Kresem dolnym tej pary jest $\gcd(a,b)$, a górnym
$\operatorname{lcm}(a,b)$. Dla $a=12$, $b=18$
mamy $\inf\{12,18\}=6$ i $\sup\{12,18\}=36$
\emph{w porządku podzielności}, mimo że w zwykłym
porządku liczb naturalnych minimum tej pary to $12$,
a maksimum to $18$.
\end{przyklad}

\section{Kraty i działania kresowe}\label{sec:order-lattices}
\begin{definicja}[Krata]\label{def:lattice}
Zbiór częściowo uporządkowany $(L,\preceq)$ jest
\emph{kratą}, jeśli każda para $a,b\in L$ ma kres
dolny $a\wedge b$ i kres górny $a\vee b$ w $L$.
Krata jest \emph{zupełna}, jeśli \emph{każdy} podzbiór
$A\subseteq L$, także pusty, ma kres dolny i kres górny.
\end{definicja}

\begin{twierdzenie}[Pusta rodzina i krańcowe elementy kraty zupełnej]
\label{thm:complete-lattice-bottom-top-xliii}
Jeśli $(L,\preceq)$ jest kratą zupełną, to $L$ ma
jedyny element najmniejszy $0_L$ oraz jedyny element
największy $1_L$. Ponadto
$0_L=\sup_L\varnothing$ i $1_L=\inf_L\varnothing$.
Nie każda krata jest zupełna: nawet krata, w której istnieją
kresy wszystkich skończonych niepustych podzbiorów, może
nie mieć kresów pewnych podzbiorów nieskończonych.
\end{twierdzenie}
\begin{proof}
Z definicji zupełności istnieje $s=\sup_L\varnothing$.
Każdy element $u\in L$ jest ograniczeniem górnym pustego
zbioru, ponieważ nie ma elementu $a\in\varnothing$,
który mógłby naruszyć nierówność $a\preceq u$.
Skoro $s$ jest \emph{najmniejszym} z tych ograniczeń,
to $s\preceq u$ dla wszystkich $u\in L$. Jest więc
najmniejszy. Analogicznie każdy element $\ell\in L$
jest ograniczeniem dolnym pustego zbioru, a jego największe
ograniczenie dolne $t=\inf_L\varnothing$ spełnia
$\ell\preceq t$ dla każdego $\ell\in L$.
Jest więc największy. Oba elementy są jedyne dzięki
antysymetrii. Samo istnienie $s$ pokazuje, że krata
zupełna musi być niepusta.

Aby wykazać, że nie każda krata jest zupełna, rozważmy
$(\N,\le)$. Każda para liczb naturalnych ma minimum
i maksimum, lecz cały zbiór $\N$ nie ma ograniczenia
górnego w $\N$: dla każdego $u\in\N$ zachodzi
$u<u+1$. Zatem $\sup_{\N}\N$ nie istnieje.
\end{proof}
\begin{przyklad}[Krata, która nie jest zupełna]
\label{ex:lattice-not-complete-xliii}
Rozważmy $L=\N$ ze zwykłym porządkiem (z zerem).
Każda para liczb naturalnych ma minimum i maksimum,
więc $L$ jest kratą. Co więcej, każdy niepusty
skończony podzbiór ma minimum i maksimum. Jednak zbiór
$A=\N\subseteq L$ nie ma w $L$ żadnego ograniczenia
górnego: dla dowolnego $u\in\N$ liczba $u+1\in A$
jest większa niż $u$. W konsekwencji $\sup_L A$
nie istnieje i krata $L$ nie jest zupełna.
Dodatkowo $\inf_L\varnothing$ nie istnieje, bo
$\N$ nie ma elementu największego.
\end{przyklad}


\begin{twierdzenie}[Podstawowe prawa działań w kracie]\label{thm:lattice-laws}
Niech $(L,\preceq)$ będzie kratą. Dla wszystkich
$a,b,c\in L$ zachodzą
\begin{align*}
 a\wedge b&=b\wedge a,&a\vee b&=b\vee a,\\
 (a\wedge b)\wedge c&=a\wedge(b\wedge c),&
 (a\vee b)\vee c&=a\vee(b\vee c),\\
 a\wedge a&=a,&a\vee a&=a,\\
 a\wedge(a\vee b)&=a,&a\vee(a\wedge b)&=a.
\end{align*}
\end{twierdzenie}
\begin{proof}
Ponieważ zbiór $\{a,b\}$ nie zmienia się po zamianie
$a$ i $b$, jego kresy także się nie zmieniają.
To dowodzi przemienności obu działań.

Dla łączności $\wedge$ zauważmy, że
$(a\wedge b)\wedge c$ jest największym elementem
leżącym jednocześnie poniżej $a\wedge b$ i $c$.
Każdy taki element leży poniżej $a$, $b$ i $c$.
Odwrotnie, jeśli $x\preceq a,b,c$, to z definicji
$a\wedge b$ wynika $x\preceq a\wedge b$, a ponadto
$x\preceq c$; zatem $x\preceq(a\wedge b)\wedge c$.
Pokazuje to, że $(a\wedge b)\wedge c$ jest
największym wspólnym ograniczeniem dolnym trzech
elementów. Dokładnie ten sam argument, zaczynając
od pary $(b,c)$, pokazuje, że $a\wedge(b\wedge c)$
ma tę samą własność. Jedyność kresu daje równość.
Dla $\vee$ zamieniamy we wszystkich zdaniach
„poniżej” na „powyżej” i „największe ograniczenie dolne”
na „najmniejsze ograniczenie górne”.

Dla idempotentności obserwujemy, że $a$ jest zarazem
największym elementem poniżej $a,a$ i najmniejszym
powyżej $a,a$. Stąd $a\wedge a=a\vee a=a$.
Dla absorpcji: $a\preceq a\vee b$, więc $a$ jest
dolnym ograniczeniem pary $\{a,a\vee b\}$.
Każde dolne ograniczenie tej pary musi leżeć poniżej $a$,
zatem $a$ jest jej kresem dolnym, czyli
$a\wedge(a\vee b)=a$. Analogicznie
$a\wedge b\preceq a$ i każde ograniczenie górne
pary $\{a,a\wedge b\}$ leży powyżej $a$;
stąd $a\vee(a\wedge b)=a$.
\end{proof}

\begin{uwaga}[Nie każda krata jest rozdzielna]
Prawa absorpcji są prawdziwe w każdej kracie, ale
prawa rozdzielności nie muszą zachodzić. Istotnym
kontrprzykładem jest pięcioelementowa krata z trzema
wzajemnie nieporównywalnymi elementami $a,b,c$ między
najmniejszym $0$ a największym $1$.
W niej $b\vee c=1$, więc $a\wedge(b\vee c)=a$,
a jednocześnie $a\wedge b=a\wedge c=0$, zatem
$(a\wedge b)\vee(a\wedge c)=0\ne a$.
Nie wolno więc uogólniać praw rozdzielności sum i
przecięć zbiorów z rozdziału~\ref{ch:sets} na wszystkie kraty.
\end{uwaga}

\begin{center}
\begin{tikzpicture}[scale=.8]
 \node (z) at (0,0) {$0$}; \node (o) at (0,3) {$1$};
 \node (a) at (-2,1.5) {$a$};
 \node (b) at (0,1.5) {$b$};
 \node (c) at (2,1.5) {$c$};
 \foreach \x in {a,b,c} {\draw[thick] (z)--(\x)--(o);}
\end{tikzpicture}

{\small Krata $M_3$: trzy nieporównywalne elementy między $0$ i $1$.}
\end{center}

\begin{twierdzenie}[Krata podzbiorów jest zupełna]
\label{thm:powerset-complete-lattice}
Niech $X$ będzie dowolnym zbiorem. W porządku
$(\mathcal P(X),\subseteq)$ każda rodzina
$\mathcal A\subseteq\mathcal P(X)$ ma kresy
\[
 \sup\mathcal A=\bigcup_{A\in\mathcal A}A,
 \qquad \inf\mathcal A=\bigcap_{A\in\mathcal A}A,
\]
gdzie dla $\mathcal A=\varnothing$ przyjmujemy
$\bigcap_{A\in\mathcal A}A=X$. W szczególności
$(\mathcal P(X),\subseteq)$ jest kratą zupełną.
\end{twierdzenie}
\begin{proof}
Dla każdego $A\in\mathcal A$ mamy
$A\subseteq\bigcup_{B\in\mathcal A}B$, a jeżeli
$U\subseteq X$ zawiera każdy $A\in\mathcal A$, to
z definicji sumy zawiera także całą sumę. Suma jest więc
najmniejszym ograniczeniem górnym. Analogicznie
przecięcie zawiera się w każdym $A$, a każdy
$D\subseteq X$ zawarty w każdym $A$ zawiera się
w ich przecięciu, dlatego jest największym
ograniczeniem dolnym. Jeśli $\mathcal A$ jest pusta,
każdy podzbiór $X$ jest zarówno ograniczeniem górnym,
jak i dolnym pustej rodziny: najmniejszym w
$\mathcal P(X)$ jest $\varnothing$, największym $X$.
Z definicji sumy i pustego przecięcia wynika dokładnie
$\sup\varnothing=\varnothing$ i $\inf\varnothing=X$.
\end{proof}

\section{Porządki iloczynowe i leksykograficzne}\label{sec:order-products}
\begin{definicja}[Porządek współrzędny]
Niech $(P,\preceq_P)$ i $(Q,\preceq_Q)$ będą
zbiorami częściowo uporządkowanymi. Na $P\times Q$
określamy
\[
 (p,q)\preceq_\times(p',q')\iff
 p\preceq_P p'\ \text{oraz}\ q\preceq_Q q'.
\]
Nazywamy go \emph{porządkiem współrzędnym}.
\end{definicja}

\begin{twierdzenie}[Porządek współrzędny jest częściowy]
\label{thm:product-poset}
Niech $(P,\preceq_P)$ i $(Q,\preceq_Q)$ będą zbiorami
częściowo uporządkowanymi. Na $P\times Q$ określmy relację
$(p,q)\preceq_\times(p',q')$ wtedy i tylko wtedy, gdy
$p\preceq_P p'$ oraz $q\preceq_Q q'$.
Wtedy $\preceq_\times$ jest porządkiem częściowym na $P\times Q$.
\end{twierdzenie}
\begin{proof}
Dla $(p,q)$ zwrotność obu porządków daje
$p\preceq_P p$ i $q\preceq_Q q$,
czyli $(p,q)\preceq_\times(p,q)$.
Jeśli $(p,q)\preceq_\times(p',q')$ i odwrotnie,
to antysymetria na $P$ daje $p=p'$, a na $Q$ daje
$q=q'$. Z jednoznaczności współrzędnych par
(Twierdzenie~\ref{thm:kuratowski}) wynika
$(p,q)=(p',q')$. Jeśli mamy dwa kolejne porównania
$(p,q)\preceq_\times(p',q')$ oraz
$(p',q')\preceq_\times(p'',q'')$, to przechodniość
na $P$ daje $p\preceq_P p''$, a na $Q$ daje
$q\preceq_Q q''$, czyli zachodzi przechodniość na iloczynie.
\end{proof}

\begin{przyklad}[Nieporównywalność w produkcie]
Nawet gdy $P=Q=\R$ z porządkiem liniowym,
porządek współrzędny na $\R^2$ nie jest liniowy:
$(0,1)$ i $(1,0)$ są nieporównywalne.
Pierwszy punkt ma mniejszą pierwszą współrzędną,
ale większą drugą; żaden nie dominuje drugiego
w obu współrzędnych.
\end{przyklad}

\begin{definicja}[Porządek leksykograficzny]
Niech $(P,\preceq_P)$ i $(Q,\preceq_Q)$ będą
zbiorami liniowo uporządkowanymi, a $\prec_P,\prec_Q$
ich ostrymi odpowiednikami. Na $P\times Q$ ustalamy
\[
 (p,q)\preceq_{\rm lex}(p',q')\iff
 (p\prec_P p')\ \text{lub}\ (p=p'\text{ i }q\preceq_Q q').
\]
Najpierw porównujemy pierwszą współrzędną, a dopiero
przy remisie drugą.
\end{definicja}

\begin{twierdzenie}[Porządek leksykograficzny jest liniowy]
\label{thm:lex-order}
Niech $(P,\preceq_P)$ i $(Q,\preceq_Q)$ będą zbiorami
liniowo uporządkowanymi, a $p\prec_P p'$ niech oznacza
$p\preceq_P p'$ i $p\ne p'$. Na $P\times Q$ określmy
$(p,q)\preceq_{\rm lex}(p',q')$ wtedy i tylko wtedy, gdy
$p\prec_P p'$ lub ($p=p'$ i $q\preceq_Q q'$).
Wtedy $\preceq_{\rm lex}$ jest porządkiem liniowym na $P\times Q$.
\end{twierdzenie}
\begin{proof}
Zwrotność: $p=p$ i $q\preceq_Q q$.
Dla antysymetrii załóżmy oba porównania w przeciwnych
kierunkach. Jeżeli $p\prec_P p'$, to porównanie odwrotne
nie może zachodzić: wymagałoby $p'\prec_P p$ lub
$p'=p$, a obie możliwości są sprzeczne z asymetrią
$\prec_P$. Podobnie nie może zachodzić $p'\prec_P p$.
Pozostaje więc $p=p'$, a z obu porównań dostajemy
$q\preceq_Q q'$ i $q'\preceq_Q q$. Antysymetria
daje $q=q'$, zatem równość par.

Dla przechodniości załóżmy
$(p,q)\preceq_{\rm lex}(p',q')$ i
$(p',q')\preceq_{\rm lex}(p'',q'')$.
Jeśli $p\prec_P p'$, to z $p'\preceq_P p''$
i przechodniości otrzymujemy $p\preceq_P p''$.
Równość $p=p''$ jest niemożliwa, ponieważ
$p\prec_P p'$ i $p'\preceq_P p$ dawałyby sprzeczność;
stąd $p\prec_P p''$. Jeśli $p=p'$ i
$q\preceq_Q q'$, to mamy dwa przypadki: gdy
$p'\prec_P p''$, od razu $p\prec_P p''$;
a gdy $p'=p''$ i $q'\preceq_Q q''$, wynika
$p=p''$ i $q\preceq_Q q''$. We wszystkich przypadkach
mamy porównanie leksykograficzne pierwszej i trzeciej pary.

Liniowość: dla dowolnych $p,p'$ zachodzi dokładnie jedna
z możliwości $p\prec_P p'$, $p=p'$,
$p'\prec_P p$. W dwóch skrajnych przypadkach pary
są porównywalne z definicji. W przypadku środkowym
porównujemy $q,q'$ liniowym porządkiem na $Q$.
\end{proof}


\section{Gęstość porządku: między dwoma punktami jest trzeci}
\label{sec:order-density}

\emph{Gęstość} oznacza tutaj określoną własność \emph{porządku},
a nie to, że zbiór ma dużo elementów. Na przykład liczby całkowite
są nieskończone, lecz między $1$ i $2$ nie ma żadnej liczby całkowitej.
Natomiast między dowolnymi dwiema różnymi liczbami wymiernymi
można wskazać trzecią liczbę wymierną. Te dwa przykłady motywują definicję.

\begin{definicja}[Gęsty porządek liniowy]\label{def:dense-linear-order}
Niech $(P,<)$ będzie zbiorem liniowo uporządkowanym,
gdzie $<$ oznacza ostry porządek. Mówimy, że porządek $<$
jest \emph{gęsty}, jeśli
\[
 \forall a,b\in P\quad
 \bigl(a<b\ \Longrightarrow\ \exists c\in P\ (a<c\land c<b)\bigr).
\]
Słownie: \emph{dla każdej pary elementów, z których pierwszy jest
ściśle mniejszy od drugiego, istnieje element leżący ściśle między nimi}.
Nie żądamy, żeby między dwoma elementami znajdowało się dokładnie
jedno $c$, lecz przynajmniej jedno.
\end{definicja}

\begin{przyklad}[Dlaczego $\Q$ jest gęsto uporządkowany?]
Ustalmy dwie dowolne liczby wymierne $a<b$.
Zdefiniujmy $c:=(a+b)/2$.
Ponieważ suma liczb wymiernych i iloraz liczby wymiernej
przez $2$ są wymierne, $c\in\Q$.
Ponadto $b-a>0$, a stąd
$c-a=(b-a)/2>0$ oraz $b-c=(b-a)/2>0$.
Obie nierówności dają $a<c<b$.
Konstrukcja działa dla \emph{każdej} pary $a<b$, więc
$(\Q,<)$ jest gęstym porządkiem liniowym.
Dokładnie ten sam rachunek pokazuje gęstość zwykłego
porządku na $\R$.
\end{przyklad}

\begin{przyklad}[Nieskończoność nie wystarcza]
Porządek $(\Z,<)$ nie jest gęsty. Aby obalić definicję
z kwantyfikatorem „dla każdej pary”, wystarczy \emph{jedna}
para, dla której brak elementu pośredniego.
Weźmy $a=0$ i $b=1$. Nie istnieje $c\in\Z$
z $0<c<1$. Zatem porządek $\Z$ nie jest gęsty,
choć zbiór $\Z$ jest nieskończony.
Tak samo zbiór $\{0,1\}$ ze zwykłym porządkiem
nie jest gęsty. To pokazuje, że własność nie zależy
od liczby elementów, lecz od ich rozmieszczenia w porządku.
\end{przyklad}

\begin{definicja}[Podzbiór gęsty w porządku]\label{def:order-dense-subset}
Niech $(P,<)$ będzie liniowo uporządkowanym zbiorem
i niech $D\subseteq P$. Mówimy, że $D$ jest
\emph{gęsty w porządku $P$}, jeśli
\[
 \forall a,b\in P\quad
 (a<b\ \Longrightarrow\ \exists d\in D\ (a<d\land d<b)).
\]
Różnica wobec Definicji~\ref{def:dense-linear-order} polega
na zakresie kwantyfikatorów: krańce $a,b$ wybieramy z $P$,
a element pośredni musi należeć do wybranego podzbioru $D$.
\end{definicja}

\begin{przyklad}[Gęsty podzbiór a gęsty porządek]
Zbiór liczb wymiernych $\Q$ jest gęsty w $\R$.
Przypomnijmy uzasadnienie. Dla dowolnych $a,b\in\R$
spełniających $a<b$ wybieramy liczbę
naturalną $n$ na tyle dużą, że $n(b-a)>1$.
Można to zrobić dzięki \emph{własności Archimedesa}
liczb rzeczywistych: dla każdej liczby rzeczywistej
$u>0$ istnieje $n\in\N$ z $nu>1$.
Z archimedesowości $\R$ istnieje liczba całkowita większa
od $na$. Zbiór takich liczb całkowitych jest niepusty
i ograniczony od dołu, ma więc element najmniejszy $k$.
Z minimalności $k$ wynika $k-1\le na$, stąd
$na<k\le na+1<nb$.
Zatem liczba $d:=k/n$ jest wymierna i spełnia $a<d<b$.

Zauważmy, że nie każda definicja gęstości jest taka sama:
„porządek $P$ jest gęsty” wymaga elementu pośredniego w $P$;
„podzbiór $D$ jest gęsty w $P$” żąda, aby element pośredni
należał konkretnie do $D$.
\end{przyklad}

\begin{uwaga}[Gęstość a końce porządku]
Gęsty porządek może mieć element najmniejszy lub największy.
Na przykład przedział $[0,1]$ ze zwykłym porządkiem jest
gęsty, choć $0$ jest najmniejszy, a $1$ największy.
Definicja dotyczy wyłącznie par $a<b$ i tego, co
znajduje się \emph{między} nimi, a nie tego, czy
istnieją elementy poza nimi. W zbiorze pustym oraz
jednoelementowym warunek gęstości jest spełniony
w sposób pusty, ponieważ w ogóle nie występują pary $a<b$.
\end{uwaga}



\section{Dobre porządki i zasada najmniejszego kontrprzykładu}
\label{sec:order-well-order}
\begin{definicja}[Dobry porządek]\label{def:well-order}
Porządek liniowy $\preceq$ na zbiorze $W$ nazywamy
\emph{dobrym porządkiem}, jeśli każdy niepusty podzbiór
$A\subseteq W$ ma element najmniejszy (w porządku
odziedziczonym po $W$). Zbiór $W$ z takim porządkiem
nazywamy \emph{dobrze uporządkowanym}.
\end{definicja}

\begin{twierdzenie}[Leksykograficzny iloczyn dobrych porządków]
\label{thm:lex-well-orders-xliii}
Niech $(P,\preceq_P)$ i $(Q,\preceq_Q)$ będą zbiorami
dobrze uporządkowanymi. Na $P\times Q$ określmy
porządek leksykograficzny wzorem
$(p,q)\preceq_{\rm lex}(p',q')$ wtedy i tylko wtedy, gdy
$(p\preceq_P p'$ i $p\ne p')$ lub
$(p=p'$ i $q\preceq_Q q')$.
Wówczas jest to dobry porządek na $P\times Q$.
Jeśli jeden z czynników jest pusty, iloczyn jest pusty,
a teza pozostaje prawdziwa.
\end{twierdzenie}
\begin{proof}
Dzięki Twierdzeniu~\ref{thm:lex-order} porządek
leksykograficzny jest liniowy. Niech
$A\subseteq P\times Q$ będzie niepusty.
Zbiór pierwszych współrzędnych
$S=\{p\in P:\exists q\in Q\ ((p,q)\in A)\}$
jest niepustym podzbiorem $P$, więc ma najmniejszy
element $p_0$. Włókno
$T=\{q\in Q:(p_0,q)\in A\}$ jest niepuste z definicji $S$.
Ponieważ $Q$ jest dobrze uporządkowany, $T$ ma najmniejszy
element $q_0$. Para $(p_0,q_0)$ należy do $A$.
Dla dowolnego $(p,q)\in A$ mamy $p_0\preceq_P p$.
Jeśli $p_0\ne p$, to z $p_0\preceq_P p$ i definicji
porządku leksykograficznego wynika
$(p_0,q_0)\preceq_{\rm lex}(p,q)$. Jeśli $p_0=p$,
to $q\in T$, stąd $q_0\preceq_Q q$ i ponownie
$(p_0,q_0)\preceq_{\rm lex}(p,q)$.
Zatem para $(p_0,q_0)$ jest najmniejsza w $A$.
Nie wybieraliśmy jednocześnie elementów z wielu
niepustych podzbiorów: wskazaliśmy dwa najmniejsze
elementy kolejno. Aksjomat wyboru nie jest potrzebny.
\end{proof}
\begin{przyklad}[Element bez bezpośredniego poprzednika]
\label{ex:lex-no-predecessor-xliii}
Na $\N\times\N$ weźmy zwykły porządek leksykograficzny.
Każda para $(0,n)$ jest mniejsza od $(1,0)$, ale nie
ma wśród nich największej: większą od $(0,n)$ jest
$(0,n+1)$. Nie istnieje element bezpośrednio poprzedzający $(1,0)$:
każda para poniżej $(1,0)$ ma postać $(0,n)$,
a między nią i $(1,0)$ leży $(0,n+1)$.
Mimo to cały porządek jest dobry. Pokazuje to, że
dobry porządek nie musi dopuszczać dowodzenia przez
przechodzenie tylko od bezpośredniego poprzednika
do następnego elementu; potrzebna jest pełna postać
indukcji z Twierdzenia~\ref{thm:well-order-induction}.
\end{przyklad}



\begin{przyklad}[Dobre i niedobre porządki]
$\N=\{0,1,2,\ldots\}$ w zwykłym porządku jest dobrze
uporządkowany: każdy niepusty zbiór liczb naturalnych
ma najmniejszy element. Natomiast $\Z$ nie jest dobrze
uporządkowany, bo sam nie ma najmniejszego elementu.
Również $(0,1)\subseteq\R$ nie ma elementu
najmniejszego w zwykłym porządku: dla każdego $x\in(0,1)$
liczba $x/2$ nadal należy do $(0,1)$ i jest mniejsza.
Ważne: \emph{skończony} porządek liniowy jest zawsze dobry,
co wynika z Twierdzenia~\ref{thm:finite-poset-extrema}
i porównywalności dowolnej pary.
\end{przyklad}

\begin{twierdzenie}[Niezgodność dobrego uporządkowania z gęstością]\label{thm:well-order-not-dense}
Niech $(W,\preceq)$ będzie dobrze uporządkowanym zbiorem
zawierającym co najmniej dwa elementy. Dla $a,b\in W$
niech $a<b$ oznacza, że $a\preceq b$ i $a\ne b$.
Wtedy ostry porządek $<$
nie jest gęsty.
\end{twierdzenie}
\begin{proof}
Ponieważ $W$ jest niepusty i dobrze uporządkowany,
ma element najmniejszy $m$. Zbiór $W\setminus\{m\}$
jest niepusty, ponieważ założyliśmy istnienie co
najmniej dwóch elementów. Ma więc element najmniejszy $n$.
Wówczas $m<n$. Załóżmy, że istnieje $c\in W$
takie, że $m<c<n$.
Ponieważ $c\ne m$, mamy $c\in W\setminus\{m\}$.
Nierówność $c<n$ przeczy jednak temu, że $n$ jest
najmniejszym elementem tego zbioru.
Zatem między $m$ i $n$ nie ma elementu $W$, co obala
gęstość. Zwróćmy uwagę, że nie musieliśmy
numerować wszystkich elementów $W$ ani stosować
aksjomatu wyboru.
\end{proof}

\begin{twierdzenie}[Indukcja w zbiorze dobrze uporządkowanym]
\label{thm:well-order-induction}
Niech $(W,\preceq)$ będzie dobrze uporządkowanym zbiorem,
niech $v\prec w$ oznacza $v\preceq w$ i $v\ne w$,
a $P(w)$ niech będzie własnością elementów $w\in W$.
Załóżmy, że dla każdego $w\in W$ zachodzi implikacja:
\[
 \bigl(\forall v\in W\ (v\prec w\Rightarrow P(v))\bigr)
 \quad\Longrightarrow\quad P(w).
\]
Wówczas $P(w)$ zachodzi dla wszystkich $w\in W$.
\end{twierdzenie}
\begin{proof}
Przypuśćmy przeciwnie, że istnieje element $w\in W$,
dla którego $P(w)$ jest fałszywe. Zbiór
$C:=\{w\in W:\neg P(w)\}$ jest wtedy niepusty.
Ponieważ porządek jest dobry, $C$ ma element najmniejszy
$w_0$. Jeśli $v\prec w_0$, to $v\notin C$,
bo w przeciwnym razie $w_0$ nie byłby najmniejszym
elementem $C$. Zatem $P(v)$ jest prawdziwe dla
każdego $v\prec w_0$. Hipoteza indukcyjna z tezy
zastosowana do $w_0$ daje wtedy $P(w_0)$.
Jest to sprzeczne z $w_0\in C$. Stąd $C$ musi być pusty,
a wszystkie $w\in W$ spełniają $P(w)$.
\end{proof}

\begin{uwaga}[Dobrze ufundowany porządek a brak ciągów malejących]
W dalszej części teorii mnogości spotkamy ogólniejsze pojęcie
relacji dobrze ufundowanej, w której każdy niepusty podzbiór
ma element minimalny, ale porządek nie musi być liniowy.
W ZFC własność tę można też opisać przez brak nieskończonych
ściśle malejących ciągów. W samej ZF nie należy bez dodatkowej
zasady wyboru utożsamiać tych dwóch sformułowań: wyprowadzanie
z braku elementu minimalnego nieskończonego ciągu kolejnych
wyborów wymaga ostrożności. Rozwiniemy to szczegółowo przy
aksjomacie wyboru i rekurencji pozaskończonej.
\end{uwaga}


\begin{przyklad}[Dobrze ufundowany, ale nie liniowy porządek]
Na $P=\{a,b,c\}$ określmy porządek, w którym
$a\prec c$, $b\prec c$, natomiast $a$ i $b$ są
nieporównywalne. Każdy niepusty podzbiór $P$ ma element
minimalny: jeśli zawiera $a$ lub $b$, taki element jest
minimalny; jeśli zawiera wyłącznie $c$, minimalny jest $c$.
Nie jest to dobry porządek, ponieważ nie jest liniowy.
W szczególności w $\{a,b\}$ istnieją dwa elementy
minimalne, lecz nie ma elementu najmniejszego.
\end{przyklad}

\begin{uwaga}[Indukcja nie oznacza wyboru następnego elementu]
Dowód Twierdzenia~\ref{thm:well-order-induction} nie buduje
ciągu kolejnych elementów $W$. Wybiera jedynie najmniejszy
element zbioru kontrprzykładów, jeśli taki zbiór jest niepusty.
Dlatego dowód nie wymaga aksjomatu wyboru. Natomiast próba
wybrania po jednym mniejszym elemencie na każdym etapie
w dowolnym nieuporządkowanym zbiorze może ukrywać dodatkowe
założenie o wyborze.
\end{uwaga}

\begin{wniosek}[Wybór elementów z podzbiorów dobrze uporządkowanego zbioru]
\label{cor:lxv-wellorder-minimal-selection}
Niech $(W,\preceq)$ będzie dobrze uporządkowanym zbiorem.
Dla każdej rodziny $\mathcal A$ niepustych podzbiorów $W$
istnieje funkcja $c:\mathcal A\to W$ taka, że
$c(A)\in A$ dla każdego $A\in\mathcal A$.
Nie zakłada się aksjomatu wyboru.
\end{wniosek}
\begin{proof}
Dla $A\in\mathcal A$ dobry porządek daje dokładnie jeden
element najmniejszy $m_A$. Określmy
$c=\{(A,w)\in\mathcal A\times W:\ w\in A\ \land\
\forall v\in A\ (w\preceq v)\}$.
Jest to zbiór otrzymany przez wyróżnianie w
$\mathcal A\times W$. Dla każdego $A\in\mathcal A$
istnieje dokładnie jedna para $(A,w)$ w tej relacji,
a więc $c$ jest funkcją o dziedzinie $\mathcal A$.
Jej wartość $c(A)=m_A$ należy do $A$.
Jeśli $\mathcal A=\varnothing$, otrzymujemy funkcję pustą.
Nie zastosowano aksjomatu wyboru, ponieważ najmniejszy
element jest jednoznacznie wyznaczony przez porządek.
\end{proof}
\begin{przyklad}[Dlaczego ustalony dobry porządek jest istotny]
\label{ex:lxv-wellorder-minimal-selection}
Dla $W=\{0,1,2\}$ w zwykłym porządku oraz
$\mathcal A=\{\{1,2\},\{0,2\}\}$ otrzymujemy
$c(\{1,2\})=1$ i $c(\{0,2\})=0$.
Wybór tych wartości nie wymaga rozstrzygania między
kilkoma kandydatami. Bez ustalonego dobrego porządku
sama niepustość zbiorów z dowolnej rodziny nie daje
w ZF ogólnego twierdzenia o istnieniu funkcji wyboru.
\end{przyklad}

\section{Preporządek i iloraz: powrót do równoważności}
\label{sec:order-prequotient}
\begin{twierdzenie}[Porządek na ilorazie preporządku]
\label{thm:preorder-quotient}
Niech $P$ będzie zbiorem, a $\preceq$ preporządkiem na $P$.
Dla $a,b\in P$ określmy
\[
 a\sim b\quad\Longleftrightarrow\quad
 (a\preceq b\text{ i }b\preceq a).
\]
Wtedy $\sim$ jest relacją równoważności. Co więcej,
wzór
\[
 [a]\sqsubseteq[b]\quad\Longleftrightarrow\quad a\preceq b
\]
wyznacza dobrze określony porządek częściowy na $P/{\sim}$.
\end{twierdzenie}
\begin{proof}
Zwrotność $\sim$ wynika ze zwrotności $\preceq$.
Symetria wynika z symetrii spójnika „i”: jeśli
$a\preceq b$ i $b\preceq a$, to zachodzą również
$b\preceq a$ i $a\preceq b$. Jeżeli $a\sim b$ i
$b\sim c$, to mamy $a\preceq b\preceq c$ i
$c\preceq b\preceq a$. Przechodniość $\preceq$
daje $a\preceq c$ oraz $c\preceq a$, czyli $a\sim c$.

Sprawdźmy niezależność definicji $\sqsubseteq$ od
reprezentantów. Niech $[a]=[a']$ i $[b]=[b']$.
Na mocy Twierdzenia~\ref{thm:class-equality} zachodzą
$a'\preceq a$, $a\preceq a'$ oraz
$b'\preceq b$, $b\preceq b'$.
Jeśli $a\preceq b$, to przechodniość trzech
porównań $a'\preceq a\preceq b\preceq b'$ daje
$a'\preceq b'$. Argument odwrócony (z zamianą
reprezentantów) daje implikację odwrotną.
Zatem wzór określa jednoznaczną relację na zbiorze klas.

Zwrotność i przechodniość $\sqsubseteq$ dziedziczymy
z $\preceq$: jeśli $[a]\sqsubseteq[b]$ i
$[b]\sqsubseteq[c]$, to $a\preceq b\preceq c$,
więc $[a]\sqsubseteq[c]$.
Jeśli $[a]\sqsubseteq[b]$ i $[b]\sqsubseteq[a]$,
to $a\preceq b$ i $b\preceq a$, czyli $a\sim b$;
stąd $[a]=[b]$ na mocy
Twierdzenia~\ref{thm:class-equality}. To antysymetria.
\end{proof}

\begin{przyklad}[Porównywanie wartości bez rozróżniania znaków]
Na $\Z$ rozważmy preporządek $a\preceq b\iff |a|\le|b|$.
Mamy $a\sim b$ wtedy i tylko wtedy, gdy $|a|=|b|$.
Klasa $[0]=\{0\}$, a dla $n\ge1$ klasa $[n]$
zawiera dokładnie $n$ i $-n$. Na ilorazie powstaje
porządek liniowy klas
$[0]\sqsubseteq[1]\sqsubseteq[2]\sqsubseteq\cdots$.
Tak usuwamy dokładnie ten brak antysymetrii,
który napotkaliśmy na początku rozdziału.
\end{przyklad}

\begin{uwaga}[Przygotowanie do lematu Kuratowskiego--Zorna]
Twierdzenie~\ref{thm:finite-poset-extrema} zapewnia istnienie
elementu maksymalnego tylko dla skończonych niepustych porządków.
Nie wolno wyciągać z niego analogicznego wniosku dla dowolnych
porządków nieskończonych; na przykład $(\N,\le)$ nie ma
elementu maksymalnego. Później zbadamy warunki, przy których
\emph{nieskończony} porządek ma element maksymalny. Doprowadzi
to do lematu Kuratowskiego--Zorna i aksjomatu wyboru.
\end{uwaga}

\chapter{Funkcje: jednoznaczność przyporządkowania}\label{ch:functions}
\section{Od relacji do funkcji}
Relacja może przyporządkować jednemu argumentowi wiele wyników albo żadnego. Gdy podajemy wzór $x\mapsto x^2$, oczekujemy jednak, że dla każdego dopuszczalnego argumentu istnieje \emph{dokładnie jeden} wynik. To właśnie odróżnia funkcję od dowolnej relacji.
\begin{definicja}[Funkcja między zbiorami]\label{def:function}
Niech $A,B$ będą zbiorami. Aby określić \emph{odwzorowanie}
$f:A\to B$, podajemy jego dziedzinę $A$, przeciwdziedzinę $B$
i \emph{wykres} $G_f\subseteq A\times B$, spełniający warunek
\[
\forall a\in A\ \exists!b\in B\ ((a,b)\in G_f).
\]
Równoważnie, dane odwzorowanie można formalnie traktować jako
trójkę $(A,B,G_f)$. W dalszej części tekstu dla wygody
mówimy po prostu „funkcja $f:A\to B$” i operujemy jej wykresem,
ale pamiętamy o zadanej przeciwdziedzinie.
Piszemy $f(a)=b$ wtedy i tylko wtedy, gdy $(a,b)\in G_f$.
\emph{Zbiorem wartości} jest $f[A]:=\{f(a):a\in A\}\subseteq B$.
Symbol $B^A$ oznacza zbiór wszystkich funkcji z $A$ do $B$.
\end{definicja}

\begin{uwaga}[Istnienie zbioru wszystkich funkcji]\label{rem:function-set-zf-xliv}
Zapis $B^A$ nie oznacza jedynie nieformalnej kolekcji. Ponieważ
iloczyn $A\times B$ jest zbiorem (zob. konstrukcję w
Rozdziale~\ref{ch:relations}), można utworzyć zbiór
$\mathcal P(A\times B)$ jego podzbiorów. Następnie przez wyróżnianie
otrzymujemy
\[
 B^A=\{G\in\mathcal P(A\times B):
 \forall a\in A\ \exists!b\in B\ ((a,b)\in G)\}.
\]
Jest to zapis zbioru \emph{wykresów}. Jeśli odwzorowania formalizujemy
jako trójki $(A,B,G)$, to ustalone $A,B$ pozwalają jednoznacznie
odtworzyć trójkę z wykresu; dlatego oba ujęcia dają ten sam zbiór funkcji
po naturalnym utożsamieniu. Nie stosujemy tutaj aksjomatu wyboru.
\end{uwaga}

\begin{twierdzenie}[Kiedy istnieje funkcja między zadanymi zbiorami?]
\label{thm:function-existence-xliv}
Niech $A$ i $B$ będą zbiorami. Funkcja $A\to B$ istnieje wtedy
i tylko wtedy, gdy $A=\varnothing$ lub $B\ne\varnothing$.
Gdy $A=\varnothing$, istnieje dokładnie jedna taka funkcja, niezależnie
od $B$.
\end{twierdzenie}
\begin{proof}
Jeśli $A=\varnothing$, relacja pusta jest podzbiorem
$A\times B=\varnothing$. Warunek istnienia dokładnie jednej wartości
dla każdego argumentu jest spełniony, ponieważ w $A$ nie ma
argumentów. Innej relacji zawartej w $A\times B$ nie ma.
Jeśli $A\ne\varnothing$ i $f:A\to B$ istnieje, weźmy $a_0\in A$.
Wtedy $f(a_0)\in B$, a więc $B\ne\varnothing$.
Odwrotnie, jeśli $B\ne\varnothing$, ustalmy jeden element $b_0\in B$.
Przyporządkowanie $a\mapsto b_0$ dla wszystkich $a\in A$ jest funkcją.
Wybór jednego elementu z jednego niepustego zbioru nie wymaga aksjomatu
wyboru.
\end{proof}

\begin{uwaga}
Przeciwdziedzina jest częścią danych odwzorowania, ale \emph{nie musi
wynikać z samego wykresu}. Ten sam wzór $f(x)=x^2$ i ten sam wykres
mogą być rozważane raz jako odwzorowanie $\R\to\R$, a raz jako
odwzorowanie $\R\to[0,\infty)$; tylko drugie jest surjekcją.
Jeśli ustalimy tę samą dziedzinę i przeciwdziedzinę, dwa odwzorowania
$f,g:A\to B$ są równe wtedy i tylko wtedy, gdy $f(a)=g(a)$
dla każdego $a\in A$. Jest to równoznaczne z równością ich wykresów.
W konwencji trójek odwzorowania o różnych przeciwdziedzinach
są różnymi danymi, nawet gdy mają identyczny wykres.
\end{uwaga}
\begin{przyklad}[Relacja niebędąca funkcją]
Niech $A=\{1,2\}$, $B=\{u,v\}$. Relacja $\{(1,u),(1,v),(2,u)\}$ nie jest funkcją $A\to B$, ponieważ argumentowi $1$ przypisuje dwa różne wyniki. Relacja $\{(1,u)\}$ również nie jest funkcją $A\to B$: argument $2$ nie ma obrazu. Natomiast $\{(1,u),(2,u)\}$ jest funkcją, mimo że oba argumenty mają ten sam obraz.
\end{przyklad}

\begin{uwaga}[Trzy obiekty, których nie należy mylić]
Dla $f:A\to B$ rozróżniamy: argument $a\in A$,
jego pojedynczą wartość $f(a)\in B$ oraz zbiór wszystkich
wartości $f[A]\subseteq B$. Na przykład dla $f(x)=x^2$
na $A=\{-2,-1,0,1,2\}$ mamy $f(-2)=4$ oraz
$f[A]=\{0,1,4\}$. W szczególności zbiór wartości może mieć
mniej elementów niż dziedzina i nie musi być równy przeciwdziedzinie.
\end{uwaga}

\begin{uwaga}[Zależność iniekcji i surjekcji od danych]
Iniekcyjność zależy tylko od dziedziny i przypisanych wartości,
podczas gdy surjekcyjność zależy również od przeciwdziedziny.
Zmniejszenie przeciwdziedziny do $f[A]$ zawsze daje surjekcję
$A\to f[A]$. Nie wolno jednak w ten sposób automatycznie
zmieniać funkcji bez odnotowania zmiany jej danych.
\end{uwaga}



\begin{twierdzenie}[Liczba funkcji między zbiorami skończonymi]
\label{thm:finite-function-count-xliv}
Niech $A,B$ będą zbiorami skończonymi, $n,m\in\N$ oraz
$|A|=n$, $|B|=m$. Wtedy liczba funkcji $A\to B$ wynosi $m^n$. Przyjmujemy tu konwencję
$0^0=1$, wyrażającą istnienie jednej funkcji $\varnothing\to\varnothing$.
\end{twierdzenie}
\begin{proof}
Dla $n=0$ zbiór $A$ jest pusty. Z
Twierdzenia~\ref{thm:function-existence-xliv} istnieje dokładnie jedna
funkcja $A\to B$, zgodnie z wartością $m^0=1$, również dla $m=0$.
Załóżmy, że wzór wykazano dla dziedzin $n$-elementowych.
Niech $|A|=n+1$. Wybierzmy $a_0\in A$ i oznaczmy
$A'=A\setminus\{a_0\}$, wobec czego $|A'|=n$.
Funkcję $f:A\to B$ wyznaczają dokładnie dwie dane:
jej ograniczenie $f|_{A'}:A'\to B$ oraz wartość $f(a_0)\in B$.
Odwrotnie, dowolna taka para danych określa jedną funkcję na $A$,
ponieważ $A'$ i $\{a_0\}$ są rozłączne. Istnieje $m^n$ możliwości
ograniczenia i $m$ możliwości wartości w $a_0$; razem daje to
$m^n\cdot m=m^{n+1}$. Jest to rozumowanie skończone i nie
odwołuje się do aksjomatu wyboru.
\end{proof}
\begin{przyklad}[Ręczne liczenie funkcji]
Dla $A=\{1,2\}$ i $B=\{u,v,w\}$ funkcję określa para
$(f(1),f(2))$. Możliwe pary to
$(u,u),(u,v),(u,w),(v,u),(v,v),(v,w),(w,u),(w,v),(w,w)$.
Jest ich $9=3^2$. Gdy $B=\varnothing$, nie ma żadnej funkcji
z tego niepustego $A$; gdy $A=\varnothing$, istnieje jedna
funkcja nawet dla $B=\varnothing$.
\end{przyklad}

\section{Iniekcje, surjekcje i bijekcje}
\begin{definicja}[Trzy podstawowe własności]\label{def:inj-surj-bij}
Niech $f:A\to B$ będzie funkcją. Mówimy, że $f$ jest \emph{iniekcją}, jeśli $\forall a,a'\in A\ (f(a)=f(a')\Rightarrow a=a')$. Mówimy, że jest \emph{surjekcją}, jeśli $\forall b\in B\ \exists a\in A\ (f(a)=b)$. Jest \emph{bijekcją}, jeśli jest jednocześnie iniekcją i surjekcją.
\end{definicja}
\begin{przyklad}
Funkcja $f:\Z\to\Z$, $f(n)=2n$, jest iniekcją: z $2m=2n$ wynika $m=n$. Nie jest surjekcją, gdyż liczba $1$ nie ma przeciwobrazu. Funkcja $g:\Z\to\{0,1\}$, przypisująca resztę z dzielenia przez $2$, jest surjekcją, lecz nie jest iniekcją, bo $g(0)=g(2)$. Funkcja $h:\Z\to\Z$, $h(n)=n+1$, jest bijekcją: różne argumenty mają różne obrazy, a dla każdego $k\in\Z$ zachodzi $h(k-1)=k$.
\end{przyklad}

\begin{przyklad}[Sprawdzamy warunki po kolei]
Niech $A=\{1,2,3\}$, $B=\{a,b\}$ i
$f(1)=a$, $f(2)=a$, $f(3)=b$.
Aby sprawdzić iniekcję, porównujemy wartości dla różnych
argumentów: $f(1)=f(2)$, mimo że $1\ne2$; funkcja nie
jest więc iniekcją. Aby sprawdzić surjekcję, bierzemy
po kolei elementy przeciwdziedziny: $a=f(1)$ i $b=f(3)$,
więc każdy z nich jest osiągany i funkcja jest surjekcją.
Jeśli zmienimy przeciwdziedzinę na $\{a,b,c\}$, zachowując
te same wartości, funkcja przestanie być surjekcją, bo $c$
nie ma przeciwobrazu. Nie zmieniły się wzory, zmieniły się
dane określające funkcję.
\end{przyklad}

\begin{twierdzenie}[Charakteryzacja iniekcji przez włókna]\label{thm:fibres}
Niech $A,B$ będą zbiorami, a $f:A\to B$ funkcją.
Dla $b\in B$ zdefiniujmy włókno $f^{-1}[\{b\}]:=\{a\in A:f(a)=b\}$. Funkcja $f$ jest iniekcją wtedy i tylko wtedy, gdy każde włókno ma co najwyżej jeden element. Jest surjekcją wtedy i tylko wtedy, gdy każde włókno jest niepuste.
\end{twierdzenie}
\begin{proof}
Jeśli $f$ jest iniekcją i $a,a'$ należą do tego samego włókna, to $f(a)=b=f(a')$, a więc $a=a'$. Odwrotnie, gdy $f(a)=f(a')=b$, oba argumenty leżą we włóknie nad $b$; założenie o co najwyżej jednym elemencie daje $a=a'$. Druga równoważność wynika bezpośrednio z rozwinięcia definicji surjekcji: $b$ należy do zbioru wartości dokładnie wtedy, gdy istnieje $a$ z $f(a)=b$.
\end{proof}

\begin{przyklad}[Włókna w poprzednim przykładzie]
Dla funkcji z poprzedniego przykładu
$f^{-1}[\{a\}]=\{1,2\}$ oraz $f^{-1}[\{b\}]=\{3\}$.
Włókno nad $a$ ma dwa elementy, więc funkcja nie jest iniekcją.
Obydwa włókna nad elementami $B=\{a,b\}$ są niepuste,
więc funkcja jest surjekcją. W ten sposób twierdzenie o włóknach
pozwala rozpoznać obie własności bez zgadywania.
\end{przyklad}


\begin{twierdzenie}[Włókna rozcinają dziedzinę]\label{thm:fiber-partition-xliv}
Niech $f:A\to B$ będzie funkcją i niech $C=f[A]$. Rodzina
$\{f^{-1}[\{b\}]:b\in C\}$ jest podziałem $A$: każde jej włókno
jest niepuste, dwa różne włókna są rozłączne, a ich suma jest równa $A$.
Dla $A=\varnothing$ otrzymujemy pustą rodzinę, która jest podziałem
pustego zbioru w rozumieniu Definicji~\ref{def:partition}.
\end{twierdzenie}
\begin{proof}
(1) Jeśli $b\in C$, to z definicji obrazu istnieje $a\in A$
taki, że $f(a)=b$. Stąd $a\in f^{-1}[\{b\}]$ i włókno jest niepuste.

(2) Jeśli $b,c\in C$ i $b\ne c$, to żaden $a\in A$ nie może
należeć do obu włókien: dawałoby to jednocześnie $f(a)=b$ i $f(a)=c$,
wbrew $b\ne c$.

(3) Każdy $a\in A$ należy do włókna nad $f(a)$, gdyż
$f(a)\in C$. Zatem suma włókien zawiera $A$, a wszystkie są jego
podzbiorami; jest więc równa $A$.
\end{proof}
\begin{przyklad}[Włókna reszty z dzielenia]
Niech $f:\{0,1,2,3,4,5,6,7\}\to\{0,1,2\}$ przypisuje resztę
z dzielenia przez $3$. Włókna to
$\{0,3,6\}$, $\{1,4,7\}$ i $\{2,5\}$.
Każde jest niepuste, żadne dwa nie mają wspólnego elementu i razem
dają wszystkie osiem argumentów. Funkcja jest surjekcją, ale nie
iniekcją, bo na przykład $f(0)=f(3)$.
\end{przyklad}

\begin{figure}[htbp]
\centering
\begin{tikzpicture}[>=Stealth, every node/.style={font=\small}, scale=.98]
 \node[draw,rounded corners,minimum width=1.6cm,minimum height=3.05cm] (dom) at (0,0) {};
 \node[draw,rounded corners,minimum width=1.4cm,minimum height=3.05cm] (cod) at (4.5,0) {};
 \node at (0,1.84) {$A$}; \node at (4.5,1.84) {$B$};
 \node (a1) at (0,.94) {$a_1$}; \node (a2) at (0,.18) {$a_2$};
 \node (a3) at (0,-.58) {$a_3$}; \node (a4) at (0,-1.28) {$a_4$};
 \node (b1) at (4.5,.67) {$b_1$}; \node (b2) at (4.5,-.60) {$b_2$};
 \draw[->,blue!65!black] (a1.east) -- (b1.west);
 \draw[->,blue!65!black] (a2.east) -- (b1.west);
 \draw[->,teal!65!black] (a3.east) -- (b2.west);
 \draw[->,teal!65!black] (a4.east) -- (b2.west);
 \node[align=center,anchor=north] at (2.25,-1.85) {$f^{-1}[\{b_1\}]=\{a_1,a_2\}$,\quad
 $f^{-1}[\{b_2\}]=\{a_3,a_4\}$};
\end{tikzpicture}
\caption{Włókna funkcji rozcinają dziedzinę na klasy argumentów o tej samej wartości.}
\label{fig:fiber-partition-xliv}
\end{figure}

\section{Funkcje stałe, identycznościowe i charakterystyczne}
\begin{definicja}
Dla zbioru $A$ \emph{funkcja identycznościowa} $\id_A:A\to A$ jest określona przez $\id_A(a)=a$. Funkcja $f:A\to B$ jest \emph{stała}, gdy istnieje $b_0\in B$ takie, że $f(a)=b_0$ dla wszystkich $a\in A$. Dla $S\subseteq A$ \emph{funkcja charakterystyczna} $\chi_S:A\to\{0,1\}$ jest określona przez $\chi_S(a)=1$ dla $a\in S$ oraz $\chi_S(a)=0$ dla $a\notin S$.
\end{definicja}
\begin{twierdzenie}[Zbiory i funkcje zero-jedynkowe]\label{thm:characteristic-bijection}
Dla dowolnego zbioru $A$ przyporządkowanie $S\mapsto\chi_S$ jest bijekcją między $\mathcal P(A)$ i $\{0,1\}^A$.
\end{twierdzenie}
\begin{proof}
Każdy podzbiór $S$ wyznacza jednoznaczną funkcję charakterystyczną. Jeśli $\chi_S=\chi_T$, to dla każdego $a\in A$ zachodzi $a\in S\Leftrightarrow\chi_S(a)=1\Leftrightarrow\chi_T(a)=1\Leftrightarrow a\in T$, stąd $S=T$. Jest więc iniekcja. Niech teraz $g:A\to\{0,1\}$. Definiujemy $S:=\{a\in A:g(a)=1\}$. Dla każdego $a$ mamy $\chi_S(a)=g(a)$, bo $g(a)$ może być tylko zerem lub jedynką. Zatem $g=\chi_S$, co dowodzi surjekcji.
\end{proof}
\begin{uwaga}[Przypadek pustej dziedziny]
Dla każdego $B$ istnieje dokładnie jedna funkcja $\varnothing\to B$: relacja pusta. Jeżeli $A\ne\varnothing$, funkcja $A\to\varnothing$ nie istnieje, bo dla dowolnego $a\in A$ nie można wskazać wartości w pustej przeciwdziedzinie. Funkcja $\varnothing\to B$ jest zawsze iniekcją, a surjekcją tylko gdy $B=\varnothing$. Powyższe wynika z definicji, nie z wyjątku od niej.
\end{uwaga}

\begin{uwaga}[Funkcje stałe na pustym zbiorze]
Przy definicji funkcji stałej przyjęliśmy warunek
$\exists b_0\in B\ \forall a\in A\ f(a)=b_0$. Wobec tego
jedyna funkcja $\varnothing\to\varnothing$ \emph{nie jest stała}
w tej konwencji, ponieważ nie można wskazać $b_0$.
Przy spotykanej niekiedy definicji „$f(a)=f(a')$ dla wszystkich
$a,a'\in A$” funkcja ta byłaby stała. Dla niepustego $A$
obie definicje są równoważne. Wyraźne podanie konwencji usuwa
pozorną sprzeczność.
\end{uwaga}

\section{Odwracalność i bijekcje}
\begin{definicja}[Funkcja odwrotna]
Dla $f:A\to B$ funkcję $g:B\to A$ nazywamy \emph{odwrotną} do $f$, jeżeli $g(f(a))=a$ dla każdego $a\in A$ oraz $f(g(b))=b$ dla każdego $b\in B$.
\end{definicja}
\begin{twierdzenie}[Kryterium odwracalności]\label{thm:inverse-iff-bij}
Niech $A,B$ będą zbiorami, a $f:A\to B$ funkcją. Wówczas $f$ ma funkcję odwrotną wtedy i tylko wtedy, gdy jest bijekcją. Wówczas odwrotna jest jednoznaczna i oznaczamy ją $f^{-1}:B\to A$.
\end{twierdzenie}
\begin{proof}
Załóżmy, że $g$ jest odwrotna. Jeśli $f(a)=f(a')$, zastosowanie $g$ daje $g(f(a))=g(f(a'))$, czyli $a=a'$, zatem $f$ jest iniekcją. Dla dowolnego $b\in B$ element $a=g(b)$ spełnia $f(a)=b$, więc $f$ jest surjekcją.

Odwrotnie, niech $f$ będzie bijekcją. Dla każdego $b\in B$ surjekcja zapewnia istnienie $a\in A$ z $f(a)=b$, a iniekcja zapewnia jednoznaczność takiego $a$. Definiujemy $g(b)$ jako ten jedyny element. Wtedy z konstrukcji $f(g(b))=b$, a dla każdego $a$ jednoznaczność przeciwobrazu $f(a)$ daje $g(f(a))=a$. Jeżeli $h$ jest inną odwrotną, to dla $b\in B$ mamy $h(b)=h(f(g(b)))=g(b)$. Zatem $h=g$.
\end{proof}
\begin{uwaga}
Symbol $f^{-1}$ ma dwa zastosowania: przy bijekcji oznacza funkcję odwrotną, a w zapisie $f^{-1}[C]$ dla dowolnej funkcji oznacza przeciwobraz zbioru $C$. Drugi zapis nie wymaga odwracalności. Rozwiniemy to w rozdziale~\ref{ch:function-operations}.
\end{uwaga}


\begin{uwaga}[Notacja złożenia]
Jeżeli $f:A\to B$ i $g:B\to C$, oznaczamy przez $g\circ f:A\to C$
funkcję $a\mapsto g(f(a))$. Formalną definicję i prawa składania
podamy w Rozdziale~\ref{ch:function-operations}. Zapis
$g\circ f=\id_A$ znaczy więc dokładnie, że $g(f(a))=a$
dla wszystkich $a\in A$.
\end{uwaga}

\section{Odwrotność lewostronna i prawostronna}
\begin{definicja}[Jednostronne odwrotności]
Niech $f:A\to B$. Funkcję $g:B\to A$ nazywamy
\emph{lewostronnie odwrotną} do $f$, jeżeli
$g\circ f=\id_A$. Nazywamy ją \emph{prawostronnie odwrotną},
jeżeli $f\circ g=\id_B$. Odwrotność w zwykłym sensie
spełnia oba warunki jednocześnie.
\end{definicja}
\begin{twierdzenie}[Co dają jednostronne odwrotności?]
Niech $A,B$ będą zbiorami, a $f:A\to B$ funkcją.
Jeśli $f$ ma odwrotność lewostronną, to jest iniekcją.
Jeśli ma odwrotność prawostronną, to jest surjekcją.
\end{twierdzenie}
\begin{proof}
Niech $g\circ f=\id_A$ i załóżmy $f(a)=f(a')$.
Po zastosowaniu funkcji $g$ mamy
$g(f(a))=g(f(a'))$, czyli $a=a'$. To jest definicja iniekcji.
Jeśli natomiast $f\circ g=\id_B$, to dla dowolnego $b\in B$
element $a=g(b)$ należy do $A$ i spełnia
$f(a)=f(g(b))=b$. Zatem każdy element $B$ ma przeciwobraz,
co jest definicją surjekcji.
\end{proof}
\begin{uwaga}[Dlaczego nie piszemy od razu odwrotności obu implikacji?]
W teorii mnogości bez aksjomatu wyboru nie każda surjekcja
musi mieć odwrotność prawostronną. Takie odwrócenie wymaga
wybrania po jednym elemencie z każdego włókna funkcji.
Szczegółowo wrócimy do tego problemu przy aksjomacie wyboru.
\end{uwaga}


\begin{twierdzenie}[Dokładne kryterium odwrotności lewostronnej]
\label{thm:left-inverse-converse-xliv}
Niech $A,B$ będą zbiorami, a $f:A\to B$ funkcją.
Istnieje funkcja $g:B\to A$ spełniająca $g\circ f=\id_A$
wtedy i tylko wtedy, gdy $f$ jest iniekcją oraz
$(A\ne\varnothing\ \text{lub}\ B=\varnothing)$.
\end{twierdzenie}
\begin{proof}
Konieczność iniekcyjności wynika z poprzedniego twierdzenia.
Ponadto, jeśli $A=\varnothing$ i $B\ne\varnothing$, nie istnieje
żadna funkcja $B\to A$, co wykazaliśmy w
Twierdzeniu~\ref{thm:function-existence-xliv}. Zatem w tej sytuacji
odwrotność lewostronna nie może istnieć.

Załóżmy teraz, że $f$ jest iniekcją. Jeśli $B=\varnothing$,
warunek istnienia $f:A\to B$ wymusza $A=\varnothing$.
Jedyne funkcje w obie strony mają wtedy puste wykresy,
a ich złożenie jest $\id_{\varnothing}$.
Jeżeli $A\ne\varnothing$, ustalmy $a_0\in A$ i zdefiniujmy
$g:B\to A$ przez
\[
 g(b)=\begin{cases}
 a,&\text{gdy }b=f(a)\text{ dla pewnego }a\in A,\\
 a_0,&\text{gdy }b\notin f[A].
 \end{cases}
\]
Iniekcyjność $f$ zapewnia, że w pierwszym przypadku
$a$ jest jednoznaczny. Oba przypadki się wykluczają i pokrywają
całe $B$. Dla każdego $a\in A$ mamy $g(f(a))=a$,
zatem $g\circ f=\id_A$. Do konstrukcji użyliśmy tylko jednego
ustalonego elementu $a_0$, nie aksjomatu wyboru.
\end{proof}

\begin{przyklad}[Dlaczego pustej dziedziny nie wolno pominąć]
Jedyna funkcja $f:\varnothing\to\{0\}$ jest iniekcją.
Gdyby miała odwrotność lewostronną $g:\{0\}\to\varnothing$,
to musiałaby istnieć wartość $g(0)\in\varnothing$, co jest
niemożliwe. Dla porównania, iniekcja
$f:\{0\}\to\{0,1\}$, $f(0)=0$, ma odwrotność lewostronną
$g(0)=g(1)=0$; brak surjekcyjności nie przeszkadza.
\end{przyklad}

\begin{twierdzenie}[Surjekcje na zbiór skończony mają przekrój]
\label{thm:finite-section-xliv}
Niech $f:A\to B$ będzie surjekcją, a $B$ zbiorem skończonym.
Wówczas istnieje funkcja $g:B\to A$ spełniająca
$f\circ g=\id_B$.
\end{twierdzenie}
\begin{proof}
Jeśli $B=\varnothing$, surjekcja $f:A\to B$ może istnieć
tylko przy $A=\varnothing$; jedyna funkcja pusta $g$
spełnia tezę.
Jeśli $B\ne\varnothing$, wypiszmy jego elementy jako
$B=\{b_1,\dots,b_n\}$, gdzie $n\ge1$.
Surjekcyjność daje dla każdego $i$ niepuste włókno
$f^{-1}[\{b_i\}]$. Wybieramy po jednym elemencie $a_i$
z tych \emph{skończenie wielu} zbiorów: najpierw z pierwszego,
potem z drugiego, aż do $n$-tego. Skończony proces wyboru
można formalnie uzasadnić indukcją po $n$ i nie wymaga aksjomatu
wyboru. Niech $g(b_i)=a_i$. Ponieważ wszystkie $b_i$ są różne,
 $g$ jest dobrze określoną funkcją, a
$(f\circ g)(b_i)=f(a_i)=b_i$ dla każdego $i$.
\end{proof}
\begin{przyklad}[Konkretny przekrój]
Funkcja $f:\{1,2,3,4\}\to\{u,v\}$ dana przez
$f(1)=f(2)=u$, $f(3)=f(4)=v$ jest surjekcją.
Przyporządkowanie $g(u)=2$, $g(v)=4$ jest jej prawostronną
odwrotnością, choć $g$ nie jest odwrotnością dwustronną:
$g(f(1))=2\ne1$.
\end{przyklad}


\begin{twierdzenie}[Przekroje surjekcji a aksjomat wyboru]
\label{thm:surjections-sections-ac-lxvi}
W teorii ZF następujące zdania są równoważne:
\begin{enumerate}[label=(\roman*)]
\item aksjomat wyboru: dla każdego zbioru indeksów $I$ i każdej
rodziny $(X_i)_{i\in I}$ niepustych zbiorów istnieje funkcja $c:I\to\bigcup_{i\in I}X_i$
spełniająca $c(i)\in X_i$ dla wszystkich $i\in I$;
\item dla każdej surjekcji $f:A\to B$ istnieje funkcja
$g:B\to A$ taka, że $f\circ g=\id_B$.
\end{enumerate}
\end{twierdzenie}
\begin{proof}
Załóżmy (i) i niech $f:A\to B$ będzie surjekcją.
Dla każdego $b\in B$ włókno $X_b=f^{-1}[\{b\}]$ jest
niepuste. Aksjomat wyboru zastosowany do rodziny
$(X_b)_{b\in B}$ daje $g(b)\in X_b$. Stąd
$f(g(b))=b$ dla każdego $b\in B$, czyli $f\circ g=\id_B$.
Jeśli $B=\varnothing$, z istnienia funkcji $A\to\varnothing$
wynika $A=\varnothing$ i wybieramy jedyną funkcję pustą.

Załóżmy teraz (ii). Niech $(X_i)_{i\in I}$ będzie rodziną
niepustych zbiorów. Zdefiniujmy rozłączną sumę
\[
 D=\{(i,x):i\in I,\ x\in X_i\}.
\]
Jest ona zbiorem: jej elementy należą do
$I\times\bigcup_{i\in I}X_i$, a wskazany warunek
wyznacza podzbiór przez wyróżnianie. Projekcja
$p:D\to I$, $p(i,x)=i$, jest surjekcją: dla ustalonego
$i\in I$ niepustość $X_i$ daje pewien $x\in X_i$,
a więc $p(i,x)=i$. Nie dokonujemy przy tym jednoczesnego
wyboru wszystkich $x$; dowodzimy tylko zdania
$\forall i\in I\ \exists x\in X_i$.
Z (ii) istnieje $g:I\to D$ z $p\circ g=\id_I$.
Dla każdego $i$ para $g(i)$ ma pierwszą współrzędną $i$,
więc istnieje dokładnie jeden $c(i)\in X_i$ taki, że
$g(i)=(i,c(i))$. Zbiór tych par określa funkcję $c$,
która jest funkcją wyboru. Dla $I=\varnothing$ obie
konstrukcje są puste i argument nadal działa.
\end{proof}
\begin{przyklad}[Dlaczego rozłączna suma jest potrzebna]
Niech $I=\{0,1\}$, $X_0=\{a,b\}$ i $X_1=\{b,c\}$.
Zwykła suma $X_0\cup X_1=\{a,b,c\}$ nie pamięta,
z którego zbioru pochodzi element $b$. Rozłączna suma
$D=\{(0,a),(0,b),(1,b),(1,c)\}$ zachowuje indeksy.
Przekrój $g(0)=(0,a)$, $g(1)=(1,c)$ projekcji
$p:D\to I$ wyznacza funkcję wyboru $c(0)=a$, $c(1)=c$.
Dla dowolnej rodziny niepustych zbiorów istnienie takiego
przekroju jest dokładnie treścią aksjomatu wyboru, a nie
wnioskiem z samej definicji surjekcji w ZF.
\end{przyklad}

\section{Funkcje jako szczególne relacje}
\begin{twierdzenie}[Kryterium wykresu]
Niech $A,B$ będą zbiorami, a $R\subseteq A\times B$ relacją. Relacja $R$ jest wykresem funkcji $A\to B$ wtedy i tylko wtedy, gdy dla każdego $a\in A$ zbiór $\{b\in B:(a,b)\in R\}$ jest jednoelementowy.
\end{twierdzenie}
\begin{proof}
Jeśli $R$ jest wykresem funkcji, to z definicji dla każdego $a$ istnieje dokładnie jedno $b$ z $(a,b)\in R$, czyli wskazany zbiór ma dokładnie jeden element. W drugą stronę jednoelementowość oznacza zarazem istnienie i jednoznaczność wartości dla każdego argumentu, a to jest definicja funkcji.
\end{proof}
\begin{przyklad}[Pierwiastek jako relacja]
Relacja $R=\{(x,y)\in\R^2:y^2=x\}$ nie jest wykresem funkcji $\R\to\R$: dla $x<0$ nie ma żadnego $y$, a dla $x>0$ istnieją dwa rozwiązania. Po ograniczeniu do $A=[0,\infty)$ i $B=[0,\infty)$ relacja jest wykresem funkcji pierwiastkowej: dla każdego $x\ge0$ istnieje dokładnie jeden nieujemny $y$ z $y^2=x$.
\end{przyklad}
\chapter{Obrazy, przeciwobrazy i działania na funkcjach}\label{ch:function-operations}
\section{Dlaczego trzeba rozróżniać obraz i przeciwobraz?}
Funkcja przypisuje wartości argumentom. Możemy jednak pytać nie tylko o pojedyncze $f(a)$, lecz także o wszystkie wartości przyjmowane na danym podzbiorze oraz o wszystkie argumenty, których wartości wpadają do zadanego podzbioru przeciwdziedziny. Te dwie konstrukcje zachowują się inaczej.
\begin{definicja}[Obraz i przeciwobraz]\label{def:image-preimage}
Niech $f:A\to B$, $S\subseteq A$, $T\subseteq B$. \emph{Obrazem} $S$ jest $f[S]:=\{f(a):a\in S\}\subseteq B$, a \emph{przeciwobrazem} $T$ jest $f^{-1}[T]:=\{a\in A:f(a)\in T\}\subseteq A$.
\end{definicja}
\begin{przyklad}
Dla $f:\R\to\R$, $f(x)=x^2$, mamy $f[[-2,1]]=[0,4]$, $f^{-1}[[1,4]]=[-2,-1]\cup[1,2]$ oraz $f^{-1}[(-\infty,0)]=\varnothing$. Przeciwobraz jest zdefiniowany, mimo że $f$ nie jest bijekcją.
\end{przyklad}

\begin{uwaga}[Proste testy i przypadki puste]
Dla każdego $f:A\to B$ mamy $f[\varnothing]=\varnothing$,
$f^{-1}[\varnothing]=\varnothing$, $f^{-1}[B]=A$ i
$f[A]\subseteq B$. Ostatnie zawieranie może być ostre.
Na przykład dla $f:\{1,2\}\to\{u,v,w\}$,
$f(1)=u$, $f(2)=v$, mamy $f[A]=\{u,v\}$, ale
$f^{-1}[\{w\}]=\varnothing$. Te równości wynikają z definicji,
a nie z odwracalności funkcji.
\end{uwaga}
\begin{twierdzenie}[Monotoniczność obrazów i przeciwobrazów]
\label{thm:images-monotone-xliv}
Dla funkcji $f:A\to B$ i podzbiorów $S,U\subseteq A$,
$T,V\subseteq B$ zachodzi
\[
 S\subseteq U\ \Longrightarrow\ f[S]\subseteq f[U],\qquad
 T\subseteq V\ \Longrightarrow\ f^{-1}[T]\subseteq f^{-1}[V].
\]
\end{twierdzenie}
\begin{proof}
Jeśli $b\in f[S]$, to istnieje $s\in S$ z $f(s)=b$.
Ponieważ $S\subseteq U$, mamy $s\in U$, a więc $b\in f[U]$.
Jeśli $a\in f^{-1}[T]$, to $f(a)\in T\subseteq V$,
stąd $a\in f^{-1}[V]$. W obu przypadkach dowiedliśmy
zawierania przez sprawdzenie dowolnego elementu.
\end{proof}

\section{Prawa obrazów i przeciwobrazów}
\begin{twierdzenie}[Suma, przecięcie i różnica]\label{thm:image-laws}
Niech $f:A\to B$, $S,U\subseteq A$ i $T,V\subseteq B$. Wtedy
\begin{align*}
f[S\cup U]&=f[S]\cup f[U], & f[S\cap U]&\subseteq f[S]\cap f[U],\\
f^{-1}[T\cup V]&=f^{-1}[T]\cup f^{-1}[V], &f^{-1}[T\cap V]&=f^{-1}[T]\cap f^{-1}[V],\\
f^{-1}[B\setminus T]&=A\setminus f^{-1}[T].
\end{align*}
Jeżeli $f$ jest iniekcją, to $f[S\cap U]=f[S]\cap f[U]$ oraz $f[S\setminus U]=f[S]\setminus f[U]$.
\end{twierdzenie}
\begin{proof}
Dla $b\in B$ warunek $b\in f[S\cup U]$ oznacza, że istnieje $a\in S\cup U$ z $f(a)=b$. Taki $a$ należy do $S$ lub $U$, czyli $b\in f[S]\cup f[U]$. Rozumowanie odwracalne dowodzi pierwszej równości. Jeśli $b\in f[S\cap U]$, istnieje $a\in S\cap U$ z $f(a)=b$, a więc $b$ należy do obu obrazów. To daje zawieranie.

Dla każdego $a\in A$ zachodzi
\[
a\in f^{-1}[T\cap V]\iff f(a)\in T\cap V\iff (f(a)\in T\ \land\ f(a)\in V),
\]
co jest równoważne $a\in f^{-1}[T]\cap f^{-1}[V]$. Analogicznie zastąpienie spójnika $\land$ przez $\lor$ daje równość dla sumy. Ponadto $a\in f^{-1}[B\setminus T]$ wtedy i tylko wtedy, gdy $f(a)\notin T$, czyli $a\notin f^{-1}[T]$.

Jeśli $f$ jest iniekcją i $b\in f[S]\cap f[U]$, istnieją $s\in S$, $u\in U$ z $f(s)=b=f(u)$. Iniekcja daje $s=u\in S\cap U$, stąd $b\in f[S\cap U]$. Dla różnicy: jeśli $b=f(s)$, $s\in S\setminus U$, to $b\notin f[U]$, bo $f(u)=f(s)$ wymuszałoby $u=s\in U$. Odwrotnie, jeśli $b\in f[S]\setminus f[U]$, wybieramy $s\in S$ z $f(s)=b$. Ponieważ $b\notin f[U]$, mamy $s\notin U$, a więc $b\in f[S\setminus U]$.
\end{proof}
\begin{przyklad}[Dlaczego obraz nie zachowuje przecięć]
Niech $A=\{0,1\}$, $B=\{*\}$ i $f(0)=f(1)=*$. Dla $S=\{0\}$ i $U=\{1\}$ zachodzi $S\cap U=\varnothing$, więc $f[S\cap U]=\varnothing$, podczas gdy $f[S]\cap f[U]=\{*\}$. To samo pokazuje, że bez iniekcji nie można ogólnie oczekiwać zachowania różnic.
\end{przyklad}


\begin{twierdzenie}[Dokładny warunek zachowania przecięcia przez obraz]
\label{thm:image-intersection-fibers-lxvi}
Niech $f:A\to B$ i $S,U\subseteq A$. Wówczas
$f[S\cap U]=f[S]\cap f[U]$ wtedy i tylko wtedy, gdy
dla każdego $b\in B$ spełniającego
$f^{-1}[\{b\}]\cap S\ne\varnothing$ oraz
$f^{-1}[\{b\}]\cap U\ne\varnothing$ zachodzi
$f^{-1}[\{b\}]\cap S\cap U\ne\varnothing$.
Warunek ten dotyczy ustalonych $S,U$ i nie wymaga
iniekcyjności całej funkcji.
\end{twierdzenie}
\begin{proof}
Oznaczmy $F_b=f^{-1}[\{b\}]$. Z definicji obrazu
$b\in f[S]$ wtedy i tylko wtedy, gdy istnieje $s\in S$
z $f(s)=b$, czyli $F_b\cap S\ne\varnothing$.
Podobnie $b\in f[U]$ jest równoważne
$F_b\cap U\ne\varnothing$, a $b\in f[S\cap U]$
jest równoważne $F_b\cap S\cap U\ne\varnothing$.
Zawsze $f[S\cap U]\subseteq f[S]\cap f[U]$.
Zatem równość tych zbiorów zachodzi dokładnie wtedy,
gdy każdy element prawej strony należy do lewej;
po zastąpieniu przynależności otrzymanymi właśnie
równoważnościami dostajemy warunek z tezy.
Wszystkie świadki wybierane są osobno dla ustalonego $b$;
nie używamy aksjomatu wyboru.
\end{proof}
\begin{przyklad}[Zachowanie przecięcia bez iniekcyjności]
Niech $A=\{0,1,2\}$, $B=\{u,v\}$ oraz
$f(0)=f(1)=u$, $f(2)=v$. Dla $S=\{0,2\}$ i
$U=\{0,1\}$ mamy $S\cap U=\{0\}$,
$f[S\cap U]=\{u\}$ oraz
$f[S]\cap f[U]=\{u,v\}\cap\{u\}=\{u\}$.
Funkcja nie jest iniekcją, a mimo to zachowuje to
konkretne przecięcie. Natomiast dla $S'=\{0\}$ i
$U'=\{1\}$ otrzymujemy $f[S'\cap U']=\varnothing$
i $f[S']\cap f[U']=\{u\}$: włókno $F_u=\{0,1\}$
spotyka oba zbiory, lecz nie ich przecięcie.
\end{przyklad}

\begin{twierdzenie}[Dowolne rodziny i ważny wyjątek pustej rodziny]
\label{thm:images-arbitrary-families-xliv}
Niech $f:A\to B$, $I$ będzie zbiorem indeksów,
$(S_i)_{i\in I}$ rodziną podzbiorów $A$, a
$(T_i)_{i\in I}$ rodziną podzbiorów $B$. Wówczas
\begin{align*}
 f\Big[\bigcup_{i\in I}S_i\Big]
   &=\bigcup_{i\in I}f[S_i],\\
 f^{-1}\Big[\bigcup_{i\in I}T_i\Big]
   &=\bigcup_{i\in I}f^{-1}[T_i],\\
 f^{-1}\Big[\bigcap_{i\in I}T_i\Big]
   &=\bigcap_{i\in I}f^{-1}[T_i],\\
 f\Big[\bigcap_{i\in I}S_i\Big]
   &\subseteq\bigcap_{i\in I}f[S_i].
\end{align*}
Dla pustej rodziny podzbiorów $A$ przyjmujemy
$\bigcap_{i\in\varnothing}S_i=A$, a dla pustej rodziny
podzbiorów $B$: $\bigcap_{i\in\varnothing}T_i=B$.
Jeśli $I\ne\varnothing$ i $f$ jest iniekcją, to ostatnie
zawieranie jest równością. Jeśli $I=\varnothing$, ostatnia
równość zachodzi dokładnie wtedy, gdy $f$ jest surjekcją.
\end{twierdzenie}
\begin{proof}
(1) Element $b$ należy do obrazu sumy wtedy i tylko wtedy,
gdy istnieją indeks $i\in I$ i $a\in S_i$ spełniające
$f(a)=b$. Jest to równoważne istnieniu $i\in I$ z
$b\in f[S_i]$, czyli przynależności do prawej strony.
Działa to także dla $I=\varnothing$: obie strony są puste.

(2) Element $a\in A$ należy do przeciwobrazu sumy wtedy
i tylko wtedy, gdy $f(a)\in T_i$ dla pewnego $i\in I$.
Jest to równoważne $a\in f^{-1}[T_i]$ dla pewnego $i$.

(3) Analogicznie, $a$ należy do przeciwobrazu przecięcia
wtedy i tylko wtedy, gdy $f(a)\in T_i$ dla \emph{każdego}
$i\in I$, co jest równoważne przynależności do wszystkich
$f^{-1}[T_i]$. Gdy $I$ jest pusty, oba warunki są
prawdziwe dla każdego $a\in A$, ponieważ $f(a)\in B$;
obie strony są więc równe $A$.

(4) Jeżeli $b=f(a)$ dla pewnego $a$ należącego do
wszystkich $S_i$, to dla każdego $i$ mamy $b\in f[S_i]$.
Stąd zawieranie. Jeśli $I\ne\varnothing$, $f$ jest
iniekcją i $b\in\bigcap_i f[S_i]$, to dla każdego $i$
istnieje $a_i\in S_i$ z $f(a_i)=b$. Nie musimy jednocześnie
wybierać wszystkich $a_i$: wybierzmy \emph{jeden} indeks
$i_0\in I$ i jego świadka $a\in S_{i_0}$. Dla dowolnego
$i\in I$ istnieje świadek $a_i\in S_i$; iniekcyjność
daje $a_i=a$, więc $a\in S_i$. Zatem $a$ należy do
wszystkich $S_i$ i $b\in f[\bigcap_i S_i]$.

Jeśli natomiast $I=\varnothing$, lewa strona ostatniego
zawierania jest $f[A]$, a prawa wynosi $B$. Ich równość
jest równoważna surjekcyjności $f$.
\end{proof}
\begin{przyklad}[Pusta rodzina nie może zostać pominięta]
Niech $A=\{1\}$, $B=\{u,v\}$ i $f(1)=u$.
Funkcja jest iniekcją. Dla $I=\varnothing$ otrzymujemy
$f[\bigcap_{i\in I}S_i]=f[A]=\{u\}$,
natomiast $\bigcap_{i\in I}f[S_i]=B=\{u,v\}$.
Zatem iniekcja \emph{nie wystarcza} do zachowania przecięcia
pustej rodziny, mimo że zachowuje przecięcia wszystkich
rodzin niepustych.
\end{przyklad}

\begin{twierdzenie}[Kiedy obraz zachowuje przecięcia i różnice?]
\label{thm:image-injective-characterization-xliv}
Dla funkcji $f:A\to B$ następujące trzy warunki są równoważne:
\begin{enumerate}[label=(\roman*)]
\item $f$ jest iniekcją;
\item dla wszystkich $S,U\subseteq A$ mamy
$f[S\cap U]=f[S]\cap f[U]$;
\item dla wszystkich $S,U\subseteq A$ mamy
$f[S\setminus U]=f[S]\setminus f[U]$.
\end{enumerate}
\end{twierdzenie}
\begin{proof}
Implikacje (i)$\Rightarrow$(ii) oraz (i)$\Rightarrow$(iii)
zostały udowodnione w Twierdzeniu~\ref{thm:image-laws}.
Aby wykazać (ii)$\Rightarrow$(i), załóżmy $f(a)=f(a')$.
Dla $S=\{a\}$ i $U=\{a'\}$ zbiór
$f[S]\cap f[U]$ jest niepusty. Z (ii) wynika, że
$S\cap U$ jest niepusty, ponieważ obraz pustego zbioru
jest pusty. Stąd $a=a'$.
Aby wykazać (iii)$\Rightarrow$(i), również załóżmy
$f(a)=f(a')$. Gdyby $a\ne a'$, to dla
$S=\{a\}$ i $U=\{a'\}$ mielibyśmy
$f[S\setminus U]=\{f(a)\}$, a zarazem
$f[S]\setminus f[U]=\varnothing$, co przeczy (iii).
Zatem $a=a'$ i $f$ jest iniekcją.
\end{proof}

\begin{twierdzenie}[Zależność obrazu i przeciwobrazu]\label{thm:image-preimage-adjunction}
Dla $f:A\to B$, $S\subseteq A$ i $T\subseteq B$ zachodzi
\[
f[S]\subseteq T\quad\Longleftrightarrow\quad S\subseteq f^{-1}[T].
\]
W szczególności $S\subseteq f^{-1}[f[S]]$ oraz $f[f^{-1}[T]]\subseteq T$. Pierwsze zawieranie jest równością dla wszystkich $S\subseteq A$ wtedy i tylko wtedy, gdy $f$ jest iniekcją; drugie zawieranie jest równością dla wszystkich $T\subseteq B$ wtedy i tylko wtedy, gdy $f$ jest surjekcją.
\end{twierdzenie}
\begin{proof}
Warunek $f[S]\subseteq T$ oznacza: dla każdego $a\in S$ mamy $f(a)\in T$. Jest to dokładnie warunek $a\in f^{-1}[T]$ dla każdego $a\in S$, czyli $S\subseteq f^{-1}[T]$. Podstawienie $T=f[S]$ daje pierwsze zawieranie, a podstawienie $S=f^{-1}[T]$ drugie.

Jeśli $f$ jest iniekcją i $a\in f^{-1}[f[S]]$, to $f(a)=f(s)$ dla pewnego $s\in S$, więc $a=s\in S$. Zatem zachodzi równość. Jeśli równość zachodzi dla każdego $S$, a $f(a)=f(a')$, to dla $S=\{a\}$ mamy $a'\in f^{-1}[f[S]]=S$, więc $a'=a$. Analogicznie, jeśli $f$ jest surjekcją, każdy $b\in T$ jest równy $f(a)$ dla pewnego $a$, który wtedy należy do $f^{-1}[T]$, więc $b\in f[f^{-1}[T]]$. Odwrotnie, stosując postulowaną równość do $T=B$, otrzymujemy $f[A]=B$, czyli surjekcję.
\end{proof}

\begin{twierdzenie}[Dokładne wzory na dwukrotne przechodzenie]
\label{thm:image-preimage-exact-xliv}
Niech $f:A\to B$, $S\subseteq A$ i $T\subseteq B$. Wtedy
\[
 f[f^{-1}[T]]=T\cap f[A],\qquad
 f^{-1}[f[S]]=
 \{a\in A:\exists s\in S\ \ f(a)=f(s)\}.
\]
Dla drugiej równości prawa strona jest sumą wszystkich włókien,
które spotykają $S$. W szczególności jest to najmniejszy
podzbiór $A$ będący sumą włókien $f$ i zawierający $S$.
\end{twierdzenie}
\begin{proof}
Element $b\in B$ należy do $f[f^{-1}[T]]$ wtedy i tylko
wtedy, gdy istnieje $a\in A$ spełniające $f(a)=b$ oraz
$f(a)\in T$. Pierwsza część daje $b\in f[A]$, druga
$b\in T$, zatem pierwsza równość wynika z ekstensjonalności.

Element $a\in A$ należy do $f^{-1}[f[S]]$ wtedy i tylko
wtedy, gdy $f(a)\in f[S]$, czyli istnieje $s\in S$ takie,
że $f(s)=f(a)$. Jest to druga równość.
Dla ostatniego stwierdzenia oznaczmy tę prawą stronę przez $C$.
Każdy $s\in S$ należy do $C$ (możemy wziąć tego samego $s$
jako świadka). Jeśli $a\in C$ i $f(a')=f(a)$, to ten sam
świadek $s$ dowodzi $a'\in C$, więc $C$ jest sumą włókien.
Jeśli $D\subseteq A$ jest dowolną sumą włókien zawierającą $S$,
a $a\in C$, to istnieje $s\in S\subseteq D$ z $f(a)=f(s)$.
Całe włókno zawierające $s$ leży w $D$, zatem $a\in D$.
Dowodzi to minimalności $C$.
\end{proof}
\begin{przyklad}[Przeciwobraz obrazu jest zwykle większy]
Dla $f:\R\to\R$, $f(x)=x^2$, weźmy $S=\{2\}$
i $T=\{-1,1,4\}$. Wówczas $f[S]=\{4\}$, lecz
$f^{-1}[f[S]]=\{-2,2\}$: do zbioru dochodzi argument $-2$
o tej samej wartości funkcji. Ponadto $f^{-1}[T]=\{-2,-1,1,2\}$,
zatem $f[f^{-1}[T]]=\{1,4\}=T\cap[0,\infty)$.
Liczba $-1$ nie może należeć do obrazu funkcji kwadratowej.
\end{przyklad}


\begin{figure}[htbp]
\centering
\begin{tikzpicture}[>=Stealth,every node/.style={font=\small},scale=1]
  \draw[rounded corners,draw=gray!60] (-.7,-1.25) rectangle (.75,1.25);
  \draw[rounded corners,draw=gray!60] (4.0,-1.25) rectangle (5.45,1.25);
  \node at (0,1.52) {$A$}; \node at (4.72,1.52) {$B$};
  \node (a) at (0,.75) {$-2$}; \node (b) at (0,-.1) {$2$};
  \node (c) at (0,-.89) {$0$};
  \node (d) at (4.72,.48) {$4$}; \node (e) at (4.72,-.75) {$0$};
  \draw[->,blue!70!black] (a.east) -- (d.west);
  \draw[->,blue!70!black] (b.east) -- (d.west);
  \draw[->] (c.east) -- (e.west);
  \node[anchor=north,align=center] at (2.37,-1.45)
    {$S=\{2\}$,\quad$f[S]=\{4\}$,\\$f^{-1}[f[S]]=\{-2,2\}$};
\end{tikzpicture}
\caption{Przeciwobraz obrazu może zawierać więcej argumentów niż zbiór wyjściowy.}
\label{fig:preimage-image-xliv}
\end{figure}

\section{Składanie funkcji i identyczność}
\begin{definicja}[Złożenie]\label{def:composition}
Dla funkcji $f:A\to B$ i $g:B\to C$ \emph{złożenie} $g\circ f:A\to C$ określamy przez $(g\circ f)(a)=g(f(a))$ dla $a\in A$.
\end{definicja}
\begin{twierdzenie}[Łączność i funkcje identycznościowe]\label{thm:composition-laws}
Dla $f:A\to B$, $g:B\to C$, $h:C\to D$ zachodzi $h\circ(g\circ f)=(h\circ g)\circ f$, $f\circ\id_A=f$ i $\id_B\circ f=f$.
\end{twierdzenie}
\begin{proof}
Dla każdego $a\in A$ obie strony pierwszej równości przyjmują wartość $h(g(f(a)))$. Funkcje mają tę samą dziedzinę i przeciwdziedzinę, a ich wartości są równe, więc są równe. Dla każdego $a\in A$ mamy $(f\circ\id_A)(a)=f(a)$ i $(\id_B\circ f)(a)=f(a)$, co dowodzi pozostałych równości.
\end{proof}
\begin{twierdzenie}[Iniekcja i surjekcja złożenia]\label{thm:composition-inj-surj}
Niech $f:A\to B$, $g:B\to C$. Jeśli obie funkcje są iniekcjami, to $g\circ f$ jest iniekcją; jeśli obie są surjekcjami, to $g\circ f$ jest surjekcją. Ponadto jeśli $g\circ f$ jest iniekcją, to $f$ jest iniekcją, a jeśli $g\circ f$ jest surjekcją, to $g$ jest surjekcją.
\end{twierdzenie}
\begin{proof}
Jeśli $g(f(a))=g(f(a'))$, iniekcja $g$ daje $f(a)=f(a')$, a iniekcja $f$ daje $a=a'$. Jeśli $c\in C$, surjekcja $g$ daje $b\in B$ z $g(b)=c$, a surjekcja $f$ daje $a\in A$ z $f(a)=b$. Wtedy $(g\circ f)(a)=c$.

Jeśli $g\circ f$ jest iniekcją i $f(a)=f(a')$, to $g(f(a))=g(f(a'))$, więc $a=a'$. Jeśli $g\circ f$ jest surjekcją i $c\in C$, istnieje $a\in A$ z $g(f(a))=c$; zatem $b=f(a)\in B$ jest przeciwobrazem $c$ przez $g$.
\end{proof}
\begin{przyklad}[Dlaczego odwrotne implikacje zawodzą]
Niech $A=\{0\}$, $B=\{0,1\}$, $C=\{0\}$, $f(0)=0$ i $g(0)=g(1)=0$. Złożenie $g\circ f$ jest bijekcją, choć $f$ nie jest surjekcją, a $g$ nie jest iniekcją.
\end{przyklad}

\begin{twierdzenie}[Obraz złożenia]\label{thm:composition-image-xliv}
Niech $f:A\to B$, $g:B\to C$, $S\subseteq A$. Wówczas
$(g\circ f)[S]=g[f[S]]$.
\end{twierdzenie}
\begin{proof}
Dla $c\in C$ przynależność $c\in(g\circ f)[S]$ oznacza,
że istnieje $a\in S$ spełniające $g(f(a))=c$.
Element $b=f(a)$ należy do $f[S]$ oraz spełnia $g(b)=c$,
więc $c\in g[f[S]]$.
Odwrotnie, jeśli $c\in g[f[S]]$, to istnieje $b\in f[S]$
z $g(b)=c$. Z definicji obrazu istnieje $a\in S$ z $f(a)=b$.
Wtedy $(g\circ f)(a)=g(f(a))=g(b)=c$, co daje pierwszą
przynależność. Oba zawierania dowodzą równości.
\end{proof}
\begin{przyklad}[Obraz i przeciwobraz złożenia]
Niech $A=\{-2,-1,0,1,2\}$, $B=\{0,1,4\}$,
$C=\{0,1\}$, $f(x)=x^2$, $g(0)=g(4)=0$, $g(1)=1$.
Dla $S=\{-2,-1\}$ mamy $f[S]=\{4,1\}$ i
$g[f[S]]=\{0,1\}$; bezpośrednio $(g\circ f)(-2)=0$,
$(g\circ f)(-1)=1$, zatem $(g\circ f)[S]=\{0,1\}$.
Dla $T=\{1\}\subseteq C$ otrzymujemy
$g^{-1}[T]=\{1\}$ oraz
$f^{-1}[g^{-1}[T]]=\{-1,1\}$,
co zgadza się z $(g\circ f)^{-1}[T]$.
\end{przyklad}

\begin{twierdzenie}[Przeciwobraz złożenia]\label{thm:composition-preimage}
Jeżeli $f:A\to B$, $g:B\to C$ i $T\subseteq C$, to $(g\circ f)^{-1}[T]=f^{-1}[g^{-1}[T]]$.
\end{twierdzenie}
\begin{proof}
Dla dowolnego $a\in A$ zachodzi $a\in(g\circ f)^{-1}[T]$ wtedy i tylko wtedy, gdy $g(f(a))\in T$, czyli $f(a)\in g^{-1}[T]$, co z kolei jest równoważne $a\in f^{-1}[g^{-1}[T]]$. Z ekstensjonalności otrzymujemy równość.
\end{proof}

\section{Zestawianie, sklejanie i iloczyn kartezjański funkcji}
Samo składanie funkcji nie wyczerpuje naturalnych operacji. Często chcemy
z dwóch funkcji o wspólnej dziedzinie zbudować jedną funkcję o wartościach
w iloczynie kartezjańskim, skleić funkcje określone na różnych częściach
zbioru albo połączyć niezależne funkcje działające na dwóch różnych osiach.

\begin{definicja}[Zestawianie funkcji]\label{def:pairing-functions}
Niech $f:A\to B$ i $g:A\to C$ będą funkcjami o tej samej dziedzinie.
Oznaczenie $(f,g)$ rezerwujemy w tym podrozdziale dla funkcji
zestawionej z $f$ i $g$. Gdy obie funkcje mają wspólną dziedzinę,
możemy równocześnie obliczać ich wartości i zapisywać parę wyników.
Nie chodzi o składanie funkcji: żadna funkcja nie jest tu stosowana po drugiej.
Ich \emph{zestawieniem} nazywamy funkcję
\[
(f,g):A\to B\times C,
\qquad
(f,g)(a):=(f(a),g(a)).
\]
\end{definicja}

\begin{twierdzenie}[Własność uniwersalna zestawienia]\label{thm:pairing-universal}
Niech $f:A\to B$ i $g:A\to C$. Oznaczmy przez
$\pi_B:B\times C\to B$ i $\pi_C:B\times C\to C$ rzuty,
$\pi_B(b,c)=b$, $\pi_C(b,c)=c$. Wtedy
\[
\pi_B\circ(f,g)=f,
\qquad
\pi_C\circ(f,g)=g.
\]
Ponadto, jeżeli $h:A\to B\times C$ spełnia
$\pi_B\circ h=f$ oraz $\pi_C\circ h=g$, to $h=(f,g)$.
\end{twierdzenie}
\begin{proof}
Dla każdego $a\in A$ mamy
\[
(\pi_B\circ(f,g))(a)=\pi_B(f(a),g(a))=f(a)
\]
oraz analogicznie $(\pi_C\circ(f,g))(a)=g(a)$.
Jeżeli $h$ ma tę samą własność, to dla każdego $a\in A$
zapisać możemy $h(a)=(b_a,c_a)$ dla pewnych $b_a\in B$, $c_a\in C$.
Z założeń wynika
\[
 b_a=\pi_B(h(a))=f(a),
 \qquad
 c_a=\pi_C(h(a))=g(a),
\]
a więc $h(a)=(f(a),g(a))=(f,g)(a)$. Zatem $h=(f,g)$.
\end{proof}

\begin{przyklad}
Dla $f:\R\to\R$, $f(x)=x^2$, oraz $g:\R\to\R$, $g(x)=x+1$, zestawienie
$(f,g):\R\to\R^2$ ma postać
$x\mapsto(x^2,x+1)$. Nie należy go mylić z iloczynem funkcji w sensie
mnożenia wartości ani z iloczynem kartezjańskim funkcji z Definicji~\ref{def:cartesian-product-functions}.
\end{przyklad}

\begin{definicja}[Sklejanie funkcji]\label{def:gluing-functions}
Niech $A_1,A_2\subseteq X$, $f_1:A_1\to B$, $f_2:A_2\to B$.
Mówimy, że funkcje \emph{zgadzają się na części wspólnej}, jeśli
\[
\forall x\in A_1\cap A_2\quad f_1(x)=f_2(x).
\]
Wtedy \emph{sklejeniem} $f_1$ i $f_2$ na $A_1\cup A_2$ nazywamy funkcję
$h:A_1\cup A_2\to B$ daną wzorem
\[
 h(x)=\begin{cases}
 f_1(x),&x\in A_1,\\
 f_2(x),&x\in A_2.
 \end{cases}
\]
\end{definicja}

\begin{twierdzenie}[Kryterium sklejania]\label{thm:gluing-functions}
Niech $X,B$ będą zbiorami, $A_1,A_2\subseteq X$, a
$f_1:A_1\to B$ i $f_2:A_2\to B$ funkcjami.
Sklejenie z Definicji~\ref{def:gluing-functions} jest dobrze określone wtedy i tylko wtedy,
gdy $f_1$ i $f_2$ zgadzają się na $A_1\cap A_2$. W takim przypadku jest ono jedyną funkcją
$h:A_1\cup A_2\to B$, której ograniczenia do $A_1$ i $A_2$ są równe odpowiednio $f_1$ i $f_2$.
\end{twierdzenie}
\begin{proof}
Jeżeli istnieje funkcja $h$ o podanych ograniczeniach, to dla każdego $x\in A_1\cap A_2$
mamy $f_1(x)=h(x)=f_2(x)$, więc zgodność na części wspólnej jest konieczna.
Załóżmy odwrotnie, że $f_1$ i $f_2$ zgadzają się na $A_1\cap A_2$.
Dla $x\in A_1\cup A_2$ chcemy zadać wartość $h(x)$ według wzoru z definicji.
Jeżeli $x\in A_1\setminus A_2$ albo $x\in A_2\setminus A_1$, nie ma żadnej niejednoznaczności.
Jeżeli zaś $x\in A_1\cap A_2$, oba wzory dają tę samą wartość z założenia, więc funkcja jest dobrze określona.
Dla każdego $x\in A_1$ definicja $h$ daje $h(x)=f_1(x)$,
a więc $h|_{A_1}=f_1$. Tak samo dla każdego $x\in A_2$
mamy $h(x)=f_2(x)$, co daje $h|_{A_2}=f_2$.
Równości te są zgodne również na przecięciu dziedzin,
ponieważ założyliśmy $f_1(x)=f_2(x)$ dla takich $x$.
Jeśli $k$ jest inną funkcją o tych samych ograniczeniach, to dla $x\in A_1\cup A_2$ mamy
albo $x\in A_1$, skąd $k(x)=f_1(x)=h(x)$, albo $x\in A_2$, skąd $k(x)=f_2(x)=h(x)$.
Zatem $k=h$.
\end{proof}

\begin{przyklad}[Sklejanie na półosiach]
Funkcje $f_1:(-\infty,0]\to\R$, $f_1(x)=x^2$, oraz
$f_2:[0,\infty)\to\R$, $f_2(x)=x^2$ zgadzają się w punkcie $0$,
więc sklejają się do funkcji $h:\R\to\R$, $h(x)=x^2$.
Natomiast funkcje $g_1:(-\infty,0]\to\R$, $g_1(x)=0$, oraz
$g_2:[0,\infty)\to\R$, $g_2(x)=1$, nie sklejają się do funkcji
na całej $\R$, bo na części wspólnej $\{0\}$
przypisują różne wartości.
\end{przyklad}


\begin{twierdzenie}[Sklejanie dowolnej rodziny funkcji]
\label{thm:gluing-family-xliv}
Niech $I$ będzie zbiorem indeksów, $X,B$ zbiorami oraz
$(A_i)_{i\in I}$ rodziną podzbiorów $X$. Dla każdego $i\in I$
niech $f_i:A_i\to B$ będzie funkcją. Istnieje funkcja
\[
 h:\bigcup_{i\in I} A_i\longrightarrow B,
 \qquad h|_{A_i}=f_i\quad(i\in I),
\]
wtedy i tylko wtedy, gdy dla wszystkich $i,j\in I$ oraz
$x\in A_i\cap A_j$ mamy $f_i(x)=f_j(x)$.
Wówczas $h$ jest jednoznaczna, a jej wykresem jest
$\bigcup_{i\in I}G_{f_i}$. Dla $I=\varnothing$ jest to
jedyna funkcja $\varnothing\to B$.
\end{twierdzenie}
\begin{proof}
Załóżmy, że $h$ istnieje. Jeśli $x\in A_i\cap A_j$,
warunki ograniczenia dają $f_i(x)=h(x)=f_j(x)$,
zatem zgodność jest konieczna.

Odwrotnie, załóżmy zgodność na wszystkich częściach wspólnych.
Rodzina wykresów $(G_{f_i})_{i\in I}$ jest rodziną zbiorów,
więc jej suma $G=\bigcup_{i\in I}G_{f_i}$ jest zbiorem.
Dla każdego $x\in\bigcup_i A_i$ istnieje indeks $i$
z $x\in A_i$, stąd para $(x,f_i(x))\in G$ zapewnia istnienie
wartości. Jeżeli $(x,b),(x,b')\in G$, to istnieją indeksy
$i,j$ z $b=f_i(x)$ i $b'=f_j(x)$. Punkt $x$ należy
wówczas do $A_i\cap A_j$, a więc $b=b'$ na mocy zgodności.
Zatem $G$ jest wykresem funkcji $h$.
Dla $x\in A_i$ definicja daje $h(x)=f_i(x)$,
więc $h|_{A_i}=f_i$. Jeśli $k$ ma te same ograniczenia,
to dla każdego $x$ wybieramy \emph{jeden} indeks $i$
ze $x\in A_i$ i otrzymujemy $k(x)=f_i(x)=h(x)$.
Nie dokonujemy jednoczesnego wyboru indeksu dla wszystkich $x$;
sprawdzamy równość punktowo. Zatem $k=h$.
Gdy $I=\varnothing$, dziedzina i suma wykresów są puste,
a argument pozostaje poprawny.
\end{proof}
\begin{przyklad}[Trzy funkcje na nakładających się dziedzinach]
Niech $A_1=\{1,2\}$, $A_2=\{2,3\}$, $A_3=\{3,4\}$,
$B=\{0,1\}$. Ustalmy
$f_1(1)=0$, $f_1(2)=1$,
$f_2(2)=1$, $f_2(3)=0$,
$f_3(3)=0$, $f_3(4)=1$.
Na $A_1\cap A_2=\{2\}$ wartości wynoszą $1$,
a na $A_2\cap A_3=\{3\}$ wynoszą $0$.
Przecięcie $A_1\cap A_3$ jest puste, więc nie daje
ograniczenia. Sklejona funkcja na $\{1,2,3,4\}$
przyjmuje kolejno wartości $(0,1,0,1)$.
Zmiana jedynie $f_3(3)$ na $1$ uniemożliwia sklejenie,
ponieważ punktowi $3$ przypisano by dwie różne wartości.
\end{przyklad}

\begin{definicja}[Iloczyn kartezjański funkcji]\label{def:cartesian-product-functions}
Niech $f:A\to B$ i $g:C\to D$. \emph{Iloczynem kartezjańskim funkcji}
$f$ i $g$ nazywamy funkcję
\[
 f\times g:A\times C\to B\times D,
 \qquad
 (f\times g)(a,c):=(f(a),g(c)).
\]
\end{definicja}


\begin{przyklad}[Zestawianie a iloczyn kartezjański]
Weźmy $f:\R\to\R$, $f(x)=x+1$, oraz $g:\R\to\R$, $g(x)=2x$.
Zestawienie $(f,g)$ przyjmuje \emph{jeden} argument i zwraca
parę: $(f,g)(3)=(4,6)$. Iloczyn kartezjański $f\times g$
przyjmuje natomiast \emph{parę} argumentów i zwraca parę:
$(f\times g)(3,5)=(4,10)$.
Dziedziny tych dwóch funkcji są różne: odpowiednio $\R$ i $\R^2$.
Nie wolno więc utożsamiać tych konstrukcji.
\end{przyklad}

\begin{twierdzenie}[Podstawowe własności iloczynu funkcji]\label{thm:product-functions}
Niech $f:A\to B$, $f_1:B\to E$, $g:C\to D$, $g_1:D\to F$.
Wtedy
\[
 (f_1\times g_1)\circ(f\times g)=(f_1\circ f)\times(g_1\circ g).
\]
Jeżeli $f$ i $g$ są iniekcjami, to $f\times g$ jest iniekcją.
Jeżeli $f$ i $g$ są surjekcjami, to $f\times g$ jest surjekcją.
W szczególności iloczyn kartezjański bijekcji jest bijekcją.
\end{twierdzenie}
\begin{proof}
Dla każdego $(a,c)\in A\times C$ mamy
\[
((f_1\times g_1)\circ(f\times g))(a,c)
=(f_1\times g_1)(f(a),g(c))=(f_1(f(a)),g_1(g(c))),
\]
co jest dokładnie wartością funkcji $((f_1\circ f)\times(g_1\circ g))(a,c)$.
To dowodzi pierwszej równości.

Załóżmy teraz, że $f$ i $g$ są iniekcjami oraz
$(f\times g)(a,c)=(f\times g)(a',c')$.
Wtedy $(f(a),g(c))=(f(a'),g(c'))$, a zatem $f(a)=f(a')$ i $g(c)=g(c')$.
Iniekcyjność daje $a=a'$ oraz $c=c'$, więc $(a,c)=(a',c')$.
Jeśli zaś $f$ i $g$ są surjekcjami, to dla dowolnego $(b,d)\in B\times D$
wybieramy $a\in A$ z $f(a)=b$ oraz $c\in C$ z $g(c)=d$.
Wtedy $(f\times g)(a,c)=(b,d)$, więc $f\times g$ jest surjekcją.
\end{proof}


\begin{twierdzenie}[Kiedy iniekcyjność iloczynu daje iniekcyjność czynników?]
\label{thm:product-converses-xliv}
Niech $f:A\to B$ i $g:C\to D$ będą funkcjami.
\begin{enumerate}[label=(\roman*)]
\item Jeżeli $A\ne\varnothing$, $C\ne\varnothing$ i
$f\times g$ jest iniekcją, to $f$ i $g$ są iniekcjami.
\item Jeżeli $B\ne\varnothing$, $D\ne\varnothing$ i
$f\times g$ jest surjekcją, to $f$ i $g$ są surjekcjami.
\end{enumerate}
Bez odpowiednich założeń o niepustości oba wnioski mogą być fałszywe.
\end{twierdzenie}
\begin{proof}
(i) Ustalmy $c_0\in C$. Jeżeli $f(a)=f(a')$, to
\[
 (f\times g)(a,c_0)=(f(a),g(c_0))
 =(f(a'),g(c_0))=(f\times g)(a',c_0).
\]
Iniekcja iloczynu daje $(a,c_0)=(a',c_0)$, czyli $a=a'$.
Zatem $f$ jest iniekcją. Ustalając $a_0\in A$ i
porównując $(f\times g)(a_0,c)$ z
$(f\times g)(a_0,c')$, otrzymujemy analogicznie
iniekcyjność $g$.

(ii) Ustalmy $d_0\in D$. Dla dowolnego $b\in B$
surjekcja $f\times g$ daje parę $(a,c)\in A\times C$
taką, że $(f(a),g(c))=(b,d_0)$. W szczególności
$f(a)=b$, więc $f$ jest surjekcją. Ustalając
$b_0\in B$ i rozpatrując $(b_0,d)$ dla dowolnego
$d\in D$, dostajemy surjekcyjność $g$.
\end{proof}
\begin{przyklad}[Puste czynniki iloczynu]
Niech $f:\varnothing\to\{0\}$ będzie jedyną funkcją,
a $g:\{0,1\}\to\{0\}$ będzie stała. Funkcja $g$
nie jest iniekcją, ale $f\times g$ ma pustą dziedzinę,
więc jest iniekcją. Z kolei dla
$f:\varnothing\to\{0\}$ i $g:\varnothing\to\varnothing$
iloczyn $f\times g:\varnothing\to\varnothing$ jest
surjekcją, chociaż $f$ nią nie jest. Pustość czynnika
powoduje, że iloczyn nie zachowuje informacji o czynniku.
\end{przyklad}

\section{Ograniczenia i przedłużenia}
\begin{definicja}
Jeśli $f:A\to B$ i $S\subseteq A$, to \emph{ograniczeniem} $f$ do $S$ nazywamy $f|_S:S\to B$, $f|_S(s)=f(s)$. Jeśli $A\subseteq D$ i $F:D\to B$ spełnia $F|_A=f$, nazywamy $F$ \emph{przedłużeniem} $f$.
\end{definicja}
\begin{przyklad}
Funkcja $f:(0,\infty)\to\R$, $f(x)=1/x$, jest ograniczeniem funkcji $F:\R\setminus\{0\}\to\R$ o tym samym wzorze. Nie można jednak zadać wartości w punkcie $0$ przez ten sam wzór $1/x$, ponieważ wyrażenie $1/0$ nie jest określone. Istnieje za to przedłużenie w sensie teoriomnogościowym na całą $\R$: definiujemy $\widetilde F:\R\to\R$ przez $\widetilde F(x)=1/x$ dla $x\ne0$ i $\widetilde F(0)=0$. Rozróżniamy więc przedłużenie \emph{funkcji} od możliwości stosowania jednego wzoru w powiększonej dziedzinie.
\end{przyklad}
\begin{twierdzenie}[Przedłużenie na większą dziedzinę]\label{thm:extension-choice}
Niech $A\subseteq D$ i $f:A\to B$. Jeżeli $D\setminus A=\varnothing$, to $f$ jest jedynym przedłużeniem na $D$. Jeżeli $D\setminus A\ne\varnothing$, to przedłużenie $D\to B$ istnieje wtedy i tylko wtedy, gdy $B\ne\varnothing$.
\end{twierdzenie}
\begin{proof}
Gdy $D=A$, każde przedłużenie ma wszystkie wartości zadane przez $f$, więc jest równe $f$. Gdy $D\setminus A\ne\varnothing$ i istnieje przedłużenie $F$, wybieramy $d\in D\setminus A$, a wtedy $F(d)\in B$, zatem $B$ jest niepusty. Jeśli $B\ne\varnothing$, wybieramy \emph{jeden} element $b_0\in B$ i definiujemy $F(d)=f(d)$ dla $d\in A$ oraz $F(d)=b_0$ dla $d\in D\setminus A$. Funkcja jest określona jednoznacznie w każdym punkcie i przedłuża $f$. Nie jest tu potrzebny aksjomat wyboru: wybieramy jeden element z jednego niepustego zbioru.
\end{proof}

\begin{twierdzenie}[Ograniczenie i obrazy]\label{thm:restriction-images-xliv}
Niech $f:A\to B$, $S\subseteq A$ i $T\subseteq B$.
Dla ograniczenia $f|_S:S\to B$ zachodzi
\[
 (f|_S)[S]=f[S],\qquad
 (f|_S)^{-1}[T]=S\cap f^{-1}[T].
\]
Jeśli $f$ jest iniekcją, to każde jego ograniczenie
jest iniekcją. Surjekcyjność na $B$ nie musi się zachować.
\end{twierdzenie}
\begin{proof}
Wartości funkcji $f|_S$ są dokładnie wartościami
$f(s)$ dla $s\in S$, co daje pierwszą równość.
Dla $s\in S$ zachodzi $s\in(f|_S)^{-1}[T]$
wtedy i tylko wtedy, gdy $f(s)\in T$.
Warunek ten jest równoważny $s\in S\cap f^{-1}[T]$,
więc zbiory są równe. Jeśli $s,s'\in S$ i
$f|_S(s)=f|_S(s')$, to $f(s)=f(s')$; iniekcja
$f$ daje $s=s'$.
\end{proof}
\begin{przyklad}[Ograniczenie może zniszczyć surjekcję]
Funkcja $f:\{1,2\}\to\{u,v\}$ określona przez
$f(1)=u$, $f(2)=v$ jest bijekcją.
Po ograniczeniu do $S=\{1\}$ obraz wynosi $\{u\}$,
więc $f|_S:S\to\{u,v\}$ nie jest surjekcją,
choć pozostaje iniekcją.
\end{przyklad}

\section{Funkcje na ilorazach i własność uniwersalna}
\begin{twierdzenie}[Faktoryzacja przez iloraz]\label{thm:quotient-universal}
Niech $\sim$ będzie relacją równoważności na $A$, $q:A\to A/{\sim}$ odwzorowaniem $q(a)=[a]$, a $f:A\to B$ funkcją. Następujące warunki są równoważne:
\begin{enumerate}[label=(\roman*)]
\item $a\sim a'$ implikuje $f(a)=f(a')$ dla wszystkich $a,a'\in A$;
\item istnieje funkcja $\bar f:A/{\sim}\to B$ taka, że $f=\bar f\circ q$.
\end{enumerate}
Jeśli $\bar f$ istnieje, jest jednoznaczna.
\end{twierdzenie}
\begin{proof}
Załóżmy (i). Dla każdej klasy $C\in A/{\sim}$ istnieje $a\in A$ z $C=[a]$. Definiujemy $\bar f(C)=f(a)$. Musimy wykazać, że definicja nie zależy od reprezentanta: jeżeli $C=[a']$, to $[a]=[a']$, skąd $a\sim a'$ na mocy Twierdzenia~\ref{thm:class-equality}. Z (i) otrzymujemy $f(a)=f(a')$. Zatem $\bar f$ jest funkcją, a dla każdego $a\in A$ mamy $(\bar f\circ q)(a)=\bar f([a])=f(a)$.

Załóżmy (ii). Jeśli $a\sim a'$, to $[a]=[a']$, więc $f(a)=\bar f([a])=\bar f([a'])=f(a')$. Dla jednoznaczności niech $g$ także spełnia $f=g\circ q$. Dla dowolnej klasy $C=[a]$ zachodzi $g(C)=g(q(a))=f(a)=\bar f(C)$, więc $g=\bar f$.
\end{proof}
\begin{przyklad}[Reszty modulo $m$]
Niech $m\in\N$, $m\ge2$, i niech $f:\Z\to\{0,\ldots,m-1\}$ przypisuje liczbie resztę z dzielenia przez $m$. Jeśli $a\equiv b\pmod m$, to $f(a)=f(b)$, więc istnieje dokładnie jedna funkcja $\bar f:\Z/{\equiv_m}\to\{0,\ldots,m-1\}$ taka, że $\bar f([a])=f(a)$. Jest ona bijekcją: każda reszta występuje, a różne klasy mają różne reszty.
\end{przyklad}

\begin{twierdzenie}[Iloraz przez równość wartości — wersja elementarna]
\label{thm:first-isomorphism-functions-xliv}
Niech $f:A\to B$ będzie dowolną funkcją. Określmy relację
$\sim_f$ na $A$ wzorem
\[
 a\sim_f a'\quad\Longleftrightarrow\quad f(a)=f(a').
\]
Jest to relacja równoważności. Funkcja
\[
 \widetilde f:A/{\sim_f}\longrightarrow f[A],
 \qquad \widetilde f([a])=f(a),
\]
jest dobrze określoną bijekcją. W szczególności
$A/{\sim_f}$ i zbiór wartości $f$ mają tę samą moc,
także gdy $A=\varnothing$.
\end{twierdzenie}
\begin{proof}
(1) \emph{Relacja równoważności.} Dla każdego $a\in A$
zachodzi $f(a)=f(a)$, więc $a\sim_f a$.
Z $f(a)=f(a')$ wynika $f(a')=f(a)$, więc relacja jest
symetryczna. Jeśli $f(a)=f(a')$ i $f(a')=f(a'')$,
to $f(a)=f(a'')$, co daje przechodniość.

(2) \emph{Dobrze określoność.} Każda klasa w ilorazie
ma postać $[a]$ dla pewnego $a\in A$.
Jeśli $[a]=[a']$, to $a\sim_f a'$ na mocy
Twierdzenia~\ref{thm:class-equality}, a więc $f(a)=f(a')$.
Zatem wynik nie zależy od reprezentanta klasy i należy do $f[A]$.

(3) \emph{Iniekcyjność.} Jeśli
$\widetilde f([a])=\widetilde f([a'])$, to
$f(a)=f(a')$, czyli $a\sim_f a'$.
Wobec tego $[a]=[a']$.

(4) \emph{Surjekcyjność.} Dla każdego $b\in f[A]$
istnieje $a\in A$ z $f(a)=b$. Zatem
$\widetilde f([a])=b$.

Gdy $A=\varnothing$, zarówno iloraz, jak i $f[A]$
są puste, a jedyna funkcja między nimi jest bijekcją,
więc również ten przypadek obejmuje teza.
\end{proof}
\begin{przyklad}[Iloraz funkcji o powtarzających się wartościach]
Dla $A=\{-2,-1,0,1,2\}$ i $f:A\to\R$, $f(x)=x^2$,
klasy relacji $\sim_f$ to $\{-2,2\}$, $\{-1,1\}$ i
$\{0\}$. Zbiór wartości wynosi $f[A]=\{0,1,4\}$.
Indukowana bijekcja przyporządkowuje tym klasom
odpowiednio $4$, $1$ i $0$. To pokazuje, dlaczego
odwzorowanie nieiniektywne może stać się bijekcją
po utożsamieniu argumentów o tych samych wartościach.
\end{przyklad}

\chapter{Języki formalne, termy i formuły}\label{ch:syntax-fol}

\section{Dlaczego odróżniamy napis od jego znaczenia?}
Zdanie „dla każdej liczby istnieje większa” jest zrozumiałe, lecz pozostawia wiele
rzeczy domyślnych: z jakiego zbioru pochodzą liczby, co znaczy „większa” i jaki
jest zakres kwantyfikatorów. W dowodzie zwykle radzimy sobie z tymi kwestiami,
ale gdy przedmiotem badań staje się \emph{sam dowód}, potrzebujemy języka,
którego poprawne wyrażenia można rozpoznawać według jednoznacznych reguł.
Pytanie, czy wyrażenie zostało dobrze zapisane, należy do \emph{składni};
pytanie, czy jest prawdziwe, należy do \emph{semantyki}.

Warto pamiętać o dwóch poziomach wypowiedzi. Pisząc „$\varphi$ jest formułą”,
mówimy \emph{o} pewnym napisie, czyli używamy metajęzyka. Pisząc symbol $\forall$
wewnątrz $\varphi$, tworzymy wyrażenie samego języka formalnego. W dalszych
rozdziałach rozróżnienie to okaże się niezbędne do sformułowania twierdzeń Gödla.

\section{Alfabety, słowa i sygnatury}
\begin{definicja}[Alfabet i słowo]\label{def:syntax-alphabet}
\emph{Alfabetem} nazywamy zbiór symboli $\Sigma$. \emph{Słowem} nad $\Sigma$
nazywamy skończony ciąg $(s_1,\ldots,s_n)$ symboli $s_i\in\Sigma$.
Dopuszczamy $n=0$: otrzymujemy wtedy \emph{słowo puste}, oznaczane $\varepsilon$.
Zbiór wszystkich słów nad $\Sigma$ oznaczamy przez $\Sigma^*$.
\emph{Konkatenacją} słów $u$ i $v$ jest słowo $uv$, otrzymane przez dopisanie
symboli $v$ za symbolami $u$.
\end{definicja}
\begin{przyklad}
Dla $\Sigma=\{a,b\}$ słowa $ab$, $aab$ i $\varepsilon$ należą do $\Sigma^*$.
Mamy $a\varepsilon=\varepsilon a=a$, natomiast $ab\ne ba$.
Słowo samo przez się nie ma jeszcze znaczenia matematycznego.
\end{przyklad}
\begin{stwierdzenie}[Łączność konkatenacji]\label{prop:syntax-concat}
Niech $\Sigma$ będzie alfabetem. Dla dowolnych słów $u,v,w\in\Sigma^*$ zachodzi $(uv)w=u(vw)$.
\end{stwierdzenie}
\begin{proof}
Niech $u=(u_1,\ldots,u_k)$, $v=(v_1,\ldots,v_\ell)$ i
$w=(w_1,\ldots,w_m)$. W obu sposobach łączenia otrzymujemy ciąg długości
$k+\ell+m$, którego pierwsze $k$ pozycji zajmują symbole $u_i$, następne
$\ell$ pozycji symbole $v_j$, a ostatnie $m$ pozycji symbole $w_r$.
Dwa skończone ciągi o jednakowej długości i jednakowych symbolach na każdej
pozycji są równe, co kończy dowód.
\end{proof}

\begin{twierdzenie}[Słowa tworzą monoid wolny]\label{thm:syntax-free-monoid-xlv}
Dla dowolnego alfabetu $\Sigma$ konkatenacja jest działaniem
łącznym na $\Sigma^*$, a słowo puste $\varepsilon$ jest jej
obustronnym elementem neutralnym. Ponadto funkcja długości
$|\cdot|:\Sigma^*\to\N$ spełnia $|uv|=|u|+|v|$.
Dla każdego monoidu $(M,\cdot,e)$ i każdej funkcji
$j:\Sigma\to M$ istnieje dokładnie jeden homomorfizm monoidów
$\overline j:\Sigma^*\to M$ taki, że
$\overline j((a))=j(a)$ dla każdej litery $a\in\Sigma$.
\end{twierdzenie}
\begin{proof}
Łączność wykazano w Stwierdzeniu~\ref{prop:syntax-concat}.
Jeżeli $u=(u_1,\ldots,u_k)$, to dopisanie na końcu zera symboli
nie zmienia żadnej spośród $k$ pozycji, stąd $u\varepsilon=u$.
Dopisanie zera symboli przed pierwszą pozycją daje
$\varepsilon u=u$; twierdzenie obejmuje też $k=0$.
Z definicji konkatenacji $uv$ zawiera najpierw $|u|$ symboli,
a następnie $|v|$ symboli. Te dwa kolejne bloki nie
nakładają się, więc łączna długość jest ich sumą.

Niech teraz $(M,\cdot,e)$ będzie dowolnym monoidem, a
$j:\Sigma\to M$ funkcją. Definiujemy
$\overline j(\varepsilon)=e$ oraz
$\overline j((a_1,\ldots,a_n))=
 j(a_1)\cdot\ldots\cdot j(a_n)$ dla $n\geq1$.
Łączność działania w $M$ zapewnia niezależność wyniku
od ustawienia nawiasów. Dla słów $u,v$ ich konkatenacja
polega na dopisaniu listy liter $v$ do listy liter $u$,
więc przez łączność i własność elementu neutralnego
$\overline j(uv)=\overline j(u)\cdot\overline j(v)$,
także gdy jedno ze słów jest puste.
Otrzymaliśmy homomorfizm, który na słowach jednoliterowych
zgadza się z $j$.
Jeżeli $h$ jest innym takim homomorfizmem, to
$h(\varepsilon)=e$, a dla każdego niepustego słowa
$w=(a_1)\cdots(a_n)$ mamy
$h(w)=h((a_1))\cdots h((a_n))
=j(a_1)\cdots j(a_n)=\overline j(w)$.
Zatem homomorfizm jest jednoznaczny, co uzasadnia
określenie \emph{monoid wolny na alfabecie $\Sigma$}.
\end{proof}
\begin{przyklad}[Czy kolejność dopisywania ma znaczenie?]
Dla $\Sigma=\{a,b\}$, $u=ab$, $v=b$ i $w=a$ otrzymujemy
$(uv)w=abba=u(vw)$, a $|uv|=3=2+1$.
Nie zachodzi natomiast przemienność: $uv=abb$, podczas
gdy $vu=bab$. Słowo $\varepsilon$ nie jest literą alfabetu,
lecz ciągiem o długości zero.
\end{przyklad}

\begin{definicja}[Sygnatura języka pierwszego rzędu]\label{def:syntax-signature}
\emph{Sygnaturą} $\mathcal L$ nazywamy zbiór symboli pozalogicznych
z podziałem na symbole stałych $\operatorname{Const}(\mathcal L)$,
funkcyjne $\operatorname{Fun}(\mathcal L)$ i relacyjne
$\operatorname{Rel}(\mathcal L)$. Każdemu symbolowi funkcyjnemu i relacyjnemu
przypisujemy dodatnią liczbę naturalną, zwaną \emph{arnością};
symbole stałych mają arność $0$.
Symbole należące do różnych kategorii są odrębne, nawet jeśli w zapisie
nieformalnym używamy dla nich tych samych liter.
\end{definicja}
Wybieramy ponadto nieskończony ciąg zmiennych $v_0,v_1,\ldots$ oraz znaki
$=,\neg,\land,\lor,\to,\leftrightarrow,\forall,\exists$ i symbole pomocnicze
(nawiasy, przecinki). Znak równości traktujemy jako symbol logiczny.
\emph{Językiem pierwszego rzędu o sygnaturze $\mathcal L$} nazywamy te
symbole wraz z opisanymi poniżej regułami tworzenia termów i formuł.

\begin{przyklad}[Trzy sygnatury]\label{ex:syntax-signatures}
W języku arytmetyki mamy stałą $0$, symbol następstwa $S$ arności $1$,
symbole $+$, $\cdot$ arności $2$ i symbol relacyjny $<$ arności $2$.
Język grafów może zawierać tylko relację $E$ arności $2$.
W języku grup używamy stałej $e$, działania $\ast$ arności $2$
i symbolu $\iota$ arności $1$ dla operacji odwracania.
Żaden z tych symboli nie ma jeszcze przypisanej interpretacji.
\end{przyklad}


\begin{uwaga}[Dlaczego pojawia się teraz symbol $\to$?]
W początkowych rozdziałach zapis $p\Rightarrow q$ oznaczał
implikację między zdaniami. W ścisłym języku formalnym używamy
symbolu $\to$ jako \emph{znaku budującego formułę}.
Gdy zaś w zwykłym tekście matematycznym napiszemy
„z $A$ wynika $B$”, jest to wypowiedź o związku między formułami,
a niekoniecznie kolejna formuła wewnątrz języka.
Rozróżnienie symbolu wewnątrz języka i wypowiedzi o języku
będzie szczególnie ważne w twierdzeniu o pełności.
\end{uwaga}

\section{Termy: napisy nazywające obiekty}
\begin{definicja}[Term]\label{def:syntax-term}
Zbiór termów $\operatorname{Term}(\mathcal L)$ jest najmniejszym zbiorem
napisów spełniającym następujące warunki:
\begin{enumerate}[label=(\roman*)]
\item każda zmienna $v_i$ jest termem;
\item każdy symbol stałej jest termem;
\item jeśli $f$ ma arność $n$ i $t_1,\ldots,t_n$ są termami,
      to $f(t_1,\ldots,t_n)$ jest termem.
\end{enumerate}
Innych termów nie ma. \emph{Term domknięty} to term bez wystąpień zmiennych.
W zapisie formalnym symbole działań występują przed argumentami; zapis
$+(t_1,t_2)$ zwyczajowo skracamy do $t_1+t_2$.
\end{definicja}
\begin{przyklad}[Co jest, a co nie jest termem?]
Dla sygnatury arytmetyki termami są $S(0)$, $+(v_2,S(0))$
i $\cdot(+(v_0,0),v_3)$. Napisy $S(v_0,v_1)$ i $+(,v_0)$
nie są termami: pierwszy zawiera za dużo argumentów, drugi nie zawiera
pierwszego argumentu. Napis $v_0<v_1$ nie jest termem, lecz formułą atomową.
\end{przyklad}

\begin{lemat}[Jednoznaczne odczytywanie termów]\label{lem:syntax-unique-terms}
Niech $\mathcal L$ będzie sygnaturą. Każdy term języka $\mathcal L$ jest dokładnie jednego z trzech rodzajów:
zmienną, stałą lub napisem $f(t_1,\ldots,t_n)$.
W trzecim przypadku symbol $f$ i wszystkie termy $t_1,\ldots,t_n$
są jednoznacznie wyznaczone.
\end{lemat}
\begin{proof}
Trzy rodzaje rozpoczynają się różnymi rodzajami symboli, więc nie można
pomieszać zmiennej, stałej i zastosowania symbolu funkcyjnego.
Jeśli jeden napis da się zapisać jako $f(t_1,\ldots,t_n)$ i
$g(s_1,\ldots,s_m)$, to porównanie jego pierwszych symboli daje $f=g$,
a przypisana arność wymusza $n=m$.
Przecinki rozdzielające argumenty znajdują się na poziomie pierwszej pary
nawiasów po $f$. Przecinków należących do termów zagnieżdżonych nie należy
brać pod uwagę: występują one przy dodatniej głębokości zagnieżdżenia.
Zatem miejsce zakończenia każdego argumentu jest jednoznaczne.
Porównując odpowiadające im fragmenty obu identycznych napisów, otrzymujemy
$t_i=s_i$ dla każdego $i$. Formalne uzasadnienie można prowadzić
indukcyjnie po długości termu.
\end{proof}

\begin{przyklad}[Jednoznaczny rozkład termu]
W sygnaturze z dwuargumentowym $f$, jednoargumentowym $g$
i stałą $c$ rozważmy $t=f(g(v_0),f(c,v_1))$.
Zewnętrzny symbol $f$ ma dwa argumenty. Pierwszy to
$g(v_0)$, drugi to $f(c,v_1)$; przecinek wewnątrz
drugiego argumentu znajduje się na wyższym poziomie
zagnieżdżenia i nie rozdziela argumentów zewnętrznego $f$.
Drzewo konstrukcji ma liście $v_0,c,v_1$, a jego korzeń
ma etykietę $f$. Jest więc jeden sposób odczytania termu
zgodny z arnościami ustalonej sygnatury.
\end{przyklad}

\section{Formuły: napisy wyrażające własności}
\begin{definicja}[Formuła atomowa i formuła]\label{def:syntax-formula}
\emph{Formułą atomową} jest napis $t_1=t_2$ albo $R(t_1,\ldots,t_n)$,
gdzie $t_i$ są termami, a $R$ symbolem relacyjnym arności $n$.
Zbiór formuł jest najmniejszym zbiorem zawierającym wszystkie formuły atomowe
i zamkniętym na operacje
\[
\varphi\longmapsto\neg\varphi,
\qquad
(\varphi,\psi)\longmapsto
(\varphi\circ\psi),\quad
\circ\in\{\land,\lor,\to,\leftrightarrow\},
\]
oraz na tworzenie $\forall v_i\,\varphi$ i $\exists v_i\,\varphi$
dla dowolnego $i\in\N$. Nic innego nie jest formułą.
Dla jednoznaczności parsowania formuły z głównym spójnikiem binarnym
w pełnym zapisie są objęte nawiasami. W przykładach używamy również
nawiasów grupujących dla czytelności, na przykład zapis
$\forall v_0\,(v_0=v_0)$ oznacza ściśle tę samą formułę
co $\forall v_0\,v_0=v_0$. Nadmiarowe nawiasy grupujące
nie stanowią dodatkowej reguły tworzenia formuł i nie są
częścią napisu w ścisłej gramatyce. Tak rozumiane skróty
nie wpływają na jednoznaczność jego odczytywania.
\end{definicja}
\begin{przyklad}
W języku arytmetyki napis $\forall v_0\,(v_0=v_0)$ jest zdaniem
w sensie zdefiniowanym poniżej, natomiast
$\exists v_1\,(+(v_1,v_1)=v_0)$ jest formułą z wolną zmienną $v_0$.
Napisy $\forall(v_0=v_0)$ oraz $\exists v_1\,+(v_1,v_1)$ nie są formułami.
\end{przyklad}

\begin{lemat}[Jednoznaczność budowy formuły]\label{lem:syntax-unique-formulas}
Niech $\mathcal L$ będzie sygnaturą. Dla każdej formuły języka $\mathcal L$ $\varphi$ można jednoznacznie odtworzyć jej ostatni
krok konstrukcji, a następnie wszystkie bezpośrednie podformuły.
\end{lemat}
\begin{proof}
Jeśli napis zaczyna się od symbolu relacyjnego, jego arność i nawiasy
określają argumenty. Jeśli zaczyna się od termu, to po jego końcu
w formule atomowej musi wystąpić znak równości i drugi term; jednoznaczność
termów zapewnia Lemat~\ref{lem:syntax-unique-terms}.
Jeśli napis zaczyna się od $\neg$, usunięcie pierwszego symbolu daje
jedyną podformułę. Jeśli zaczyna się od kwantyfikatora, następny symbol
wskazuje związaną zmienną, a pozostały napis jest podformułą.
Pozostają nawiasowane formuły binarne. Czytając je od lewej do prawej
i zliczając poziom otwarcia nawiasów, znajdujemy dokładnie jeden
spójnik na zewnętrznym poziomie. Dzieli on napis na lewą i prawą
podformułę. Nie może być ich dwóch: wewnętrzne spójniki znajdują
się w zagnieżdżonych nawiasach, a formuła musi być zbudowana według
podanej gramatyki. Powtarzanie procedury dla podformuł kończy się,
ponieważ są one krótsze od formuły wyjściowej.
\end{proof}

\begin{przyklad}[Drzewo budowy formuły]\label{ex:syntax-formula-tree-xlv}
Dla relacji dwuargumentowej $R$ i symbolu funkcyjnego
jednoargumentowego $f$ rozważmy
$\theta=\forall v_0\,(R(v_0,f(v_1))\to v_0=v_1)$.
Ostatnią operacją jest kwantyfikacja po $v_0$; pod nią
znajduje się implikacja z dwiema podformułami atomowymi.
Drzewo odzwierciedla kolejność konstrukcji, nie jest
rysunkiem wartościowania ani modelem relacji $R$.
\end{przyklad}
\begin{figure}[htbp]
\centering
\begin{tikzpicture}[>=Latex, every node/.style={draw,rounded corners=2pt,inner sep=4pt,font=\small}]
  \node (q) at (0,1.9) {$\forall v_0$};
  \node (imp) at (0,0.95) {$\to$};
  \node (rel) at (-2.45,0) {$R(v_0,f(v_1))$};
  \node (eq) at (2.15,0) {$v_0=v_1$};
  \draw[->] (q)--(imp);
  \draw[->] (imp)--(rel);
  \draw[->] (imp)--(eq);
\end{tikzpicture}
\caption{Drzewo składniowe formuły $\theta$. Strzałki prowadzą
od konstrukcji do jej bezpośrednich podformuł.}
\label{fig:syntax-tree-xlv}
\end{figure}

\section{Zmienne wolne, związane i zakres kwantyfikatora}
\begin{definicja}[Zmienne wolne]\label{def:syntax-free-vars}
Zbiór $\operatorname{FV}(t)$ zmiennych termu definiujemy przez
$\operatorname{FV}(v_i)=\{v_i\}$, $\operatorname{FV}(c)=\varnothing$
i
$\operatorname{FV}(f(t_1,\ldots,t_n))=\bigcup_{j=1}^{n}\operatorname{FV}(t_j)$.
Dla formuł definiujemy
\begin{align*}
\operatorname{FV}(t_1=t_2)
  &=\operatorname{FV}(t_1)\cup\operatorname{FV}(t_2),\\
\operatorname{FV}(R(t_1,\ldots,t_n))
  &=\bigcup_{j=1}^{n}\operatorname{FV}(t_j),\\
\operatorname{FV}(\neg\varphi)&=\operatorname{FV}(\varphi),\\
\operatorname{FV}(\varphi\circ\psi)
  &=\operatorname{FV}(\varphi)\cup\operatorname{FV}(\psi),\\
\operatorname{FV}(\forall v_i\,\varphi)
  &=\operatorname{FV}(\varphi)\setminus\{v_i\},\\
\operatorname{FV}(\exists v_i\,\varphi)
  &=\operatorname{FV}(\varphi)\setminus\{v_i\}.
\end{align*}
W tych wzorach $\circ$ oznacza dowolny spójnik binarny.
Formułę bez zmiennych wolnych nazywamy \emph{zdaniem}.
\end{definicja}
\begin{przyklad}[To samo oznaczenie może wystąpić wolno i związanie]
W formule $(\exists v_1\,(v_0=v_1))\land(v_1=v_2)$
pierwsze wystąpienie $v_1$ po kwantyfikatorze jest związane, ale
w drugim członie $v_1$ występuje wolno. Zatem
$\operatorname{FV}(\varphi)=\{v_0,v_1,v_2\}$.
Kwantyfikator nie oddziałuje na wyrażenie znajdujące się poza
zakresem swojej podformuły.
\end{przyklad}

\begin{przyklad}[Obliczanie zmiennych wolnych etapami]
Niech
$\rho=(\exists v_1\,R(v_0,v_1))\land
 (\forall v_0\,(R(v_0,v_2)\lor R(v_1,v_0)))$,
gdzie $R$ jest symbolem relacyjnym arności $2$.
Pierwszy człon ma zmienne wolne $\{v_0\}$: $v_1$ zostało
związane. Podformuły atomowe drugiego członu mają
zbiory zmiennych odpowiednio $\{v_0,v_2\}$ i
$\{v_1,v_0\}$, więc ich alternatywa ma zbiór
$\{v_0,v_1,v_2\}$. Kwantyfikator $\forall v_0$
usuwa z niego $v_0$ i pozostawia $\{v_1,v_2\}$.
Ostatecznie $\operatorname{FV}(\rho)=\{v_0,v_1,v_2\}$.
Wystąpienie zmiennej $v_1$ jest związane w pierwszym
członie, lecz wolne w drugim; podobnie $v_0$ jest wolne
w pierwszym członie, ale związane w drugim.
\end{przyklad}

\section{Podstawianie bez przechwytywania zmiennych}
Chcemy wstawić do formuły term opisujący konkretny obiekt.
Nie wystarczy jednak zastąpić każdego wystąpienia danej litery:
kwantyfikator może przez przypadek związać zmienną podstawianego termu.

\begin{definicja}[Podstawienie termu w termie]\label{def:syntax-term-substitution-xlv}
Dla termów $u,t$ języka $\mathcal L$ i zmiennej $x$
term $t[x:=u]$ definiujemy rekurencyjnie:
\begin{enumerate}[label=(\roman*)]
\item $x[x:=u]=u$, a $y[x:=u]=y$ dla każdej zmiennej $y\ne x$;
\item $c[x:=u]=c$ dla każdego symbolu stałej $c$;
\item $f(t_1,\ldots,t_n)[x:=u]
 =f(t_1[x:=u],\ldots,t_n[x:=u])$.
\end{enumerate}
Termy nie zawierają kwantyfikatorów, więc żadna zmienna
podstawianego termu nie może zostać w ten sposób związana.
\end{definicja}
\begin{lemat}[Zmienne termu po podstawieniu]
\label{lem:syntax-term-fv-subst-xlv}
Niech $\mathcal L$ będzie sygnaturą. Dla dowolnych termów
$t,u$ języka $\mathcal L$ i zmiennej $x$ zachodzi
\[
\operatorname{FV}(t[x:=u])=
\begin{cases}
(\operatorname{FV}(t)\setminus\{x\})\cup\operatorname{FV}(u),
 &x\in\operatorname{FV}(t),\\
\operatorname{FV}(t),&x\notin\operatorname{FV}(t).
\end{cases}
\]
\end{lemat}
\begin{proof}
W przypadku $t=x$ wynikiem jest $u$, a po prawej
otrzymujemy $\operatorname{FV}(u)$, więc równość zachodzi.
Dla $t=y\ne x$ nie zmieniamy termu ani jego jedynej
zmiennej. Dla stałej oba zbiory zmiennych są puste.
Przy $t=f(t_1,\ldots,t_n)$ podstawienie wykonuje się
oddzielnie w każdym argumencie. Jeżeli $x$ nie występuje
w żadnym $t_i$, nie zmienia się żaden argument.
W przeciwnym razie dla argumentów zawierających $x$
założenie indukcyjne usuwa tę zmienną i dodaje cały zbiór
$\operatorname{FV}(u)$, a pozostałe argumenty zachowują swoje
zmienne. Po zsumowaniu zbiorów zmiennych wszystkich
argumentów otrzymujemy dokładnie pierwszy przypadek wzoru.
\end{proof}
\begin{przyklad}[Podstawienie w zagnieżdżonym termie]
Niech $t=f(x,g(x,y))$ i $u=h(y)$ dla symboli funkcyjnych
$f$ arności $2$, $g$ arności $2$ i $h$ arności $1$.
Wtedy $t[x:=u]=f(h(y),g(h(y),y))$.
Obie początkowe kopie $x$ są zastąpione; nowy term ma
zbiór zmiennych $\{y\}$, zgodnie z lematem.
\end{przyklad}

\begin{definicja}[Dopuszczalne podstawienie]\label{def:syntax-free-for}
Niech $x$ będzie zmienną, $t$ termem, a $\varphi$ formułą.
Mówimy, że $t$ jest \emph{dopuszczalny do podstawienia za $x$ w $\varphi$},
jeśli żadne wolne wystąpienie $x$ nie znajduje się w zakresie
kwantyfikatora wiążącego zmienną z $\operatorname{FV}(t)$.
Wtedy $\varphi[x:=t]$ oznacza formułę powstałą przez zastąpienie
\emph{wolnych} wystąpień $x$ termem $t$. Formalnie:
na atomach podstawiamy term w każdym jego argumencie,
a dla negacji i spójników binarnych wykonujemy operację
w bezpośrednich podformułach. Przy $Qy\,\psi$, gdzie
$Q\in\{\forall,\exists\}$, jeżeli $y=x$, pozostawiamy
całą formułę bez zmian. Jeżeli $y\ne x$, podstawiamy
rekurencyjnie w $\psi$ i zachowujemy kwantyfikator $Qy$.
Ta ostatnia reguła jest dozwolona tylko przy spełnieniu
powyższego warunku braku przechwycenia; w przeciwnym razie
najpierw trzeba poprawnie przemianować odpowiednią
zmienną związaną. Jeżeli $x\notin\operatorname{FV}(\varphi)$,
podstawienie nie zmienia formuły.
\end{definicja}
\begin{przyklad}[Przechwycenie zmiennej]
Rozważmy $\varphi=\forall v_1\,(v_0=v_1)$ i term $t=v_1$.
Bezpośrednia zamiana $v_0$ na $v_1$ dałaby
$\forall v_1\,(v_1=v_1)$, w którym nowe wystąpienie $v_1$
stało się związane. Podstawienie nie jest dopuszczalne.
Możemy jednak najpierw zmienić nazwę zmiennej związanej na świeżą
zmienną $v_2$, otrzymując $\forall v_2\,(v_0=v_2)$, a następnie
wykonać dopuszczalne podstawienie i otrzymać
$\forall v_2\,(v_1=v_2)$.
\end{przyklad}
\begin{lemat}[Zmienne wolne po podstawieniu]\label{lem:syntax-fv-subst}
Niech $\mathcal L$ będzie sygnaturą, $x$ zmienną,
$t$ termem języka $\mathcal L$, a $\varphi$ formułą tego języka.
Załóżmy, że $t$ jest dopuszczalny do podstawienia za $x$ w $\varphi$.
Jeśli $x\in\operatorname{FV}(\varphi)$, to
\[
\operatorname{FV}(\varphi[x:=t])
=(\operatorname{FV}(\varphi)\setminus\{x\})
\cup\operatorname{FV}(t).
\]
Jeśli $x\notin\operatorname{FV}(\varphi)$, podstawienie nie zmienia
formuły ani zbioru jej zmiennych wolnych.
\end{lemat}
\begin{proof}
Dowodzimy tezy indukcją po budowie $\varphi$.
Dla atomów równościowych i relacyjnych operacja dotyczy
wyłącznie termów. Wzór wynika wtedy z
Lematu~\ref{lem:syntax-term-fv-subst-xlv} po wzięciu sumy
zbiorów zmiennych argumentów. W przypadku negacji
$\operatorname{FV}(\neg\psi)=\operatorname{FV}(\psi)$,
więc teza jest dokładnie założeniem indukcyjnym.
Dla spójnika binarnego stosujemy to założenie osobno
w dwóch podformułach i sumujemy otrzymane zbiory.
Jeżeli w całej formule $x$ nie jest wolna, definicja
podstawienia nie zmienia żadnego wolnego wystąpienia $x$;
formuła pozostaje ta sama.

Rozważmy przypadek $\varphi=Qy\,\psi$, gdzie
$Q\in\{\forall,\exists\}$ i $x$ jest wolna w $\varphi$.
Wówczas $x\ne y$ oraz $x\in\operatorname{FV}(\psi)$.
Dopuszczalność podstawienia daje
$y\notin\operatorname{FV}(t)$: każde wolne wystąpienie $x$
w podformule $\psi$ jest bowiem w zakresie kwantyfikatora
$Qy$. Zatem, korzystając kolejno z definicji wolnych
zmiennych, założenia indukcyjnego i $y\notin\operatorname{FV}(t)$,
otrzymujemy
\begin{align*}
\operatorname{FV}((Qy\,\psi)[x:=t])
 &=\operatorname{FV}(\psi[x:=t])\setminus\{y\}\\
 &=((\operatorname{FV}(\psi)\setminus\{x\})
    \cup\operatorname{FV}(t))\setminus\{y\}\\
 &=(\operatorname{FV}(\varphi)\setminus\{x\})
    \cup\operatorname{FV}(t).
\end{align*}
To zamyka również oba przypadki kwantyfikatorów.
\end{proof}

\begin{uwaga}[Nazwy zmiennych związanych]
Zmiana nazwy zmiennej związanej na nową, niewystępującą
w rozważanej formule, nie powinna zmieniać jej znaczenia.
Takie formuły nazywa się \emph{$\alpha$-równoważnymi}.
W następnym rozdziale uzasadnimy to na podstawie definicji spełniania,
a w rozdziale o rachunku dowodowym będziemy korzystać wyłącznie
z podstawień dopuszczalnych albo wcześniej zmienimy nazwę zmiennej związanej.
\end{uwaga}

\begin{uwaga}[Dlaczego warunek dopuszczalności jest konieczny?]
Formuła $\forall y\,(x=y)$ ma wolną zmienną $x$.
Próba wstawienia termu $y$ bez przemianowania zmieniłaby
ją w $\forall y\,(y=y)$, czyli w zdanie. Pierwsza formuła
jest fałszywa przy dowolnym wartościowaniu w strukturze
co najmniej dwuelementowej; druga jest w tej strukturze
prawdziwa. Problem nie jest więc tylko notacyjny.
\end{uwaga}

\begin{lemat}[Podstawienie i kwantyfikator zasłaniający]\label{lem:syntax-shadowing-lxvii}
Niech $x,y$ będą zmiennymi, $t$ termem, a $\psi$ formułą.
Jeżeli $x=y$, to $(Qx\,\psi)[x:=t]=Qx\,\psi$ dla
$Q\in\{\forall,\exists\}$. Jeżeli $x\ne y$ i
$x\notin\operatorname{FV}(\psi)$, to również
$(Qy\,\psi)[x:=t]=Qy\,\psi$, niezależnie od zmiennych termu $t$.
Jeżeli $x\ne y$ i $x\in\operatorname{FV}(\psi)$, to
$t$ jest dopuszczalny do podstawienia za $x$ w $Qy\,\psi$
wtedy i tylko wtedy, gdy jest dopuszczalny w $\psi$
oraz $y\notin\operatorname{FV}(t)$.
\end{lemat}
\begin{proof}
W pierwszym przypadku zewnętrzny kwantyfikator wiąże
wszystkie wystąpienia $x$ w swoim zakresie, które mogłyby
być wolne w $\psi$; reguła podstawiania pozostawia całą
formułę bez zmian. W drugim przypadku w $\psi$ nie ma
wolnych wystąpień $x$, więc nie ma czego zastępować.
W trzecim przypadku każde wolne wystąpienie $x$ w $\psi$
leży w zakresie nowego kwantyfikatora $Qy$.
Dopuszczalność w $Qy\,\psi$ wymaga zatem, po pierwsze,
aby żaden kwantyfikator wewnątrz $\psi$ nie przechwytywał
zmiennej z $\operatorname{FV}(t)$, czyli dopuszczalności
w $\psi$, a po drugie, aby $y$ nie była wolna w $t$.
Oba warunki razem wykluczają wszystkie możliwe
przechwycenia i są wystarczające.
\end{proof}
\begin{przyklad}[Wystąpienie wolne i zasłonięte]
Niech $\psi=(x=z)\land\forall x\,(x=y)$, a $t=f(y)$.
W formule $\psi$ wolne wystąpienie $x$ znajduje się tylko
w pierwszym członie. Podstawienie daje
$\psi[x:=t]=(f(y)=z)\land\forall x\,(x=y)$;
wystąpienie $x$ pod wewnętrznym kwantyfikatorem pozostaje
nietknięte. Natomiast w $\forall y\,\psi$ to samo
podstawienie nie jest dopuszczalne, ponieważ $y$ z termu
$f(y)$ zostałaby przechwycona przez zewnętrzny kwantyfikator.
W $\forall x\,\psi$ podstawienie nie zmienia niczego.
\end{przyklad}

\section{Indukcja po budowie formuły}
Definicje termów i formuł podają również sposób prowadzenia dowodów.
Aby wykazać własność wszystkich formuł, wystarczy udowodnić ją
najpierw dla formuł atomowych, a potem sprawdzić, że przechodzi
przez każdy spójnik i oba kwantyfikatory. Nie jest to nowe
założenie logiczne, lecz zastosowanie zasady indukcji do skończonych
drzew konstrukcji napisów.

\begin{twierdzenie}[Zasada indukcji strukturalnej]
\label{thm:syntax-structural-induction-xlv}
Niech $P$ będzie własnością termów języka $\mathcal L$.
Jeśli $P$ zachodzi dla każdej zmiennej i każdej stałej,
a z $P(t_1),\ldots,P(t_n)$ wynika $P(f(t_1,\ldots,t_n))$
dla każdego symbolu funkcyjnego $f$ arności $n$, to
$P$ zachodzi dla wszystkich termów.
Analogicznie niech $Q$ będzie własnością formuł.
Jeśli $Q$ zachodzi dla wszystkich atomów oraz zachowuje się
przy negacji, czterech spójnikach binarnych i obu
kwantyfikatorach, to $Q$ zachodzi dla wszystkich formuł.
\end{twierdzenie}
\begin{proof}
Każda formuła i każdy term powstają w skończenie wielu
krokach, ponieważ zbiory z Definicji~\ref{def:syntax-term}
i Definicji~\ref{def:syntax-formula} są najmniejszymi
zbiorami zamkniętymi na wskazane konstrukcje.
Niech $T_P$ oznacza zbiór termów spełniających $P$.
Na mocy pierwszych dwóch założeń zawiera on wszystkie
zmienne i stałe, a na mocy ostatniego jest zamknięty
na tworzenie termów $f(\vec t)$. Minimalność zbioru
wszystkich termów daje $T_P=\operatorname{Term}(\mathcal L)$.
Dla formuł zbiór $F_Q$ formuł spełniających $Q$
zawiera wszystkie atomy i jest zamknięty na wszystkie
operacje z ich definicji; przez tę samą minimalność
jest więc zbiorem wszystkich formuł.
Tak otrzymujemy zasadę indukcji bez zakładania z góry
rozstrzygalności dowolnej własności $P$ lub $Q$.
\end{proof}

\begin{przyklad}[Liczba wolnych zmiennych jest skończona]
Każdy term i każda formuła zawierają skończenie wiele symboli,
zatem mają skończenie wiele wystąpień zmiennych.
Zbiór różnych zmiennych wolnych jest podzbiorem zbioru wszystkich
wystąpień zmiennych, a więc jest skończony. Wynika stąd, że
zawsze możemy wybrać \emph{świeżą} zmienną spośród
$v_0,v_1,\ldots$, niewystępującą w danym skończonym napisie.
\end{przyklad}
\chapter{Struktury, wartościowania i prawda}\label{ch:semantics-fol}

\section{Co oznacza nadanie znaczenia napisowi?}
Ten sam napis $R(v_0,v_1)$ może oznaczać „$v_0$ jest mniejszy od $v_1$”,
„$v_0$ i $v_1$ są sąsiadami w grafie” lub zupełnie inną relację.
Składnia z rozdziału~\ref{ch:syntax-fol} nie rozstrzyga, która z interpretacji
jest właściwa. Znaczenie przypisujemy dopiero przez wskazanie zbioru obiektów
i interpretacji symboli pozalogicznych. Następnie definiujemy, kiedy
konkretna formuła jest spełniona. Ta druga definicja jest indukcyjna:
wartości formuł złożonych zależą od wartości ich części.

\section{Struktury i ich uniwersa}
\begin{definicja}[Struktura dla sygnatury]\label{def:sem-structure}
Niech $\mathcal L$ będzie sygnaturą z Definicji~\ref{def:syntax-signature}.
\emph{Strukturą} $\mathfrak A$ dla $\mathcal L$ nazywamy niepusty zbiór $A$,
zwany jej \emph{uniwersum}, razem z następującymi danymi:
\begin{enumerate}[label=(\roman*)]
\item dla każdego $c\in\operatorname{Const}(\mathcal L)$ element
      $c^{\mathfrak A}\in A$;
\item dla każdego symbolu funkcyjnego $f$ arności $n$ funkcja
      $f^{\mathfrak A}:A^n\to A$;
\item dla każdego symbolu relacyjnego $R$ arności $n$ relacja
      $R^{\mathfrak A}\subseteq A^n$.
\end{enumerate}
Symbol równości $=$ jest zawsze interpretowany jako rzeczywista równość
w $A$; nie jest elementem sygnatury pozalogicznej.
\end{definicja}
W tym skrypcie przyjmujemy standardową konwencję, że uniwersum jest
niepuste. Nie ogranicza to rozważań o pustych zbiorach jako obiektach:
na przykład uniwersum modelu teorii mnogości może zawierać zbiór pusty,
ale samo uniwersum modelu nie jest puste.

\begin{przyklad}[Dwie interpretacje tej samej sygnatury]
Weźmy język z jednym binarnym symbolem $R$.
Niech $\mathfrak A$ ma uniwersum $\{0,1\}$ i interpretację
$R^{\mathfrak A}=\{(0,1)\}$.
Niech $\mathfrak B$ ma to samo uniwersum, lecz
$R^{\mathfrak B}=\{(0,0),(1,1)\}$.
Wówczas zdanie $\exists v_0\,R(v_0,v_0)$ jest fałszywe
w $\mathfrak A$, a prawdziwe w $\mathfrak B$.
Mamy tę samą składnię, ale różne interpretacje.
\end{przyklad}
\begin{przyklad}[Liczby naturalne i grafy]
W języku arytmetyki standardową strukturą jest
$\mathfrak N=(\N,0,S,+,\cdot,<)$, gdzie $S(n)=n+1$.
W języku grafów struktura ma postać $(V,E)$, gdzie $V\ne\varnothing$
i $E\subseteq V\times V$. Formuła $\forall v_0\,\neg E(v_0,v_0)$
wyraża brak pętli: żaden wierzchołek nie jest połączony ze sobą.
\end{przyklad}


\begin{uwaga}[Stałe, funkcje i relacje pełnią inne role]
W strukturze arytmetycznej symbol $0$ wskazuje pojedynczy element,
$S$ i $+$ wskazują funkcje przypisujące wyniki argumentom,
a $<$ wskazuje relację, której pary mogą być prawdziwe albo fałszywe.
Dlatego napis $S(0)$ jest termem (oznacza obiekt), lecz $0<S(0)$
jest formułą (wyraża własność dwóch obiektów).
Gdy w następnym kroku przypiszemy zmiennej $v_0$ wartość $2$,
term $S(v_0)$ będzie oznaczał $3$, a formuła $v_0<S(v_0)$
będzie prawdziwa. Są to dwa różne rodzaje obliczeń.
\end{uwaga}

\section{Wartościowania i interpretacja termów}
\begin{definicja}[Wartościowanie]\label{def:sem-assignment}
Niech $\mathfrak A$ będzie strukturą o uniwersum $A$.
\emph{Wartościowaniem zmiennych} nazywamy funkcję
$s:\{v_0,v_1,\ldots\}\to A$.
Dla $a\in A$ i zmiennej $x$ symbol $s[x\mapsto a]$
oznacza wartościowanie przypisujące $x$ wartość $a$ i zgadzające
się z $s$ na wszystkich pozostałych zmiennych.
\end{definicja}
\begin{uwaga}[Dlaczego wartościowania istnieją?]
Przyjmujemy $A\ne\varnothing$. Po wybraniu jednego
elementu $a_0\in A$ możemy przypisać każdej zmiennej
wartość $a_0$; otrzymujemy wartościowanie stałe.
Nie wymaga to wyboru nieskończonej rodziny elementów:
jedno ustalone $a_0$ wystarcza dla wszystkich zmiennych.
Przy pustym uniwersum takiego wartościowania dla
nieskończonego zbioru zmiennych nie byłoby.
\end{uwaga}

\begin{definicja}[Wartość termu]\label{def:sem-term-value}
Wartość termu $t$ w strukturze $\mathfrak A$ względem wartościowania $s$
oznaczamy przez $t^{\mathfrak A}[s]$ i określamy indukcyjnie:
\begin{enumerate}[label=(\roman*)]
\item $x^{\mathfrak A}[s]=s(x)$ dla zmiennej $x$;
\item $c^{\mathfrak A}[s]=c^{\mathfrak A}$ dla stałej $c$;
\item jeśli $t=f(t_1,\ldots,t_n)$, to
      $t^{\mathfrak A}[s]=f^{\mathfrak A}(t_1^{\mathfrak A}[s],\ldots,t_n^{\mathfrak A}[s])$.
\end{enumerate}
\end{definicja}
\begin{przyklad}
Przy $s(v_0)=3$, $s(v_1)=5$ w $\mathfrak N$ term
$t=+(S(v_0),v_1)$ ma wartość
$t^{\mathfrak N}[s]=S^{\mathfrak N}(3)+5=9$.
Term domknięty $S(S(0))$ ma wartość $2$ niezależnie od $s$.
\end{przyklad}

\begin{lemat}[Zależność wartości termu od jego zmiennych]\label{lem:sem-term-coincidence}
Niech $t$ będzie termem, a $s,r$ wartościowaniami w tej samej
strukturze $\mathfrak A$. Jeżeli $s(x)=r(x)$ dla każdej zmiennej
$x\in\operatorname{FV}(t)$, to
$t^{\mathfrak A}[s]=t^{\mathfrak A}[r]$.
W szczególności wartość termu domkniętego nie zależy od wartościowania.
\end{lemat}
\begin{proof}
Przeprowadzimy indukcję po konstrukcji termu. Jeśli $t=x$,
założenie daje $s(x)=r(x)$, a więc równość wartości. Jeśli $t=c$
jest stałą, obie wartości są równe $c^{\mathfrak A}$.
Jeśli $t=f(t_1,\ldots,t_n)$, to z definicji zbioru zmiennych
wolnych $\operatorname{FV}(t_i)\subseteq\operatorname{FV}(t)$.
Zatem $s$ i $r$ zgadzają się na zmiennych każdego $t_i$.
Z założenia indukcyjnego dostajemy
$t_i^{\mathfrak A}[s]=t_i^{\mathfrak A}[r]$ dla $i=1,\ldots,n$.
Przyłożenie tej samej funkcji $f^{\mathfrak A}$ do identycznych
$n$-tek daje równość wartości $t$.
Jeśli $t$ jest domknięty, zbiór $\operatorname{FV}(t)$ jest pusty,
więc warunek zgodności wartościowań jest spełniony automatycznie.
\end{proof}

\section{Relacja spełniania: definicja rekurencyjna}
\begin{definicja}[Spełnianie formuły]\label{def:sem-satisfaction}
Niech $\mathcal L$ będzie sygnaturą, $\mathfrak A$ strukturą
języka $\mathcal L$ o uniwersum $A$, a $s$ wartościowaniem.
Dla każdej formuły $\varphi$ języka $\mathcal L$ napis
$\mathfrak A,s\models\varphi$ oznacza, że $\mathfrak A$ spełnia
formułę $\varphi$ przy $s$. Definiujemy tę relację indukcyjnie:
\begin{enumerate}[label=(\roman*)]
\item $\mathfrak A,s\models t_1=t_2$ wtedy i tylko wtedy, gdy
      $t_1^{\mathfrak A}[s]=t_2^{\mathfrak A}[s]$;
\item $\mathfrak A,s\models R(t_1,\ldots,t_n)$ wtedy i tylko wtedy, gdy
      $(t_1^{\mathfrak A}[s],\ldots,t_n^{\mathfrak A}[s])\in R^{\mathfrak A}$;
\item $\mathfrak A,s\models\neg\varphi$ wtedy i tylko wtedy, gdy
      nie zachodzi $\mathfrak A,s\models\varphi$;
\item $\mathfrak A,s\models\varphi\land\psi$ wtedy i tylko wtedy,
      gdy spełnione są obie formuły; dla $\lor$, $\to$ i $\leftrightarrow$
      używamy tablic prawdy z rozdziału~\ref{ch:propositional};
\item $\mathfrak A,s\models\forall x\,\varphi$ wtedy i tylko wtedy,
      gdy dla każdego $a\in A$ zachodzi
      $\mathfrak A,s[x\mapsto a]\models\varphi$;
\item $\mathfrak A,s\models\exists x\,\varphi$ wtedy i tylko wtedy,
      gdy dla pewnego $a\in A$ zachodzi
      $\mathfrak A,s[x\mapsto a]\models\varphi$.
\end{enumerate}
\end{definicja}
W szczególności formuła $\varphi\to\psi$ jest fałszywa dokładnie wtedy,
gdy $\varphi$ jest prawdziwa i $\psi$ fałszywa przy tym samym wartościowaniu.
Dla kwantyfikatorów zmieniamy wartość wyłącznie związanej zmiennej.

\begin{stwierdzenie}[Jednoznaczność spełniania]
\label{prop:sem-uniqueness-xlv}
Dla każdej sygnatury $\mathcal L$, struktury $\mathfrak A$
o niepustym uniwersum i wartościowania $s$ reguły
Definicji~\ref{def:sem-satisfaction} wyznaczają jednoznacznie,
które formuły $\mathcal L$ są spełnione przy $s$.
\end{stwierdzenie}
\begin{proof}
Najpierw Lematu~\ref{lem:syntax-unique-terms} używamy
rekurencyjnie do wyznaczenia wartości wszystkich termów:
zmienne i stałe mają jedną wskazaną wartość, a wartość
aplikacji $f$ jest zadana przez funkcję $f^{\mathfrak A}$.
Z Lematu~\ref{lem:syntax-unique-formulas} wynika, że
każda formuła ma jeden z określonych kroków końcowych.
Dla atomu reguły równości albo przynależności do relacji
wyznaczają wartość logiczną. Dla negacji i spójników
binarnych wartość wynika jednoznacznie z wartości
podformuł. Dla kwantyfikatora $\forall x$ rozstrzygamy,
czy wskazana podformuła jest spełniona przy
$s[x\mapsto a]$ dla wszystkich $a\in A$; dla $\exists x$
sprawdzamy, czy zachodzi to dla pewnego $a\in A$.
Indukcja po skończonej budowie formuły dowodzi zatem,
że żaden jej krok nie dopuszcza dwóch różnych ocen.
\end{proof}

\begin{przyklad}[Parzystość jako formuła]
Niech $\varphi(v_0)$ będzie formułą
$\exists v_1\,(+(v_1,v_1)=v_0)$.
W modelu $\mathfrak N$ przy wartościowaniu z $s(v_0)=8$
formuła jest spełniona: wybieramy wartość $4$ dla $v_1$.
Przy wartościowaniu $r(v_0)=7$ nie jest spełniona,
bo $n+n=7$ nie ma rozwiązania w $\N$.
Nie pytamy tu o „prawdę formuły bez kontekstu”: wynik zależy
od wartości zmiennej wolnej $v_0$.
\end{przyklad}

\begin{przyklad}[Kolejność kwantyfikatorów]
W języku z relacją $R$ rozważmy strukturę z uniwersum $\N$,
w której $R^{\mathfrak A}$ oznacza zwykły porządek $<$.
Wtedy $\forall x\,\exists y\,R(x,y)$ jest prawdziwe:
dla danego $n$ wystarczy wziąć $n+1$.
Natomiast $\exists y\,\forall x\,R(x,y)$ jest fałszywe:
dla dowolnego kandydata $m$ można wziąć $x=m$ i otrzymać
fałszywe $m<m$. Te dwa napisy różnią się kolejnością kwantyfikatorów
— i naprawdę wyrażają odmienne twierdzenia.
\end{przyklad}

\begin{przyklad}[Obliczenie spełniania w skończonej strukturze]
Niech $A=\{0,1\}$, a $R^{\mathfrak A}=\{(0,0),(1,1)\}$.
Dla każdego $a\in A$ istnieje dokładnie jeden $b\in A$
z własnością $(a,b)\in R^{\mathfrak A}$: jest nim $b=a$.
Stąd $\mathfrak A\models\forall x\,\exists y\,R(x,y)$.
Nie istnieje jednak jeden element $b\in A$, dla którego
jednocześnie $(0,b)$ i $(1,b)$ należałyby do tej relacji.
Zatem $\mathfrak A\not\models\exists y\,\forall x\,R(x,y)$.
Rozpisując świadków, dla $x=0$ wybieramy $y=0$,
a dla $x=1$ wybieramy $y=1$. W drugim zdaniu $y$ musiałoby
być wybrane \emph{przed} $x$, więc te dwa wybory nie wystarczą.
\end{przyklad}


\begin{twierdzenie}[Prawa kwantyfikatorów w semantyce pierwszego rzędu]
\label{thm:sem-quantifier-laws-xlv}
Niech $\mathcal L$ będzie sygnaturą, $\mathfrak A$ jej
strukturą o niepustym uniwersum $A$, $s$ wartościowaniem,
$\varphi$ formułą, a $x,y$ zmiennymi. Wtedy:
\begin{align*}
\mathfrak A,s\models\neg\forall x\,\varphi
 &\iff\mathfrak A,s\models\exists x\,\neg\varphi,\\
\mathfrak A,s\models\neg\exists x\,\varphi
 &\iff\mathfrak A,s\models\forall x\,\neg\varphi,\\
\mathfrak A,s\models\forall x\,\forall y\,\varphi
 &\iff\mathfrak A,s\models\forall y\,\forall x\,\varphi,\\
\mathfrak A,s\models\exists x\,\exists y\,\varphi
 &\iff\mathfrak A,s\models\exists y\,\exists x\,\varphi.
\end{align*}
Nie twierdzimy, że można dowolnie zamieniać
kwantyfikator ogólny z egzystencjalnym.
\end{twierdzenie}
\begin{proof}
Pierwsza formuła po lewej oznacza, że nie jest prawdą,
iż dla każdego $a\in A$ zachodzi
$\mathfrak A,s[x\mapsto a]\models\varphi$.
Ponieważ negujemy kwantyfikator ogólny, istnieje
$a\in A$ takie, że
$\mathfrak A,s[x\mapsto a]\not\models\varphi$.
Z definicji negacji jest to dokładnie warunek
$\mathfrak A,s[x\mapsto a]\models\neg\varphi$,
czyli prawdziwość $\exists x\,\neg\varphi$ przy $s$.
Odwrotny kierunek czytamy w tej samej równoważności
od końca do początku.

Podobnie zaprzeczenie $\exists x\,\varphi$ oznacza,
że dla żadnego $a\in A$ formuła $\varphi$ nie jest
spełniona przy $s[x\mapsto a]$. To znaczy, że
$\neg\varphi$ jest spełniona przy każdym takim
wartościowaniu, co jest definicją $\forall x\,\neg\varphi$.

Rozpiszmy trzecią równoważność: lewa strona wymaga,
aby dla każdego $a\in A$ i każdego $b\in A$ zachodziło
$\mathfrak A,s[x\mapsto a][y\mapsto b]\models\varphi$.
Gdy $x\ne y$, operacje przypisania wartości różnym
zmiennym komutują, więc otrzymujemy dokładnie warunek
po prawej, z $b$ wybranym przed $a$. Gdy $x=y$,
obydwa zdania są identyczne i równoważność jest oczywista.
Dla dwóch kwantyfikatorów egzystencjalnych zamiast
warunku dla wszystkich par $(a,b)$ otrzymujemy istnienie
przynajmniej jednej pary; kolejność jej współrzędnych
nie wpływa na istnienie świadka. Tak samo obsługujemy
przypadek $x=y$.
\end{proof}
\begin{przyklad}[Zakres prawa de Morgana dla kwantyfikatorów]
Niech $A=\{0,1\}$, a $P^{\mathfrak A}=\{0\}$ dla
jednoargumentowego symbolu relacyjnego $P$.
Zdanie $\forall x\,P(x)$ jest fałszywe, gdyż
$1\notin P^{\mathfrak A}$. Jego negacja jest więc
prawdziwa. Zarazem $1$ świadczy o prawdziwości
$\exists x\,\neg P(x)$.
Nie wynika natomiast równoważność
$\forall x\,\exists y\,R(x,y)$ i
$\exists y\,\forall x\,R(x,y)$; kontrprzykład
ze skończonej struktury podano bezpośrednio wyżej.
\end{przyklad}

\section{Lemat o zgodności wartościowań}
\begin{twierdzenie}[Spełnienie zależy tylko od zmiennych wolnych]
\label{thm:sem-coincidence}
Niech $\mathcal L$ będzie sygnaturą, $\varphi$ formułą
języka $\mathcal L$, $\mathfrak A$ strukturą tego języka
o niepustym uniwersum $A$, a $s,r$ jej wartościowaniami
zgadzającymi się na
$\operatorname{FV}(\varphi)$. Wtedy
\[
\mathfrak A,s\models\varphi
\quad\Longleftrightarrow\quad
\mathfrak A,r\models\varphi.
\]
\end{twierdzenie}
\begin{proof}
Dowód przeprowadzamy indukcją po budowie $\varphi$.
Dla formuł atomowych wystąpienia zmiennych wolnych pochodzą
z termów. Zgodność wartościowań na tych zmiennych oraz
Lemat~\ref{lem:sem-term-coincidence} dają identyczne wartości
wszystkich termów. Stąd równoważność spełniania obu typów atomów.

Dla negacji teza przechodzi natychmiast: jeśli $\psi$ jest
prawdziwa przy obu wartościowaniach albo fałszywa przy obu,
to tak samo jest z $\neg\psi$. Dla każdego spójnika binarnego
stosujemy założenie indukcyjne do obu podformuł, ponieważ
ich zmienne wolne należą do $\operatorname{FV}(\varphi)$;
tablica prawdy daje ten sam wynik.

Pozostał przypadek $\varphi=\forall x\,\psi$.
Załóżmy, że $s$ i $r$ zgadzają się na
$\operatorname{FV}(\psi)\setminus\{x\}$.
Dla dowolnego $a\in A$ wartościowania
$s[x\mapsto a]$ i $r[x\mapsto a]$ zgadzają się na
wszystkich zmiennych wolnych $\psi$: na $x$ mają tę samą
wartość $a$, a na pozostałych wymaganą zgodność zapewnia założenie.
Z założenia indukcyjnego wynika więc
\[
\mathfrak A,s[x\mapsto a]\models\psi
\quad\Longleftrightarrow\quad
\mathfrak A,r[x\mapsto a]\models\psi.
\]
Ponieważ zachodzi to dla \emph{każdego} $a\in A$,
definicja $\forall$ daje tezę. W przypadku
$\exists x\,\psi$ ten sam argument stosujemy do istnienia
odpowiedniej wartości $a$, zamiast do wszystkich $a$.
\end{proof}

\begin{wniosek}[Prawda zdania jest niezależna od wartościowania]
\label{cor:sem-sentence}
Jeżeli $\sigma$ jest zdaniem i $\mathfrak A$ strukturą,
to albo $\mathfrak A,s\models\sigma$ dla wszystkich wartościowań $s$,
albo dla żadnego. Możemy więc pisać $\mathfrak A\models\sigma$.
\end{wniosek}
\begin{proof}
Zdanie nie ma zmiennych wolnych, zatem każde dwa wartościowania
spełniają założenie Twierdzenia~\ref{thm:sem-coincidence}.
\end{proof}

\section{Podstawienie termu a jego wartość}
Lemat z tego paragrafu tłumaczy, dlaczego reguły podstawiania w dowodach
muszą zabraniać przechwytywania zmiennych. Bez tego formuła
po podstawieniu mogłaby wyrażać coś innego niż zamierzona instancja.

W zapisie $t[x:=u]$ stosujemy
Definicję~\ref{def:syntax-term-substitution-xlv}; w szczególności
podstawienie termu nie wymaga warunku unikania przechwycenia,
bo termy nie zawierają kwantyfikatorów.

\begin{lemat}[Podstawienie w termie]\label{lem:sem-term-subst}
Niech $u,t$ będą termami, $x$ zmienną, $\mathfrak A$ strukturą
i $s$ wartościowaniem. Wówczas
\[
(t[x:=u])^{\mathfrak A}[s]
=t^{\mathfrak A}[s[x\mapsto u^{\mathfrak A}[s]]].
\]
\end{lemat}
\begin{proof}
Prowadzimy indukcję po $t$. Jeżeli $t=x$, lewa strona jest
$u^{\mathfrak A}[s]$, a prawa wartością $x$ przy wartościowaniu
przypisującym $x$ dokładnie ten sam element uniwersum $A$.
Dla innej zmiennej $y$ podstawienie niczego nie zmienia,
a oba wartościowania przypisują $y$ tę samą wartość.
Dla stałej obie strony są jej interpretacją.
Jeżeli $t=f(t_1,\ldots,t_n)$, stosujemy założenie indukcyjne
do każdego $t_i$, a następnie do obu otrzymanych $n$-tek
przykładamy funkcję $f^{\mathfrak A}$.
\end{proof}

\begin{twierdzenie}[Lemat o podstawieniu]\label{thm:sem-substitution}
Niech $\varphi$ będzie formułą, $x$ zmienną, a $t$ termem
dopuszczalnym do podstawienia za $x$ w $\varphi$.
Dla każdej struktury $\mathfrak A$ i wartościowania $s$ zachodzi
\[
\mathfrak A,s\models\varphi[x:=t]
\quad\Longleftrightarrow\quad
\mathfrak A,s[x\mapsto t^{\mathfrak A}[s]]\models\varphi.
\]
\end{twierdzenie}
\begin{proof}
Prowadzimy indukcję po budowie $\varphi$. Najpierw
rozpatrzmy atom równościowy $u_1=u_2$. Po lewej
podstawiamy term $t$ do obu argumentów, więc zgodnie
z Lematem~\ref{lem:sem-term-subst} ich wartości przy $s$
równe są wartościom $u_1,u_2$ przy
$s[x\mapsto t^{\mathfrak A}[s]]$. Definicja równości
atomowej daje żądaną równoważność. Dla atomu relacyjnego
wszystkie argumenty mają odpowiednio te same wartości,
więc należą lub nie należą do tej samej relacji
$R^{\mathfrak A}$. Dla negacji i każdego spójnika
binarnego stosujemy założenie indukcyjne do odpowiednich
podformuł i reguły spełniania z
Definicji~\ref{def:sem-satisfaction}.

Pozostały kwantyfikatory. Niech $\varphi=Qy\,\psi$,
gdzie $Q\in\{\forall,\exists\}$. Rozważamy trzy przypadki.
\begin{enumerate}[label=(\roman*)]
\item Gdy $y=x$, kwantyfikator wiąże wszystkie wystąpienia
$x$ wolne w $\psi$. Stąd $x\notin\operatorname{FV}(\varphi)$,
a $\varphi[x:=t]=\varphi$. Wartościowania
$s$ i $s[x\mapsto t^{\mathfrak A}[s]]$ zgadzają się
na zmiennych wolnych $\varphi$; zgodnie z
Twierdzeniem~\ref{thm:sem-coincidence} spełniają ją
jednocześnie albo oba nie spełniają.
\item Gdy $y\ne x$, lecz $x\notin\operatorname{FV}(\psi)$,
podstawienie nic nie zmienia. Ponieważ wtedy również
$x\notin\operatorname{FV}(\varphi)$, tezę ponownie
daje twierdzenie o zgodności wartościowań.
\item Gdy $y\ne x$ oraz $x\in\operatorname{FV}(\psi)$,
każde wolne wystąpienie $x$ w $\psi$ pozostaje w zakresie
$Qy$. Dopuszczalność podstawienia implikuje więc
$y\notin\operatorname{FV}(t)$. Dla dowolnego $a\in A$
założenie indukcyjne zastosowane do $\psi$ i
$s[y\mapsto a]$ daje
\[
\mathfrak A,s[y\mapsto a]\models\psi[x:=t]
\iff
\mathfrak A,s[y\mapsto a]
  [x\mapsto t^{\mathfrak A}[s[y\mapsto a]]]
  \models\psi.
\]
Lemat~\ref{lem:sem-term-coincidence} pozwala zastąpić
$t^{\mathfrak A}[s[y\mapsto a]]$ przez
$t^{\mathfrak A}[s]$, gdyż $y$ nie jest zmienną termu $t$.
Dla różnych zmiennych $x,y$ operacje aktualizacji
wartościowania komutują, zatem ostatnie wartościowanie to
$s[x\mapsto t^{\mathfrak A}[s]][y\mapsto a]$.
Równoważność zachodzi dla każdego $a\in A$.
Jeśli $Q=\forall$, zastosowanie warunku \emph{dla wszystkich}
$a$ daje tezę; jeśli $Q=\exists$, warunku
\emph{dla pewnego} $a$.
\end{enumerate}
Każdy przypadek tworzenia formuł został rozpatrzony.
\end{proof}
\begin{przyklad}[Zmiana nazwy zmiennej związanej]
Formuły $\forall v_0\,(v_0=v_1)$ i
$\forall v_2\,(v_2=v_1)$ mają takie samo znaczenie w każdej
strukturze przy każdym wartościowaniu: w obu przypadkach
sprawdzamy równość wszystkich elementów uniwersum z wartością $v_1$.
Nie wolno natomiast dowolnie przemianowywać zmiennych tak,
aby zmienić zakres kwantyfikatora lub przechwycić wolne wystąpienie.
\end{przyklad}

\begin{przyklad}[Dlaczego lemat o podstawieniu wymaga dopuszczalności]
Niech $A=\{0,1\}$, $\varphi=\forall y\,(x=y)$,
$t=y$ oraz $s(x)=0$, $s(y)=1$. Naiwna zamiana wolnego
$x$ na $y$ daje $\forall y\,(y=y)$, prawdziwą w $A$.
Natomiast po prawej stronie wzoru z
Twierdzenia~\ref{thm:sem-substitution} występuje
$\varphi$ przy $s[x\mapsto t^{\mathfrak A}[s]]$,
czyli przy $x=1$. Formuła $\forall y\,(1=y)$ jest
fałszywa, ponieważ dla $y=0$ otrzymujemy $1\ne0$.
Wobec tego równoważność z lematu o podstawieniu bez
warunku dopuszczalności jest fałszywa. Po przemianowaniu
zmiennej związanej na świeżą $z$ otrzymujemy poprawne
$\forall z\,(y=z)$, fałszywe przy $s(y)=1$.
\end{przyklad}

\begin{twierdzenie}[Poprawność przemianowania zmiennej związanej]
\label{thm:sem-alpha-renaming-xlv}
Niech $\mathfrak A$ będzie strukturą języka $\mathcal L$,
$\psi$ formułą tego języka, a $x,y$ dwiema różnymi
zmiennymi. Załóżmy, że $y$ \emph{nie występuje w ogóle}
w napisie $\psi$ (ani wolno, ani związanie).
Wtedy dla każdego wartościowania $s$ oraz
$Q\in\{\forall,\exists\}$ zachodzi
\[
\mathfrak A,s\models Qx\,\psi
\quad\Longleftrightarrow\quad
\mathfrak A,s\models Qy\,\psi[x:=y].
\]
Podstawienie po prawej jest dopuszczalne.
\end{twierdzenie}
\begin{proof}
Ponieważ $y$ nie występuje w $\psi$, żaden kwantyfikator
w jej wnętrzu nie może przechwycić $y$ podstawianego za $x$.
Dla dowolnego $a\in A$ stosujemy
Twierdzenie~\ref{thm:sem-substitution} do formuły $\psi$,
termu $y$ i wartościowania $s[y\mapsto a]$:
\[
\mathfrak A,s[y\mapsto a]\models\psi[x:=y]
\iff
\mathfrak A,s[y\mapsto a][x\mapsto a]\models\psi.
\]
Wartościowania po prawej oraz $s[x\mapsto a]$ różnią się
co najwyżej na zmiennej $y$. Ponieważ
$y\notin\operatorname{FV}(\psi)$, twierdzenie o zgodności
wartościowań daje
\[
\mathfrak A,s[y\mapsto a]\models\psi[x:=y]
\iff
\mathfrak A,s[x\mapsto a]\models\psi.
\]
Dla $Q=\forall$ wymagamy tej równoważności dla każdego
$a\in A$, a dla $Q=\exists$ wystarczy odpowiedni
świadek $a$. Otrzymujemy oba kierunki w każdym przypadku.
\end{proof}
\begin{przyklad}[Kiedy przemianowanie jest błędne?]
Niech $\psi$ ma postać $x=y$. Wtedy $y$ nie jest świeża.
Nie wolno zastąpić $\forall x\,(x=y)$ przez
$\forall y\,(y=y)$: w strukturze
$A=\{0,1\}$ pierwsza formuła jest fałszywa przy każdym
wartościowaniu, natomiast druga jest prawdziwa.
Jeżeli zamiast $y$ weźmiemy świeżą zmienną $z$,
formuła $\forall z\,(z=y)$ jest równoważna wyjściowej
przy każdym wartościowaniu, zgodnie z twierdzeniem.
\end{przyklad}

\section{Modele, kontrmodele i konsekwencja semantyczna}
\begin{definicja}[Model i konsekwencja semantyczna]\label{def:sem-entailment}
Niech $\mathcal L$ będzie sygnaturą; wszystkie rozważane
zdania i struktury należą do tego samego języka $\mathcal L$.
Jeżeli $\sigma$ jest zdaniem, to $\mathfrak A\models\sigma$
oznacza jego prawdziwość w $\mathfrak A$.
Dla zbioru zdań $\Gamma$ piszemy $\mathfrak A\models\Gamma$,
jeśli $\mathfrak A$ spełnia każde zdanie z $\Gamma$;
wówczas $\mathfrak A$ nazywamy \emph{modelem} $\Gamma$.
Dla zdania $\sigma$ zapis $\Gamma\models\sigma$ oznacza, że
każdy model $\Gamma$ jest również modelem $\sigma$.
Jeżeli $\Gamma=\varnothing$, piszemy $\models\sigma$
i nazywamy $\sigma$ \emph{logicznie prawdziwym} (ważnym).
\emph{Kontrmodel} dla $\Gamma\models\sigma$ to struktura,
która spełnia $\Gamma$, lecz nie spełnia $\sigma$.
\end{definicja}
W tym rozdziale wystarczają zbiory \emph{zdań}. W dalszej części
skryptu rozszerzymy tę relację na formuły ze zmiennymi
wolnymi, zachowując zgodność z powyższą definicją.

\begin{przyklad}[Kontrmodel dla błędnej zmiany kwantyfikatorów]
Rozważmy formułę, której głównym spójnikiem jest implikacja
\[
\bigl(\forall x\,\exists y\,R(x,y)\bigr)
\to\bigl(\exists y\,\forall x\,R(x,y)\bigr).
\]
Nie jest ona logicznie prawdziwa. W strukturze o uniwersum $\N$,
w której $R$ oznacza $<$, poprzednik jest prawdziwy (można wziąć $y=x+1$),
a następnik fałszywy (dla dowolnego $y$ wybieramy $x=y$).
Jedna taka struktura wystarcza, aby obalić powszechną ważność implikacji.
\end{przyklad}
\begin{przyklad}[Ważność a prawdziwość w jednym modelu]
Zdanie $\forall x\,(x=x)$ jest logicznie prawdziwe:
każdy element jest równy samemu sobie, niezależnie od struktury.
Zdanie $\exists x\,\exists y\,\neg(x=y)$ jest prawdziwe
w strukturze o co najmniej dwóch elementach, ale fałszywe
w strukturze jednoelementowej. Nie jest więc logicznie prawdziwe.
\end{przyklad}

\begin{twierdzenie}[Monotoniczność konsekwencji]
\label{thm:sem-entailment-monotone-xlv}
Dla dowolnych zbiorów zdań $\Gamma\subseteq\Delta$
tej samej sygnatury i dowolnego zdania $\sigma$ zachodzi:
$\Gamma\models\sigma$ pociąga $\Delta\models\sigma$.
Jeżeli $\Gamma$ nie ma modelu, to $\Gamma\models\sigma$
dla każdego zdania $\sigma$ tej sygnatury.
\end{twierdzenie}
\begin{proof}
Niech $\mathfrak A$ będzie dowolnym modelem $\Delta$.
Ponieważ każdy element $\Gamma$ należy do $\Delta$,
struktura $\mathfrak A$ spełnia wszystkie zdania $\Gamma$.
Z $\Gamma\models\sigma$ otrzymujemy
$\mathfrak A\models\sigma$. Wybór modelu $\Delta$
był dowolny, więc $\Delta\models\sigma$.
Jeżeli $\Gamma$ nie ma modeli, nie istnieje kontrmodel:
żadna struktura nie spełnia wszystkich przesłanek
$\Gamma$ i jednocześnie nie spełnia $\sigma$.
Warunek definiujący $\Gamma\models\sigma$ zachodzi
wówczas dla każdego $\sigma$.
\end{proof}
\begin{przyklad}[Konsekwencja z niespełnialnych przesłanek]
Niech $\Gamma=\{\forall x\,\neg(x=x)\}$.
Każde uniwersum jest niepuste, więc wybieramy $a\in A$.
Atom $a=a$ jest prawdziwy, stąd
$\forall x\,\neg(x=x)$ jest fałszywe w każdej strukturze.
Zatem $\Gamma$ nie ma modelu i
$\Gamma\models\exists x\,\neg(x=x)$, chociaż zdanie
$\exists x\,\neg(x=x)$ nie jest logicznie prawdziwe.
Nie należy mylić \emph{wynikania z przesłanek}
z \emph{ważnością bez przesłanek}.
\end{przyklad}

\section{Rozszerzenie języka: logika drugiego rzędu}\label{sec:second-order-introduction}
\subsection{Jakie pytania są niedostępne w pierwszym rzędzie?}
W logice pierwszego rzędu formuła $\forall x\,\varphi(x)$ mówi o
\emph{wszystkich elementach} uniwersum $A$. Nawet jeśli dla dowolnie
wybranego podzbioru $D\subseteq A$ potrafimy opisać, że pewne zdanie
dotyczy $D$, sam kwantyfikator $\forall x$ nie obejmuje zbiorów $D$.
Aby móc powiedzieć na przykład „każdy zbiór zawierający zero
i zamknięty na następnik zawiera wszystkie liczby”, potrzebujemy
zmiennych, których wartościami są \emph{zbiory elementów}, a nie pojedyncze
elementy. Jest to główna motywacja logiki drugiego rzędu.

\begin{definicja}[Zmienne i formuły drugiego rzędu]\label{def:second-order-syntax}
Ustalmy język pierwszego rzędu $\mathcal L$ i jego zmienne
indywidualne $x,y,z,\ldots$. Rozszerzamy alfabet o nieskończenie
wiele \emph{zmiennych relacyjnych} $X^{(n)},Y^{(n)},\ldots$ dla
każdej dodatniej liczby naturalnej $n$ (jej \emph{arności}).
Jeśli $X^{(n)}$ ma arność $n$, a $t_1,\ldots,t_n$ są termami, to
$X^{(n)}(t_1,\ldots,t_n)$ jest formułą atomową. Oprócz spójników
i kwantyfikatorów pierwszego rzędu dopuszczamy
$\forall X^{(n)}\,\varphi$ i $\exists X^{(n)}\,\varphi$.
Formuły otrzymujemy przez skończone iterowanie tych konstrukcji;
nic innego nie uznajemy za formułę. Zmienną relacyjną nazywamy
\emph{wolną}, jeśli jej wystąpienie nie pozostaje w zasięgu
wiążącego ją kwantyfikatora.
Dla czytelności zmienną jednoargumentową zapisujemy często po
prostu jako $X$, a $X(x)$ odczytujemy „$x$ należy do zbioru $X$”.
\end{definicja}

\begin{przyklad}[Relacja a zbiór]
Napis $X(x)$ nie jest tym samym co $x=X$: pierwszy jest formułą
mówiącą, że element $x$ należy do zbioru reprezentowanego przez $X$;
drugi nie jest w opisanej składni poprawnym atomem, bo $x$ i $X$
są zmiennymi różnych rodzajów. Zapis $\exists X\,\forall x\,X(x)$
mówi, że istnieje zbiór zawierający każdy element uniwersum.
\end{przyklad}

\subsection{Semantyka pełna i semantyka Henkina}
\begin{definicja}[Spełnianie w pełnej semantyce drugiego rzędu]
\label{def:second-order-full-semantics}
Niech $\mathfrak A$ będzie strukturą pierwszego rzędu o
niepustym uniwersum $A$. Wartościowanie $s$ przypisuje zmiennym
indywidualnym elementy $A$, a zmiennej relacyjnej $X^{(n)}$
relację $s(X^{(n)})\subseteq A^n$. Dla atomu przyjmujemy
\[
 \mathfrak A,s\models X^{(n)}(t_1,\ldots,t_n)
 \quad\Longleftrightarrow\quad
 (t_1^{\mathfrak A}[s],\ldots,t_n^{\mathfrak A}[s])
 \in s(X^{(n)}).
\]
Pozostałe atomy, spójniki i kwantyfikatory indywidualne
interpretujemy jak w Definicji~\ref{def:sem-satisfaction}.
Dla kwantyfikatorów nowego rodzaju określamy:
\[
\begin{aligned}
\mathfrak A,s\models\forall X^{(n)}\,\varphi
 &\iff \text{dla każdego }D\subseteq A^n\text{ zachodzi }
 \mathfrak A,s[X^{(n)}\mapsto D]\models\varphi,\\
\mathfrak A,s\models\exists X^{(n)}\,\varphi
 &\iff \text{dla pewnego }D\subseteq A^n\text{ zachodzi }
 \mathfrak A,s[X^{(n)}\mapsto D]\models\varphi.
\end{aligned}
\]
Ta interpretacja jest \emph{pełną semantyką drugiego rzędu}:
zmienne relacyjne obejmują \emph{wszystkie} relacje odpowiedniej
arności na $A$.
\end{definicja}

\begin{definicja}[Semantyka Henkina]\label{def:second-order-henkin-semantics}
Niech $\mathfrak A$ będzie strukturą pierwszego rzędu
o niepustym uniwersum $A$. W \emph{semantyce Henkina} ustalamy
niepustą rodzinę $\mathcal S_n\subseteq\mathcal P(A^n)$ dla
każdej rozważanej arności $n$. Wartościowanie zmiennych
indywidualnych przyjmuje wartości w $A$, natomiast
wartościowanie każdej zmiennej relacyjnej $X^{(n)}$
musi spełniać $s(X^{(n)})\in\mathcal S_n$,
również gdy $X^{(n)}$ jest zmienną wolną.
Kwantyfikatory $\forall X^{(n)}$ i $\exists X^{(n)}$
przebiegają wówczas tylko po $D\in\mathcal S_n$.
Atom $X^{(n)}(\vec t)$ ma tę samą interpretację co powyżej. Nie zakładamy, że
$\mathcal S_n=\mathcal P(A^n)$; ewentualne dodatkowe warunki
na rodziny $\mathcal S_n$ (na przykład aksjomaty
\emph{rozumienia}) należy wysłowić oddzielnie.
\end{definicja}

\begin{przyklad}[Jedno zdanie, różne zakresy kwantyfikatora]
Niech $A=\N$ i $S(n)=n+1$. Formuła
\[
\forall X\,\Bigl(\bigl(X(0)\land
\forall x\,(X(x)\to X(S(x)))\bigr)
\to\forall x\,X(x)\Bigr)
\tag{$\mathrm{Ind}_2$}
\]
wyraża indukcję \emph{dla każdego zbioru} $X$.
Przy pełnej semantyce $X$ przebiega cały zbiór potęgowy
$\mathcal P(\N)$: założenie oznacza, że $X$ zawiera zero
i jest zamknięty na następnik, a wniosek, że $X=\N$.
W semantyce Henkina $X$ przebiega tylko rodzinę
$\mathcal S_1$. Wówczas formuła może być spełniona,
nawet jeżeli istnieją inne zewnętrzne podzbiory $A$
niedostępne dla kwantyfikatora.
\end{przyklad}

\begin{przyklad}[Model, w którym dwa odczytania indukcji się różnią]
Niech $A=\Z$, $0^{\mathfrak A}=0$ i
$S^{\mathfrak A}(z)=z+1$. Przy pełnej semantyce rozważmy
podzbiór $D=\{z\in\Z:z\geq0\}$. Zawiera on $0$
i jest zamknięty na funkcję $S^{\mathfrak A}$,
lecz nie zawiera $-1$. Podstawiając $D$ za $X$,
sprawdzamy, że poprzednik implikacji $\mathrm{Ind}_2$
jest prawdziwy, a jej następnik fałszywy.
Tak więc $\mathfrak A\not\models\mathrm{Ind}_2$
w pełnej semantyce.

Ustalmy teraz semantykę Henkina z
$\mathcal S_1=\{\Z\}$ (a dla pozostałych arności
niepustymi rodzinami, np. $\mathcal S_n=\{\Z^n\}$).
Jedyną wartością zmiennej jednoargumentowej $X$
jest wtedy całe $\Z$. Dla $X=\Z$ implikacja
$\mathrm{Ind}_2$ jest prawdziwa: spełniony jest zarówno
jej poprzednik, jak i następnik.
Zatem zdanie indukcji drugiego rzędu jest spełnione
w tej strukturze Henkina. Model nie spełnia jednak
odpowiedniego warunku pełnej semantyki.
Ten przykład nie twierdzi, że otrzymana struktura
spełnia wszystkie pozostałe aksjomaty arytmetyki Peana;
ilustruje wyłącznie rolę \emph{zakresu kwantyfikatora $X$}.
\end{przyklad}

\begin{uwaga}[Dlaczego to rozróżnienie jest ważne?]
Twierdzenia Gödla o pełności oraz Löwenheima--Skolema
w naszych dalszych rozdziałach dotyczą \emph{logiki pierwszego
rzędu}. Nie wolno przenosić ich bez dodatkowych założeń
na pełną semantykę drugiego rzędu. W
późniejszym twierdzeniu o kategoryczności arytmetyki Peana
zastosujemy właśnie formułę $\mathrm{Ind}_2$ i wykażemy, dlaczego
pełna semantyka zapewnia kategoryczność arytmetyki Peana
(nie ma niestandardowych modeli spełniających pełną indukcję
drugiego rzędu). Kategoryczność oznacza tu, że każde dwa modele
są izomorficzne, a nie że istnieje algorytm rozstrzygający
wszystkie zdania teorii.
\end{uwaga}

\section{Dlaczego pełność nie wynika z samej definicji?}
Zdefiniowaliśmy prawdziwość w strukturze. Nie rozstrzyga to,
czy każde zdanie prawdziwe we wszystkich modelach ma formalny dowód.
W kolejnym rozdziale określimy reguły dowodzenia, a następnie
wykażemy ich \emph{poprawność}: dowód formalny gwarantuje
wynikanie semantyczne. Odwrotność, czyli \emph{pełność Gödla},
wymaga odrębnej konstrukcji.
\chapter{Dowody formalne i dedukcja naturalna}\label{ch:natural-deduction}

\section{Od dowodu pisanego słowami do obiektu matematycznego}
W rozdziale~\ref{ch:quantifiers} poznaliśmy typowe techniki dowodzenia,
a w rozdziałach~\ref{ch:syntax-fol} i~\ref{ch:semantics-fol}
oddzieliliśmy formuły od ich interpretacji. Teraz uczynimy
samą procedurę dowodzenia przedmiotem definicji.
Potrzebujemy skończonego obiektu, którego każdy krok da się sprawdzić
bez zgadywania, co autor „miał na myśli”. Taką rolę spełnia
\emph{formalna dedukcja}: drzewo formuł budowane z założeń
za pomocą jawnie wypisanych reguł wnioskowania.

Wybierzemy \emph{dedukcję naturalną}, ponieważ jej reguły
odpowiadają znanym metodom argumentowania. Inne systemy,
na przykład systemy Hilberta i rachunek sekwentów, omówimy
porównawczo w dalszej części.

\section{Jak czytać zapis reguły wnioskowania?}\label{sec:nd-how-to-read}
Zanim wypiszemy formalne reguły, nauczymy się czytać ich zapis.
To ważne: pozioma kreska, którą za chwilę zobaczymy, \emph{nie jest kreską ułamkową}
i nie oznacza dzielenia liczb. Oznacza przejście od podanych zdań do nowego zdania.

\begin{definicja}[Schemat reguły wnioskowania]\label{def:nd-rule-scheme}
\emph{Reguła wnioskowania} jest umową określającą, kiedy z już uzyskanych
formuł wolno zapisać kolejną formułę. Jej ogólny schemat ma postać
\[
 \frac{\text{przesłanka 1}\quad\cdots\quad\text{przesłanka }n}
      {\text{wniosek}}\;\text{(nazwa reguły)}.
\]
To, co zapisano \emph{nad kreską}, jest wymagane przed zastosowaniem reguły;
to, co zapisano \emph{pod kreską}, jest otrzymanym wnioskiem.
Przez \emph{przesłankę} rozumiemy tutaj formułę wcześniej uzyskaną w dowodzie
lub przyjętą jako założenie. Reguła może mieć jedną przesłankę, kilka przesłanek,
a czasem żadnej. Schemat opisuje wszystkie dopuszczalne użycia jednej reguły,
nie tylko jeden konkretny dowód.
\end{definicja}

\begin{przyklad}[Dwie przesłanki nad kreską]
Jeśli umiemy uzasadnić zdanie $A$ oraz zdanie $B$, to mamy prawo
wywnioskować ich koniunkcję $A\land B$. Zapis
\[
\frac{A\qquad B}{A\land B}\;(\land I)
\]
czytamy słownie: „\emph{z A oraz B wyprowadzamy A i B}”.
Wcześniej moglibyśmy np. ustalić, że liczba $6$ jest parzysta i że jest
dodatnia. Wtedy z tych dwóch ustaleń otrzymujemy zdanie
„liczba $6$ jest parzysta i dodatnia”. Nie sprawdzamy w tym kroku
na nowo obu faktów: korzystamy z tego, że zostały już uzasadnione.
Puste miejsce między $A$ i $B$ nad kreską oznacza tylko odstęp
w drzewie: \emph{obie} przesłanki są wymagane, to nie jest symbol
alternatywy. Gdybyśmy mieli tylko jedną przesłankę, tej reguły
zastosować nie wolno.
\end{przyklad}

\begin{uwaga}[Co oznaczają litery i skróty przy regułach?]
Litery $A,B,P,Q$ lub greckie $\varphi,\psi,\theta$ występujące w schematach
są \emph{metazmiennymi}: symbolizują dowolne formuły.
Nie należy ich utożsamiać ze zmiennymi $x,y$ występującymi wewnątrz formuł.
W nazwie $(\land I)$ litera $I$ pochodzi od angielskiego
\emph{introduction} („wprowadzenie” spójnika do wniosku),
a $E$ w $(\land E)$ od \emph{elimination} („eliminacja” spójnika).
Indeks $1$ lub $2$ przy $E$ mówi, którą część koniunkcji wybieramy:
$\land E_1$ daje pierwszy człon, $\land E_2$ drugi.
Analogicznie $\lor I_1$ i $\lor I_2$ wprowadzają alternatywę,
odwołując się odpowiednio do jej pierwszego lub drugiego członu.
\end{uwaga}

\begin{przyklad}[Reguła z jedną przesłanką]
Jeśli mamy już dowód formuły $A\land B$, to zapis
\[
\frac{A\land B}{A}\;(\land E_1)
\]
pozwala nam uzyskać $A$. To nie jest równoważność obu formuł:
mając $A$, nie możemy tą samą regułą uzyskać $A\land B$,
bo nie wiemy nic o $B$. Kreskę czytamy od góry do dołu,
a nie dowolnie w obu kierunkach. Na przykład z „$6$ jest liczbą
parzystą i dodatnią” można wywnioskować „$6$ jest liczbą parzystą”,
ale z samej parzystości rozważanej liczby nie wynika jej
dodatniość: na przykład $-6$ jest parzyste i ujemne.
Odwrotne przejście byłoby nieuprawnione.
\end{przyklad}

\begin{uwaga}[Drzewo dowodu]
Kilka reguł można zestawić w \emph{drzewo dowodu}.
Liście drzewa są u góry, a ostateczny wniosek na dole.
Kiedy jedna reguła kończy się na formule $A$, tę formułę można
wykorzystać niżej jako przesłankę innej reguły.
Symbol $\vdots$ (pionowe kropki) oznacza w schematach
pewien \emph{ciąg dalszych kroków}, które na razie pomijamy na rysunku;
nie jest osobną regułą ani formułą.
Nawiasy kwadratowe $[A]^i$ oznaczają założenie tymczasowe
o \emph{etykiecie} $i$. Nie orzekają, że $A$ jest prawdziwe.
Reguła z tym samym indeksem na dole może \emph{rozładować}
wybrane wystąpienia założenia, czyli sprawić, że końcowy wniosek
nie będzie już od niego zależał jako od otwartego założenia.
Wyjaśnimy to bezpośrednio na przykładach poniżej.
\end{uwaga}


\section{Co kryją trzy pionowe kropki?}\label{sec:nd-vertical-dots}
Najbardziej tajemniczym elementem drzewa dowodu może być symbol
$\vdots$. Przypomina trzy pionowo ustawione kropki i jest
\emph{skrótem graficznym}, a nie nowym spójnikiem, kwantyfikatorem,
regułą wnioskowania ani twierdzeniem. Informuje, że pomiędzy
wypisanymi formułami znajduje się fragment formalnego dowodu,
którego kolejnych kroków nie pokazano.

\begin{definicja}[Znaczenie symbolu $\vdots$ w drzewie dowodu]
Jeżeli zapisujemy fragment
\[
\begin{array}{c}
\mathcal D\\[-2pt]
\vdots\\[-2pt]
\varphi
\end{array}
\]
to symbol $\mathcal D$ oznacza wcześniej rozpoczętą dedukcję,
a $\vdots$ zastępuje skończone drzewo poprawnych zastosowań
reguł prowadzące do formuły $\varphi$.
Skrót jest dopuszczalny tylko wtedy, gdy takie drzewo
rzeczywiście istnieje, a jego założenia i warunki użycia
reguł są zgodne z otaczającą częścią dowodu.
Nie istnieje reguła pozwalająca z dowolnej formuły $\psi$
przejść za pomocą samych kropek do dowolnej formuły $\varphi$.
\end{definicja}

\begin{przyklad}[Od skrótu do wszystkich kroków]
Weźmy formułę $A\land B$. Jeżeli zamierzamy wyprowadzić
$B\land A$, możemy skrótowo napisać
\[
\begin{array}{c}
[A\land B]^1\\
\vdots\\
B\land A
\end{array}
\]
Jednak czytelnik ma prawo zapytać, \emph{jakimi regułami}
doszliśmy do ostatniej formuły. Rozwinięty fragment wygląda tak:
\[
\frac{
  \displaystyle\frac{[A\land B]^1}{B}\;(\land E_2)
  \qquad
  \displaystyle\frac{[A\land B]^1}{A}\;(\land E_1)
}{B\land A}\;(\land I).
\]
Zatem trzy kropki w pierwszym zapisie zastępowały
trzy konkretne kroki: uzyskanie $B$, uzyskanie $A$
i ich połączenie w $B\land A$.
\emph{W żadnym miejscu nie założyliśmy dodatkowo $B\land A$}.
Aby otrzymać dowód bez otwartych założeń,
można dopiero teraz rozładować $[A\land B]^1$
regułą $(\to I)_1$, otrzymując $(A\land B)\to(B\land A)$.
\end{przyklad}

\begin{uwaga}[Kropki nie zmieniają statusu założenia]
Na każdym etapie pominiętego fragmentu trzeba odróżnić
\emph{wnioski wyprowadzone} od \emph{założeń nadal otwartych}.
Jeżeli we fragmencie powyższego przykładu wykorzystaliśmy
$[A\land B]^1$, to otrzymane $B\land A$ wciąż zależy
od tego założenia. Sam symbol $\vdots$ go nie rozładowuje.
Rozładowanie następuje dopiero przez regułę z wyraźnym
oznaczeniem indeksu, taką jak $(\to I)_1$.
\end{uwaga}

\begin{przyklad}[Czego kropki \emph{nie} usprawiedliwiają]
Gdybyśmy bez dodatkowych założeń napisali
\[
\begin{array}{c}
A\\
\vdots\\
B
\end{array}
\]
nie wynikałoby z tego, że $A\vdash_{\mathrm{ND}}B$.
Na przykład z samego „$2$ jest liczbą parzystą”
nie da się w logice wyprowadzić dowolnego zdania
„$3$ jest liczbą parzystą”. Aby uznać taki zapis
za rzeczywisty dowód, trzeba podać pominięte kroki
lub odesłać do \emph{wcześniej udowodnionego} twierdzenia.
W szczególności $\vdots$ nie jest odpowiednikiem
przyzwolenia „tu dzieje się coś oczywistego”.
\end{przyklad}

\begin{przyklad}[Kropki w dwóch niezależnych gałęziach]
W regule eliminacji alternatywy pojawiają się dwa
poddowody, na ogół o \emph{różnej} liczbie kroków:
\[
\frac{A\lor B\qquad
  \begin{array}{c}[A]^i\\\vdots\\ C\end{array}
  \qquad
  \begin{array}{c}[B]^j\\\vdots\\ C\end{array}
}{C}\;(\lor E)_{i,j}.
\]
Pierwsza para kropek zastępuje dowód $C$ przy założeniu $A$,
a druga --- dowód tego samego $C$ przy założeniu $B$.
Nie zakładamy, że oba fragmenty są identyczne.
Na końcu rozładowujemy odpowiednio $A$ i $B$;
pozostałe założenia z obu gałęzi nadal pozostają otwarte.
Warunkiem stosowania reguły jest istnienie \emph{obu}
fragmentów, nie tylko jednej gałęzi.
Dlaczego? Z $A\lor B$ nie wiemy, który przypadek zachodzi.
Jeżeli tylko z $A$ udowodniliśmy $C$, rzeczywisty przypadek
może być $B$ i wówczas $C$ nie musi zachodzić.
Przykładowo przy wartościach $A=1$, $B=0$, $C=0$
alternatywa $A\lor B$ jest prawdziwa, ale z samego
warunkowego rozumowania „gdyby $B$, to $C$”
nie można otrzymać $C$. Oba poddowody są potrzebne.
Natomiast nie zakładamy, że \emph{równocześnie} zachodzą $A$ i $B$:
to są dwa oddzielne, warunkowe rozumowania.
\end{przyklad}

\begin{uwaga}[Trzy różne rodzaje wielokropków]
Warto rozróżniać $\vdots$ (pionowe kropki, oznaczające
pominięty fragment \emph{dedukcji}), $\cdots$ (poziomy
wielokropek, np. między kilkoma przesłankami lub elementami
ciągu) oraz $\ddots$ (kropki ukośne, używane np.
w macierzach). To tylko konwencje \emph{rysowania};
żaden z tych symboli nie zastępuje logicznej reguły.
\end{uwaga}

\section{Założenia, wnioski i rozładowywanie założeń}
\begin{definicja}[Dedukcja i dowodliwość]\label{def:nd-derivation}
Ustalmy język pierwszego rzędu $\mathcal L$.
\emph{Dedukcją} formuły $\varphi$ z założeń $\Gamma$ nazywamy
skończone drzewo formuł, którego liście są dozwolonymi
założeniami (z $\Gamma$ albo tymczasowymi) lub instancjami
reguł bez przesłanek, a każda pozostała formuła wynika
z formuł bezpośrednio ponad nią na mocy jednej z reguł
podanych poniżej. Tymczasowe założenia mogą być
\emph{rozładowane} przez odpowiednią regułę: od tej chwili
nie należą do otwartych założeń dedukcji.

Zapis $\Gamma\vdash_{\mathrm{ND}}\varphi$ oznacza, że istnieje
taka dedukcja z wnioskiem $\varphi$, której wszystkie
\emph{nierozładowane} założenia należą do $\Gamma$.
Dla $\Gamma=\varnothing$ piszemy $\vdash_{\mathrm{ND}}\varphi$.
\end{definicja}

\begin{przyklad}[Co oznacza rozładowanie założenia?]\label{ex:nd-discharge-explained}
Załóżmy, że chcemy udowodnić implikację „jeśli $A$, to $A$”.
Wewnątrz dowodu możemy na próbę założyć $A$. Skoro założyliśmy $A$,
mamy już $A$ jako wniosek tego krótkiego poddowodu. Teraz kończymy
okres obowiązywania założenia i zapisujemy implikację $A\to A$:
\[
\frac{[A]^1}{A\to A}\;(\to I)_1.
\]
Liczba $1$ przy nawiasie i przy nazwie reguły wiąże te dwa miejsca.
Nie znaczy ona, że formuła $A$ została udowodniona bez założeń.
Bez założeń otrzymaliśmy dopiero \emph{implikację} $A\to A$.
W dłuższym dowodzie wewnątrz rozładowywanego fragmentu
mogą znajdować się inne kroki i kilka użyć tego samego założenia.
\end{przyklad}


\subsection{Co oznacza nawias przy założeniu?}\label{sec:nd-assumption-scope}
\begin{uwaga}[Założenie nie jest stwierdzeniem prawdziwości]\label{rem:nd-assumption-not-truth}
Napis $[A]^i$ czytamy: „\emph{na potrzeby tego poddowodu zakładamy $A$;
założenie oznaczamy numerem $i$}”. Nawias kwadratowy nie jest spójnikiem
logicznym, a $i$ nie jest wykładnikiem potęgi. To tylko oznaczenie
\emph{zależności dowodu od przyjętego założenia}.

Nie twierdzimy, że $A$ jest prawdziwe „wewnątrz nawiasu” ani że
staje się fałszywe po wyjściu z nawiasu. W ogóle nie ustalamy tutaj
wartości logicznej $A$. W poddowodzie \emph{wolno posługiwać się $A$}
bez wcześniejszego uzasadnienia; po jego rozładowaniu nie wolno
powoływać się na \emph{to konkretne tymczasowe założenie}.
Jeżeli jednak $A$ zostało niezależnie dowiedzione albo pozostaje
jednym z trwałych założeń $\Gamma$, to możemy nadal z niego korzystać.
\end{uwaga}

\begin{przyklad}[Założenie może być nawet fałszywe]\label{ex:nd-false-assumption}
Niech $A$ oznacza zdanie „liczba $2$ jest nieparzysta”.
Wiemy, że $A$ jest fałszywe, ale nadal wolno rozważyć
\emph{hipotetyczny} dowód rozpoczynający się od $[A]^3$.
Z tego założenia otrzymujemy samo $A$; rozładowanie daje
\[
\frac{[A]^3}{A\to A}\;(\to I)_3.
\]
Nie udowodniliśmy w ten sposób, że $2$ jest nieparzyste.
Udowodniliśmy jedynie zdanie „\emph{jeśli $2$ jest nieparzyste,
to $2$ jest nieparzyste}”, czyli implikację $A\to A$.
Jest ona prawdziwa również wtedy, gdy $A$ jest fałszywe.
Właśnie dlatego rozładowanie założenia nie może oznaczać
zmiany jego wartości logicznej.
\end{przyklad}

\begin{przyklad}[Jak zmienia się zbiór otwartych założeń]\label{ex:nd-open-assumptions}
Ustalmy dowolne formuły $A$ i $B$.
Zademonstrujemy kolejno dwa zagnieżdżone poddowody:
\[
\frac{
 \displaystyle\frac{
   \displaystyle\frac{[A]^1\qquad[B]^2}{A\land B}\;(\land I)
 }{B\to(A\land B)}\;(\to I)_2
}{A\to\bigl(B\to(A\land B)\bigr)}\;(\to I)_1.
\]
Odczytujemy drzewo \emph{od góry do dołu}.
Na początku założyliśmy jednocześnie $A$ (etykieta $1$)
i $B$ (etykieta $2$), dlatego wolno nam utworzyć $A\land B$.
W tym momencie wniosek zależy od \emph{obu} tymczasowych założeń.

Teraz stosujemy $(\to I)_2$. Zamykamy założenie oznaczone $2$,
więc otrzymujemy $B\to(A\land B)$, ale nadal korzystamy z $[A]^1$:
\emph{otwarte pozostaje tylko założenie $A$}.
Dopiero zastosowanie $(\to I)_1$ zamyka także założenie $A$.
Końcowa formuła $A\to(B\to(A\land B))$ nie zależy już
od żadnego otwartego założenia, czyli
$\vdash_{\mathrm{ND}}A\to(B\to(A\land B))$.

Nie otrzymujemy natomiast ani $\vdash_{\mathrm{ND}}A$,
ani $\vdash_{\mathrm{ND}}B$. Po zakończeniu poddowodu
nie wolno po prostu „wyciągnąć” $A$ albo $B$ jako wniosku
niezależnego od założeń.
\end{przyklad}

\begin{uwaga}[Rozładowanie przy dowodzie nie wprost]
Podobnie działa reguła wprowadzania negacji:
\[
\frac{[A]^i\quad\vdots\quad\bot}{\neg A}\;(\neg I)_i.
\]
Jeżeli przyjęcie $A$ prowadzi do sprzeczności, to po
zamknięciu tego założenia możemy wnioskować $\neg A$.
To kolejny argument, że zapis $[A]^i$ \emph{nie ogłasza
prawdziwości $A$}: wewnątrz hipotetycznego poddowodu
możemy dojść właśnie do wniosku, że $A$ jest niemożliwe
(przy pozostałych otwartych założeniach).
\end{uwaga}

Przejście od poddowodu do implikacji streszcza zatem schemat
\[
\Gamma\cup\{A\}\vdash_{\mathrm{ND}}B
\quad\Longrightarrow\quad
\Gamma\vdash_{\mathrm{ND}}A\to B,
\]
gdzie $\Gamma$ oznacza zbiór pozostałych, dostępnych założeń.
Jego pełne uzasadnienie podamy w
Twierdzeniu~\ref{thm:nd-deduction}.

Ściśle rzecz biorąc, w drzewie trzeba odnotować, które
wystąpienia założeń zostają rozładowane. Będziemy używać indeksów
$[\varphi]^i$ oraz zapisu „rozładowujemy $i$”.

Dla wygody symbol $\bot$ będzie oznaczał ustaloną zamkniętą
formułę fałszywą w każdej strukturze, na przykład
$\neg\forall v_0\,(v_0=v_0)$.
Nie dodajemy więc nowego symbolu do języka: $\bot$ jest skrótem.
Znaczenie tego skrótu jest uzasadnione
Definicją~\ref{def:sem-satisfaction} i zwykłą własnością równości.

\section{Reguły dla koniunkcji i implikacji}
\begin{definicja}[Reguły $\land$ oraz $\to$]\label{def:nd-rules-conj-imp}
W dedukcji naturalnej dopuszczamy następujące schematy:
\[
\frac{\varphi\qquad\psi}{\varphi\land\psi}\;(\land I),
\qquad
\frac{\varphi\land\psi}{\varphi}\;(\land E_1),
\qquad
\frac{\varphi\land\psi}{\psi}\;(\land E_2),
\]
\[
\frac{[\varphi]^i\quad\vdots\quad\psi}
     {\varphi\to\psi}\;(\to I)_i,
\qquad
\frac{\varphi\to\psi\qquad\varphi}{\psi}\;(\to E).
\]
Reguła $(\to I)_i$ rozładowuje oznaczone wystąpienia
założenia $\varphi$ użyte w poddedukcji $\psi$.
\end{definicja}
\begin{przyklad}[Pierwszy formalny dowód]\label{ex:nd-conj-swap}
Udowodnimy $\vdash_{\mathrm{ND}}(P\land Q)\to(Q\land P)$.
Rozpoczynamy od tymczasowego założenia $P\land Q$.
Reguła $\land E_2$ daje $Q$, a reguła $\land E_1$ daje $P$.
Reguła $\land I$ pozwala połączyć te dwa wnioski w $Q\land P$.
Następnie $\to I$ rozładowuje początkowe założenie,
pozostawiając bezzałożeniowy dowód implikacji.
Pełne drzewo (a nie tylko jego fragment) ma postać
\[
\frac{\displaystyle
 \frac{
 \frac{[P\land Q]^1}{Q}\;(\land E_2)
 \qquad
 \frac{[P\land Q]^1}{P}\;(\land E_1)
 }{Q\land P}\;(\land I)
}{(P\land Q)\to(Q\land P)}\;(\to I)_1.
\]
Sprawdzamy je od góry: dwa górne liście mają tę samą postać $P\land Q$.
Z lewego odczytujemy $Q$, z prawego $P$. Łączymy je w $Q\land P$.
Na końcu zamykamy tymczasowe założenie o numerze $1$,
co pozwala wprowadzić implikację. Pod kreską końcową nie ma
już żadnego nierozładowanego założenia, stąd znak $\vdash_{\mathrm{ND}}$.
\end{przyklad}

\begin{twierdzenie}[Twierdzenie o dedukcji dla dedukcji naturalnej]
\label{thm:nd-deduction}
Dla dowolnego zbioru formuł $\Gamma$ i dowolnych formuł
$\varphi,\psi$ zachodzi
\[
\Gamma\cup\{\varphi\}\vdash_{\mathrm{ND}}\psi
\quad\Longleftrightarrow\quad
\Gamma\vdash_{\mathrm{ND}}\varphi\to\psi.
\]
\end{twierdzenie}
\begin{proof}
Załóżmy najpierw, że istnieje dedukcja $\psi$,
której wszystkie otwarte założenia należą do
$\Gamma\cup\{\varphi\}$. Jeśli niektóre otwarte liście
są wystąpieniami $\varphi$, oznaczamy je wspólnym indeksem
$i$. Zastosowanie reguły $\to I_i$ do całej dedukcji
rozładowuje te wystąpienia i daje $\varphi\to\psi$.
Pozostałe otwarte założenia pochodzą z $\Gamma$.
Nawet jeśli w poddedukcji nie było otwartego wystąpienia
$\varphi$, reguła $\to I$ jest dozwolona:
implikacja z dodatkowym, nieużytym założeniem też ma dowód.

W przeciwną stronę załóżmy, że
$\Gamma\vdash_{\mathrm{ND}}\varphi\to\psi$.
Do otrzymanego drzewa dopisujemy dodatkowy liść
$\varphi$ i stosujemy $\to E$. Wniosek to $\psi$,
a otwarte założenia należą do $\Gamma\cup\{\varphi\}$.
\end{proof}

\section{Alternatywa, równoważność i negacja}
\begin{definicja}[Pozostałe reguły zdaniowe]\label{def:nd-rules-prop}
Dla alternatywy przyjmujemy
\[
\frac{\varphi}{\varphi\lor\psi}\;(\lor I_1),
\quad
\frac{\psi}{\varphi\lor\psi}\;(\lor I_2),
\quad
\frac{\varphi\lor\psi\qquad
[\varphi]^i\;\vdots\;\theta\qquad
[\psi]^j\;\vdots\;\theta}{\theta}\;(\lor E)_{i,j}.
\]
Ostatnia reguła rozładowuje założenie $\varphi$
w pierwszej poddedukcji oraz założenie $\psi$
w drugiej. Jej sens jest następujący: rozpatrujemy oba
możliwe przypadki i dochodzimy do tego samego wniosku.

Dla negacji i sprzeczności przyjmujemy
\[
\frac{\varphi\qquad\neg\varphi}{\bot}\;(\neg E),
\quad
\frac{[\varphi]^i\;\vdots\;\bot}{\neg\varphi}\;(\neg I)_i,
\quad
\frac{\bot}{\theta}\;(\bot E).
\]
Ponieważ używamy logiki klasycznej, dodajemy regułę
\emph{dowodu nie wprost}:
\[
\frac{[\neg\varphi]^i\;\vdots\;\bot}{\varphi}\;(\mathrm{RAA})_i.
\]
Dla równoważności reguły zapisujemy jawnie:
\[
\frac{\varphi\to\psi\qquad\psi\to\varphi}
     {\varphi\leftrightarrow\psi}\;(\leftrightarrow I),
\quad
\frac{\varphi\leftrightarrow\psi}{\varphi\to\psi}\;(\leftrightarrow E_1),
\quad
\frac{\varphi\leftrightarrow\psi}{\psi\to\varphi}\;(\leftrightarrow E_2).
\]
Pierwsza reguła wymaga obu implikacji, druga i trzecia wydobywają
z równoważności odpowiednio pierwszy lub drugi kierunek.
\end{definicja}
\begin{uwaga}[Klasyczność ma znaczenie]
Reguła $\mathrm{RAA}$ nie jest regułą intuicjonistycznej
dedukcji naturalnej. Odłożenie jej prowadzi do słabszego
pojęcia wyprowadzalności. Nie będziemy zamieniać tych
dwóch rachunków bez wyraźnego zaznaczenia zmiany założeń.
\end{uwaga}

\begin{twierdzenie}[Prawo wyłączonego środka]
\label{thm:nd-lem}
Dla każdej formuły $\varphi$ zachodzi
$\vdash_{\mathrm{ND}}\varphi\lor\neg\varphi$.
\end{twierdzenie}
\begin{proof}
Załóżmy tymczasowo $\neg(\varphi\lor\neg\varphi)$.
Jeśli dodatkowo założymy $\varphi$, to przez $\lor I_1$
otrzymamy $\varphi\lor\neg\varphi$.
Razem z pierwszym założeniem daje to $\bot$ przez $\neg E$.
Rozładowując założenie $\varphi$ regułą $\neg I$,
otrzymujemy $\neg\varphi$.
Stosując teraz $\lor I_2$, otrzymujemy ponownie
$\varphi\lor\neg\varphi$, sprzeczne z pierwszym założeniem.
Zatem z $\neg(\varphi\lor\neg\varphi)$ wyprowadziliśmy $\bot$.
Reguła $\mathrm{RAA}$ rozładowuje to założenie
i daje $\varphi\lor\neg\varphi$ bez żadnych otwartych założeń.
\end{proof}

\section{Kwantyfikatory i warunki na zmienne}
Na tym etapie ujawnia się ważna różnica między rachunkiem zdań
a logiką pierwszego rzędu. Nie każda zmienna może być
uogólniana, a podstawienie nie może przechwytywać zmiennych.
Zlekceważenie tych warunków prowadziłoby do fałszywych twierdzeń.

\begin{definicja}[Reguły kwantyfikatorów]\label{def:nd-rules-quant}
Niech $t$ będzie termem dopuszczalnym do podstawienia
za $x$ w $\varphi$ (Definicja~\ref{def:syntax-free-for}).
Stosujemy reguły
\[
\frac{\forall x\,\varphi}{\varphi[x:=t]}\;(\forall E),
\qquad
\frac{\varphi}{\forall x\,\varphi}\;(\forall I),
\]
\[
\frac{\varphi[x:=t]}{\exists x\,\varphi}\;(\exists I).
\]
Reguła $\forall I$ jest dozwolona tylko wtedy, gdy $x$
nie występuje wolno w \emph{żadnym otwartym założeniu}
poddedukcji kończącej się na $\varphi$.

W regule eliminacji kwantyfikatora egzystencjalnego
stosujemy zmienną $y$, niewystępującą w ogóle
w $\varphi$ ani w końcowym wniosku $\theta$,
a ponadto niewystępującą wolno w żadnym innym
otwartym założeniu rzeczywiście użytym w obu
gałęziach dowodu (poza rozładowywanym założeniem
$[\varphi[x:=y]]^i$):
\[
\frac{\exists x\,\varphi\qquad
[\varphi[x:=y]]^i\;\vdots\;\theta}
{\theta}\;(\exists E)_i.
\]
Reguła rozładowuje tymczasowe założenie
$\varphi[x:=y]$. Warunek świeżości dotyczy obu gałęzi:
$y$ nie może występować wolno ani w pozostałych otwartych
założeniach poddedukcji, ani w otwartych założeniach dowodu
przesłanki $\exists x\,\varphi$.
Nie wymagamy natomiast, aby $y$ nie występowała w żadnej
formule całego zbioru $\Gamma$: zbiór ten może być nieskończony,
a konkretna dedukcja korzysta tylko ze skończenie wielu
otwartych założeń. W przeciwnym razie połączenie obu gałęzi
mogłoby przenieść informację o konkretnym $y$ do wniosku,
chociaż świadek kwantyfikatora egzystencjalnego nie musi
być wartością $y$. Przed zastosowaniem reguły w razie
potrzeby zmieniamy nazwy zmiennych związanych,
aby podstawienie było dopuszczalne.
\end{definicja}

\begin{przyklad}[Jak czytać reguły kwantyfikatorów]
Jeżeli znamy $\forall x\,P(x)$, wolno wstawić np. term $0$ za $x$
i otrzymać $P(0)$; taki krok nosi nazwę $(\forall E)$.
Jeżeli natomiast ustaliliśmy $P(0)$, otrzymujemy przez $(\exists I)$
formułę $\exists x\,P(x)$: liczba oznaczona przez term $0$ jest
konkretnym świadkiem istnienia. Te kroki pokazuje zapis
\[
\frac{\forall x\,P(x)}{P(0)}\;(\forall E),
\qquad
\frac{P(0)}{\exists x\,P(x)}\;(\exists I).
\]
Nie wolno jednak odwrócić pierwszej kreski: z jednego przypadku
$P(0)$ nie wynika twierdzenie o wszystkich elementach.
\end{przyklad}

\begin{przyklad}[Co znaczą kropki w eliminacji egzystencjalnej]
Załóżmy, że wiadomo tylko $\exists x\,P(x)$ i chcemy pokazać
zdanie $\theta$ niezawierające zmiennej $y$.
Tymczasowo przyjmujemy $P(y)$, gdzie $y$ jest nową nazwą,
nieobecną w innych otwartych założeniach i we wniosku $\theta$.
Jeżeli z tego założenia potrafimy wykonać jakieś kroki i dojść
do $\theta$, w schemacie oznaczamy je pionowymi kropkami:
\[
\frac{\exists x\,P(x)\qquad [P(y)]^2\;\vdots\;\theta}
{\theta}\;(\exists E)_2.
\]
Kropki nie pozwalają dowolnie przeskakiwać kroków: w rzeczywistym dowodzie
muszą odpowiadać konkretnemu, sprawdzalnemu drzewu
(zob. sekcję~\ref{sec:nd-vertical-dots}). Indeks $2$ wskazuje,
że założenie $P(y)$ po użyciu reguły nie pozostaje otwarte.
\end{przyklad}

\begin{przyklad}[Dlaczego ograniczenie dla $\forall I$ jest konieczne]
Z założenia $P(x)$ nie wolno wywnioskować $\forall x\,P(x)$.
W strukturze o uniwersum $\{0,1\}$ niech $P$ oznacza zbiór $\{0\}$.
Przy wartościowaniu $s(x)=0$ założenie $P(x)$ jest prawdziwe,
a $\forall x\,P(x)$ jest fałszywe. Niepoprawny krok
polegałby na zastosowaniu $\forall I$ do zmiennej wolnej
w otwartym założeniu $P(x)$.
\end{przyklad}

\begin{przyklad}[Dlaczego świadek egzystencjalny jest lokalny]
Z $\exists x\,P(x)$ nie możemy po prostu wyprowadzić
$P(y)$, jeśli $y$ jest ustaloną wcześniej zmienną.
W strukturze o dwóch elementach, przy $P$ prawdziwym
tylko dla jednego z nich, poprzednik jest prawdziwy,
a $P(y)$ może być fałszywe dla wartościowania przypisującego
$y$ drugi element. Reguła $\exists E$ pozwala użyć
świeżej zmiennej \emph{wewnątrz poddedukcji}, ale nie
pozwala jej przemycić do końcowego wniosku.
\end{przyklad}

\begin{twierdzenie}[Przykład dowodzenia z kwantyfikatorami]
\label{thm:nd-univ-chain}
Dla formuł jednoargumentowych $P(x)$ i $Q(x)$ zachodzi
\[
\bigl\{\forall x\,(P(x)\to Q(x)),\ \forall x\,P(x)\bigr\}
\vdash_{\mathrm{ND}}\forall x\,Q(x).
\]
\end{twierdzenie}
\begin{proof}
Oznaczmy oba założenia odpowiednio przez $H_1$ i $H_2$.
Zastosujmy $\forall E$, podstawiając za $x$ samą
zmienną $x$. Takie podstawienie jest zawsze dopuszczalne,
nawet gdy $P(x)$ lub $Q(x)$ zawierają wewnętrzne
kwantyfikatory wiążące inne wystąpienia zmiennych.
Z $H_1$ otrzymujemy $P(x)\to Q(x)$,
a z $H_2$ otrzymujemy $P(x)$.
Reguła $\to E$ daje $Q(x)$.
W żadnym z otwartych założeń $H_1,H_2$ zmienna $x$
nie występuje wolno: wszystkie jej wystąpienia
w tych zdaniach są związane przez zewnętrzne kwantyfikatory.
Możemy zatem zastosować $\forall I$ i otrzymać
bezpośrednio $\forall x\,Q(x)$.
\end{proof}

\begin{twierdzenie}[Od istnienia przypadku do istnienia własności]
\label{thm:nd-exists-projection}
Niech $P(x)$ i $Q(x)$ będą formułami o jedynej zmiennej wolnej $x$.
Wtedy
\[
\exists x\,(P(x)\land Q(x))
\vdash_{\mathrm{ND}}\exists x\,P(x).
\]
\end{twierdzenie}
\begin{proof}
Przyjmijmy założenie $H:=\exists x\,(P(x)\land Q(x))$.
Wybierzmy nową zmienną $y$, niewystępującą w $H$ ani
we wniosku. W poddedukcji przyjmijmy tymczasowo
$P(y)\land Q(y)$. Reguła $\land E_1$ daje $P(y)$,
a następnie $\exists I$ daje $\exists x\,P(x)$.
Zastosowanie $\exists E$ rozładowuje założenie
$P(y)\land Q(y)$; wniosek nie zawiera wolnego $y$,
więc jest dozwolony. Jedynym pozostałym założeniem jest $H$.
\end{proof}

\begin{przyklad}[Dlaczego świeżość obejmuje także przesłankę główną]
\label{ex:nd-exists-major-premise-fresh-46}
Najpierw rozważmy błędne użycie zmiennej występującej wolno
w \emph{innym otwartym założeniu}. Niech $P$ będzie
predykatem jednoargumentowym, a $R$ skrótem zdania
$\forall z\,P(z)$. Z założeń
$\exists x\,P(x)$ oraz $P(y)\to R$ nie wolno
wywnioskować $R$ przez poddedukcję od $[P(y)]^i$:
chociaż z $P(y)$ i drugiego założenia otrzymujemy $R$,
zmienna $y$ pozostaje wolna w $P(y)\to R$.

W uniwersum $\{0,1\}$ z $P^{\mathfrak A}=\{0\}$
oraz $s(y)=1$ obydwa założenia są spełnione,
lecz $R=\forall z\,P(z)$ jest fałszywe.
To kontrmodel dla takiego niepoprawnego wnioskowania.

Osobno sprawdźmy konieczność świeżości względem
\emph{przesłanki głównej}. Niech teraz
$H:=\exists x\,\neg(x=y)$, a jako żądany wniosek
weźmy fałszywe zdanie
$\theta:=\forall z\,\neg(z=z)$.
Gdyby wolno było użyć $y$ jako zmiennej świadka,
otrzymalibyśmy tymczasowe założenie
$[\neg(y=y)]^i$. Reguła $=I$ daje $y=y$,
więc $\neg E$ daje $\bot$, a następnie $\bot E$
pozwala wyprowadzić $\theta$.
Taki schemat jest \emph{niedozwolony}, ponieważ $y$
występuje wolno w $H$.
W dwuelementowym uniwersum $H$ jest prawdziwe
przy każdym wartościowaniu $y$ (istnieje inny element),
lecz $\theta$ jest fałszywe.
Oba kontrprzykłady pokazują, że świeżość
jest konieczna do zachowania prawdy.
\end{przyklad}

\section{Równość i podstawianie równych obiektów}
\begin{definicja}[Reguły równości]\label{def:nd-equality}
Dopuszczamy regułę $=I$: dla dowolnego termu $t$
można bez założeń wyprowadzić $t=t$. Zapisujemy ją jako
$\frac{}{t=t}\;(=I)$. Puste miejsce nad kreską znaczy,
że ta reguła nie wymaga żadnej wcześniejszej przesłanki.
Reguła $=E$ ma postać: z $t=u$ i
$\theta[x:=t]$ wolno wyprowadzić $\theta[x:=u]$,
o ile obydwa podstawienia są dopuszczalne.
\end{definicja}
\begin{przyklad}[Symetria równości]
Z założenia $t=u$ chcemy wyprowadzić $u=t$.
Najpierw przez $=I$ otrzymujemy $t=t$.
Weźmy formułę $\theta(x)$ postaci $x=t$.
Jej instancją $\theta[x:=t]$ jest $t=t$,
a instancją $\theta[x:=u]$ jest $u=t$.
Z przesłanek $t=u$ i $t=t$ reguła $=E$ daje $u=t$.
\end{przyklad}
\begin{przyklad}[Niesprzeczność wynika z prawdziwości równości]
Dla każdego termu $t$ mamy bezzałożeniowy dowód $t=t$.
Gdyby z jakichś założeń uzyskano $\neg(t=t)$,
reguła $\neg E$ po połączeniu z $t=t$ dawałaby $\bot$.
Jest to prosty przykład konfliktu formalnego z podstawowym
prawem logiki z równością.
\end{przyklad}

\section{Własność skończonego użycia założeń}
\begin{twierdzenie}[Skończona liczba użytych przesłanek]
\label{thm:nd-finite-support}
Jeżeli $\Gamma\vdash_{\mathrm{ND}}\varphi$, to istnieje
skończony zbiór $\Gamma_0\subseteq\Gamma$ taki, że
$\Gamma_0\vdash_{\mathrm{ND}}\varphi$.
\end{twierdzenie}
\begin{proof}
Dowód formalny jest z definicji skończonym drzewem.
Ma więc tylko skończenie wiele liści. Wybierzmy spośród
nich te, które pozostają otwartymi założeniami, i zbierzmy
występujące w nich formuły do zbioru $\Gamma_0$.
Zbiór ten jest skończony i zawiera się w $\Gamma$.
To samo drzewo, z tymi samymi regułami i rozładowaniami,
jest dedukcją $\varphi$ z $\Gamma_0$.
\end{proof}

\begin{uwaga}[Czego jeszcze nie dowiedliśmy?]
Sam fakt, że zapisaliśmy reguły, nie rozstrzyga,
czy zawsze prowadzą one od zdań prawdziwych do prawdziwych.
Tę własność — \emph{poprawność} rachunku dowodowego —
udowodnimy w następnym rozdziale.
Nie udowodniliśmy także, że z każdej konsekwencji
semantycznej można otrzymać formalny dowód.
Taką tezę nazywamy \emph{pełnością} i będzie ona przedmiotem
osobnego, znacznie trudniejszego argumentu.
\end{uwaga}
\chapter[Poprawność rachunku pierwszego rzędu]{Poprawność rachunku dowodowego i konsekwencja logiczna}
\label{ch:soundness-fol}

\section{Dwa rodzaje wynikania}
Zapis $\Gamma\vdash_{\mathrm{ND}}\varphi$ z
rozdziału~\ref{ch:natural-deduction} oznacza istnienie
skończonego, formalnego dowodu. Zapis $\Gamma\models\varphi$
z rozdziału~\ref{ch:semantics-fol} mówi natomiast o
prawdziwości formuły we wszystkich interpretacjach spełniających założenia.
Są to pojęcia zbudowane odmiennymi metodami.
Pierwsze dotyczy skończonych napisów i reguł, drugie
— wszystkich struktur i ich wartościowań.

\begin{definicja}[Konsekwencja semantyczna dla formuł]
\label{def:soundness-entailment-formulas}
Niech $\mathcal L$ będzie językiem pierwszego rzędu,
$\Gamma$ zbiorem formuł języka $\mathcal L$,
a $\varphi$ formułą tego samego języka.
Piszemy $\Gamma\models\varphi$, jeżeli dla każdej
struktury $\mathfrak A$ języka $\mathcal L$
o niepustym uniwersum $A$ i każdego wartościowania $s$ zmiennych języka
$\mathcal L$ w zbiorze $A$ zachodzi implikacja
\[
\bigl(\forall\gamma\in\Gamma\ \mathfrak A,s\models\gamma\bigr)
\quad\Longrightarrow\quad
\mathfrak A,s\models\varphi.
\]
Gdy $\Gamma=\varnothing$, piszemy $\models\varphi$.
\end{definicja}
Jeżeli $\Gamma$ składa się ze zdań, a $\varphi$ jest zdaniem,
to dzięki Wnioskowi~\ref{cor:sem-sentence}
powyższa definicja zgadza się z
Definicją~\ref{def:sem-entailment}.
Nie można jednak pominąć wartościowania, gdy w formułach występują
zmienne wolne. Przykładowo $P(x)$ nie wynika semantycznie
z samego istnienia elementu spełniającego $P$.

\begin{przyklad}[Różnica między pojedynczym modelem a wszystkimi modelami]
W strukturze $(\N,<)$ zdanie $\forall x\,\exists y\,(x<y)$
jest prawdziwe. Nie oznacza to, że jest prawdziwe
w każdej strukturze z binarnym symbolem relacyjnym $<$.
Jeśli za uniwersum przyjmiemy $\{a\}$, a relację $<$
za pustą, to zdanie jest fałszywe. Dla wynikania
semantycznego potrzebna jest kwantyfikacja po \emph{wszystkich}
strukturach rozważanego języka.
\end{przyklad}


\begin{przyklad}[Semantyczna wersja prostego kroku dowodowego]
Sprawdzając semantyczną poprawność reguły, nie twierdzimy,
że jej przesłanki są prawdziwe w każdej strukturze.
Pytamy jedynie: \emph{jeżeli} w ustalonej strukturze
przesłanki są spełnione, czy musi być spełniony wniosek?
To pozwala zrozumieć, dlaczego nad kreską zapisujemy
kilka wymaganych przesłanek, nie utożsamiając ich
z dodatkowymi aksjomatami całej teorii.

Niech $P$ i $Q$ oznaczają dowolne formuły tego samego języka.
Mamy $\{P\land Q\}\models P$: w dowolnej strukturze i przy
dowolnym wartościowaniu, w którym $P\land Q$ jest prawdziwa,
prawdziwe jest także $P$. To jest fakt o \emph{znaczeniu}
spójnika $\land$. W rozdziale~\ref{ch:natural-deduction}
ten sam krok zapisaliśmy formalną regułą $(\land E_1)$.
Twierdzenie o poprawności pokaże, że tak wygląda każdy
legalny krok systemu dowodowego: nie może zamienić
prawdziwych przesłanek na fałszywy wniosek.
\end{przyklad}

\section{Trzy pomocnicze fakty semantyczne}
Dwa pierwsze fakty udowodniliśmy już w rozdziale 10,
ale przytoczymy je w dokładnej postaci potrzebnej poniżej.

\begin{lemat}[Zmiana jednej nieistotnej zmiennej]
\label{lem:soundness-fresh}
Niech $\mathcal L$ będzie językiem pierwszego rzędu,
$\mathfrak A$ dowolną strukturą tego języka o uniwersum $A$,
$\varphi$ formułą, a $y$ zmienną niewystępującą wolno
w $\varphi$. Dla każdego wartościowania $s$ i $a\in A$ zachodzi
\[
\mathfrak A,s\models\varphi
\quad\Longleftrightarrow\quad
\mathfrak A,s[y\mapsto a]\models\varphi.
\]
\end{lemat}
\begin{proof}
Wartościowania $s$ i $s[y\mapsto a]$ różnią się
co najwyżej na zmiennej $y$. Ponieważ $y$ nie jest
wolną zmienną $\varphi$, zgadzają się na
$\operatorname{FV}(\varphi)$. Teza wynika z
Twierdzenia~\ref{thm:sem-coincidence}.
\end{proof}

\begin{lemat}[Interpretacja podstawienia]\label{lem:soundness-subst}
Niech $\mathcal L$ będzie językiem pierwszego rzędu,
$\mathfrak A$ jego strukturą o uniwersum $A$,
$\varphi$ formułą, $x$ zmienną, $t$ termem dopuszczalnym
do podstawienia za $x$ w $\varphi$, a $s$ wartościowaniem.
Wówczas
\[
\mathfrak A,s\models\varphi[x:=t]
\quad\Longleftrightarrow\quad
\mathfrak A,s[x\mapsto t^{\mathfrak A}[s]]\models\varphi.
\]
\end{lemat}
\begin{proof}
To Twierdzenie~\ref{thm:sem-substitution},
którego dowód indukcyjny obejmował wszystkie rodzaje formuł,
w szczególności przypadki obu kwantyfikatorów.
\end{proof}

\begin{lemat}[Zastępowanie równych termów]
\label{lem:soundness-equality}
Niech $\mathcal L$ będzie językiem pierwszego rzędu
z równością, $\mathfrak A$ strukturą języka $\mathcal L$
o niepustym uniwersum $A$, $s$ wartościowaniem,
$x$ zmienną, $t,u$ termami języka $\mathcal L$,
a $\theta$ formułą tego języka.
Załóżmy, że $t$ i $u$ są dopuszczalne do podstawienia
za zmienną $x$ w $\theta$.
Jeżeli zachodzą $\mathfrak A,s\models t=u$ oraz
$\mathfrak A,s\models\theta[x:=t]$, to
$\mathfrak A,s\models\theta[x:=u]$.
\end{lemat}
\begin{proof}
Z prawdziwości atomu $t=u$ otrzymujemy
$t^{\mathfrak A}[s]=u^{\mathfrak A}[s]=a$
dla pewnego $a\in A$.
Lemat~\ref{lem:soundness-subst} przekształca przesłankę
$\mathfrak A,s\models\theta[x:=t]$ w
$\mathfrak A,s[x\mapsto a]\models\theta$.
Ponieważ $u^{\mathfrak A}[s]=a$, ponowne zastosowanie
tego lematu daje
$\mathfrak A,s\models\theta[x:=u]$.
\end{proof}

\section{Twierdzenie o poprawności}
\begin{twierdzenie}[Poprawność dedukcji naturalnej]
\label{thm:soundness}
Niech $\mathcal L$ będzie dowolnym językiem pierwszego rzędu
z równością, $\Gamma$ zbiorem formuł tego języka,
a $\varphi$ formułą. Jeśli
\[
\Gamma\vdash_{\mathrm{ND}}\varphi,
\]
to
\[
\Gamma\models\varphi.
\]
\end{twierdzenie}
\begin{proof}
Ustalmy dowolną strukturę $\mathfrak A$ języka $\mathcal L$
i dowolne wartościowanie $s$ spełniające wszystkie
formuły z $\Gamma$. Pokażemy, że spełniona jest także $\varphi$.
Dowód prowadzimy indukcją po strukturze skończonego drzewa dedukcji.
Hipoteza indukcyjna mówi dokładniej, że każda poddedukcja
jest semantycznie poprawna względem swoich \emph{otwartych}
założeń. Rozładowane założenia rozpatrujemy tylko na czas
odpowiedniego kroku.

\smallskip
\emph{1. Liście i reguły bez przesłanek.}
Jeśli liść jest otwartym założeniem $\gamma$
danej poddedukcji, to każde wartościowanie
spełniające wszystkie jej otwarte założenia spełnia też
$\gamma$. Dotyczy to również założeń tymczasowych,
które będą rozładowane w dalszych krokach.
Dla całego dowodu jego otwarte założenia należą do $\Gamma$.
Reguła $=I$ kończąca się na $t=t$ jest poprawna,
ponieważ $t^{\mathfrak A}[s]=t^{\mathfrak A}[s]$.

\smallskip
\emph{2. Koniunkcja i implikacja.}
Jeżeli obie przesłanki reguły $\land I$ są spełnione,
to prawdziwa jest ich koniunkcja; jeśli spełniona jest
koniunkcja, każdy jej człon jest spełniony, co uzasadnia
$\land E_1$ i $\land E_2$.
Dla $\to E$ spełnianie $\alpha\to\beta$ oraz $\alpha$
wymusza spełnianie $\beta$ zgodnie z tablicą implikacji.
Dla $\to I$ mamy poddedukcję $\beta$ z dodatkowym,
tymczasowym założeniem $\alpha$. Rozważmy dowolną
strukturę i wartościowanie spełniające pozostałe otwarte
założenia. Jeśli $\alpha$ jest fałszywa, implikacja
$\alpha\to\beta$ jest prawdziwa niezależnie od $\beta$.
Jeśli $\alpha$ jest prawdziwa, hipoteza indukcyjna
zastosowana do poddedukcji daje $\beta$.
W obu przypadkach implikacja jest prawdziwa.

\smallskip
\emph{3. Alternatywa, równoważność i negacja.}
Reguły $\lor I_1$ i $\lor I_2$ są poprawne, bo
spełnianie jednego członu wystarcza dla alternatywy.
W $\lor E$ prawdziwa alternatywa oznacza prawdziwość
co najmniej jednego członu. W pierwszym przypadku
hipoteza indukcyjna dla pierwszej poddedukcji daje $\theta$;
w drugim przypadku daje ją druga poddedukcja.
Reguły równoważności są poprawne, ponieważ
$\alpha\leftrightarrow\beta$ jest prawdziwa dokładnie wtedy,
gdy prawdziwe są obie implikacje wzajemne.

Jeśli prawdziwe byłyby zarazem $\alpha$ i $\neg\alpha$,
otrzymalibyśmy sprzeczność z definicją negacji.
Dlatego reguła $\neg E$ może prowadzić do $\bot$
jedynie z przesłanek, których nie można jednocześnie spełnić.
Jeśli z dodatkowego założenia $\alpha$ otrzymano $\bot$,
to przy spełnieniu pozostałych założeń $\alpha$ nie może być
prawdziwa: inaczej otrzymalibyśmy spełnianie $\bot$.
Zatem $\neg\alpha$ jest prawdziwa, co uzasadnia $\neg I$.
Reguła $\bot E$ jest poprawna, ponieważ $\bot$ nie jest
spełniona przy żadnym wartościowaniu; zatem nie istnieje
wartościowanie, dla którego jej przesłanka byłaby prawdziwa,
a wniosek fałszywy. Wreszcie reguła klasycznego dowodu
nie wprost $\mathrm{RAA}$ jest poprawna: jeśli założenie
$\neg\alpha$ prowadziłoby do sprzeczności, nie może być
prawdziwe, a klasyczna semantyka dwuwartościowa
zapewnia prawdziwość $\alpha$.

\smallskip
\emph{4. Eliminacja kwantyfikatora ogólnego.}
Przypuśćmy, że $\mathfrak A,s\models\forall x\,\alpha$.
Niech $a=t^{\mathfrak A}[s]$. Definicja spełniania
kwantyfikatora daje
$\mathfrak A,s[x\mapsto a]\models\alpha$.
Lemat~\ref{lem:soundness-subst} zamienia ten warunek
na $\mathfrak A,s\models\alpha[x:=t]$.
To dokładnie wniosek reguły $\forall E$.

\smallskip
\emph{5. Wprowadzanie kwantyfikatora ogólnego.}
Niech poddedukcja z otwartymi założeniami $\Delta$
kończy się na $\alpha$, a zmienna $x$ nie występuje
wolno w żadnej formule z $\Delta$.
Załóżmy, że $\mathfrak A,s$ spełnia $\Delta$.
Dla dowolnego $a\in A$ Lemat~\ref{lem:soundness-fresh}
zastosowany oddzielnie do każdego elementu $\Delta$
pokazuje, że $\mathfrak A,s[x\mapsto a]$ również spełnia
wszystkie otwarte założenia. Hipoteza indukcyjna daje więc
$\mathfrak A,s[x\mapsto a]\models\alpha$.
Ponieważ $a$ było dowolne, definicja $\forall$ daje
$\mathfrak A,s\models\forall x\,\alpha$.
W tym właśnie miejscu konieczny jest warunek o niewystępowaniu
$x$ w otwartych założeniach.

\smallskip
\emph{6. Wprowadzanie kwantyfikatora egzystencjalnego.}
Jeśli $\mathfrak A,s\models\alpha[x:=t]$,
to dla $a=t^{\mathfrak A}[s]$ Lemat~\ref{lem:soundness-subst}
daje $\mathfrak A,s[x\mapsto a]\models\alpha$.
Wartość $a$ jest świadkiem spełniania $\exists x\,\alpha$.

\smallskip
\emph{7. Eliminacja kwantyfikatora egzystencjalnego.}
Załóżmy, że ostatnia reguła ma przesłanki
$\exists x\,\alpha$ oraz poddedukcję $\theta$ z dodatkowego
założenia $\alpha[x:=y]$, gdzie $y$ jest świeżą zmienną
nieobecną w $\alpha$ i we wniosku $\theta$ oraz niewystępującą
wolno w żadnym otwartym założeniu obu gałęzi poza
rozładowywanym $\alpha[x:=y]$. Przez $\Delta$ oznaczmy
sumę tych pozostałych otwartych założeń; w szczególności
obejmuje ona założenia potrzebne do wyprowadzenia
przesłanki głównej $\exists x\,\alpha$.
Niech $\mathfrak A,s$ spełnia $\Delta$.
Hipoteza indukcyjna zastosowana do dowodu przesłanki głównej
zapewnia $\mathfrak A,s\models\exists x\,\alpha$.
Z tego wynika istnienie
$a\in A$ takiego, że
$\mathfrak A,s[x\mapsto a]\models\alpha$.
Połóżmy $r=s[y\mapsto a]$.
Ponieważ $y$ nie występuje wolno w $\alpha$,
wartościowania $r[x\mapsto a]$ i $s[x\mapsto a]$
zgadzają się na wszystkich jej zmiennych wolnych.
Stąd na mocy Twierdzenia~\ref{thm:sem-coincidence}
mamy $\mathfrak A,r[x\mapsto a]\models\alpha$.
Z Lematu~\ref{lem:soundness-subst} dla termu $y$ wynika
$\mathfrak A,r\models\alpha[x:=y]$.
Świeżość $y$ oraz Lemat~\ref{lem:soundness-fresh}
zapewniają też, że $\mathfrak A,r$ spełnia $\Delta$.
Hipoteza indukcyjna dla poddedukcji daje zatem
$\mathfrak A,r\models\theta$. Zmienna $y$ nie występuje
wolno w $\theta$, więc ponowne użycie lematu o świeżej
zmiennej daje $\mathfrak A,s\models\theta$.
Reguła $\exists E$ jest więc poprawna.

\smallskip
\emph{8. Eliminacja równości.}
Jest to dokładnie treść Lematu~\ref{lem:soundness-equality}.

Wyczerpaliśmy wszystkie reguły z
Definicji~\ref{def:nd-rules-conj-imp},
Definicji~\ref{def:nd-rules-prop},
Definicji~\ref{def:nd-rules-quant} i
Definicji~\ref{def:nd-equality}.
Indukcja po drzewie dowodu jest zakończona.
Dla korzenia drzewa otwarte założenia zawierają się w $\Gamma$,
więc przyjęte na początku $\mathfrak A,s$ spełnia także wniosek
$\varphi$. Ponieważ $\mathfrak A$ i $s$ były dowolne,
wynika $\Gamma\models\varphi$.
\end{proof}


\section{Jak poprawność działa przy pominiętych krokach?}
Dla czytelnika przyzwyczajonego do zapisu $\vdots$
warto oddzielić trzy poziomy rozumowania.

\begin{przyklad}[Pełne uzasadnienie krótkiego drzewa]
Rozważmy udowodnioną w rozdziale 11 implikację
$(P\land Q)\to(Q\land P)$. Jeżeli w dowolnej strukturze
przy danym wartościowaniu prawdziwa jest formuła $P\land Q$,
to na mocy definicji koniunkcji prawdziwe są osobno $P$ i $Q$.
Wobec tego prawdziwa jest $Q\land P$. Gdy $P\land Q$
nie jest prawdziwa, implikacja jest prawdziwa ze względu na
swoją tabelę prawdy. Oznacza to, że cała implikacja jest
prawdziwa przy każdym wartościowaniu i w każdej strukturze.

W drzewie formalnym pomiędzy $[P\land Q]^1$ i $Q\land P$
można było użyć skrótu $\vdots$, ale jego poprawność
weryfikujemy przez sprawdzenie trzech konkretnych zastosowań
reguł: $(\land E_2)$, $(\land E_1)$ i $(\land I)$.
Twierdzenie~\ref{thm:soundness} mówi, że taka weryfikacja
może być przeprowadzona dla \emph{każdego} skończonego drzewa.
\end{przyklad}

\begin{uwaga}[Poprawność nie mówi, że założenia są prawdziwe]
Jeżeli $\Gamma\vdash_{\mathrm{ND}}\varphi$, poprawność orzeka:
\emph{w każdej interpretacji, w której spełnione są wszystkie
formuły z $\Gamma$, spełniona jest także $\varphi$}.
Jeżeli $\Gamma$ ma jakieś fałszywe założenie w danej
interpretacji, twierdzenie nie daje w tej interpretacji
informacji o wartości $\varphi$. Otwieranie i rozładowywanie
założeń regułą $(\to I)$ ma więc sens \emph{warunkowy},
nie zmienia wartości logicznych formuł.
\end{uwaga}

\section{Konsekwencje poprawności}
\begin{wniosek}[Brak formalnego dowodu zdania fałszywego w pewnym modelu]
\label{cor:soundness-countermodel}
Niech $\mathcal L$ będzie językiem pierwszego rzędu,
$\Gamma$ zbiorem jego formuł, a $\varphi$ formułą
tego języka. Jeśli istnieją struktura $\mathfrak A$
języka $\mathcal L$ i wartościowanie $s$
spełniające wszystkie formuły z $\Gamma$, ale niespełniające
$\varphi$, to $\Gamma\nvdash_{\mathrm{ND}}\varphi$.
\end{wniosek}
\begin{proof}
Gdyby $\Gamma\vdash_{\mathrm{ND}}\varphi$,
Twierdzenie~\ref{thm:soundness} dawałoby
$\Gamma\models\varphi$. To przeczyłoby istnieniu
opisanego kontrmodelu.
\end{proof}
\begin{przyklad}[Dlaczego $\forall x\,P(x)$ nie wynika z $P(y)$]
Niech uniwersum struktury wynosi $\{0,1\}$,
a $P^{\mathfrak A}=\{0\}$. Przy wartościowaniu
$s(y)=0$ spełniona jest formuła $P(y)$,
ale nie jest spełnione $\forall x\,P(x)$.
Z Wniosku~\ref{cor:soundness-countermodel} otrzymujemy
$\{P(y)\}\nvdash_{\mathrm{ND}}\forall x\,P(x)$.
To samo można było zauważyć bezpośrednio na poziomie
reguł: warunek dla $\forall I$ jest naruszony, jeśli
uogólniamy zmienną wolną w nieusuniętym założeniu.
\end{przyklad}

\begin{wniosek}[Model wyklucza dowód sprzeczności]
\label{cor:soundness-consistency}
Niech $\mathcal L$ będzie językiem pierwszego rzędu,
a $\Gamma$ zbiorem formuł tego języka.
Jeśli istnieją struktura $\mathfrak A$ języka $\mathcal L$
i wartościowanie $s$ spełniające wszystkie
formuły z $\Gamma$, to
$\Gamma\nvdash_{\mathrm{ND}}\bot$.
\end{wniosek}
\begin{proof}
Formuła $\bot$ jest fałszywa w każdej strukturze.
Model $\Gamma$ jest więc kontrmodelem dla
$\Gamma\models\bot$. Z poprzedniego wniosku otrzymujemy
$\Gamma\nvdash_{\mathrm{ND}}\bot$.
\end{proof}

\begin{uwaga}[Ostrożnie z odwracaniem implikacji]
Udowodniliśmy: istnienie modelu daje brak formalnego dowodu
sprzeczności. Odwrotność — że każdy niesprzeczny zbiór zdań
ma model — nie wynika z samej poprawności.
Wymaga twierdzenia o pełności, a dla nieskończonych zbiorów zdań
stanowi jego ważną postać.
\end{uwaga}

\section{Twierdzenia o dedukcji: syntaktyczne i semantyczne}
\begin{twierdzenie}[Semantyczne twierdzenie o dedukcji]
\label{thm:soundness-semantic-deduction}
Niech $\mathcal L$ będzie językiem pierwszego rzędu,
$\Gamma$ zbiorem formuł tego języka, a $\alpha,\beta$
formułami języka $\mathcal L$. Wtedy
\[
\Gamma\cup\{\alpha\}\models\beta
\quad\Longleftrightarrow\quad
\Gamma\models\alpha\to\beta.
\]
\end{twierdzenie}
\begin{proof}
Załóżmy, że $\Gamma\cup\{\alpha\}\models\beta$.
Ustalmy strukturę $\mathfrak A$ i wartościowanie $s$
spełniające $\Gamma$. Jeśli nie spełniają $\alpha$,
implikacja $\alpha\to\beta$ jest prawdziwa.
Jeśli spełniają $\alpha$, to spełniają wszystkie założenia
z $\Gamma\cup\{\alpha\}$, a więc także $\beta$.
Implikacja jest znów prawdziwa. Zatem
$\Gamma\models\alpha\to\beta$.

W drugą stronę załóżmy
$\Gamma\models\alpha\to\beta$.
Niech $\mathfrak A,s$ spełnia wszystkie formuły z
$\Gamma\cup\{\alpha\}$. Spełnia więc $\Gamma$
i $\alpha$. Z założenia wynika spełnianie implikacji
$\alpha\to\beta$. Ponieważ jej poprzednik jest prawdziwy,
na mocy tablicy implikacji prawdziwy jest następnik $\beta$.
To dowodzi $\Gamma\cup\{\alpha\}\models\beta$.
\end{proof}

Zestawmy to z Twierdzeniem~\ref{thm:nd-deduction},
które dotyczyło dowodliwości syntaktycznej. Obie równoważności
mają podobną postać, ale zostały udowodnione na podstawie
zupełnie różnych definicji: jedna za pomocą reguł drzewa dowodu,
druga przez rozważenie dowolnej struktury i wartościowania.

\section{Skończoność formalnych dowodów a zwartość semantyczna}
\begin{twierdzenie}[Syntaktyczna zwartość wyprowadzalności]
\label{thm:soundness-syntactic-finitary}
Niech $\mathcal L$ będzie językiem pierwszego rzędu,
$\Gamma$ zbiorem formuł tego języka, a $\varphi$
formułą języka $\mathcal L$.
Jeżeli $\Gamma\vdash_{\mathrm{ND}}\varphi$, to istnieje
skończony zbiór $\Gamma_0\subseteq\Gamma$ taki, że
$\Gamma_0\vdash_{\mathrm{ND}}\varphi$.
\end{twierdzenie}
\begin{proof}
Twierdzenie~\ref{thm:nd-finite-support} dowiodło tej tezy
bez odwoływania się do semantyki. Powodem jest skończoność
liczby liści drzewa formalnego dowodu.
\end{proof}

\begin{uwaga}[Czego \emph{nie} wolno z tego wywnioskować]
Z samej powyższej własności i poprawności nie możemy jeszcze
stwierdzić, że jeśli $\Gamma\models\varphi$, to istnieje
skończony $\Gamma_0\subseteq\Gamma$ z
$\Gamma_0\models\varphi$.
Do takiego przejścia trzeba mieć także twierdzenie,
że konsekwencja semantyczna jest zawsze dowodliwa
(pełność pierwszego rzędu).
Nie należy więc utożsamiać \emph{syntaktycznej skończoności
dowodów} z \emph{semantycznym twierdzeniem o zwartości}.
\end{uwaga}

\section{Co dalej?}
Dysponujemy teraz jednym kierunkiem pomostu między składnią
a semantyką:
\[
\Gamma\vdash_{\mathrm{ND}}\varphi
\quad\Longrightarrow\quad
\Gamma\models\varphi.
\]
Przed nami kierunek odwrotny. Jego dowód wymaga skonstruowania
modelu dla każdego odpowiednio niesprzecznego zbioru zdań.
W późniejszych rozdziałach wprowadzimy rozszerzenia teorii
o stałe świadczące o istnieniu, zbudujemy model termowy
i przeprowadzimy szczegółowy dowód twierdzenia Gödla o pełności.
Dopiero potem wyprowadzimy semantyczne twierdzenie o zwartości
i dalej rozwiniemy teorię modeli.
\chapter{Twierdzenie Gödla o pełności i konstrukcja Henkina}
\label{ch:completeness}

\section{O co naprawdę chodzi w pełności?}
W rozdziale~\ref{ch:soundness-fol} udowodniliśmy, że każdy formalny dowód prowadzi do konsekwencji semantycznej. Pozostało pytanie odwrotne: czy każda formuła, której nie sposób obalić żadnym modelem spełniającym przesłanki, ma formalny dowód? Okazuje się, że w klasycznej logice pierwszego rzędu odpowiedź jest twierdząca. Jest to \emph{twierdzenie Gödla o pełności}, nie należy go jednak mylić z późniejszymi twierdzeniami Gödla \emph{o niezupełności} teorii arytmetycznych. Pierwsze dotyczy mocy reguł czystej logiki, drugie ograniczeń efektywnie aksjomatyzowanych teorii o odpowiedniej sile.

\begin{definicja}[Niesprzeczność syntaktyczna]\label{def:henkin-consistency}
Niech $\mathcal L$ będzie językiem pierwszego rzędu z równością, a $T$ zbiorem zdań języka $\mathcal L$. Jak w poprzednich rozdziałach, wszystkie rozważane struktury mają niepuste uniwersa. Mówimy, że $T$ jest \emph{syntaktycznie niesprzeczny}, jeżeli $T\nvdash_{\mathrm{ND}}\bot$. W przeciwnym przypadku nazywamy $T$ sprzecznym. Symbol $\bot$ jest skrótem ustalonego zdania sprzecznego zdefiniowanym w rozdziale~\ref{ch:natural-deduction}.
\end{definicja}

\begin{uwaga}[Dlaczego zaczynamy od zbiorów zdań?]
Dla zdań prawda w strukturze nie zależy od wartościowania. W dowodzie konstrukcyjnym będziemy dodawać nowe symbole stałych, a następnie tworzyć \emph{model termowy}. Pod koniec rozdziału wyjaśnimy starannie, jak wynik dla zdań przenieść na dowolne formuły ze zmiennymi wolnymi.
\end{uwaga}

\begin{twierdzenie}[Twierdzenie o istnieniu modelu]\label{thm:henkin-model-existence}
Niech $\mathcal L$ będzie dowolnym językiem pierwszego rzędu z równością, a $T$ zbiorem zdań języka $\mathcal L$. Jeśli $T$ jest syntaktycznie niesprzeczny, to istnieje struktura $\mathfrak M$ o niepustym uniwersum taka, że $\mathfrak M\models T$.
\end{twierdzenie}

W kolejnych sekcjach skonstruujemy taką strukturę. To będzie główny wynik, z którego wyprowadzimy pełność. Pracujemy w zwykłej metateorii zbiorów z aksjomatem wyboru (ZFC); dla języków przeliczalnych wskażemy także elementarną konstrukcję niewymagającą lematu Kuratowskiego--Zorna.

\section{Cztery fakty o dowodach, które będą potrzebne}
\begin{lemat}[Skończone wsparcie sprzeczności]\label{lem:henkin-finite-contradiction}
Jeśli $T\vdash_{\mathrm{ND}}\bot$, to istnieje skończony zbiór $T_0\subseteq T$ taki, że $T_0\vdash_{\mathrm{ND}}\bot$.
\end{lemat}
\begin{proof}
Jest to szczególny przypadek Twierdzenia~\ref{thm:nd-finite-support}. Drzewo dowodu ma skończenie wiele liści, zatem tylko skończenie wiele różnych zdań z $T$ może wystąpić w nim jako nierozładowane założenia. Te same reguły tworzą dowód sprzeczności już z tych zdań.
\end{proof}

\begin{lemat}[Dodanie zdania i negacji]\label{lem:henkin-extension-neg}
Niech $S$ będzie niesprzecznym zbiorem zdań, a $\sigma$ zdaniem tego samego języka. Jeśli $S\cup\{\sigma\}$ jest sprzeczny, to $S\vdash_{\mathrm{ND}}\neg\sigma$. Jeżeli ponadto $S\cup\{\neg\sigma\}$ jest sprzeczny, to również $S$ jest sprzeczny.
\end{lemat}
\begin{proof}
Pierwsza część: z $S$ oraz dodatkowego otwartego założenia $\sigma$ istnieje dowód $\bot$. Reguła $\neg I$ rozładowuje $\sigma$ i daje $S\vdash\neg\sigma$. Dla drugiej części analogicznie otrzymujemy $S\vdash\neg\neg\sigma$. Z $\neg\neg\sigma$ przez klasyczny dowód nie wprost uzyskujemy $\sigma$, a zestawiając $\sigma$ i $\neg\sigma$ otrzymujemy $\bot$. W ten sposób z samego $S$ powstał dowód sprzeczności.
\end{proof}

\begin{lemat}[Negacja kwantyfikatora]\label{lem:henkin-quant-duality}
Dla każdej formuły $\varphi(x)$, w której $x$ jest wyróżnioną zmienną, w dedukcji naturalnej z rozdziału~\ref{ch:natural-deduction} dają się bez założeń wyprowadzić oba kierunki równoważności
\[
\neg\forall x\,\varphi(x)\quad\longleftrightarrow\quad\exists x\,\neg\varphi(x).
\]
Wszystkie użyte podstawienia wymagają zwykłych warunków unikania przechwycenia zmiennych.
\end{lemat}
\begin{proof}
Załóżmy $\exists x\,\neg\varphi(x)$ i tymczasowo $\forall x\,\varphi(x)$. Wybieramy zmienną $y$ niewystępującą w ogóle w $\varphi$ ani w pozostałych otwartych założeniach. W poddedukcji reguły $\exists E$ przyjmujemy $\neg\varphi[x:=y]$. Przez $\forall E$ z drugiego założenia otrzymujemy $\varphi[x:=y]$, a następnie przez $\neg E$ dostajemy $\bot$. Ponieważ $y$ nie występuje we wniosku $\bot$ ani w pozostałych otwartych założeniach, możemy zastosować $\exists E$ i rozładować założenie o świadku. Dopiero teraz reguła $\neg I$ rozładowuje $\forall x\,\varphi$ i daje $\neg\forall x\,\varphi$ z jedynego pozostałego założenia $\exists x\,\neg\varphi$.

W przeciwnym kierunku załóżmy $\neg\forall x\,\varphi(x)$ i chwilowo $\neg\exists x\,\neg\varphi(x)$. Wybieramy $y$ niewystępującą w ogóle w napisie $\varphi$ ani wolno w żadnym z tych dwóch założeń. Z tymczasowego $\neg\varphi[x:=y]$ reguła $\exists I$ daje $\exists x\,\neg\varphi(x)$, a z przyjętą negacją prowadzi do $\bot$. Reguła $\mathrm{RAA}$ rozładowuje założenie $\neg\varphi[x:=y]$ i daje $\varphi[x:=y]$. Zmienna $y$ nie występuje wolno w żadnym z pozostałych otwartych założeń, więc $\forall I$ daje $\forall y\,\varphi[x:=y]$. Jest to dopuszczalne przemianowanie $\forall x\,\varphi(x)$: stosując $\forall E$ z termem $x$, a potem $\forall I$ dla $x$ (które nie jest wolne w otwartych założeniach), uzyskujemy $\forall x\,\varphi(x)$. Sprzeczność z pierwszym założeniem pozwala rozładować $\neg\exists x\,\neg\varphi(x)$ regułą $\mathrm{RAA}$, co daje $\exists x\,\neg\varphi(x)$.

Rozładowując po kolei pozostałe założenia regułą $\to I$, otrzymujemy oba kierunki żądanej równoważności.
\end{proof}

\begin{lemat}[Zastąpienie nowej stałej świeżą zmienną]\label{lem:henkin-constant-to-variable}
Niech $c$ będzie nowym symbolem stałej, niewystępującym w języku zdań zbioru $S$. Załóżmy, że w języku rozszerzonym istnieje skończona dedukcja z $S$ zdania $\theta(c)$. Wówczas dla zmiennej $y$, wybranej jako nieobecna w całej tej dedukcji, można skonstruować dedukcję z $S$ formuły $\theta(y)$ otrzymanej przez zamianę wystąpień $c$ na $y$.
\end{lemat}
\begin{proof}
Wybieramy $y$ niewystępującą nigdzie w skończonym drzewie dowodu, także jako zmienna związana i jako parametr którejś reguły kwantyfikatora. Jest to możliwe, bo zmiennych mamy nieskończenie wiele. W każdym termie i formule drzewa zamieniamy $c$ na $y$. Ponieważ żaden kwantyfikator nie wiąże $y$, nie dochodzi do przechwycenia zmiennej. Kroki dla spójników i równości pozostają przykładami tych samych reguł. W przypadku $\forall I$ uogólniana zmienna nie jest $y$, a zamiana $c$ na $y$ nie dodaje jej wystąpień do otwartych założeń, więc warunek uboczny pozostaje spełniony. W przypadku $\exists E$ zmienna świeża użyta w oryginalnym kroku różni się od $y$, więc również nie traci świeżości po tej zamianie. To uzasadnia indukcyjnie każdy węzeł nowego drzewa. Ponieważ żadne zdanie z $S$ nie zawierało $c$, wszystkie nierozładowane założenia nadal są zdaniami z $S$.
\end{proof}

\section{Świadkowie Henkina: od istnienia do nazwy}
Zdanie $\exists x\,\varphi(x)$ mówi, że istnieje pewien element, ale nie musi podawać termu, który go nazywa. Chcemy dodać symbol stałej $c_{\varphi}$ wraz z warunkiem: „jeśli taki element istnieje, to $c_{\varphi}$ jest jednym z nich”. To nie oznacza, że sami wybieramy element już istniejącego modelu --- model dopiero zbudujemy.

\begin{definicja}[Aksjomat świadka]\label{def:henkin-witness-axiom}
Niech $\exists x\,\varphi(x)$ będzie \emph{zdaniem} w danym języku, przy czym jedyną ewentualną wolną zmienną w $\varphi$ jest $x$. Dodajemy nową stałą $c_{\varphi}$ i zdanie
\[
H_{\varphi}:\qquad (\exists x\,\varphi(x))\to\varphi(c_{\varphi}).
\]
Nazywamy je \emph{aksjomatem świadka Henkina}. Różnym zdaniom egzystencjalnym przypisujemy różne nowe symbole stałych.
\end{definicja}

\begin{lemat}[Dodanie jednego świadka zachowuje niesprzeczność]\label{lem:henkin-one-witness}
Jeśli $S$ jest niesprzecznym zbiorem zdań języka $\mathcal L$, $\exists x\,\varphi(x)$ zdaniem tego języka i $c$ nowym symbolem stałej, to $S\cup\{(\exists x\,\varphi(x))\to\varphi(c)\}$ jest niesprzeczny w języku $\mathcal L\cup\{c\}$.
\end{lemat}
\begin{proof}
Niech $A:=\exists x\,\varphi(x)$ i $B:=\varphi[x:=c]$. Załóżmy przeciwnie, że z $S\cup\{A\to B\}$ można wyprowadzić $\bot$ w języku rozszerzonym. Reguła $\neg I$ daje wtedy dowód $S\vdash_{\mathrm{ND}}\neg(A\to B)$ w tym języku. Wybierzmy zmienną $y$ nieobecną w całym otrzymanym drzewie oraz w napisie $\varphi$. Lemat~\ref{lem:henkin-constant-to-variable} daje w języku wyjściowym $\mathcal L$ dowód
\[
S\vdash_{\mathrm{ND}}\neg\bigl(A\to\varphi[x:=y]\bigr).
\]
Pokażemy już wyłącznie regułami dedukcji naturalnej, że z $S$ wynika sprzeczność.

Najpierw przyjmijmy tymczasowo $\neg A$. Z dodatkowego założenia $A$ otrzymujemy $\bot$ przez $\neg E$, stąd $\varphi[x:=y]$ przez $\bot E$; rozładowując $A$ przez $\to I$, dostajemy $A\to\varphi[x:=y]$. Przeczy to wyprowadzonej negacji implikacji, więc $\mathrm{RAA}$ rozładowuje $\neg A$ i daje $S\vdash_{\mathrm{ND}}A$.

Następnie przyjmijmy tymczasowo $\varphi[x:=y]$. Z pomocniczego założenia $A$ przez $\to I$ otrzymujemy $A\to\varphi[x:=y]$, znowu sprzeczne z jego negacją. Reguła $\neg I$ rozładowuje $\varphi[x:=y]$ i daje $S\vdash_{\mathrm{ND}}\neg\varphi[x:=y]$. Ponieważ $y$ nie występuje wolno w żadnym zdaniu z $S$, reguła $\forall I$ daje $S\vdash_{\mathrm{ND}}\forall y\,\neg\varphi[x:=y]$.

Wybieramy nową zmienną $z$ nieobecną w $\varphi$ ani w otwartych założeniach. W poddedukcji z tymczasowego $\varphi[x:=z]$ eliminacja $\forall E$ daje $\neg\varphi[x:=z]$, więc $\neg E$ daje $\bot$. Wniosek i pozostałe otwarte założenia są niezależne od $z$, dlatego $\exists E$, zastosowane do $A$, rozładowuje założenie świadka i daje $S\vdash_{\mathrm{ND}}\bot$ w pierwotnym języku $\mathcal L$. Jest to sprzeczne z niesprzecznością $S$.
\end{proof}

\begin{lemat}[Jednoczesne dodanie świadków dla jednego języka]\label{lem:henkin-all-witnesses}
Niech $S$ będzie niesprzecznym zbiorem zdań języka $\mathcal L$. Dla \emph{każdego} egzystencjalnego zdania $\exists x\,\varphi(x)$ tego języka wprowadźmy różną stałą $c_{\varphi}$ spoza $\mathcal L$ i aksjomat $H_{\varphi}$ z Definicji~\ref{def:henkin-witness-axiom}. Wtedy zbiór $S$ powiększony o wszystkie te aksjomaty jest niesprzeczny.
\end{lemat}
\begin{proof}
Gdyby rozszerzony zbiór był sprzeczny, Lemat~\ref{lem:henkin-finite-contradiction} pozwoliłby wskazać skończoną listę użytych nowych aksjomatów $H_{\varphi_1},\ldots,H_{\varphi_m}$. Dodajmy je po jednym. W chwili dodawania $H_{\varphi_k}$ symbol $c_{\varphi_k}$ nie występuje ani w $S$, ani w poprzednich aksjomatach, ponieważ wszystkie formuły $\varphi_j$ były napisane w \emph{starym} języku i symbole $c_{\varphi_j}$ są parami różne. Lemat~\ref{lem:henkin-one-witness} mówi zatem, że każdy krok zachowuje niesprzeczność. Po $m$ krokach wciąż mamy zbiór niesprzeczny, co przeczy założeniu o istniejącym dowodzie $\bot$.
\end{proof}


\begin{przyklad}[Świadek nie jest wcześniej wybranym elementem]\label{ex:henkin-witness-47}
Rozważmy język z jednoargumentowym symbolem relacyjnym $P$ i teorię
$S=\{\exists x\,P(x)\}$. Dołączamy nową stałą $c$ i zdanie
$H=(\exists x\,P(x))\to P(c)$. W każdej strukturze spełniającej
$S\cup\{H\}$ zachodzi $P(c)$, ponieważ przesłanka implikacji jest
prawdziwa. Nie znaczy to jednak, że $S$ już przed rozszerzeniem
wskazywała konkretny element: w strukturze o uniwersum $\{a,b\}$
i $P^{\mathfrak A}=\{a,b\}$ możemy interpretować $c$ jako $a$ albo
jako $b$. Obie interpretacje spełniają $H$.
Jeżeli natomiast $P^{\mathfrak B}=\varnothing$, to $H$ jest prawdziwe
niezależnie od interpretacji $c$, gdyż jego poprzednik jest fałszywy.
W tym drugim przypadku sam aksjomat świadka nie wymusza istnienia
elementu o własności $P$.
\end{przyklad}

\section{Dlaczego potrzebujemy wielu etapów?}\label{sec:henkin-stages}
W aksjomatach dodanych na pierwszym etapie pojawiają się nowe stałe, a przy ich użyciu można zbudować \emph{nowe} zdania egzystencjalne. One też potrzebują świadków. Dlatego jednorazowe dołączenie nazw nie wystarczy.

\begin{definicja}[Łańcuch rozszerzeń Henkina]\label{def:henkin-chain}
Dla dowolnego języka $\mathcal L$ i niesprzecznego zbioru jego zdań $T$ wykonujemy następującą konstrukcję:
\begin{enumerate}[label=(\roman*)]
\item Wprowadzamy jeden nowy symbol stałej $d$ i kładziemy $\mathcal L_0=\mathcal L\cup\{d\}$, $T_0=T$.
\item Znając $\mathcal L_n$ i $T_n$, dla każdego zdania egzystencjalnego języka $\mathcal L_n$ dodajemy osobną świeżą stałą i odpowiadający jej aksjomat świadka. Otrzymujemy język $\mathcal L_{n+1}$ i zbiór zdań $T_{n+1}$.
\item Definiujemy $\mathcal L_{\infty}=\bigcup_{n\in\N}\mathcal L_n$ oraz $T_{\infty}=\bigcup_{n\in\N}T_n$.
\end{enumerate}
Każdy skończony napis języka $\mathcal L_{\infty}$ korzysta tylko ze skończenie wielu symboli, zatem należy do pewnego $\mathcal L_n$.
\end{definicja}

\begin{figure}[htbp]
\centering
\begin{tikzpicture}[
  stage/.style={draw,rounded corners=2pt,inner xsep=9pt,inner ysep=5pt},
  >=Stealth,node distance=13mm
]
\node[stage] (l0) {$({\mathcal L}_0,T_0)$};
\node[stage,right=27mm of l0] (l1) {$({\mathcal L}_1,T_1)$};
\node[stage,right=27mm of l1] (l2) {$({\mathcal L}_2,T_2)$};
\node[right=12mm of l2] (dots) {$\cdots$};
\draw[->] (l0) -- node[above=2mm,font=\scriptsize]{świadkowie} (l1);
\draw[->] (l1) -- node[above=2mm,font=\scriptsize]{nowe świadki} (l2);
\draw[->,dotted] (l2) -- (dots);
\node[below=11mm of l1,font=\small] {$\mathcal L_{\infty}=\bigcup_{n\in\N}\mathcal L_n,\qquad
T_{\infty}=\bigcup_{n\in\N}T_n$};
\end{tikzpicture}
\caption{Schemat kolejnych rozszerzeń języka i teorii; strzałka kropkowana oznacza przejście do sumy wszystkich etapów.}
\label{fig:henkin-chain}
\end{figure}

\begin{twierdzenie}[Niesprzeczność i własność świadków]\label{thm:henkin-chain-properties}
Dla powyższej konstrukcji zachodzą dwa fakty:
\begin{enumerate}[label=(\roman*)]
\item $T_{\infty}$ jest syntaktycznie niesprzeczny.
\item Dla każdego egzystencjalnego zdania $\exists x\,\varphi(x)$ języka $\mathcal L_{\infty}$ istnieje term domknięty $c$ taki, że $T_{\infty}$ zawiera zdanie $(\exists x\,\varphi(x))\to\varphi(c)$.
\end{enumerate}
\end{twierdzenie}
\begin{proof}
Zbiór $T_0$ jest niesprzeczny: samo dołączenie nieużywanej stałej $d$ nie wprowadza nowego dowodu sprzeczności, bo z takiego skończonego dowodu można byłoby zastąpić $d$ świeżą zmienną jak w Lemacie~\ref{lem:henkin-constant-to-variable}. Indukcyjnie, jeśli $T_n$ jest niesprzeczny, to Lemat~\ref{lem:henkin-all-witnesses} zapewnia niesprzeczność $T_{n+1}$. Przypuśćmy teraz, że z $T_{\infty}$ da się wyprowadzić $\bot$. Dowód jest skończony, więc używa tylko skończenie wielu zdań $T_{\infty}$; wszystkie te zdania należą do jednego $T_n$, co dawałoby sprzeczność $T_n$. To dowodzi (i).

Niech $\exists x\,\varphi(x)$ będzie zdaniem w $\mathcal L_{\infty}$. Ma skończoną długość, zatem wszystkie użyte w nim symbole należą już do pewnego $\mathcal L_n$. Na etapie $n+1$ dodaliśmy dla niego nową stałą $c$ i aksjomat świadka. Ten aksjomat należy do $T_{n+1}$, a więc także do $T_{\infty}$. Otrzymujemy (ii).
\end{proof}

\begin{przyklad}[Dwa etapy nadawania nazw świadkom]
Rozważmy język z jednoargumentowymi symbolami relacyjnymi $P,Q$ oraz
$T=\{\exists x\,P(x),\ \forall x\,(P(x)\to Q(x))\}$.
Teoria jest niesprzeczna: w strukturze o uniwersum $\{a\}$ interpretujemy oba predykaty jako $\{a\}$.
Na pierwszym etapie możemy dołączyć stałą $c_P$ i zdanie
$(\exists x\,P(x))\to P(c_P)$. Z $T$ i tego aksjomatu przez $\to E$
otrzymujemy $P(c_P)$, a z $\forall x\,(P(x)\to Q(x))$ przez $\forall E$ otrzymujemy
$P(c_P)\to Q(c_P)$, więc także $Q(c_P)$.
W skonstruowanym później modelu termowym klasa $[c_P]$ spełnia oba predykaty.
Nie wybieraliśmy wcześniej konkretnego elementu: stała nazywa świadka dopiero w konstruowanym modelu.

Po pierwszym etapie w języku występuje również zdanie
$\exists x\,(x=c_P)$, którego nie było w języku wyjściowym.
Jest ono prawdziwe w każdej strukturze z interpretacją $c_P$, lecz w opisanej pełnej konstrukcji Henkina otrzymuje własny aksjomat świadka w następnym etapie.
To pokazuje, dlaczego jeden etap nie zapewnia własności świadków dla wszystkich zdań języka rozszerzonego.
\end{przyklad}

\section{Uzupełnienie do teorii maksymalnej}\label{sec:henkin-lindenbaum}
Samo $T_{\infty}$ może nie rozstrzygać każdego zdania. Chcemy powiększyć je do niesprzecznego zbioru, w którym dla każdego zdania $\sigma$ występuje dokładnie jedna z formuł $\sigma$, $\neg\sigma$. Ten krok nie jest dodatkową interpretacją prawdy: to własność syntaktycznego zbioru formuł.

\begin{definicja}[Maksymalny zbiór niesprzeczny]\label{def:henkin-max-consistent}
Zbiór zdań $\Delta$ danego języka jest \emph{maksymalnie niesprzeczny}, gdy jest niesprzeczny i każde jego właściwe rozszerzenie przez choć jedno nowe zdanie jest sprzeczne.
\end{definicja}

\begin{lemat}[Lemat Lindenbauma dla języka przeliczalnego]\label{lem:henkin-lindenbaum-countable}
Niech język $\mathcal L$ ma przeliczalnie wiele symboli, a $S$ będzie niesprzecznym zbiorem jego zdań. Wtedy istnieje maksymalnie niesprzeczny zbiór zdań $\Delta\supseteq S$.
\end{lemat}
\begin{proof}
Wszystkich skończonych napisów nad przeliczalnym alfabetem jest przeliczalnie wiele: dla każdej długości zbiór napisów jest przeliczalny, a łączymy przeliczalnie wiele długości. W szczególności można wypisać wszystkie zdania języka jako $\sigma_0,\sigma_1,\ldots$ (powtórzenia są nieszkodliwe). Kładziemy $S_0=S$. Na etapie $n$ badamy $\sigma_n$: jeśli $S_n\cup\{\sigma_n\}$ jest niesprzeczny, ustalamy $S_{n+1}=S_n\cup\{\sigma_n\}$; w przeciwnym wypadku ustalamy $S_{n+1}=S_n\cup\{\neg\sigma_n\}$. W tym drugim przypadku Lemat~\ref{lem:henkin-extension-neg} gwarantuje niesprzeczność wybranej gałęzi. Zbiór $\Delta=\bigcup_n S_n$ jest niesprzeczny, ponieważ każdy ewentualny skończony dowód sprzeczności byłby zawarty już w pewnym $S_n$. Dla każdego zdania $\sigma_n$ zbiór $\Delta$ zawiera $\sigma_n$ albo $\neg\sigma_n$, więc dołączenie zdania spoza $\Delta$ wprowadziłoby parę sprzecznych zdań. To jest żądana maksymalność.
\end{proof}

\begin{uwaga}[Języki o dowolnej liczbie symboli i aksjomat wyboru]\label{rem:henkin-zorn}
Dla dowolnego (także nieprzeliczalnego) języka stosujemy \emph{lemat Kuratowskiego--Zorna}: jeżeli częściowo uporządkowany zbiór ma własność, że każdy jego łańcuch ma ograniczenie górne w tym zbiorze, to istnieje w nim element maksymalny. Jest to konsekwencja aksjomatu wyboru, którego pełną teorię rozwiniemy w części zbiorowej skryptu. Tu zastosowanie jest konkretne: porządkujemy przez zawieranie wszystkie niesprzeczne zbiory zdań zawierające $T_{\infty}$. Rozważmy najpierw niepusty łańcuch $\mathcal C$. Jego suma $U=\bigcup\mathcal C$ zawiera $T_{\infty}$. Gdyby $U\vdash_{\mathrm{ND}}\bot$, skończony dowód używałby skończenie wielu zdań $\sigma_1,\ldots,\sigma_r$; każde z nich należy do któregoś $S_i\in\mathcal C$. Skończenie wiele elementów łańcucha ma największy element względem zawierania, zatem jeden zbiór z $\mathcal C$ zawiera wszystkie $\sigma_i$ i byłby sprzeczny. Stąd $U$ jest niesprzecznym ograniczeniem górnym. Dla pustego łańcucha ograniczeniem górnym jest $T_{\infty}$, nie zaś suma pustej rodziny. Lemat Kuratowskiego--Zorna daje więc maksymalny element $\Delta$. W ten sposób dalsza konstrukcja działa dla \emph{każdego} języka. W przypadku przeliczalnym nie trzeba korzystać z tego lematu, ponieważ poprzedni dowód podał bezpośredni przepis budowania $\Delta$. Jeżeli początkowy język był przeliczalny, każdy $\mathcal L_n$ i ich suma $\mathcal L_{\infty}$ pozostają przeliczalne, więc możemy zastosować właśnie ten przepis.
\end{uwaga}

\begin{lemat}[Własności maksymalnej teorii]\label{lem:henkin-max-properties}
Niech $\Delta$ będzie maksymalnie niesprzecznym zbiorem zdań. Wtedy dla każdego zdania $\sigma$:
\begin{enumerate}[label=(\roman*)]
\item dokładnie jedno z $\sigma$, $\neg\sigma$ należy do $\Delta$;
\item jeśli $\Delta\vdash_{\mathrm{ND}}\sigma$, to $\sigma\in\Delta$;
\item jeśli $\sigma\to\tau\in\Delta$ i $\sigma\in\Delta$, to $\tau\in\Delta$.
\end{enumerate}
\end{lemat}
\begin{proof}
Jeśli $\sigma\notin\Delta$, to na mocy maksymalności $\Delta\cup\{\sigma\}$ jest sprzeczny. Z Lematu~\ref{lem:henkin-extension-neg} wynika $\Delta\vdash\neg\sigma$. Dołączenie zdania, które ma już dowód z $\Delta$, nie może stworzyć nowej sprzeczności: w hipotetycznym dowodzie $\bot$ każde użycie nowej przesłanki zastąpilibyśmy jej dowodem z $\Delta$. Z maksymalności wynika więc $\neg\sigma\in\Delta$. Gdyby należały obie formuły, przez $\neg E$ mielibyśmy $\Delta\vdash\bot$, sprzecznie z założeniem. To dowodzi (i). To samo rozumowanie o zastępowaniu nowej przesłanki jej dowodem dowodzi (ii). W (iii) reguła $\to E$ daje $\Delta\vdash\tau$, po czym używamy (ii).
\end{proof}

\section{Model termowy i równość}\label{sec:henkin-term-model}
Niech $\Delta\supseteq T_{\infty}$ będzie maksymalnie niesprzecznym zbiorem zdań języka $\mathcal L_{\infty}$. Każdy obiekt modelu będziemy reprezentować przez \emph{term domknięty} tego języka. Musimy jednak utożsamiać termy, których równość jest dowodliwa z $\Delta$.

\begin{definicja}[Równoważność termów]\label{def:henkin-term-equivalence}
Dla termów domkniętych $t,u$ języka $\mathcal L_{\infty}$ definiujemy
\[
t\sim_{\Delta}u\quad\Longleftrightarrow\quad\Delta\vdash_{\mathrm{ND}}t=u.
\]
Przez $[t]$ oznaczamy klasę równoważności termu $t$ względem tej relacji.
\end{definicja}

\begin{lemat}[Równość termów jest zgodna z operacjami]\label{lem:henkin-congruence}
Relacja $\sim_{\Delta}$ jest równoważnością. Jeżeli $t_i\sim_{\Delta}u_i$ dla $i=1,\ldots,n$, to dla każdego $n$-argumentowego symbolu funkcyjnego $f$ zachodzi
$f(t_1,\ldots,t_n)\sim_{\Delta}f(u_1,\ldots,u_n)$.
Dla symbolu relacyjnego $R$ mamy równoważność
\[
R(t_1,\ldots,t_n)\in\Delta
\quad\Longleftrightarrow\quad
R(u_1,\ldots,u_n)\in\Delta.
\]
\end{lemat}
\begin{proof}
Zwrotność wynika z reguły $=I$, a symetrię wyprowadziliśmy już z $=E$ w rozdziale~\ref{ch:natural-deduction}. Dla przechodniości niech $t,u,v$ będą termami domkniętymi i $z$ zmienną, która nie występuje w tych termach. Regułę $=E$ stosujemy do $u=v$ oraz $\theta[z:=u]$, gdzie $\theta(z)$ ma postać $t=z$. Z przesłanek $t=u$ i $u=v$ otrzymujemy $\theta[z:=v]$, czyli $t=v$. Zatem $\sim_{\Delta}$ jest równoważnością.

Niech $F_i:=f(u_1,\ldots,u_i,t_{i+1},\ldots,t_n)$ dla $i=0,\ldots,n$. Dla każdego $i=1,\ldots,n$ dobieramy zmienną $z$ nieobecną w żadnym z termów domkniętych i rozważamy formułę
\[
\theta_i(z):\quad F_{i-1}=f(u_1,\ldots,u_{i-1},z,t_{i+1},\ldots,t_n).
\]
Refleksyjność daje równość $F_{i-1}=F_{i-1}$, która jest dokładnie formułą $\theta_i[z:=t_i]$. Z $t_i=u_i$ i reguły $=E$ otrzymujemy $\theta_i[z:=u_i]$, czyli $F_{i-1}=F_i$. Składając te równości przez przechodniość, dostajemy $F_0=F_n$, a więc zgodność symbolu funkcyjnego z $\sim_{\Delta}$. Dla symbolu relacyjnego bierzemy w $i$-tym kroku $\rho_i(z):=R(u_1,\ldots,u_{i-1},z,t_{i+1},\ldots,t_n)$. Z $R(t_1,\ldots,t_n)$ i kolejnych $t_i=u_i$ reguła $=E$ daje $R(u_1,\ldots,u_n)$. Zamkniętość $\Delta$ na dowody wynika z Lematu~\ref{lem:henkin-max-properties}; w przeciwnym kierunku stosujemy symetrię każdej równości i wykonujemy analogiczne kroki w odwrotnej kolejności.
\end{proof}

\begin{definicja}[Konstrukcja modelu termowego]\label{def:henkin-term-model}
Niech $\operatorname{CTerm}(\mathcal L_{\infty})$ oznacza zbiór wszystkich termów domkniętych języka $\mathcal L_{\infty}$. Definiujemy strukturę $\mathfrak M$ języka $\mathcal L_{\infty}$ następująco:
\begin{enumerate}[label=(\roman*)]
\item jej uniwersum jest $M=\operatorname{CTerm}(\mathcal L_{\infty})/{\sim_{\Delta}}$;
\item dla symbolu stałej $c$ kładziemy $c^{\mathfrak M}=[c]$;
\item dla symbolu funkcyjnego $f$ arności $n$ kładziemy
\[
f^{\mathfrak M}([t_1],\ldots,[t_n])=[f(t_1,\ldots,t_n)];
\]
\item dla symbolu relacyjnego $R$ arności $n$ ustalamy
\[
([t_1],\ldots,[t_n])\in R^{\mathfrak M}
\quad\Longleftrightarrow\quad R(t_1,\ldots,t_n)\in\Delta.
\]
\end{enumerate}
\end{definicja}

\begin{uwaga}[Dlaczego definicja ma sens?]\label{rem:henkin-well-defined}
Zbiór $M$ jest niepusty, ponieważ w $\mathcal L_0$ umieściliśmy stałą $d$ i klasa $[d]$ jest jego elementem. Jeżeli przedstawimy tę samą klasę przez inny term, to Lemat~\ref{lem:henkin-congruence} gwarantuje identyczny wynik funkcji i identyczną ocenę prawdziwości atomu relacyjnego. Bez tego sprawdzenia wzory nie definiowałyby funkcji i relacji na \emph{klasach} termów.
\end{uwaga}

\begin{przyklad}[Dlaczego trzeba ilorazować termy]
Załóżmy, że język zawiera stałe $a,b$, jednoargumentowy symbol funkcyjny $f$
oraz jednoargumentowy symbol relacyjny $P$, a $a=b\in\Delta$.
Reguła $=E$ daje $f(a)=f(b)$ oraz oba kierunki
$P(a)\leftrightarrow P(b)$. Dlatego $[a]=[b]$ i $[f(a)]=[f(b)]$.
Gdyby uniwersum składało się z samych napisów $a,b,f(a),f(b)$,
a równość była identycznością napisów, zdanie $a=b$ byłoby fałszywe
mimo przynależności do $\Delta$.
Iloraz przez $\sim_{\Delta}$ usuwa tę przeszkodę.
\end{przyklad}

\section{Lemat o prawdzie: najważniejszy krok}\label{sec:henkin-truth-lemma}
W całej tej sekcji zapis $\mathfrak M\models\varphi[s]$ oznacza,
że formuła $\varphi$ jest spełniona w strukturze $\mathfrak M$
przy wartościowaniu $s$, zgodnie z Definicją~\ref{def:sem-satisfaction}.
Zapis $\varphi[\bar x:=\bar t]$ jest skrótem \emph{jednoczesnego}
podstawienia termów $\bar t$ za odpowiednie \emph{wolne}
wystąpienia zmiennych $\bar x$. Tutaj podstawiane termy są domknięte,
więc nie mogą ulec związaniu przez kwantyfikator.

\begin{przyklad}[Jak czytać lemat o prawdzie]
Jeżeli $P$ jest jednoargumentowym symbolem relacyjnym,
$a$ jest nową stałą, a $s(x)=[a]$, to równoważność
$\mathfrak M\models P(x)[s]\iff P(a)\in\Delta$
ma bardzo konkretny sens. Lewa strona mówi:
„element modelu reprezentowany przez term $a$ należy do
interpretacji relacji $P$”. Prawa mówi: „zdanie $P(a)$
znajduje się w maksymalnej teorii $\Delta$”.
Lemat o prawdzie orzeka, że obie metody odczytania zgadzają się.
\end{przyklad}
Uzasadnimy, że skonstruowany model interpretuje formuły dokładnie tak, jak przewiduje teoria $\Delta$.

\begin{lemat}[Wartości termów w modelu termowym]\label{lem:henkin-term-values}
Niech $t$ będzie termem ze zmiennymi wolnymi wśród $x_1,\ldots,x_k$, a $s$ wartościowaniem w $\mathfrak M$ takim, że $s(x_i)=[u_i]$ dla termów domkniętych $u_i$. Wtedy
\[
t^{\mathfrak M}[s]=[t[x_1:=u_1,\ldots,x_k:=u_k]].
\]
\end{lemat}
\begin{proof}
Indukcja po budowie termu. Dla zmiennej $x_i$ prawa strona wynosi $[u_i]=s(x_i)$; dla stałej $c$ obie strony są równe $[c]$. Gdy $t=f(t_1,\ldots,t_n)$, stosujemy założenie indukcyjne do wszystkich $t_i$, a następnie definicję $f^{\mathfrak M}$. Otrzymujemy dokładnie klasę termu powstałego przez wskazane podstawienie. Wszystkie podstawiane termy są domknięte, więc nie występuje problem przechwycenia zmiennych.
\end{proof}

\begin{twierdzenie}[Lemat o prawdzie]\label{thm:henkin-truth}
Niech $\varphi$ będzie formułą języka $\mathcal L_{\infty}$ o zmiennych wolnych wśród $x_1,\ldots,x_k$. Niech $u_1,\ldots,u_k$ będą termami domkniętymi i $s$ wartościowaniem w $\mathfrak M$ spełniającym $s(x_i)=[u_i]$. Wtedy
\[
\mathfrak M\models\varphi[s]
\quad\Longleftrightarrow\quad
\varphi[x_1:=u_1,\ldots,x_k:=u_k]\in\Delta.
\]
W szczególności dla każdego zdania $\sigma$ zachodzi $\mathfrak M\models\sigma$ wtedy i tylko wtedy, gdy $\sigma\in\Delta$.
\end{twierdzenie}
\begin{proof}
Dowodzimy tezy indukcyjnie po budowie $\varphi$. Dla atomu $t=v$ Lemat~\ref{lem:henkin-term-values} pokazuje, że wartości termów są równe wtedy i tylko wtedy, gdy ich podstawione wersje należą do tej samej klasy, czyli gdy $\Delta\vdash t[\bar u]=v[\bar u]$. Z zamkniętości $\Delta$ na dowody dostajemy równoważność z przynależnością atomu $t[\bar u]=v[\bar u]$ do $\Delta$. Dla atomu $R(t_1,\ldots,t_n)$ wynik jest bezpośrednią konsekwencją definicji $R^{\mathfrak M}$.

Rozważmy negację $\neg\psi$. Z hipotezy indukcyjnej
$\mathfrak M\models\psi[s]$ dokładnie wtedy, gdy
$\psi[\bar u]\in\Delta$. Z Lematu~\ref{lem:henkin-max-properties}
dokładnie jedna z formuł $\psi[\bar u]$ i
$\neg\psi[\bar u]$ należy do $\Delta$.
Zatem semantyczna negacja odpowiada syntaktycznej negacji.

Sprawdźmy wszystkie spójniki binarne, aby nie ukrywać
żadnego kroku za słowem „analogicznie”. Dla koniunkcji
$\psi\land\theta$ przynależność do $\Delta$ wymusza
przynależność obu składników przez $\land E_1$ i $\land E_2$;
na odwrót, dwa składniki z $\Delta$ dają koniunkcję przez
$\land I$. Dla alternatywy: jeśli w $\Delta$ jest jeden
składnik, to jest także ich alternatywa przez $\lor I$.
Jeżeli zaś alternatywa należy do $\Delta$, lecz żaden składnik
nie należy, to z maksymalności $\Delta$ obie negacje
$\neg\psi$ i $\neg\theta$ należą do $\Delta$;
reguła $\lor E$, stosowana do obu gałęzi,
prowadziłaby wtedy do sprzeczności. To niemożliwe.
Zatem alternatywa należy do $\Delta$ dokładnie wtedy,
gdy należy co najmniej jeden z jej członów.

Dla implikacji $\psi\to\theta$ rozpatrzmy dwa przypadki.
Jeżeli $\psi\notin\Delta$, to $\neg\psi\in\Delta$,
a z $\neg\psi$ możemy formalnie wyprowadzić
$\psi\to\theta$: tymczasowe założenie $\psi$
daje sprzeczność z $\neg\psi$,
ze sprzeczności otrzymujemy $\theta$ i rozładowujemy
$\psi$. Jeśli $\theta\in\Delta$, to również
otrzymujemy $\psi\to\theta$, wprowadzając implikację
z niewykorzystanym założeniem $\psi$.
Odwrotnie, jeśli $\psi\to\theta\in\Delta$
i $\psi\in\Delta$, reguła $\to E$ daje
$\theta\in\Delta$. Stąd implikacja należy do $\Delta$
dokładnie wtedy, gdy poprzednik nie należy do $\Delta$
lub należy do niego następnik.
Wreszcie $\psi\leftrightarrow\theta$ należy do $\Delta$
dokładnie wtedy, gdy należą do niego obie implikacje
$\psi\to\theta$ i $\theta\to\psi$, co wynika z
reguł $\leftrightarrow I$ i $\leftrightarrow E$.
W każdym przypadku odtworzyliśmy dokładnie klasyczną
tabelę prawdziwości odpowiedniego spójnika;
hipoteza indukcyjna kończy więc dowód kroku dla formuł złożonych.

Niech teraz $\varphi$ ma postać $\exists x\,\psi(x,\bar z)$, gdzie $\bar z$ symbolizuje pozostałe zmienne wolne. Po zastąpieniu $\bar z$ termami domkniętymi $\bar u$ otrzymujemy zdanie $\exists x\,\psi(x,\bar u)$. Jeśli $\mathfrak M\models(\exists x\,\psi(x,\bar z))[s]$, istnieje $a\in M$ spełniający formułę $\psi$ przy zmianie wartościowania $x\mapsto a$. Każdy $a$ ma postać $[t]$ dla pewnego termu domkniętego $t$. Założenie indukcyjne daje $\psi(t,\bar u)\in\Delta$. Reguła $\exists I$ i zamkniętość $\Delta$ na dowody dają $\exists x\,\psi(x,\bar u)\in\Delta$. Odwrotnie, jeśli to zdanie egzystencjalne należy do $\Delta$, własność świadków z Twierdzenia~\ref{thm:henkin-chain-properties} dostarcza stałej $c$ i zdania $(\exists x\,\psi(x,\bar u))\to\psi(c,\bar u)$ w $\Delta$. Reguła $\to E$ daje $\psi(c,\bar u)\in\Delta$. Z hipotezy indukcyjnej $\mathfrak M\models\psi[s[x\mapsto[c]]]$, skąd $\mathfrak M\models(\exists x\,\psi)[s]$.

Rozpatrzmy teraz formułę $\forall x\,\psi(x,\bar z)$.
Jeśli $\forall x\,\psi(x,\bar u)\in\Delta$, to dla
każdego termu domkniętego $t$ reguła $\forall E$
daje $\psi(t,\bar u)\in\Delta$.
Z hipotezy indukcyjnej $\mathfrak M\models\psi[s[x\mapsto[t]]]$.
Ponieważ \emph{każdy} element $M$ jest klasą jakiegoś
termu domkniętego, otrzymujemy
$\mathfrak M\models(\forall x\,\psi(x,\bar z))[s]$.

Odwrotnie, przypuśćmy, że formuła uniwersalna jest
spełniona w modelu, ale zdanie
$\forall x\,\psi(x,\bar u)$ \emph{nie} należy do $\Delta$.
Maksymalność $\Delta$ daje
$\neg\forall x\,\psi(x,\bar u)\in\Delta$.
Na mocy Lematu~\ref{lem:henkin-quant-duality}
w $\Delta$ należy wtedy również
$\exists x\,\neg\psi(x,\bar u)$.
Własność świadków pozwala wybrać nową stałą $c$
i otrzymać $\neg\psi(c,\bar u)\in\Delta$.
Z już sprawdzonego przypadku negacji oraz hipotezy
indukcyjnej wynika
$\mathfrak M\not\models\psi[s[x\mapsto[c]]]$.
To przeczy założeniu
$\mathfrak M\models(\forall x\,\psi)[s]$.
Zatem zdanie uniwersalne musi należeć do $\Delta$.
Przypadki atomowe, wszystkie spójniki i oba kwantyfikatory
zostały rozpatrzone; indukcja jest zakończona.
\end{proof}

\begin{proof}[Dowód Twierdzenia~\ref{thm:henkin-model-existence}]
Z niesprzecznego $T$ otrzymaliśmy niesprzeczne $T_{\infty}$, następnie maksymalnie niesprzeczną teorię $\Delta\supseteq T_{\infty}$, a na końcu model termowy $\mathfrak M$. Dla każdego $\sigma\in T$ mamy $\sigma\in\Delta$. Z Lematu o prawdzie wynika $\mathfrak M\models\sigma$. Ponieważ $\sigma$ było dowolne, $\mathfrak M\models T$. Struktura języka rozszerzonego po zapomnieniu nowych stałych daje model w języku wyjściowym $\mathcal L$.
\end{proof}

\begin{przyklad}[Jak lemat o prawdzie działa dla formuły z kwantyfikatorem]
Niech $R$ będzie jednoargumentowym symbolem relacyjnym.
Rozważmy zdanie $\forall x\,R(x)$ w języku $\mathcal L_{\infty}$.
Jeżeli należy do $\Delta$, to dla każdego termu domkniętego $t$
również $R(t)\in\Delta$ przez $\forall E$ i zamkniętość $\Delta$ na dowody.
Każdy element modelu jest klasą $[t]$, więc $\mathfrak M\models\forall x\,R(x)$.
Jeżeli natomiast $\forall x\,R(x)\notin\Delta$, to
$\neg\forall x\,R(x)\in\Delta$, a stąd
$\exists x\,\neg R(x)\in\Delta$.
Własność świadków dostarcza stałej $c$ z $\neg R(c)\in\Delta$;
klasa $[c]$ obala zdanie uniwersalne w modelu.
Oba kierunki są rozpatrzone oddzielnie, a fakt, że każdy element ma nazwę,
wynika z konstrukcji uniwersum $M$.
\end{przyklad}

\section{Od istnienia modelu do pełności}\label{sec:henkin-completeness}
\begin{twierdzenie}[Pełność dla zbiorów zdań]\label{thm:henkin-completeness-sentences}
Niech $T$ będzie zbiorem zdań, a $\sigma$ zdaniem tego samego dowolnego języka pierwszego rzędu z równością. Jeśli $T\models\sigma$, to $T\vdash_{\mathrm{ND}}\sigma$.
\end{twierdzenie}
\begin{proof}
Przypuśćmy, że $T\nvdash\sigma$. Wówczas $T\cup\{\neg\sigma\}$ jest niesprzeczny: gdyby był sprzeczny, reguła klasycznego dowodu nie wprost (RAA) pozwoliłaby rozładować $\neg\sigma$ i uzyskać $T\vdash\sigma$. Na mocy Twierdzenia~\ref{thm:henkin-model-existence} istnieje model $\mathfrak M$ zbioru $T\cup\{\neg\sigma\}$. Spełnia on każde zdanie z $T$, ale nie spełnia $\sigma$. To przeczy $T\models\sigma$. Zatem $T\vdash\sigma$.
\end{proof}

\begin{lemat}[Zamiana wolnych zmiennych na parametry]\label{lem:henkin-parameters}
Niech $\Gamma$ będzie zbiorem formuł, a $\varphi$ formułą języka $\mathcal L$. Dla każdej zmiennej wolnej w którymś z tych napisów wprowadźmy osobną nową stałą $d_x$. Przez $\Gamma^d$ i $\varphi^d$ oznaczmy wyniki \emph{wyłącznie wolnych} wystąpień $x$ przez $d_x$. Wtedy
\[
\Gamma^d\vdash_{\mathrm{ND}}\varphi^d
\quad\Longrightarrow\quad
\Gamma\vdash_{\mathrm{ND}}\varphi.
\]
\end{lemat}
\begin{proof}
Niech $D$ będzie skończonym drzewem dowodu $\Gamma^d\vdash_{\mathrm{ND}}\varphi^d$ w języku rozszerzonym. Oznaczmy przez $C=\{d_{x_1},\ldots,d_{x_m}\}$ zbiór nowych stałych rzeczywiście występujących w $D$; jest on skończony. Wszystkie wykorzystane otwarte przesłanki należą do $\Gamma^d$, czyli są zdaniami, i w $D$ występuje jedynie skończenie wiele pomocniczych zmiennych.

Najpierw uzasadnijmy dopuszczalność przemianowania zmiennych, na którą będziemy się powoływać. Z reguł $\forall E$, $\forall I$ można, przy świeżej zmiennej $z$, przeprowadzić dowody obu kierunków między $\forall x\,\alpha$ oraz $\forall z\,\alpha[x:=z]$: w jedną stronę instancjujemy pierwszy kwantyfikator przez $z$ i uogólniamy na $z$; w drugą stronę analogicznie podstawiamy $x$ i uogólniamy na $x$. Dla $\exists$ analogiczne dwa kierunki uzyskujemy, otwierając świadka regułą $\exists E$, a następnie stosując $\exists I$, po czym rozładowując świadka. W każdym z tych przekształceń zmienną nową wybieramy poza wszystkimi formułami i otwartymi założeniami danego fragmentu. Indukcja po budowie kontekstu formuły z użyciem pozostałych reguł pozwala uzyskać dowody obu kierunków dla dowolnej dopuszczalnej zmiany nazwy zmiennej związanej. Nie przyjmujemy więc przemianowania jako dodatkowej reguły bez uzasadnienia.

W drzewie $D$ kolejno zmieniamy na nowe nazwy wszystkie zmienne związane i własne kroków $\forall I$ oraz $\exists E$, które mogłyby kolidować z $x_1,\ldots,x_m$. Przemianowanie zmiennej własnej w całej odpowiedniej poddedukcji jest poprawne, o ile nowa nazwa nie występuje w żadnym jej pozostałym otwartym założeniu, we wniosku $\exists E$ ani w przesłance głównej; można to zapewnić, wybierając nazwy spoza całego skończonego drzewa. Przy $\forall I$ nowa nazwa pozostaje niewolna w otwartych założeniach. Gdy nazwa występuje w podformule wiązanej innym kwantyfikatorem, stosujemy poprzednio uzasadnione przemianowanie zmiennej związanej. Po wykonaniu tych skończenie wielu operacji żaden kwantyfikator ani warunek świeżości nie koliduje z $x_i$.

Zastępujemy teraz równocześnie w całym drzewie każde wystąpienie stałej $d_{x_i}$ zmienną $x_i$. Dla podstawień termów zachowujemy dopuszczalność, ponieważ żadna ze zmiennych $x_i$ nie zostaje przechwycona. W krokach $\forall I$ uogólniana zmienna nie jest żadną z $x_i$, a w $\exists E$ zmienna świadka nie jest żadną z nich; zatem nowe wolne wystąpienia $x_i$ nie naruszają warunków ubocznych. Reguły zdaniowe i reguły równości są zachowane po jednolitej zamianie w całych formułach i termach. Otrzymujemy poprawny dowód z przesłanek będących dopuszczalnymi przemianowaniami pierwotnych formuł z $\Gamma$, którego wniosek jest dopuszczalnym przemianowaniem $\varphi$. Wprowadzając w miejscach użycia przesłanek dowody odpowiednich kierunków równoważności z poprzedniego akapitu i postępując tak samo przy wniosku, dostajemy dosłownie dowód $\Gamma\vdash_{\mathrm{ND}}\varphi$.
\end{proof}

\begin{twierdzenie}[Twierdzenie Gödla o pełności logiki pierwszego rzędu]\label{thm:godel-completeness}
Dla dowolnego języka pierwszego rzędu z równością, dowolnego zbioru formuł $\Gamma$ i dowolnej formuły $\varphi$ zachodzi
\[
\Gamma\models\varphi\quad\Longleftrightarrow\quad\Gamma\vdash_{\mathrm{ND}}\varphi.
\]
Z lewej strony $\models$ dotyczy wszystkich struktur \emph{i wartościowań}, zgodnie z Definicją~\ref{def:soundness-entailment-formulas}.
\end{twierdzenie}
\begin{proof}
Implikacja od prawej do lewej jest Twierdzeniem o poprawności~\ref{thm:soundness}. Załóżmy teraz $\Gamma\models\varphi$. Dodajmy nowe stałe $d_x$ i utwórzmy zdania $\Gamma^d,\varphi^d$ z Lematu~\ref{lem:henkin-parameters}. W każdej strukturze rozszerzonego języka interpretacje $d_x$ wyznaczają wartościowanie zmiennych $x$. Dlatego spełnianie $\Gamma^d$ pociąga spełnianie $\varphi^d$; inaczej po przypisaniu każdej zmiennej $x$ wartości $d_x$ otrzymalibyśmy kontrprzykład dla $\Gamma\models\varphi$. Zatem $\Gamma^d\models\varphi^d$. Pełność dla zdań z Twierdzenia~\ref{thm:henkin-completeness-sentences} daje $\Gamma^d\vdash\varphi^d$. Lemat~\ref{lem:henkin-parameters} pozwala wrócić do języka wyjściowego: $\Gamma\vdash\varphi$.
\end{proof}

\begin{przyklad}[Wniosek z dwóch kwantyfikatorów]
Nie wolno z samego $\exists x\,P(x)$ wywnioskować
$\exists x\,Q(x)$: konieczna jest dodatkowa przesłanka,
że każdy obiekt mający cechę $P$ ma również cechę $Q$.
Wyjaśnimy osobno rolę obu założeń i użycie świeżej zmiennej.
Rozważmy zbiór założeń\, $T=\{\forall x(P(x)\to Q(x)),\ \exists x P(x)\}$.\par\noindent\ignorespaces
Zbiór $T$ semantycznie implikuje $\exists x Q(x)$. W dowolnej strukturze, w której oba założenia są prawdziwe, drugi warunek daje świadka $a$ z $P(a)$, pierwszy daje $P(a)\to Q(a)$, więc zachodzi $Q(a)$. Twierdzenie o pełności gwarantuje istnienie formalnego dowodu; możemy go też rozpisać wprost: z $\exists x P(x)$ bierzemy w poddedukcji świeże $P(y)$, z $\forall x(P(x)\to Q(x))$ otrzymujemy $P(y)\to Q(y)$, przez $\to E$ uzyskujemy $Q(y)$, przez $\exists I$ uzyskujemy $\exists x Q(x)$ i na końcu rozładowujemy $P(y)$ przez $\exists E$.
\end{przyklad}

\begin{uwaga}[Dlaczego nie jest to twierdzenie o niezupełności?]
W pełności mówimy o wszystkich \emph{logicznych konsekwencjach} ustalonych przesłanek. Nie twierdzimy, że każda teoria $T$ rozstrzyga każde zdanie $\sigma$ (tj. że $T\vdash\sigma$ lub $T\vdash\neg\sigma$). Teorie arytmetyczne mogą nie rozstrzygać zdań, ponieważ istnieją różne modele tej samej teorii. Własności tej nie należy mylić z niedoskonałością rachunku logicznego.
\end{uwaga}
\chapter[Algorytmy i obliczalność]{Algorytmy, obliczalność i nierozstrzygalność}
\label{ch:computability-intro}

\section{Pytanie inne niż prawdziwość}
We wcześniejszych rozdziałach poznaliśmy zdania prawdziwe
w strukturach i zdania posiadające formalne dowody.
Teraz zadamy inne pytanie: czy istnieje \emph{jedna skończona procedura},
która dla dowolnych danych rozstrzyga określoną własność?
Możemy na przykład sprawdzić, czy dany napis jest poprawnie zbudowanym
termem języka pierwszego rzędu, analizując go od lewej do prawej.
To jest zadanie algorytmiczne. Tymczasem nie każde pytanie
o obiektach matematycznych posiada algorytm rozstrzygający.

Nie chcemy opierać się na niejasnym określeniu
„zadanie da się wykonać mechanicznie”. Wprowadzimy więc prosty model
maszyny wykonującej pojedyncze, jednoznacznie określone kroki.
Pokażemy następnie za pomocą argumentu przekątniowego,
że istnieje konkretne pytanie, na które nie odpowie
żaden algorytm działający poprawnie dla wszystkich danych.

\section{Napisy, kodowanie i problemy decyzyjne}
\begin{definicja}[Skończony alfabet i zbiór słów]
Niech $\Sigma$ będzie skończonym, niepustym zbiorem symboli.
Przez $\Sigma^*$ oznaczamy zbiór wszystkich słów skończonych
nad $\Sigma$, łącznie ze słowem pustym $\varepsilon$.
Długość słowa $w$ oznaczamy przez $|w|$.
\emph{Problemem decyzyjnym} nazywamy podzbiór $L\subseteq\Sigma^*$.
Pytanie brzmi: dla podanego słowa $w$ ustalić, czy $w\in L$.
\end{definicja}

\begin{przyklad}[Prosty problem decyzyjny]
Niech $\Sigma=\{0,1\}$, a $L$ składa się ze słów
zawierających parzystą liczbę jedynek.
Słowo $10101$ nie należy do $L$, bo ma trzy jedynki,
a słowo $1010$ należy, bo ma ich dwie.
Możemy rozstrzygnąć przynależność do $L$,
przechodząc po słowie raz i zapamiętując tylko,
czy dotychczas napotkaliśmy parzystą, czy nieparzystą liczbę jedynek.
\end{przyklad}

\begin{definicja}[Kodowanie par słów]
Dla słów binarnych $u,v$ określmy
\[
\operatorname{code}(u,v):=0^{|u|}1uv.
\]
Zapis $0^{|u|}$ oznacza ciąg $|u|$ zer stojących obok siebie;
nie jest potęgą liczby zero. Kod zaczyna się od tylu zer,
ile bitów ma pierwszy składnik, po czym zapisujemy jedynkę,
a następnie oba słowa.
\end{definicja}

\begin{przyklad}[Jak odczytać kod pary]
Dla $u=10$ i $v=01$ otrzymujemy
$\operatorname{code}(u,v)=0011001$.
Dwie początkowe zera oznaczają, że pierwszy składnik
ma długość dwa. Po pierwszej jedynce oddzielającej
odczytujemy więc dwa następne bity $10$; pozostałe bity $01$
tworzą drugi składnik. Kod daje się jednoznacznie odczytać.
\end{przyklad}

\begin{lemat}[Jednoznaczność kodowania par]\label{lem:computability-pair-code}
Dla dowolnych słów $u,v,u',v'$ nad $\{0,1\}$ zachodzi
\[
\operatorname{code}(u,v)=\operatorname{code}(u',v')
\quad\Longrightarrow\quad u=u'\ \text{oraz}\ v=v'.
\]
\end{lemat}
\begin{proof}
Każdy kod rozpoczyna się jednoznacznie określonym
blokiem zer zakończonym pierwszą jedynką.
Liczba tych zer jest długością pierwszego składnika.
Jeżeli dwa kody są równe, mają taki sam początkowy blok,
a więc $|u|=|u'|$. Kolejne $|u|$ bitów muszą być takie same,
stąd $u=u'$. Pozostała część kodu to odpowiednio $v$ i $v'$, 
więc również te słowa są równe. Opis ten daje zarazem algorytm
odkodowania: liczymy początkowe zera do pierwszej jedynki,
po czym odczytujemy tyle bitów pierwszego składnika.
Brak jedynki rozdzielającej albo zbyt mało dalszych bitów
oznacza, że napis nie jest kodem żadnej pary.
\end{proof}

\section{Maszyna Turinga: elementarne instrukcje}
\begin{definicja}[Deterministyczna maszyna Turinga]\label{def:computability-tm}
\emph{Deterministyczną maszyną Turinga} nazywamy skończony opis
złożony z:
\begin{enumerate}[label=(\roman*)]
\item skończonego zbioru stanów $Q$, zawierającego
stan początkowy $q_0$ oraz dwa różne stany końcowe
$q_{\mathrm{tak}}$ i $q_{\mathrm{nie}}$, przy czym
$q_0\notin\{q_{\mathrm{tak}},q_{\mathrm{nie}}\}$;
\item skończonego alfabetu taśmy $\Omega$ z wyróżnionym
symbolem pustym $\square$ oraz niepustym podzbiorem symboli wejściowych
$\Sigma\subseteq\Omega\setminus\{\square\}$;
\item funkcji przejścia
\[
\delta:(Q\setminus\{q_{\mathrm{tak}},q_{\mathrm{nie}}\})\times\Omega
\longrightarrow Q\times\Omega\times\{L,P,S\},
\]
której wartość $\delta(q,a)=(q',b,d)$ ma następujący sens:
„gdy maszyna jest w stanie $q$ i czyta $a$, zapisuje $b$,
przechodzi do stanu $q'$ i przesuwa głowicę
w lewo ($L$), prawo ($P$) albo pozostawia ją w miejscu ($S$)”.
\end{enumerate}
Maszyna ma taśmę o komórkach ponumerowanych liczbami
całkowitymi. Na początku wejściowe słowo jest zapisane
w kolejnych komórkach zaczynających się od komórki zero,
a wszystkie pozostałe komórki zawierają $\square$.
Głowica startuje nad komórką zero. Odczytanie stanu
$q_{\mathrm{tak}}$ oznacza \emph{akceptację}, a
$q_{\mathrm{nie}}$ \emph{odrzucenie}. W obu przypadkach
maszyna zatrzymuje się. Jeżeli nie osiągnie żadnego
ze stanów końcowych, wykonuje nieskończenie wiele kroków.
\end{definicja}

\begin{uwaga}[Ważne rozróżnienie: skończona maszyna i nieograniczony czas]
Opis maszyny, a więc wykaz stanów oraz tabela przejść,
jest skończony. Nie znaczy to jednak, że dla każdego wejścia
liczba wykonanych kroków ma ustaloną wspólną granicę.
Maszyna może poruszać się po dowolnie długich fragmentach
taśmy, a dla niektórych wejść może w ogóle się nie zatrzymać.
To właśnie możliwość nieograniczonego obliczania będzie istotna
w problemie stopu.
\end{uwaga}

\begin{przyklad}[Sprawdzenie pierwszego symbolu]
Niech alfabet wejściowy będzie równy $\{0,1\}$.
Maszyna może w stanie $q_0$ odczytać pierwszy symbol
wejściowego słowa. Jeśli jest to $1$, przechodzi
bezpośrednio do $q_{\mathrm{tak}}$, a jeśli jest to $0$
lub symbol pusty, przechodzi do $q_{\mathrm{nie}}$.
Każde obliczenie kończy się po jednym kroku.
Jest to prosty algorytm rozstrzygający własność
„słowo zaczyna się od jedynki”.
\end{przyklad}

\section{Decydowanie i rozpoznawanie}
\begin{definicja}[Zbiór rozstrzygalny i rozpoznawalny]\label{def:computability-decidable-c.e.}
Niech $L\subseteq\Sigma^*$ będzie problemem decyzyjnym.
\begin{enumerate}[label=(\roman*)]
\item Zbiór $L$ nazywamy \emph{rozstrzygalnym}, jeśli istnieje
maszyna Turinga, która zatrzymuje się dla \emph{każdego} wejścia $w$
i akceptuje dokładnie wtedy, gdy $w\in L$.
\item Zbiór $L$ nazywamy \emph{rozpoznawalnym}
(lub \emph{rekurencyjnie przeliczalnym}), jeśli istnieje
maszyna Turinga akceptująca dokładnie słowa z $L$.
Dla słów spoza $L$ maszyna może je odrzucić albo
nie zatrzymać się nigdy.
\end{enumerate}
\end{definicja}

\begin{przyklad}[Różnica między dwiema obietnicami]
Jeżeli algorytm \emph{rozstrzyga} własność,
możemy uruchomić go na dowolnych danych i otrzymać
po skończonym czasie odpowiedź „tak” lub „nie”.
Jeżeli algorytm jedynie \emph{rozpoznaje} własność,
to dla przypadku pozytywnego w końcu zobaczymy „tak”,
ale dla negatywnego możemy nie otrzymać żadnej odpowiedzi.
Z samego długiego oczekiwania nie wynika,
że odpowiedź powinna być „nie”.
\end{przyklad}

\begin{lemat}[Zbiór rozstrzygalny jest rozpoznawalny]\label{lem:computability-decidable-recognizable}
Każdy rozstrzygalny zbiór słów jest rozpoznawalny.
\end{lemat}
\begin{proof}
Niech maszyna $M$ rozstrzyga $L$.
Dla słowa należącego do $L$ maszyna zatrzymuje się
i akceptuje. Dla słowa nienależącego zatrzymuje się
i odrzuca. Są to dokładnie warunki wymagane
od maszyny rozpoznającej $L$, która może, lecz nie musi,
zatrzymywać się także na przypadkach negatywnych.
Dlatego ta sama maszyna $M$ rozpoznaje $L$.
\end{proof}

\section{Program też można zapisać jako słowo}
Każda maszyna Turinga ma skończenie wiele stanów,
symboli i wpisów w tabeli przejść. Możemy więc
zapisać jej cały opis jako skończone słowo binarne.
Ustalamy numerację stanów i symboli, zapisujemy ich liczby
oraz kolejne wpisy tabeli przejść, oddzielając pola kodami
ich długości (na przykład przez prefiks długościowy,
jak w Lemacie~\ref{lem:computability-pair-code}).
Zapis obejmuje także stan początkowy, dwa stany końcowe,
symbol pusty i alfabet wejściowy. Kodowanie to jest
jednoznacznie odczytywalne; skończona procedura sprawdza,
czy liczby symboli i stanów są poprawne, czy tabela jest
zupełna i jednoznaczna dla par $(q,a)$ ze stanu niekońcowego
oraz czy przejścia mają dozwolone kierunki.
W ten sposób każdej maszynie przypisujemy pewien napis,
który nazywamy \emph{kodem maszyny}. Przyjmujemy odtąd
jedną ustaloną, efektywną metodę takiego kodowania.
Nie każdy napis musi kodować poprawną maszynę;
niepoprawny opis rozpoznajemy podczas sprawdzania kodu.

W pozostałych rozważaniach o kodach maszyn ograniczamy się do
maszyn o alfabecie wejściowym $\{0,1\}$; wszystkie programy
konstruowane poniżej korzystają wyłącznie z takich wejść.
To wystarcza do dowodu nierozstrzygalności problemu stopu.
Zasadniczo można również kodować słowa nad innymi skończonymi
alfabetami przez słowa binarne, ustalając uprzednio ich
efektywnie odczytywalne kodowanie.

\begin{twierdzenie}[Maszyna uniwersalna]\label{thm:computability-universal}
Istnieje maszyna Turinga $U$ o binarnym alfabecie wejściowym
taka, że dla każdego poprawnego kodu $e$ maszyny $M_e$
o binarnym alfabecie wejściowym i każdego $w\in\{0,1\}^*$
maszyna $U$ uruchomiona na kodzie pary $(e,w)$
symuluje obliczenie $M_e$ na $w$ i zatrzymuje się
z taką samą odpowiedzią wtedy i tylko wtedy, gdy $M_e$
zatrzyma się na $w$.
\end{twierdzenie}
\begin{proof}
Przedstawiamy konstrukcję symulatora i uzasadniamy, dlaczego
wykonuje ją \emph{jedna} maszyna niezależna od $e$.
Najpierw $U$ odkodowuje parę $(e,w)$ według
Lematu~\ref{lem:computability-pair-code} i sprawdza poprawność
kodu $e$. Jeśli kod pary lub maszyny jest błędny, $U$ może
odrzucić wejście; te przypadki nie są objęte tezą twierdzenia.
Dla poprawnego kodu zapisuje jego tabelę przejść
oraz konfigurację początkową $M_e$ na $w$.

Konfigurację przedstawiamy skończonym napisem: obejmuje
stan, numer komórki pod głowicą (z zapisem znaku dla indeksów
ujemnych) i listę par „numer odwiedzonej komórki -- zapisany
symbol”. Indeksy zapisujemy w ustalonym kanonicznym kodzie binarnym
ze znakiem, bez nieistotnych zer wiodących; symbole kodujemy
binarnie, a kolejne pola jednoznacznie rozdzielamy.
Poza tą listą na taśmie symulowanej maszyny stoją symbole
puste. W każdej rundzie $U$ odczytuje symbol spod
symulowanej głowicy (z listy albo jako symbol pusty),
wyszukuje w zakodowanej tabeli odpowiedni wpis
$\delta(q,a)=(q',b,d)$, zmienia zapis wskazanej komórki,
stan i pozycję głowicy. Jeśli głowica wejdzie na nową
komórkę, dodaje do listy jej numer i dotychczasowy symbol pusty.

Każdą z tych czynności wykonuje się przez skończone
przeszukiwanie i przepisywanie zakodowanych napisów.
Są to operacje dostępne dla maszyny Turinga: skanowanie
skończonej listy, porównywanie zakodowanych numerów,
aktualizacja pola i obliczenie numeru sąsiedniej komórki
przez dodanie lub odjęcie jedności. W szczególności
porównanie dwóch skończonych zapisów binarnych wykonujemy
przez porównanie znaków i długości, a następnie odpowiednich
bitów, zaś dodanie lub odjęcie jedności -- przez skończony
przebieg z przeniesieniem albo pożyczką. Usunięcie lub dopisanie
pola do listy wymaga tylko przesunięcia skończonego fragmentu
zapisu na taśmie symulatora. Wszystkie te czynności można
złożyć w skończony zestaw podprogramów, niezależny od
liczby symboli zapisanych w aktualnej konfiguracji.
Kod $e$ jest \emph{danymi},
a procedura realizująca te operacje ma stałą, skończoną
tabelę instrukcji; nie zależy od konkretnej symulowanej maszyny.

Indukcyjnie po liczbie rund zapis $U$ reprezentuje dokładnie
konfigurację $M_e$ po tej samej liczbie kroków: początek
jest poprawny, a opisany odczyt i aktualizacja realizują
jedno przejście $\delta$. Jeśli $M_e$ wejdzie do stanu
$q_{\mathrm{tak}}$ lub $q_{\mathrm{nie}}$, $U$ kończy pracę
z tą samą odpowiedzią. Jeśli $M_e$ nie zatrzyma się, każda
jej kolejna konfiguracja zostanie zasymulowana,
więc także $U$ się nie zatrzyma.
\end{proof}

\begin{uwaga}[Co zostało tutaj pokazane?]
Nie twierdzimy, że każdy program ma krótki kod ani że
maszyna uniwersalna działa szybko. Istotna jest jedna
własność logiczna: istnieje jeden skończony program
potrafiący na podstawie zapisu innego programu naśladować
jego zachowanie. W szczególności programy mogą przyjmować
jako dane \emph{własne kody}.
\end{uwaga}

\section{Problem stopu}\label{sec:computability-halting}
\begin{definicja}[Problem stopu]\label{def:computability-halting}
W ustalonym binarnym kodowaniu definiujemy
$H\subseteq\{0,1\}^*$ jako zbiór słów
$\operatorname{code}(e,w)$, dla których $e$ jest poprawnym
kodem maszyny $M_e$ i maszyna $M_e$ zatrzymuje się
na wejściu $w$, niezależnie od tego, czy akceptuje,
czy odrzuca. Napisy, które nie kodują pary lub mają
niepoprawny kod maszyny, nie należą do $H$.
Pytanie problemu stopu brzmi: mając zakodowaną parę
$(e,w)$, czy możemy zawsze w skończonym czasie ustalić,
czy $M_e$ zatrzyma się na $w$?
\end{definicja}

\begin{twierdzenie}[Nierozstrzygalność problemu stopu]\label{thm:computability-halting-undecidable}
Nie istnieje maszyna Turinga rozstrzygająca zbiór $H$.
\end{twierdzenie}
\begin{proof}
Załóżmy przeciwnie, że istnieje maszyna $D$ rozstrzygająca
$H$: na każdym słowie binarnym zatrzymuje się i poprawnie
odpowiada, czy słowo to należy do $H$.
W szczególności na wejściu $\operatorname{code}(e,w)$,
gdzie $e$ jest poprawnym kodem maszyny, rozstrzyga,
czy $M_e$ zatrzyma się na $w$.

Zbudujmy program $K$ przyjmujący jako dane pojedynczy
napis $x$. Jeśli $x$ nie jest poprawnym kodem maszyny,
program $K$ może od razu zakończyć pracę.
Jeśli $x$ jest poprawnym kodem, program $K$ uruchamia
$D$ na $\operatorname{code}(x,x)$, aby ustalić,
czy maszyna o kodzie $x$ zatrzymuje się na wejściu $x$.
Następnie \emph{postępuje odwrotnie} do otrzymanej odpowiedzi:
\begin{enumerate}[label=(\roman*)]
\item jeśli $D$ przewiduje, że $M_x$ się zatrzyma,
      program $K$ zaczyna nieskończoną pętlę;
\item jeśli $D$ przewiduje, że $M_x$ się nie zatrzyma,
      program $K$ natychmiast kończy pracę.
\end{enumerate}
Każdy z tych kroków da się opisać przez maszynę Turinga:
wykonujemy obliczenie $D$ (które zawsze się kończy),
a potem albo zatrzymujemy się, albo powtarzamy jeden krok bez końca.
Zatem $K$ jest poprawnie opisanym programem i ma własny kod $k$.

Uruchommy teraz $K$ na jego własnym kodzie $k$.
Jeśli $D$ na wejściu $\operatorname{code}(k,k)$ odpowie,
że $M_k=K$ zatrzyma się na $k$, to z definicji
$K$ zapętli się i \emph{nie} zatrzyma.
Jeśli $D$ odpowie, że $K$ \emph{nie} zatrzyma się na $k$,
to $K$ natychmiast się zatrzyma.
W obu przypadkach przewidywanie $D$ jest błędne.
To przeczy założeniu, że $D$ zawsze udziela poprawnej odpowiedzi.
Zatem taki program nie istnieje.
\end{proof}

\begin{figure}[htbp]
\centering
\begin{tikzpicture}[
  >=Stealth,
  font=\small,
  pytanie/.style={draw,rounded corners,align=center,minimum width=38mm,minimum height=9mm},
  przypadek/.style={draw,rounded corners,align=center,text width=40mm,minimum height=16mm},
  line width=.55pt]
\node[pytanie] (d) {$D(\operatorname{code}(k,k))$};
\node[przypadek,below left=11mm and 8mm of d] (tak)
  {$D$ mówi: „zatrzyma się”\\ $K(k)$ nie zatrzymuje się};
\node[przypadek,below right=11mm and 8mm of d] (nie)
  {$D$ mówi: „nie zatrzyma się”\\ $K(k)$ zatrzymuje się};
\draw[->] (d.south west) -- (tak.north east);
\draw[->] (d.south east) -- (nie.north west);
\end{tikzpicture}
\caption{Argument przekątniowy: w obu możliwych przypadkach
  odpowiedź $D$ jest sprzeczna z zachowaniem programu $K$.}
\label{fig:computability-diagonal}
\end{figure}

\begin{uwaga}[Dlaczego to nie jest paradoks językowy?]
W dowodzie nie zakładaliśmy, że jakiś program
ma jednocześnie się zatrzymać i nie zatrzymać.
Przypuściliśmy jedynie istnienie \emph{uniwersalnego
rozstrzygającego} programu $D$. Na tej podstawie
skonstruowaliśmy nowy program $K$, którego zachowanie
przeczyłoby odpowiedzi $D$ dla jednego wybranego wejścia.
Wniosek jest zatem taki, że \emph{nie istnieje} program
$D$ działający poprawnie we wszystkich przypadkach.
\end{uwaga}

\section{Problem stopu jest mimo to rozpoznawalny}
\begin{twierdzenie}[Rozpoznawalność problemu stopu]\label{thm:computability-halting-recognizable}
Zbiór $H$ jest rozpoznawalny, ale nie jest rozstrzygalny.
\end{twierdzenie}
\begin{proof}
Nierozstrzygalność udowodniliśmy w
Twierdzeniu~\ref{thm:computability-halting-undecidable}.
Aby rozpoznać $H$, bierzemy dowolne słowo binarne.
Sprawdzamy, czy jest kodem pary $(e,w)$; jeśli nie,
odrzucamy. Następnie sprawdzamy, czy $e$ opisuje
poprawną maszynę; w przeciwnym razie także odrzucamy.
Dla poprawnego kodu symulujemy działanie $M_e$ na $w$
za pomocą maszyny uniwersalnej z
Twierdzenia~\ref{thm:computability-universal}.
Jeżeli symulowana maszyna zatrzyma się,
rozpoznający akceptuje. Jeżeli się nie zatrzyma,
symulacja trwa w nieskończoność, ale definicja
rozpoznawalności na to pozwala.
\end{proof}

\begin{twierdzenie}[Rozstrzygalność przez dwa rozpoznawania]\label{thm:computability-two-recognizers}
Dla dowolnego $L\subseteq\{0,1\}^*$ zbiór $L$
jest rozstrzygalny wtedy i tylko wtedy, gdy zarówno
$L$, jak i jego dopełnienie $\{0,1\}^*\setminus L$
są rozpoznawalne.
\end{twierdzenie}
\begin{proof}
Jeśli $L$ jest rozstrzygalny, to rozpoznajemy $L$
maszyną rozstrzygającą. Aby rozpoznać dopełnienie,
wystarczy zamienić jej odpowiedzi „tak” i „nie”.

Odwrotnie, niech $R_1$ rozpoznaje $L$, a $R_0$
rozpoznaje jego dopełnienie. Dla zadanego słowa $w$
nie uruchamiamy $R_1$ i nie oczekujemy bez końca,
a dopiero potem $R_0$. Zamiast tego przeplatamy obliczenia:
wykonujemy jeden krok $R_1$, jeden krok $R_0$,
następny krok $R_1$, następny krok $R_0$ itd.
Jeżeli któryś z programów zaakceptuje $w$,
zatrzymujemy się i zwracamy odpowiadającą mu decyzję.
Ewentualne odrzucenie przez jedną z maszyn ignorujemy:
z punktu widzenia rozpoznawalności negatywny wynik
nie jest świadectwem przynależności do jej zbioru.
Zatrzymanej bez akceptacji symulacji nie wznawiamy,
a drugą kontynuujemy.
Każde $w$ należy albo do $L$, albo do jego dopełnienia,
więc przynajmniej jeden rozpoznający musi je po skończonej
liczbie kroków zaakceptować. Opisana procedura
zawsze się zatem zatrzymuje i poprawnie rozstrzyga $L$.
\end{proof}

\begin{wniosek}[Dopełnienie problemu stopu nie jest rozpoznawalne]
Zbiór $\{0,1\}^*\setminus H$ nie jest rozpoznawalny.
\end{wniosek}
\begin{proof}
Gdyby dopełnienie $H$ było rozpoznawalne,
to na mocy Twierdzenia~\ref{thm:computability-two-recognizers}
zbiór $H$ byłby rozstrzygalny, ponieważ sam $H$
jest rozpoznawalny. Przeczy to
Twierdzeniu~\ref{thm:computability-halting-undecidable}.
\end{proof}

\section{Dowody formalne jako dane dla algorytmu}
Zwróćmy uwagę na ważne powiązanie z rozdziałem~\ref{ch:natural-deduction}.
Poprawność pojedynczego kroku dedukcji naturalnej
jest sprawdzalna, ponieważ mamy skończony zapis reguły
oraz jej warunków ubocznych. Przy ustalonym efektywnie
opisanym, przeliczalnym języku i obliczalnym systemie aksjomatów
możemy więc przeglądać kolejno skończone drzewa
i sprawdzać, które są dowodami.

\begin{twierdzenie}[Rozpoznawalność dowodliwości w efektywnym systemie]
\label{thm:computability-proof-recognizable}
Niech $\mathcal L$ będzie językiem pierwszego rzędu,
w którym poprawność kodów symboli, ich rodzaje i arności
można rozstrzygać algorytmicznie. Zakładamy, że termy,
formuły oraz zdania języka $\mathcal L$ mają skończone
kody binarne, których poprawność można sprawdzać
algorytmicznie. Niech $T$ będzie rozstrzygalnym zbiorem zdań
tego języka, a sprawdzanie wszystkich reguł dedukcji naturalnej
(wraz z warunkami ubocznymi) niech będzie algorytmiczne.
Wtedy zbiór kodów zdań $\sigma$ spełniających
$T\vdash_{\mathrm{ND}}\sigma$ jest rozpoznawalny.
\end{twierdzenie}
\begin{proof}
Na wejściu najpierw sprawdzamy, czy otrzymany napis
jest kodem zdania $\sigma$; w przeciwnym razie go odrzucamy.
Następnie wypisujemy kolejno wszystkie skończone słowa
binarne, na przykład według długości i leksykograficznie.
Każdy dowód formalny jest skończonym drzewem o skończonym
kodzie binarnym, więc jego kod pojawi się w tym przeglądzie.
Dla każdego słowa sprawdzamy,
czy koduje poprawne drzewo: czy wszystkie węzły
zawierają poprawne formuły, czy każdy krok używa
jednej z dozwolonych reguł z zachowaniem warunków ubocznych
i czy wszystkie pozostałe otwarte założenia należą do $T$.
Ostatni warunek jest rozstrzygalny z założenia o $T$;
pozostałe kontrole wynikają z jawnych założeń o składni
i regułach.
Jeżeli znajdziemy drzewo z wnioskiem $\sigma$, akceptujemy.
Jeśli takie drzewo istnieje, będzie jednym ze skończonych
napisów w naszym przeglądzie, więc w końcu je znajdziemy.
Jeśli nie istnieje, przegląd może trwać bez końca;
to dokładnie odpowiada definicji rozpoznawalności.
\end{proof}

\begin{przyklad}[Nie mylmy sprawdzenia dowodu ze znalezieniem dowodu]
Gdy ktoś poda nam \emph{konkretne} drzewo dowodu,
możemy krok po kroku skontrolować jego poprawność.
Nie oznacza to jednak, że istnieje algorytm,
który dla dowolnego zdania zawsze rozstrzygnie,
czy \emph{jakieś} poprawne drzewo dowodu istnieje.
Sprawdzanie gotowego certyfikatu i rozstrzyganie
istnienia certyfikatu to dwa różne problemy.
\end{przyklad}


\section{Przejście od maszyn Turinga do rachunku na liczbach}
W rozdziale~\ref{ch:arithmetization} wprowadzimy funkcje
pierwotnie rekurencyjne. Są one szczególną klasą funkcji
obliczalnych, ale nie należy utożsamiać tej klasy ze wszystkimi
funkcjami obliczalnymi przez maszyny Turinga.
W tym miejscu chodzi o mechanizm kodowania: liczby naturalne
mogą reprezentować słowa i konfiguracje maszyny, a pojedynczy
krok obliczenia daje się opisać działaniami arytmetycznymi.

\begin{przyklad}[Dwa fragmenty taśmy jako dwie liczby]
Niech alfabet taśmy ma $B\geqslant 2$ symboli oznaczonych
liczbami $0,\ldots,B-1$, przy czym symbolowi pustemu
przypisujemy numer $0$. Niepusty skończony napis
$a_0a_1\cdots a_k$ przedstawiamy liczbą
$A=a_0+a_1B+\cdots+a_kB^k$ i osobno zapisujemy jego
długość $k+1$. Pusty napis reprezentujemy przez parę $(0,0)$.
Dla napisu niepustego pierwszy symbol jest równy $A\bmod B$,
a kod pozostałej części wynosi $\lfloor A/B\rfloor$.

Zakodujmy tak dwa skończone fragmenty taśmy, leżące po lewej
i po prawej stronie głowicy, w kolejności \emph{od komórek
najbliższych głowicy}. Niech ich kody i długości będą równe
$(L,\ell)$ i $(R,r)$, a $a$ oznacza symbol aktualnie czytany.
Załóżmy, że funkcja przejścia nakazuje
$\delta(q,a)=(q',c,d)$, a więc maszyna \emph{zapisuje symbol $c$}
(w ogólności różny od $a$) i wykonuje ruch $d$.

Jeśli $d=P$, to zapisany symbol $c$ staje się pierwszym
symbolem po lewej stronie nowego położenia głowicy.
Otrzymujemy $L'=c+BL$ i $\ell'=\ell+1$.
Jeśli $r>0$, nowo odczytany symbol jest równy
$a'=R\bmod B$, zaś $R'=\lfloor R/B\rfloor$ i $r'=r-1$.
Jeśli $r=0$, otrzymujemy $a'=0$, $R'=0$ i $r'=0$.

Analogicznie, jeśli $d=L$, to $R'=c+BR$ i $r'=r+1$.
Dla $\ell>0$ otrzymujemy $a'=L\bmod B$,
$L'=\lfloor L/B\rfloor$ i $\ell'=\ell-1$,
a dla $\ell=0$ otrzymujemy $a'=0$, $L'=0$ i $\ell'=0$.
Wreszcie jeśli $d=S$, głowica się nie przesuwa:
$L'=L$, $\ell'=\ell$, $R'=R$, $r'=r$ oraz $a'=c$.
We wszystkich trzech przypadkach nowy stan wynosi $q'$.

Wzory te realizują zapis i ruch głowicy także wtedy,
gdy głowica przechodzi na dotąd nieodwiedzoną komórkę.
Oddzielne długości zapobiegają utracie końcowych zer
w kodach fragmentów taśmy. W ten sposób \emph{cały krok}
maszyny (zapis, przesunięcie i zmiana stanu) zostaje wyrażony
przez elementarne operacje arytmetyczne.
\end{przyklad}

\begin{uwaga}[Znaczenie dla dalszych dowodów]
Weryfikacja \emph{podanego} skończonego śladu maszyny
różni się od wyszukiwania śladu zakończonego zatrzymaniem.
W pierwszym przypadku wystarczy skończona liczba
kontroli sąsiednich konfiguracji. W drugim liczba
prób nie ma z góry ograniczenia. Właśnie ta różnica
umożliwi w dalszej części jednoczesne kodowanie
sprawdzalności dowodu i nierozstrzygalności dowodliwości.
\end{uwaga}

\section{Związek z twierdzeniami Gödla}
Twierdzenie o pełności~\ref{thm:godel-completeness} utożsamia
konsekwencję semantyczną z dowodliwością w logice pierwszego rzędu.
Nie zapewnia jednak algorytmu, który dla każdej formuły
rozstrzyga istnienie dowodu. Systematyczny przegląd dowodów
wykrywa dowód, jeśli istnieje; w pozostałych przypadkach
może nie zakończyć pracy.

W dalszych rozdziałach zakodujemy formuły i dowody liczbami
naturalnymi oraz wykorzystamy to w twierdzeniach Gödla
o niezupełności. Przed włączeniem istniejącego opracowania
sprawdzimy jego dowody i spójność notacji.
\chapter{Arytmetyka formalna i kodowanie składni}\label{ch:arithmetization}

\section{Od maszyn i zdań do rachunku arytmetycznego}
W rozdziale~\ref{ch:computability-intro} zobaczyliśmy, że program
i dowód mają skończony opis. Chcemy wykonać jeszcze jeden krok:
zamienić opis na liczbę naturalną, a następnie wyrazić
\emph{samym językiem arytmetyki} własność
„ta liczba koduje poprawny dowód”.
Wtedy o napisach i obliczeniach będziemy mogli mówić wewnątrz
teorii liczb. Jest to podstawa metody Gödla; samego twierdzenia
o niezupełności nie będziemy jeszcze dowodzić.

\begin{definicja}[Język arytmetyki i numerały]\label{def:arith-language}
Niech $\mathcal L_{\mathrm{ar}}=\{0,S,+,\cdot,<\}$
oznacza język pierwszego rzędu, w którym $0$ jest stałą,
$S$ symbolem funkcji jednoargumentowej, $+$ i $\cdot$
symbolami funkcji dwuargumentowych, zaś $<$ symbolem relacji
dwuargumentowej. W modelu standardowym jego uniwersum to
$\N=\{0,1,2,\ldots\}$, a wymienione znaki mają zwykłe znaczenie.
Dla $n\in\N$ \emph{numerał} $\overline n$ jest termem
domkniętym definiowanym przez
$\overline0:=0$ oraz $\overline{n+1}:=S(\overline n)$.
Zatem $\overline3$ to napis $S(S(S(0)))$, a nie cyfra „3”.
\end{definicja}

\begin{uwaga}[Jak czytać notację metamatematyczną?]\label{rem:arith-meta}
Piszemy $T\vdash\varphi$, jeśli istnieje skończony formalny
dowód $\varphi$ z aksjomatów teorii $T$.
Piszemy $\mathfrak N\models\varphi$, gdy zdanie jest prawdziwe
w zwykłej strukturze liczb naturalnych.
Symbol $\mathfrak N$ oznacza strukturę liczb naturalnych, a $\N$ jej uniwersum.
Teoria $Q$ poniżej to arytmetyka Robinsona; nie należy jej mylić ze zbiorem $\Q$ liczb wymiernych.
Równość $3=3$ jest wypowiedzią metamatematyczną
o liczbach; równość $\overline3=\overline3$ jest natomiast
zdaniem języka formalnego. Oba poziomy należy odróżniać.
\end{uwaga}

\section{Arytmetyka Peana i Robinsona}\label{sec:arytmetyka-peana-robinsona}
Zanim podamy aksjomaty, wyjaśnijmy, po co one są. Chcemy zapisać w języku $\mathcal L_{\mathrm{ar}}$ (Definicji~\ref{def:arith-language}) na tyle dużo prawd o liczbach naturalnych, by teoria potrafiła „liczyć” --- a jednocześnie zrobić to skończoną, mechanicznie sprawdzalną listą reguł. Arytmetyka Robinsona $Q$ to taki możliwie minimalny zestaw: opisuje tylko podstawowe własności zera, następnika, dodawania, mnożenia i porządku.

\begin{definicja}[Arytmetyka Robinsona $Q$]\label{res:4:1}
\emph{Arytmetyką Robinsona} $Q$ nazywamy teorię o następujących ośmiu aksjomatach. Nazwy (Q1)--(Q8) są stałe, niezależne od numeracji wzorów:
\begin{align*}
\forall x\ S(x)\ne0, &\tag{Q1}\\
\forall x\,\forall y\,(S(x)=S(y)\to x=y), &\tag{Q2}\\
\forall x\,(x\ne0\to\exists y\ x=S(y)), &\tag{Q3}\\
\forall x\ x+0=x, &\tag{Q4}\\
\forall x\,\forall y\ x+S(y)=S(x+y), &\tag{Q5}\\
\forall x\ x\cdot0=0, &\tag{Q6}\\
\forall x\,\forall y\ x\cdot S(y)=x\cdot y+x, &\tag{Q7}\\
\forall x\,\forall y\ (x<y\leftrightarrow\exists z\,(z+S(x)=y)). &\tag{Q8}
\end{align*}
Przyjmujemy $x\le y$ jako skrót $x<y\lor x=y$. Kolejność składników w (Q8) ma znaczenie: w obliczeniach z aksjomatów $Q$ nie wolno bez osobnego uzasadnienia przestawiać składników dodawania. W Lemacie~\ref{res:4:4} wyprowadzimy z podanych aksjomatów opis odcinka poniżej każdego ustalonego standardowego numerału.
\end{definicja}

\begin{uwaga}[Intuicyjne znaczenie aksjomatów]\label{res:4:2}
Czytajmy $S(x)$ jako „$x+1$”, a $0$ jako zero. Wtedy:
\begin{itemize}
\item \textup{(Q1)} $S(x)\ne0$: \emph{zero nie jest niczyim następnikiem} --- żadna liczba $+1$ nie daje $0$. To odróżnia $\mathbb N$ od np. liczb modulo, gdzie „zawinięcie” pozwoliłoby wrócić do zera.
\item \textup{(Q2)} $S(x)=S(y)\to x=y$: \emph{następnik jest różnowartościowy} --- jeśli $x+1=y+1$, to $x=y$. Gwarantuje, że kolejne liczby się nie „sklejają”.
\item \textup{(Q3)} $x\ne0\to\exists y\ x=S(y)$: \emph{każda niezerowa liczba ma poprzednik} --- nie ma „dziur” ani dziwnych elementów bez $-1$.
\item \textup{(Q4)} $x+0=x$: dodanie zera nic nie zmienia (baza rekurencji dodawania).
\item \textup{(Q5)} $x+S(y)=S(x+y)$: \emph{definiuje dodawanie przez rekurencję} --- dodać $y+1$ to dodać $y$ i potem $+1$. Aksjomaty (Q4)--(Q5) wspólnie mówią, jak liczyć $x+y$ krok po kroku.
\item \textup{(Q6)} $x\cdot0=0$: mnożenie przez zero daje zero (baza).
\item \textup{(Q7)} $x\cdot S(y)=x\cdot y+x$: \emph{definiuje mnożenie przez rekurencję} --- pomnożyć przez $y+1$ to pomnożyć przez $y$ i dodać $x$.
\item \textup{(Q8)} $x<y\leftrightarrow\exists z\,(z+S(x)=y)$: w modelu standardowym mówi, że do $x+1$ można dodać nieujemną liczbę i otrzymać $y$. Nie zakłada ogólnych praw porządku, których $Q$ mogłaby nie dowodzić.
\end{itemize}
\end{uwaga}

\begin{przyklad}[$Q$ liczy konkretne przykłady]\label{res:4:3}
Z aksjomatów (Q4)--(Q5) $Q$ dowodzi np. $\overline2+\overline1=\overline3$: istotnie $\overline2+S(0)=S(\overline2+0)=S(\overline2)=\overline3$. Ogólnie $Q$ poprawnie dowodzi każdej prawdziwej równości i nierówności między \emph{konkretnymi} numerałami. To wystarczy do reprezentowalności (Tw.~\ref{res:7:6}).
\end{przyklad}

\begin{lemat}[Skończone odcinki i rachunek na numerałach]\label{res:4:4}
Dla każdego ustalonego standardowego $n$ teoria $Q$ dowodzi
\begin{align}
\forall x\,(x<\overline n\leftrightarrow(x=\overline0\lor\cdots\lor x=\overline{n-1})),\label{eq:4:1}\\
\forall x\,(x<\overline n\lor x=\overline n\lor\overline n<x).\label{eq:4:2}
\end{align}
Dla $n=0$ alternatywa po prawej w~\eqref{eq:4:1} oznacza fałsz. Ponadto $Q$ rozstrzyga równości i nierówności między termami domkniętymi.
\end{lemat}

\begin{proof}
Z (Q4)--(Q7) obliczamy każdy term domknięty, rozwijając dodawanie i mnożenie skończenie wiele razy, zawsze względem drugiego argumentu. Dla dwóch różnych numerałów $\overline i,\overline j$ z $i<j$, równość $\overline i=\overline j$ po $i$ zastosowaniach (Q2) prowadzi do $0=S^{j-i}(0)$, wbrew (Q1). Przez podstawianie równych termów rozstrzygamy więc dowolną równość między termami domkniętymi.

Z (Q3) i prawa wyłączonego środka otrzymujemy $x=0\lor\exists u\ x=S(u)$. W drugiej gałęzi stosujemy ten sam podział do $u$: albo $u=0$ i $x=\overline1$, albo $u=S(v)$ i $x=S^2(v)$. Powtarzamy to dokładnie $n+1$ razy, otrzymując w $Q$
\[
x=\overline0\lor x=\overline1\lor\cdots\lor x=\overline n\lor\exists u\ x=S^{n+1}(u).
\]
Zapis $S^k(u)$ oznacza $k$ kolejnych zastosowań następnika. Jeśli $x=S^{n+1}(u)$, to $z+S(x)=S^{n+2}(z+u)\ne\overline n$ na mocy (Q1)--(Q2) i (Q5): domniemana równość po $n$ skróceniach następnika dawałaby $S^2(z+u)=0$. Zatem $x\not<\overline n$. Dla $x=\overline j$ z $j\le n$ zachodzi $z+S(\overline j)=S^{j+1}(z)$. Przy $j<n$ świadkiem równości z $\overline n$ jest $z=\overline{n-j-1}$. Przy $j=n$ taka równość po $n$ skróceniach dawałaby $S(z)=0$, więc jest niemożliwa. To dowodzi obu kierunków~\eqref{eq:4:1}. W ostatnim członie alternatywy $u+S(\overline n)=S^{n+1}(u)=x$, więc $\overline n<x$; pozostałe gałęzie dają~\eqref{eq:4:2}. Jest to indukcja w metateorii budująca osobny skończony dowód dla każdego $n$, a nie zastosowanie nieistniejącego w $Q$ schematu indukcji.
\end{proof}

\begin{uwaga}[Jak połączyć skończenie wiele sprawdzeń]\label{res:4:5}
Niech $n$ będzie ustaloną liczbą standardową, a $\theta(x)$ dowolną formułą. Jeśli $Q$ dowodzi wszystkich $\theta(\overline0),\ldots,\theta(\overline{n-1})$, to dowodzi też $\forall x<\overline n\,\theta(x)$. Ten ostatni zapis jest skrótem $\forall x(x<\overline n\to\theta(x))$. Istotnie, z $x=\overline j$ oraz $\theta(\overline j)$ otrzymujemy $\theta(x)$ przez podstawienie równych termów. Łączymy skończenie wiele takich implikacji przez alternatywę i stosujemy~\eqref{eq:4:1}. Następnie generalizujemy $x$. Dla $n=0$ teza wynika z $\forall x\,\neg(x<0)$. Tak powstaje \emph{jeden skończony dowód dla ustalonej granicy} $n$, a nie dowód $\forall x\,\theta(x)$. Tego kroku użyjemy między innymi w dowodzie Rossera.
\end{uwaga}

\begin{definicja}[Formuły ograniczone i $\Sigma_1$]\label{res:4:6}
Formuła $\Delta_0$ powstaje z formuł atomowych przez spójniki oraz kwantyfikatory $\exists x<t$, $\forall x<t$, gdzie $t$ nie zawiera $x$. Są to skróty $\exists x(x<t\land\cdots)$ i $\forall x(x<t\to\cdots)$. Kwantyfikatory z $x\le t$ również są ograniczone: oznaczają odpowiednio $(\exists x<t\,\varphi(x))\lor\varphi(t)$ oraz $(\forall x<t\,\varphi(x))\land\varphi(t)$, po podstawieniu bez przechwytywania zmiennych. Nie potrzeba tu żadnych ogólnych praw porządku. Formuła $\Sigma_1$ ma postać $\exists u_1\cdots\exists u_r\,\delta$, gdzie $\delta$ jest $\Delta_0$; dopuszczamy $r=0$.
\end{definicja}

\begin{lemat}[$\Sigma_1$-zupełność $Q$]\label{res:4:7}
$Q$ dowodzi każdego prawdziwego zdania $\Sigma_1$ oraz negacji każdego fałszywego zdania $\Delta_0$.
\end{lemat}
\begin{proof}
Dla każdego termu domkniętego $t$ Lemat~\ref{res:4:4} daje numerał $\overline n$ taki, że $Q\vdash t=\overline n$. Indukcję po budowie zamkniętej formuły $\Delta_0$ rozpoczynamy od atomów (równość i $<$): $Q$ dowodzi ich wartości logicznej dla obliczonych numerałów. Spójniki rozpatrujemy klasycznie. W przypadku $\exists x<t\,\psi(x)$, po ustaleniu $Q\vdash t=\overline n$, równoważność~\eqref{eq:4:1} i podstawianie równych termów sprowadzają kwantyfikator do skończonej alternatywy $\psi(\overline0)\lor\cdots\lor\psi(\overline{n-1})$; analogicznie $\forall x<t$ sprowadza się do skończonej koniunkcji. Stosujemy założenie indukcyjne do każdej zamkniętej instancji $\psi(\overline j)$, również gdy w jej wnętrzu pozostają dalsze kwantyfikatory ograniczone. W ten sposób $Q$ dowodzi prawdziwej zamkniętej formuły $\Delta_0$ i negacji fałszywej. Jeśli zdanie $\exists\vec u\,\delta(\vec u)$ klasy $\Sigma_1$ jest prawdziwe, wybieramy standardowe wartości świadków $\vec m$, dowodzimy w $Q$ zdania $\delta(\overline{m_1},\ldots,\overline{m_r})$, a następnie wprowadzamy kolejno kwantyfikatory egzystencjalne. Przy $r=0$ korzystamy z przypadku $\Delta_0$. Nie wynika stąd rozstrzygalność fałszywych zdań $\Sigma_1$.
\end{proof}

\begin{uwaga}[Czego nie wolno zakładać w $Q$]\label{res:4:8}
Teoria $Q$ nie zawiera schematu indukcji. Z samych równań
$x+0=x$ i $x+S(y)=S(x+y)$ możemy wyliczyć $x+\overline n$
dla \emph{ustalonego numerału} $\overline n$, lecz nie uzyskujemy
przez to automatycznie dowodu $\forall y\,(0+y=y)$ ani
przemienności dodawania. Aby otrzymać twierdzenie o \emph{każdym}
$y$, trzeba podać osobny argument; nie należy go zastępować
nieuprawnionym użyciem indukcji. Brak aksjomatu nie jest jednak
sam w sobie dowodem niezależności takiego twierdzenia od $Q$.
W dalszej części konsekwentnie używamy jedynie rachunku na
\emph{konkretnych} numerałach i skończonych odcinkach.
\end{uwaga}

\begin{definicja}[Arytmetyka Peana PA]\label{res:4:9}
\emph{Arytmetyką Peana} $\mathrm{PA}$ nazywamy teorię w języku $\mathcal L_{\mathrm{ar}}$ zawierającą aksjomaty (Q1)--(Q8) oraz \emph{schemat indukcji}. Aksjomat (Q8) można równoważnie traktować jako definicję symbolu $<$ przez pozostałe symbole. Dla każdej formuły $\varphi(x,\vec y)$, gdzie $\vec y=(y_1,\ldots,y_m)$ jest skończoną listą różnych zmiennych innych niż $x$, a wszystkie zmienne wolne formuły należą do $\{x,y_1,\ldots,y_m\}$, przyjmujemy zdanie (zapis $\forall\vec y$ skraca kolejne kwantyfikatory po $y_1,\ldots,y_m$)
\[
\forall\vec y\,\Bigl[\bigl(\varphi(0,\vec y)\land\forall x\,(\varphi(x,\vec y)\to\varphi(S(x),\vec y))\bigr)\to\forall x\,\varphi(x,\vec y)\Bigr].
\]
\end{definicja}

\begin{uwaga}[Co mówi schemat indukcji]\label{res:4:10}
To formalna wersja „zasady domina”: jeśli własność $\varphi$ zachodzi dla $0$ (\emph{baza}) i z jej zachodzenia dla $x$ wynika zachodzenie dla $x+1$ (\emph{krok}), to zachodzi dla \emph{wszystkich} liczb. Słowo „schemat” oznacza, że jest to \emph{nieskończenie wiele} aksjomatów --- po jednym dla każdej formuły $\varphi$. Mimo nieskończoności zbiór ten jest \emph{rozstrzygalny} (można mechanicznie sprawdzić, czy dane zdanie ma kształt aksjomatu indukcji), więc PA pozostaje efektywnie aksjomatyzowalna --- co jest niezbędne dla twierdzeń Gödla.
\end{uwaga}

\begin{uwaga}\label{res:4:11}
$Q$ jest bardzo słaba: nie dowodzi nawet przemienności dodawania. Aby tę ostatnią niedowodliwość uzasadnić, skonstruujmy model. Do $\N$ dołączmy dwa różne elementy $a,b$, niech $S(a)=a$, $S(b)=b$, a na $\N$ następnik pozostanie standardowy. Dodawanie na $\N$ jest zwykłe; dla $c\in\{a,b\}$ i $n\in\N$ ustalmy $c+n=c$, natomiast dla dowolnego $z$ i $c\in\{a,b\}$ połóżmy $z+c=c$. Dla mnożenia zachowujemy zwykłe działanie na $\N$, przyjmujemy $c\cdot0=0$, $c\cdot n=c$ dla $n\ge1$, a gdy $d\in\{a,b\}$, ustalamy $n\cdot d=d$ dla $n\in\N$ i $c\cdot d=c$ dla $c\in\{a,b\}$. Relację $<$ interpretujemy dokładnie przez prawą stronę (Q8). Sprawdźmy zwłaszcza kroki rekurencyjne. Dla $y=n\in\N$ równania (Q5) i (Q7) są standardowe, gdy $x\in\N$; dla $x=c\in\{a,b\}$ mamy $c+S(n)=c=S(c+n)$, a $c\cdot S(0)=c=0+c$ oraz, dla $n\ge1$, $c\cdot S(n)=c=c+c$. Dla $y=d\in\{a,b\}$ zachodzi $S(d)=d$: po obu stronach (Q5) otrzymujemy $d$, natomiast w (Q7) otrzymujemy $d$ przy $x\in\N$ i $x$ przy $x\in\{a,b\}$. Równania (Q4) i (Q6) są spełnione na mocy definicji działań; (Q1)--(Q3) wynikają z opisanej funkcji $S$, a (Q8) zachodzi z definicji $<$. Mamy jednak $a+b=b\ne a=b+a$. Poprawność rachunku dowodowego wyklucza więc dowód ogólnej przemienności z $Q$. Mimo tej słabości $Q$ wystarcza do reprezentowania funkcji obliczalnych na ustalonych numerałach (Tw.~\ref{res:7:6}). Pierwsze twierdzenie Gödla w wersji Rossera obejmuje każde niesprzeczne efektywnie aksjomatyzowalne rozszerzenie $Q$. Drugie twierdzenie dowodzimy dalej dla rozszerzeń PA; sama reprezentowalność nie zastępuje warunków wyprowadzalności.
\end{uwaga}

% Rozdziały 5--6, strony 9--13.
\section{Funkcje i zbiory rekurencyjne}\label{sec:funkcje-rekurencyjne}
\begin{definicja}[Funkcje pierwotnie rekurencyjne]\label{res:5:1}
Przyjmujemy, że $\N^0$ jest zbiorem jednoelementowym, więc funkcja $\N^0\to\N$ oznacza stałą zeroargumentową. Klasa \emph{funkcji pierwotnie rekurencyjnych} (p.r.) to najmniejsza klasa funkcji totalnych $f:\mathbb N^k\to\mathbb N$ dla wszystkich arności $k\ge0$, zawierająca:
\begin{enumerate}[label=(\roman*)]
\item funkcje zerowe $Z^k(x_1,\ldots,x_k)=0$ dla każdego $k\ge0$ (w tym stałą $Z^0=0$);
\item następnik $S(x)=x+1$;
\item rzutowania $P_i^k(x_1,\ldots,x_k)=x_i$ dla $k\ge1$ i $1\le i\le k$;
\end{enumerate}
i zamknięta ze względu na:
\begin{enumerate}[label=(\alph*)]
\item \emph{złożenie}: jeśli $g:\mathbb N^m\to\mathbb N$ oraz $h_1,\ldots,h_m:\mathbb N^k\to\mathbb N$ są p.r. ($k,m\ge0$), to $f(\vec x)=g(h_1(\vec x),\ldots,h_m(\vec x))$ jest p.r.; dla $m=0$ wynik jest stały;
\item \emph{rekursję prostą}: jeśli $g:\N^k\to\N$ i $h:\N^{k+2}\to\N$ są p.r. ($k\ge0$), to funkcja $f:\N^{k+1}\to\N$ określona przez
\[
f(\vec x,0)=g(\vec x),\qquad f(\vec x,S(y))=h(\vec x,y,f(\vec x,y))
\]
jest p.r.
\end{enumerate}
Przypadek $k=0$ nie jest pomijany: pozwala stosować rekursję prostą do funkcji jednej zmiennej bez dodatkowych parametrów, np. do poprzednika. Wszystkie wartości pozostają liczbami naturalnymi.
\end{definicja}

\begin{definicja}[Funkcja charakterystyczna zbioru]\label{res:5:2}
\emph{Funkcją charakterystyczną} zbioru (relacji) $R\subseteq\mathbb N^k$ nazywamy funkcję $\chi_R:\mathbb N^k\to\mathbb N$ zdefiniowaną wzorem
\[
\chi_R(\vec x)=\begin{cases}1,&\vec x\in R\quad(\text{„warunek zachodzi”}),\\0,&\vec x\notin R\quad(\text{„warunek nie zachodzi”}).\end{cases}
\]
Innymi słowy $\chi_R$ zamienia \emph{pytanie} „czy $\vec x$ należy do $R$?” na \emph{liczbę} ($1$ = tak, $0$ = nie). Dzięki temu „rozstrzyganie”, czy coś zachodzi, staje się po prostu obliczaniem wartości funkcji --- a o funkcjach umiemy już mówić, że są (lub nie są) rekurencyjne.
\end{definicja}

\begin{uwaga}[Po co funkcja charakterystyczna?]\label{res:5:3}
Cała teoria obliczalności mówi wprost tylko o funkcjach $\mathbb N^k\to\mathbb N$. Zbiory i relacje (np. „$x$ jest liczbą pierwszą”, „$p$ jest dowodem formuły $f$”) to obiekty innego rodzaju --- to warunki prawda/fałsz. Funkcja charakterystyczna jest mostem: pozwala przełożyć „rozstrzygalność zbioru” na „obliczalność funkcji”. Mówimy, że zbiór jest \emph{rozstrzygalny z definicji} wtedy i tylko wtedy, gdy istnieje mechaniczna procedura licząca jego $\chi_R$.
\end{uwaga}

\begin{przyklad}[Funkcje charakterystyczne konkretnych zbiorów]\label{res:5:4}
Przyjmijmy $\operatorname{sgn}(0)=0$ i $\operatorname{sgn}(S(n))=1$ dla $n\in\N$.
Funkcja $\operatorname{sgn}$ jest p.r. przez rekursję prostą; w poniższych wzorach
jej wartością jest zawsze $0$ albo $1$.
\begin{itemize}
\item Zbiór liczb parzystych: $\chi_{\mathrm{parz}}(x)=1$ dla $x\in\{0,2,4,\ldots\}$, w przeciwnym razie $0$. Jest pierwotnie rekurencyjna, bo $\chi_{\mathrm{parz}}(x)=1-(x\bmod 2)$, a reszta z dzielenia jest p.r. (Lemat~\ref{res:5:8}).
\item Relacja równości $E=\{(x,y):x=y\}$: $\chi_E(x,y)=1-\operatorname{sgn}((x\dotminus y)+(y\dotminus x))$; wynik jest $1$ dokładnie wtedy, gdy obie obcięte różnice są zerowe, czyli gdy $x=y$.
\item Relacja $\le$: $\chi_{\le}(x,y)=1-\operatorname{sgn}(x\dotminus y)$, bo $x\dotminus y=0$ dokładnie wtedy, gdy $x\le y$.
\item Zbiór liczb pierwszych: $\chi_{\mathrm{pr}}(x)=1$, gdy $x>1$ i żadna liczba $2\le d<x$ nie dzieli $x$; warunek wyrażalny ograniczonym kwantyfikatorem $\forall d<x$, więc też p.r.
\end{itemize}
Widać wspólny wzorzec: „tak/nie” kodujemy jako $1/0$, a logiczne spójniki odwzorowujemy na operacje arytmetyczne ($\land$ odpowiada mnożeniu, $\neg$ odpowiada $1-(\cdot)$).
\end{przyklad}

\begin{definicja}[Funkcje częściowo rekurencyjne i rekurencyjne]\label{res:5:5}
Funkcja częściowa $f:\mathbb N^k\rightharpoonup\mathbb N$ może nie mieć wartości dla niektórych argumentów. Piszemy $f(\vec n)\downarrow=m$, gdy wartość istnieje i wynosi $m$, oraz $f(\vec n)\uparrow$, gdy nie istnieje.

Klasa funkcji \emph{częściowo rekurencyjnych} powstaje z funkcji bazowych z Definicji~\ref{res:5:1} przez złożenie, rekursję prostą i \emph{minimalizację}:
\[
(\mu g)(\vec x)\downarrow=y\quad\Longleftrightarrow\quad
g(\vec x,y)\downarrow=0\ \text{oraz}\ (\forall z<y)\ g(\vec x,z)\downarrow\ne0.
\]
Przy złożeniu i rekursji wynik istnieje tylko wtedy, gdy istnieją wszystkie wartości potrzebne do jego obliczenia. Funkcję częściowo rekurencyjną określoną dla wszystkich argumentów nazywamy \emph{rekurencyjną} (\emph{totalną obliczalną}).

Zbiór $R\subseteq\mathbb N^k$ jest \emph{rekurencyjny} (\emph{rozstrzygalny}), gdy $\chi_R$ jest funkcją totalną rekurencyjną. Jest \emph{rekurencyjnie przeliczalny} (\emph{efektywnie wyliczalny}), gdy stanowi dziedzinę funkcji częściowo rekurencyjnej.
\end{definicja}

\begin{uwaga}[Czym $\mu g$ różni się od $\min\{\cdots\}$]\label{res:5:6}
Jeśli $\mu y\,[g(\vec x,y)=0]$ jest określone, jego wartość równa się $\min\{y\in\mathbb N:g(\vec x,y)=0\}$ --- to faktycznie najmniejszy argument zerujący $g$. Dla częściowego $g$ samo istnienie tego minimum nie gwarantuje jednak określoności $\mu g$. Znaczenie ma również sposób wyznaczania wyniku:
\begin{itemize}
\item Zwykłe $\min\{\cdots\}$ to statyczny zapis teoriomnogościowy: „weź najmniejszy element tego zbioru”. Milczy o tym, jak ów element znaleźć, i ma sens także wtedy, gdy zbioru nie da się efektywnie przeszukać.
\item Operator $\mu$ to przepis obliczeniowy (procedura): $y\gets0$; dopóki $g(\vec x,y)\ne0$, wykonuj $y\gets y+1$; zwróć $y$. Sprawdzamy kolejno $y=0,1,2,\ldots$ i zatrzymujemy się na pierwszym trafieniu. Dlatego $\mu$ zachowuje obliczalność.
\end{itemize}
\emph{Totalność a ograniczoność przeszukiwania.} Gdy $g$ jest totalna, pętla kończy się na danym wejściu dokładnie wtedy, gdy istnieje zero. Gdy $g$ jest częściowa, może utknąć już przy obliczaniu wcześniejszej wartości $g(\vec x,z)$, nawet jeśli późniejsze $g(\vec x,y)$ wynosi zero. Na przykład przy $g(\vec x,0)\uparrow$ i $g(\vec x,1)=0$ zbiór zer ma minimum $1$, ale procedura $\mu g$ utknie już na argumencie $0$.

Jeśli test jest totalny i zero istnieje na każdym wejściu, otrzymana funkcja jest totalna obliczalna. Samo zakończenie każdego obliczenia nie gwarantuje pierwotnej rekurencyjności. Jeśli test i z góry znane ograniczenie przeszukiwania są p.r., ograniczona minimalizacja również jest p.r. (Lemat~\ref{res:5:8}). Nieograniczona minimalizacja pozwala wyjść poza klasę p.r., także wtedy, gdy wynikowa funkcja jest totalna. Brak zakończenia nie jest jedyną różnicą między tymi klasami.

\emph{Umowa notacyjna:} gdy piszemy $\mu y\,[W(y)]$ dla warunku $W$, rozumiemy przez to $\mu y\,[\chi_{\neg W}(y)=0]$, tj. najmniejsze $y$ spełniające $W$ (funkcja charakterystyczna, Def.~\ref{res:5:2}, tłumaczy „warunek” na „zerowanie funkcji”).
\end{uwaga}

\begin{przyklad}[Operator $\mu$ w działaniu]\label{res:5:7}
Funkcja „część całkowita pierwiastka” $\lfloor\sqrt x\rfloor$ to $\mu y\,[(y+1)^2>x]$, czyli najmniejsze $y$, dla którego $(y+1)^2$ przekracza $x$. Dla $x=10$ dostajemy $y=3$ (bo $4^2=16>10$, a $3^2=9\le10$). Ten przykład pokazuje typową rolę $\mu$: „szukaj najmniejszego świadka spełniającego warunek”. Tutaj można ograniczyć przeszukiwanie przez $y\le x$, bo $(x+1)^2>x$. Test i ograniczenie są p.r., więc funkcja jest p.r. Nie wystarczyłoby samo stwierdzenie, że dla każdego wejścia przeszukiwanie kiedyś się kończy.
\end{przyklad}

\begin{lemat}[Zamkniętość p.r.]\label{res:5:8}
Klasa funkcji pierwotnie rekurencyjnych zawiera dodawanie, mnożenie, potęgowanie, funkcję poprzednika, obciętą różnicę $\dotminus$, oraz jest zamknięta ze względu na definicje przez przypadki z warunkami p.r., ograniczone kwantyfikatory oraz ograniczoną minimalizację relacji p.r.
\end{lemat}
\begin{proof}
Dodawanie definiują $x+0=x$, $x+S(y)=S(x+y)$; mnożenie --- $x\cdot0=0$, $x\cdot S(y)=x\cdot y+x$; potęgowanie --- $x^0=1$, $x^{S(y)}=x^y\cdot x$. Są to kolejne zastosowania rekursji prostej. Poprzednik jest przykładem rekursji prostej z $k=0$: $\operatorname{pd}(0)=Z^0=0$ i $\operatorname{pd}(S(y))=P_1^2(y,\operatorname{pd}(y))=y$. Stała $1$ jest złożeniem $S\circ Z^0$. Obcięta różnica $x\dotminus0=x$, $x\dotminus S(y)=\operatorname{pd}(x\dotminus y)$.

Niech $r(\vec x,y)\in\{0,1\}$ będzie p.r. funkcją charakterystyczną relacji $R$. Ograniczone sumy i iloczyny są p.r. przez rekursję po liczbie składników. Stąd
\[
\chi_{\forall y<n\,R(\vec x,y)}=\prod_{y<n}r(\vec x,y),\qquad
\chi_{\exists y<n\,R(\vec x,y)}=1-\prod_{y<n}(1-r(\vec x,y)).
\]
Iloczyn pusty wynosi $1$. Definicję przez dwa przypadki zapisujemy jako $\chi_R f+(1-\chi_R)g$. Najmniejsze $y\le n$ spełniające $R$, z wartością zastępczą $n+1$ przy braku świadka, jest równe
\[
\sum_{j=0}^{n}\prod_{i=0}^{j}(1-r(\vec x,i)).
\]
Jeśli pierwszym trafieniem jest $y$, dokładnie pierwszych $y$ iloczynów wynosi $1$; jeśli trafienia nie ma, wszystkie $n+1$ iloczynów wynoszą $1$. To daje ograniczoną minimalizację bez użycia nieograniczonego $\mu$.

Dla $d>0$ relacja $d\mid x$ ma postać $\exists q\le x\,(dq=x)$, a iloraz i resztę znajdujemy przez ograniczone przeszukiwanie; dla $d=0$ ustalamy wartość zastępczą $0$. Własność bycia liczbą pierwszą jest zatem p.r. Również $i$-ta liczba pierwsza jest p.r. dla $i\ge1$: zaczynamy od $p_1=2$, a $p_{i+1}$ wyznaczamy jako najmniejszą liczbę pierwszą w przedziale od $p_i+1$ do $p_i!+1$. Ten przedział zawiera liczbę pierwszą, ponieważ każdy dzielnik pierwszy $p_i!+1$ jest większy od $p_i$ i nie przekracza $p_i!+1$. Test pierwszości i ograniczone przeszukiwanie są p.r., a silnia jest p.r. jako ograniczony iloczyn. Zatem rekursja wyznaczająca $p_{i+1}$ jest pierwotna rekurencyjna.
\end{proof}

\section{Numeracja Gödla}\label{sec:numeracja-godla}
\begin{definicja}[Kod Gödla]\label{res:6:1}
Ustalamy skończony alfabet symboli innych niż zmienne (dla języka $\mathcal L_{\mathrm{ar}}$) i efektywne, różnowartościowe kodowanie: tym symbolom przypisujemy kody $1,\ldots,b-1$, gdzie $b\ge2$, a zmiennej $v_i$ kod $b+i$ dla $i\in\N$. Dzięki temu test poprawności kodu pojedynczego symbolu, rozpoznawanie jego rodzaju i odzyskiwanie indeksu zmiennej są p.r. Ciągowi symboli $s_1\ldots s_k$ przypisujemy liczbę
\[
\ulcorner s_1\ldots s_k\urcorner=\prod_{i=1}^k p_i^{\operatorname{kod}(s_i)+1},
\]
gdzie $p_i$ jest $i$-tą liczbą pierwszą. Kod pustego napisu wynosi $1$, a $0$ nie jest kodem napisu. Poprawny kod ma dodatnie wykładniki dokładnie przy początkowym odcinku liczb pierwszych, bez luk. Dowolny ciąg liczb kodujemy przez
\[
\langle a_1,\ldots,a_k\rangle=\prod_{i=1}^k p_i^{a_i+1}.
\]
Przesunięcie o $1$ pozwala kodować także zera. Te konstrukcje iterujemy dla ciągów wierszy dowodu i konfiguracji obliczenia. Operacje na niepoprawnych kodach mają ustalone wartości zastępcze, np. $0$, dzięki czemu są funkcjami totalnymi.
\end{definicja}

\begin{definicja}[Skończone obliczenie rejestrowe]\label{res:6:2}
Program używa ustalonej skończonej liczby $d\ge1$ rejestrów $r_0,\ldots,r_{d-1}$ z liczbami naturalnymi i niepustej, skończonej listy numerowanych instrukcji. Wskazujemy instrukcję początkową; wszystkie numery rejestrów i adresy skoków muszą należeć do odpowiednich zakresów. Dopuszczamy:
\begin{itemize}
\item $\operatorname{Inc}(i,j)$: zwiększ $r_i$ o $1$ i przejdź do instrukcji $j$;
\item $\operatorname{Dec}(i,j,k)$: jeśli $r_i=0$, przejdź do $j$; w przeciwnym razie zmniejsz $r_i$ o $1$ i przejdź do $k$;
\item $\operatorname{Halt}$: zakończ, zwracając zawartość $r_0$.
\end{itemize}
Konfiguracja to numer bieżącej instrukcji i zawartości wszystkich $d$ rejestrów. Po zatrzymaniu umownie konfiguracja pozostaje niezmieniona. Na wejściu $\vec x\in\N^k$, gdzie $k\le d$, rejestry $r_0,\ldots,r_{k-1}$ zawierają odpowiednie argumenty, a pozostałe zawierają $0$. Ślad obliczenia to skończony ciąg konfiguracji zaczynający się od tej konfiguracji początkowej, w którym każda następna wynika z poprzedniej zgodnie z programem. Ślad zakończonego obliczenia kończy się w stanie $\operatorname{Halt}$. Zapis listy instrukcji i konfiguracji kodujemy liczbami naturalnymi jak w Definicji~\ref{res:6:1}.
\end{definicja}

\begin{lemat}[Programy i funkcje częściowo rekurencyjne]\label{res:6:3}
Funkcje obliczane przez powyższe programy są dokładnie funkcjami częściowo rekurencyjnymi. Weryfikacja podanego skończonego śladu jest p.r.
\end{lemat}
\begin{proof}
\emph{Od definicji rekurencyjnej do programu.} Wyzerowanie rejestru wykonujemy przez pętlę zmniejszania. Zawartość rejestru kopiujemy, przenosząc ją do dwóch rejestrów pomocniczych i odtwarzając oryginał z jednego z nich. Daje to funkcje bazowe. Złożenie wykonujemy przez kolejne wywołania programów dla argumentów, z rozłącznymi rejestrami roboczymi, a następnie programu funkcji zewnętrznej. Przy rekursji prostej najpierw obliczamy $g(\vec x)$, potem wykonujemy pętlę z licznikiem $i=0,\ldots,y-1$, zastępując bieżący wynik $z$ przez $h(\vec x,i,z)$. W minimalizacji wywołujemy kolejno program $g(\vec x,0),g(\vec x,1),\ldots$ i kończymy na pierwszym zerze. Jeśli wywołanie nie kończy się, nie przechodzimy do następnego argumentu. To dokładnie konwencja z Definicji~\ref{res:5:5}.

\emph{Od programu do funkcji rekurencyjnej.} Jednokrokowe przejście zakodowanej konfiguracji jest p.r.: najpierw odczytujemy numer instrukcji i zawartość wskazanego rejestru, wykonujemy jeden ze skończenie wielu rodzajów operacji, a następnie ponownie kodujemy konfigurację. Dla ustalonego programu jest to złożenie funkcji p.r.; wersję jednostajną, w której kod programu jest dodatkowym argumentem, otrzymujemy przez odczyt instrukcji z jej kodu. Konfigurację po $t$ krokach definiujemy rekursją prostą po $t$. Minimalizacja wybiera pierwszy moment zatrzymania, a odczyt $r_0$ zwraca wynik. Gdy zatrzymania nie ma, funkcja pozostaje nieokreślona. Weryfikacja podanego śladu sprawdza kod programu, konfigurację początkową, każdą parę kolejnych konfiguracji i końcowy stan; liczba sprawdzeń jest ograniczona długością zakodowanego śladu. Uzasadnijmy p.r. charakter kodowania niezależnie od późniejszego Lematu~\ref{res:6:5}: dla kodu $n>0$ wykładnik $p_i$ w $n$ znajdujemy przez ograniczone sprawdzenie podzielności $p_i^k\mid n$ dla $k\le n$; liczbę składników ciągu i poprawność kodu również rozpoznajemy przez ograniczone przeszukiwanie $i\le n$ i odtworzenie iloczynu. Liczby pierwsze $p_i$, ograniczone iloczyny i test podzielności są p.r. na mocy Lematu~\ref{res:5:8}. Zmiana jednego rejestru, wybór następnej instrukcji i zapis zmienionej krotki sprowadzają się zatem do skończonych złożeń tych funkcji. Sprawdzanie kolejnych pozycji śladu jest ograniczonym iloczynem wartości odpowiedniego testu. Dzięki temu nie korzystamy z Lematu~\ref{res:6:5} jako przesłanki dowodu równoważności programów i funkcji częściowo rekurencyjnych; ten późniejszy lemat rozwija niezależnie kontrolę składni i gotowych dowodów. Nie poszukujemy tutaj czasu zatrzymania.
\end{proof}

\begin{uwaga}[Wyliczanie a półrozstrzyganie]\label{res:6:4}
Jeżeli $A$ jest dziedziną programu, symulujemy etapami program na wszystkich wejściach: w etapie $s$ wykonujemy $s$ kroków na wejściach $0,\ldots,s$ i wypisujemy te, na których nastąpiło zatrzymanie. Każdy element $A$ pojawi się po skończonym czasie. Odwrotnie, program na wejściu $n$ może czekać, aż wyliczenie wypisze $n$. Zatrzyma się dokładnie dla $n\in A$. Dla krotek używamy kodowania ciągów. To uzasadnia równoważność dwóch opisów zbioru rekurencyjnie przeliczalnego, także gdy zbiór jest pusty.
\end{uwaga}

\begin{lemat}[Rekurencyjność składni]\label{res:6:5}
Następujące zbiory i funkcje są pierwotnie rekurencyjne:
\begin{enumerate}[label=(\alph*)]
\item zbiór liczb Gödla poprawnych termów;
\item zbiór liczb Gödla poprawnych formuł;
\item totalna funkcja $\operatorname{sub}(f,x,n)$ zwracająca kod formuły powstałej z poprawnej formuły o kodzie $f$ przez zastąpienie wolnych wystąpień zmiennej o poprawnym kodzie $x$ numerałem $\overline n$; dla niepoprawnych danych jej wartością jest $0$;
\item dla efektywnie aksjomatyzowalnej teorii $T$ relacja $\operatorname{Prov}_T(p,f)$: „$p$ koduje poprawny dowód formuły o kodzie $f$ wraz ze świadectwami pochodzenia użytych aksjomatów”.
\end{enumerate}
\end{lemat}
\begin{proof}
\emph{Odczyt i sklejanie.} Dla $n>0$ wykładnik $p_i$ w $n$ jest najmniejszym $k\le n$, dla którego $p_i^{k+1}\nmid n$. Dla poprawnego kodu długość otrzymujemy jako liczbę kolejnych początkowych liczb pierwszych dzielących $n$. Dla dowolnego $n$ sprawdzamy dodatkowo, czy iloczyn potęg odtworzony z odczytanych wykładników jest równy $n$; wyklucza to luki i nieoczekiwane dalsze czynniki. Wyszukiwanie wykładników i długości ograniczamy przez $n$ (dla $n=1$ otrzymujemy kod pustego ciągu), natomiast dla $n=0$ stosujemy wartość zastępczą. Wszystkie te operacje są p.r. na mocy Lematu~\ref{res:5:8}. Konkatenacja kodów $m,n$ to odpowiednie przesunięcie indeksów liczb pierwszych w rozkładzie $n$, możliwe przez ograniczone iloczyny. Podciąg, dołączenie elementu i odczyt wiersza dowodu mają ten sam charakter.

\emph{Rozpoznawanie gramatyki.} Dla napisu długości $\ell$ tworzymy tablicę wszystkich jego podnapisów. Przechodzimy kolejno po długościach $1,\ldots,\ell$. Dla danego podnapisu sprawdzamy skończenie wiele podziałów na bezpośrednie składniki z Definicji~\ref{def:syntax-term} lub~\ref{def:syntax-formula}, korzystając wyłącznie z wcześniej obliczonych wpisów tablicy. Jest co najwyżej $\ell^2$ podnapisów, a liczba podziałów jest ograniczona funkcją p.r. od $\ell$. Kod tablicy aktualizujemy rekursją prostą. To ważne uzupełnienie: sama rekurencyjna definicja „podformuły są formułami” nie dowodzi jeszcze pierwotnej rekurencyjności testu poprawności.

\emph{Podstawienie.} Korzystając z drzewa rozbioru, przechodzimy po pozycjach napisu i zaznaczamy zasięgi kwantyfikatorów wiążących daną zmienną. Wolne wystąpienia zastępujemy napisem $S^n(0)$, inne pozostawiamy. Numerał jest domknięty, więc nie grozi przechwycenie jego zmiennych. Długość wyniku jest ograniczona przez $\ell(3n+1)$ przy zapisie $S(\cdot)$ z nawiasami; każda czynność i sklejenie kodów są p.r. Tak samo otrzymujemy funkcje $\operatorname{neg}(f)$ i $\operatorname{imp}(f,g)$.

\emph{Dowody z certyfikatami.} Wyliczenie aksjomatów zastępujemy programem, który na wejściu $a$ czeka, aż pojawi się aksjomat o kodzie $a$, i wtedy zwraca $1$ (Uwaga~\ref{res:6:4}). Przez \emph{efektywnie aksjomatyzowalną} rozumiemy tutaj teorię o rekurencyjnie przeliczalnym zbiorze aksjomatów pozalogicznych, przy ustalonym, efektywnym systemie reguł logiki pierwszego rzędu. Jest to program częściowo rekurencyjny, więc może mieć postać rejestrową z Definicji~\ref{res:6:2}. Wiersz dowodu zawiera kod formuły, rodzaj uzasadnienia i odpowiednie dane: instancję schematu logicznego, numery wcześniejszych przesłanek albo skończony ślad akceptacji aksjomatu pozalogicznego przez ten program. Nie rozstrzygamy, czy program kiedyś się zatrzyma: sprawdzamy już podany ślad. To jest nasze „świadectwo pochodzenia z wyliczenia”. Ślady, modus ponens, generalizacja i warunki dotyczące zmiennych wymagają wyłącznie ograniczonych testów na zakodowanym drzewie lub liście dowodu, łącznie ze sprawdzeniem zakresów rozładowanych założeń, jeśli używamy dedukcji naturalnej. Liczba wierszy oraz długości dołączonych świadectw nie przekraczają danych odczytanych z całego kodu $p$. Testy mają ograniczenia i funkcje aktualizacji p.r.; dlatego $\operatorname{Prov}_T(p,f)$ jest p.r. dla przyjętego kodu dowodu.

Każdy zwykły dowód można uzupełnić o certyfikaty: występuje w nim tylko skończenie wiele aksjomatów, a każdy zostaje kiedyś wypisany. Usunięcie certyfikatów daje zwykły dowód. Mamy więc dokładnie
\[
T\vdash\varphi\quad\Longleftrightarrow\quad\exists p\in\mathbb N\ \operatorname{Prov}_T(p,\ulcorner\varphi\urcorner).
\]
Bez certyfikatów rozstrzygalne, lecz nie p.r. sprawdzanie aksjomatów dawałoby jedynie rekurencyjność relacji dowodu. Ograniczone kwantyfikatory nie zmieniają funkcji rekurencyjnej w funkcję p.r.
\end{proof}

\begin{uwaga}[Dlaczego świadectwo akceptacji jest częścią kodu?]
Dla teorii z rozstrzygalnym zbiorem aksjomatów można po prostu
sprawdzić, czy dany wiersz jest aksjomatem. Dla teorii o aksjomatach
\emph{jedynie wyliczalnych} taki test mógłby nie kończyć się na
napisie, który aksjomatem nie jest. Dlatego do \emph{gotowego dowodu}
dołączamy skończony zapis obliczenia potwierdzający, że aksjomat
został wypisany. Weryfikator nie uruchamia wtedy nieograniczonego
poszukiwania: sprawdza konkretne przejścia skończonego śladu.
W szczególności zapis „relacja dowodu jest pierwotnie rekurencyjna”
dotyczy przyjętego \emph{pełnego kodu dowodu}, zawierającego te
certyfikaty, a nie dowolnie ustalonego surowego zapisu wierszy.
\end{uwaga}

\begin{przyklad}[Sprawdzenie dowodu a poszukiwanie dowodu]\label{res:6:6}
Sprawdzenie $\operatorname{Prov}_T(p,f)$ jest jak sprawdzanie gotowego rozwiązania: mając na papierze cały dowód (kod $p$), mechanicznie weryfikujemy linijka po linijce, że każdy krok jest dozwolony --- to skończona, ograniczona robota, stąd rekurencyjność. Natomiast pytanie $\exists p\,\operatorname{Prov}_T(p,f)$ każe znaleźć dowód wśród nieskończenie wielu kandydatów $p$. Nieograniczony kwantyfikator egzystencjalny nie zachowuje w ogólności rozstrzygalności: gdy dowód istnieje, przeszukiwanie kiedyś go znajdzie (stąd rekurencyjna przeliczalność), ale gdy nie istnieje --- szukanie może trwać w nieskończoność. To rozróżnienie „ograniczone a nieograniczone przeszukiwanie” jest sednem różnicy między zbiorem rekurencyjnym a rekurencyjnie przeliczalnym.
\end{przyklad}

\begin{definicja}[Predykat dowodliwości]\label{res:6:7}
Niech $T$ będzie efektywnie aksjomatyzowalną teorią w języku arytmetyki. Definiujemy $\operatorname{Bew}_T(x):=\exists p\,\operatorname{Prf}_T(p,x)$, gdzie $\operatorname{Prf}_T$ jest formułą arytmetyczną reprezentującą relację $\operatorname{Prov}_T$; jej standardowy wybór skonstruujemy po dowodzie reprezentowalności (Lemat~\ref{res:7:7}). Odróżniamy zatem relację na rzeczywistych liczbach $\operatorname{Prov}_T$ od formuły $\operatorname{Prf}_T$, którą można umieścić pod znakiem $T\vdash$. Obie wersje zgadzają się w $\mathfrak N$. Zbiór kodów zdań dowodliwych jest rekurencyjnie przeliczalny, jako rzut relacji rekurencyjnej; nierozstrzygalność w rozpatrywanych niesprzecznych rozszerzeniach $Q$ uzasadnimy osobno w dalszym rozdziale o niezupełności.
\end{definicja}

% Rozdział 7, strony 14--17.


\section{Jak odróżnić kod napisu od samego napisu?}
\begin{przyklad}[Pełny rachunek kodu dla trzech znaków]
Niech $a,b$ będą dwoma różnymi symbolami o kodach odpowiednio $1,2$,
zgodnie z Definicją~\ref{res:6:1}. Wówczas napis $aba$ otrzymuje kod
\[
\ulcorner aba\urcorner=2^{1+1}\cdot3^{2+1}\cdot5^{1+1}=2700.
\]
Z rozkładu $2700=2^2\cdot3^3\cdot5^2$ odzyskujemy
kolejno wykładniki $2,3,2$; odejmując od każdego $1$,
otrzymujemy kody znaków $1,2,1$. Jest to więc
operacja jednoznaczna, a nie nadanie napisowi znaczenia
arytmetycznego: liczba $2700$ nie \emph{jest} napisem $aba$,
lecz ustalonym kodem, z którego napis można odtworzyć.
\end{przyklad}
\begin{uwaga}[Przejście do rozdziału następnego]
W tym rozdziale skonstruowaliśmy arytmetyczne kody
formuł i dowodów oraz pokazaliśmy, w jakim sensie
kontrola gotowego kodu jest algorytmiczna. W rozdziale
\ref{ch:representability-diagonal} pokażemy, jak
umieszczać te kontrole wewnątrz formuł arytmetycznych.
\end{uwaga}
\chapter{Reprezentowalność i diagonalizacja}\label{ch:representability-diagonal}

\section{Co dokładnie będziemy udowadniać?}
W rozdziale~\ref{ch:arithmetization} przypisaliśmy
formułom i skończonym obliczeniom liczby. Sam kod jeszcze
nie wystarcza: potrzebujemy \emph{formuł arytmetycznych},
które opowiadają o tych kodach. Warto pamiętać, że
weryfikowanie danego obliczenia oraz udowodnienie w teorii
zdania o \emph{wszystkich} obliczeniach to różne zadania.

\begin{przyklad}[Różne poziomy wypowiedzi]
Obliczamy na kartce $2+3=5$, czyli w języku arytmetyki
$\overline2+\overline3=\overline5$.
Inna wypowiedź głosi, że \emph{istnieje dowód} tego zdania.
Po zakodowaniu dowodów drugą wypowiedź można zapisać
wewnątrz arytmetyki jako zdanie o pewnej liczbie $p$.
Jeszcze inna wypowiedź pyta, czy istnieje algorytm
rozstrzygający dla dowolnego zdania, czy ma dowód.
Te trzy pytania nie są tym samym.
\end{przyklad}

\section{Reprezentowalność}\label{sec:reprezentowalnosc}

\begin{definicja}[Reprezentowalność relacji]\label{res:7:1}
Zapis $\overline{\vec n}$ oznacza krotkę numerałów
$(\overline{n_1},\ldots,\overline{n_k})$; nie jest operacją
algebraiczną na krotce liczb.
Niech $T$ będzie teorią w języku $\mathcal L_{\mathrm{ar}}$ z Definicji~\ref{def:arith-language}, a $k\in\N$ ustaloną arnością. Relacja $R\subseteq\N^k$ jest \emph{reprezentowalna} w $T$, jeśli istnieje formuła $\rho(x_1,\ldots,x_k)$ bez innych zmiennych wolnych taka, że dla każdej krotki $\vec n\in\N^k$ zachodzi
\[
R(\vec n)\Longrightarrow T\vdash\rho(\overline{\vec n}),\qquad
\neg R(\vec n)\Longrightarrow T\vdash\neg\rho(\overline{\vec n}).
\]
\end{definicja}

\begin{definicja}[Reprezentowalność funkcji]\label{res:7:2}
Niech $T$ będzie teorią języka $\mathcal L_{\mathrm{ar}}$, a $k\in\N$. Funkcja częściowa $f:\N^k\rightharpoonup\N$ jest \emph{reprezentowalna} w $T$, jeśli istnieje formuła $\varphi_f(\vec x,y)$, której wszystkie zmienne wolne należą do $x_1,\ldots,x_k,y$, taka, że dla każdego standardowego wejścia $\vec n$, dla którego $f(\vec n)\downarrow=m$, zachodzi
\[
T\vdash\forall y\,\bigl(\varphi_f(\overline{\vec n},y)
\leftrightarrow y=\overline m\bigr).
\]
Na wejściach spoza dziedziny $f$ nie wymagamy dowodu nieistnienia wyniku. Dla funkcji totalnej warunek dotyczy każdego standardowego wejścia.
\end{definicja}

\begin{uwaga}[Metajęzyk i numerały pod znakami dowodliwości]
Liczba $g=\ulcorner G\urcorner$ jest w metateorii kodem napisu $G$.
Natomiast wewnątrz formuły arytmetycznej piszemy
$\operatorname{Bew}_T(\overline g)$, czyli podstawiamy \emph{term}
$\overline g=S^g(0)$. Dla czytelności w klasycznym zapisie
$\operatorname{Bew}_T(\ulcorner G\urcorner)$ numerał bywa domyślny;
w tym skrypcie w dowodach podkreślamy różnicę między liczbą
$g$ a jej numerałem $\overline g$.
\end{uwaga}

\begin{uwaga}[Obliczenie konkretne a totalność]\label{res:7:3}
Definicja~\ref{res:7:2} wymaga osobnego dowodu dla każdej ustalonej krotki numerałów; nie wynika z niej $T\vdash\forall\vec x\,\exists!y\,\varphi_f(\vec x,y)$. Pierwszy warunek wystarcza do diagonalizacji. Drugi jest istotnie
silniejszym \emph{wymogiem formalnym} i nie wolno go przyjmować
bez osobnego dowodu wewnątrz PA.
\end{uwaga}


\section{Przygotowanie do kodowania: chińskie twierdzenie o resztach}
Aby jedna liczba mogła pamiętać wiele kolejnych wartości,
potrzebujemy procedury, która równocześnie spełnia kilka
warunków na reszty z dzielenia. Przypomnimy i udowodnimy
potrzebną wersję twierdzenia, zanim użyjemy jej w kodowaniu.

\begin{definicja}[Względna pierwszość]
Liczby dodatnie $m,n$ nazywamy \emph{względnie pierwszymi},
jeśli ich jedynym wspólnym dodatnim dzielnikiem jest $1$.
Liczby $m_0,\ldots,m_s$ nazywamy \emph{parami względnie pierwszymi},
jeśli każde dwie różne są względnie pierwsze.
\end{definicja}

\begin{lemat}[Tożsamość Bézouta w potrzebnym przypadku]\label{lem:arith-bezout}
Jeśli dodatnie liczby całkowite $m,n$ są względnie pierwsze,
to istnieją $u,v\in\Z$ takie, że $um+vn=1$.
\end{lemat}
\begin{proof}
Rozważmy zbiór wszystkich dodatnich liczb postaci $am+bn$,
gdzie $a,b\in\Z$. Jest on niepusty, bo zawiera $m$.
Niech $d=um+vn$ będzie jego najmniejszym elementem.
Dzieląc $m$ przez $d$ z resztą, otrzymujemy
$m=qd+r$, gdzie $0\le r<d$. Ponieważ
$r=m-q(um+vn)=(1-qu)m-qvn$, reszta $r$ również
jest całkowitą kombinacją $m$ i $n$. Gdyby $r>0$,
sprzeciwiałaby się minimalności $d$. Zatem $r=0$
i $d\mid m$. Identyczny argument daje $d\mid n$.
Ponieważ $m,n$ są względnie pierwsze, $d=1$.
Stąd $um+vn=1$, zgodnie z tezą.
\end{proof}

\begin{twierdzenie}[Chińskie twierdzenie o resztach]\label{thm:arith-crt}
Niech $m_0,\ldots,m_s$ będą dodatnimi liczbami całkowitymi
parami względnie pierwszymi, a $c_0,\ldots,c_s\in\Z$.
Niech $M:=\prod_{i=0}^s m_i$. Istnieje dokładnie jedna
liczba $a\in\{0,\ldots,M-1\}$ taka, że
\[
a\equiv c_i\pmod{m_i}\qquad(i=0,\ldots,s).
\]
\end{twierdzenie}
\begin{proof}
Dla każdego $i$ połóżmy $M_i=M/m_i$.
Ponieważ $m_i$ jest względnie pierwsza z każdym innym czynnikiem
iloczynu, jest względnie pierwsza również z $M_i$.
Uzasadnia to wielokrotne zastosowanie Lematu~\ref{lem:arith-bezout}:
jeśli $au+bv=1$ i $ar+cw=1$, to po pomnożeniu tych
równości otrzymujemy $a(aur+cuw+brv)+bc(vw)=1$,
więc $a$ jest względnie pierwsza także z $bc$.
Powtarzamy ten argument dla kolejnych czynników $M_i$.
Z Lematu~\ref{lem:arith-bezout} istnieją zatem całkowite
$u_i,v_i$ takie, że $u_iM_i+v_im_i=1$.
Oznaczmy $a_0=\sum_{i=0}^s c_i u_i M_i$.
Ustalmy indeks $j$. Jeśli $i\ne j$, liczba $M_i$ jest
podzielna przez $m_j$, więc składnik $c_i u_i M_i$ daje
resztę zero z dzielenia przez $m_j$. Dla $i=j$ mamy
$u_j M_j\equiv1\pmod{m_j}$. Wynika stąd
$a_0\equiv c_j\pmod{m_j}$. Ponieważ $j$ był dowolny,
$a_0$ spełnia wszystkie wymagane kongruencje.

Niech $a$ będzie resztą z dzielenia $a_0$ przez dodatnie $M$.
Ma wtedy $0\le a<M$ i nadal spełnia te same kongruencje,
bo każde $m_i$ dzieli $M$. Dla jednoznaczności załóżmy,
że $a'$ jest drugim rozwiązaniem w przedziale $[0,M)$.
Wówczas każde $m_i$ dzieli $a-a'$. Dla dwóch względnie pierwszych liczb $m,n$ z
$m\mid z$ i $n\mid z$ wynika $mn\mid z$:
napiszmy $z=mk$, a z $um+vn=1$ otrzymamy
$k=u(mk)+v(nk)$ i stąd $n\mid k$. Powtarzając ten
argument dla kolejnych parami względnie pierwszych
modułów, widzimy, że $M\mid a-a'$. Mamy jednak $|a-a'|<M$, zatem $a=a'$.
\end{proof}

\begin{przyklad}[Dwie reszty w jednym numerze]
Chcemy znaleźć liczbę $a$ z resztą $2$ modulo $3$ i
resztą $3$ modulo $5$. Liczba $8$ spełnia oba warunki:
$8=2\cdot3+2=1\cdot5+3$. Ponadto każde inne
rozwiązanie różni się od $8$ o wielokrotność $15$.
Przechodząc do przedziału $[0,15)$, otrzymujemy więc
jednoznaczny wynik $8$. W Lemacie następnym będziemy
tak kodować nie dwie, lecz dowolnie wiele pozycji ciągu.
\end{przyklad}

\begin{lemat}[Kodowanie ciągów funkcją $\beta$]\label{res:7:4}
Funkcja $\beta(a,b,i)=a\bmod(1+(i+1)b)$ jest pierwotnie rekurencyjna. Dla każdego skończonego ciągu $c_0,\ldots,c_s$ istnieją $a,b\in\N$ takie, że $\beta(a,b,i)=c_i$ dla $i\le s$. Jej graf w modelu standardowym opisuje formuła $\Delta_0$
\[
B(a,b,i,r):=r<1+(i+1)b\ \land\
\exists q\le a\,[a=q(1+(i+1)b)+r].
\]
Jeśli wejściem jest kod $w=\langle c_0,\ldots,c_s\rangle$ niepustego ciągu
w numeracji z Definicji~\ref{res:6:1}, można wybrać parę $(a,b)$
funkcją pierwotnie rekurencyjną argumentu $w$.
\end{lemat}
\begin{proof}
Pierwotna rekurencyjność $\beta$ wynika z Lematu~\ref{res:5:8}. Wybierzmy $b=(s+1)!(1+\max_{i\le s}c_i)$ i oznaczmy $m_i=1+(i+1)b$. Mamy $c_i<m_i$. Moduły $m_0,\ldots,m_s$ są parami względnie pierwsze. Gdyby dwa z nich miały wspólny dzielnik większy od $1$, najmniejszy dzielnik większy od $1$ byłby liczbą pierwszą $p$ dzielącą oba. Dla $i<j$ otrzymalibyśmy $p\mid(m_j-m_i)=(j-i)b$. Ponieważ $m_i=1+(i+1)b$, liczba $p$ nie dzieli $b$. Pierwszość $p$ wymusza zatem $p\mid(j-i)$; ponieważ $1\le j-i\le s$, liczba $j-i$ dzieli $(s+1)!$, a więc dzieli $b$. To sprzeczność.

Z chińskiego twierdzenia o resztach istnieje $a$ spełniające $a\equiv c_i\pmod{m_i}$ dla każdego $i\le s$. Możemy wybrać $0\le a<M$, gdzie $M=\prod_{i\le s}m_i$; wtedy $\beta(a,b,i)=c_i$. Przypomnijmy, że konstrukcja reszt jest efektywna: dla $M_i=M/m_i$ wyznaczamy współczynniki Bézouta $u_iM_i+v_im_i=1$ i bierzemy resztę liczby $\sum_i c_i u_iM_i$ modulo $M$.

Wybór kodu jest nawet p.r. Z wejściowego kodu $w$
odczytujemy pierwotnie rekurencyjnie jego długość $s+1$ oraz
składniki $c_0,\ldots,c_s$ (Lemat~\ref{res:6:5}).
Obliczenie $b$ i $M$ wymaga tylko skończonych iloczynów, silni
i maksimum; są to operacje p.r.
Dla każdego $a<M$ sprawdzenie wszystkich kongruencji jest p.r.,
a na mocy udowodnionego istnienia i jednoznaczności z
Twierdzenia~\ref{thm:arith-crt} ograniczona minimalizacja
wybiera jedno $a$. Wartość na niepoprawnych kodach $w$
ustalamy jako $(0,0)$, zachowując totalność procedury.
Formuła $B$ jest $\Delta_0$, bo dzielnik jest dodatni, a iloraz $q$
nie przekracza $a$.
\end{proof}

\begin{lemat}[Arytmetyczny opis skończonego obliczenia]\label{res:7:5}
Dla ustalonego programu rejestrowego $e$ z Definicji~\ref{res:6:2} istnieje formuła $\Delta_0$ $C_e(\vec x,y,t,\vec a,\vec b)$ taka, że
\[
e(\vec n)\downarrow=m\quad\Longleftrightarrow\quad
\mathfrak N\models\exists t\,\exists\vec a\,\exists\vec b\,
C_e(\overline{\vec n},\overline m,t,\vec a,\vec b).
\]
Ponadto PA dowodzi jednoznaczności wyniku: dwa ślady o tym samym wejściu, spełniające $C_e$, mają ten sam wynik.
\end{lemat}
\begin{proof}
Konfiguracja programu składa się ze stanu i zawartości $d$ rejestrów, czyli $d+1$ liczb. Ustalamy kolejność współrzędnych $u=(u_0,u_1,\ldots,u_d)$: $u_0$ jest numerem stanu, natomiast $u_{h+1}$ jest zawartością rejestru $r_h$, $0\le h<d$. Każdą składową śladu kodujemy osobną parą $(a_h,b_h)$ z Lematu~\ref{res:7:4}, gdzie $0\le h\le d$. Przyjmijmy skróty
$B_i(u):=\bigwedge_{h=0}^d B(a_h,b_h,i,u_h)$ i
$A:=1+\sum_{h=0}^d a_h$; zapis $\exists u<A$ oznacza ograniczenie wszystkich współrzędnych krotki $u$ przez $A$.
Niech $I_e(\vec x,u)$ oznacza poprawną konfigurację początkową, $H_e(u)$ --- stan zatrzymania, a $S_e(u,v)$ --- wykonanie jednej instrukcji. Są to formuły bez kwantyfikatorów. Dla każdej z ustalonych skończenie wielu instrukcji wzrost rejestru zapisujemy równością $v_{h+1}=S(u_{h+1})$, a zmniejszenie dodatniej zawartości równością $u_{h+1}=S(v_{h+1})$; przypadek zerowy rozdzielamy warunkiem $u_{h+1}=0$. Pozostałe rejestry zachowują wartości, a numer stanu zmienia się zgodnie z instrukcją. Alternatywa tych przypadków definiuje $S_e(u,v)$ i wyraża deterministyczny pojedynczy krok bez użycia odejmowania. Formuła $I_e$ ustala stan początkowy i rejestry wejściowe, a $H_e$ rozpoznaje końcowy numer stanu; $u_1$ jest zawartością rejestru wynikowego $r_0$.

Definiujemy
\begin{align*}
C_e(\vec x,y,t,\vec a,\vec b):={}&
\exists u<A\,[B_0(u)\land I_e(\vec x,u)]\\
&\land\forall i<t\ \exists u<A\ \exists v<A\,
[B_i(u)\land B_{S(i)}(v)\land\neg H_e(u)\land S_e(u,v)]\\
&\land\exists u<A\,[B_t(u)\land H_e(u)\land u_1=y].
\end{align*}
Jest to formuła $\Delta_0$. Dla rzeczywistego śladu wszystkie odczyty $\beta(a_h,b_h,i)$ są nie większe od $a_h<A$, więc kodowanie z poprzedniego lematu dostarcza świadków. Odwrotnie, w $\N$ jednoznaczność reszty sprawia, że formuła opisuje kolejne konfiguracje tego samego obliczenia. Warunek $\neg H_e(u)$ dla $i<t$ oznacza, że ślad kończy się przy pierwszym zatrzymaniu.

Uzasadnijmy ostatnie stwierdzenie \emph{wewnątrz PA}. Arytmetyka ta dowodzi istnienia i jednoznaczności reszty przy dzieleniu przez dowolną dodatnią liczbę: istnienie otrzymujemy przez indukcję po dzielnej, a jednoznaczność z własności porządku i mnożenia. Stąd PA dowodzi, że każdy kod $(a_h,b_h)$ wyznacza dokładnie jedną wartość składowej w dowolnej chwili. Dla dwóch śladów o tym samym wejściu indukcja po wspólnym odcinku ich czasów wykazuje równość kolejnych konfiguracji: baza wynika z $I_e$, a krok z deterministyczności $S_e$. Gdyby jeden ślad kończył się wcześniej, drugi wymagałby w tej samej chwili $\neg H_e$, sprzecznie z $H_e$ w pierwszym. Czasy zatrzymania, a więc i wyniki, są równe. Argument wykorzystuje indukcję w PA; nie przenosimy go bez zmian do $Q$.
\end{proof}

\begin{twierdzenie}[Reprezentowalność funkcji rekurencyjnych]\label{res:7:6}
Każda funkcja częściowo rekurencyjna jest reprezentowalna w $Q$ w sensie Definicji~\ref{res:7:2}; dotyczy to w szczególności funkcji totalnych rekurencyjnych. Każda relacja rekurencyjna jest reprezentowalna w $Q$ w sensie Definicji~\ref{res:7:1}. Formuły te reprezentują odpowiednie funkcje i relacje także w każdym rozszerzeniu $T\supseteq Q$.
\end{twierdzenie}
\begin{proof}
Na mocy Lematu~\ref{res:6:3} wybieramy program $e$ obliczający $f$. Nie możemy powtórzyć rozumowania z Lematu~\ref{res:7:5}, ponieważ $Q$ nie ma indukcji i nie dowodzi ogólnej jednoznaczności dzielenia z resztą. Wzmocnimy więc formułę kodującą ślad o warunki potrzebne do prześledzenia \emph{każdego konkretnego} obliczenia.

Niech
\begin{align*}
D(t)&:=\forall u\le t\ \forall j<u\,(j<t),\\
U(\vec a,\vec b,t)&:=\forall i\le t\ \bigwedge_{h=0}^d\forall v\,\forall w\,
\bigl((B(a_h,b_h,i,v)\land B(a_h,b_h,i,w))\to v=w\bigr).
\end{align*}
Pierwszy warunek jest dodatkową hipotezą o parametrze czasu $t$:
jeżeli $u\le t$ i $j<u$, możemy wprost wywnioskować $j<t$,
bez przyjmowania ogólnego twierdzenia o przechodniości $<$ w $Q$.
Drugi gwarantuje jednoznaczność odczytów również wtedy,
gdy parametry kodu w pewnym modelu $Q$ są niestandardowe. Kładziemy
\[
\varphi_f(\vec x,y):=\exists t\,\exists\vec a\,\exists\vec b\,
\bigl(C_e(\vec x,y,t,\vec a,\vec b)\land D(t)\land U(\vec a,\vec b,t)\bigr).
\]
Nie zakładamy, że formuła $\varphi_f$ jest $\Sigma_1$; wymaganie to nie występuje w Definicji~\ref{res:7:2}.

Ustalmy $\vec n$ takie, że $f(\vec n)\downarrow=m$, i niech $s$ będzie rzeczywistą liczbą kroków do zatrzymania. Wybieramy standardowe kody tego śladu. Zamknięte instancje $C_e$ i $D(\overline s)$ są prawdziwymi zdaniami $\Delta_0$, zatem $Q$ je dowodzi na mocy Lematu~\ref{res:4:7}. Dla wybranych numerałów $\overline{\vec a},\overline{\vec b},\overline s$
teoria $Q$ dowodzi również $U$. Ustalmy składową $h$ i
standardową chwilę $j\le s$. Z Lematu~\ref{res:4:4} możemy
rozwinąć $i\le\overline s$ na skończoną listę wartości
$i=\overline0,\ldots,\overline s$. Po podstawieniu
$i=\overline j$ każdy warunek
$B(\overline{a_h},\overline{b_h},\overline j,v)$ implikuje
$v<\overline{M_{h,j}}$, gdzie $M_{h,j}=1+(j+1)b_h>0$.
Ten sam lemat sprowadza $v$ do jednego ze standardowych
numerałów $\overline0,\ldots,\overline{M_{h,j}-1}$.
Analogicznie świadek ilorazu spełnia $q\le\overline{a_h}$,
więc ma jedną ze skończenie wielu wartości standardowych.
Dla każdej pary $(q,v)$ rozstrzygamy w $Q$ instancję
$\overline{a_h}=q\overline{M_{h,j}}+v$ przez rachunek na
numerałach. Dla ustalonego $a_h$ i dodatniego $M_{h,j}$
dokładnie jedna para spośród par z $0\le v<M_{h,j}$
spełnia tę równość. Po złożeniu skończenie wielu dowodów
$Q$ dowodzi zatem, że dowolne dwa odczyty $v,w$ są równe.
Tak otrzymujemy zamkniętą instancję $U$ bez korzystania
z indukcji wewnątrz $Q$. Ponieważ rzeczywisty ślad spełnia
$C_e$, uzyskujemy $Q\vdash\varphi_f(\overline{\vec n},\overline m)$.

Pozostało wykazać, że \emph{dowolny} wynik wskazany przez
$\varphi_f(\overline{\vec n},y)$ jest równy $\overline m$,
choć świadkowie $t,\vec a,\vec b$ mogą być niestandardowi.
Niech $u^{(0)},\ldots,u^{(s)}$ będą konkretnymi standardowymi
konfiguracjami rzeczywistego obliczenia, zatrzymanego po
$s$ krokach. W $Q$ rozpatrujemy dowolnych świadków
$C_e\land D\land U$ i stosujemy lemat~\ref{res:4:4}
do rozłącznych przypadków $t<\overline s$,
$t=\overline s$ i $\overline s<t$.

Wyjaśnijmy wspólny mechanizm: dla każdej \emph{ustalonej}
standardowej chwili $j$, do której obowiązują warunki śladu,
można w $Q$ otrzymać $B_j(\overline{u^{(j)}})$. W chwili $0$
wynika to z $I_e(\overline{\vec n},u)$ i $B_0(u)$
w $C_e$: początkowa konfiguracja jest wymuszona przez
równości definiujące $I_e$. Jeśli mamy odczyt rzeczywistej
konfiguracji w chwili $j$ oraz warunek przejścia $\overline j<t$,
$U$ utożsamia współrzędne użyte w obu odczytach kodu
$B_j$. Trzeba jeszcze uzasadnić zastosowanie $U$ do
następnej chwili. Dla ustalonego standardowego $j$ teoria $Q$
dowodzi $(\overline j<t)\to(\overline{j+1}\le t)$:
z (Q8) otrzymujemy $t=z+S(\overline j)$, a (Q3)
daje $z=0$ albo $z=S(w)$. W pierwszym przypadku
rachunek na ustalonym numerale prowadzi do
$t=\overline{j+1}$; w drugim do
$t=S^{j+2}(w)=w+S(\overline{j+1})$, a więc
$\overline{j+1}<t$ na mocy (Q8). Zatem $U$ obejmuje
także odczyt $B_{j+1}$. Ponieważ $u^{(j)}$ jest ustaloną
standardową konfiguracją, po podstawieniu jej numerałów
skończona alternatywa instrukcji $S_e$ redukuje się w $Q$
do jednej właściwej gałęzi: z (Q1)--(Q7) rozstrzygamy
przypadek zerowy lub dodatni i obliczamy wszystkie
nowe współrzędne. Otrzymujemy dokładnie
$u^{(j+1)}$. Powtarzamy tę \emph{skończoną} liczbę kroków
w metateorii; budujemy odpowiedni formalny dowód w $Q$
dla każdego potrzebnego standardowego $j$, nie stosując
indukcji po dowolnym elemencie modelu $Q$.

Jeżeli $t<\overline s$, Lemat~\ref{res:4:4} daje skończoną
alternatywę $t=\overline0\lor\cdots\lor t=\overline{s-1}$.
Dla każdego jej składnika rozwijamy dokładnie tyle kroków
opisanej konstrukcji, ile wymaga dany numerał. W rozpatrywanej gałęzi $t=\overline j$ otrzymujemy
$B_j(\overline{u^{(j)}})$, a rzeczywista konfiguracja w tej chwili
nie jest końcowa, bo $s$ jest pierwszym czasem zatrzymania.
Końcowy warunek $H_e$ z $C_e$ daje więc sprzeczność.

Jeżeli $t=\overline s$, ten sam skończony argument odtwarza
konfigurację $\overline{u^{(s)}}$ w chwili zatrzymania. Ponieważ
$u_1=y$ w ostatnim warunku $C_e$, otrzymujemy
$y=\overline m$.

Wreszcie jeśli $\overline s<t$, warunek $D(t)$ pozwala
uzyskać $\overline j<t$ dla każdego standardowego $j<s$:
bierzemy $u=\overline s\le t$ i korzystamy ze standardowego
porównania $\overline j<\overline s$. Sama nierówność
$\overline s<t$ pozwala ponadto zastosować warunek
przejścia również w chwili $i=\overline s$. Skończenie
wiele wcześniejszych kroków daje odczyt konfiguracji
$\overline{u^{(s)}}$, już końcowej, natomiast przejście w tej chwili
wymaga $\neg H_e(\overline{u^{(s)}})$. Sprzeczność. Gdy $s=0$,
nie ma kroków pośrednich i wystarczy konfiguracja początkowa.

Otrzymaliśmy formalnie w $Q$ implikację
$\varphi_f(\overline{\vec n},y)\to y=\overline m$.
Ponieważ dowiedliśmy też $\varphi_f(\overline{\vec n},\overline m)$,
podstawianie równych termów daje implikację odwrotną
$y=\overline m\to\varphi_f(\overline{\vec n},y)$.
Generalizując $y$, otrzymujemy
$Q\vdash\forall y\,(\varphi_f(\overline{\vec n},y)\leftrightarrow y=\overline m)$.

Dla relacji rekurencyjnej $R$ bierzemy jej totalną funkcję
charakterystyczną i kładziemy
$\rho_R(\vec x):=\varphi_{\chi_R}(\vec x,\overline1)$.
Jeśli $R(\vec n)$, otrzymujemy dowód
$\rho_R(\overline{\vec n})$; w przeciwnym razie
reprezentowalność wartości $0$ daje dowód jej negacji,
ponieważ $Q\vdash\overline0\ne\overline1$.
Dowody z $Q$ są też dowodami w każdym rozszerzeniu $Q$.
\end{proof}

\begin{przyklad}[Dlaczego samo istnienie śladu nie wystarcza do silnej reprezentacji]
Załóżmy, że program na wejściu $2$ wykonuje dwa kroki i zwraca $5$.
Sama formuła „istnieje ślad kończący się wynikiem $y$” pokazuje,
że $Q$ dowodzi jej instancji dla $y=\overline5$.
Definicja~\ref{res:7:2} wymaga jednak więcej: $Q$ musi jeszcze
dowiedzieć, że \emph{każdy} wynik opisany przez tę formułę
jest równy $\overline5$, nawet jeśli wartości zmiennych kodujących
ślad są niestandardowe w pewnym modelu $Q$.
Z tego właśnie powodu w Twierdzeniu~\ref{res:7:6}
dodaliśmy $U$ (jednoznaczność odczytu kodu) i $D$
(legalność odwołań do wcześniejszych chwil). Rozumowanie
rozdziela trzy przypadki czasu $t$: mniejszy od rzeczywistego
czasu $2$, równy mu i większy. W pierwszym przypadku nie może
jeszcze wystąpić zatrzymanie; w drugim wynik jest równy $5$;
w trzecim kod żąda kolejnego kroku po zatrzymaniu, co daje
sprzeczność. Nie zakładamy indukcji po zmiennej $t$ w $Q$.
\end{przyklad}

\begin{lemat}[Standardowy arytmetyczny predykat dowodu]\label{res:7:7}
Dla efektywnie aksjomatyzowalnej teorii $T$ można wybrać kodowanie pełnych dowodów tak, aby formuła $\operatorname{Prf}_T(p,x)$ była klasy $\Delta_0$ i poprawnie rozstrzygała relację dowodu w modelu standardowym. Dla każdej pary numerałów $\overline p,\overline x$ teoria $Q$ dowodzi odpowiedniej instancji $\operatorname{Prf}_T$ albo jej negacji. Wówczas $\operatorname{Bew}_T(x):=\exists p\,\operatorname{Prf}_T(p,x)$ jest formułą $\Sigma_1$.
\end{lemat}
\begin{proof}
Z Lematu~\ref{res:6:5} mamy pierwotnie rekurencyjny weryfikator $V_T(d,x)$ pełnego dowodu $d$, zawierającego certyfikaty pochodzenia aksjomatów. Ustalamy program obliczający $V_T$ i dołączamy do $d,x$ kod całego śladu jego obliczenia, kończącego się wynikiem $1$.

Dane $(d,t,\vec a,\vec b)$, przy ustalonych długościach krotek zależnych od wybranego programu weryfikującego,
pakujemy iterowanym parowaniem
$\langle u,v\rangle=(u+v)^2+u$. Jest ono p.r. i różnowartościowe:
przy ustalonej sumie $u+v=s$ przyjmuje wartości w przedziale
$[s^2,s^2+s]$, rozłącznym z przedziałami dla innych sum.
Każda składowa jest nie większa od pełnego kodu $p$,
gdyż $\langle u,v\rangle\ge u,v$ dla $u,v\in\N$;
własność zachowuje iteracja. Dla ustalonej liczby składowych
wielokrotnie złożone parowanie jest określonym termem
języka $\mathcal L_{\mathrm{ar}}$: każdą parę zastępujemy
wyrażeniem $(u+v)(u+v)+u$. Niech
$\operatorname{Pack}(p;d,t,\vec a,\vec b)$
oznacza równość $p$ z tym termem, a więc formułę
bez kwantyfikatorów. Ograniczenie we wszystkich następujących
kwantyfikatorach rozumiemy współrzędnie i niezależnie:
każda współrzędna $d,t,\vec a,\vec b$ jest nie większa od $p$. Definiujemy
\[
\operatorname{Prf}_T(p,x):=\exists d,t,\vec a,\vec b\le p\,
\bigl(\operatorname{Pack}(p;d,t,\vec a,\vec b)\land
C_{V_T}(d,x,\overline1,t,\vec a,\vec b)\bigr).
\]
Ponieważ $\operatorname{Pack}$ nie ma kwantyfikatorów, $C_{V_T}$
jest $\Delta_0$, a wszystkie dodatkowe kwantyfikatory są
ograniczone przez term $p$, cała formuła jest $\Delta_0$.
Ustalone urządzenie $V_T$ jest totalnym weryfikatorem z Lematu~\ref{res:6:5}:
na wejściu $(d,x)$ zwraca $1$ dokładnie wtedy, gdy $d$
jest pełnym, poprawnym dowodem formuły o kodzie $x$.
Każdy taki dowód daje więc zakończone obliczenie weryfikatora
i kod jego śladu; z kolei ślad zakończony wynikiem $1$
może istnieć jedynie dla zaakceptowanej pary $(d,x)$.
Powstaje równoważność w \emph{modelu standardowym}; nie
przyjmujemy z niej bez dowodu ogólnej własności w modelach $Q$.
Otrzymany $p$ koduje zarówno dowód, jak i świadectwo
obliczenia, dlatego relacja $\operatorname{Prf}_T$ ma tę samą
projekcję egzystencjalną na kod zdania co
$\operatorname{Prov}_T$ z Lematu~\ref{res:6:5}.
Dla dwóch konkretnych numerałów prawdziwość lub fałszywość
zamkniętej formuły $\Delta_0$ jest dowodliwa w $Q$ na mocy
Lematu~\ref{res:4:7}. Egzystencjalne domknięcie po $p$
jest formułą $\Sigma_1$.
\end{proof}

% Schemat diagonalizacji i lemat o punkcie stałym.

\section{Schemat operacji diagonalnej}
W poniższym rysunku każda strzałka oznacza
\emph{ściśle określoną operację na skończonym napisie albo jego kodzie}.
Prawdziwość formuły nie jest tutaj jeszcze rozpatrywana.

\begin{center}
\begin{tikzpicture}[node distance=5mm,>=Stealth,
  every node/.style={font=\small},box/.style={draw,rounded corners,align=center,inner sep=5pt}]
\node[box] (a) {$\theta(x)$\\formuła};
\node[box,right=of a] (b) {$m=\ulcorner\theta\urcorner$\\kod formuły};
\node[box,right=of b] (c) {$\theta(\overline m)$\\podstawienie};
\node[box,right=of c] (d) {$d(m)$\\kod wyniku};
\draw[->] (a)--(b);
\draw[->] (b)--(c);
\draw[->] (c)--(d);
\end{tikzpicture}
\end{center}

Dla przykładu, jeśli $\theta(x)$ oznacza $x=x$, to kod $m$
oznacza napis $x=x$, a $d(m)$ oznacza kod zdania
$\overline m=\overline m$. W następnym lemacie
przeniesiemy taki zabieg z poziomu rachunku na napisach
do poziomu dowodliwych zdań arytmetycznych.

\section{Lemat o punkcie stałym (diagonalizacja)}\label{sec:diagonalizacja}
\begin{lemat}[Lemat diagonalny / o punkcie stałym]\label{res:8:1}
Niech $\psi(x)$ będzie formułą języka $\mathcal L_{\mathrm{ar}}$,
której jedyną możliwą zmienną wolną jest $x$.
Istnieje zdanie $G$ takie, że dla liczby $g:=\ulcorner G\urcorner$
zachodzi $Q\vdash G\leftrightarrow\psi(\overline g)$.
Ta sama równoważność jest dowodliwa w każdej teorii
$T\supseteq Q$ tego języka, niezależnie od efektywności jej aksjomatów. Zdanie $G$ „mówi o sobie”, że ma własność $\psi$; nie jest to dosłowna równość napisów, lecz dowodliwa równoważność.
\end{lemat}
\begin{proof}
Ustalamy wyróżnioną zmienną $x$ i definiujemy funkcję diagonalizacji
\[
d(n):=\ulcorner\theta(\overline n)\urcorner,
\]
gdy $n$ koduje formułę $\theta(x)$ bez innych zmiennych wolnych; na pozostałych wejściach kładziemy $d(n)=0$. Jest to funkcja p.r. na mocy Lematu~\ref{res:6:5}(c): podstawiamy numerał $\overline n$ w formule o kodzie $n$. Twierdzenie~\ref{res:7:6} daje formułę $\varphi_d(x,y)$ z
\[
Q\vdash\forall y\,(\varphi_d(\overline n,y)\leftrightarrow y=\overline{d(n)})
\]
dla każdego standardowego $n$.

Najpierw przemianowujemy zmienne związane w $\psi(x)$
i $\varphi_d$, jeśli mogłyby kolidować ze zmienną $y$;
następnie zastępujemy wolne wystąpienia $x$ w $\psi$ przez $y$.
Dzięki tej kolejności żadne podstawienie nie przechwytuje zmiennej.
Definiujemy $\delta(x):=\exists y\,(\varphi_d(x,y)\land\psi(y))$,
$m:=\ulcorner\delta(x)\urcorner$ i $G:=\delta(\overline m)$.
Formuła $\delta$ ma co najwyżej zmienną wolną $x$, toteż
$G$ jest zdaniem. Oznaczmy $g:=\ulcorner G\urcorner$.
Z definicji $d(m)=\ulcorner\delta(\overline m)\urcorner=g$,
więc reprezentowalność daje
\[
Q\vdash\forall y\,(\varphi_d(\overline m,y)\leftrightarrow y=\overline g).
\]
Korzystając z równoważności dla $\varphi_d(\overline m,y)$
i z reguł logiki pierwszego rzędu z równością, otrzymujemy
\[
\begin{aligned}
Q&\vdash G\leftrightarrow\exists y\,(y=\overline g\land\psi(y)),\\
Q&\vdash \exists y\,(y=\overline g\land\psi(y))
\leftrightarrow\psi(\overline g).
\end{aligned}
\]
Z przechodniości równoważności otrzymujemy
$Q\vdash G\leftrightarrow\psi(\overline g)$.
Pierwsza równoważność wynika przez zastąpienie formuły
$\varphi_d(\overline m,y)$ równoważną formułą $y=\overline g$
pod kwantyfikatorem, druga przez eliminację kwantyfikatora
egzystencjalnego i zastępowanie równych.
To jest szukana równoważność. Nie użyto żadnej hipotezy o aksjomatach $T$.
\end{proof}



\section{Dlaczego lemat diagonalny nie jest zwykłym podstawieniem?}
\begin{przyklad}[Instrukcja „podstaw własny numer”]
Rozważmy formułę $\theta(x)$ z jedną zmienną wolną.
Jej kod $m=\ulcorner\theta\urcorner$ jest liczbą naturalną.
Operacja $d(m)$ tworzy kod napisu $\theta(\overline m)$.
Dostajemy więc zależność między \emph{dwiema liczbami}:
kodem formuły i kodem jej instancji.
Lemat~\ref{res:8:1} pozwala opisać tę zależność
\emph{formułą arytmetyczną}, a następnie użyć jej
wewnątrz teorii $Q$.
Nie stwierdzamy, że zdanie jest dosłownie swoim
kodem; stwierdzamy, że pewna równoważność jest
\emph{formalnie dowodliwa}.
\end{przyklad}
\begin{uwaga}[Granica tej iteracji]
W kolejnych rozdziałach wykorzystamy zbudowaną relację
dowodu oraz lemat diagonalny do pierwszego twierdzenia
o niezupełności. Będziemy musieli oddzielnie rozważyć
niesprzeczność, $\omega$-niesprzeczność i wzmocnienie Rossera;
bez tych warunków nie wolno formułować jednej
bezwarunkowej wersji twierdzenia.
\end{uwaga}
\chapter{Pierwsze twierdzenie Gödla o niezupełności}\label{ch:first-incompleteness}

\section{Dlaczego mechaniczne wyliczanie dowodów nie wystarcza?}
W rozdziałach~\ref{ch:arithmetization}--\ref{ch:representability-diagonal}
zakodowaliśmy dowody jako liczby naturalne, a potem skonstruowaliśmy
formułę arytmetyczną mówiącą, że pewna liczba jest kodem poprawnego
dowodu. Pozwala nam to zadać bardzo szczególne pytanie: czy teoria
potrafi udowodnić zdanie, które po zakodowaniu głosi o sobie, że
\emph{nie ma dowodu w tej teorii}?

Musimy precyzyjnie odróżniać trzy kwestie. Zdanie może być
prawdziwe w standardowym zbiorze liczb naturalnych; może mieć
dowód w wybranej teorii; i może być prawdziwe w jakimś
\emph{niestandardowym} modelu tej teorii. W tym rozdziale tych
pojęć nie utożsamiamy.

\begin{definicja}[Niesprzeczność i zupełność teorii]\label{def:incomp-properties}
Niech $T$ będzie teorią pierwszego rzędu o ustalonym klasycznym
rachunku dowodowym. Teorię $T$ nazywamy \emph{niesprzeczną},
jeśli nie istnieje zdanie $\sigma$ takie, że
$T\vdash\sigma$ oraz $T\vdash\neg\sigma$.
Nazywamy ją \emph{zupełną}, jeśli dla każdego zdania $\sigma$
jej języka zachodzi $T\vdash\sigma$ lub $T\vdash\neg\sigma$.
Zdanie, dla którego nie zachodzi żadna z tych dowodliwości,
nazywamy \emph{nierozstrzygalnym w $T$}.
\end{definicja}

\begin{uwaga}[Niezupełność a sprzeczność]
Teoria sprzeczna w klasycznym rachunku dowodzi każdego zdania,
a więc jest zupełna w powyższym znaczeniu, lecz nie interesuje
nas taka „zupełność”. Pierwsze twierdzenie Gödla dotyczy
\emph{niesprzecznych} teorii o efektywnie wyliczalnych aksjomatach.
Na przykład pełna teoria ustalonego skończonego grafu jest
zupełna i niesprzeczna; potrzebna jest odpowiednia siła arytmetyczna.
\end{uwaga}

\begin{definicja}[$\omega$-niesprzeczność]\label{def:omega-consistency-incomp}
Niech $T$ będzie teorią w języku pierwszego rzędu
zawierającym język arytmetyki $\mathcal L_{\mathrm{ar}}$.
Teoria $T$ jest \emph{$\omega$-niesprzeczna},
jeżeli nie istnieje formuła $\theta(x)$ jej języka,
której jedyną zmienną wolną jest $x$, taka, że jednocześnie
\[
 T\vdash\exists x\,\theta(x)
 \quad\text{oraz}\quad
 (\forall n\in\N)\;T\vdash\neg\theta(\overline n).
\]
To znaczy: teoria nie może udowodnić, że istnieje
świadek pewnej własności, oraz osobno wykluczyć
\emph{każdego konkretnego} numerału jako świadka.
\end{definicja}

\begin{lemat}[$\omega$-niesprzeczność implikuje niesprzeczność]
Każda $\omega$-niesprzeczna teoria arytmetyczna jest niesprzeczna.
\end{lemat}
\begin{proof}
Gdyby $T$ była sprzeczna, dowodziłaby dowolnego zdania
oraz jego negacji. Biorąc $\theta(x):=(x=x)$, otrzymalibyśmy
$T\vdash\exists x\,(x=x)$ i dla każdego $n\in\N$
$T\vdash\neg(\overline n=\overline n)$. Jest to dokładnie
sytuacja wykluczona w Definicji~\ref{def:omega-consistency-incomp}.
\end{proof}

\section{Jak formuła mówi o dowodach?}
Przypomnijmy konstrukcję z Lematu~\ref{res:7:7}.
Jeżeli aksjomaty teorii $T$ w efektywnie kodowanym języku
można algorytmicznie wyliczać (bez zakładania rozstrzygalności
samego zbioru aksjomatów), przyjmujemy kod pełnego dowodu $p$,
zawierający również certyfikaty pochodzenia jego aksjomatów. Wtedy
$\operatorname{Prf}_T(p,x)$ jest formułą arytmetyczną
$\Delta_0$, która w strukturze standardowej opisuje
weryfikację dowodu o kodzie $p$ i ostatnim zdaniu o kodzie $x$.
Definiujemy
\[
\operatorname{Bew}_T(x):=\exists p\,\operatorname{Prf}_T(p,x).
\]
Przez $g=\ulcorner G\urcorner$ oznaczamy \emph{liczbę}
kodującą zdanie $G$, a przez $\overline g$ jej \emph{numerał},
czyli term języka arytmetyki. W zapisie formuły $\operatorname{Bew}_T(\overline g)$
występuje zatem numerał, a nie liczba metateoretyczna $g$.
Podobnie w zamkniętych instancjach $\operatorname{Prf}_T$
kody dowodu i zdania zastępujemy ich numerałami
$\overline p$ i $\overline g$; w formułach otwartych
pozostają odpowiednie zmienne.

\begin{lemat}[Dwie podstawowe własności predykatu dowodu]\label{lem:incomp-proof-check}
Niech $T\supseteq Q$ będzie teorią o efektywnie wyliczalnych
aksjomatach i $g\in\N$ liczbą kodującą zdanie.
\begin{enumerate}[label=(\roman*)]
\item Jeśli $p\in\N$ rzeczywiście koduje pełny dowód zdania
o kodzie $g$, to $Q\vdash\operatorname{Prf}_T(\overline p,\overline g)$.
\item Jeśli $p\in\N$ nie jest takim kodem, to
$Q\vdash\neg\operatorname{Prf}_T(\overline p,\overline g)$.
\end{enumerate}
W pierwszym przypadku ponadto $T\vdash\operatorname{Bew}_T(\overline g)$.
\end{lemat}
\begin{proof}
Formuła $\operatorname{Prf}_T$ jest $\Delta_0$,
a jej prawdziwość w $\mathfrak N$ jest równoważna
rzeczywistemu przejściu algorytmicznej weryfikacji.
Lemat~\ref{res:4:7} mówi, że $Q$ dowodzi każdej prawdziwej
instancji zamkniętej formuły $\Delta_0$ i zaprzeczenia
każdej fałszywej. To dowodzi (i) i (ii).
W sytuacji (i) wprowadzamy kwantyfikator egzystencjalny,
otrzymując $Q\vdash\exists u\,\operatorname{Prf}_T(u,\overline g)$.
Każdy dowód w $Q$ jest również dowodem w $T$, ponieważ
$T$ zawiera aksjomaty $Q$.
\end{proof}

\begin{przyklad}[Co dokładnie daje gotowy dowód?]
Rozważmy zdanie $\sigma:=(0=0)$. Z aksjomatu zwrotności
równości istnieje jego skończony dowód w $Q$, a więc także
w dowolnej teorii $T\supseteq Q$. Wybierzmy rzeczywisty
pełny kod $p$ tego dowodu, zawierający certyfikaty
pochodzenia aksjomatów, i połóżmy
$g=\ulcorner\sigma\urcorner$. Weryfikator akceptuje parę
$(p,g)$, zatem Lemat~\ref{lem:incomp-proof-check} daje
$Q\vdash\operatorname{Prf}_T(\overline p,\overline g)$.
Dopiero po wprowadzeniu kwantyfikatora egzystencjalnego
i użyciu $T\supseteq Q$ otrzymujemy
$T\vdash\operatorname{Bew}_T(\overline g)$.
Nie wynika stąd, że każdy kod oznacza poprawny dowód:
formalnie potwierdziliśmy \emph{jeden} konkretny dowód.
\end{przyklad}

\section{Zdanie Gödla}
\begin{twierdzenie}[Pierwsze twierdzenie Gödla --- postać pierwotna]\label{thm:godel-first-original}
Niech $T$ będzie efektywnie aksjomatyzowalną teorią
klasycznej logiki pierwszego rzędu w efektywnie kodowanym
języku zawierającym $\mathcal L_{\mathrm{ar}}$, przy czym
$T\supseteq Q$.
Istnieje zdanie arytmetyczne $G_T$ o kodzie
$g=\ulcorner G_T\urcorner$ takie, że
\[
 Q\vdash G_T\leftrightarrow
  \neg\operatorname{Bew}_T(\overline g).
\]
Zdanie to ma następujące własności:
\begin{enumerate}[label=(\alph*)]
\item jeśli $T$ jest niesprzeczna, to $T\nvdash G_T$;
\item jeśli $T$ jest $\omega$-niesprzeczna, to $T\nvdash\neg G_T$.
\end{enumerate}
W szczególności $\omega$-niesprzeczna teoria o podanych
własnościach jest niezupełna.
\end{twierdzenie}
\begin{proof}
Lemat diagonalny (Lemat~\ref{res:8:1}) stosujemy do
formuły $\psi(x):=\neg\operatorname{Bew}_T(x)$.
Otrzymujemy zdanie $G_T$ i dowodliwą w $Q$
równoważność z tezy. Oznaczmy $g=\ulcorner G_T\urcorner$.

\emph{Punkt (a).} Przypuśćmy, że $T\vdash G_T$.
Istnieje wtedy rzeczywisty, skończony dowód tego zdania,
więc możemy w metateorii wskazać jego kod $p\in\N$.
Na mocy Lematu~\ref{lem:incomp-proof-check}
$T\vdash\operatorname{Bew}_T(\overline g)$.
Z drugiej strony $T$ zawiera $Q$ i dowodzi $G_T$;
stosując eliminację równoważności, otrzymujemy
$T\vdash\neg\operatorname{Bew}_T(\overline g)$.
Są to dwa przeciwne zdania, więc $T$ byłaby sprzeczna.
Otrzymaliśmy sprzeczność z założeniem.

\emph{Punkt (b).} Zakładamy teraz $\omega$-niesprzeczność.
Z Lematu poprzedniego wynika niesprzeczność, a punkt (a)
daje $T\nvdash G_T$. Zatem \emph{żadna} standardowa
liczba $p\in\N$ nie koduje dowodu $G_T$.
Dla każdego ustalonego $p$ stosujemy część (ii) Lematu
\ref{lem:incomp-proof-check}, uzyskując oddzielny dowód
\[
T\vdash\neg\operatorname{Prf}_T(\overline p,\overline g).
\]
Gdyby jednak $T\vdash\neg G_T$, to z dowodliwej
równoważności dla $G_T$ i klasycznej eliminacji podwójnej
negacji wynikałoby
$T\vdash\operatorname{Bew}_T(\overline g)$, czyli
\[
T\vdash\exists u\,\operatorname{Prf}_T(u,\overline g).
\]
To przeczy $\omega$-niesprzeczności: teoria dowodziłaby
istnienia pewnego $u$ z daną własnością, a jednocześnie
wykluczałaby każdy standardowy numerał jako świadka.
Wniosek o niezupełności jest natychmiastowy z definicji.
\end{proof}

\begin{uwaga}[Czy $G_T$ jest prawdziwe?]\label{rem:godel-true}
Jeśli $T$ jest niesprzeczna, punkt (a) pokazuje, że
$G_T$ nie ma żadnego rzeczywistego dowodu w $T$.
Dlatego w modelu standardowym $\mathfrak N$ zdanie
$\operatorname{Bew}_T(\overline g)$ jest fałszywe.
Ponieważ $\mathfrak N\models Q$, obowiązuje też w nim
dowodliwa równoważność dla $G_T$, a więc
$\mathfrak N\models G_T$.
Nie założyliśmy tutaj, że wszystkie aksjomaty $T$ są
prawdziwe w $\mathfrak N$: wystarczy niesprzeczność,
aby ustalić prawdziwość \emph{tego konkretnego} zdania.
\end{uwaga}

\begin{przyklad}[Dlaczego nie wystarcza seria osobnych sprawdzeń?]
Dla każdej liczby $p$ możemy po skończonym obliczeniu
sprawdzić, czy jest kodem dowodu $G_T$.
Z niesprzeczności wynika, że każdy taki test daje
odpowiedź przeczącą. Są to jednak wyniki dla
$0,1,2,\ldots$ rozpatrywanych oddzielnie.
Nie wolno stąd bez dodatkowego dowodu wnioskować,
że $T\vdash\forall p\,\neg\operatorname{Prf}_T(p,\overline g)$.
Właśnie ta luka między \emph{wszystkimi konkretnymi testami}
a \emph{jednym zdaniem ogólnym} tłumaczy, dlaczego
punkt (b) wymaga silniejszego założenia.
\end{przyklad}

\section{Rosser: jak osłabić założenie do zwykłej niesprzeczności?}
Idea polega na porównaniu \emph{dwóch} rodzajów dowodów:
dowodów rozważanego zdania i dowodów jego negacji.
Rosser nie pyta wyłącznie, czy istnieje dowód zdania.
Pyta, czy istnieje dowód tego zdania bez żadnego
\emph{wcześniejszego pod względem kodu} dowodu negacji.

\begin{definicja}[Relacja Rossera]\label{def:rosser-predicate}
Niech $T\supseteq Q$ będzie efektywnie aksjomatyzowalną
teorią w efektywnie kodowanym języku zawierającym
$\mathcal L_{\mathrm{ar}}$.
Funkcja $\operatorname{neg}(n)$, zwracająca kod negacji
formuły o kodzie $n$ (a na niepoprawnych kodach ustaloną
wartość zastępczą), jest pierwotnie rekurencyjna.
Wybierzmy formułę $N(x,z)$ \emph{silnie reprezentującą}
tę funkcję w $Q$ w sensie Definicji~\ref{res:7:2}.
Definiujemy formułę
\[
\operatorname{RosPr}_T(x):=
\exists z\,\exists p\,
\left(N(x,z)\land\operatorname{Prf}_T(p,x)
\land\forall q<p\,\neg\operatorname{Prf}_T(q,z)\right).
\]
Kwantyfikator $\forall q<p$ jest ograniczony:
przebiega tylko po liczbach mniejszych od $p$.
\end{definicja}

\begin{uwaga}[Co daje reprezentacja negacji?]
Jeśli $r$ jest rzeczywistym kodem zdania $R$,
a $r^-:=\operatorname{neg}(r)=\ulcorner\neg R\urcorner$,
to $Q$ dowodzi
$\forall z\,(N(\overline r,z)\leftrightarrow z=\overline{r^-})$.
Dlatego instancja $\operatorname{RosPr}_T(\overline r)$
jest w $Q$ równoważna formułce
\[
\exists p\,\bigl(\operatorname{Prf}_T(p,\overline r)
\land\forall q<p\,\neg\operatorname{Prf}_T(q,\overline{r^-})\bigr).
\]
Jest to równoważność dowodliwa w $Q$, nie tylko
równoważność zachodząca w modelu standardowym:
silna reprezentowalność daje dla ustalonego $r$
zarówno istnienie właściwego $z$, jak i wykluczenie
wszystkich innych wartości $z$. Symbol
$\operatorname{neg}$ pozostaje wyłącznie w metajęzyku.
\end{uwaga}

\begin{twierdzenie}[Wzmocnienie Rossera]\label{thm:rosser-first}
Niech $T\supseteq Q$ będzie efektywnie aksjomatyzowalną,
niesprzeczną teorią klasycznej logiki pierwszego rzędu
w efektywnie kodowanym języku zawierającym
$\mathcal L_{\mathrm{ar}}$.
Istnieje zdanie arytmetyczne $R_T$ takie, że
$T\nvdash R_T$ i $T\nvdash\neg R_T$.
Do tego wniosku nie jest potrzebna $\omega$-niesprzeczność.
\end{twierdzenie}
\begin{proof}
Stosujemy Lemat~\ref{res:8:1} do
$\neg\operatorname{RosPr}_T(x)$.
Otrzymujemy zdanie $R_T$ z
$Q\vdash R_T\leftrightarrow\neg\operatorname{RosPr}_T(\overline r)$,
gdzie $r=\ulcorner R_T\urcorner$; niech
$r^-=\ulcorner\neg R_T\urcorner$.
Z poprzedniej uwagi dostajemy dowodliwą w $Q$
równoważność
\begin{equation}\label{eq:rosser-instance}
R_T\leftrightarrow\neg\exists p\,
\left(\operatorname{Prf}_T(p,\overline r)\land
 \forall q<p\,\neg\operatorname{Prf}_T(q,\overline{r^-})\right).
\end{equation}

\emph{Pierwsza możliwość:} przypuśćmy $T\vdash R_T$.
Niech $a$ koduje rzeczywisty dowód $R_T$.
Z niesprzeczności wynika $T\nvdash\neg R_T$,
więc żaden kod $q<a$ nie koduje dowodu $\neg R_T$.
Dla każdego z \emph{faktycznie skończenie wielu} $q<a$
Lemat~\ref{lem:incomp-proof-check} daje
$Q\vdash\neg\operatorname{Prf}_T(\overline q,\overline{r^-})$.
Na mocy Lematu~\ref{res:4:4}, stosowanego do ustalonego
standardowego $a$, mamy w $Q$ równoważność
$q<\overline a$ ze skończoną alternatywą
$q=\overline0\lor\cdots\lor q=\overline{a-1}$.
W każdej jej gałęzi zastępujemy $q$ odpowiednim numerałem
i korzystamy z jednej z uzyskanych negacji
$\operatorname{Prf}_T(\overline q,\overline{r^-})$.
Generalizacja po $q$ daje
$Q\vdash\forall q<\overline a\,\neg\operatorname{Prf}_T(q,\overline{r^-})$.
Gdy $a=0$, alternatywa jest pusta, a wynik pochodzi
z $Q\vdash\forall q\,\neg(q<\overline0)$.
Ponadto $Q\vdash\operatorname{Prf}_T(\overline a,\overline r)$.
Silna reprezentowalność funkcji $\operatorname{neg}$ daje
$Q\vdash N(\overline r,\overline{r^-})$.
Wstawiając do definicji Rossera świadków
$z=\overline{r^-}$ i $p=\overline a$, otrzymujemy
$Q\vdash\operatorname{RosPr}_T(\overline r)$.
Równoważność~\eqref{eq:rosser-instance} wraz z
$T\vdash R_T$ daje zarazem
$T\vdash\neg\operatorname{RosPr}_T(\overline r)$.
Teoria byłaby sprzeczna, zatem $T\nvdash R_T$.

\emph{Druga możliwość:} przypuśćmy $T\vdash\neg R_T$.
Niech $b$ koduje rzeczywisty dowód $\neg R_T$.
Pierwszy przypadek wykazał $T\nvdash R_T$,
więc żaden kod $p$ nie koduje dowodu $R_T$.
W szczególności dla każdej standardowej liczby $p\le b$
teoria $Q$ rozstrzyga negatywnie konkretny test
$\operatorname{Prf}_T(\overline p,\overline r)$.
Ponieważ $p\le\overline b$ jest skrótem
$p<\overline b\lor p=\overline b$, Lemat~\ref{res:4:4}
rozwija tę przesłankę w $Q$ na skończoną alternatywę
$p=\overline0\lor\cdots\lor p=\overline b$.
Każdy jej składnik daje jeden z wcześniej
udowodnionych przypadków negatywnych. Po generalizacji:
\[
Q\vdash\forall p\le\overline b\,
\neg\operatorname{Prf}_T(p,\overline r).
\]
Rzeczywisty dowód negacji o kodzie $b$ daje natomiast,
na mocy Lematu~\ref{lem:incomp-proof-check},
$Q\vdash\operatorname{Prf}_T(\overline b,\overline{r^-})$.

Teraz dowodzimy w $Q$ negacji całego warunku
pod kwantyfikatorem egzystencjalnym w~\eqref{eq:rosser-instance}.
Niech $p$ będzie dowolne. Z Lematu~\ref{res:4:4}
wynika rozbicie na dwie sytuacje
$p\le\overline b$ albo $\overline b<p$.
W pierwszej sytuacji przeczy mu pierwszy czynnik:
$\neg\operatorname{Prf}_T(p,\overline r)$.
W drugiej mamy konkretnego świadka $q=\overline b$
dla zaprzeczenia drugiego czynnika, gdyż
$q<p$ i $\operatorname{Prf}_T(q,\overline{r^-})$.
Formalnie dla dowolnego $p$ w $Q$ otrzymujemy
\[
\operatorname{Prf}_T(p,\overline r)
\longrightarrow
\exists q<p\,\operatorname{Prf}_T(q,\overline{r^-}).
\]
W pierwszej gałęzi poprzednik jest fałszywy; w drugiej
świadkiem następnika jest $q=\overline b$. Zatem $Q$ dowodzi
negacji koniunkcji
$\operatorname{Prf}_T(p,\overline r)\land
\forall q<p\,\neg\operatorname{Prf}_T(q,\overline{r^-})$
dla każdego $p$. Po generalizacji i zanegowaniu kwantyfikatora
istnienia dostajemy
\[
Q\vdash\neg\exists p\,\bigl(
\operatorname{Prf}_T(p,\overline r)\land
\forall q<p\,\neg\operatorname{Prf}_T(q,\overline{r^-})\bigr).
\]
Dopiero teraz, korzystając z dowodliwej w $Q$ eliminacji
zmiennej $z$ za pomocą formuły $N$ (poprzednia uwaga),
otrzymujemy $Q\vdash\neg\operatorname{RosPr}_T(\overline r)$.
Stosując~\eqref{eq:rosser-instance}, uzyskujemy
$Q\vdash R_T$, a więc $T\vdash R_T$.
Jest to sprzeczne z $T\vdash\neg R_T$ i niesprzecznością.
Zatem $T\nvdash\neg R_T$.
\end{proof}

\begin{uwaga}[Dlaczego dowód Rossera jest mocniejszy?]
W dowodzie z pierwszej części użyliśmy jedynie \emph{skończenie wielu}
ujemnych testów dla $q<a$. W drugiej części wystarczyła
skończona kontrola $p\le b$ oraz \emph{jeden} rzeczywisty dowód
negacji o kodzie $b$, który obala pozostałe kandydatury.
Nie potrzebowaliśmy twierdzenia $T\vdash\forall p\,\neg\operatorname{Prf}_T(p,\overline r)$
ani $\omega$-niesprzeczności. Samo zdanie Rossera nie musi
mieć tak prostej interpretacji słownej jak zdanie Gödla.
\end{uwaga}

\section{Dwie konsekwencje obliczeniowe}
\begin{wniosek}[Zbiór zdań prawdziwych nie jest efektywnie wyliczalny]
Zbiór kodów wszystkich zdań języka arytmetyki prawdziwych
w standardowej strukturze $\mathfrak N$ nie jest rekurencyjnie przeliczalny.
\end{wniosek}
\begin{proof}
Gdyby kody prawdziwych zdań można było efektywnie wyliczać,
przyjęlibyśmy te zdania jako aksjomaty teorii $T$.
Byłaby to efektywnie aksjomatyzowalna teoria w ustalonym,
efektywnie kodowanym języku arytmetyki.
Struktura $\mathfrak N$ byłaby modelem $T$, zatem $T$
byłaby niesprzeczna. Ponadto dla każdego zdania $\sigma$
jedno ze zdań $\sigma,\neg\sigma$ jest prawdziwe w $\mathfrak N$,
więc należałoby do aksjomatów: teoria byłaby zupełna.
Ponieważ wszystkie aksjomaty $Q$ są prawdziwe w $\mathfrak N$,
$T$ zawierałaby $Q$. Przeczyłoby to Twierdzeniu~\ref{thm:rosser-first}.
\end{proof}

\begin{wniosek}[Zbiór zdań dowodliwych jest nierozstrzygalny]\label{cor:undecidability-theorems}
Dla każdej niesprzecznej i efektywnie aksjomatyzowalnej
teorii $T\supseteq Q$ zbiór kodów jej twierdzeń nie jest rozstrzygalny.
\end{wniosek}
\begin{proof}
Przypuśćmy, że istnieje algorytm rozstrzygający,
czy dane zdanie $\sigma$ jest dowodliwe w $T$.
Ponieważ język $T$ jest efektywnie kodowany,
algorytmicznie rozpoznajemy poprawne kody zdań.
Definiujemy zatem totalną relację rekurencyjną $D(n)$:
ma ona wartość prawdziwą dokładnie wtedy, gdy $n$
jest kodem zdania dowodliwego w $T$; dla kodów
niepoprawnych jest fałszywa.
Z Twierdzenia~\ref{res:7:6} istnieje formuła arytmetyczna
$\rho_D(x)$, dla której $Q$ dowodzi
$\rho_D(\overline n)$, gdy $D(n)$ zachodzi, oraz
$\neg\rho_D(\overline n)$, gdy $D(n)$ nie zachodzi.
Użycie twierdzenia jest uprawnione właśnie dlatego,
że $D$ jest totalna i rozstrzygalna.
Lemat diagonalny daje zdanie $H$ takie, że
\[
Q\vdash H\leftrightarrow\neg\rho_D(\overline h),
\qquad h=\ulcorner H\urcorner.
\]
Jeśli $T\vdash H$, to $D(h)$ zachodzi, więc
$Q\vdash\rho_D(\overline h)$, a z
$T\vdash H$ i dowodliwej równoważności wynika
$T\vdash\neg\rho_D(\overline h)$. Niesprzeczność
$T\supseteq Q$ wyklucza tę możliwość.
Jeśli natomiast $T\nvdash H$, to $D(h)$ nie zachodzi,
więc $Q\vdash\neg\rho_D(\overline h)$.
Z równoważności $Q\vdash H\leftrightarrow
\neg\rho_D(\overline h)$ wnioskujemy
$T\vdash H$, wbrew założeniu. Obie możliwości
prowadzą do sprzeczności.
\end{proof}

\begin{uwaga}[Czego nie dowiedliśmy]
Pierwsze twierdzenie nie twierdzi, że \emph{każde}
zdanie nierozstrzygalne w $T$ jest prawdziwe w $\mathfrak N$,
ani że istnieje algorytm rozpoznający prawdę zdań.
Nie wyprowadza też jeszcze tego, że $T$ nie potrafi
udowodnić własnej niesprzeczności: to osobne twierdzenie,
do którego potrzebujemy dodatkowej analizy predykatu dowodliwości.
\end{uwaga}


\section{Porządek kodów dowodów w argumencie Rossera}
W formule Rossera ograniczenie $q<p$ dotyczy \emph{kodów dowodów},
a nie długości zdań ani numerów formuł. Przy ustalonym standardowym
$p$ relacja $q<p$ obejmuje skończenie wiele liczb standardowych;
$Q$ potrafi rozstrzygnąć każdą z odpowiednich instancji
predykatu dowodu dzięki jego reprezentacji $\Delta_0$
(Lemat~\ref{lem:incomp-proof-check}). Nie wolno natomiast wyciągać z tego
wniosku, że $Q$ dowodzi ogólnego zdania o wszystkich kodach.
W szczególności przejście od metamatematycznego braku dowodu do
wewnętrznego zdania o braku dowodów wymaga osobnego argumentu;
w dowodzie Rossera używa się konkretnych świadków i skończonych
ograniczeń, nie zaś zasady refleksji $\Box\varphi\to\varphi$.

\section{Co naprawdę znaczy niezależność zdania?}
\begin{uwaga}[Prawdziwość, niedowodliwość i rozstrzygnięcie przez nowy aksjomat]
Niech $T$ będzie niesprzeczną teorią, a $\rho$ zdaniem niezależnym od $T$,
czyli $T\nvdash\rho$ oraz $T\nvdash\neg\rho$.
Nie oznacza to, że zdanie $\rho$ jest fałszywe, ani że nie ma żadnego modelu,
w którym zachodzi. Przeciwnie, przez twierdzenie Gödla o pełności
(tego nie mylić z jego twierdzeniem o niezupełności)
teorie $T+\rho$ i $T+\neg\rho$ mają modele:
gdyby pierwsza była sprzeczna, to z twierdzenia o dedukcji
uzyskalibyśmy $T\vdash\neg\rho$; analogicznie dla drugiej
(w drugim przypadku stosujemy także klasyczną negację podwójną).

Możemy dołączyć $\rho$ jako nowy aksjomat i wtedy rozstrzygniemy ten jeden przypadek.
Jednak każde niesprzeczne efektywnie aksjomatyzowalne rozszerzenie $Q$
pozostaje niezupełne, więc znajdzie się inne zdanie niezależne.
Nawet nieskończone efektywne dokładanie aksjomatów nie omija tego ograniczenia,
jeżeli ostateczna teoria wciąż jest niesprzeczna i zawiera $Q$.
Zbiór wszystkich prawdziwych zdań standardowej arytmetyki
jest wprawdzie zupełny, lecz nie jest efektywnie aksjomatyzowalny.
\end{uwaga}

\begin{przyklad}[Dwa różne modele tej samej teorii]
Jeżeli $T$ nie rozstrzyga $\rho$, istnieje model $\mathfrak A\models T+\rho$
i model $\mathfrak B\models T+\neg\rho$.
Modele te nie są izomorficzne, ponieważ zdania pierwszego rzędu
mają tę samą wartość w strukturach izomorficznych.
To nie oznacza sprzeczności w $T$: teoria nie określiła,
która z dwóch sytuacji musi występować w każdym jej modelu.
\end{przyklad}
\chapter{Prawda arytmetyczna i twierdzenie Tarskiego}\label{ch:tarski-truth}

\section{Prawda to nie to samo co dowodliwość}
W rozdziale~\ref{ch:first-incompleteness} udowodniliśmy,
że dostatecznie silna, niesprzeczna, efektywnie
aksjomatyzowalna teoria arytmetyczna nie rozstrzyga
wszystkich zdań. Możemy jednak zapytać o jeszcze
silniejszy pomysł: czy istnieje \emph{jedna formuła
arytmetyczna}, która dla każdej liczby kodującej
zdanie rozpoznaje, czy to zdanie jest prawdziwe
w standardowej strukturze liczb naturalnych?

Odpowiedź jest przecząca. Zanim ją udowodnimy,
oddzielmy trzy poziomy języka:
\begin{itemize}
\item Zdanie arytmetyczne $\sigma$ jest \emph{napisem}.
\item Liczba $s=\ulcorner\sigma\urcorner$ jest jego kodem,
      a $\overline s$ termem będącym numerałem kodu.
\item Zapis $\mathfrak N\models\sigma$ jest
      \emph{metamatematycznym} stwierdzeniem prawdziwości
      tego zdania w zwykłych liczbach naturalnych.
\end{itemize}
Sama notacja $\mathfrak N\models\sigma$ nie jest dodatkowym
predykatem języka $\{0,S,+,\cdot,<\}$.

\begin{przyklad}[Dwie wypowiedzi o liczbie $4$]
Zdanie $\overline2+\overline2=\overline4$ jest prawdziwe
w $\mathfrak N$, czyli
$\mathfrak N\models(\overline2+\overline2=\overline4)$.
Niech $T\supseteq Q$ będzie efektywnie aksjomatyzowalną
teorią i niech $g$ oznacza kod tego zdania. Inne zdanie,
$\operatorname{Bew}_T(\overline g)$, mówi o istnieniu
\emph{dowodu} danej równości w $T$.
Ponieważ $Q$ dowodzi tej równości, istnieje jej rzeczywisty,
skończony dowód w $T$. Zatem oba zdania są prawdziwe
w $\mathfrak N$, niezależnie od tego, czy $T$ jest poprawna
semantycznie. Poprzedni rozdział pokazał jednak, że taka
zgodność nie zachodzi w ogólności dla wszystkich zdań
arytmetycznych.
\end{przyklad}

\section{Co miałoby znaczyć „formuła definiująca prawdę”?}
\begin{definicja}[Definiowalność prawdy arytmetycznej]
Formuła $\operatorname{Tr}(x)$ języka arytmetyki,
w której jedyną zmienną mogącą występować wolno jest $x$,
\emph{definiowałaby prawdę arytmetyczną w $\mathfrak N$},
gdyby dla każdego \emph{zdania} $\sigma$ tego języka
zachodziła równoważność metamatematyczna
\[
\mathfrak N\models\operatorname{Tr}(\overline{\ulcorner\sigma\urcorner})
\quad\Longleftrightarrow\quad
\mathfrak N\models\sigma.
\]
Nie narzucamy żadnego wymagania na wartości
$\operatorname{Tr}(x)$ dla liczb niekodujących zdań.
Równoważność dotyczy więc wyłącznie kodów \emph{zdań},
a nie dowolnych napisów lub formuł z wolnymi zmiennymi.
\end{definicja}

\begin{uwaga}[Dwa różne rodzaje równoważności]
Jeśli $s=\ulcorner\sigma\urcorner$, zdanie
$\sigma\leftrightarrow\operatorname{Tr}(\overline s)$
można \emph{zapisać} w języku arytmetyki.
Natomiast stwierdzenie
$\mathfrak N\models(\sigma\leftrightarrow\operatorname{Tr}(\overline s))$
jest wypowiedzią metamatematyczną o prawdziwości tej równoważności.
Teoria $T$ może również \emph{dowodzić}
$T\vdash\sigma\leftrightarrow\operatorname{Tr}(\overline s)$.
Są to dwa różne wymagania. Dopiero gdy wszystkie twierdzenia $T$
są prawdziwe w $\mathfrak N$ (czyli gdy $T$ jest poprawna
względem $\mathfrak N$), dowodliwość pociąga za sobą
prawdziwość w $\mathfrak N$. W ogólności takiego przejścia
nie wolno wykonywać.
\end{uwaga}

\section{Twierdzenie Tarskiego}
\begin{twierdzenie}[Niedefiniowalność prawdy --- wersja semantyczna]\label{thm:tarski-semantic}
Nie istnieje formuła $\operatorname{Tr}(x)$ języka
$\mathcal L_{\mathrm{ar}}=\{0,S,+,\cdot,<\}$,
która definiuje prawdę wszystkich zdań tego języka
w standardowej strukturze $\mathfrak N$.
\end{twierdzenie}
\begin{proof}
Przypuśćmy przeciwnie, że taka formuła istnieje.
Stosujemy lemat diagonalny
(Lemat~\ref{res:8:1}) do formuły
$\psi(x):=\neg\operatorname{Tr}(x)$.
Dostajemy zdanie arytmetyczne $L$ takie, że
\[
Q\vdash L\leftrightarrow\neg\operatorname{Tr}(\overline\ell),
\qquad \ell:=\ulcorner L\urcorner.
\]
Ponieważ $\mathfrak N$ jest modelem $Q$, w modelu
standardowym zachodzi
\begin{equation}\label{eq:tarski-anti}
\mathfrak N\models
L\leftrightarrow\neg\operatorname{Tr}(\overline\ell).
\end{equation}
Z założenia, że $\operatorname{Tr}$ definiuje prawdę,
dla \emph{tego samego} zdania $L$ dostajemy również
\begin{equation}\label{eq:tarski-truth}
\mathfrak N\models
L\leftrightarrow\operatorname{Tr}(\overline\ell).
\end{equation}

W klasycznej logice prawdziwość~\eqref{eq:tarski-truth}
oznacza, że $L$ i $\operatorname{Tr}(\overline\ell)$
mają tę samą wartość logiczną. Tymczasem
\eqref{eq:tarski-anti} mówi, że wartości te są przeciwne.
Jeśli $L$ jest prawdziwe, to z~\eqref{eq:tarski-anti}
wynika fałszywość $\operatorname{Tr}(\overline\ell)$,
a z~\eqref{eq:tarski-truth} jej prawdziwość.
Jeśli $L$ jest fałszywe, z pierwszej równoważności
wynika prawdziwość $\operatorname{Tr}(\overline\ell)$,
a z drugiej jej fałszywość. Oba przypadki są niemożliwe.
Zatem formuła $\operatorname{Tr}$ nie istnieje.
\end{proof}

\begin{uwaga}[Dlaczego to nie jest zwykły paradoks kłamcy?]
Nie wpisaliśmy do języka arytmetycznego nieformalnego
zdania „to zdanie jest fałszywe”.
Najpierw kodujemy skończone napisy, następnie
konstruujemy funkcję podstawienia na kodach,
a na koniec używamy \emph{udowodnionego}
Lematu~\ref{res:8:1}. Zdanie $L$ i jego kod $\ell$
to różne obiekty. Sprzeczność powstaje dopiero
po dołożeniu hipotezy, że jedna wewnętrzna formuła
rozpoznaje prawdę \emph{wszystkich} zdań.
\end{uwaga}

\begin{twierdzenie}[Wersja syntaktyczna Tarskiego]\label{thm:tarski-syntactic}
Niech $T$ będzie dowolną niesprzeczną teorią
klasycznej logiki pierwszego rzędu zawierającą $Q$.
Nie istnieje formuła arytmetyczna $\operatorname{Tr}(x)$
taka, że dla \emph{każdego} zdania arytmetycznego
$\sigma$ teoria $T$ dowodzi
\[
\sigma\leftrightarrow
\operatorname{Tr}(\overline{\ulcorner\sigma\urcorner}).
\]
Nie zakładamy efektywnego wyliczania aksjomatów $T$.
\end{twierdzenie}
\begin{proof}
Załóżmy, że wszystkie takie równoważności są dowodliwe.
Dla $\psi(x)=\neg\operatorname{Tr}(x)$ lemat diagonalny
daje zdanie $L$ i liczbę $\ell=\ulcorner L\urcorner$ z
\[
T\vdash L\leftrightarrow\neg\operatorname{Tr}(\overline\ell),
\]
ponieważ $T\supseteq Q$. Założenie zastosowane
\emph{do zdania $L$} daje jednocześnie
$T\vdash L\leftrightarrow\operatorname{Tr}(\overline\ell)$.
Złożenie obu równoważności daje
$T\vdash L\leftrightarrow\neg L$.
Ten krok jest czysto syntaktyczny: nie zakłada, że zdania
dowodliwe w $T$ są prawdziwe w $\mathfrak N$.
Wyprowadźmy jawnie sprzeczność.
Z kierunku $L\to\neg L$ i klasycznego rozumowania
przez sprzeczność otrzymujemy $T\vdash\neg L$.
Z drugiego kierunku $\neg L\to L$ otrzymujemy następnie
$T\vdash L$. Zatem $T$ dowodzi zdania i jego negacji,
co przeczy niesprzeczności.
\end{proof}

\begin{przyklad}[Dlaczego problem dotyczy jednego uniwersalnego predykatu?]
Dla ustalonej skończonej, niepustej listy zdań
$\sigma_1,\ldots,\sigma_n$ (gdzie $n\ge1$)
możemy ustalić kody $g_i=\ulcorner\sigma_i\urcorner$ i zapisać
\[
\operatorname{Tr}_{\mathrm{fin}}(x):=
\bigvee_{i=1}^{n}(x=\overline{g_i}\land\sigma_i).
\]
Dla każdego $i\in\{1,\ldots,n\}$ w $\mathfrak N$ zachodzi
$\operatorname{Tr}_{\mathrm{fin}}(\overline{g_i})
\leftrightarrow\sigma_i$.
Jeżeli na liście powtórzono to samo zdanie, jego powtórzenia
mają identyczną wartość logiczną. Natomiast różne zdania
mają różne kody dzięki różnowartościowości numeracji;
dla danego $g_i$ wszystkie człony o innych kodach są fałszywe.
Nie otrzymaliśmy jednak \emph{jednej} formuły
obsługującej równocześnie wszystkie zdania arytmetyczne:
po dodaniu nowych zdań do listy musimy zmienić formułę.
Twierdzenie Tarskiego nie przeczy takiej skończonej konstrukcji.
\end{przyklad}

\section{Dlaczego prawda nie jest tym samym co dowodliwość?}
\begin{twierdzenie}[Niesprzeczna teoria nie dowodzi równoważności prawdy i dowodliwości]
Niech $T\supseteq Q$ będzie niesprzeczna i efektywnie
aksjomatyzowalna. Wówczas nie jest prawdą, że dla każdego
zdania $\sigma$ zachodzi
$T\vdash\sigma\leftrightarrow
\operatorname{Bew}_T(\overline{\ulcorner\sigma\urcorner})$.
\end{twierdzenie}
\begin{proof}
Gdyby wszystkie te równoważności były dowodliwe,
formuła $\operatorname{Tr}(x):=\operatorname{Bew}_T(x)$
spełniałaby warunek wykluczony w
Twierdzeniu~\ref{thm:tarski-syntactic}. Ponieważ $T$
jest niesprzeczna i zawiera $Q$, otrzymujemy sprzeczność.
\end{proof}

\begin{uwaga}[Dwa rodzaje ograniczeń]
Pierwsze twierdzenie Gödla mówi o granicach
\emph{dowodzenia} w pojedynczej efektywnie aksjomatyzowalnej,
niesprzecznej teorii rozszerzającej $Q$.
Semantyczne twierdzenie Tarskiego mówi o niemożliwości
\emph{zdefiniowania pełnej prawdy} arytmetycznej
jedną formułą tego samego języka i nie wymaga wyboru
teorii $T$. W jego wersji syntaktycznej można wybrać
dowolną niesprzeczną teorię $T\supseteq Q$,
bez założenia efektywnej aksjomatyzowalności. Nie są to identyczne twierdzenia,
choć oba korzystają z tego samego mechanizmu diagonalizacji.
\end{uwaga}

\section{Zapowiedź warunków wyprowadzalności}
Do badania drugiego twierdzenia o niezupełności
nie wystarczy znajomość tego, które konkretne
instancje predykatu $\operatorname{Prf}_T$
są prawdziwe w $\mathfrak N$.
Potrzebujemy formalnych relacji między zdaniami
mówiącymi o \emph{istnieniu dowodów}.

\begin{definicja}[Skrót $\Box\varphi$]\label{def:box-provability}
Ustalmy efektywnie aksjomatyzowalną teorię
$T\supseteq\mathrm{PA}$ i jej predykat
$\operatorname{Bew}_T(x)$.
Dla \emph{zdania} arytmetycznego $\varphi$ oznaczmy
$g=\ulcorner\varphi\urcorner$.
Zapis $\Box\varphi$ jest skrótem pełnego zdania
$\operatorname{Bew}_T(\overline g)$.
Jeśli chcemy zapisać $\Box\Box\varphi$, najpierw
tworzymy zdanie $\Box\varphi$, potem bierzemy
\emph{jego} kod $h=\ulcorner\Box\varphi\urcorner$
i zapisujemy $\operatorname{Bew}_T(\overline h)$.
Nie podnosimy liczby $g$ do kwadratu: drugi znak $\Box$
oznacza nowy poziom wypowiedzi o dowodzie.
Jest to skrót \emph{metamatematyczny}, a nie nowy
spójnik dopisany do alfabetu arytmetyki.
\end{definicja}

\begin{definicja}[Warunki Hilberta--Bernaysa--Löba]\label{def:hilbert-bernays-loeb}
Niech $T\supseteq\mathrm{PA}$ będzie efektywnie
aksjomatyzowalną teorią klasycznej logiki pierwszego rzędu,
a $\Box$ ustalonym predykatem dowodliwości,
rozumianym jak w poprzedniej definicji.
Mówimy, że spełnia on następujące warunki, jeśli dla
wszystkich zdań arytmetycznych $\varphi,\psi$ zachodzą:
\begin{align*}
\textup{(D1)}\quad&T\vdash\varphi\ \Longrightarrow\ T\vdash\Box\varphi,\\
\textup{(D2)}\quad&T\vdash\Box(\varphi\to\psi)\to(\Box\varphi\to\Box\psi),\\
\textup{(D3)}\quad&T\vdash\Box\varphi\to\Box\Box\varphi.
\end{align*}
Pierwsza linia jest regułą o \emph{rzeczywistych dowodach};
pozostałe dwie linie są \emph{formułami dowodliwymi w $T$}.
\end{definicja}

\begin{lemat}[Warunek (D1) dla standardowego predykatu]\label{lem:d1-provability}
Niech $T\supseteq\mathrm{PA}$ będzie efektywnie
aksjomatyzowalna, a $\Box$ pochodzi z predykatu
$\operatorname{Prf}_T$ z Lematu~\ref{res:7:7}.
Wtedy zachodzi warunek \textup{(D1)}.
\end{lemat}
\begin{proof}
Jeżeli $T\vdash\varphi$, wybieramy jeden
rzeczywisty kod $p$ dowodu $\varphi$.
Niech $g=\ulcorner\varphi\urcorner$.
Ponieważ $p$ jest rzeczywistym kodem poprawnego dowodu,
zdanie $\operatorname{Prf}_T(\overline p,\overline g)$
jest prawdziwą zamkniętą instancją formuły $\Delta_0$.
Z części (i) Lematu~\ref{lem:incomp-proof-check}
wynika więc $Q\vdash\operatorname{Prf}_T(\overline p,\overline g)$.
Wprowadzając kwantyfikator egzystencjalny,
dostajemy $Q\vdash\Box\varphi$.
Ponieważ $T\supseteq\mathrm{PA}\supseteq Q$,
ten sam wniosek jest dowodliwy w $T$.
\end{proof}

\begin{przyklad}[Dlaczego (D1) nie działa w drugą stronę?]
Warunek (D1) mówi: jeśli \emph{mamy konkretny dowód}
$\varphi$, to możemy w teorii udowodnić,
że jakiś dowód $\varphi$ istnieje. Nie mówi,
że z $T\vdash\Box\varphi$ zawsze wolno wyciągnąć
wniosek $T\vdash\varphi$.
Z samego dowodu zdania egzystencjalnego
$\exists p\,\operatorname{Prf}_T(p,\overline g)$
nie wynika możliwość wskazania \emph{standardowej liczby
naturalnej} kodującej dowód $\varphi$: w niestandardowym
modelu arytmetyki świadkiem tej formuły może być element,
który nie jest kodem rzeczywistego, skończonego dowodu.
Dodatkowe założenie może zmienić sytuację. Jeśli $T$ jest
$\Sigma_1$-poprawna względem $\mathfrak N$ (to znaczy:
każde zdanie $\Sigma_1$ dowodliwe w $T$ jest prawdziwe
w $\mathfrak N$), to z
$T\vdash\Box\varphi$ wynika
$\mathfrak N\models\Box\varphi$, gdyż
$\Box\varphi$ jest zdaniem $\Sigma_1$.
Wtedy istnieje standardowa liczba $p$, dla której
$\mathfrak N\models\operatorname{Prf}_T(\overline p,\overline g)$,
a więc rzeczywisty dowód $\varphi$ w $T$.
Takiej $\Sigma_1$-poprawności nie zakłada jednak (D1);
jego odwrócenie nie wynika z samego warunku wyprowadzalności.
\end{przyklad}

\begin{uwaga}[Co pozostaje do udowodnienia?]
Reguły \textup{(D2)} i \textup{(D3)} wymagają
\emph{wewnętrznego} kodowania operacji łączenia dowodów
i weryfikacji kodu dowodu.
Nie wynika z (D1), że są automatycznie spełnione.
W rozdziale~\ref{ch:derivability-proof} udowodnimy je
dla odpowiedniego standardowego predykatu w arytmetyce
Peana, zamiast przyjmować je bez uzasadnienia.
\end{uwaga}

\section{Jak czytać trzy warunki na konkretnym przykładzie?}
\begin{przyklad}[Od rachunku $1+1=2$ do zdania o dowodzie]
Niech $\varphi$ będzie zdaniem
$S(0)+S(0)=S(S(0))$ i niech $T\supseteq\mathrm{PA}$
będzie efektywnie aksjomatyzowalna. Z aksjomatów
dodawania otrzymujemy kolejno
\[
S(0)+S(0)=S(S(0)+0)=S(S(0)),
\]
a więc $T\vdash\varphi$. Możemy wybrać kod $p$
\emph{jednego konkretnego} dowodu tej równości.
Weryfikacja $p$ daje
$Q\vdash\operatorname{Prf}_T(\overline p,\overline g)$,
gdzie $g=\ulcorner\varphi\urcorner$.
Po wprowadzeniu kwantyfikatora egzystencjalnego
uzyskujemy $T\vdash\Box\varphi$.
Tak działa \textup{(D1)}. Nie używaliśmy tu wniosku,
że z istnienia kodu dowodu wynika prawdziwość dowodzonego
zdania w każdym możliwym modelu; posłużyliśmy się
\emph{rzeczywistym dowodem}, który już mieliśmy.
\end{przyklad}

\begin{uwaga}[Dlaczego (D2) jest formalizacją modus ponens?]
Reguła \emph{modus ponens} pozwala z dowodów
$\varphi\to\psi$ oraz $\varphi$ otrzymać dowód $\psi$.
Na poziomie kodów odpowiada jej obliczalna operacja,
która skleja dwa skończone dowody i dopisuje
ostatni wiersz. Warunek \textup{(D2)} mówi,
że teoria $T$ potrafi uzasadnić istnienie
odpowiedniego kodu \emph{wewnątrz arytmetyki}:
\[
\Box(\varphi\to\psi)\to(\Box\varphi\to\Box\psi).
\]
Zwróćmy uwagę na nawiasy: jest to \emph{jedno zdanie}
$T$ i nie należy go mylić z regułą o \emph{dwóch
gotowych dowodach}.
\end{uwaga}

\begin{uwaga}[Dlaczego (D3) nie znaczy, że teoria jest nieomylna?]
Zdanie $\Box\varphi$ głosi, że istnieje kod
dowodu $\varphi$. Warunek \textup{(D3)} głosi,
że z takiego istnienia w teorii wynika istnienie
kodu dowodu \emph{tego właśnie zdania o dowodzie}:
\[
\Box\varphi\to\Box\Box\varphi.
\]
Nie jest to implikacja $\Box\varphi\to\varphi$.
Ta ostatnia byłaby \emph{wewnętrznym stwierdzeniem
poprawności dowodów} $\varphi$, a nie jedynie zdaniem
o kolejnych dowodach. Nie wolno jej utożsamiać
z metamatematycznym warunkiem prawdziwości $\varphi$
w $\mathfrak N$.
Warunek (D3) dotyczy wyłącznie dalszego
poziomu \emph{dowodliwości}. W następnym rozdziale
przedstawimy formalne uzasadnienie (D2) i (D3);
sama definicja warunków nie gwarantuje ich spełnienia.
\end{uwaga}


\section{Zakres twierdzenia Tarskiego}
Twierdzenie Tarskiego dotyczy \emph{jednej} formuły arytmetycznej,
która miałaby poprawnie rozpoznawać wszystkie zdania pierwszego
rzędu w standardowym modelu. Nie wyklucza definiowania prawdy
dla formuł o ustalonej, ograniczonej złożoności ani zewnętrznej
definicji spełniania w bogatszym metajęzyku. Również predykat
$\operatorname{Bew}_T$ nie jest predykatem prawdy: koduje istnienie
formalnego dowodu, którego poprawność jest syntaktyczna.

\section{Definiowalność, wyliczalność i rozstrzygalność to różne własności}
\begin{uwaga}
\emph{Definiowalność} podzbioru $A\subseteq\N$
oznacza istnienie formuły arytmetycznej $\theta(x)$
bez innych zmiennych wolnych, takiej że dla każdego $n\in\N$
zachodzi $n\in A$ wtedy i tylko wtedy, gdy
$\mathfrak N\models\theta(\overline n)$. \emph{Efektywna wyliczalność}
oznacza istnienie programu wypisującego wszystkie jego elementy i tylko je.
\emph{Rozstrzygalność} oznacza istnienie programu, który dla każdego
argumentu kończy działanie z poprawną odpowiedzią „tak” albo „nie”.
Nie są to pojęcia równoważne. Zbiór twierdzeń efektywnie aksjomatyzowalnej
teorii można efektywnie wyliczać, lecz na mocy wyników z rozdziału
\ref{ch:first-incompleteness} dla efektywnie aksjomatyzowalnych, niesprzecznych
rozszerzeń $Q$ nie można go rozstrzygać.
Pełna prawda arytmetyczna nie jest nawet definiowalna jedną formułą
tego samego języka, co wynika z twierdzenia Tarskiego.

Dla zbioru nieskończonego przeliczalność oznacza istnienie
bijekcji z $\N$, natomiast określenie „co najwyżej przeliczalny”
obejmuje również zbiory skończone i odpowiada istnieniu
iniekcji do $\N$. Żaden z tych warunków sam w sobie
nie dostarcza algorytmu wyznaczającego taką bijekcję lub iniekcję. W kolejnych rozdziałach
\ref{ch:downward-ls} i \ref{ch:upward-ls} zobaczymy,
że konstruowanie modeli przeliczalnych może wymagać wyborów
świadków, których prawdziwości nie umiemy algorytmicznie sprawdzić.
\end{uwaga}

\begin{przyklad}[Dlaczego arytmetyzacja nie tworzy predykatu prawdy?]
Kodujemy napis zdania jako liczbę naturalną $n$,
a więc możemy pytać o jego składnię i o istnienie formalnego
dowodu tego zdania. Nie otrzymujemy jednak automatycznie
algorytmu ani formuły mówiącej, czy zakodowane zdanie
jest prawdziwe w standardowej arytmetyce.
Należy odróżniać badanie skończonego kodu składni
od prawdziwości zdania w standardowej strukturze liczb naturalnych.
Dlatego diagonalizacja nie czyni arytmetyki sprzeczną:
zdanie Gödla mówi o dowodliwości, natomiast
niedozwolony uniwersalny predykat Tarskiego miałby
wyrażać prawdę własnego języka.
\end{przyklad}
\chapter{Operacje na dowodach i warunki wyprowadzalności}\label{ch:derivability-proof}

\section{Dlaczego nie wystarczy sprawdzać pojedynczych dowodów?}
W rozdziale~\ref{ch:tarski-truth} wprowadziliśmy trzy warunki
Hilberta--Bernaysa--Löba, oznaczone (D1)--(D3).
Pierwszy z nich mówi, że do \emph{konkretnego} dowodu możemy
przypisać arytmetyczny dowód zdania o jego istnieniu. Pozostałe
dwa dotyczą natomiast \emph{dowolnych kodów dowodów}, także tych,
które przyjmują niestandardowe wartości w modelach teorii. Musimy więc
w samej arytmetyce uzasadnić poprawność operacji na kodach.

Przez cały rozdział ustalamy efektywnie aksjomatyzowalną teorię
$T\supseteq\mathrm{PA}$ w języku arytmetyki oraz \emph{konkretny}
predykat $\operatorname{Prf}_T(p,x)$ z Lematu~\ref{res:7:7}.
Kod $p$ zawiera dowód, świadectwa rozpoznania aksjomatów i
certyfikat poprawnego przebiegu jego weryfikatora. Zgodnie
z przyjętą tam konstrukcją $\operatorname{Prf}_T$ jest formułą $\Delta_0$.
Stąd $\operatorname{Bew}_T(x):=\exists p\,\operatorname{Prf}_T(p,x)$
jest formułą $\Sigma_1$.

\begin{uwaga}[Dlaczego własność $\Delta_0$ wymaga dodatkowego uzasadnienia?]
Z faktu, że pewna relacja jest pierwotnie rekurencyjna, \emph{nie wynika
bez dalszego argumentu}, że jej wykres można przedstawić formułą
$\Delta_0$ w samym języku $\{0,S,+,\cdot\}$. W niniejszym
predykacie kod $p$ zawiera nie tylko wiersze dowodu, lecz także
świadectwa weryfikacji, a ich składniki i ślady obliczenia są
ograniczone przez rozmiar kodu. Weryfikacja takiego
\emph{certyfikowanego} napisu sprowadza się do ograniczonych
kwantyfikatorów po jego częściach. Właśnie z tego dodatkowego
sposobu kodowania, a nie ze słowa „obliczalny”, pochodzi
przyjęta postać $\Delta_0$. Alternatywnie wystarcza reprezentacja
$\Sigma_1$ predykatu dowodu, jeżeli dla niej osobno udowodnimy
wymagane wewnętrzne operacje na kodach.
\end{uwaga}

Przy wyborze prezentacji $T$ ustalamy dodatkowo, że aksjomaty
$\mathrm{PA}$ otrzymują osobny, łatwo sprawdzalny rodzaj
certyfikatu. Jest to dozwolone: jeśli $T\supseteq\mathrm{PA}$
w sensie zbioru twierdzeń, dołączenie wszystkich aksjomatów
PA do wyliczenia aksjomatów $T$ nie zmienia dowodliwych zdań.
Dzięki oznaczeniu ich typu kod dowodu PA można przekształcić
w kod dowodu $T$ za pomocą funkcji pierwotnie rekurencyjnej.
Ta konwencja jest częścią \emph{wyboru standardowego
predykatu dowodu}; nie zakładamy, że każda formuła
rozpoznająca dowody spełnia analogiczne własności.

\begin{uwaga}[Liczba, numerał i kod składni]
Jeżeli $a=\ulcorner\varphi\urcorner$ jest \emph{liczbą naturalną},
to $\overline a$ jest \emph{termem} arytmetyki. Zapis
$\operatorname{Bew}_T(\overline a)$ jest zdaniem, a
$T\vdash\operatorname{Bew}_T(\overline a)$ stwierdzeniem
metamatematycznym, że istnieje jego dowód. Gdy w poniższych
wzorach występuje funkcja obliczeniowa $m(p,q)$, rozumiemy ją
jako skrót definicyjny dla reprezentującej ją formuły;
nie dopisujemy nowego symbolu funkcyjnego do
$\mathcal L_{\mathrm{ar}}$ bez definicji.
\end{uwaga}

\section{Kodowanie ciągów o zmiennej długości w PA}
Przy czytaniu kolejnych dowodów warto pamiętać zasadniczą różnicę:
możliwość zakodowania \emph{każdego jednego} skończonego ciągu
w metateorii nie daje automatycznie formuły, która
w $\mathrm{PA}$ zapewnia kod \emph{dla każdej długości}.
Potrzebny jest dowód indukcyjny.

\begin{lemat}[Dołączanie jednego wyrazu do kodu ciągu]
\label{lem:xi-append-beta}
Niech $\beta(a,b,i)$ oznacza funkcję odczytu z
Lematu~\ref{res:7:4}. Istnieje formuła arytmetyczna
opisująca dopisanie wyrazu $c$ do ciągu o długości $n$,
dla której $\mathrm{PA}$ dowodzi
\[
\forall a\,\forall b\,\forall n\,\forall c\;\exists a'\,\exists b'
\left(\left[\forall i<n\;\beta(a',b',i)=\beta(a,b,i)\right]
\land\beta(a',b',n)=c\right).
\]
Równość $\beta(u,v,i)=w$ jest skrótem formuły
$B(u,v,i,w)$ z Lematu~\ref{res:7:4}. W~PA istnienie i
jednoznaczność takiego $w$ wynikają z dzielenia z resztą;
nie traktujemy $\beta$ jako pierwotnego symbolu języka.
\end{lemat}
\begin{proof}
W $\mathrm{PA}$ mamy dzielenie z resztą, istnienie
wspólnych wielokrotności skończonej rodziny liczb i
chińskie twierdzenie o resztach; ich dowody uzyskuje się
z algorytmu Euklidesa i indukcji. Wyjaśnimy konkretną
konstrukcję, a nie tylko istnienie kodu.

Dla $i<n$ niech $r_i:=\beta(a,b,i)$, a $r_n:=c$.
Istnienie odczytów $r_i$ i ich jednoznaczność są dowodliwe
w PA, więc przez indukcję i rekurencję na skończonych listach
możemy utworzyć ich maksimum $R$ oraz liczbę
$b':=(n+1)!(R+1)$. Arytmetyka PA dowodzi, że $b'>r_i$
i że każda dodatnia liczba nie większa od $n+1$ dzieli
$(n+1)!$, a zatem również $b'$. Ustalamy
$m_i:=1+(i+1)b'$ dla $i\le n$. Liczby $m_i$ są parami
względnie pierwsze: wspólny dzielnik $m_i$ i $m_j$ dzieli
$(j-i)b'$ i jest względnie pierwszy z $b'$, a więc dzieli
$j-i$; ponieważ $j-i$ dzieli $b'$, musi być równy $1$.
Dodatkowo $r_i<m_i$, ponieważ $r_i\le R$ oraz
$m_i\ge 1+b'>R$.

Chińskie twierdzenie o resztach zapewnia $a'$ takie, że
$a'\equiv r_i\pmod {m_i}$ dla wszystkich $i\le n$.
Formalizacja nie ogranicza się do ustalonych numerałów:
prowadzimy w PA indukcję po liczbie $k$ już uwzględnionych
kongruencji. Jeśli $a_k$ spełnia pierwsze $k$ kongruencji,
a $P_k=\prod_{i<k}m_i$, to $P_k$ i $m_k$ są względnie
pierwsze (każdy $m_i$ jest względnie pierwszy z $m_k$).
W PA algorytm Euklidesa uzasadnia istnienie współczynników
Bézouta, a więc odwrotności $P_k$ modulo $m_k$.
Wybieramy resztę $0\le v<m_k$ spełniającą
$P_kv\equiv r_k-a_k\pmod{m_k}$ i kładziemy
$a_{k+1}:=a_k+P_kv$, redukując wynik modulo $P_km_k$.
Wtedy $a_{k+1}$ spełnia wszystkie pierwsze $k+1$
kongruencje. Istnienie odpowiednich reszt, współczynników
i iloczynów dla każdego $k\le n+1$ jest dowodliwe przez
indukcję PA, a nie przez nieskończoną listę instancji.
Otrzymujemy kod $a'$ dla \emph{dowolnego} $n$.
Warunki na reszty dają tezę z definicji funkcji $\beta$.
\end{proof}

\begin{uwaga}[Co było tu naprawdę istotne?]
Indukcja odbywała się w teorii $\mathrm{PA}$ po \emph{zmiennej}
liczbie kongruencji. W $Q$ wykazywaliśmy tylko konstrukcje
na ustalonych numerałach. Tego dodatkowego użycia indukcji
nie wolno przenosić do $Q$ bez dowodu.
\end{uwaga}

\begin{lemat}[Funkcje pierwotnie rekurencyjne w PA]
\label{lem:xi-pr-in-pa}
Dla każdej funkcji pierwotnie rekurencyjnej
$f:\N^k\to\N$ można wybrać formułę $F_f(\vec x,y)$
taką, że $\mathrm{PA}\vdash\forall\vec x\,\exists!y\,F_f(\vec x,y)$.
Arytmetyka Peana dowodzi także równań definiujących $f$.
\end{lemat}
\begin{proof}
Dowód prowadzimy indukcją po konstrukcji funkcji
pierwotnie rekurencyjnej. Funkcje bazowe (zero, następnik
i rzuty) są reprezentowane przez równania, których
istnienie i jednoznaczność wartości są bezpośrednie.
Jeżeli $f(\vec x)=g(h_1(\vec x),\ldots,h_r(\vec x))$,
to z dowodliwej totalności każdego $h_i$ otrzymujemy
wartości $z_i$, a z totalności $g$ wartość $y$;
jednoznaczność kolejnych wartości daje jednoznaczność $y$.

Pozostaje rekursja prosta:
$f(\vec x,0)=g(\vec x)$,
$f(\vec x,S(n))=h(\vec x,n,f(\vec x,n))$.
Nie używamy w formule symboli funkcyjnych $g,h$, lecz
ich reprezentacji $G(\vec x,z)$ i $H(\vec x,i,z,w)$.
Formuła $F_f(\vec x,n,y)$ stwierdza istnienie kodu ciągu
o długości $n+1$, którego wyraz początkowy $z_0$
spełnia $G(\vec x,z_0)$, ostatni jest równy $y$, a
$H(\vec x,i,z_i,z_{i+1})$ zachodzi dla każdego $i<n$.
Wszystkie odczyty $z_i$ wyrażamy formułą $B$; warunek
$\forall i<n$ jest ograniczony. Przy ustalonym $\vec x$
stosujemy schemat indukcji PA do zdania, że taki kod
istnieje dla długości $n+1$. Dla $n=0$ wystarcza
udowodniona totalność $G$. W kroku indukcyjnym totalność
$H$ dostarcza wyraz $z_{n+1}$, a
Lemat~\ref{lem:xi-append-beta} zapewnia jego dopisanie.
Dla dwóch kodów spełniających te same warunki indukcja
po $i\le n$ i jednoznaczność $G,H$ dają równość ich
$i$-tych wyrazów; w szczególności równe są wyniki $y$.
Tak PA dowodzi istnienia, jednoznaczności i równań
rekurencji dla formuły $F_f$.
\end{proof}

\section{Formalne łączenie dwóch dowodów}
Gdy mamy dowód $\varphi\to\psi$ oraz osobny dowód
$\varphi$, reguła eliminacji implikacji pozwala dopisać
wiersz $\psi$. Na kartce wygląda to prosto. Musimy
jednak upewnić się, że \emph{ta sama operacja działa
na kodach dowolnej długości}, i że $\mathrm{PA}$
może potwierdzić jej poprawność.

\begin{lemat}[Formalizacja modus ponens]\label{lem:xi-formal-mp}
Niech $T\supseteq\mathrm{PA}$ będzie efektywnie
aksjomatyzowalna. Dla ustalonych zdań $\varphi,\psi$
można skonstruować funkcję pierwotnie rekurencyjną
$m_{\varphi,\psi}:\N^2\to\N$, reprezentowaną w PA
przez formułę $M_{\varphi,\psi}(p,q,r)$, taką że
$\mathrm{PA}\vdash\forall p\,\forall q\,\exists!r\,
M_{\varphi,\psi}(p,q,r)$ oraz
\[
\mathrm{PA}\vdash\forall p\,\forall q\,\forall r\;\bigl(
 M_{\varphi,\psi}(p,q,r)\land
\operatorname{Prf}_T(p,\overline{\ulcorner\varphi\to\psi\urcorner})
\land\operatorname{Prf}_T(q,\overline{\ulcorner\varphi\urcorner})
\to\operatorname{Prf}_T(r,
\overline{\ulcorner\psi\urcorner})\bigr).
\]
Zapis $\operatorname{Prf}_T(m_{\varphi,\psi}(p,q),x)$
można stosować jedynie w rozszerzeniu definicyjnym PA;
w pierwotnym języku oznacza podstawienie wyniku opisanego
przez $M_{\varphi,\psi}$.
\end{lemat}
\begin{proof}
\emph{Konstrukcja algorytmu.} Dekodujemy z $p$ i $q$
kompletne dowody, łącznie z certyfikatami aksjomatów.
Następnie: (1) kopiujemy wszystkie wiersze pierwszego
dowodu; (2) kopiujemy wiersze drugiego, przesuwając
o stałą długość pierwszego dowodu numery wszystkich
odwołań do poprzednich wierszy; (3) dopisujemy
wiersz $\psi$, wskazując dwa końcowe wiersze jako
podstawy reguły eliminacji implikacji; (4) tworzymy
certyfikat akceptacji nowego kodu przez weryfikator.
Na niepoprawnych wejściach algorytm oddaje ustalony
kod zastępczy. Odczyty, przepisanie wierszy oraz
przesunięcia wskaźników wykonuje się przez rekursję
ograniczoną długością wejściowych list; są to operacje
pierwotnie rekurencyjne. Weryfikator $V_T$ z Lematu~\ref{res:6:5} jest
funkcją pierwotnie rekurencyjną. \emph{Wybieramy}
realizujący go program otrzymany konstrukcyjnie z jego
schematu pierwotnie rekurencyjnego, a nie dowolny
program obliczający tę samą funkcję. Dla funkcji
bazowych czas działania takiej realizacji ma oszacowanie
pierwotnie rekurencyjne; przy składaniu sumujemy
oszacowania wykonywanych podprogramów, a przy rekursji
prostej tworzymy ograniczenie kolejnych czasów przez
rekursję prostą po jej liczniku. Indukcja po \emph{tym
konkretnym schemacie} daje więc pierwotnie rekurencyjne
ograniczenie liczby kroków, a w konsekwencji pierwotnie
rekurencyjne wytwarzanie pełnego śladu. Tak powstaje
certyfikat akceptacji nowego dowodu. Dla arbitralnego
programu obliczającego funkcję pierwotnie rekurencyjną
czas działania nie musi być pierwotnie rekurencyjny;
wybór realizacji jest zatem istotny.

\emph{Dowód wewnątrz PA.} Nie wystarcza stwierdzenie,
że operacja działa na wszystkich zwykłych liczbach.
Definiujemy w arytmetyce następujący niezmiennik
dla $j$ kopiowanych wierszy: „każdy z pierwszych $j$
wierszy jest poprawnie uzasadniony, a wskaźniki
wcześniejszych przesłanek są przesunięte zgodnie
z długością poprzedniego fragmentu”. Przypadek $j=0$
jest pusty. W kroku indukcyjnym rozróżniamy, czy
kopiowany wiersz jest aksjomatem, czy wynikiem
reguły. W pierwszym przypadku zachowujemy
jego świadectwo pochodzenia z wyliczenia aksjomatów,
a w drugim wykorzystujemy założoną poprawność
jego przesłanek i przesunięcie ich wskaźników.
Ostatni nowy wiersz przechodzi test eliminacji implikacji.
Lematy o kodowaniu ciągów i funkcji pierwotnie
rekurencyjnych gwarantują, że PA potrafi zbudować
kod całej listy oraz ślad jej weryfikacji.
W ten sposób indukcja \emph{po zmiennej liczbie
wierszy} daje w PA dokładnie podaną implikację.
\end{proof}

\begin{przyklad}[Jak odczytać twierdzenie o kodach?]
Załóżmy, że $p$ koduje dowód „jeśli $A$, to $B$”,
a $q$ koduje dowód $A$. Funkcja $m(p,q)$ tworzy
\emph{nowy} kod zawierający oba fragmenty i dodatkowy
krok prowadzący do $B$. Lemat~\ref{lem:xi-formal-mp}
nie twierdzi, że $A$ jest prawdziwe wyłącznie dlatego,
że ktoś zapisał symbol $p$: mówi, że \emph{jeżeli oba
kody przechodzą formalne testy poprawności}, to
kod otrzymany przez ich połączenie również przechodzi
taki test. To odróżnia regułę od jej dowolnego użycia.
\end{przyklad}

\begin{twierdzenie}[Warunek (D2)]\label{thm:xi-d2}
Dla standardowego predykatu dowodliwości $\Box$
teorii $T\supseteq\mathrm{PA}$ zachodzi dla
wszystkich zdań $\varphi,\psi$:
\[
T\vdash\Box(\varphi\to\psi)\to
(\Box\varphi\to\Box\psi).
\]
\end{twierdzenie}
\begin{proof}
Niech $g=\ulcorner\varphi\urcorner$,
$h=\ulcorner\psi\urcorner$ oraz
$k=\ulcorner\varphi\to\psi\urcorner$.
Lemat~\ref{lem:xi-formal-mp} daje w $\mathrm{PA}$
implikację, że jeśli istnieją $p$ i $q$ spełniające
$\operatorname{Prf}_T(p,\overline k)$ i
$\operatorname{Prf}_T(q,\overline g)$, to
istnieje kod dowodu zdania o kodzie $h$.
Po zastosowaniu wprowadzenia kwantyfikatora egzystencjalnego
i twierdzenia o dedukcji otrzymujemy
\[
\mathrm{PA}\vdash
(\exists p\,\operatorname{Prf}_T(p,\overline k))\to
\bigl((\exists q\,\operatorname{Prf}_T(q,\overline g))\to
\exists r\,\operatorname{Prf}_T(r,\overline h)\bigr).
\]
Po rozwinięciu definicji $\Box$ jest to dokładnie (D2).
Ponieważ $T\supseteq\mathrm{PA}$, ten sam dowód
jest dostępny w $T$.
\end{proof}

\section{Jednolita wersja rozstrzygania formuł ograniczonych}
Z (D1) nie otrzymamy (D3) przez zwykłe podstawienie.
(D1) dotyczy rzeczywistego dowodu wybranego zdania,
a (D3) musi być \emph{jednym dowodem w PA} implikacji
zawierającej kwantyfikator egzystencjalny po kodach.

\begin{lemat}[Sformalizowana prawdziwość $\Delta_0$]
\label{lem:xi-delta-proofcompiler}
Niech $\delta(\vec x)$ będzie ustaloną formułą $\Delta_0$.
Istnieje funkcja pierwotnie rekurencyjna $b_\delta^+$,
która dla standardowego $\vec n$ z
$\mathfrak N\models\delta(\overline{\vec n})$
wytwarza kod dowodu zdania $\delta(\overline{\vec n})$.
Ponadto dla tego samego, \emph{jednego} algorytmu
konstrukcji $\mathrm{PA}$ dowodzi zależności jednolitej
\[
\forall\vec x\;\bigl(\delta(\vec x)\to
\operatorname{Prf}_T\bigl(b_\delta^+(\vec x),
\operatorname{inst}_\delta(\vec x)\bigr)\bigr),
\]
gdzie $\operatorname{inst}_\delta(\vec n)$ jest kodem
napisu $\delta(\overline{\vec n})$. Funkcja podstawiania
numerałów jest pierwotnie rekurencyjna na mocy
Lematu~\ref{res:6:5}; funkcja $b_\delta^+$ jest totalna,
także na wejściach, dla których $\delta$ jest fałszywa
(wtedy jej wynik nie musi być kodem poprawnego dowodu).
Funkcje w wyświetlonym wzorze są oznaczeniami
w rozszerzeniu definicyjnym PA. Dokładniej: jeśli
$F_\delta^+(\vec x,p)$ i $I_\delta(\vec x,z)$
reprezentują odpowiednio wykresy $b_\delta^+$ oraz
$\operatorname{inst}_\delta$, to w pierwotnym języku
arytmetyki twierdzenie przyjmuje postać
\[
\mathrm{PA}\vdash\forall\vec x\,\forall p\,\forall z\;\bigl(
[\delta(\vec x)\land F_\delta^+(\vec x,p)\land
 I_\delta(\vec x,z)]\to\operatorname{Prf}_T(p,z)\bigr).
\]
Istnienie i jednoznaczność wyników reprezentowanych funkcji
są dowodliwe w PA na mocy Lematu~\ref{lem:xi-pr-in-pa}.
\end{lemat}
\begin{proof}
W metateorii stosujemy indukcję \emph{po budowie
ustalonej formuły $\delta$}. Wzmocnimy niezmiennik
indukcji: konstruujemy równocześnie \emph{dwa} totalne
kompilatory pierwotnie rekurencyjne $b_\delta^+$
i $b_\delta^-$, dla których PA dowodzi
\[
\begin{aligned}
\delta(\vec x)&\to
 \operatorname{Prf}_T(b_\delta^+(\vec x),
              \operatorname{inst}_\delta(\vec x)),\\
\neg\delta(\vec x)&\to
 \operatorname{Prf}_T(b_\delta^-(\vec x),
              \operatorname{inst}_{\neg\delta}(\vec x)).
\end{aligned}
\]
Obie implikacje rozumiemy jako formuły rozszerzenia
definicyjnego PA; w pierwotnym języku zastępujemy każdy
występujący symbol funkcji przez reprezentującą go formułę
wykresu i używamy dowodliwej jednoznaczności jego wartości.
Wprowadzenie drugiego kompilatora jest konieczne: przy negacji
potrzebny jest dowód fałszywości podformuły, a nie tylko
dowód podformuły prawdziwej. Dla obu polaryzacji funkcje
pozostają totalne, także poza odpowiednią gałęzią warunku.
W każdym kroku uzasadniamy jednolicie poprawność programu
\emph{wewnątrz} PA.

\emph{Formuły atomowe.} Dla każdego podstawienia
numerałów do termów obliczamy ich wartości,
a następnie produkujemy skończone uzasadnienie
z aksjomatów dodawania, mnożenia i następnika.
Jeżeli atomem jest równość, rozstrzygamy ją
przez porównanie wynikowych numerałów; jeśli
jest relacją $<$, korzystamy z definicji porządku.
Dla nierównej pary wartości program tworzy dowód
nierówności, a dla nieprawdziwego porównania
program tworzy dowód jego negacji. W PA poprawność
tych \emph{obu} procedur dla wszystkich parametrów
wynika z indukcji po liczbie wykonywanych operacji,
a nie z osobnych sprawdzeń.

\emph{Spójniki.} Przy negacji zamieniamy role
kompilatorów dodatniego i ujemnego, dołączając
odpowiedni krok klasycznego rachunku zdań. Dla
koniunkcji prawdziwej łączymy dwa dowody składników;
przy jej fałszywości wystarczy dowód negacji jednego
składnika. Analogicznie dla alternatywy. Dla implikacji
$\eta\to\theta$ gałąź dodatnia korzysta z dowodu
$\neg\eta$ albo $\theta$, a gałąź ujemna łączy dowody
$\eta$ i $\neg\theta$ i stosuje odpowiednią tautologię
rachunku zdań. Wybór gałęzi opiera się na pierwotnie
rekurencyjnej procedurze rozstrzygania ustalonej formuły
$\Delta_0$. Wszystkie operacje na kodach dowodów są
pierwotnie rekurencyjne, a PA dowodzi ich poprawności
przez indukcję po długości przepisanych dowodów,
jak w Lemacie~\ref{lem:xi-formal-mp}.

\emph{Kwantyfikatory ograniczone.} Kwantyfikatory
z ograniczeniem $u\le t$ najpierw rozwijamy zgodnie
z Definicją~\ref{res:4:6} do przypadku $u<t$ oraz
jednej dodatkowej instancji $u=t$; dlatego wystarcza
rozpatrzyć ograniczenie ostre. Jeśli
$\delta(\vec x)=\exists u<t(\vec x)\,\eta(\vec x,u)$,
to dla danych $\vec n$ obliczamy
$k=t(\overline{\vec n})$ i badamy kolejno
$u=0,\ldots,k-1$. Jeśli istnieje świadek,
wybieramy najmniejszy i dołączamy odpowiedni
dowód instancji $\eta(\overline{\vec n},\overline u)$,
dowód nierówności $\overline u<t(\overline{\vec n})$
oraz wprowadzenie kwantyfikatora. Dla
$\forall u<t(\vec x)\,\eta(\vec x,u)$ w przypadku
prawdziwym łączymy $k$ dowodów poszczególnych
instancji. W przypadku fałszywym znajdujemy
konkretny indeks $j<k$ i dołączamy dowód
$\neg\eta(\overline{\vec n},\overline j)$.
Dla fałszywej formuły z kwantyfikatorem egzystencjalnym
łączymy natomiast dowody negacji \emph{wszystkich}
$k$ instancji. W ten sposób otrzymujemy także
kompilator ujemny dla obu typów kwantyfikatorów.
W każdym przypadku kompilator dołącza również dowód
równości $t(\overline{\vec n})=\overline k$,
skonstruowany z obliczenia wartości termu $t$.
Bez tego kroku dowód ograniczony przez numerał
$\overline k$ nie byłby jeszcze dowodem formuły
z oryginalnym ograniczeniem $t(\overline{\vec n})$.
Podstawienie równych termów zapewnia odpowiednie
przeniesienie ograniczenia.
Przejście od listy instancji do zdania z ograniczonym
kwantyfikatorem jest formalizowane w PA przez
indukcję po \emph{zmiennym} $k$. Dla konkretnego $k$
można wyprowadzić z aksjomatów PA zdanie
$u<\overline k\to(u=\overline0\lor\cdots\lor
u=\overline{k-1})$ (dla $k=0$ otrzymujemy sprzeczność
z $u<0$). Z dowodów $k$ instancji, po podstawieniu
równości w każdym przypadku, tworzymy dowód
$\forall u<\overline k\,\eta$; ten sam schemat
z dowodami negacji daje dowód
$\forall u<\overline k\,\neg\eta$, a przez prawa
logiki klasycznej --- zaprzeczenie ograniczonego
kwantyfikatora egzystencjalnego. Dla fałszywej formuły
ograniczonej uniwersalnie lub prawdziwej egzystencjalnie
wystarcza jeden świadek oraz odpowiednia reguła logiczna.
Operacja tworzenia dowodu rozkładu skończonego przedziału
ma kod wyznaczany pierwotnie rekurencyjnie z $k$.
Indukcja w PA formalizuje jej poprawność i sprawdzanie
całego wygenerowanego kodu dla zmiennego $k$;
nie zastępujemy jej nieskończoną rodziną dowodów
dla oddzielnych numerałów. Dla ustalonego $k$
wykorzystujemy także Lemat~\ref{res:4:4}, lecz sam ten lemat
nie zastępuje indukcji po zmiennym $k$ w PA.
Lemat~\ref{lem:xi-append-beta} gwarantuje możliwość
wydłużania kodu listy. W ten sposób dowód jednolity
nie jest tylko nieskończoną rodziną oddzielnych
dowodów dla kolejnych numerałów.

Cała procedura jest pierwotnie rekurencyjna,
ponieważ dla ustalonej $\delta$ stosuje jedynie
skończoną liczbę zagnieżdżonych pętli ograniczonych
i pierwotnie rekurencyjne operacje na kodach.
Uzyskany dowód jest początkowo dowodem w PA;
ponieważ $T\supseteq\mathrm{PA}$, można
jednolicie przeliczyć go na pełny kod dowodu $T$
z certyfikatami jego aksjomatów. PA formalizuje
tę konwersję. Otrzymujemy tezę.
\end{proof}

\begin{lemat}[Sformalizowana $\Sigma_1$-zupełność]
\label{lem:xi-formal-sigma1}
Dla ustalonej formuły $\Sigma_1$ postaci
$\sigma(\vec x)=\exists\vec u\,\delta(\vec x,\vec u)$,
gdzie $\delta$ jest $\Delta_0$, zachodzi
\[
\mathrm{PA}\vdash\forall\vec x\;
\bigl(\sigma(\vec x)\to
\operatorname{Bew}_T(\operatorname{inst}_\sigma(\vec x))\bigr).
\]
W szczególności dla zdania $\sigma$:
$\mathrm{PA}\vdash\sigma\to
\operatorname{Bew}_T(\overline{\ulcorner\sigma\urcorner})$.
\end{lemat}
\begin{proof}
Dla danych $\vec x,\vec u$ załóżmy
$\delta(\vec x,\vec u)$. Poprzedni lemat
dostarcza kod poprawnego dowodu
$\delta(\overline{\vec x},\overline{\vec u})$.
Kolejno dołączamy do niego wprowadzenia
kwantyfikatorów egzystencjalnych po wszystkich
zmiennych $\vec u$, otrzymując kod dowodu
$\sigma(\overline{\vec x})$. Ta operacja kodowa
jest pierwotnie rekurencyjna, a PA dowodzi jej
poprawności, zatem
\[
\mathrm{PA}\vdash
\forall\vec x\,\forall\vec u\;\left(
\delta(\vec x,\vec u)\to\exists p\,
\operatorname{Prf}_T(p,\operatorname{inst}_\sigma(\vec x))
\right).
\]
Wprowadzając kwantyfikatory egzystencjalne
po $\vec u$ po lewej stronie implikacji,
uzyskujemy $\sigma(\vec x)\to
\operatorname{Bew}_T(\operatorname{inst}_\sigma(\vec x))$.
To jest \emph{jedno} twierdzenie PA, również wtedy,
gdy wartości parametrów w modelu są niestandardowe.
\end{proof}

\begin{twierdzenie}[Warunki (D1)--(D3) dla $T$]
\label{thm:xi-hbl-verified}
Dla każdej efektywnie aksjomatyzowalnej teorii
$T\supseteq\mathrm{PA}$ ustalony standardowy
predykat dowodliwości z Lematu~\ref{res:7:7}
spełnia wszystkie warunki (D1)--(D3) z
Definicji~\ref{def:hilbert-bernays-loeb}.
\end{twierdzenie}
\begin{proof}
(D1) został udowodniony w Lemacie
\ref{lem:d1-provability}.
Warunek (D2) to Twierdzenie~\ref{thm:xi-d2}.
Dla (D3) ustalmy zdanie $\varphi$ i niech
$g:=\ulcorner\varphi\urcorner$. Ponieważ
$\operatorname{Bew}_T(x)$ jest $\Sigma_1$,
zdanie $\sigma:=\operatorname{Bew}_T(\overline g)$
ma postać $\Sigma_1$. Z Lematu
\ref{lem:xi-formal-sigma1} otrzymujemy
\[
\mathrm{PA}\vdash
\operatorname{Bew}_T(\overline g)\to
\operatorname{Bew}_T\left(
\overline{\ulcorner\operatorname{Bew}_T(\overline g)\urcorner}
\right).
\]
Pierwsza strona to $\Box\varphi$,
a druga to $\Box\Box\varphi$. Zatem (D3)
jest dowodliwe już w PA, więc także w $T$.
\end{proof}


\begin{przyklad}[Czym różni się dowód od jego kodu?]
Ustalmy dwie formalne dedukcje: pierwsza kończy się zdaniem $P$,
a druga kończy się zdaniem $P\to Q$. Ich teksty są obiektami
skończonymi, więc można je zakodować liczbami $q$ i $p$.
Metamatematycznie potrafimy połączyć teksty i dopisać regułę
modus ponens. To daje kod $m(p,q)$. Warunek (D2) mówi więcej:
\emph{ta sama operacja łączenia przechodzi dowód poprawności
wewnątrz arytmetyki}, ze zmiennymi $p,q$, a nie tylko po
wybraniu dwóch konkretnych liczb. Rozdzielenie tych poziomów
jest konieczne: skończona lista sprawdzonych przykładów nie
jest dowodem jednego zdania z dwoma kwantyfikatorami ogólnymi.
\end{przyklad}

\section{Czego nie wolno wyciągać z (D1)--(D3)?}
\begin{przyklad}[Trzy twierdzenia, trzy znaczenia]
Niech $\varphi$ oznacza konkretną równość
$\overline1+\overline1=\overline2$.
(D1) mówi, że mając formalny dowód tej równości,
można skonstruować dowód zdania $\Box\varphi$.
(D2) mówi, że teoria uzasadnia łączenie dowodów
implikacji i jej poprzednika. (D3) mówi,
że teoria potrafi formalnie potwierdzać istnienie
dowodu zdania o istnieniu dowodu.
Żaden z tych warunków nie jest zasadą
$\Box\varphi\to\varphi$ dla każdego $\varphi$.
\end{przyklad}

\begin{uwaga}[Dlaczego wybór predykatu jest istotny?]
Formuły zgodne na liczbach \emph{standardowych}
nie muszą być dowodliwie równoważne w PA.
Dlatego (D2) i (D3) są twierdzeniami
\emph{o ustalonej reprezentacji kodu dowodu},
a nie o dowolnej formule z odpowiednim zbiorem
wartości prawdziwych w $\mathfrak N$.
W szczególności predykat Rossera z
Definicji~\ref{def:rosser-predicate} nie otrzymuje
tych własności automatycznie.
\end{uwaga}
\chapter{Drugie twierdzenie Gödla i twierdzenie Löba}
\label{ch:second-incompleteness}

\section{Czy teoria może dowieść własnej niesprzeczności?}
Pierwsze twierdzenie Gödla, udowodnione w rozdziale~\ref{ch:first-incompleteness},
wyznacza ograniczenia rozstrzygania zdań w efektywnie aksjomatyzowalnych
teoriach arytmetycznych. Teraz pytamy o jedno szczególne zdanie:
czy teoria może dowieść, że nie istnieje dowód sprzeczności
\emph{w tej samej teorii}? Aby pytanie miało ścisły sens,
ustalamy sposób kodowania dowodów.

\begin{definicja}[Formalna niesprzeczność]\label{def:con-t}
Niech $T$ będzie efektywnie aksjomatyzowalną teorią klasycznej logiki
pierwszego rzędu w języku arytmetyki, zawierającą $\mathrm{PA}$.
Niech $\operatorname{Prf}_T(p,x)$ będzie standardowym predykatem
z Lematu~\ref{res:7:7}, a
$\operatorname{Bew}_T(x):=\exists p\,\operatorname{Prf}_T(p,x)$.
Ustalmy zdanie $\bot:=(0=S(0))$ i jego kod $b:=\ulcorner\bot\urcorner$.
\emph{Zdaniem formalnej niesprzeczności} teorii $T$ nazywamy
\[
\operatorname{Con}_T:=
\neg\operatorname{Bew}_T(\overline b)
\qquad\text{(równoważnie: }\neg\Box\bot\text{)}.
\]
Jest to zdanie języka arytmetyki, a nie wypowiedź
metamatematyczna. Ponieważ $\operatorname{Bew}_T$ jest $\Sigma_1$,
zdanie $\operatorname{Con}_T$ jest $\Pi_1$.
\end{definicja}

\begin{przyklad}[Co kwantyfikuje zdanie niesprzeczności?]
Po rozwinięciu definicji $\operatorname{Con}_T$ otrzymujemy
$\forall p\,\neg\operatorname{Prf}_T(p,\overline b)$.
Dla pojedynczej standardowej liczby $n$ zdanie
$\neg\operatorname{Prf}_T(\overline n,\overline b)$ wyklucza jedynie
jeden kod. Nawet gdy dla każdego standardowego $n$ mamy osobny
dowód tej negacji, nie oznacza to, że istnieje jeden dowód
zdania z kwantyfikatorem $\forall p$. To rozróżnienie będzie
kluczowe przy modelach niestandardowych.
\end{przyklad}

\begin{lemat}[Znaczenie niesprzeczności w modelu standardowym]
\label{lem:con-standard}
Dla $T$ z Definicji~\ref{def:con-t} zachodzi
$\mathfrak N\models\operatorname{Con}_T$ wtedy i tylko wtedy,
gdy $T$ jest niesprzeczna.
\end{lemat}
\begin{proof}
Z definicji wybranego predykatu $\operatorname{Prf}_T$
wynika, że $\mathfrak N\models
\operatorname{Bew}_T(\overline b)$ dokładnie wtedy,
gdy istnieje rzeczywista, skończona dedukcja zdania
$\bot$ z aksjomatów $T$. Ponieważ $T$ zawiera
$\mathrm{PA}$, dowodzi $\neg\bot$.
Jeżeli więc dowodzi także $\bot$, jest sprzeczna.
Odwrotnie, jeżeli $T$ jest sprzeczna, to w klasycznej
logice pierwszego rzędu dowodzi dowolnego zdania,
w szczególności $\bot$ (zasada eksplozji).
Brak dowodu $\bot$ jest zatem równoważny
metamatematycznej niesprzeczności $T$.
\end{proof}

\begin{uwaga}[Niesprzeczność nie jest poprawnością semantyczną]
Prawdziwość $\operatorname{Con}_T$ w $\mathfrak N$
nie oznacza, że wszystkie aksjomaty $T$ są prawdziwe
w $\mathfrak N$. Niesprzeczna teoria może mieć modele,
ale nie mieć modelu standardowego. W dalszym ciągu
nie zakładamy poprawności semantycznej $T$ ani
$\omega$-niesprzeczności.
\end{uwaga}

\section{Od dowodliwości zdania do dowodliwości sprzeczności}
W Twierdzeniu~\ref{thm:xi-hbl-verified} wykazano,
że ustalony predykat dowodliwości spełnia warunki
(D1)--(D3) z Definicji~\ref{def:hilbert-bernays-loeb}.
Wewnętrzne operacje na kodach uzasadniają te warunki;
nie wolno dodawać nieudowodnionej implikacji
$\Box\varphi\to\varphi$.

\begin{lemat}[Wewnętrzne rozumowanie o dowodach]
\label{lem:second-box-g}
Niech $T$ będzie efektywnie aksjomatyzowalną teorią
rozszerzającą $\mathrm{PA}$, a $\Box$ jej standardowym
predykatem dowodliwości. Jeśli zdanie arytmetyczne
$G$ spełnia $T\vdash G\leftrightarrow\neg\Box G$,
to $T\vdash\Box G\to\Box\bot$.
\end{lemat}
\begin{proof}
Z kierunku $G\to\neg\Box G$ równoważności diagonalnej
mamy $T\vdash G\to\neg\Box G$.
Warunek (D1) stosujemy do \emph{całej gotowej implikacji},
a potem (D2):
\begin{equation}\label{eq:second-imp-neg}
T\vdash\Box(G\to\neg\Box G),
\qquad T\vdash\Box G\to\Box\neg\Box G.
\end{equation}
Niezależnie (D3) daje
\begin{equation}\label{eq:second-pos}
T\vdash\Box G\to\Box\Box G.
\end{equation}
Z tautologii $\neg\Box G\to(\Box G\to\bot)$
na mocy (D1) otrzymujemy
$T\vdash\Box(\neg\Box G\to(\Box G\to\bot))$.
Zastosowanie (D2) dla zdań $\neg\Box G$ oraz
$\Box G\to\bot$ daje
$T\vdash\Box\neg\Box G\to\Box(\Box G\to\bot)$.
Drugie zastosowanie (D2), tym razem dla
$\Box G$ oraz $\bot$, daje
$T\vdash\Box(\Box G\to\bot)\to
(\Box\Box G\to\Box\bot)$.
Składając te dwie implikacje, uzyskujemy
\[
T\vdash\Box\neg\Box G\to
(\Box\Box G\to\Box\bot).
\]
Wewnątrz dowodu w $T$ załóżmy tymczasowo
$\Box G$. Z~\eqref{eq:second-imp-neg}
wynika $\Box\neg\Box G$,
z~\eqref{eq:second-pos} wynika $\Box\Box G$,
a z ostatniej implikacji --- $\Box\bot$.
Po usunięciu założenia $\Box G$ za pomocą
twierdzenia o dedukcji otrzymujemy
$T\vdash\Box G\to\Box\bot$, zgodnie z tezą.
Nigdzie nie korzystaliśmy z $\Box G\to G$.
\end{proof}

\begin{twierdzenie}[Drugie twierdzenie Gödla o niezupełności]
\label{thm:godel-second}
Niech $T$ będzie niesprzeczną, efektywnie
aksjomatyzowalną teorią klasycznej logiki pierwszego
rzędu w języku arytmetyki, zawierającą $\mathrm{PA}$.
Dla standardowego predykatu dowodliwości
z Definicji~\ref{def:con-t} zachodzi
$T\nvdash\operatorname{Con}_T$.
\end{twierdzenie}
\begin{proof}
Z Lematu diagonalnego~\ref{res:8:1}, zastosowanego
do formuły $\neg\operatorname{Bew}_T(x)$,
otrzymujemy zdanie $G$, dla którego
$Q\vdash G\leftrightarrow\neg\Box G$.
Ponieważ $T\supseteq\mathrm{PA}\supseteq Q$,
ta sama równoważność jest dowodliwa w $T$.
Na mocy Lematu~\ref{lem:second-box-g} mamy
$T\vdash\Box G\to\Box\bot$.
Kontrapozycja, definicja $\operatorname{Con}_T$
oraz drugi kierunek równoważności diagonalnej dają
\begin{equation}\label{eq:con-implies-g}
T\vdash\operatorname{Con}_T\to\neg\Box G,
\qquad T\vdash\neg\Box G\to G,
\qquad T\vdash\operatorname{Con}_T\to G.
\end{equation}
Przypuśćmy, że $T\vdash\operatorname{Con}_T$.
Wówczas~\eqref{eq:con-implies-g} daje $T\vdash G$.
Warunek (D1) zastosowany do tego \emph{rzeczywistego}
dowodu $G$ daje $T\vdash\Box G$.
Jednocześnie kierunek $G\to\neg\Box G$ punktu
stałego daje $T\vdash\neg\Box G$, wbrew
niesprzeczności $T$. Zatem
$T\nvdash\operatorname{Con}_T$.
Nie potrzebowaliśmy $\omega$-niesprzeczności.
\end{proof}

\begin{przyklad}[Arytmetyka Peana]
Jeśli $\mathrm{PA}$ jest niesprzeczna, to
$\mathrm{PA}\nvdash\operatorname{Con}_{\mathrm{PA}}$.
Nie oznacza to, że formalne zdanie o jej
niesprzeczności jest fałszywe: z
Lematu~\ref{lem:con-standard} wynika wówczas
$\mathfrak N\models\operatorname{Con}_{\mathrm{PA}}$.
Wniosek dotyczy ograniczeń dowodzenia
\emph{wewnątrz tej samej teorii}.
\end{przyklad}

\section{Twierdzenie Löba}
Drugie twierdzenie ma następujące uogólnienie:
jeśli teoria potrafi wykazać, że dowodliwość
pewnego zdania pociąga samo to zdanie, to już
potrafi udowodnić to zdanie bez tego założenia.

\begin{twierdzenie}[Twierdzenie Löba]\label{thm:loeb}
Niech $T\supseteq Q$ będzie teorią klasycznej
logiki pierwszego rzędu w języku arytmetyki.
Niech $\operatorname{Bew}_T(x)$ będzie ustaloną
formułą arytmetyczną, a $\Box\varphi$
skrótowym zapisem
$\operatorname{Bew}_T(\overline{\ulcorner\varphi\urcorner})$
dla każdego zdania $\varphi$. Załóżmy, że dla wszystkich
zdań $\varphi,\psi$ spełnione są warunki (D1)--(D3)
z Definicji~\ref{def:hilbert-bernays-loeb}.
Wtedy dla każdego zdania $\varphi$ zachodzi
\[
T\vdash\Box\varphi\to\varphi
\quad\Longrightarrow\quad T\vdash\varphi.
\]
Nie zakładamy niesprzeczności $T$.
\end{twierdzenie}
\begin{proof}
Ustalmy zdanie $\varphi$ i załóżmy
$T\vdash\Box\varphi\to\varphi$.
Z Lematu diagonalnego~\ref{res:8:1}
dla formuły
$\psi(x):=(\operatorname{Bew}_T(x)\to\varphi)$
otrzymujemy zdanie $L$ takie, że
\begin{equation}\label{eq:loeb-fixed}
T\vdash L\leftrightarrow(\Box L\to\varphi).
\end{equation}
Zmienna $x$ jest liczbowa; zapis $\Box x$
nie byłby formułą naszego języka.

\emph{Krok 1: wykazujemy $T\vdash\Box L\to\Box\varphi$.}
Z pierwszego kierunku~\eqref{eq:loeb-fixed}
mamy $T\vdash L\to(\Box L\to\varphi)$.
Korzystając z (D1) i (D2), otrzymujemy
\[
T\vdash\Box(L\to(\Box L\to\varphi)),
\qquad
T\vdash\Box L\to\Box(\Box L\to\varphi).
\]
Warunek (D2) zastosowany do zdań $\Box L$
oraz $\varphi$ daje
$T\vdash\Box(\Box L\to\varphi)\to
(\Box\Box L\to\Box\varphi)$.
Z kolei (D3) daje $T\vdash\Box L\to\Box\Box L$.
Załóżmy pomocniczo $\Box L$.
Pierwsza implikacja dostarcza
$\Box(\Box L\to\varphi)$, a (D3) --- $\Box\Box L$.
Podstawiając je do powyższej instancji (D2),
otrzymujemy $\Box\varphi$. Usunięcie
założenia pomocniczego daje
\begin{equation}\label{eq:loeb-first}
T\vdash\Box L\to\Box\varphi.
\end{equation}

\emph{Krok 2: dowodzimy $L$.}
Z założenia twierdzenia
$T\vdash\Box\varphi\to\varphi$
i~\eqref{eq:loeb-first}
wynika $T\vdash\Box L\to\varphi$.
Drugi kierunek~\eqref{eq:loeb-fixed}
daje $T\vdash(\Box L\to\varphi)\to L$.
Stosując modus ponens, otrzymujemy $T\vdash L$.

\emph{Krok 3: kończymy dowód.}
Dopiero teraz zastosowanie (D1)
do gotowego dowodu $L$ daje $T\vdash\Box L$.
Na mocy~\eqref{eq:loeb-first} uzyskujemy
$T\vdash\Box\varphi$, a z hipotezy
$T\vdash\Box\varphi\to\varphi$
otrzymujemy $T\vdash\varphi$.
\end{proof}

\begin{wniosek}[Drugie twierdzenie z twierdzenia Löba]
\label{cor:godel-from-loeb}
Niech $T$ będzie niesprzeczną teorią klasyczną
zawierającą $Q$, a jej predykat $\Box$
spełnia (D1)--(D3). Wtedy
$T\nvdash\neg\Box\bot$.
\end{wniosek}
\begin{proof}
Przypuśćmy, że $T\vdash\neg\Box\bot$.
Tautologia klasyczna
$\neg\Box\bot\to(\Box\bot\to\bot)$
daje $T\vdash\Box\bot\to\bot$.
Z Twierdzenia~\ref{thm:loeb}, zastosowanego
do $\varphi=\bot$, wynika $T\vdash\bot$,
wbrew niesprzeczności.
\end{proof}

\section{Modele, które „widzą” dowód sprzeczności}
\begin{twierdzenie}[Niestandardowe modele teorii z zaprzeczeniem niesprzeczności]
\label{thm:models-not-con}
Niech $T\supseteq\mathrm{PA}$ będzie niesprzeczną,
efektywnie aksjomatyzowalną teorią klasycznej logiki
pierwszego rzędu. Wtedy
$T\cup\{\neg\operatorname{Con}_T\}$ jest
niesprzeczna i ma model. Żaden model tego rozszerzenia
nie jest izomorficzny ze standardową strukturą $\mathfrak N$.
\end{twierdzenie}
\begin{proof}
Gdyby $T\cup\{\neg\operatorname{Con}_T\}$
była sprzeczna, twierdzenie o dedukcji dawałoby
$T\vdash\neg\operatorname{Con}_T\to\bot$,
a klasyczna logika --- $T\vdash\operatorname{Con}_T$.
Przeczy to Twierdzeniu~\ref{thm:godel-second}.
Rozszerzenie jest więc niesprzeczne.
Z twierdzenia Gödla o pełności,
Twierdzenie~\ref{res:5:3}, wynika istnienie modelu.
Z Lematu~\ref{lem:con-standard} mamy
$\mathfrak N\models\operatorname{Con}_T$,
więc $\mathfrak N$ nie spełnia
$\neg\operatorname{Con}_T$.
Ponieważ prawdziwość zdań jest zachowywana
przez izomorfizmy struktur, żaden model
rozszerzenia nie jest izomorficzny z $\mathfrak N$.
\end{proof}

\begin{uwaga}[Skąd bierze się niestandardowy świadek?]
Niech $\mathfrak M\models
T\cup\{\neg\operatorname{Con}_T\}$.
Z definicji zaprzeczenia niesprzeczności
istnieje $p\in M$ takie, że
$\mathfrak M\models\operatorname{Prf}_T(p,\overline b)$.
Nie może to być żadna standardowa liczba
$\overline n^{\mathfrak M}$.
Ponieważ $T$ jest niesprzeczna,
dla każdego standardowego $n$ relacja
$\operatorname{Prf}_T(n,b)$ jest fałszywa.
Z rozstrzygania zamkniętych instancji
$\Delta_0$ w $Q$
(Lemat~\ref{lem:incomp-proof-check}) wynika
$Q\vdash\neg\operatorname{Prf}_T(\overline n,\overline b)$.
Model $\mathfrak M$ spełnia $Q$, więc żaden
$\overline n^{\mathfrak M}$ nie jest świadkiem.
Element $p$ jest niestandardowym kodem w
rozumieniu $\mathfrak M$, a nie rzeczywistą,
skończoną dedukcją sprzeczności w metateorii.
\end{uwaga}

\begin{uwaga}[Zależność od predykatu dowodu]
Wniosek $T\nvdash\operatorname{Con}_T$
dotyczy określonej formalizacji dowodliwości,
dla której udowodniono (D1)--(D3).
Formuły zgodne w standardowym modelu nie muszą
być dowodliwie równoważne w $T$.
Nie należy bez dodatkowego dowodu stosować
tego twierdzenia do arbitralnego predykatu,
w szczególności do predykatu Rossera
z Definicji~\ref{def:rosser-predicate}.
\end{uwaga}

\section{Granice interpretacji}
\begin{uwaga}[Czego twierdzenia nie mówią?]
Nie udowodniliśmy niesprzeczności $\mathrm{PA}$
ani tego, że zdanie
$\operatorname{Con}_{\mathrm{PA}}$
nie może zostać udowodnione w teorii silniejszej.
Jeśli teoria $U$ dowodzi takiego zdania, jej
argument jest względny wobec aksjomatów $U$.
Drugie twierdzenie Gödla nie podważa poprawnych
algorytmów weryfikacji dowodów ani nie uzasadnia
twierdzeń o nieomylności ludzkiego rozumowania.
\end{uwaga}

\chapter{Aksjomaty ZF, ZFC i konstrukcje zbiorowe}\label{ch:zf-axioms}
\chaptermark{Aksjomaty ZF i ZFC}

\section{Dlaczego nie każdy opis wyznacza zbiór?}
We wcześniejszych rozdziałach pracowaliśmy ze zbiorami w elementarnym
języku matematyki. Teraz pytamy, co dokładnie wolno w tym języku
\emph{przyjąć} i jakie zbiory możemy \emph{skonstruować} z aksjomatów.
Naiwna zasada „każda własność wyznacza zbiór wszystkich obiektów,
które ją mają” prowadzi do sprzeczności. Gdyby istniał zbiór
$R=\{x:x\notin x\}$, mielibyśmy
$R\in R\iff R\notin R$. Nie jest to problem z prawami logiki;
problemem jest nieograniczone tworzenie zbiorów.

\begin{definicja}[Język teorii mnogości]\label{def:zf-language}
Językiem teorii Zermela--Fraenkla jest \emph{język pierwszego rzędu}
z jednym dwuargumentowym symbolem relacyjnym $\in$ i równością.
Jego zmienne $x,y,z,\ldots$ przebiegają wyłącznie po \emph{zbiorach}.
Zdanie $x\in y$ oznacza, że zbiór $x$ jest elementem zbioru $y$.
Teoria $\mathrm{ZF}$ składa się z opisanych niżej aksjomatów
i schematów aksjomatów. Teoria $\mathrm{ZFC}$ to $\mathrm{ZF}$
wraz z \emph{aksjomatem wyboru}, który wprowadzimy i omówimy
w rozdziale~\ref{ch:axiom-choice}.
\end{definicja}


\section{Jak czytać aksjomaty: znaczenie zmiennych i skrótów}\label{sec:zf-reading-notation}

\emph{Najważniejsza różnica względem zwykłych podręczników z analizy:}
w formalnym języku ZF nie ma osobnej kategorii „elementów”
i osobnej kategorii „zbiorów”. \emph{Wszystkie zmienne przebiegają po zbiorach}.
Litery $x,a,b,p,u$ nie oznaczają innych rodzajów obiektów
niż wielkie litery $A,B,X$. Zmienne pisane małą literą
są tylko \emph{konwencją redakcyjną}. Dlatego
\[
             x\in a
\]
znaczy dokładnie: „zbiór oznaczony przez $x$ jest elementem
zbioru oznaczonego przez $a$”. Nie oznacza to, że
$x$ jest małym obiektem, a $a$ małym zbiorem.
Zamiast $x\in a$ moglibyśmy napisać $x\in A$,
$u\in B$ albo $A\in B$, gdyby symbole były odpowiednio
zdefiniowane i związane kwantyfikatorami. Wielkość liter
nie wpływa na prawdziwość zdania.

\begin{przyklad}[Element sam jest zbiorem]
Niech $a=\{\varnothing,\{\varnothing\}\}$
i $x=\{\varnothing\}$. Wtedy $x\in a$, ponieważ
$x$ jest jednym z dwóch elementów zbioru $a$.
Jednocześnie $x$ \emph{też jest zbiorem}, a jego
jedynym elementem jest $\varnothing$.
Mamy więc $\varnothing\in x$ oraz $x\in a$,
również $x\subseteq a$, ponieważ $\varnothing\in a$; jednak
te dwa fakty wynikają z osobnych sprawdzeń. W ogólności
$x\in a$ nie implikuje $x\subseteq a$.
Dla kontrprzykładu do implikacji $x\in a\Rightarrow x\subseteq a$
weźmy $a=\{\{\{\varnothing\}\}\}$ i $x=\{\{\varnothing\}\}$.
Wtedy $x\in a$, ale $\{\varnothing\}\in x$ oraz
$\{\varnothing\}\notin a$, więc $x\not\subseteq a$.
\end{przyklad}

\begin{definicja}[Słownik notacji aksjomatów]\label{def:zf-notation-conventions}
W całym rozdziale używamy następujących skrótów:
\begin{enumerate}[label=(\roman*)]
\item $\forall x\,\psi(x)$ oznacza: „dla każdego zbioru $x$
      zachodzi $\psi(x)$”; $\exists x\,\psi(x)$: „istnieje zbiór $x$”.
\item $\forall x\in a\;\psi(x)$ jest skrótem formuły
      $\forall x\,(x\in a\to\psi(x))$.
      Analogicznie $\exists x\in a\;\psi(x)$ oznacza
      $\exists x\,(x\in a\land\psi(x))$.
\item $x\subseteq a$ oznacza
      $\forall t\,(t\in x\to t\in a)$; nie jest osobnym
      symbolem pierwotnego języka ZF.
\item $\exists!y\,\psi(y)$ oznacza „istnieje dokładnie
      jeden $y$ spełniający $\psi$”, czyli
      $\exists y\,[\psi(y)\land\forall z\,(\psi(z)\to z=y)]$.
\item $\vec p=(p_1,\ldots,p_n)$ oznacza skończony ciąg
      \emph{parametrów}, które występują w rozważanej formule;
      $\varphi(x,\vec p)$ nie jest dodatkowym spójnikiem,
      lecz sposobem pokazania, od jakich zmiennych zależy formuła.
\item $\{x\in a:\varphi(x)\}$, $\bigcup a$, $\mathcal P(a)$
      oraz $\varnothing$ są nazwami \emph{skonstruowanych zbiorów}.
      W pierwotnym języku mającym jedynie $\in$ i równość
      trzeba je rozwinąć do formuł określających ich elementy.
\end{enumerate}
\end{definicja}

\begin{przyklad}[Czytanie kwantyfikatorów od lewej do prawej]
Rozważmy formułę
\[
\forall a\,\forall b\,\exists p\,\forall x\;
\bigl(x\in p\leftrightarrow(x=a\lor x=b)\bigr).
\]
Przeczytajmy każdy fragment. „Dla \emph{dowolnego} zbioru $a$
i dla \emph{dowolnego} zbioru $b$ istnieje zbiór $p$,
taki że dla \emph{każdego} zbioru $x$ zachodzi:
$x$ jest elementem $p$ dokładnie wtedy, gdy $x=a$ lub $x=b$”.
W szczególności $p=\{a,b\}$. Warto zwrócić uwagę, że
$\forall x$ jest objęte zakresem $\exists p$: po ustaleniu
$a,b$ wybieramy \emph{jeden} zbiór $p$ działający dla
wszystkich $x$, a nie osobne $p$ dla każdego $x$.
\end{przyklad}

\begin{przyklad}[Zamiana liter nie zmienia aksjomatu]
Zapis $\forall x\,\forall y\,\forall z\;\Theta(x,y,z)$
i zapis $\forall a\,\forall b\,\forall t\;\Theta(a,b,t)$
wyrażają to samo zdanie, jeżeli w drugim zapisie
w całej formule jednocześnie i konsekwentnie zmienimy
nazwy \emph{związanych} zmiennych. To jest zmiana nazw
zmiennych związanych, a nie zmiana treści twierdzenia.
Dlatego w części formalnej zwykle stosujemy małe litery,
a wielkie litery pozostawiamy na intuicyjne nazwy
ustalonych zbiorów w przykładach.
\end{przyklad}

\begin{uwaga}[Zbiory w teorii a uniwersum modelu]
W modelu $\mathfrak M=(M,E)$ teorii mnogości relacja $E$ interpretuje
symbol $\in$. Elementy $M$ są obiektami nazywanymi „zbiorami”
\emph{wewnątrz modelu}; nie należy automatycznie utożsamiać $E$
z zewnętrzną relacją przynależności używaną w metateorii.
Tę różnicę zilustrowaliśmy w Twierdzeniu~\ref{thm:skolem-paradox}.
\end{uwaga}

\section{Ekstensjonalność, pustka i pary}
\begin{definicja}[Aksjomat ekstensjonalności]\label{def:zf-extensionality}
Postulujemy zdanie
\[
\forall x\,\forall y\Bigl(\bigl(\forall z\,(z\in x\leftrightarrow z\in y)\bigr)
\to x=y\Bigr).
\]
Dwa zbiory o tych samych elementach są więc równe.
\end{definicja}

\begin{definicja}[Aksjomat istnienia zbioru pustego]\label{def:zf-empty}
Postulujemy $\exists e\,\forall x\,(x\notin e)$.
Taki zbiór jest jedyny na mocy ekstensjonalności;
oznaczamy go $\varnothing$.
\end{definicja}

\begin{definicja}[Aksjomat pary]\label{def:zf-pairing}
Dla dowolnych zbiorów $a,b$ istnieje zbiór $p$ taki, że
\[
\forall x\bigl(x\in p\leftrightarrow(x=a\lor x=b)\bigr).
\]
Jedyny taki zbiór zapisujemy $\{a,b\}$. W szczególności
$\{a\}:=\{a,a\}$ oznacza singleton.
\end{definicja}


\begin{przyklad}[Rozwinięcie aksjomatu ekstensjonalności]
Aksjomat~\ref{def:zf-extensionality} ma budowę:
\[
\forall x\,\forall y\;
\underbrace{\Bigl(\forall z\,(z\in x\leftrightarrow z\in y)\Bigr)}_{\text{te same elementy}}
\to\underbrace{(x=y)}_{\text{równe zbiory}}.
\]
Najpierw ustalamy dwa dowolne zbiory $x$ i $y$.
Następnie rozpatrujemy \emph{warunek}: każdy obiekt $z$
należy do $x$ dokładnie wtedy, gdy należy do $y$.
Dopiero z tego warunku wnioskujemy, że $x=y$.
Implikacja nie głosi, że dowolne dwa zbiory są równe.
Przykładowo $x=\{\varnothing\}$ i $y=\varnothing$
nie spełniają poprzednika, bo $\varnothing\in x$, lecz
$\varnothing\notin y$.
\end{przyklad}

\begin{przyklad}[Nieuporządkowana para a para uporządkowana]
Zbiór $\{a,b\}$ nie rozróżnia kolejności, bo $\{a,b\}=\{b,a\}$.
Natomiast para Kuratowskiego
$(a,b):=\{\{a\},\{a,b\}\}$ rozróżnia współrzędne:
jak wykazaliśmy w rozdziale~\ref{ch:relations},
$(a,b)=(c,d)$ implikuje $a=c$ i $b=d$.
Aksjomat pary zapewnia, że oba zbiory występujące w tym zapisie
rzeczywiście istnieją.
\end{przyklad}


\begin{przyklad}[Aksjomat pary, krok po kroku]
Ustalmy $a=\varnothing$ i $b=\{\varnothing\}$.
Aksjomat pary gwarantuje istnienie $p$ spełniającego
\[
\forall x\quad(x\in p\leftrightarrow
  (x=\varnothing\lor x=\{\varnothing\})).
\]
Sprawdźmy, co to oznacza:
\begin{enumerate}[label=(\roman*)]
\item biorąc $x=\varnothing$, widzimy, że $x\in p$;
\item biorąc $x=\{\varnothing\}$, także otrzymujemy $x\in p$;
\item jeżeli $x$ nie jest żadnym z tych dwóch zbiorów,
      prawa strona równoważności jest fałszywa,
      więc $x\notin p$.
\end{enumerate}
Zatem $p=\{\varnothing,\{\varnothing\}\}$.
Aksjomat nie wymaga, aby $a\ne b$; dla $a=b$
daje zbiór jednoelementowy $\{a\}$.
\end{przyklad}

\section{Suma zbioru, zbiór potęgowy i separacja}
\begin{definicja}[Aksjomat sumy]\label{def:zf-union}
Dla każdego zbioru $a$ istnieje zbiór $u$ spełniający
\[
\forall x\bigl(x\in u\leftrightarrow
\exists b\,(b\in a\land x\in b)\bigr).
\]
Piszemy $u=\bigcup a$. Dla $a=\{A,B\}$ otrzymujemy
zwykłą sumę dwóch zbiorów $A\cup B$.
\end{definicja}

\begin{definicja}[Aksjomat zbioru potęgowego]\label{def:zf-power}
Dla każdego zbioru $a$ istnieje zbiór $p$ taki, że
\[
\forall x\bigl(x\in p\leftrightarrow x\subseteq a\bigr),\qquad
x\subseteq a\ \Longleftrightarrow\ \forall t\,(t\in x\to t\in a).
\]
Oznaczamy go $\mathcal P(a)$; zawiera wszystkie \emph{zbiory}
będące podzbiorami $a$.
\end{definicja}


\begin{przyklad}[Aksjomat sumy na konkretnym zbiorze]
Niech $a=\{\{\varnothing\},\{\{\varnothing\}\}\}$.
Jego elementami są dwa zbiory:
$b_1=\{\varnothing\}$ i $b_2=\{\{\varnothing\}\}$.
Aksjomat sumy mówi, że
\[
\bigcup a=\{x:\exists b\in a\ (x\in b)\}.
\]
Jedyny element $b_1$ to $\varnothing$, a jedyny
element $b_2$ to $\{\varnothing\}$. Zatem
$\bigcup a=\{\varnothing,\{\varnothing\}\}$.
Nie chodzi o sumowanie liczb: symbol $\bigcup$ oznacza
zebranie \emph{wszystkich elementów zbiorów należących do $a$}.
\end{przyklad}

\begin{definicja}[Schemat wyróżniania (separacji)]\label{def:zf-separation}
Niech $\varphi(x,\vec p)$ będzie dowolną formułą języka teorii
mnogości, której parametry $\vec p$ są ustalonymi zbiorami.
Dla każdego zbioru $a$ istnieje zbiór $b$ o własności
\[
\forall x\bigl(x\in b\leftrightarrow
(x\in a\land\varphi(x,\vec p))\bigr).
\]
Oznaczamy go $\{x\in a:\varphi(x,\vec p)\}$.
Jedno \emph{wystąpienie schematu} powstaje po ustaleniu jednej
formuły $\varphi$; mówimy o schemacie, a nie pojedynczym
zdaniu, ponieważ takich formuł jest nieskończenie wiele.
\end{definicja}

\begin{przyklad}[Dlaczego paradoks Russella już nie powstaje?]
Dla dowolnego \emph{wcześniej danego} zbioru $a$ wolno zbudować
$R_a:=\{x\in a:x\notin x\}$. Gdyby $R_a\in a$,
to z definicji otrzymalibyśmy
$R_a\in R_a\iff R_a\notin R_a$, sprzeczność.
Wniosek jest więc inny: $R_a\notin a$. Tym samym
nie może istnieć zbiór uniwersalny $V$ zawierający wszystkie
zbiory, bo musiałby zawierać $R_V$. Separacja nie pozwala
wskazać „zbioru wszystkich”, zanim nie dostaliśmy zbioru
ograniczającego $a$.
\end{przyklad}

\begin{twierdzenie}[Przecięcie i różnica są zbiorami]\label{thm:zf-intersections}
Dla każdych zbiorów $A,B$ istnieją zbiory $A\cap B$ i
$A\setminus B$. Dla każdej niepustej rodziny $\mathcal F$
będącej zbiorem istnieje $\bigcap\mathcal F$.
\end{twierdzenie}
\begin{proof}
Z separacji dla formuły $x\in B$ otrzymujemy
$A\cap B=\{x\in A:x\in B\}$. Dla formuły $x\notin B$
otrzymujemy $A\setminus B=\{x\in A:x\notin B\}$.
Jeżeli $\mathcal F\ne\varnothing$, wybieramy jeden element
$F_0\in\mathcal F$. Wówczas zbiór
$\{x\in F_0:\forall F\in\mathcal F\ (x\in F)\}$
istnieje dzięki separacji i jest równy przecięciu rodziny.
Nie potrzebujemy aksjomatu wyboru: wybieramy jeden element
z jednego niepustego zbioru.
\end{proof}

\begin{twierdzenie}[Iloczyn kartezjański jest zbiorem]\label{thm:zf-product}
Dla dowolnych zbiorów $A,B$ istnieje zbiór
$A\times B=\{(a,b):a\in A,\ b\in B\}$.
\end{twierdzenie}
\begin{proof}
Dla $(a,b)=\{\{a\},\{a,b\}\}$ i $a\in A$, $b\in B$
wszystkie elementy tej pary należą do $\mathcal P(A\cup B)$.
Stąd $(a,b)\in\mathcal P(\mathcal P(A\cup B))$.
Korzystając dwukrotnie z aksjomatu zbioru potęgowego,
a następnie z separacji dla formuły
$\exists a\in A\,\exists b\in B\ (z=(a,b))$,
wyodrębniamy wszystkie i tylko takie pary.
\end{proof}


\begin{przyklad}[Schemat wyróżniania jest ograniczony zbiorem]
Niech $a=\{\varnothing,\{\varnothing\},\{\{\varnothing\}\}\}$
i niech $\varphi(x)$ oznacza formułę „$\varnothing\in x$”.
Wówczas separacja tworzy zbiór
\[
 b=\{x\in a:\varnothing\in x\}=\{\{\varnothing\}\}.
\]
Sprawdźmy kolejno trzy elementy $a$:
$\varnothing$ nie ma elementów, więc odpada;
$\{\varnothing\}$ zawiera $\varnothing$, więc zostaje;
$\{\{\varnothing\}\}$ zawiera wyłącznie $\{\varnothing\}$,
więc również odpada. Widzimy, że \emph{najpierw istnieje}
zbiór $a$, a dopiero potem wolno wyodrębnić jego elementy.
\end{przyklad}

\section{Zastępowanie i nieskończoność}
\begin{definicja}[Schemat zastępowania]\label{def:zf-replacement}
Dla każdej formuły $\varphi(x,y,\vec p)$ zachodzi następujący
aksjomat: jeśli dla każdego $x\in a$ istnieje dokładnie jeden
$y$ taki, że $\varphi(x,y,\vec p)$, to istnieje zbiór
\[
 b=\{y:\exists x\in a\ \varphi(x,y,\vec p)\}.
\]
Ściśle: dla każdej formuły $\varphi$ schemat zawiera domknięcie ogólne
implikacji $[\forall x\in a\;\exists!y\,\varphi(x,y,\vec p)]\Rightarrow
[\exists b\;\forall y\,(y\in b\leftrightarrow\exists x\in a\;\varphi(x,y,\vec p))]$.
Formuła może używać parametrów, ale musi wyznaczać
\emph{jednoznaczną} wartość na każdym elemencie $a$.
\end{definicja}

\begin{przyklad}[Obraz zbioru przez funkcję]
Jeśli $f:A\to B$ jest funkcją i $C\subseteq A$, to
$f[C]=\{f(x):x\in C\}$ istnieje dzięki zastępowaniu,
bo przyporządkowanie $x\mapsto f(x)$ jest jednoznaczne.
W tym prostym przypadku można też użyć separacji w $B$:
$f[C]=\{y\in B:\exists x\in C\ (f(x)=y)\}$.
Schemat zastępowania jest szczególnie użyteczny wtedy,
gdy nie znamy z góry zbioru obejmującego wszystkie wyniki.
\end{przyklad}

\begin{definicja}[Aksjomat nieskończoności]\label{def:zf-infinity}
Istnieje zbiór $I$ taki, że $\varnothing\in I$ i dla
każdego $x\in I$ zachodzi $x\cup\{x\}\in I$.
Zapis $S(x):=x\cup\{x\}$ nazywamy tu \emph{następnikiem zbioru}.
Zbiór spełniający te dwa warunki nazywamy \emph{indukcyjnym}.
\end{definicja}

\begin{twierdzenie}[Istnienie najmniejszego zbioru indukcyjnego]
\label{thm:zf-omega}
W ZF istnieje zbiór $\omega$ zawierający $\varnothing$
i zamknięty na $S$, zawarty w każdym zbiorze indukcyjnym.
Ponadto zachodzi zasada indukcji: jeśli $X\subseteq\omega$
zawiera $\varnothing$ i jest zamknięty na $S$, to $X=\omega$.
\end{twierdzenie}
\begin{proof}
Niech $I$ będzie zbiorem indukcyjnym danym aksjomatem
nieskończoności. Zbiór wszystkich indukcyjnych podzbiorów
$I$ istnieje przez separację w $\mathcal P(I)$:
\[
 \mathcal D=\{J\subseteq I:\varnothing\in J\land
 \forall x\in J\ (S(x)\in J)\}.
\]
Rodzina $\mathcal D$ nie jest pusta, bo $I\in\mathcal D$.
Przez Twierdzenie~\ref{thm:zf-intersections} istnieje
$\omega:=\bigcap\mathcal D$. Ponieważ wszystkie elementy
$\mathcal D$ zawierają $\varnothing$ i są zamknięte na $S$,
również ich przecięcie ma obie własności. Jeżeli $J$ jest
dowolnym zbiorem indukcyjnym, to $J\cap I\in\mathcal D$,
a więc $\omega\subseteq J\cap I\subseteq J$.
Wreszcie każdy indukcyjny podzbiór $X$ zbioru $\omega$
jest również indukcyjnym zbiorem, więc z minimalności
$\omega\subseteq X$. Ponieważ $X\subseteq\omega$, mamy równość.
\end{proof}

\begin{uwaga}[Liczby naturalne jako zbiory]
Definiując $0:=\varnothing$, $1:=S(0)=\{0\}$,
$2:=S(1)=\{0,1\}$ itd., otrzymujemy standardową
konstrukcję von Neumanna. W tym sensie zwykła $\N$
z poprzednich rozdziałów może być reprezentowana przez $\omega$.
Nie jest to nowy aksjomat, lecz konstrukcja oparta na ZF.
\end{uwaga}

\section{Fundowanie i twierdzenie Cantora}
\begin{definicja}[Aksjomat ufundowania (regularności)]
\label{def:zf-foundation}
Każdy niepusty zbiór $a$ ma element $b$, który nie ma
żadnego wspólnego elementu z $a$:
\[
 \forall a\bigl(a\ne\varnothing\to
 \exists b\in a\ (a\cap b=\varnothing)\bigr).
\]
Jest to aksjomat ufundowania. W szczególności wyklucza
$x\in x$: w przeciwnym razie zbiór $\{x\}$
nie miałby elementu rozłącznego z $\{x\}$.
\end{definicja}

\begin{twierdzenie}[Twierdzenie Cantora]\label{thm:zf-cantor}
Dla każdego zbioru $A$ nie istnieje surjekcja
$f:A\to\mathcal P(A)$. Jednocześnie istnieje iniekcja
$\iota:A\to\mathcal P(A)$ zadana przez $\iota(a)=\{a\}$.
\end{twierdzenie}
\begin{proof}
Przypuśćmy, że $f:A\to\mathcal P(A)$ jest surjekcją.
Z separacji istnieje zbiór
\[
 D=\{a\in A:a\notin f(a)\}\subseteq A.
\]
Ponieważ $D\in\mathcal P(A)$, z surjektywności istnieje
$d\in A$ z $f(d)=D$. Wtedy
$ d\in D\iff d\notin f(d)=D$, sprzeczność.
Dla iniekcji zauważmy, że jeśli $\{a\}=\{b\}$,
to z definicji singletonu $a=b$. Aksjomat pary
zapewnia istnienie singletonów. Zastępowanie zastosowane do
jednoznacznej formuły $y=\{x\}$ gwarantuje istnienie
zbioru $\{\{a\}:a\in A\}$, będącego obrazem tej iniekcji.
Dowód obejmuje także $A=\varnothing$: surjekcja
$\varnothing\to\mathcal P(\varnothing)=\{\varnothing\}$
nie istnieje, zaś pusta funkcja jest iniekcją.
\end{proof}


\section{Liczby naturalne von Neumanna: konstrukcja i aksjomaty Peana}

Samo istnienie zbioru $\omega$ z Twierdzenia~\ref{thm:zf-omega}
nie daje jeszcze definicji arytmetycznych działań.
Skonstruujemy je w ZF, używając indukcji i zastępowania,
bez aksjomatu wyboru.

\begin{definicja}[Liczby naturalne von Neumanna]
\label{def:zf-vn-natural}
Liczbami naturalnymi nazywamy elementy $\omega$.
Przyjmujemy $0:=\varnothing$, $1:=S(0)$, $2:=S(1)$
oraz $3:=S(2)$, gdzie $S(n):=n\cup\{n\}$. Stąd
$0=\varnothing$, $1=\{0\}$, $2=\{0,1\}$ i $3=\{0,1,2\}$.
\end{definicja}

\begin{lemat}[Własności następnika]\label{lem:zf-vn-successor}
Dla wszystkich $m,n\in\omega$:
\begin{enumerate}[label=(\roman*)]
\item zbiór $n$ jest przechodni, tj. $x\in y\in n$
      implikuje $x\in n$; ponadto $n\subseteq\omega$ i $n\notin n$;
\item $S(n)\ne0$, zaś $S(m)=S(n)$ implikuje $m=n$;
\item każdy $n\in\omega$ jest zerem albo ma postać $S(k)$,
      gdzie $k\in\omega$;
\item zbiór $\omega$ jest przechodni.
\end{enumerate}
\end{lemat}
\begin{proof}
Przechodniość $n$ dowodzimy indukcją po $n$.
Zbiór $0$ jest przechodni. Jeśli $n$ jest przechodni,
a $x\in y\in S(n)$, to albo $y=n$,
albo $y\in n$. W pierwszym przypadku $x\in n\subseteq S(n)$,
a w drugim z przechodniości $n$ także $x\in n\subseteq S(n)$.
Ufundowanie daje $n\notin n$: w przeciwnym razie
jedyny element zbioru $\{n\}$ przecinałby $\{n\}$.

Ponieważ $n\in S(n)$, zbiór $S(n)$ jest niepusty.
Jeśli $S(m)=S(n)$ i $m\ne n$, to
$m\in S(n)$ i $n\in S(m)$ prowadzą do
$m\in n$ oraz $n\in m$.
Żaden z elementów zbioru $\{m,n\}$ nie byłby
z nim rozłączny, wbrew aksjomatowi ufundowania.

Dla (iii) schemat zastępowania daje zbiór $S[\omega]$.
Zbiór $J:=\{0\}\cup S[\omega]$ zawiera $0$
i jest zamknięty na $S$: jeżeli $k\in J\subseteq\omega$,
to $S(k)\in S[\omega]\subseteq J$.
Z minimalności $\omega$ wynika $J=\omega$.

Wreszcie przez indukcję wykażemy $n\subseteq\omega$
dla wszystkich $n\in\omega$. Dla $n=0$ jest to jasne;
jeśli $n\subseteq\omega$, to
$S(n)=n\cup\{n\}\subseteq\omega$, ponieważ $n\in\omega$.
Jest to również przechodniość całego $\omega$.

Sprawdźmy jeszcze własność używaną w konstrukcji rekurencyjnej:
$0\in S(n)$ dla każdego $n\in\omega$.
Dla $n=0$ mamy $S(0)=\{0\}$.
Jeśli $0\in S(n)$, to
$S(n)\subseteq S(S(n))$, zatem $0\in S(S(n))$.
Indukcja na $\omega$ kończy uzasadnienie.
Ponadto $n\in S(n)$ bezpośrednio z definicji następnika.
\end{proof}

\begin{wniosek}[Aksjomaty Peana dla $(\omega,0,S)$]
\label{cor:zf-peano}
W strukturze $(\omega,0,S)$ zachodzą:
$0\in\omega$, $S[\omega]\subseteq\omega$,
$0\notin S[\omega]$, $S$ jest iniekcją,
a każdy podzbiór $X\subseteq\omega$ zawierający
$0$ i zamknięty na $S$ jest równy $\omega$.
Podzbiór oznacza tu zbiór w sensie ZF;
nie twierdzimy, że w dowolnym modelu ZF
indukcja obejmuje wszystkie zewnętrzne podzbiory
jego interpretacji $\omega$.
\end{wniosek}
\begin{proof}
Pierwsze dwa warunki daje konstrukcja
zbioru indukcyjnego. Dwa następne
wynikają z Lematu~\ref{lem:zf-vn-successor}(ii),
a ostatni z Twierdzenia~\ref{thm:zf-omega}.
Nie wymaga to kwantyfikacji po wszystkich
zewnętrznych podzbiorach uniwersum modelu.
\end{proof}

\begin{przyklad}[Pierwsze następniki]
\label{ex:zf-vn-steps}
Z definicji mamy $S(2)=2\cup\{2\}=\{0,1,2\}=3$.
Zachodzi $2\in3$ i $2\subseteq3$,
natomiast $3\notin3$ na mocy ufundowania.
\end{przyklad}

\begin{figure}[htbp]
\centering
$\displaystyle
0=\varnothing\ \xrightarrow{\ S\ }\ 
1=\{0\}\ \xrightarrow{\ S\ }\ 
2=\{0,1\}\ \xrightarrow{\ S\ }\ 
3=\{0,1,2\}.
$
\caption{Pierwsze liczby naturalne von Neumanna i operacja następnika.}
\label{fig:zf-vn-successor}
\end{figure}

\section{Twierdzenie o rekursji na $\omega$ w ZF}

\begin{twierdzenie}[Rekursja na liczbach naturalnych]
\label{thm:zf-omega-recursion}
Niech $B$ będzie niepustym zbiorem,
$b\in B$ oraz $G:\omega\times B\to B$ funkcją.
Istnieje dokładnie jedna funkcja $F:\omega\to B$
taka, że dla każdego $n\in\omega$
\[
F(0)=b,\qquad F(S(n))=G(n,F(n)).
\]
Twierdzenie jest dowodliwe w ZF, bez AC.
\end{twierdzenie}
\begin{proof}
Dla $n\in\omega$ nazwijmy \emph{$n$-przybliżeniem}
funkcję $f:S(n)\to B$ spełniającą $f(0)=b$
oraz, dla każdego $k\in S(n)$, warunek
$S(k)\in S(n)\Longrightarrow
f(S(k))=G(k,f(k))$.
Warunek $f(0)=b$ ma sens, ponieważ
$0\in S(n)$ dla każdego $n\in\omega$,
co wykazaliśmy w dowodzie
Lematu~\ref{lem:zf-vn-successor}.
Ponieważ $S(n)\times B$ jest zbiorem,
wszystkie takie funkcje tworzą zbiór
przez separację w $\mathcal P(S(n)\times B)$.
W szczególności nie zakładamy wcześniej
istnienia funkcji określonej na całym $\omega$.

Indukcyjnie po $n$ dowodzimy, że istnieje dokładnie
jedno $n$-przybliżenie $f_n$.
Dla $n=0$ jest nim $\{(0,b)\}$.
Jeśli istnieje dokładnie jedno $f_n$, dołączamy
parę $(S(n),G(n,f_n(n)))$.
Wartość $f_n(n)$ jest określona,
gdyż $n\in S(n)$ zgodnie z definicją następnika.
Nowy argument nie należał do $S(n)$,
gdyż $S(n)\notin S(n)$,
a $S(S(n))=S(n)\cup\{S(n)\}$.
Powstaje $(S(n))$-przybliżenie $f_{S(n)}$.
Dowolne inne przybliżenie na $S(S(n))$
po ograniczeniu do $S(n)$ musi równać się $f_n$;
jego jedyna nowa wartość jest wyznaczona równaniem
rekurencyjnym, więc również jest jedyna.

Jednoznacznie określona wartość $f_n(n)$
jest definiowalna formułą teorii mnogości
z parametrami $B,b,G$: piszemy, że istnieje
funkcja $f$ o dziedzinie $S(n)$ spełniająca
warunek przybliżenia oraz $f(n)=y$.
Poprzednia indukcja gwarantuje dla każdego
$n\in\omega$ istnienie dokładnie jednego takiego $y$.
Schemat zastępowania, zastosowany do jednoznacznego
przyporządkowania $n\mapsto (n,y)$,
daje zbiór par
$F:=\{(n,f_n(n)):n\in\omega\}$.
To funkcja $\omega\to B$, przy czym
$F(0)=b$ oraz
$F(S(n))=f_{S(n)}(S(n))=G(n,f_n(n))=G(n,F(n))$.
Jeżeli $F,H$ spełniają oba równania,
to zbiór $\{n\in\omega:F(n)=H(n)\}$
istnieje przez separację, zawiera $0$
i jest zamknięty na $S$. Indukcja w $\omega$
daje $F=H$.
\end{proof}

\begin{uwaga}[Rekursja nie wymaga wyboru]
\label{rem:zf-recursion-not-choice}
Każda wartość przybliżenia jest jednoznacznie
wyznaczona, a schemat zastępowania tworzy
jeden zbiór wszystkich par. Nie wybieramy
dowolnie elementów z rodziny niepustych zbiorów.
\end{uwaga}

\section{Dodawanie: definicja i wyprowadzenie praw}

\begin{definicja}[Dodawanie w $\omega$]\label{def:zf-addition}
Dla każdego $a\in\omega$ zastosujmy
Twierdzenie~\ref{thm:zf-omega-recursion}
z $B=\omega$, wartością początkową $a$
i funkcją $G(n,z)=S(z)$.
Definiujemy $a+n$ przez
\begin{equation}\label{eq:zf-add}
a+0=a,\qquad a+S(n)=S(a+n).
\end{equation}
Schemat zastępowania zapewnia,
że wykres funkcji $+:\omega\times\omega\to\omega$
jest zbiorem.
\end{definicja}

\begin{lemat}[Zero i następnik w dodawaniu]
\label{lem:zf-add-basic}
Dla wszystkich $a,n\in\omega$:
$0+n=n$ oraz $S(a)+n=S(a+n)$.
\end{lemat}
\begin{proof}
Pierwszą tożsamość dowodzimy indukcją po $n$.
Dla $n=0$ jest jasna; z $0+n=n$
wynika $0+S(n)=S(0+n)=S(n)$.
Druga tożsamość dla $n=0$ ma postać
$S(a)=S(a+0)$.
Jeśli $S(a)+n=S(a+n)$, to
\[
S(a)+S(n)=S(S(a)+n)=S(S(a+n))=S(a+S(n)).
\]
Każde przejście wynika z definicji dodawania
albo z hipotezy indukcyjnej.
\end{proof}

\begin{twierdzenie}[Łączność i przemienność dodawania]
\label{thm:zf-add-assoc-comm}
Dla wszystkich $a,b,c\in\omega$:
$(a+b)+c=a+(b+c)$ oraz $a+b=b+a$.
\end{twierdzenie}
\begin{proof}
Ustalmy $a,b$ i indukujmy po $c$.
Dla $c=0$ obie strony pierwszej równości
to $a+b$. W kroku indukcyjnym:
\[
(a+b)+S(c)=S((a+b)+c)=S(a+(b+c))
=a+S(b+c)=a+(b+S(c)).
\]
Przemienność dowodzimy indukcją po $b$ przy
ustalonym $a$. Dla $b=0$ korzystamy z
$a+0=a=0+a$. Jeśli $a+b=b+a$, to
\[
a+S(b)=S(a+b)=S(b+a)=S(b)+a,
\]
gdzie ostatnia równość pochodzi
z Lematu~\ref{lem:zf-add-basic}.
\end{proof}

\begin{lemat}[Suma zerowa i skracanie]
\label{lem:zf-add-cancel}
Dla $a,b,c\in\omega$:
\begin{enumerate}[label=(\roman*)]
\item $a+b=0$ wtedy i tylko wtedy, gdy $a=b=0$;
\item $a+c=b+c$ implikuje $a=b$.
\end{enumerate}
\end{lemat}
\begin{proof}
Jeśli $b\ne0$, to z Lematu~\ref{lem:zf-vn-successor}(iii)
mamy $b=S(k)$, a wtedy
$a+b=S(a+k)\ne0$. Gdy $b=0$,
równość $a+b=0$ oznacza $a=0$.
Kierunek odwrotny jest oczywisty.
Dla skracania indukujemy po $c$.
Przy $c=0$ teza jest bezpośrednia.
Jeśli $a+S(c)=b+S(c)$, to
$S(a+c)=S(b+c)$, więc
iniekcyjność $S$ daje $a+c=b+c$,
a hipoteza indukcyjna --- $a=b$.
\end{proof}

\section{Mnożenie, rozdzielność i potęgowanie}

\begin{definicja}[Mnożenie w $\omega$]\label{def:zf-multiplication}
Dla każdego $a\in\omega$ z
Twierdzenia~\ref{thm:zf-omega-recursion},
przy $B=\omega$, $b=0$ i $G(n,z)=z+a$,
otrzymujemy jednoznaczną funkcję spełniającą
\begin{equation}\label{eq:zf-mult}
a\cdot0=0,\qquad a\cdot S(n)=a\cdot n+a.
\end{equation}
Jej dwuzmienny wykres istnieje przez zastępowanie.
\end{definicja}

\begin{lemat}[Zero i jedynka w mnożeniu]
\label{lem:zf-mult-basic}
Dla każdego $a\in\omega$ mamy
$a\cdot0=0$, $0\cdot a=0$,
$a\cdot1=a$ i $1\cdot a=a$.
\end{lemat}
\begin{proof}
Pierwsza równość jest definicją,
drugą uzyskujemy przez indukcję po $a$:
$0\cdot S(a)=0\cdot a+0=0$.
Dalej
$a\cdot1=a\cdot S(0)=0+a=a$.
Dla ostatniej równości indukujemy po $a$:
$1\cdot0=0$, a z $1\cdot a=a$ dostajemy
$1\cdot S(a)=1\cdot a+1=a+S(0)=S(a)$.
\end{proof}

\begin{twierdzenie}[Rozdzielność mnożenia względem dodawania]
\label{thm:zf-distributivity}
Dla wszystkich $a,b,c\in\omega$:
$a(b+c)=ab+ac$ i $(a+b)c=ac+bc$.
\end{twierdzenie}
\begin{proof}
W obu przypadkach indukujemy po $c$.
Dla $c=0$ stosujemy równości zerowe.
Jeśli $a(b+c)=ab+ac$, to
\[
a(b+S(c))=aS(b+c)=a(b+c)+a
=(ab+ac)+a=ab+(ac+a)=ab+aS(c).
\]
Druga tożsamość w kroku indukcyjnym
wynika z równości
\[
(a+b)S(c)=(a+b)c+(a+b)
=(ac+bc)+(a+b)=(ac+a)+(bc+b)=aS(c)+bS(c).
\]
W przestawianiu składników używamy
Twierdzenia~\ref{thm:zf-add-assoc-comm}.
\end{proof}

\begin{twierdzenie}[Przemienność i łączność mnożenia]
\label{thm:zf-mult-assoc-comm}
Dla wszystkich $a,b,c\in\omega$:
$ab=ba$ oraz $(ab)c=a(bc)$.
\end{twierdzenie}
\begin{proof}
Dla przemienności ustalmy $a$ i indukujmy po $b$.
Baza to $a0=0=0a$. W kroku:
\[
aS(b)=ab+a=ba+1a=(b+1)a=S(b)a,
\]
przy czym korzystamy z rozdzielności
i Lematu~\ref{lem:zf-mult-basic}.
Dla łączności ustalmy $a,b$
i indukujmy po $c$. Baza jest zerowa,
a krok ma postać
\[
(ab)S(c)=(ab)c+ab=a(bc)+ab
=a(bc+b)=a(bS(c)).
\]
Użyliśmy wyłącznie już wykazanej rozdzielności.
\end{proof}

\begin{definicja}[Potęgowanie]
\label{def:zf-exponentiation}
Dla $a\in\omega$ rekurencja z
Twierdzenia~\ref{thm:zf-omega-recursion},
przy $B=\omega$, $b=1$ i $G(n,z)=z\cdot a$,
wyznacza $a^n$ przez
$a^0=1$ i $a^{S(n)}=a^n\cdot a$.
W szczególności $0^0=1$ jest konwencją
tej konstrukcji.
\end{definicja}

\begin{twierdzenie}[Prawo wykładników]
\label{thm:zf-exponent-add}
Dla każdego $a,m,n\in\omega$
zachodzi $a^{m+n}=a^m\cdot a^n$.
\end{twierdzenie}
\begin{proof}
Indukujemy po $n$. Dla $n=0$
obie strony wynoszą $a^m$,
gdyż $a^0=1$.
Z hipotezy indukcyjnej:
\[
a^{m+S(n)}=a^{S(m+n)}=a^{m+n}\cdot a
=(a^m a^n)a=a^m(a^n a)=a^m a^{S(n)}.
\]
Łączność mnożenia pochodzi z
Twierdzenia~\ref{thm:zf-mult-assoc-comm}.
\end{proof}

\begin{przyklad}[Obliczanie z definicji]
\label{ex:zf-arithmetic}
Mamy $2+1=2+S(0)=S(2)=3$,
a $2\cdot2=2\cdot S(1)=2\cdot1+2=2+2=4$.
Kolejno $2^0=1$, $2^1=2$, $2^2=4$, $2^3=8$.
Ostatnie symbole oznaczają odpowiednie
iteracje następnika, nie nowe obiekty pierwotne.
\end{przyklad}

\section{Porządek naturalny, minimum i indukcja silna}

\begin{definicja}[Porządek arytmetyczny]
\label{def:zf-natural-order}
Dla $m,n\in\omega$ definiujemy
$m\le n$ wtedy i tylko wtedy, gdy
istnieje $k\in\omega$ z $m+k=n$.
Piszemy $m<n$, gdy $m\le n$ i $m\ne n$.
Relacji tej nie utożsamiamy z $\in$
bez osobnego dowodu.
\end{definicja}

\begin{lemat}[Odcinki początkowe $\omega$]
\label{lem:zf-initial-segments}
Dla każdego $n\in\omega$ zachodzi
$\{m\in\omega:m\le n\}=S(n)=n\cup\{n\}$.
W szczególności $m<n$ wtedy i tylko wtedy,
gdy $m\in n$.
\end{lemat}
\begin{proof}
Indukujemy po $n$. Przy $n=0$ równość
$m+k=0$ wymusza $m=k=0$
(Lemat~\ref{lem:zf-add-cancel}),
zatem jedynym $m\le0$ jest $0$.
Jeśli teza zachodzi dla $n$,
a $m+k=S(n)$, to przy $k=0$
mamy $m=S(n)$. Przy $k\ne0$
zapisujemy $k=S(t)$, zatem
$S(m+t)=S(n)$. Iniekcyjność
następnika daje $m+t=n$,
czyli $m\le n$ i na mocy
hipotezy indukcyjnej $m\in S(n)$.
W obu sytuacjach $m\in S(S(n))$.
Odwrotnie, jeśli $m\in S(S(n))$,
to albo $m=S(n)$, i wtedy
$m+0=S(n)$, albo $m\in S(n)$.
W drugim przypadku istnieje $t$
z $m+t=n$, więc $m+S(t)=S(n)$.
To daje $m\le S(n)$.
Ostatnią równoważność otrzymujemy,
usuwając z $S(n)$ jego element $n$.
\end{proof}

\begin{twierdzenie}[Dobry porządek $\omega$]
\label{thm:zf-wellordered-omega}
Relacja $\le$ z Definicji~\ref{def:zf-natural-order}
jest dobrym porządkiem $\omega$:
każdy niepusty podzbiór $\omega$
ma element najmniejszy.
\end{twierdzenie}
\begin{proof}
Zwrotność wynika z $m+0=m$,
a przechodniość z łączności dodawania.
Jeśli $m+k=n$ i $n+t=m$,
to $m+(k+t)=m+0$, zatem
z prawa skracania $k+t=0$.
Stąd $k=t=0$ i $m=n$:
porządek jest antysymetryczny.

Liniowość dowodzimy indukcją po $n$
dla własności $\forall m\in\omega\,
(m\le n\lor n\le m)$.
Przy $n=0$ mamy $0+m=m$.
Jeżeli $m\le n$, to
$m+k=n$ daje $m+S(k)=S(n)$,
zatem $m\le S(n)$. Jeśli $n\le m$,
mamy $m=n+k$. Przy $k=0$
otrzymujemy $m=n\le S(n)$,
a przy $k=S(t)$:
$m=n+S(t)=S(n)+t$, więc $S(n)\le m$.

Niech $A\subseteq\omega$ będzie niepusty.
Z aksjomatu ufundowania wybieramy $m\in A$
z $m\cap A=\varnothing$. Dla każdego $n\in A$
liniowość daje $m\le n$ lub $n\le m$.
Druga możliwość z $n\ne m$
oznacza $n\in m$ na mocy
Lematu~\ref{lem:zf-initial-segments},
sprzecznie z $m\cap A=\varnothing$.
Wobec tego $m\le n$ dla wszystkich $n\in A$.
\end{proof}

\begin{wniosek}[Indukcja silna]\label{cor:zf-strong-induction}
Niech $P(n)$ będzie dowolną formułą ZF
z parametrami i wyróżnioną zmienną $n$.
Jeżeli dla każdego $n\in\omega$
z założenia $\forall k\in n\,P(k)$
wynika $P(n)$, to $P(n)$ zachodzi
dla wszystkich $n\in\omega$.
\end{wniosek}
\begin{proof}
Gdyby zbiór $C:=\{n\in\omega:\neg P(n)\}$
był niepusty, miałby najmniejszy element $n_0$.
Wtedy $P(k)$ dla każdego $k\in n_0$,
co na mocy założenia daje $P(n_0)$,
wbrew $n_0\in C$.
Zbiór $C$ istnieje przez separację.
\end{proof}

\begin{przyklad}[Brak największej liczby naturalnej]
\label{ex:zf-no-largest}
Każde $n\in\omega$ ma następnik
$S(n)\in\omega$, a
$n\in S(n)$ oznacza na mocy
Lematu~\ref{lem:zf-initial-segments},
że $n<S(n)$. Zatem żadne $n$
nie jest największym elementem $\omega$.
\end{przyklad}

\begin{lemat}[Zgodność porządku z działaniami]
\label{lem:zf-natural-monotone}
Dla $a,b,c\in\omega$:
\begin{enumerate}[label=(\roman*)]
\item $a<b$ implikuje $a+c<b+c$;
\item $a<b$ i $c\ne0$ implikują $ac<bc$;
\item $ab=0$ implikuje $a=0$ lub $b=0$;
\item $c\ne0$ i $ac=bc$ implikują $a=b$.
\end{enumerate}
\end{lemat}
\begin{proof}
Jeśli $a<b$, to $b=a+d$ dla pewnego $d\ne0$.
Stąd $b+c=(a+c)+d$, a prawo skracania
wyklucza $b+c=a+c$, co daje (i).
Dla (iii), jeśli $a,b\ne0$,
to $b=S(t)$ i $a=S(k)$,
więc $ab=at+a=at+S(k)=S(at+k)\ne0$.
Dla (ii) z $b=a+d$ mamy
$bc=ac+dc$, gdzie $d\ne0$.
Z (iii) i $c\ne0$ wynika $dc\ne0$,
zatem $ac<bc$.
Dla (iv) liniowość porządku
daje $a=b$, $a<b$ albo $b<a$.
Dwa ostatnie przypadki przeczą (ii)
i równości $ac=bc$.
\end{proof}

\begin{twierdzenie}[Dzielenie z resztą]
\label{thm:zf-division}
Dla wszystkich $a\in\omega$ i
$b\in\omega\setminus\{0\}$
istnieje dokładnie jedna para
$(q,r)\in\omega\times\omega$
taka, że $a=bq+r$ oraz $r<b$.
\end{twierdzenie}
\begin{proof}
Ustalmy $b\ne0$ i zastosujmy indukcję silną
po $a$. Jeśli $a<b$, kładziemy
$q=0$ i $r=a$. Wtedy $a=b\cdot0+a$.
Jeśli $b\le a$, istnieje $t\in\omega$
z $a=b+t$. Ponieważ $0+b=b$,
mamy $0\le b$, a z $b\ne0$ wynika $0<b$.
Z Lematu~\ref{lem:zf-natural-monotone}(i),
po dodaniu $t$, otrzymujemy
$t=t+0<t+b=a$.
Hipoteza indukcyjna daje
$t=bq+r$, gdzie $r<b$,
a wtedy
$a=b+(bq+r)=bS(q)+r$.
Istnienie zostało udowodnione.

Dla jednoznaczności niech
$bq+r=bq'+r'$ i $r,r'<b$.
Jeśli $q=q'$, skracanie dodawania
daje $r=r'$. Przy $q<q'$
możemy zapisać $q'=q+d$
dla $d\ne0$; wtedy $d=S(k)$ i
\[
bq'=bq+bd=bq+(bk+b)\ge bq+b.
\]
Istotnie, $0\le bk$, więc z definicji porządku
$bq+b\le bq+bk+b=bq'$.
Podobnie $0\le r'$, stąd
$bq'\le bq'+r'=a$.
Przechodniość daje $bq+b\le a$.
Z drugiej strony $r<b$ i
Lemat~\ref{lem:zf-natural-monotone}(i)
dają $a=bq+r<bq+b$, sprzeczność.
Przypadek $q'<q$ jest symetryczny.
\end{proof}

\begin{przyklad}[Dzielenie $17$ przez $5$]
\label{ex:zf-division}
Dla liczb von Neumanna $17$ i $5$
mamy $17=5\cdot3+2$ i $2<5$,
ponieważ $2+3=5$.
Twierdzenie~\ref{thm:zf-division}
gwarantuje jednoznaczność $(q,r)=(3,2)$.
\end{przyklad}

\begin{uwaga}[Zakres konstrukcji arytmetyki]
\label{rem:zf-arithmetic-scope}
Dodawanie, mnożenie, potęgowanie
i porządek na $\omega$ zostały skonstruowane
za pomocą aksjomatów ZF, a ich wykresy
są zbiorami. Ani aksjomaty Peana dla
$(\omega,0,S)$, ani twierdzenie o rekursji
nie wymagają aksjomatu wyboru.
Ogólniejszą rekursję pozaskończoną
omówimy przy liczbach porządkowych.
\end{uwaga}

\begin{uwaga}[Czego jeszcze nie przyjęliśmy?]
Żaden z dotychczas wymienionych aksjomatów nie orzeka,
że dla \emph{dowolnej} rodziny niepustych zbiorów można
jednocześnie wybrać po jednym elemencie z każdego z nich.
To będzie osobny aksjomat. Rozdział~\ref{ch:axiom-choice}
wykaże jego równoważność z twierdzeniem o dobrym uporządkowaniu
i z lematem Kuratowskiego--Zorna. Zaznaczmy też, że
mówimy tu o \emph{ZF} (z aksjomatem ufundowania), a nie
o teorii z klasami właściwymi jako osobnymi obiektami.
\end{uwaga}
\chapter{Aksjomat wyboru i jego równoważne postacie}\label{ch:axiom-choice}
\chaptermark{Aksjomat wyboru i jego równoważne postacie}

\section{Wybrać jeden element czy wybierać nieskończenie wiele razy?}
Z jednego niepustego zbioru $A$ możemy wskazać element $a\in A$;
nie wymaga to osobnego aksjomatu. Co innego, jeśli mamy rodzinę
$(A_i)_{i\in I}$ niepustych zbiorów i chcemy \emph{jednej funkcji},
która dla każdego $i\in I$ podaje element $A_i$.
Nawet jeśli każde żądanie z osobna da się spełnić,
nie znamy jeszcze reguły gwarantującej dokonanie
wszystkich wyborów jednocześnie.

\begin{definicja}[Aksjomat wyboru, AC]\label{def:choice}
Dla każdego zbioru indeksów $I$ i każdej rodziny
$(A_i)_{i\in I}$ zbiorów niepustych istnieje funkcja
$c:I\to\bigcup_{i\in I}A_i$ taka, że
\[
 \forall i\in I\quad c(i)\in A_i.
\]
Funkcję $c$ nazywamy \emph{funkcją wyboru} tej rodziny.
Teoria $\mathrm{ZFC}$ to $\mathrm{ZF}+\mathrm{AC}$.
\end{definicja}

\begin{przyklad}[Dlaczego rodzina skończona nie wymaga AC?]
Dla dwóch niepustych zbiorów $A_0,A_1$ istnieją elementy
$a_0\in A_0$ i $a_1\in A_1$, więc funkcja
$c(0)=a_0$, $c(1)=a_1$ jest funkcją wyboru.
Przez zwykłą indukcję po liczbie zbiorów otrzymamy tę samą
własność dla każdej rodziny \emph{skończonej}.
AC jest nowym postulatem dopiero w przypadku rodzin dowolnych,
niekoniecznie skończonych.
\end{przyklad}

\begin{przyklad}[Kiedy wybór jest dany konkretnym wzorem?]
Jeżeli wszystkie $A_i$ są niepustymi podzbiorami $\N$,
możemy wybrać $c(i)=\min A_i$. Minimum istnieje dzięki
zwykłemu dobremu uporządkowaniu $\N$.
Taką rodzinę obsługujemy bez dodatkowego aksjomatu.
Nie oznacza to jednak, że dowolny zbiór ma z góry
wyróżniony najmniejszy element: nie każdy zbiór ma
wcześniej ustalony dobry porządek.
\end{przyklad}

\section{Dobre uporządkowanie i jego początkowe odcinki}
\begin{definicja}[Twierdzenie o dobrym uporządkowaniu]
\label{def:wellorder-principle}
Przez \emph{zasadę dobrego uporządkowania} rozumiemy zdanie:
\emph{każdy zbiór $A$ można wyposażyć w porządek liniowy $\le_A$,
w którym każdy niepusty podzbiór ma element najmniejszy}.
Dla dobrego porządku $<$ na $W$ zbiór
$W_{<w}:=\{v\in W:v<w\}$ nazywamy \emph{właściwym odcinkiem
początkowym} określonym przez $w\in W$.
\end{definicja}

\begin{lemat}[Dobre porządki można jednoznacznie porównać]
\label{lem:choice-wellorder-comparison}
Niech $(U,<_{U})$ i $(V,<_{V})$ będą zbiorami dobrze
uporządkowanymi. Wówczas dokładnie jedna z trzech
możliwości zachodzi: są izomorficzne jako zbiory uporządkowane;
$U$ jest izomorficzny z właściwym odcinkiem początkowym $V$;
$V$ jest izomorficzny z właściwym odcinkiem początkowym $U$.
\end{lemat}
\begin{proof}
\emph{Odcinkiem początkowym} nazywamy także cały zbiór;
dla odcinka właściwego wymagamy, aby był postaci $U_{<u}$.
Najpierw zauważmy, że każdy izomorfizm porządkowy między
odcinkami początkowymi jest jednoznaczny: dwa takie
izomorfizmy nie mogą różnić się na wspólnej dziedzinie.
Gdyby się różniły, wybierzmy najmniejszy argument $x$,
na którym dają różne wyniki. Oba przekształcenia muszą
wtedy przenosić zbiór wcześniejszych elementów $U_{<x}$
na zbiór wcześniejszych elementów od odpowiednich obrazów.
Ponieważ na $U_{<x}$ są takie same, mają ten sam obraz
odcinka, a w dobrym porządku element jest wyznaczony
przez swój odcinek poprzedników. Sprzeczność.

Rozważmy zbiór wszystkich izomorfizmów między odcinkami
początkowymi $U$ i $V$. Jest to \emph{zbiór}, ponieważ
składa się z podzbiorów $U\times V$.
Z poprzedniego akapitu izomorfizmy te są zgodne na
częściach wspólnych, więc ich suma $f$ jest izomorfizmem
między pewnymi odcinkami początkowymi $I\subseteq U$
i $J\subseteq V$. Jeśli oba byłyby właściwe,
wybierzmy najmniejsze elementy $u\in U\setminus I$
i $v\in V\setminus J$. Rozszerzenie $f\cup\{(u,v)\}$
byłoby jeszcze izomorfizmem odcinków początkowych,
wbrew definicji $f$ jako sumy wszystkich. Stąd
$I=U$ lub $J=V$. Jeśli oba są całymi zbiorami,
mamy izomorfizm $U\cong V$. W pozostałych przypadkach
jeden zbiór jest izomorficzny z właściwym odcinkiem
drugiego. Te trzy możliwości są rozłączne.
Rzeczywiście, żaden zbiór dobrze uporządkowany $W$
nie jest izomorficzny z własnym właściwym odcinkiem
początkowym $W_{<w}$.
Dla dowolnego ściśle rosnącego zanurzenia porządkowego
$f:W\to W$ pokażemy przez indukcję po dobrym porządku,
że $x\le f(x)$ dla każdego $x\in W$.
Wystarczy zauważyć, że dla wszystkich $y<x$
hipoteza indukcyjna daje $y\le f(y)<f(x)$,
zatem $f(x)$ jest większy od wszystkich $y<x$.
Gdyby $f(x)<x$, sam element $f(x)$
należałby do $W_{<x}$, a przy $y=f(x)$ hipoteza indukcyjna
wymusza $f(y)\ge y=f(x)$.
Z drugiej strony $y<x$ i ścisła monotoniczność
$f$ dają $f(y)<f(x)$, sprzeczność. Stąd $x\le f(x)$.
Gdyby obraz $f[W]$ zawierał się w $W_{<w}$,
mielibyśmy $w\le f(w)<w$, sprzeczność.
Nie wykorzystaliśmy aksjomatu wyboru.
\end{proof}

\section{Twierdzenie Hartogsa bez używania liczb porządkowych}
Do wykazania, że AC daje dobre uporządkowanie \emph{każdego}
zbioru, potrzebujemy specjalnego zbioru dobrze uporządkowanego,
którego nie da się różnowartościowo odwzorować do wyjściowego
zbioru. Zwykle opisuje się go za pomocą liczb porządkowych.
Ponieważ zdefiniujemy je później, użyjemy tu klas
izomorfizmu dobrych porządków; odpowiednia teoria ilorazów
została wprowadzona w rozdziale~\ref{ch:equivalence}.

\begin{twierdzenie}[Twierdzenie Hartogsa]\label{thm:choice-hartogs}
Dla każdego zbioru $A$ istnieje zbiór dobrze uporządkowany
$(H(A),<)$, dla którego nie istnieje iniekcja $H(A)\to A$.
Dowód nie wymaga aksjomatu wyboru.
\end{twierdzenie}
\begin{proof}
Niech $\mathcal W$ będzie zbiorem wszystkich par $(S,R)$,
gdzie $S\subseteq A$, a $R\subseteq S\times S$ jest dobrym
porządkiem na $S$. Ponieważ $R\subseteq S\times S\subseteq A\times A$,
wszystkie rozważane pary należą do
$\mathcal P(A)\times\mathcal P(A\times A)$.
Zbiór $\mathcal W$ jest jego podzbiorem określonym
przez formułę wyrażającą, że $R$ jest dobrym porządkiem na $S$;
istnieje więc przez separację
(Definicje~\ref{def:zf-power} i~\ref{def:zf-separation}
oraz Twierdzenie~\ref{thm:zf-product}).
Definiujemy relację równoważności: dwa takie porządki
są równoważne, jeśli istnieje między nimi izomorfizm
porządkowy. Jest to formuła ZF z parametrami:
izomorfizmy są podzbiorami produktów zbiorów uporządkowanych.
Dla każdego $w\in\mathcal W$ jego klasa $[w]$
istnieje przez separację w $\mathcal W$.
Zbiór wszystkich klas $H(A)$ istnieje przez zastępowanie
zastosowane do jednoznacznego przyporządkowania
$w\mapsto[w]$; można go też otrzymać przez separację
w $\mathcal P(\mathcal W)$.

Oznaczając przez $U,V$ reprezentantów z $\mathcal W$
wraz z ich relacjami dobrego porządku, dla dwóch klas
$[U],[V]\in H(A)$ piszemy $[U]<[V]$,
jeżeli $U$ jest izomorficzny z \emph{właściwym} odcinkiem
początkowym $V$. Jest to określenie niezależne
od wyboru przedstawicieli: izomorfizm porządków przenosi
właściwe odcinki początkowe na właściwe odcinki.
Lemat~\ref{lem:choice-wellorder-comparison} pokazuje,
że jest to ścisły porządek liniowy. Jest on dobry:
dla niepustego $K\subseteq H(A)$ weźmy $h\in K$.
Jeżeli nie ma elementu $K$ mniejszego od $h$, to $h$
jest najmniejszy. W przeciwnym razie podzbiór
$K\cap H(A)_{<h}$ ma element najmniejszy, ponieważ
odcinek $H(A)_{<h}$ jest izomorficzny z dobrze
uporządkowanym przedstawicielem klasy $h$:
elementowi $u$ przedstawiciela przypisujemy klasę
porządkową jego odcinka poprzedników. To przyporządkowanie
jest iniekcją: dwa różne odcinki początkowe jednego
dobrego porządku nie są izomorficzne. Jest surjekcją
na $H(A)_{<h}$ z definicji relacji $<$ i porównania
dobrych porządków. Zachowuje również porządek, gdyż
odcinek poprzedników wcześniejszego elementu jest
właściwym odcinkiem odcinka poprzedników późniejszego.

Przypuśćmy, że istnieje iniekcja $j:H(A)\to A$.
Przenieśmy porządek $H(A)$ na podzbiór $j[H(A)]\subseteq A$.
Ten podzbiór, z przeniesionym porządkiem, ma klasę
$h\in H(A)$, zatem całe $H(A)$ jest izomorficzne
z odcinkiem $H(A)_{<h}$. To niemożliwe dla dobrego
porządku, jak wykazaliśmy w Lemacie~\ref{lem:choice-wellorder-comparison}.
Wniosek: taka iniekcja nie istnieje.
\end{proof}

\begin{przyklad}[Zbiór Hartogsa dla dwóch elementów]
Jeśli $A=\{a,b\}$, to podzbiory $A$ mają tylko
moce $0$, $1$ albo $2$. Każdy dobrze uporządkowany
podzbiór ma więc jeden z trzech typów porządkowych:
pusty, jednoelementowy albo dwuelementowy.
Zatem $H(A)$ składa się z \emph{trzech} elementów,
uporządkowanych w kolejności odpowiadającej długościom
$0<1<2$. Nie istnieje iniekcja trzyelementowego
$H(A)$ do dwuelementowego $A$. Twierdzenie Hartogsa
rozszerza ten pomysł również na zbiory nieskończone,
dla których nie wystarcza liczenie elementów.
\end{przyklad}

\section{AC implikuje dobre uporządkowanie}
\begin{twierdzenie}[Twierdzenie o dobrym uporządkowaniu Zermela]
\label{thm:choice-wellorder}
W ZF aksjomat wyboru implikuje zasadę dobrego uporządkowania.
\end{twierdzenie}
\begin{proof}
Niech $A$ będzie dowolnym zbiorem. Dla $A=\varnothing$
teza jest oczywista; załóżmy $A\ne\varnothing$.
Stosujemy AC do rodziny wszystkich niepustych podzbiorów $A$:
istnieje $c:\mathcal P(A)\setminus\{\varnothing\}\to A$
taka, że $c(B)\in B$.
Niech $H:=H(A)$ z Twierdzenia~\ref{thm:choice-hartogs}.

Zbierzmy wszystkie funkcje $f:I\to A$, gdzie $I$ jest
odcinkiem początkowym $H$, i które spełniają dla każdego
$t\in I$ warunek
\[
 f(t)=c\bigl(A\setminus f[I_{<t}]\bigr).
\tag{$*$}
\]
Wymagamy przy tym, aby argument funkcji $c$ był niepusty.
Wszystkie takie funkcje tworzą zbiór: jest to
podzbiór $\mathcal P(H\times A)$ wyznaczony przez
formułę o wykresie funkcji, odcinku początkowym
i równaniu ($*$), z parametrami $H,A,c$.
Powstaje przez separację. Rodzina nie jest pusta,
ponieważ zawiera funkcję pustą o dziedzinie
pustego odcinka początkowego.
Dwie takie funkcje muszą zgadzać się na wspólnej dziedzinie:
gdyby nie, w najmniejszym miejscu rozbieżności miałyby
identyczne wcześniejsze wartości i przez ($*$) dostałyby
tę samą nową wartość. Ich suma $f$ jest więc funkcją
o dziedzinie $I$ będącej odcinkiem początkowym $H$
i nadal spełnia ($*$).

Gdyby $A\setminus f[I]$ było niepuste, a $I\ne H$,
można byłoby przedłużyć $f$ o najmniejszy element
$H\setminus I$, przypisując mu $c(A\setminus f[I])$.
To przeczyłoby temu, że $f$ jest sumą \emph{wszystkich}
dopuszczalnych funkcji. Natomiast jeśli $I=H$,
to $f$ jest iniekcją $H\to A$, gdyż każde nowe
wybrane $f(t)$ leży poza zbiorem wcześniejszych wartości,
co jest sprzeczne z twierdzeniem Hartogsa. Wobec tego
$A\setminus f[I]=\varnothing$, czyli $f$ jest bijekcją
$I\to A$. Przenosimy dobry porządek $I$ przez $f$
na $A$ i kończymy dowód.
\end{proof}

\section{Lemat Kuratowskiego--Zorna i jego dowód}
\begin{definicja}[Łańcuch i element maksymalny]\label{def:choice-chain}
W częściowo uporządkowanym zbiorze $(P,\le)$
\emph{łańcuchem} jest podzbiór $C\subseteq P$, w którym
dowolne dwa elementy są porównywalne. \emph{Ograniczenie
górne} łańcucha $C$ to element $u\in P$ spełniający
$\forall c\in C\ (c\le u)$. Element $m\in P$ jest
\emph{maksymalny}, jeśli nie istnieje $p\in P$ z $m<p$.
Są to pojęcia z rozdziału~\ref{ch:orders};
\emph{maksymalny} nie musi oznaczać \emph{największy}.
\end{definicja}

\begin{twierdzenie}[Lemat Kuratowskiego--Zorna]\label{thm:choice-zorn}
Niech $(P,\le)$ będzie \emph{zbiorem częściowo uporządkowanym}.
Załóżmy, że \emph{każdy} łańcuch $C\subseteq P$
ma ograniczenie górne \emph{należące do $P$}.
Założenie dotyczy również łańcucha pustego $C=\varnothing$.
Wówczas istnieje element maksymalny $m\in P$, czyli
\[
\exists m\in P\quad\forall p\in P\;(m\le p\ \Longrightarrow\ p=m).
\]
Nie żądamy istnienia elementu \emph{największego},
który musiałby być większy lub równy \emph{każdemu} $p\in P$.
\end{twierdzenie}
\begin{proof}
Założymy zasadę dobrego uporządkowania,
którą uzyskaliśmy z AC w Twierdzeniu~\ref{thm:choice-wellorder}.
Rozdzielmy argument na pięć kroków.

\emph{Krok 1: co właściwie chcemy skonstruować?}
Zdobędziemy łańcuch $C$ \emph{maksymalny względem zawierania}:
nie ma w $P$ żadnego większego łańcucha $C'\supsetneq C$.
Gdy go już znajdziemy, z założenia twierdzenia dostaniemy
górne ograniczenie $u\in P$ dla $C$.
Pokażemy wtedy, że $u$ jest szukanym elementem maksymalnym $P$.

\emph{Krok 2: kolejność rozpatrywania elementów.}
Z założenia lematu Zorna zbiór $P$ jest niepusty:
łańcuch pusty musi mieć ograniczenie górne w $P$.
Ustalmy dowolny dobry porządek $\triangleleft$ na $P$.
Nie musimy znać jego konkretnej postaci; wiemy jedynie,
że każdy niepusty podzbiór $P$ ma względem niego
pierwszy element. Porządek $\triangleleft$ służy tylko
do \emph{wyboru kolejności} i nie należy go mylić
z oryginalnym porządkiem częściowym $\le$.

\emph{Krok 3: przyjmowanie i odrzucanie kandydatów.}
Przechodzimy przez elementy $p\in P$ w kolejności
$\triangleleft$. Oznaczamy $p$ jako \emph{przyjęty},
jeśli dla każdego wcześniej przyjętego $q$
zachodzi $p\le q$ lub $q\le p$.
W przeciwnym razie odrzucamy $p$.
Na pierwszym etapie nie ma wcześniejszych kandydatów,
więc pierwszy element zostaje przyjęty.

Dla ścisłości: rozpatrujemy funkcje
$s:I\to\{0,1\}$, których dziedziny $I$ są odcinkami
początkowymi porządku $\triangleleft$, i które spełniają
\[
s(p)=1\ \Longleftrightarrow\
\forall q\in I\,(q\triangleleft p\land s(q)=1
\ \Longrightarrow\ (q\le p\lor p\le q)).
\]
Rodzina wszystkich takich funkcji jest zbiorem:
otrzymujemy ją przez separację
w $\mathcal P(P\times\{0,1\})$.
Zawiera funkcję pustą, która spełnia regułę
na pustym odcinku początkowym.
Dwie funkcje spełniające tę regułę muszą zgadzać się
na wspólnej dziedzinie: w pierwszym punkcie potencjalnej
różnicy mają identyczne wcześniej przyjęte elementy,
a zatem reguła narzuca tę samą wartość $0$ lub $1$.
Suma wszystkich takich funkcji jest zatem funkcją;
jej dziedzina jest odcinkiem początkowym.
Gdyby nie obejmowała całego $P$, moglibyśmy dodać
najmniejszy pominięty element (w porządku $\triangleleft$),
wyznaczając mu wartość przez powyższą regułę.
Przeczyłoby to temu, że wzięliśmy sumę wszystkich
możliwych poprawnych funkcji.
Otrzymujemy więc funkcję $s:P\to\{0,1\}$.

\emph{Krok 4: przyjęte elementy tworzą maksymalny łańcuch.}
Niech $C=\{p\in P:s(p)=1\}$.
Jeśli $p,q\in C$ i $p\triangleleft q$, to w chwili
przyjmowania $q$ element $p$ był już przyjęty.
Stąd $p\le q$ lub $q\le p$. Zatem $C$ jest łańcuchem.

Załóżmy, że istnieje większy łańcuch $C'$ z
$C\subsetneq C'$. Wybierzmy $v\in C'\setminus C$.
Każdy element $q\in C$ jest porównywalny z $v$,
bo $v,q\in C'$ i $C'$ jest łańcuchem.
Dotyczy to w szczególności każdego elementu
przyjętego \emph{przed} rozpatrzeniem $v$.
Reguła nakazywałaby zatem przyjąć $v$,
czyli $v\in C$. Sprzeczność. Łańcuch $C$
jest maksymalny względem zawierania.

\emph{Krok 5: od maksymalnego łańcucha do elementu maksymalnego.}
Z założenia twierdzenia istnieje $u\in P$ takie, że
$c\le u$ dla każdego $c\in C$.
Wtedy $C\cup\{u\}$ też jest łańcuchem:
element $u$ jest porównywalny ze wszystkimi $c\in C$.
Maksymalność $C$ wymusza $u\in C$.

Przypuśćmy teraz, że $u$ nie jest elementem maksymalnym
w $P$. Istniałby element $v\in P$ z $u<v$.
Ponieważ $c\le u<v$ dla każdego $c\in C$,
zbiór $C\cup\{v\}$ byłby łańcuchem ściśle
większym od $C$. Sprzeczność.
Zatem $u$ jest maksymalny w $P$.
\end{proof}

\begin{przyklad}[Dwie moce można porównać za pomocą iniekcji]
Niech $A,B$ będą zbiorami. Niech $P$ będzie zbiorem
wszystkich \emph{częściowych iniekcji} $f:D\to B$, gdzie
$D\subseteq A$, a porządkiem będzie przedłużanie funkcji.
Zbiór $P$ istnieje przez separację w $\mathcal P(A\times B)$:
każdy jego element jest wykresem funkcji różnowartościowej
o dziedzinie zawartej w $A$.
Wyjaśnijmy najpierw, dlaczego każdy łańcuch $C\subseteq P$
ma ograniczenie górne: suma $u=\bigcup C$ jest funkcją,
bo elementy łańcucha zgadzają się na wspólnych dziedzinach.
Jest też iniekcją: jeżeli $u(a)=u(a')$, to obie wartości
są zadane przez pewne funkcje z $C$; jedna z nich
przedłuża drugą, więc obie pary należą do tej samej
funkcji różnowartościowej. Stąd $a=a'$.
Wynika więc $u\in P$, a przy $C=\varnothing$ górnym
ograniczeniem jest funkcja pusta.

Lemat Zorna daje maksymalną częściową iniekcję
$f:D\to B$. Gdyby równocześnie $A\setminus D\ne\varnothing$
i $B\setminus f[D]\ne\varnothing$, moglibyśmy wybrać
po jednym elemencie $a\in A\setminus D$ i
$b\in B\setminus f[D]$ i przedłużyć $f$ przez
$f(a):=b$. Nadal mielibyśmy iniekcję, wbrew
maksymalności. Zatem albo $D=A$, i mamy iniekcję
$A\to B$, albo $f[D]=B$, i funkcja odwrotna $f^{-1}$
jest iniekcją $B\to D\subseteq A$.

Wniosek jest nieintuicyjnie silny: przy założeniu AC
dowolne dwa zbiory można porównać pod względem
istnienia iniekcji w jedną albo w drugą stronę.
Bez AC taka porównywalność nie musi zachodzić.
\end{przyklad}


\begin{przyklad}[Co naprawdę daje lemat Zorna?]
Weźmy $P=\{a,b\}$, w którym jedynymi porównaniami
są $a\le a$ i $b\le b$. Elementy $a,b$ są
nieporównywalne. Wszystkie łańcuchy to
$\varnothing$, $\{a\}$ i $\{b\}$.
Łańcuch pusty ma ograniczenie górne, np. $a$;
łańcuchy jednoelementowe mają ograniczenia
odpowiednio $a$ i $b$. Zatem założenia Zorna są spełnione.
Oba elementy $a,b$ są \emph{maksymalne},
ale żaden nie jest \emph{największy}.
Gdyby ktoś wywnioskował z lematu Zorna istnienie
elementu największego, ten przykład natychmiast
obaliłby taki wniosek.
\end{przyklad}

\section{Zasada maksymalności Hausdorffa}\label{sec:choice-hausdorff}

W dowodzie Twierdzenia~\ref{thm:choice-zorn}
zaczęliśmy od konstruowania maksymalnego \emph{łańcucha},
a dopiero potem znaleźliśmy element maksymalny.
Ten pośredni krok jest treścią osobnej zasady,
noszącej nazwisko Hausdorffa. Najpierw dokładnie
rozróżnijmy, co znaczy maksymalność \emph{łańcucha}
i maksymalność \emph{elementu} porządku.

\begin{definicja}[Łańcuch maksymalny]\label{def:hausdorff-maximal-chain}
Niech $(P,\le)$ będzie zbiorem częściowo uporządkowanym.
Łańcuch $C\subseteq P$ nazywamy \emph{maksymalnym}, jeśli
\[
\forall D\subseteq P\quad
 (D\text{ jest łańcuchem i }C\subseteq D\ \Longrightarrow\ D=C).
\]
Jest to maksymalność w porządku \emph{zawierania łańcuchów},
a nie istnienie jednego największego elementu w $C$.
\end{definicja}

\begin{definicja}[Zasada maksymalności Hausdorffa]\label{def:hausdorff-principle}
\emph{Zasada maksymalności Hausdorffa} głosi:
\emph{każdy zbiór częściowo uporządkowany ma
łańcuch maksymalny względem zawierania}.
Dla $P=\varnothing$ jedynym łańcuchem jest
$\varnothing$ i jest on maksymalny. Nie zakładamy
istnienia ograniczeń górnych łańcuchów.
\end{definicja}

\begin{przyklad}[Dwa różne łańcuchy maksymalne]
W zbiorze $P=\{a,b,c\}$ przyjmijmy $a<c$ i $b<c$,
a $a,b$ niech będą nieporównywalne.
Wówczas $C_1=\{a,c\}$ oraz $C_2=\{b,c\}$
są łańcuchami maksymalnymi. Nie możemy do $C_1$
dołożyć $b$, ponieważ $a$ i $b$ nie dają się porównać.
Analogicznie nie możemy do $C_2$ dołożyć $a$.
Łańcuch $\{c\}$ \emph{nie} jest maksymalny:
możemy powiększyć go do $C_1$.
Przykład pokazuje, że łańcuch maksymalny
nie musi być jednoznaczny.
\end{przyklad}

\begin{twierdzenie}[Równoważność Hausdorffa i Zorna nad ZF]\label{thm:hausdorff-zorn-equivalent}
W teorii ZF następujące dwa zdania są równoważne:
\begin{enumerate}[label=(\roman*)]
\item zasada maksymalności Hausdorffa;
\item lemat Kuratowskiego--Zorna.
\end{enumerate}
\end{twierdzenie}
\begin{proof}
\emph{Kierunek Zorn $\Rightarrow$ Hausdorff.}
Niech $(P,\le)$ będzie \emph{dowolnym} zbiorem
częściowo uporządkowanym. Zdefiniujmy
\[
\mathcal C:=\{C\subseteq P:C\text{ jest łańcuchem}\}.
\]
Rodzina $\mathcal C$ jest zbiorem, ponieważ jest podzbiorem
$\mathcal P(P)$ określonym przez własność łańcucha.
Jest także niepusta: $\varnothing\in\mathcal C$.
Uporządkujmy ją przez zawieranie $\subseteq$.
Chcemy zastosować Zorna do $(\mathcal C,\subseteq)$,
więc musimy wykazać jego założenie \emph{dla każdego
łańcucha w $\mathcal C$}.

Niech $\mathcal D\subseteq\mathcal C$ będzie łańcuchem
względem zawierania. To znaczy, że dla dowolnych
$C_1,C_2\in\mathcal D$ zachodzi
$C_1\subseteq C_2$ lub $C_2\subseteq C_1$.
Połóżmy $U:=\bigcup\mathcal D$.
Dla dowolnych $x,y\in U$ istnieją łańcuchy
$C_x,C_y\in\mathcal D$ takie, że $x\in C_x$ i $y\in C_y$.
Ponieważ $\mathcal D$ jest łańcuchem w porządku zawierania,
jeden z tych dwóch łańcuchów zawiera drugi.
Zatem zarówno $x$, jak i $y$ należą do jednego
łańcucha z $\mathcal D$, więc są porównywalne.
Wynika stąd $U\in\mathcal C$.
Ponadto $C\subseteq U$ dla każdego $C\in\mathcal D$,
czyli $U$ jest \emph{ograniczeniem górnym} $\mathcal D$
w zbiorze $\mathcal C$. Dla $\mathcal D=\varnothing$
mamy $U=\varnothing\in\mathcal C$, więc warunek
obejmuje również łańcuch pusty.

Stosujemy Zorna i otrzymujemy element
maksymalny $C_\ast\in\mathcal C$ względem zawierania.
Z definicji jest to dokładnie łańcuch maksymalny
w $P$. Otrzymaliśmy zasadę Hausdorffa.

\emph{Kierunek Hausdorff $\Rightarrow$ Zorn.}
Niech $(P,\le)$ spełnia założenie lematu Zorna:
każdy łańcuch ma ograniczenie górne należące do $P$.
Z zasady Hausdorffa otrzymujemy łańcuch maksymalny
$C\subseteq P$. Według założenia Zorna istnieje
$u\in P$ takie, że $c\le u$ dla każdego $c\in C$.
Rozważmy $C\cup\{u\}$.
Wszystkie pary elementów $C$ są porównywalne,
a każdy $c\in C$ spełnia $c\le u$.
Dlatego $C\cup\{u\}$ jest łańcuchem.
Z maksymalności $C$ wynika $u\in C$.

Gdyby $u$ nie był maksymalny w $P$, znaleźlibyśmy
$v\in P$ z $u<v$. Z przechodniości wynika
$c\le u<v$ dla każdego $c\in C$, a zatem
$C\cup\{v\}$ byłby łańcuchem właściwie
zawierającym $C$. To sprzeczne z maksymalnością $C$.
Wnioskujemy, że $u$ jest maksymalnym elementem $P$.
Dowiedliśmy lematu Zorna.
\end{proof}

\begin{uwaga}[Co odróżnia zasady Hausdorffa i Zorna?]
Zasada Hausdorffa mówi o \emph{maksymalnym łańcuchu}
bez zakładania, że łańcuchy mają ograniczenia górne.
Lemat Zorna mówi o \emph{maksymalnym elemencie}
pod dodatkowym założeniem istnienia takich ograniczeń.
To nie są identyczne zdania, ale właśnie udowodniliśmy,
że nad ZF mają dokładnie tę samą siłę.
\end{uwaga}

\section{Powrót od lematu Zorna do aksjomatu wyboru}
\begin{twierdzenie}[Równoważność głównych zasad]\label{thm:choice-equivalence}
W teorii ZF następujące cztery zdania są równoważne:
(i) aksjomat wyboru z Definicji~\ref{def:choice};
(ii) każde zbiory można dobrze uporządkować;
(iii) lemat Kuratowskiego--Zorna;
(iv) zasada maksymalności Hausdorffa z Definicji~\ref{def:hausdorff-principle}.
\end{twierdzenie}
\begin{proof}
Dowody (i)$\Rightarrow$(ii) i (ii)$\Rightarrow$(iii)
przeprowadziliśmy odpowiednio w Twierdzeniach
\ref{thm:choice-wellorder} i~\ref{thm:choice-zorn}.
Wystarczy wykazać (iii)$\Rightarrow$(i);
równoważność (iii)$\Leftrightarrow$(iv) została udowodniona
w Twierdzeniu~\ref{thm:hausdorff-zorn-equivalent}.

Niech $(A_i)_{i\in I}$ będzie dowolną rodziną niepustych
zbiorów. Rozważmy zbiór $P$ wszystkich \emph{częściowych
funkcji wyboru}: do $P$ należy funkcja $f:D\to\bigcup_{i\in I}A_i$
z dziedziną $D\subseteq I$, jeżeli $f(i)\in A_i$
dla każdego $i\in D$. Jest to zbiór istniejący przez
separację w $\mathcal P(I\times\bigcup_{i\in I}A_i)$:
warunek bycia funkcją i warunek $f(i)\in A_i$
są wyrażalne w języku ZF.
Uporządkujmy $P$ przez
przedłużanie: $f\le g$ wtedy i tylko wtedy, gdy
dziedzina $g$ zawiera dziedzinę $f$ i $g$ zgadza się
z $f$ na jej dziedzinie. Funkcja pusta należy do $P$.

Jeśli $C\subseteq P$ jest łańcuchem, to $u:=\bigcup C$
jest funkcją: dwie funkcje porównywalne przez
przedłużanie nie przypiszą różnych wartości temu samemu
indeksowi. Ponadto $u(i)\in A_i$ dla każdego
$i\in\dom u$, zatem $u\in P$ i $u$ jest
ograniczeniem górnym $C$ (także dla $C=\varnothing$).
Lemat Zorna daje maksymalną funkcję $f\in P$.
Gdyby $\dom f\ne I$, można byłoby wybrać \emph{jeden}
indeks $i_0\in I\setminus\dom f$ i \emph{jeden}
element $a_0\in A_{i_0}$, a następnie przedłużyć
$f$ o parę $(i_0,a_0)$. To przeczyłoby maksymalności.
Zatem $\dom f=I$ i $f$ jest funkcją wyboru.
\end{proof}

\begin{uwaga}[Dlaczego nie jest to dowód wyboru bez wyboru?]
Dowodząc AC z Zorna, \emph{zakładaliśmy} lemat Zorna.
Dowodząc Zorna z zasady dobrego uporządkowania,
\emph{zakładaliśmy} istnienie dobrego porządku na całym $P$.
Dowodząc zasady dobrego uporządkowania, zaczęliśmy
od AC. Nie wyprowadziliśmy więc żadnej z tych zasad
z samej ZF; wykazaliśmy jedynie, że w obecności ZF
każdą z nich można zastąpić dwiema pozostałymi.
\end{uwaga}

\begin{przyklad}[Równoległa i warunkowa postać wyboru]
Dla rodziny $(A_i)_{i\in I}$ \emph{iloczynem kartezjańskim rodziny}
nazywamy zbiór
\[
\prod_{i\in I}A_i:=\{c:I\to\bigcup_{i\in I}A_i:
\forall i\in I\ (c(i)\in A_i)\}.
\]
Zbiór tych funkcji istnieje przez separację
w $\mathcal P(I\times\bigcup_{i\in I}A_i)$;
dla $I=\varnothing$ zawiera funkcję pustą.
Czasem aksjomat wyboru przedstawia się w postaci:
„iloczyn kartezjański dowolnej rodziny niepustych
zbiorów jest niepusty”. Rzeczywiście, jego element
jest dokładnie funkcją wyboru $c$.
Ta równoważność wynika z samej definicji iloczynu
rodziny; nie należy mylić jej z iloczynem
\emph{dwóch} zbiorów, który istnieje w ZF bez AC.
\end{przyklad}

\section{Kiedy zasada wyboru jest rzeczywiście potrzebna?}

\begin{twierdzenie}[Wybór ze skończonej rodziny w ZF]
\label{thm:choice-finite-family}
Niech $I$ będzie zbiorem skończonym,
a $(A_i)_{i\in I}$ rodziną zbiorów niepustych.
W ZF istnieje funkcja $c$ o dziedzinie $I$
spełniająca $c(i)\in A_i$ dla wszystkich $i\in I$.
Teza obejmuje także $I=\varnothing$.
\end{twierdzenie}
\begin{proof}
Najpierw udowodnimy przez indukcję po $n\in\omega$
zdanie: \emph{każda rodzina $(B_k)_{k\in n}$
zbiorów niepustych ma funkcję wyboru
$d:n\to\bigcup_{k\in n}B_k$}.
Dla $n=0$ wystarcza funkcja pusta.
Załóżmy tezę dla $n$ i weźmy dowolną rodzinę
$(B_k)_{k\in S(n)}$ niepustych zbiorów.
Jej ograniczenie $(B_k)_{k\in n}$ ma
funkcję wyboru $d_n$ na mocy hipotezy indukcyjnej.
Ponieważ $B_n$ jest niepusty,
istnieje \emph{jeden} element $a\in B_n$.
Dołączamy parę $(n,a)$ do $d_n$.
Jest to funkcja na $S(n)=n\cup\{n\}$,
gdyż $n\notin n$
(Lemat~\ref{lem:zf-vn-successor}),
a każda jej wartość należy do odpowiedniego $B_k$.
Otrzymaliśmy tezę dla $S(n)$.

Teraz wróćmy do danej rodziny $(A_i)_{i\in I}$.
Ponieważ $I$ jest skończony, istnieją
$n\in\omega$ i bijekcja $e:n\to I$.
Zastosujmy przed chwilą dowiedzioną tezę
do rodziny $(B_k)_{k\in n}$ określonej przez
$B_k:=A_{e(k)}$. Otrzymujemy funkcję wyboru
$d:n\to\bigcup_{i\in I}A_i$.
Definiujemy $c(i):=d(e^{-1}(i))$ dla każdego
$i\in I$. Wtedy $c(i)\in A_i$.
W żadnym kroku nie zastosowaliśmy AC
do rodziny nieskończonej.
\end{proof}

\begin{przyklad}[Rodzina pusta i powtarzające się zbiory]
\label{ex:choice-empty-family}
Dla $I=\varnothing$ funkcja pusta
$\varnothing\to\varnothing$ jest funkcją wyboru.
Jeżeli $I=\{0,1,2\}$,
$A_0=A_1=\{a,b\}$ i $A_2=\{b\}$,
możemy wybrać $c(0)=a$, $c(1)=b$,
$c(2)=b$. Wartości funkcji wyboru
nie muszą być parami różne.
\end{przyklad}

\begin{twierdzenie}[Aksjomat wyboru a przekroje surjekcji]
\label{thm:choice-sections}
Nad ZF następujące warunki są równoważne:
\begin{enumerate}[label=(\roman*)]
\item obowiązuje aksjomat wyboru;
\item dla każdej surjekcji $p:E\to I$
istnieje funkcja $s:I\to E$ spełniająca
$p\circ s=\operatorname{id}_I$.
\end{enumerate}
Oba warunki obejmują puste zbiory.
\end{twierdzenie}
\begin{proof}
Załóżmy AC i niech $p:E\to I$ będzie surjekcją.
Dla $i\in I$ włókno
$E_i:=\{e\in E:p(e)=i\}$ jest niepuste.
Każde włókno istnieje przez separację,
a funkcja $i\mapsto E_i$ jest zbiorem przez
zastępowanie. Z AC istnieje
$s:I\to E$ takie, że $s(i)\in E_i$.
Zatem $p(s(i))=i$ dla każdego $i\in I$.

Odwrotnie, ustalmy rodzinę $(A_i)_{i\in I}$
zbiorów niepustych i jej sumę rozłączną
\[
E:=\{(i,a):i\in I,\ a\in A_i\}
\subseteq I\times\bigcup_{i\in I}A_i.
\]
Jest to zbiór zbudowany w ZF.
Rzut $p:E\to I$, $p(i,a)=i$,
jest surjekcją, ponieważ każde $A_i$ jest niepuste.
Przekrój $s$ tego rzutu musi mieć postać
$s(i)=(i,c(i))$ dla jednoznacznego
$c(i)\in A_i$. Schemat zastępowania daje
wykres funkcji $c$, a więc funkcję wyboru.
Dla $I=\varnothing$ zarówno suma rozłączna,
jak i przekrój są funkcjami pustymi.
\end{proof}

\begin{wniosek}[Transwersale relacji równoważności]
\label{cor:choice-transversals}
Nad ZF aksjomat wyboru jest równoważny temu,
że dla każdej relacji równoważności $\sim$
na zbiorze $X$ istnieje podzbiór $T\subseteq X$
zawierający dokładnie jeden element każdej klasy
równoważności ze zbioru ilorazowego $X/{\sim}$.
\end{wniosek}
\begin{proof}
Przy AC stosujemy Twierdzenie~\ref{thm:choice-sections}
do surjekcji ilorazowej $q:X\to X/{\sim}$.
Jej przekrój $s$ daje
$T:=s[X/{\sim}]$, przy czym
$q(s(C))=C$ dla każdej klasy $C$.
Zatem $T$ ma dokładnie jednego reprezentanta
z każdej klasy.

W drugą stronę dla dowolnej rodziny
$(A_i)_{i\in I}$ zbiorów niepustych
bierzemy sumę rozłączną
$E:=\{(i,a):i\in I,\ a\in A_i\}$
i relację $(i,a)\sim(j,b)$
wtedy i tylko wtedy, gdy $i=j$.
Jeżeli $T$ zawiera po jednym elemencie
każdej klasy, to jego jedyny element
we włóknie nad $i$ ma postać
$(i,c(i))$ z $c(i)\in A_i$.
Otrzymujemy funkcję wyboru $c$.
Dla $E=\varnothing$ transwersalą jest
zbiór pusty.
\end{proof}

\begin{przyklad}[Dlaczego suma musi być rozłączna?]
\label{ex:choice-disjoint-index}
Jeżeli $A_0=A_1=\{a\}$,
zwykła suma to tylko $\{a\}$,
która nie odnotowuje dwóch indeksów.
Suma rozłączna
$E=\{(0,a),(1,a)\}$ ma dwa elementy,
a jej włókna pamiętają, do którego
indeksu należy wybrany element.
\end{przyklad}

\section{Przykład zastosowania: istnienie bazy przestrzeni wektorowej}

\begin{twierdzenie}[Istnienie bazy z lematu Zorna]
\label{thm:choice-hamel-basis}
Niech $K$ będzie ciałem, a $V$ przestrzenią
wektorową nad $K$. Zakładając AC, istnieje
$B\subseteq V$ takie, że każdy $v\in V$
ma dokładnie jedno przedstawienie
\[
v=\sum_{b\in B}\lambda(b)b,
\]
gdzie $\lambda:B\to K$ jest funkcją
o skończonym nośniku
$\{b\in B:\lambda(b)\ne0\}$,
a sumowanie przebiega jedynie
po tym skończonym nośniku.
Zbiór $B$ nazywamy bazą przestrzeni $V$.
\end{twierdzenie}
\begin{proof}
Niech $P$ będzie zbiorem liniowo
niezależnych podzbiorów $C\subseteq V$.
Jest to zbiór otrzymany przez separację
w $\mathcal P(V)$, gdyż warunek
\[
\sum_{j=1}^r\lambda_jc_j=0
\Longrightarrow \lambda_1=\cdots=\lambda_r=0
\]
dla wszystkich skończonych list parami
różnych $c_j\in C$ i skalarów
$\lambda_j\in K$ można wyrazić w ZF.
Zbiór $P$ jest niepusty, bo $\varnothing\in P$.
Uporządkujmy go przez inkluzję.

Dla dowolnego łańcucha $\mathcal C\subseteq P$
zbiór $U:=\bigcup\mathcal C$ jest liniowo
niezależny. Istotnie, dowolne skończenie
wiele elementów $u_1,\ldots,u_r\in U$
należy łącznie do jednego zbioru
$C_*\in\mathcal C$: wybieramy
skończenie wiele zbiorów łańcucha
zawierających poszczególne $u_j$,
a spośród nich największy względem inkluzji.
Niezależność $C_*$ daje niezależność
tej skończonej rodziny. Dla pustego łańcucha
mamy $U=\varnothing\in P$.
Zatem każdy łańcuch ma górne ograniczenie w $P$.

Lemat Kuratowskiego--Zorna,
Twierdzenie~\ref{thm:choice-zorn},
daje maksymalny element $B\in P$.
Jeśli istniałby $v\in V$ poza
powłoką liniową $B$, to $B\cup\{v\}$
byłby liniowo niezależny.
Rzeczywiście, w każdej relacji
$\lambda v+\sum_{j=1}^r\mu_jb_j=0$
współczynnik $\lambda$ musi być zerowy,
gdyż w przeciwnym razie
$v=-\lambda^{-1}\sum_j\mu_jb_j$
należałby do powłoki liniowej $B$.
Przy $\lambda=0$ niezależność $B$
wymusza $\mu_j=0$ dla wszystkich $j$.
Byłoby to sprzeczne z maksymalnością $B$.
Stąd $B$ generuje $V$.

Jeżeli dwie skończone kombinacje
elementów $B$ przedstawiają ten sam wektor,
odejmujemy je i grupujemy współczynniki
przy tych samych wektorach.
Otrzymujemy skończoną relację liniową
równą zeru. Niezależność $B$
wymusza zerowanie wszystkich współczynników
różnicy, czyli jednoznaczność funkcji $\lambda$.
\end{proof}

\begin{przyklad}[Przestrzeń zerowa]\label{ex:choice-zero-space}
Jeśli $V=\{0\}$, zbiór pusty jest
liniowo niezależny i generuje $V$:
jego jedyna skończona kombinacja liniowa
to pusta suma, czyli wektor zerowy.
Zatem bazą $V$ jest $\varnothing$.
\end{przyklad}

\chapter{Liczby kardynalne i arytmetyka mocy}\label{ch:cardinals}
\chaptermark{Liczby kardynalne i arytmetyka mocy}

\section{Dlaczego porównujemy zbiory za pomocą funkcji?}
Dwa zbiory mają tę samą liczbę elementów, gdy istnieje między nimi bijekcja. Definicja działa również dla zbiorów nieskończonych, których elementów nie można po prostu policzyć. W całym rozdziale pracujemy w ZFC, chyba że zaznaczono inaczej. Rozróżniamy twierdzenia prawdziwe już w ZF od tych, które wymagają aksjomatu wyboru.
\begin{definicja}[Równoliczność i porównanie mocy]\label{def:card-equipotent}
Dla zbiorów $A,B$ piszemy $A\approx B$, gdy istnieje bijekcja $A\to B$, oraz $|A|\leq |B|$, gdy istnieje iniekcja $A\to B$. Zapis $|A|<|B|$ oznacza $|A|\leq |B|$ i $A\not\approx B$. Symbole $|A|$ są na razie nazwami typów równoliczności; formalną reprezentację przez liczby kardynalne podamy po wprowadzeniu liczb porządkowych w rozdziale~\ref{ch:ordinals}.
\end{definicja}
\begin{twierdzenie}[Zawieranie zbiorów a porównanie mocy]\label{thm:card-subset-monotonicity}
Niech $A,B$ będą dowolnymi zbiorami. Jeżeli $A\subseteq B$, to $|A|\leq |B|$.
Nie można na ogół zastąpić znaku $\leq$ znakiem $<$, nawet jeśli $A\subsetneq B$.
Jeżeli jednak $B$ jest skończony i $A\subsetneq B$, to $|A|<|B|$.
\end{twierdzenie}
\begin{proof}
\emph{Krok 1: skonstruowanie iniekcji.}
Ponieważ $A\subseteq B$, dla każdego $a\in A$ element $a$ należy również do $B$.
Możemy zatem określić funkcję $i:A\to B$ wzorem $i(a)=a$.
Nazywamy ją \emph{włączeniem}. Jest iniekcją: jeżeli $i(a)=i(a')$,
to $a=a'$. Z definicji porównania mocy wynika $|A|\leq |B|$.
W tym kroku nie użyliśmy ani skończoności, ani aksjomatu wyboru.

\emph{Krok 2: dlaczego nierówność nie musi być ostra.}
Dla $A=\N\setminus\{0\}$ i $B=\N$ mamy $A\subsetneq B$, lecz
funkcja $f:\N\to A$, $f(n)=n+1$, jest bijekcją.
Istotnie, jest iniekcją, a każdy $m\in A$ jest postacią
$m=f(m-1)$. Dlatego $|A|=|B|$ mimo właściwego zawierania.

\emph{Krok 3: dlaczego sytuacja skończona jest inna.}
Niech $B$ ma dokładnie $n$ elementów, a $A\subsetneq B$.
Ponieważ $A$ jest właściwym podzbiorem $B$, mamy $n\ge1$.
Wybierzmy $b\in B\setminus A$. Wówczas $A\subseteq B\setminus\{b\}$.
Po ponumerowaniu elementów $B$ od $1$ do $n$ zbiór $B\setminus\{b\}$
ma dokładnie $n-1$ elementów. Zatem $A$ ma co najwyżej $n-1$ elementów,
a $B$ dokładnie $n$. Nie może istnieć bijekcja $A\to B$:
obrazy co najwyżej $n-1$ różnych elementów nie pokryją $n$-elementowego zbioru.
Łącznie z krokiem pierwszym daje to $|A|<|B|$.
\end{proof}
\begin{przyklad}[Co mówi znak nierówności?]
Dla $A=\{1,3\}$ i $B=\{1,2,3\}$ zachodzi $A\subsetneq B$ oraz
$|A|=2<3=|B|$. Natomiast dla $A=\N\setminus\{0\}$ i $B=\N$
mamy również $A\subsetneq B$, ale $|A|=|B|=\aleph_0$.
W szczególności stwierdzenie „zbiór ma mniejszą moc niż jego podzbiór”
ma odwrócony kierunek: podzbiór nie może mieć \emph{większej}
mocy niż zbiór, którego jest częścią.
\end{przyklad}

\section{Twierdzenie Cantora--Bernsteina}
\begin{twierdzenie}[Cantora--Bernsteina]\label{thm:card-cb}
Niech $A,B$ będą zbiorami. Jeżeli istnieją iniekcje $f:A\to B$ i $g:B\to A$, to istnieje bijekcja $A\to B$. Dowód jest poprawny w ZF, bez aksjomatu wyboru.
\end{twierdzenie}
\begin{proof}
Niech $A_0=A\setminus g[B]$. Na zbiorze potęgowym $\mathcal P(A)$
określmy funkcję $T(X)=g[f[X]]$. Jest ona dobrze określona, bo
$f[X]\subseteq B$, a $g[f[X]]\subseteq A$.
Twierdzenie o rekurencji na $\omega$
(Twierdzenie~\ref{thm:zf-omega-recursion}), zastosowane do
$B_0=\mathcal P(A)$, wartości początkowej $A_0$
i funkcji $(n,X)\mapsto T(X)$, daje jednoznaczny ciąg
$(A_n)_{n\in\N}$ spełniający $A_{n+1}=g[f[A_n]]$.
Połóżmy $C=\bigcup_{n\in\N}A_n$.
Jest to zbiór dzięki zastępowaniu i aksjomatowi sumy.
Ponieważ obraz sumy przez funkcję jest sumą obrazów,
\[
g[f[C]]=\bigcup_{n\in\N}g[f[A_n]]
=\bigcup_{n\in\N}A_{n+1}\subseteq C.
\]
Ponadto $A_0\subseteq C$ oraz każdy składnik
$A_{n+1}$ należy do $g[f[C]]$, stąd
\[
C=A_0\cup g[f[C]].
\]
Na $C$ zdefiniujmy $h(a)=f(a)$. Jeśli $a\in A\setminus C$, to $a\notin A_0$, a więc $a\in g[B]$. Istnieje dokładnie jedno $b\in B$ z $g(b)=a$; kładziemy $h(a)=b$. W obu przypadkach $h$ jest dobrze określona.

Sprawdzimy, że jej dwie części mają rozłączne obrazy. Obraz pierwszej części to $f[C]$. Obraz drugiej to $g^{-1}[A\setminus C]$. Gdyby $b=f(c)$ należało do obu obrazów, mielibyśmy $g(b)=g(f(c))\in g[f[C]]\subseteq C$, sprzecznie z definicją drugiej części. Każda część jest iniekcją, więc $h$ jest iniekcją.

Dla surjektywności weźmy $b\in B$. Jeśli $b\in f[C]$, to należy do obrazu pierwszej części. W przeciwnym razie $g(b)\notin C$: istotnie, gdyby $g(b)\in C$, to nie mógłby należeć do $A_0$, a z równości $C=A_0\cup g[f[C]]$ wynikałoby $g(b)\in g[f[C]]$; iniektywność $g$ dałaby $b\in f[C]$. Zatem $g(b)\in A\setminus C$ i $h(g(b))=b$. Otrzymaliśmy surjekcję.
\end{proof}
\begin{wniosek}[Antysymetria porównania mocy]\label{cor:card-antisym}
Dla dowolnych zbiorów $A,B$, jeśli
$|A|\le |B|$ i $|B|\le |A|$,
to $A\approx B$.
\end{wniosek}
\begin{proof}
Pierwsza nierówność oznacza istnienie
iniekcji $f:A\to B$, a druga
iniekcji $g:B\to A$.
Twierdzenie~\ref{thm:card-cb}
daje z tych dwóch iniekcji bijekcję $A\to B$,
co jest definicją $A\approx B$.
\end{proof}
\begin{uwaga}[Od iniekcji do jawnej bijekcji]
Twierdzenie Cantora--Bernsteina mówi, że wystarczą iniekcje w obu kierunkach.
W dalszej części rozdziału skonstruujemy nawet \emph{jawne} bijekcje
$\R\to(0,1)$ i $(0,1)\to(0,1]$, z kontrolą ich odwrotności.
Pozwoli to zobaczyć, jak nieskończony zbiór może być równoliczny
ze swoim właściwym podzbiorem.
\end{uwaga}

\section{Sumy, iloczyny i potęgowanie mocy}
\begin{definicja}[Operacje na mocach]\label{def:card-operations}
Dla zbiorów $A,B$ określamy sumę rozłączną $A\sqcup B=(A\times\{0\})\cup(B\times\{1\})$. Piszemy
\[
|A|+|B|:=|A\sqcup B|,\qquad |A|\cdot|B|:=|A\times B|,
\qquad |B|^{|A|}:=|B^A|,
\]
gdzie $B^A$ oznacza zbiór wszystkich funkcji $A\to B$.
Wykres każdej takiej funkcji jest podzbiorem $A\times B$;
stąd zbiór $B^A$ istnieje w ZF przez separację
w $\mathcal P(A\times B)$.
Podobnie istnieją $A\sqcup B$ i $A\times B$.

Definicje zależą tylko od typów równoliczności.
Istotnie, jeśli $\alpha:A\to A'$ oraz
$\beta:B\to B'$ są bijekcjami, to
$(a,0)\mapsto(\alpha(a),0)$,
$(b,1)\mapsto(\beta(b),1)$
wyznaczają bijekcję sum rozłącznych,
a $(a,b)\mapsto(\alpha(a),\beta(b))$
--- bijekcję iloczynów kartezjańskich.
Wreszcie funkcji $h:A\to B$ przypisujemy
$\beta\circ h\circ\alpha^{-1}:A'\to B'$.
Odwrotne przyporządkowanie ma postać
$k\mapsto\beta^{-1}\circ k\circ\alpha$,
więc również zbiory funkcji są równoliczne.
Nie użyto aksjomatu wyboru.
\end{definicja}
\begin{twierdzenie}[Podstawowe prawa arytmetyki mocy]\label{thm:card-laws}
Dla dowolnych zbiorów $A,B,C$ zachodzą równości
\[
|A|+|B|=|B|+|A|,\quad (|A|+|B|)+|C|=|A|+(|B|+|C|),
\]
\[
|A|\cdot(|B|+|C|)=|A|\cdot|B|+|A|\cdot|C|,
\quad |C|^{|A|+|B|}=|C|^{|A|}\cdot|C|^{|B|}.
\]
\end{twierdzenie}
\begin{proof}
Każde prawo wykażemy przez konkretną bijekcję, bez AC.

\emph{Przemienność sumy.}
Funkcja $\sigma:A\sqcup B\to B\sqcup A$
określona przez $\sigma(a,0)=(a,1)$
oraz $\sigma(b,1)=(b,0)$ zamienia znaczniki.
Jej ponowne zastosowanie daje funkcję identycznościową,
więc jest bijekcją.

\emph{Łączność sumy.}
Zdefiniujmy odwzorowanie
$(A\sqcup B)\sqcup C\to A\sqcup(B\sqcup C)$ przez
\[
((a,0),0)\longmapsto(a,0),\qquad
((b,1),0)\longmapsto((b,0),1),\qquad
(c,1)\longmapsto((c,1),1).
\]
Trzy rodzaje elementów dziedziny są rozłączne.
Każdy element przeciwdziedziny ma dokładnie jedną
z trzech przedstawionych postaci, zatem odwzorowanie
jest bijekcją.

\emph{Rozdzielność.}
Odwzorowanie $d:A\times(B\sqcup C)\to
(A\times B)\sqcup(A\times C)$
definiujemy wzorami
$d(a,(b,0))=((a,b),0)$ oraz
$d(a,(c,1))=((a,c),1)$.
Jest iniekcją dzięki jednoznaczności par uporządkowanych
i rozłącznym znacznikom. Jest surjekcją, ponieważ
każdy element przeciwdziedziny ma jedną z podanych postaci.

\emph{Potęgowanie sumy.}
Funkcji $u:A\sqcup B\to C$ przypisujemy parę
$(u_A,u_B)$ funkcji $u_A:A\to C$, $u_B:B\to C$
określonych przez $u_A(a)=u(a,0)$ i $u_B(b)=u(b,1)$.
Ta para jednoznacznie wyznacza $u$,
więc odwzorowanie $u\mapsto(u_A,u_B)$ jest iniekcją.
Dla dowolnych funkcji $v:A\to C$ i $w:B\to C$
definiujemy $u(a,0)=v(a)$ i $u(b,1)=w(b)$.
To funkcja na rozłącznej sumie, będąca przeciwobrazem
$(v,w)$; otrzymaliśmy również surjekcję.
Argument działa także wtedy, gdy któryś
ze zbiorów jest pusty.
\end{proof}

\begin{twierdzenie}[Prawo potęgowania potęgi]\label{thm:card-exponent-exponent}
Dla dowolnych zbiorów $A,B,C$ zachodzi
\[
\bigl(|C|^{|B|}\bigr)^{|A|}
=|C|^{|A|\cdot|B|}.
\]
Dotyczy to również przypadków, w których
co najmniej jeden z tych zbiorów jest pusty.
\end{twierdzenie}
\begin{proof}
Lewa strona oznacza moc zbioru funkcji
$A\to C^B$, a prawa --- moc zbioru funkcji $A\times B\to C$.
Niech $\Phi:(C^B)^A\to C^{A\times B}$
będzie określona przez
$\Phi(u)(a,b)=u(a)(b)$.
Jeśli $\Phi(u)=\Phi(v)$, to dla dowolnych
$a\in A$ oraz $b\in B$ mamy $u(a)(b)=v(a)(b)$.
Stąd $u(a)=v(a)$ jako funkcje o dziedzinie $B$,
a następnie $u=v$: jest to iniekcja.
Dla każdej funkcji $w:A\times B\to C$ połóżmy
$u_w(a)(b)=w(a,b)$. Wtedy
$u_w:A\to C^B$ i $\Phi(u_w)=w$, a więc $\Phi$ jest surjekcją.
Również dla pustych dziedzin rozumowanie pozostaje poprawne,
bo równość funkcji wyznaczają ich wartości na wszystkich
argumentach i istnieje dokładnie jedna funkcja z pustej
dziedziny w dowolny zbiór.
\end{proof}

\begin{przyklad}[Pusty zbiór w potęgowaniu mocy]
\label{ex:card-zero-exponents}
Niech $0=|\varnothing|$ i $1=|\{\varnothing\}|$.
Istnieje dokładnie jedna funkcja
$\varnothing\to\varnothing$ --- funkcja pusta,
więc $0^0=1$. Nie istnieje funkcja z singletonu
do zbioru pustego, zatem $0^1=0$.
Dla dowolnego zbioru $A$ istnieje
dokładnie jedna funkcja $\varnothing\to A$,
więc $|A|^0=1$.
Jeśli $A=\{a,b\}$, $B=\{0,1,2\}$,
to $|B^A|=3^2=9$, gdyż każda funkcja
jest jednoznacznie określona przez dwa
niezależnie wybrane obrazy $a$ i $b$.
W tym skończonym obliczeniu nie używamy AC.
\end{przyklad}

\begin{twierdzenie}[Twierdzenie Cantora w języku mocy]\label{thm:card-cantor}
Dla każdego zbioru $A$ zachodzi $|A|<2^{|A|}$, gdzie $2=|\{0,1\}|$.
\end{twierdzenie}
\begin{proof}
Każdemu $S\subseteq A$ odpowiada funkcja charakterystyczna $\chi_S:A\to\{0,1\}$; przyporządkowanie jest bijekcją $\mathcal P(A)\to\{0,1\}^A$. Twierdzenie~\ref{thm:zf-cantor} daje iniekcję $A\to\mathcal P(A)$ i brak surjekcji w drugą stronę. Gdyby zbiory były równoliczne, istniałaby taka surjekcja. Zatem nierówność jest ścisła.
\end{proof}

\section{Przeliczalność i moc kontinuum}
\begin{definicja}[Zbiory przeliczalne]\label{def:card-countable}
Zbiór $A$ nazywamy przeliczalnym, jeśli jest skończony albo istnieje bijekcja $\N\to A$. Zbiór nieskończony i przeliczalny nazywamy przeliczalnie nieskończonym. Moc $\N$ oznaczamy $\aleph_0$.
\end{definicja}
\begin{definicja}[Moc kontinuum]\label{def:card-continuum}
\emph{Mocą kontinuum} nazywamy moc zbioru liczb rzeczywistych. Oznaczamy ją
przez $\mathfrak c$, czyli z definicji
\[
  \mathfrak c:=|\R|.
\]
Symbol $\mathfrak c$ oznacza \emph{moc zbioru}, a nie sam zbiór $\R$.
Analogicznie symbol $\aleph_0:=|\N|$ oznacza moc zbioru liczb naturalnych.
Sama definicja $\mathfrak c$ nie rozstrzyga jeszcze, czy jest to ta sama
moc co $2^{\aleph_0}$; tę równość udowodnimy poniżej.
\end{definicja}
\begin{twierdzenie}[Przeliczalność par liczb naturalnych]\label{thm:card-n2}
Istnieje bijekcja $\N\times\N\to\N$. W konsekwencji suma dwóch zbiorów przeliczalnych oraz iloczyn dwóch zbiorów przeliczalnych są przeliczalne.
\end{twierdzenie}
\begin{proof}
\emph{Krok 1: liczby trójkątne.}
Dla każdego $s\in\N$ określmy
$t_s:=\sum_{k\in s}(k+1)$, gdzie $t_0=0$
(jest to pusta suma, bo $0=\varnothing$).
Zbiór $s$ zawiera dokładnie liczby $0,\ldots,s-1$
w konstrukcji von Neumanna.
Z definicji $t_{s+1}=t_s+s+1$.
Indukcja po $s$ dowodzi $2t_s=s(s+1)$:
dla $s=0$ jest to $0=0$, a z założenia dla $s$
wynika
\[
2t_{s+1}=2t_s+2(s+1)
=s(s+1)+2(s+1)=(s+1)(s+2).
\]
Zatem $t_s=s(s+1)/2$ jest liczbą naturalną,
$t_{s+1}-t_s=s+1>0$ i $t_s\ge s$,
więc ciąg $(t_s)$ jest ściśle rosnący
i nieograniczony.

\emph{Krok 2: parowanie.}
Dla $m,n\in\N$ połóżmy $s=m+n$ i
\[
\pi(m,n):=t_s+n=\frac{s(s+1)}2+n.
\]
Przy ustalonym $s$ zachodzi $n\in\{0,\ldots,s\}$,
a $m=s-n$. Otrzymujemy dokładnie wszystkie
liczby naturalne z przedziału
$\{t_s,\ldots,t_s+s\}=\{t_s,\ldots,t_{s+1}-1\}$.
Te przedziały są rozłączne, następują po sobie
bez luk i pokrywają $\N$, ponieważ $t_0=0$
i $t_s$ jest nieograniczony.
Każda liczba $k$ określa więc jednoznacznie
$s$ oraz $n=k-t_s$ i $m=s-n$.
To dowodzi bijektywności $\pi$.

\emph{Krok 3: dowolne dwa zbiory przeliczalne.}
Jeśli $A$ jest przeliczalny, istnieje iniekcja
$j_A:A\to\N$: dla zbioru skończonego wynika to
z jego numerowania, a dla przeliczalnie
nieskończonego --- z odwrotności bijekcji $\N\to A$.
Dla dwóch ustalonych zbiorów $A,B$ wybieramy
ich dwie iniekcje $j_A,j_B$, co nie wymaga AC.
Wtedy $(a,b)\mapsto\pi(j_A(a),j_B(b))$
jest iniekcją $A\times B\to\N$.
Suma rozłączna $A\sqcup B$ również zanurza się w
$\N\times\{0,1\}$, a następnie w $\N$.

Pozostaje sprawdzić, że każdy podzbiór $D\subseteq\N$
jest skończony albo przeliczalnie nieskończony.
Jeżeli $D$ jest nieskończony, wybieramy kolejno
\emph{najmniejsze} niewypisane elementy,
które istnieją przez dobry porządek $\omega$
(Twierdzenie~\ref{thm:zf-wellordered-omega}):
$e(0)=\min D$ oraz
\[
e(k+1):=\min\bigl(D\setminus
\{e(0),\ldots,e(k)\}\bigr).
\]
Zbiór w nawiasie jest niepusty, bo z $D$
usuwamy tylko skończenie wiele elementów.
Każde minimum jest jednoznaczne,
więc stosujemy rekurencję na $\omega$
z Twierdzenia~\ref{thm:zf-omega-recursion}
do zbioru stanów $\mathcal P(D)$.
Zaczynamy od pustego zbioru wybranych elementów;
w kroku dołączamy najmniejszy element $D$
spoza stanu. Dla formalnej totalności funkcji
przejścia można pozostawić stan niezmieniony,
gdy wyczerpałby całe $D$; na rozważanej
orbicie to nie zachodzi, bo $D$ jest nieskończony.
Z kolejnych stanów jednoznacznie odczytujemy
nowe elementy i otrzymujemy funkcję
$e:\N\to D$ bez użycia wyboru.
Funkcja jest ściśle rosnąca i iniektywna.
Jeśli $d\in D$ nie byłby wartością $e$,
wszystkie wybrane elementy musiałyby być
mniejsze od $d$ (inaczej w pierwszym kroku
przekraczającym $d$ wybralibyśmy samo $d$).
Nie istnieje jednak nieskończenie wiele
różnych liczb naturalnych mniejszych od $d$.
Zatem $e$ jest surjekcją i $D$ jest przeliczalny.
To zamyka oba wnioski.
\end{proof}

\begin{przyklad}[Ręczne odwrócenie parowania Cantora]
\label{ex:card-unpair}
Dla $(m,n)=(2,1)$ mamy $s=m+n=3$ i
$\pi(2,1)=3\cdot4/2+1=7$.
Odwrotnie, dla $k=7$ szukamy jedynego $s\in\N$ z
\[
\frac{s(s+1)}2\le k<\frac{(s+1)(s+2)}2.
\]
Ponieważ $6\le7<10$, mamy $s=3$.
Stąd $n=7-6=1$ i $m=s-n=2$.
Przedziały dla kolejnych $s$ są rozłączne,
więc wyznaczona para jest jednoznaczna.
\end{przyklad}

\begin{figure}[htbp]
\centering
\begin{tikzpicture}[x=1.17cm,y=.72cm,font=\small]
  \foreach \m in {0,1,2,3}{
    \foreach \n in {0,1,2,3}{
      \node[inner sep=2pt] at (\m,-\n) {$(\m,\n)$};
    }
  }
  \foreach \m/\n in {0/3,1/2,2/1,3/0}{
    \node[rounded corners=2pt,fill=blue!12,draw=blue!45!black,
          inner sep=3pt] at (\m,-\n) {$(\m,\n)$};
  }
  \node[align=left,anchor=west] at (3.7,-1.5)
      {$m+n=3$\\$\pi(2,1)=7$};
\end{tikzpicture}
\caption{Fragment siatki $\N^2$: wyróżnione pary
z przekątnej $m+n=3$.}
\label{fig:card-pairing-diagonal}
\end{figure}

\begin{uwaga}[Gdzie pojawia się aksjomat wyboru?]
Dowód dla dwóch ustalonych zbiorów nie wymaga AC: dwie egzystencjalne przesłanki pozwalają wybrać dwa świadki bez postulatu jednoczesnego wyboru z nieskończonej rodziny. Natomiast zdanie „przeliczalna suma przeliczalnych zbiorów jest przeliczalna” może wymagać słabszej zasady wyboru: aby z przeliczalnej rodziny zbiorów przeliczalnych wybrać jednocześnie ich enumeracje. Wrócimy do tego w rozdziale o słabszych zasadach wyboru.
\end{uwaga}
\begin{twierdzenie}[Moc zbioru liczb rzeczywistych]\label{thm:card-continuum}
Istnieje bijekcja $\R\to\mathcal P(\N)$; zatem $\mathfrak c=2^{\aleph_0}>\aleph_0$.
\end{twierdzenie}
\begin{proof}
Najpierw skonstruujmy iniekcję $\mathcal P(\N)\to[0,1]$ przez
\[
S\longmapsto \sum_{n\in S}\frac{2}{3^{n+1}}.
\]
Dla dwóch różnych zbiorów niech $k$ będzie najmniejszym indeksem różnicy. Różnica wkładów na pozycji $k$ ma moduł $2/3^{k+1}$, a suma modułów wszystkich późniejszych wkładów wynosi $1/3^{k+1}$. Wartości są więc różne.

W drugą stronę każdej liczbie rzeczywistej $x$
przypiszmy jej przekrój
$D_x:=\{q\in\Q:q<x\}$.
Zbiór ten istnieje przez separację.
Jeśli $x<y$, to z gęstości $\Q$ w $\R$
istnieje $q\in\Q$ z $x<q<y$,
a więc $q\in D_y\setminus D_x$.
Stąd przyporządkowanie $x\mapsto D_x$
jest iniekcją $\R\to\mathcal P(\Q)$.

Zbudujmy teraz konkretną iniekcję $j:\Q\to\N$.
Każda liczba wymierna ma dokładnie jedno
przedstawienie $q=a/b$, gdzie $a\in\Z$,
$b\in\N\setminus\{0\}$ i $\gcd(|a|,b)=1$.
Zdefiniujmy iniekcję $\gamma:\Z\to\N$
przez $\gamma(a)=2a$ dla $a\ge0$
i $\gamma(a)=-2a-1$ dla $a<0$.
Następnie połóżmy
$j(q):=\pi(\gamma(a),b-1)$,
gdzie $\pi$ jest bijekcją
z Twierdzenia~\ref{thm:card-n2}.
Jednoznaczność zredukowanego ułamka,
iniektywność $\gamma$ i $\pi$
dowodzą iniektywności $j$.
Przyporządkowanie $S\mapsto j[S]$
jest zatem iniekcją
$\mathcal P(\Q)\to\mathcal P(\N)$,
bez potrzeby wyboru dowolnej enumeracji.

Złożyliśmy iniekcję
$\R\to\mathcal P(\Q)\to\mathcal P(\N)$.
W przeciwną stronę skonstruowaliśmy
iniekcję $\mathcal P(\N)\to[0,1]\subseteq\R$.
Twierdzenie~\ref{thm:card-cb} daje bijekcję
$\R\to\mathcal P(\N)$, zgodnie z tezą.
\end{proof}
\begin{uwaga}[Czego nie rozstrzyga twierdzenie Cantora?]
Wiemy, że $\aleph_0<\mathfrak c=2^{\aleph_0}$. Pytanie, czy istnieje moc ściśle pośrednia między $\aleph_0$ a $\mathfrak c$, to hipoteza continuum. Sam dowód diagonalny nie rozstrzyga tej hipotezy.
\end{uwaga}

\section{Dlaczego $\N$ i $\R$ nie są równoliczne?}\label{sec:card-n-vs-r}
\begin{twierdzenie}[Nieprzeliczalność prostej rzeczywistej]\label{thm:card-n-not-r}
Nie istnieje surjekcja $e:\N\to\R$.
W szczególności nie istnieje bijekcja $\N\to\R$, a ponieważ
funkcja $n\mapsto n$ jest iniekcją $\N\to\R$, zachodzi
\[
|\N|<|\R|.
\]
\end{twierdzenie}
\begin{proof}
\emph{Krok 1: zakodowanie podzbiorów liczb naturalnych liczbami rzeczywistymi.}
Dla $S\subseteq\N$ połóżmy
\[
 c(S):=\sum_{n=0}^{\infty}\frac{2\chi_S(n)}{3^{n+1}}\in[0,1],
\]
gdzie $\chi_S(n)=1$ dla $n\in S$ i $\chi_S(n)=0$ w przeciwnym przypadku.
Szereg jest zbieżny, ponieważ jego wyrazy są nieujemne i
$\sum_{n=0}^{\infty}2/3^{n+1}=1$.

Pokażemy, że $c$ jest iniekcją. Niech $S\ne T$ i niech $k$
będzie najmniejszą liczbą naturalną należącą do dokładnie
jednego z tych zbiorów. Przyjmijmy dla ustalenia uwagi, że
$k\in S\setminus T$. Wówczas składniki szeregu dla $n<k$ są identyczne,
a składnik dla $n=k$ jest w $c(S)$ większy o $2/3^{k+1}$.
Cała możliwa różnica późniejszych składników ma moduł
nie większy niż
\[
\sum_{n=k+1}^{\infty}\frac{2}{3^{n+1}}
=\frac{1}{3^{k+1}}.
\]
Stąd $c(S)-c(T)\ge 1/3^{k+1}>0$, więc $c(S)\ne c(T)$.

\emph{Krok 2: sprzeczność z domniemaną enumeracją liczb rzeczywistych.}
Załóżmy, że istnieje surjekcja $e:\N\to\R$.
Dla każdego $n\in\N$ określmy zbiór $S_n$ następująco:
jeśli $e(n)$ jest postaci $c(S)$ dla jakiegoś $S\subseteq\N$,
to $S_n$ oznacza ten \emph{jedyny} zbiór $S$;
w przeciwnym razie kładziemy $S_n=\varnothing$.
Jednoznaczność w pierwszym przypadku wynika z kroku pierwszego.
Nie dokonujemy tu żadnego nieuzasadnionego wyboru z wielu możliwości.

\emph{Krok 3: diagonalizacja.}
Zbudujmy zbiór
\[
D:=\{n\in\N:n\notin S_n\}.
\]
Jest to podzbiór $\N$. Liczba $c(D)$ jest rzeczywista,
a $e$ jest — z założenia — surjekcją, więc dla pewnego $k\in\N$
zachodzi $e(k)=c(D)$. Wtedy definicja $S_k$ daje $S_k=D$.
Sprawdzamy teraz, czy $k$ należy do $D$:
\[
 k\in D\quad\Longleftrightarrow\quad k\notin S_k
 \quad\Longleftrightarrow\quad k\notin D.
\]
Otrzymaliśmy sprzeczność. Dlatego surjekcja $\N\to\R$
nie istnieje, a zatem $\N$ i $\R$ nie są równoliczne.
\end{proof}
\begin{uwaga}[Dlaczego użyliśmy cyfr $0$ i $2$ w bazie trójkowej?]
Zwykłe rozwinięcia dziesiętne mają niejednoznaczności,
na przykład $0{,}999\ldots=1$. W powyższej konstrukcji
nie musimy zakładać jednoznaczności ogólnego rozwinięcia
trójkowego: oszacowanie pierwszego różniącego się składnika
samodzielnie dowodzi, że różnym zbiorom $S$ odpowiadają różne liczby.
To właśnie sprawdzenie zapobiega pominięciu istotnej trudności.
\end{uwaga}

\section{Jawne bijekcje między przedziałami a prostą}\label{sec:card-interval-bijections}
\begin{twierdzenie}[Otwarty przedział i prosta są równoliczne]\label{thm:card-open-interval-real}
Istnieje bijekcja $F:\R\to(0,1)$. Jedną z nich jest
\[
F(x):=\frac12\left(1+\frac{x}{1+|x|}\right).
\]
\end{twierdzenie}
\begin{proof}
\emph{Krok 1: sprawdzenie wartości.}
Dla każdego $x\in\R$ zachodzi $|x|<1+|x|$,
a więc $-1<x/(1+|x|)<1$. Dodając $1$ i dzieląc przez $2$,
otrzymujemy $0<F(x)<1$. Wzór rzeczywiście określa funkcję
$F:\R\to(0,1)$.

\emph{Krok 2: wyprowadzenie funkcji odwrotnej.}
Niech $t\in(0,1)$, a $u:=2t-1$. Wtedy $-1<u<1$, zatem
$1-|u|>0$. Połóżmy
\[
G(t):=\frac{2t-1}{1-|2t-1|}=\frac{u}{1-|u|}.
\]
Wzór jest dobrze określony dla każdego $t\in(0,1)$.
Wykażemy dwie równości, które oznaczają, że $F$ i $G$
są funkcjami wzajemnie odwrotnymi.

\emph{Krok 3: sprawdzenie $F(G(t))=t$.}
Dla $x=G(t)=u/(1-|u|)$ mamy
\[
1+|x|=1+\frac{|u|}{1-|u|}=\frac{1}{1-|u|},
\qquad \frac{x}{1+|x|}=u.
\]
Stąd $F(G(t))=(1+u)/2=t$.

\emph{Krok 4: sprawdzenie $G(F(x))=x$.}
Dla dowolnego $x\in\R$ połóżmy $u=x/(1+|x|)$.
Wtedy $2F(x)-1=u$ i
$1-|u|=1/(1+|x|)$. Zatem
\[
G(F(x))=\frac{u}{1-|u|}
=\frac{x/(1+|x|)}{1/(1+|x|)}=x.
\]
Dwie tożsamości dowodzą bijektywności $F$.
\end{proof}
\begin{przyklad}[Jak działa ta bijekcja?]
Dla $x=0$ otrzymujemy $F(0)=1/2$; dla $x=1$ mamy
$F(1)=3/4$, a dla $x=-1$ mamy $F(-1)=1/4$.
Gdy $x$ rośnie bez ograniczenia, wartości $F(x)$
zbliżają się do $1$ od dołu; gdy $x$ maleje bez ograniczenia,
zbliżają się do $0$ od góry. Żaden argument nie daje
wartości $0$ ani $1$, zgodnie z przeciwdziedziną $(0,1)$.
\end{przyklad}

\begin{twierdzenie}[Przedział otwarty i półotwarty są równoliczne]\label{thm:card-open-half-open}
Istnieje bijekcja $H:(0,1)\to(0,1]$, mimo że $(0,1)$
jest właściwym podzbiorem $(0,1]$.
\end{twierdzenie}
\begin{proof}
\emph{Krok 1: wskazanie punktów do przesunięcia.}
Rozważmy ciąg różnych elementów przedziału $(0,1)$:
\[
\frac12,\ \frac13,\ \frac14,\ \frac15,\ \ldots
\]
Chcemy przesunąć każdą z tych liczb na poprzednie miejsce,
a pierwszą z nich przenieść do punktu $1$, którego brakuje
w dziedzinie.

\emph{Krok 2: ścisła definicja funkcji.}
Dla $x\in(0,1)$ określmy
\[
H(x):=
\begin{cases}
\dfrac{1}{n-1},&\text{gdy }x=\dfrac1n\text{ dla }n\in\N,\ n\ge2,\\[3pt]
x,&\text{gdy }x\notin\{\dfrac1n:n\in\N,\ n\ge2\}.
\end{cases}
\]
Ponieważ liczby $1/n$ są parami różne, pierwszy przypadek
nigdy nie przypisuje temu samemu argumentowi różnych wartości.
Pierwszy przypadek daje element $(0,1]$ (dla $n=2$ jest to $1$),
a drugi daje element $(0,1)$. Funkcja jest dobrze określona.

\emph{Krok 3: zbudowanie odwrotności.}
Zdefiniujmy $K:(0,1]\to(0,1)$ przez
\[
K(y):=
\begin{cases}
\dfrac12,&\text{gdy }y=1,\\[3pt]
\dfrac{1}{n+1},&\text{gdy }y=\dfrac1n\text{ dla }n\in\N,\ n\ge2,\\[3pt]
y,&\text{gdy }y\notin\{1\}\cup\{\dfrac1n:n\in\N,\ n\ge2\}.
\end{cases}
\]
Trzy przypadki są rozłączne i obejmują wszystkie $y\in(0,1]$.
Dla $x=1/2$ mamy $K(H(x))=K(1)=1/2$.
Dla $x=1/n$, $n\ge3$, dostajemy
$K(H(1/n))=K(1/(n-1))=1/n$.
Jeśli $x$ nie jest żadnym $1/n$, obie funkcje pozostawiają go bez zmian.
Zatem $K\circ H=\operatorname{id}_{(0,1)}$.

W drugą stronę: $H(K(1))=H(1/2)=1$;
dla $y=1/n$, $n\ge2$, zachodzi
$H(K(y))=H(1/(n+1))=1/n=y$;
pozostałe wartości również pozostają bez zmian.
Stąd $H\circ K=\operatorname{id}_{(0,1]}$.
Otrzymaliśmy funkcje wzajemnie odwrotne, więc $H$ jest bijekcją.
\end{proof}
\begin{przyklad}[Co dzieje się z konkretnymi punktami?]
$H(1/2)=1$, $H(1/3)=1/2$, $H(1/4)=1/3$,
natomiast $H(2/3)=2/3$.
Wszystkie punkty poza przeliczalnym ciągiem $1/2,1/3,\ldots$
pozostają na swoich miejscach. Jednocześnie żaden punkt nie ginie:
na miejsce $1/2$ trafia $1/3$, na miejsce $1/3$ trafia $1/4$
i tak dalej. Na tym polega „zrobienie miejsca” dla nowego punktu $1$.
\end{przyklad}

\begin{wniosek}[Trzy zbiory tej samej mocy]\label{cor:card-interval-real}
Zachodzi $(0,1)\approx(0,1]\approx\R$.
W szczególności $|(0,1)|=|(0,1]|=|\R|$.
\end{wniosek}
\begin{proof}
Twierdzenie~\ref{thm:card-open-interval-real} daje bijekcję
$F:\R\to(0,1)$, a Twierdzenie~\ref{thm:card-open-half-open}
daje bijekcję $H:(0,1)\to(0,1]$.
Złożenie $H\circ F:\R\to(0,1]$ jest bijekcją,
ponieważ ma funkcję odwrotną $F^{-1}\circ H^{-1}$.
Jest to wniosek z definicji równoliczności, bez stosowania
jakiegokolwiek dodatkowego aksjomatu wyboru.
\end{proof}
\begin{uwaga}[Nie mylmy liczności z długością przedziału]
Zbiory $(0,1)$, $(0,1]$ i $\R$ mają tę samą \emph{moc},
ale nie tę samą długość. Równoliczność wymaga bijekcji,
nie zaś zachowania odległości, długości czy ciągłości.
Funkcja $H$ nie jest ciągła w punktach $1/n$, lecz
ciągłość nie występuje w definicji równoliczności.
\end{uwaga}

\section{Rozróżnienie wyboru i porównywalności mocy}
\begin{twierdzenie}[Porównywalność mocy a aksjomat wyboru]
\label{thm:card-comparability-choice}
W teorii ZF następujące warunki są równoważne:
\begin{enumerate}[label=(\roman*)]
\item obowiązuje aksjomat wyboru;
\item dla dowolnych zbiorów $A,B$ istnieje
iniekcja $A\to B$ albo iniekcja $B\to A$.
\end{enumerate}
\end{twierdzenie}
\begin{proof}
Załóżmy AC. Z Twierdzenia~\ref{thm:choice-wellorder}
każdy z dwóch ustalonych zbiorów $A,B$
ma dobry porządek. Lemat o porównywaniu dobrych
porządków (Lemat~\ref{lem:choice-wellorder-comparison})
daje izomorfizm jednego z nich z całym drugim
lub z jego właściwym odcinkiem początkowym.
Złożenie z włączeniem odcinka początkowego
do całego zbioru daje wymaganą iniekcję.
Dwa ustalone dobre porządki można wskazać
bez dodatkowej zasady wyboru.

Odwrotnie, załóżmy porównywalność wszystkich
zbiorów i ustalmy dowolny zbiór $A$.
Twierdzenie Hartogsa~\ref{thm:choice-hartogs}
daje dobrze uporządkowany zbiór $H(A)$,
dla którego nie istnieje iniekcja $H(A)\to A$.
Z porównywalności wynika więc iniekcja
$j:A\to H(A)$. Przenieśmy porządek z $H(A)$
na $A$ przez
$a\prec b$ wtedy i tylko wtedy, gdy
$j(a)<_{H(A)}j(b)$.
Dla każdego niepustego $X\subseteq A$
jego obraz $j[X]$ ma element najmniejszy;
ponieważ $j$ jest iniekcją, element ten ma
dokładnie jeden przeciwobraz w $X$, który
jest najmniejszym elementem $X$.
Zatem $\prec$ jest dobrym porządkiem na $A$.
Ponieważ $A$ był dowolny, obowiązuje zasada
dobrego uporządkowania, a z
Twierdzenia~\ref{thm:choice-equivalence}
wynika AC.
\end{proof}
\begin{uwaga}[Unikanie błędnego koła]
\label{rem:card-no-circular-choice}
Używany tutaj zbiór $H(A)$ skonstruowano w ZF
w rozdziale~\ref{ch:axiom-choice}, bez
liczb porządkowych i bez aksjomatu wyboru.
Dopiero w rozdziale~\ref{ch:ordinals}
zdefiniujemy ordynały i opiszemy je
jako kanoniczne reprezentacje dobrych porządków.
Samo twierdzenie Cantora--Bernsteina
nie zapewnia porównywalności wszystkich mocy.
\end{uwaga}

\chapter{Liczby porządkowe i rekurencja pozaskończona}\label{ch:ordinals}
\chaptermark{Liczby porządkowe i rekurencja pozaskończona}

\section{Od dobrych porządków do obiektów kanonicznych}
Dwa dobrze uporządkowane zbiory mogą być izomorficzne, mimo że mają różne elementy. Chcemy reprezentować ich typy porządkowe konkretnymi zbiorami. W tym rozdziale pracujemy w ZF, o ile nie wskazano inaczej; przy konstrukcji liczb kardynalnych jako najmniejszych ordynałów danej mocy korzystamy z AC.
\begin{definicja}[Zbiór przechodni i liczba porządkowa]\label{def:ordinal}
Zbiór $X$ jest \emph{przechodni}, jeżeli $y\in x\in X$ implikuje $y\in X$. \emph{Liczbą porządkową} (ordynałem) nazywamy zbiór przechodni $\alpha$, na którym relacja $\in$ jest ścisłym dobrym porządkiem: jest przechodnia, dowolne dwa różne elementy są porównywalne i każdy niepusty podzbiór ma element najmniejszy względem $\in$.
\end{definicja}
\begin{przyklad}[Pierwsze ordynały]
$0=\varnothing$, $1=\{0\}$, $2=\{0,1\}$, $3=\{0,1,2\}$ są ordynałami. Również $\omega$, skonstruowane w Twierdzeniu~\ref{thm:zf-omega},
jest ordynałem: jego elementy to skończone ordynały von Neumanna.
Lemat~\ref{lem:zf-vn-successor} dowodzi przechodniości $\omega$,
a Twierdzenie~\ref{thm:zf-wellordered-omega} wraz
z Lematem~\ref{lem:zf-initial-segments} pokazuje,
że relacja $\in$ dobrze porządkuje $\omega$. Ordynał $\omega+1=\omega\cup\{\omega\}$ ma element największy, podczas gdy $\omega$ go nie ma. Dlatego te dwa porządki nie są izomorficzne, choć ich zbiory są równoliczne.
\end{przyklad}
\begin{lemat}[Elementy ordynału i odcinki początkowe]\label{lem:ordinal-segments}
Jeżeli $\alpha$ jest ordynałem i $\beta\in\alpha$, to $\beta$ jest ordynałem oraz $\beta=\{\gamma\in\alpha:\gamma\in\beta\}$. Każdy właściwy odcinek początkowy $\alpha$ ma postać $\beta$ dla dokładnie jednego $\beta\in\alpha$.
\end{lemat}
\begin{proof}
Ponieważ $\alpha$ jest przechodni, $\beta\subseteq\alpha$. Dla $\gamma\in\beta$ i $\delta\in\gamma$ przechodniość $\in$ na $\alpha$ daje $\delta\in\beta$; zatem $\beta$ jest przechodni. Ograniczenie dobrego porządku $\in$ z $\alpha$ do $\beta$ pozostaje dobrym porządkiem, więc $\beta$ jest ordynałem. Równość zbiorów wynika z $\beta\subseteq\alpha$. Jeśli $I$ jest właściwym odcinkiem początkowym, zbiór $\alpha\setminus I$ ma najmniejszy element $\beta$; wtedy $I=\{\gamma\in\alpha:\gamma\in\beta\}=\beta$. Jednoznaczność wynika z antysymetrii porządku liniowego.
\end{proof}
\begin{twierdzenie}[Porównanie ordynałów]\label{thm:ordinal-trichotomy}
Dla dowolnych ordynałów $\alpha,\beta$ zachodzi dokładnie jedna z relacji $\alpha\in\beta$, $\alpha=\beta$, $\beta\in\alpha$.
\end{twierdzenie}
\begin{proof}
Lemat~\ref{lem:choice-wellorder-comparison} porównuje dobre porządki $(\alpha,\in)$ i $(\beta,\in)$. Pokażemy najpierw, że izomorfizm $f$ ordynału $\alpha$ na odcinek początkowy ordynału $\beta$ musi być identycznością. Gdyby nie, weźmy najmniejsze $\gamma\in\alpha$ z $f(\gamma)\ne\gamma$. Dla każdego $\delta\in\gamma$ mamy $f(\delta)=\delta$, a izomorfizm przenosi zbiór poprzedników $\gamma$ na zbiór poprzedników $f(\gamma)$. Z Lematu~\ref{lem:ordinal-segments} otrzymujemy $\gamma=f(\gamma)$, sprzeczność. Izomorficzne ordynały są więc równe, a izomorfizm z właściwym odcinkiem daje odpowiednio $\alpha\in\beta$ lub $\beta\in\alpha$. Rozłączność wynika z asymetrii $\in$ na ordynałach i aksjomatu ufundowania.
\end{proof}
\begin{wniosek}[Następnik i supremum ordynałów]\label{cor:ordinal-sup}
Dla ordynału $\alpha$ zbiór $\alpha+1:=\alpha\cup\{\alpha\}$ jest ordynałem i jest najmniejszym ordynałem większym od $\alpha$. Dla każdego zbioru $S$ ordynałów jego suma $\sup S:=\bigcup S$ jest ordynałem i najmniejszym wspólnym ograniczeniem górnym w porządku $\le$, gdzie $\alpha\le\beta$ oznacza $\alpha\subseteq\beta$.
\end{wniosek}
\begin{proof}
\emph{Krok 1: zagnieżdżenie i suma.}
Lemat~\ref{lem:ordinal-segments}
i Twierdzenie~\ref{thm:ordinal-trichotomy}
dają dla dowolnych ordynałów
$\alpha,\beta$ zawieranie
$\alpha\subseteq\beta$ lub
$\beta\subseteq\alpha$.
Połóżmy $\gamma:=\bigcup S$.
Jeśli $x\in y\in\gamma$, to
dla pewnego $\delta\in S$
mamy $y\in\delta$; stąd
$x\in\delta\subseteq\gamma$.
Zatem $\gamma$ jest przechodni.
Dla różnych $x,y\in\gamma$
istnieją $\delta,\varepsilon\in S$
z $x\in\delta$, $y\in\varepsilon$.
Zagnieżdżenie tych dwóch ordynałów
daje wspólny ordynał zawierający
$x,y$, gdzie porównuje je relacja $\in$.
Relacja $\in$ na $\gamma$
jest więc ścisłym porządkiem liniowym.

\emph{Krok 2: dobry porządek.}
Jeśli $X\subseteq\gamma$ jest niepusty,
ustalmy $x\in X$ i $\delta\in S$
z $x\in\delta$. Niepusty zbiór
$X\cap\delta$ ma najmniejszy
element $m$, gdyż $\delta$ jest
dobrze uporządkowany.
Jeżeli $y\in X$ i $y\in m$,
to $m\in\delta$ i przechodniość
$\delta$ dają $y\in X\cap\delta$,
wbrew minimalności $m$.
Zatem $m$ jest najmniejszym
elementem całego $X$.
Dla $S=\varnothing$ mamy
$\gamma=\varnothing$, również ordynał.

\emph{Krok 3: supremum.}
Z definicji sumy zbiorów
każde $\delta\in S$ zawiera się
w $\gamma$, a każdy ordynał
zawierający wszystkie $\delta\in S$
zawiera także $\gamma$.
To dowodzi $\gamma=\sup S$.

\emph{Krok 4: następnik.}
Zbiór $\alpha+1:=\alpha\cup\{\alpha\}$
jest przechodni: jeśli
$y\in x\in\alpha+1$, to
albo $x\in\alpha$ i $y\in\alpha$
przez przechodniość $\alpha$,
albo $x=\alpha$ i także
$y\in\alpha$.
Relacja $\in$ porządkuje liniowo
$\alpha+1$: elementy $\alpha$
są porównywalne i wszystkie
poprzedzają $\alpha$.
Niepusty podzbiór przecinający
$\alpha$ ma minimum w $\alpha$,
a jedyny niepusty podzbiór
nieprzecinający $\alpha$ to
$\{\alpha\}$.
Zatem $\alpha+1$ jest ordynałem.
Jeśli $\alpha<\eta$ dla ordynału
$\eta$, to $\alpha\in\eta$,
a z przechodniości $\eta$
wynika $\alpha+1\subseteq\eta$.
Nie ma więc ordynału
ściśle pomiędzy $\alpha$ i $\alpha+1$.
\end{proof}
\begin{definicja}[Ordynały graniczne]\label{def:ordinal-limit}
Ordynał $\lambda$ jest \emph{graniczny}, jeśli $\lambda\ne0$ i nie jest następnikiem żadnego ordynału. Wtedy $\lambda=\sup\{\beta:\beta<\lambda\}$. Przykładem jest $\omega$, podczas gdy $\omega+1$ jest następnikiem.
\end{definicja}


\begin{lemat}[Ordynały graniczne jako sumy poprzedników]
\label{lem:ordinal-limit-union}
Jeżeli $\lambda$ jest niezerowym ordynałem granicznym, to
$\lambda=\bigcup_{\beta<\lambda}\beta$.
Jeżeli $\lambda=\gamma+1$ jest następnikiem, to
$\bigcup_{\beta<\lambda}\beta=\gamma$.
\end{lemat}
\begin{proof}
Każdy $\beta\in\lambda$ jest podzbiorem $\lambda$
na mocy Lematu~\ref{lem:ordinal-segments},
zatem $U:=\bigcup_{\beta<\lambda}\beta\subseteq\lambda$.
Załóżmy, że $\lambda$ jest graniczny i
istnieje $x\in\lambda\setminus U$.
Gdyby istniał $y\in\lambda$ z $x\in y$,
to $x\in y\subseteq U$, sprzeczność.
Zatem $x$ jest największym elementem
dobrze uporządkowanego $\lambda$
(właściwa liniowość relacji $\in$ wynika
z definicji ordynału).
Każdy $y\in\lambda$ jest więc równy $x$
lub należy do $x$, skąd
$\lambda=x\cup\{x\}=x+1$,
wbrew założeniu graniczności.
Stąd $\lambda=U$.

Jeżeli $\lambda=\gamma+1$, każdy $\beta\in\lambda$
należy do $\gamma$ lub jest równy $\gamma$,
a więc $\beta\subseteq\gamma$.
Z drugiej strony $\gamma\in\lambda$,
więc $U=\gamma$.
\end{proof}
\begin{przyklad}[Pusta suma i pierwszy ordynał graniczny]
\label{ex:ordinal-suprema}
Dla pustego zbioru ordynałów
$\sup\varnothing=\bigcup\varnothing=0$.
Dla $S=\{1,3\}$ otrzymujemy $\sup S=3$,
ponieważ $1\subseteq3$ i $1\cup3=3$.
Ponadto $\sup\{n:n<\omega\}=\omega$:
dla każdego $m\in\omega$ istnieje
następnik $m+1\in\omega$ zawierający $m$.
Natomiast
$\sup\{\beta:\beta<\omega+1\}=\omega$,
a nie $\omega+1$, zgodnie z lematem.
\end{przyklad}

\section{Indukcja pozaskończona}
\begin{twierdzenie}[Zasada indukcji pozaskończonej]\label{thm:ordinal-induction}
Niech $\alpha$ będzie ordynałem, a $P(\beta)$
formułą języka ZF z ustalonymi parametrami
i zmienną wolną $\beta$, rozważaną dla $\beta<\alpha$. Jeśli dla każdego $\beta<\alpha$ z założenia $P(\gamma)$ dla wszystkich $\gamma<\beta$ wynika $P(\beta)$, to $P(\beta)$ zachodzi dla wszystkich $\beta<\alpha$.
\end{twierdzenie}
\begin{proof}
Gdyby istniał kontrprzykład, zbiór
$E=\{\beta\in\alpha:\neg P(\beta)\}$
istniałby przez separację i byłby niepusty. Niech $\beta$ będzie jego najmniejszym elementem. Wszystkie $\gamma<\beta$ spełniają $P$, a założenie kroku daje $P(\beta)$, sprzeczność.
\end{proof}
\begin{uwaga}[Trzy przypadki w praktyce]
Wygodnie rozdziela się dowód na przypadek $0$, następnik $\beta+1$ i ordynał graniczny $\lambda$. Nie są to trzy niezależne aksjomaty: stanowią sposób sprawdzenia jednego warunku indukcji z poprzedniego twierdzenia.
\end{uwaga}

\section{Rekurencja pozaskończona}
\begin{twierdzenie}[Rekurencja na ustalonym ordynale]\label{thm:ordinal-recursion}
Niech $\alpha$ będzie ordynałem, a $F$ operacją klasową zdefiniowaną formułą teorii mnogości z ustalonymi parametrami, która każdej funkcji $g$ o dziedzinie będącej ordynałem $\beta<\alpha$ przypisuje dokładnie jeden zbiór $F(g)$. Wówczas istnieje dokładnie jedna funkcja $h$ o dziedzinie $\alpha$, dla której
\[
\forall\beta<\alpha\quad h(\beta)=F(h|_\beta).
\]
W dowodzie korzystamy ze schematu zastępowania; nie potrzeba aksjomatu wyboru.
\end{twierdzenie}
\begin{proof}
\emph{Krok 1: przybliżenia i ich zgodność.}
Poprawnym przybliżeniem długości
$\beta\le\alpha$ nazywamy funkcję
$g$ o dziedzinie $\beta$ taką, że
$g(\gamma)=F(g|_\gamma)$ dla każdego $\gamma<\beta$.
Jest to własność wyrażalna formułą
języka ZF z ustalonymi parametrami.
Jeśli dwa przybliżenia $g,k$
różniłyby się na wspólnej dziedzinie,
istniałby najmniejszy punkt $\gamma$ różnicy.
Wtedy $g|_\gamma=k|_\gamma$, więc
\[
g(\gamma)=F(g|_\gamma)=F(k|_\gamma)=k(\gamma),
\]
sprzeczność.
Dla danej dziedziny istnieje co najwyżej
jedno przybliżenie.

\emph{Krok 2: indukcyjne istnienie.}
Stosujemy indukcję pozaskończoną
do ordynału $\alpha+1$, aby uwzględnić
wszystkie $\beta\le\alpha$.
Dla $\beta=0$ funkcja pusta jest przybliżeniem.
Jeśli $\beta=\delta+1\le\alpha$,
to z hipotezy indukcyjnej mamy
jednoznaczne przybliżenie $g$
o dziedzinie $\delta$.
Ponieważ $\delta<\alpha$,
operacja $F(g)$ jest określona
i jest zbiorem.
Funkcja $g\cup\{(\delta,F(g))\}$
jest poprawnym przybliżeniem długości $\beta$.

Niech $\beta\le\alpha$ będzie
niezerowym ordynałem granicznym.
Dla każdego $\delta<\beta$
istnieje dokładnie jedno
przybliżenie $g_\delta$ długości $\delta$.
Ponieważ jest to przyporządkowanie
jednoznacznie definiowalne formułą,
schemat zastępowania daje zbiór
$\mathcal G=\{g_\delta:\delta<\beta\}$.
Poprzedni krok pokazuje zgodność
przybliżeń, więc $g:=\bigcup\mathcal G$
jest funkcją.
Na mocy Lematu~\ref{lem:ordinal-limit-union}
jej dziedzina jest równa
$\bigcup_{\delta<\beta}\delta=\beta$.
Dla $\gamma<\beta$ mamy
$\gamma+1<\beta$ (gdyż $\beta$ jest graniczny),
zatem $g|_{\gamma+1}=g_{\gamma+1}$.
Poprawność tego przybliżenia
daje $g(\gamma)=F(g|_\gamma)$.
Wobec tego $g$ jest poprawne
na całej dziedzinie $\beta$.

\emph{Krok 3: zakończenie.}
Indukcja daje poprawne przybliżenie
$h$ długości $\alpha$.
Spełnia ono wymagane równanie,
a zgodność dowolnych poprawnych
przybliżeń daje jego jednoznaczność.
Ani w etapie następnikowym, ani granicznym
nie dokonuje się nieskończonego
wyboru spośród wielu dopuszczalnych wartości.
\end{proof}
\begin{przyklad}[Dodawanie ordynałów]\label{ex:ordinal-addition}
Dla ustalonego ordynału $\alpha$ definiujemy
\emph{sumę porządkową} rekurencyjnie
po drugim argumencie:
\[
\alpha+0:=\alpha,\qquad
\alpha+(\beta+1):=(\alpha+\beta)+1,
\]
a dla niezerowego ordynału granicznego $\lambda$:
\[
\alpha+\lambda:=\sup_{\beta<\lambda}(\alpha+\beta).
\]
Symbol $+1$ po prawej oznacza następnik ordynału
z Wniosku~\ref{cor:ordinal-sup}.
Są to trzy przypadki jednoznacznej operacji
na dotychczas wyliczonych wartościach.
Dla dowolnego ustalonego ordynału $\theta$
Twierdzenie~\ref{thm:ordinal-recursion}
konstruuje funkcję $\beta\mapsto\alpha+\beta$
na $\beta<\theta$:
operację klasową na \emph{wszystkich}
funkcjach o dziedzinie $\beta$
można zdefiniować tą samą regułą,
a w przypadku niewłaściwych wartości
zwracać umownie $0$.
Indukcja zapewnia, że na konstruowanej funkcji
wszystkie wartości są ordynałami.
Jeśli $\theta<\theta'$, obie konstrukcje
zgadzają się na wspólnej dziedzinie
przez jednoznaczność rekurencji.

Dla każdego $n\in\omega$ indukcja
daje $1+n=n+1$:
dla $n=0$ mamy $1+0=1$,
a z $1+n=n+1$ otrzymujemy
$1+(n+1)=(n+1)+1$.
Wobec tego
\[
1+\omega=\sup_{n<\omega}(1+n)
=\sup_{n<\omega}(n+1)=\omega.
\]
Natomiast $\omega+1>\omega$
z definicji następnika.
Dodawanie ordynałów nie jest więc przemienne.
\end{przyklad}


\begin{przyklad}[Hierarchia von Neumanna]
\label{ex:ordinal-vn-hierarchy}
Niech $\alpha$ będzie ordynałem z $4\le\alpha$.
Definiujemy rekurencyjnie $V_\beta$ dla $\beta<\alpha$:
$V_0:=\varnothing$,
$V_{\delta+1}:=\mathcal P(V_\delta)$
wtedy, gdy $\delta+1<\alpha$,
oraz $V_\lambda:=\bigcup_{\beta<\lambda}V_\beta$
dla niezerowych ordynałów granicznych $\lambda<\alpha$.
Warunki na indeksy gwarantują, że
każdy symbol po lewej stronie oznacza
poziom należący do rozważanej rodziny.
Twierdzenie~\ref{thm:ordinal-recursion}
gwarantuje istnienie funkcji
$\beta\mapsto V_\beta$ na ustalonym zbiorze
indeksów $\alpha$.
W każdym kroku otrzymujemy jednoznacznie zbiór,
więc nie używamy aksjomatu wyboru.
W szczególności
$V_1=\{\varnothing\}$,
$V_2=\{\varnothing,\{\varnothing\}\}$,
a $V_3=\mathcal P(V_2)$ ma cztery elementy.
W ZF nie zakładamy istnienia jednego zbioru
zawierającego poziomy $V_\beta$ dla
\emph{wszystkich} ordynałów $\beta$.
\end{przyklad}

\section{Mnożenie ordynałów i różnica względem arytmetyki mocy}

\begin{definicja}[Iloczyn porządkowy ordynałów]
\label{def:ordinal-multiplication}
Dla dowolnych ordynałów $\alpha,\beta$
iloczyn porządkowy $\alpha\cdot\beta$
definiujemy rekurencyjnie po $\beta$:
\[
\alpha\cdot0:=0,\qquad
\alpha\cdot(\beta+1):=(\alpha\cdot\beta)+\alpha,
\]
a dla niezerowego ordynału granicznego $\lambda$:
\[
\alpha\cdot\lambda:=\sup_{\beta<\lambda}
(\alpha\cdot\beta).
\]
Wzór zawiera dodawanie \emph{ordynałów}
z Przykładu~\ref{ex:ordinal-addition},
a nie dodawanie mocy.
Na każdym ustalonym odcinku ordynałów
istnienie i jednoznaczność funkcji
$\beta\mapsto\alpha\cdot\beta$
wynikają z Twierdzenia~\ref{thm:ordinal-recursion}.
Aby zastosować twierdzenie dokładnie w jego postaci,
dla niepoprawnych przybliżeń (np. z wartościami
niebędącymi ordynałami) definiujemy wartość
pomocniczej operacji klasowej jako $0$.
Na właściwych przybliżeniach rekurencja
stosuje wypisane powyżej równania
i indukcja zachowuje przynależność do ordynałów.
Dzięki jednoznaczności wynik nie zależy
od wybranego odcinka indeksów.
\end{definicja}
\begin{przyklad}[Nieprzemienność mnożenia ordynałów]
\label{ex:ordinal-mul-noncommutative}
Dla każdego $n\in\omega$ rekurencja
daje $2\cdot n=2n$ w arytmetyce
skończonych ordynałów, zatem
$2\cdot\omega=\sup_{n<\omega}2n=\omega$.
Natomiast
$\omega\cdot2=(\omega\cdot1)+\omega
=\omega+\omega>\omega$.
Istotnie, $\omega+1$ należy do rodziny
$\{\omega+n:n<\omega\}$,
której supremum wynosi $\omega+\omega$,
więc $\omega+\omega\ge\omega+1>\omega$.
Odpowiada to ustawieniu drugiej kopii
porządku $\omega$ po wszystkich punktach pierwszej.
Wobec tego $2\cdot\omega\ne\omega\cdot2$,
chociaż oba ordynały mają moc $\aleph_0$.
W arytmetyce mocy
$2\cdot\aleph_0=\aleph_0\cdot2=\aleph_0$
(Twierdzenie~\ref{thm:card-n2}).
\end{przyklad}

\section{Liczby kardynalne jako ordynały początkowe}
\begin{definicja}[Ordynał początkowy]\label{def:initial-ordinal}
Ordynał $\kappa$ jest \emph{początkowy}, jeżeli żaden ordynał $\beta<\kappa$ nie jest z nim równoliczny. W ZFC \emph{liczbą kardynalną} zbioru $A$ nazywamy najmniejszy ordynał równoliczny z $A$ i oznaczamy $|A|$.
\end{definicja}
\begin{twierdzenie}[Istnienie reprezentanta kardynalnego]\label{thm:initial-exists}
Przy założeniu ZFC każdy zbiór $A$ jest równoliczny z dokładnie jednym ordynałem początkowym.
\end{twierdzenie}
\begin{proof}
\emph{Krok 1: dobry porządek na $A$.}
Na mocy Twierdzenia~\ref{thm:choice-wellorder}
istnieje dobry porządek $\prec$ na $A$.
Jest to jedyne miejsce w konstrukcji,
w którym wykorzystujemy AC.
Dalszy argument działa dla dowolnego ustalonego
zbioru dobrze uporządkowanego $(W,\prec)$ w ZF.

\emph{Krok 2: rekursja wzdłuż ustalonego dobrego porządku.}
Dla $w\in W$ oznaczmy
$W_{<w}:=\{v\in W:v\prec w\}$ oraz
$W_{\le w}:=W_{<w}\cup\{w\}$.
Chcemy skonstruować funkcję $r$ na $W$
spełniającą
\[
r(w)=\{r(v):v\prec w\}.
\tag{$*$}
\]
Poprawnym $w$-przybliżeniem nazwijmy funkcję
$f$ o dziedzinie $W_{\le w}$,
która dla każdego $t\preceq w$
(tj. $t\prec w$ albo $t=w$) spełnia
$f(t)=\{f(v):v\prec t\}$.
Nie zakładamy istnienia wspólnego z góry
zbioru $V$ dla wszystkich przybliżeń:
warunek jest formułą ZF o wykresie funkcji.

Dwa poprawne przybliżenia są zgodne
na wspólnej dziedzinie.
Gdyby się różniły, zbiór punktów niezgodności
miałby najmniejszy element $t$.
Wartości przy wszystkich $v\prec t$
byłyby równe, toteż równanie ($*$)
narzucałoby tę samą wartość w $t$,
co przeczy wyborowi $t$.
Na ustalonej dziedzinie istnieje więc
co najwyżej jedno poprawne przybliżenie.

Pokażemy indukcją po dobrym porządku $W$,
że dla każdego $w\in W$ istnieje
poprawne $w$-przybliżenie $f_w$.
Załóżmy, że istnieje ono dla wszystkich $v\prec w$.
Dzięki jednoznaczności schemat zastępowania
daje zbiór
$\mathcal F_w:=\{f_v:v\prec w\}$.
Suma $u:=\bigcup\mathcal F_w$ jest funkcją:
wszystkie składniki są zgodne.
Jej dziedziną jest dokładnie $W_{<w}$,
ponieważ każdy $t\prec w$ należy do
dziedziny $f_t$, a żadna dziedzina
$f_v$ nie zawiera elementów spoza $W_{<w}$.
Dla każdego $t\prec w$ funkcja $u$
zgadza się z $f_t$ na $W_{\le t}$,
zatem spełnia równanie ($*$) w punkcie $t$.
Kładziemy
\[
f_w:=u\cup\{(w,u[W_{<w}])\}.
\]
Nowy argument $w$ nie należał do
dziedziny $u$, więc jest to funkcja
na $W_{\le w}$ spełniająca ($*$).
Dla najmniejszego $w$ mamy
$\mathcal F_w=\varnothing$ i $u=\varnothing$,
a więc konstrukcja obejmuje przypadek początkowy.

Indukcja jest zasadna bez AC:
jeżeli istniałby kontrprzykład w $W$,
miałby element najmniejszy, dla którego
wszystkie wcześniejsze przybliżenia
już istnieją, co daje sprzeczność.
Zastępowanie zbiera rodzinę
$\{f_w:w\in W\}$ w zbiór.
Jej suma $r:=\bigcup_{w\in W}f_w$
jest funkcją o dziedzinie $W$
spełniającą ($*$).
Gdy $W=\varnothing$, bierzemy funkcję pustą.

\emph{Krok 3: obraz funkcji jest ordynałem.}
Dla $u\prec v$ zachodzi $r(u)\in r(v)$
na mocy ($*$). Jeśli $r(u)=r(v)$
dla różnych $u,v\in W$, to jeden
z nich poprzedza drugi i otrzymujemy
$r(u)\in r(u)$, wbrew aksjomatowi ufundowania.
Zatem $r$ jest iniekcją.
W drugą stronę, jeśli $r(u)\in r(v)$,
to na mocy ($*$) istnieje $t\prec v$
z $r(u)=r(t)$. Iniektywność daje
$u=t\prec v$.
Wykazaliśmy więc
\[
u\prec v\quad\Longleftrightarrow\quad
r(u)\in r(v).
\tag{$**$}
\]

Dla $w\in W$ zbiór
$r(w)=\{r(v):v\prec w\}$ jest przechodni:
jeżeli $x\in r(v)\in r(w)$,
to dla pewnego $t\prec v\prec w$
mamy $x=r(t)\in r(w)$.
Na mocy ($**$) relacja $\in$
na $r(w)$ jest dobrym porządkiem,
przeniesionym z $W_{<w}$ przez iniekcję $r$.
Zatem każde $r(w)$ jest ordynałem.
Zbiór $\gamma:=r[W]$ istnieje
przez zastępowanie, jest przechodni
z równania ($*$) i jest dobrze
uporządkowany przez $\in$
na mocy ($**$).
Jest więc ordynałem, a
$r:W\to\gamma$ izomorfizmem
porządkowym.
Jednoznaczność $\gamma$ wynika
z Twierdzenia~\ref{thm:ordinal-trichotomy}:
jedyny izomorfizm dwóch ordynałów
jest identycznością.

\emph{Krok 4: najmniejszy reprezentant mocy.}
Wracając do $W=A$, mamy ordynał
$\alpha$ równoliczny z $A$.
Zbiór
\[
E:=\{\beta\in\alpha+1:\beta\approx A\}
\]
istnieje przez separację i jest niepusty,
ponieważ $\alpha\in E$.
Dobry porządek na $\alpha+1$
daje najmniejszy element $\kappa$ zbioru $E$.
Jest on początkowy: gdyby
$\beta<\kappa$ był równoliczny z $\kappa$,
byłby również równoliczny z $A$,
a zatem $\beta\in E$, wbrew
minimalności $\kappa$.

Jeżeli $\kappa'$ byłby innym ordynałem
początkowym równolicznym z $A$,
to Twierdzenie~\ref{thm:ordinal-trichotomy}
dawałoby $\kappa<\kappa'$ lub $\kappa'<\kappa$.
Mniejszy byłby równoliczny z większym,
co przeczyłoby początkowości większego.
Stąd reprezentant kardynalny
istnieje i jest jednoznaczny.
\end{proof}
\begin{uwaga}[Dwie różne miary nieskończoności]
Ordynał mierzy \emph{typ dobrego porządku}, liczba kardynalna — \emph{liczbę elementów}. Zbiory $\omega$ i $\omega+1$ są równoliczne:
bijekcję $\omega\to\omega+1$ wyznaczają
wzory $f(0)=\omega$ oraz $f(n+1)=n$.
Nie są jednak izomorficzne jako dobre porządki,
gdyż $\omega$ nie ma elementu największego,
a $\omega+1$ ma.
Najmniejszą nieskończoną liczbą kardynalną
jest $\omega$, oznaczany również $\aleph_0$:
każdy ordynał mniejszy od $\omega$
jest skończony, a $\omega$ jest nieskończony,
bo każdy jego element ma większy następnik.
\end{uwaga}

\section{Dlaczego zbiór wszystkich ordynałów nie istnieje?}
\begin{twierdzenie}[Paradoks Burali-Forti]\label{thm:burali-forti}
Nie istnieje zbiór $O$, którego elementami są dokładnie wszystkie ordynały.
\end{twierdzenie}
\begin{proof}
Przypuśćmy, że taki zbiór $O$ istnieje. Suma $\gamma=\bigcup O$ jest ordynałem na mocy Wniosku~\ref{cor:ordinal-sup}. Z definicji $O$ wynika $\gamma\in O$, zatem $\gamma\subseteq\bigcup O=\gamma$, co jeszcze nie daje sprzeczności. Ale $\gamma+1$ również jest ordynałem, więc $\gamma+1\in O$ i $\gamma+1\subseteq\gamma$. Tymczasem $\gamma\in\gamma+1$, skąd $\gamma\in\gamma$, wbrew aksjomatowi ufundowania. Zatem $O$ nie jest zbiorem.
\end{proof}
\begin{uwaga}
Wyrażenie „klasa wszystkich ordynałów” oznacza tu własność definiowalną formułą języka ZF, a nie dodatkowy obiekt, po którym wolno kwantyfikować w języku pierwszego rzędu.
\end{uwaga}
\chapter{Słabsze zasady wyboru i ich ograniczenia}\label{ch:weak-choice}
\chaptermark{Słabsze zasady wyboru}
\section{Dlaczego rozróżniamy postacie wyboru?}
Aksjomat wyboru (AC) dotyczy dowolnej rodziny niepustych zbiorów. W wielu dowodach potrzebujemy jednak wyboru tylko z rodziny przeliczalnej albo kolejnych wyborów zależnych od poprzednich. W tym rozdziale podstawową teorią jest ZF; wszystkie dodatkowe założenia podajemy jawnie.
\begin{definicja}[Wybór przeliczalny]\label{def:ac-omega}
Zasada $\mathrm{AC}_{\omega}$ głosi, że dla każdej funkcji $n\mapsto A_n$ o dziedzinie $\omega$, której wartości są niepustymi zbiorami, istnieje funkcja $f$ o dziedzinie $\omega$ taka, że $f(n)\in A_n$ dla każdego $n\in\omega$.
\end{definicja}
\begin{definicja}[Wybór zależny]\label{def:dc}
Zasada wyboru zależnego $\mathrm{DC}$ głosi: dla każdego niepustego zbioru $X$, każdej relacji $R\subseteq X\times X$ spełniającej $\forall x\in X\ \exists y\in X\ (xRy)$ i każdego $x_0\in X$ istnieje funkcja $f:\omega\to X$ taka, że $f(0)=x_0$ i $f(n)Rf(n+1)$ dla wszystkich $n\in\omega$.
\end{definicja}
\begin{twierdzenie}[Łańcuch implikacji]\label{thm:choice-chain}
W ZF zachodzą implikacje $\mathrm{AC}\Rightarrow\mathrm{DC}\Rightarrow\mathrm{AC}_{\omega}$.
\end{twierdzenie}
\begin{proof}
Załóżmy AC. Dla każdego $x\in X$ niepusty jest zbiór $R[x]=\{y\in X:xRy\}$. AC daje funkcję $c:X\to X$ z $c(x)\in R[x]$. Zwykła rekurencja na $\omega$ z $f(0)=x_0$ i $f(n+1)=c(f(n))$ daje DC.

Załóżmy DC i rodzinę $(A_n)_{n\in\omega}$ niepustych zbiorów. Niech $X$ będzie zbiorem wszystkich skończonych funkcji $s$ o dziedzinie pewnego $n\in\omega$ spełniających $s(k)\in A_k$ dla $k<n$. Zbiór $X$ istnieje w ZF: wszystkie takie funkcje są podzbiorami $\omega\times\bigcup_{n\in\omega}A_n$, więc $X$ otrzymujemy przez wyróżnianie w zbiorze potęgowym. Nie zakładamy istnienia elementu dowolnego nieskończonego produktu. Jest niepusty, bo zawiera funkcję pustą. Niech $sRt$ oznacza, że $t$ przedłuża $s$ o dokładnie jedną parę $(n,a)$ z $a\in A_n$, gdzie $n=\dom s$. Każdy $s$ ma następnika, ponieważ $A_n\ne\varnothing$. Z DC istnieje ciąg $s_0,s_1,\ldots$ zaczynający się od funkcji pustej i spełniający $s_nRs_{n+1}$. Wówczas $f=\bigcup_{n\in\omega}s_n$ jest funkcją wyboru dla całej rodziny.
\end{proof}
\begin{uwaga}[Nie mylić dwóch rodzajów wyboru]
Z faktu, że każdy $A_n$ jest niepusty, wynika dla każdego ustalonego $n$ istnienie elementu $a_n$. Nie wynika z samej logiki pierwszego rzędu istnienie jednej funkcji $n\mapsto a_n$. Właśnie tę lukę wypełnia $\mathrm{AC}_{\omega}$.
\end{uwaga}
\section{Zastosowanie: przeliczalna suma zbiorów przeliczalnych}
\begin{twierdzenie}\label{thm:countable-union-cc}
W ZF+$\mathrm{AC}_{\omega}$ dla każdej rodziny $(A_n)_{n\in\omega}$ zbiorów co najwyżej przeliczalnych suma $\bigcup_n A_n$ jest co najwyżej przeliczalna.
\end{twierdzenie}
\begin{proof}
Dla każdego $n$ niech $E_n$ będzie zbiorem wszystkich iniekcji $A_n\to\omega$. Jest on niepusty z założenia. Z $\mathrm{AC}_{\omega}$ wybieramy $e_n\in E_n$ jednocześnie dla wszystkich $n$. Każdemu $a\in\bigcup_n A_n$ przypiszmy parę $(n,e_n(a))$, gdzie $n$ jest najmniejszym indeksem z $a\in A_n$. Jest to iniekcja do $\omega\times\omega$, który jest przeliczalny na mocy Twierdzenia~\ref{thm:card-n2}. Nie potrzeba wybierać elementów z pustych zbiorów.
\end{proof}
\section{Słabe zasady a zwykłe konstrukcje}
\begin{przyklad}[Skończenie wiele wyborów]
Jeżeli $A$ i $B$ są niepuste, to istnieją $a\in A$ i $b\in B$, a więc para $(a,b)\in A\times B$. Ten argument wykorzystuje dwukrotnie regułę wprowadzania kwantyfikatora egzystencjalnego, nie AC. Dla nieskończenie wielu zbiorów analogiczne rozumowanie nie tworzy automatycznie funkcji wyboru.
\end{przyklad}
\begin{uwaga}[Granice twierdzeń]
Implikacje odwrotne do łańcucha z Twierdzenia~\ref{thm:choice-chain} nie są twierdzeniami ZF (przy standardowych założeniach niesprzeczności). Wykazanie tego wymaga budowania modeli teorii mnogości; nie traktujemy samego braku znanego dowodu jako dowodu niezależności. Metody modelowe wprowadzimy w dalszych rozdziałach.
\end{uwaga}

\section{Różnica między wyborem a rekurencją}
\begin{przyklad}[Jawna reguła nie wymaga wyboru]
Niech $A_n=\{n,n+1\}$ dla $n\in\omega$. Funkcja $f(n)=n$ jest funkcją wyboru już w ZF: wskazaliśmy jeden wzór określający wszystkie wartości. Zasada $\mathrm{AC}_\omega$ staje się istotna dopiero wtedy, gdy nie ma z góry zadanej reguły wyboru.
\end{przyklad}
\begin{uwaga}[Wybór dla rodzin indeksowanych]
Zasada $\mathrm{AC}_\omega$ nie mówi, że dowolna rodzina zbiorów przeliczalnych ma funkcję wyboru; wymaga, aby indeksowanie rodziny było funkcją o dziedzinie $\omega$. Samo stwierdzenie, że każdy zbiór rodziny jest przeliczalny, nie dostarcza indeksowania całej rodziny.
\end{uwaga}


\section{Wybór zależny i twierdzenie Baire'a dla przestrzeni zupełnych}
\label{sec:weak-choice-baire}
W przestrzeni Baire'a $\omega^\omega$ dopuszczalne
przedłużenia skończonych ciągów można wybierać
według ustalonego porządku liczb naturalnych.
W ogólnej zupełnej przestrzeni metrycznej nie ma jednak
kanonicznego sposobu wskazania środka następnej kuli.
Poniższy dowód pokazuje, gdzie wystarcza $\mathrm{DC}$.

\begin{twierdzenie}[Twierdzenie Baire'a z DC]
\label{thm:weak-choice-baire-dc}
W teorii ZF+$\mathrm{DC}$ niech $(X,d)$ będzie
niepustą zupełną przestrzenią metryczną.
Jeśli $(D_n)_{n\in\omega}$ jest ciągiem
zbiorów otwartych i gęstych w $X$,
to $\bigcap_{n\in\omega}D_n$ jest gęsty w $X$.
W szczególności żaden niepusty zbiór otwarty
w $X$ nie jest pierwszej kategorii w $X$.
\end{twierdzenie}
\begin{proof}
\emph{Krok 1: zbiory przybliżeń.}
Ustalmy niepusty zbiór otwarty $U\subseteq X$.
Piszemy
$B(z,r):=\{x\in X:d(x,z)<r\}$
oraz
$\overline B(z,r):=\{x\in X:d(x,z)\le r\}$.
Przez $\overline B(z,r)$ rozumiemy tu
\emph{kulę domkniętą}, która w ogólnej
przestrzeni metrycznej nie musi być
domknięciem topologicznym $B(z,r)$.
Dodatnie promienie wybieramy ze zbioru
$\mathbb Q_{>0}$.

Niech $P$ będzie zbiorem skończonych ciągów
$s=((z_0,r_0),\ldots,(z_{m-1},r_{m-1}))$
elementów $X\times\mathbb Q_{>0}$ spełniających
$r_i<2^{-(i+2)}$ oraz
\[
\overline B(z_0,r_0)\subseteq U\cap D_0,\qquad
\overline B(z_{i+1},r_{i+1})
 \subseteq B(z_i,r_i)\cap D_{i+1}
\quad(i+1<m).
\]
Dopuszczamy także ciąg pusty.
Wszystkie skończone ciągi par są elementami
zbioru
$\bigcup_{m\in\omega}(X\times\mathbb Q_{>0})^m$,
istniejącego w ZF przez rekursję,
zastępowanie i sumę.
Zatem $P$ istnieje przez separację
i zawiera ciąg pusty.

\emph{Krok 2: każde przybliżenie można przedłużyć.}
Jeśli $s$ jest pusty,
otwarty zbiór $U\cap D_0$ jest niepusty
(gdyż $D_0$ jest gęsty).
Wybieramy jeden punkt $z_0$ tego zbioru.
Otwartość daje $\varepsilon>0$ z
$B(z_0,\varepsilon)\subseteq U\cap D_0$.
Wybieramy jedną wymierną liczbę
$0<r_0<\min(\varepsilon,1/4)$;
wtedy
$\overline B(z_0,r_0)\subseteq U\cap D_0$.

Jeśli $s$ ma długość $m>0$, zbiór
$B(z_{m-1},r_{m-1})\cap D_m$
jest niepusty i otwarty.
Jego niepustość wynika z gęstości $D_m$
i niepustości kuli otwartej.
Wybieramy jeden punkt $z_m$ i dodatni
promień wymierny
$r_m<2^{-(m+2)}$, tak mały, aby
\[
\overline B(z_m,r_m)
\subseteq B(z_{m-1},r_{m-1})\cap D_m.
\]
Taki promień istnieje dzięki otwartości
zbioru po prawej i gęstości $\mathbb Q$ w $\mathbb R$.
Każdy z powyższych pojedynczych wyborów
jest zwykłym użyciem kwantyfikatora egzystencjalnego.

\emph{Krok 3: zastosowanie DC.}
Dla $s,t\in P$ niech $sRt$ oznacza, że
$t$ powstaje z $s$ przez dołączenie dokładnie
jednej pary. Krok 2 daje
$\forall s\in P\;\exists t\in P\;(sRt)$.
Zasada $\mathrm{DC}$ zastosowana do $P,R$
i ciągu pustego daje funkcję
$k\mapsto s_k$ o dziedzinie $\omega$,
zaczynającą się od ciągu pustego,
taką że $s_kRs_{k+1}$.
Suma tych zgodnych przybliżeń wyznacza
jednoznaczny ciąg par
$(z_n,r_n)_{n\in\omega}$.

Dla $j>i$ z zagnieżdżenia kul mamy
$z_j\in B(z_i,r_i)$, a więc
$d(z_i,z_j)<r_i<2^{-(i+2)}$.
Zatem $(z_n)$ jest ciągiem Cauchy'ego.
Zupełność daje jego granicę $x\in X$.

\emph{Krok 4: przynależność granicy.}
Dla każdego $i$ i wszystkich $j>i$ zachodzi
$z_j\in\overline B(z_i,r_i)$.
Ponieważ ta kula jest domknięta,
$x\in\overline B(z_i,r_i)$.
Dla $i=0$ mamy zatem $x\in U\cap D_0$,
a dla $i>0$:
\[
x\in\overline B(z_i,r_i)
 \subseteq B(z_{i-1},r_{i-1})\cap D_i.
\]
Wynika stąd
$x\in U\cap\bigcap_{i\in\omega}D_i$.
Ponieważ niepusty otwarty $U$
był dowolny, przekrój zbiorów $D_i$
jest gęsty.

\emph{Krok 5: kategoria pierwsza.}
Jeśli $U$ byłby zawarty w przeliczalnej
sumie zbiorów nigdziegęstych $(F_n)$,
to $D_n:=X\setminus\overline{F_n}$
byłyby otwarte i gęste.
Z poprzedniej części otrzymalibyśmy punkt
\[
x\in U\cap\bigcap_nD_n
\subseteq U\setminus\bigcup_nF_n,
\]
sprzecznie z $U\subseteq\bigcup_nF_n$.
\end{proof}

\begin{przyklad}[Dlaczego wymagana jest zupełność?]
\label{ex:weak-choice-baire-noncomplete}
Przestrzeń $(\mathbb Q,|\cdot|)$ nie jest zupełna:
ciąg wymiernych przybliżeń do $\sqrt2$ jest
ciągiem Cauchy'ego bez granicy w $\mathbb Q$.
Ustalmy bijekcję $e:\omega\to\mathbb Q$,
której istnienie w ZF wynika z konstrukcji
w Twierdzeniu~\ref{thm:card-n2}
i z jawnego kodowania liczb wymiernych
w dowodzie Twierdzenia~\ref{thm:card-continuum}.
Połóżmy
$D_n:=\mathbb Q\setminus
\{e(0),\ldots,e(n)\}$.
Każdy $D_n$ jest otwarty w $\mathbb Q$,
gdyż usuwamy skończenie wiele punktów.
Jest też gęsty, ponieważ
każdy niepusty przedział wymierny zawiera
nieskończenie wiele liczb wymiernych.
Jednak $\bigcap_nD_n=\varnothing$,
ponieważ każdy element $\mathbb Q$
należy do obrazu $e$.
Zatem bez zupełności twierdzenie może zawieść.
\end{przyklad}

\begin{uwaga}[Dokładna rola zasad wyboru]
\label{rem:weak-choice-baire-dc}
W dowodzie Twierdzenia~\ref{thm:weak-choice-baire-dc}
kolejna dopuszczalna kula zależała od poprzedniej.
Jest to właśnie sytuacja opisywana przez DC,
a nie niezależny wybór z gotowej rodziny.
W przestrzeni $\omega^\omega$
dla odpowiednio kodowanych przybliżeń
można wybierać najmniejsze dopuszczalne
przedłużenia, co daje osobny argument w ZF.
Nie wyprowadzamy stąd nieudowodnionej
równoważności DC z twierdzeniem Baire'a
dla wszystkich klas przestrzeni.
\end{uwaga}

\section{Algebry Boole'a, filtry i twierdzenie o ideale pierwszym}
\label{sec:weak-choice-bpi}
W aksjomacie wyboru możemy wskazywać elementy z dowolnej rodziny
niepustych zbiorów. Istnieje jednak słabsza zasada wystarczająca
do wielu zastosowań topologicznych: pozwala \emph{przedłużyć już
spójny układ decyzji} dotyczących zbiorów do decyzji całkowitej.
Zaczniemy od wyjaśnienia wszystkich występujących tu pojęć.

\begin{definicja}[Algebra Boole'a i filtr]\label{def:bpi-boolean-filter}
\emph{Algebrą Boole'a} nazywamy niepusty
zbiór $B$ z dwuargumentowymi działaniami
$\wedge,\vee:B^2\to B$,
jednoargumentowym $\neg:B\to B$
oraz wyróżnionymi elementami $0_B,1_B\in B$.
Dla wszystkich $a,b,c\in B$ zakładamy
łączność i przemienność obu działań,
prawa pochłaniania
\[
a\wedge(a\vee b)=a,\qquad
a\vee(a\wedge b)=a,
\]
oba prawa rozdzielności
\[
a\wedge(b\vee c)=(a\wedge b)\vee(a\wedge c),
\quad
a\vee(b\wedge c)=(a\vee b)\wedge(a\vee c),
\]
oraz $a\wedge\neg a=0_B$
i $a\vee\neg a=1_B$.
W szczególności wynikają stąd
$a\wedge1_B=a$, $a\vee0_B=a$
i idempotentność obu działań.
Definiujemy $a\leq b$ wtedy i tylko wtedy, gdy $a\wedge b=a$.
W całej sekcji zakładamy $0_B\ne1_B$.
\emph{Filtr właściwy} $F\subseteq B$ spełnia:
$1_B\in F$, $0_B\notin F$,
$b,c\in F\implies b\wedge c\in F$ oraz
$b\in F,\ b\leq c\implies c\in F$.
Filtr $U$ jest \emph{ultrafiltrem}, jeżeli jest właściwy i
maksymalny wśród właściwych filtrów względem zawierania.
\emph{Ideał pierwszy} to zbiór $I\subseteq B$, którego
dopełnienie $B\setminus I$ jest ultrafiltrem.
\end{definicja}
\begin{przyklad}[Co oznaczają operacje?]
Dla dowolnego niepustego zbioru $S$
algebra $\mathcal P(S)$ ma
$X\wedge Y=X\cap Y$, $X\vee Y=X\cup Y$,
$\neg X=S\setminus X$, $0_B=\varnothing$, $1_B=S$.
Dla ustalonego $s\in S$ rodzina
$U_s=\{X\subseteq S:s\in X\}$ jest ultrafiltrem.
Filtr zbiorów współskończonych w $\N$ zawiera każdy zbiór
o skończonym dopełnieniu; jest właściwy, ale nie jest
ultrafiltrem, gdyż ani liczby parzyste, ani ich dopełnienie
nie są współskończone.
\end{przyklad}

\begin{lemat}[Kryterium ultrafiltru]\label{lem:bpi-ultrafilter-criterion}
Dla każdego właściwego filtru $U$ na algebrze Boole'a $B$
następujące warunki są równoważne:
(i) $U$ jest ultrafiltrem;
(ii) dla każdego $b\in B$ dokładnie jeden z elementów
$b,\neg b$ należy do $U$.
\end{lemat}
\begin{proof}
Najpierw oba elementy nie mogą należeć do $U$, gdyż
ich iloczyn to $0_B$. Załóżmy, że $U$ jest maksymalny i $b\notin U$.
Filtr generowany przez $U\cup\{b\}$ byłby
właściwym filtrem ściśle większym od $U$,
gdyby wszystkie przekroje $u\wedge b$ ($u\in U$)
były niezerowe. Maksymalność $U$ wyklucza tę
możliwość, więc istnieje $u\in U$ z
$u\wedge b=0_B$. Wtedy $u\leq\neg b$,
a domknięcie filtru ku górze daje $\neg b\in U$.
Dokładnie jeden spośród $b,\neg b$ należy do $U$.
Dla (ii)$\Rightarrow$(i) niech $V\supsetneq U$ będzie
filtrem i wybierzmy $b\in V\setminus U$. Wtedy
$\neg b\in U\subseteq V$, więc $0_B=b\wedge\neg b\in V$.
Zatem $V$ nie jest właściwy i $U$ jest maksymalny.
\end{proof}

\begin{definicja}[BPI i zasada ultrafiltru]\label{def:bpi-principles}
\emph{Twierdzenie o ideale pierwszym w algebrze Boole'a}
(ang. \emph{Boolean Prime Ideal Theorem}, w skrócie BPI)
głosi: w każdej nietrywialnej algebrze Boole'a każdy
właściwy filtr można przedłużyć do ultrafiltru.
\emph{Zasada ultrafiltru} (UF) głosi: dla każdego zbioru
$S$ każdy właściwy filtr podzbiorów $S$ można
przedłużyć do ultrafiltru na $\mathcal P(S)$.
Oba zdania formułujemy nad ZF, nie zakładając AC.
\end{definicja}

\begin{twierdzenie}[Równoważność BPI i UF w ZF]
\label{thm:bpi-ultrafilter-equivalence}
W teorii ZF zasada BPI jest równoważna zasadzie UF.
\end{twierdzenie}
\begin{proof}
\emph{BPI $\Rightarrow$ UF.} Jeśli $S=\varnothing$,
nie istnieje właściwy filtr na $\mathcal P(S)$,
więc twierdzenie jest puste.
Dla niepustego $S$ stosujemy BPI
do nietrywialnej algebry $\mathcal P(S)$. Otrzymane przedłużenie
właściwego filtru jest dokładnie ultrafiltrem podzbiorów.

\emph{UF $\Rightarrow$ BPI, etap 1: zwartość kostki.}
Niech $I$ będzie dowolnym zbiorem. Z UF pokażemy, że
$2^I=\{0,1\}^I$ z topologią produktu dyskretnej pary
jest zwarte. Załóżmy, że rodzina $\mathcal C$ domkniętych
podzbiorów $2^I$ ma własność przecięcia skończonego,
czyli każdy skończony przekrój jest niepusty.
Generuje ona właściwy filtr na $2^I$; UF przedłuża go
do ultrafiltru $U$. Dla każdego $i\in I$ dokładnie jeden
z dwóch cylindrów $\{x:x(i)=0\}$ i $\{x:x(i)=1\}$
należy do $U$. Definiujemy $z(i)$ jako tę jedyną wartość,
której cylinder należy do $U$. Zastępowanie tworzy jedną
funkcję $z:I\to\{0,1\}$, bez żadnych wyborów
z nieokreślonych rodzin. Każde podstawowe otoczenie $z$
to skończony przekrój cylindrów należących do $U$,
więc również należy do $U$. Jeśli $C\in\mathcal C$ i
$z\notin C$, otwarte $2^I\setminus C$ zawierałoby takie
otoczenie $z$, a więc i ono należałoby do $U$, choć
$C\in U$: sprzeczność. Stąd $z\in\bigcap\mathcal C$.
To jest kryterium zwartości przez rodzinę domkniętą.

\emph{Etap 2: zapis aksjomatów algebry w kostce.}
Ustalmy nietrywialną algebrę $B$ i właściwy filtr $F$.
Rozważmy kostkę $2^B$, której punktem jest dowolna
funkcja $h:B\to\{0,1\}$ (niekoniecznie homomorfizm).
Dla $b,c\in B$ narzućmy warunki
\[
 h(0_B)=0,\quad h(1_B)=1,\quad
 h(\neg b)=1-h(b),\quad
 h(b\wedge c)=h(b)h(c),\quad h(b)=1\ (b\in F).
\]
Każda równość zależy od skończenie wielu współrzędnych
kostki, więc wyznacza podzbiór domknięty (nawet
otwarto-domknięty). Pokażemy, że cała rodzina tych
warunków ma własność przecięcia skończonego.

\emph{Etap 3: spełnialność skończonej części.}
Ustalmy skończenie wiele warunków i weźmy skończoną
podalgebrę $D\subseteq B$ generowaną przez wszystkie
występujące w nich elementy. Jest skończona, gdyż jej
elementy są kombinacjami boolowskimi skończenie wielu
generatorów; każda taka kombinacja jest sumą pewnych
spośród skończenie wielu elementarnych koniunkcji.
Niech $f_1,\dots,f_m\in F$ będą elementami,
którym nakazano wartość $1$. Ich iloczyn
$b=f_1\wedge\cdots\wedge f_m$ jest niezerowy,
bo $F$ jest właściwy (dla $m=0$ bierzemy $b=1_B$).
W skończonej algebrze $D$ istnieje atom $a\leq b$:
wśród niezerowych elementów $D$ leżących poniżej $b$
wybieramy element minimalny.
Jeżeli $d\in D$, atomowość $a$ oznacza,
że $a\wedge d$ jest równe $0_B$ lub $a$;
w pierwszym przypadku $a\le\neg d$,
a w drugim $a\le d$.
Nie mogą zachodzić oba przypadki, gdyż wtedy
$a\le d\wedge\neg d=0_B$, sprzeczność.
Dla $d\in D$ połóżmy
$h_D(d)=1$ dokładnie wtedy, gdy $a\leq d$.
Z wykazanej alternatywy wynika
$h_D(\neg d)=1-h_D(d)$.
Ponieważ $a\le d\wedge e$
wtedy i tylko wtedy, gdy
$a\le d$ oraz $a\le e$, mamy
$h_D(d\wedge e)=h_D(d)h_D(e)$.
Ponadto $h_D(0_B)=0$, $h_D(1_B)=1$
i $h_D(f_j)=1$, bo $a\le b\le f_j$. Rozszerzamy $h_D$
do punktu $2^B$, przyjmując na współrzędnych poza $D$
dowolną ustaloną wartość, na przykład zero.
Ten punkt spełnia wszystkie \emph{rozważane skończenie}
warunki; nie zakładaliśmy, że rozszerzenie jest
homomorfizmem całej algebry.

\emph{Etap 4: przejście do wszystkich warunków.}
Zwartość $2^B$ daje funkcję $h$ spełniającą wszystkie
równości jednocześnie. Z pierwszych trzech rodzajów
warunków wynika, że $h:B\to\{0,1\}$ jest homomorfizmem
algebr Boole'a (suma wynika z dopełnień i iloczynu).
Zbiór $U=h^{-1}[\{1\}]$ jest właściwym filtrem;
ponieważ $h(\neg b)=1-h(b)$, Lemat
\ref{lem:bpi-ultrafilter-criterion} mówi, że jest
ultrafiltrem. Ostatnie warunki dają $F\subseteq U$.
To kończy dowód UF$\Rightarrow$BPI.
\end{proof}

\begin{przyklad}[Ultrafiltr niegłówny na $\N$]
\label{ex:bpi-free-ultrafilter}
Filtr wszystkich zbiorów współskończonych w $\N$
jest właściwy. BPI przedłuża go do ultrafiltru $U$.
Taki ultrafiltr nie jest postaci
$U_n=\{A\subseteq\N:n\in A\}$: dla każdego $n$
zbiór $\N\setminus\{n\}$ należy do początkowego filtru,
a zatem i do $U$, podczas gdy nie należy do $U_n$.
Dostaliśmy ultrafiltr niegłówny bez podawania
żadnego wzoru wymieniającego jego elementy.
\end{przyklad}

\section{Dwie wersje twierdzenia Tychonoffa}
\label{sec:weak-choice-tychonoff}
\begin{definicja}[Zwartość i topologia produktu]
\label{def:tychonoff-topology}
Przestrzeń topologiczna $X$ jest \emph{zwarta}, jeżeli
z każdego jej pokrycia zbiorami otwartymi można wybrać
skończone podpokrycie. \emph{Hausdorffa} oznacza, że
każde dwa różne punkty mają rozłączne otoczenia otwarte.
Dla rodziny $(X_i)_{i\in I}$ \emph{topologia produktu}
na $X=\prod_{i\in I}X_i$ ma bazę złożoną ze zbiorów
$\bigcap_{i\in J}\pi_i^{-1}[V_i]$, gdzie
$J\subseteq I$ jest skończony, $V_i\subseteq X_i$ jest
otwarty, a $\pi_i(x)=x(i)$ jest projekcją współrzędnej.
Nie zakładamy, że sam produkt jest niepusty.
\end{definicja}

\begin{twierdzenie}[Tychonoff dla zwartych przestrzeni Hausdorffa]
\label{thm:bpi-tychonoff-hausdorff}
W ZF zasada BPI jest równoważna następującemu zdaniu:
\emph{każdy produkt dowolnej rodziny zwartych przestrzeni
Hausdorffa jest zwarty}.
\end{twierdzenie}
\begin{proof}
\emph{BPI $\Rightarrow$ zwartość produktów.}
Na mocy Twierdzenia~\ref{thm:bpi-ultrafilter-equivalence}
możemy korzystać z UF. Niech
$X=\prod_{i\in I}X_i$; jeśli $X=\varnothing$, to jest
zwarty. W przeciwnym razie rozważmy dowolny ultrafiltr
$\mathcal U$ podzbiorów $X$. Jego obraz przez $\pi_i$
to ultrafiltr $\mathcal U_i$ na $X_i$:
$A\subseteq X_i$ należy do $\mathcal U_i$
wtedy i tylko wtedy, gdy $\pi_i^{-1}[A]\in\mathcal U$.
Zwartość $X_i$ zapewnia istnienie punktu skupienia
$\mathcal U_i$, a Hausdorffowość jego jednoznaczność.
Dokładniej, domknięcia wszystkich zbiorów należących
do $\mathcal U_i$ mają własność przecięcia skończonego;
ich przekrój jest niepusty ze zwartości. Jeśli $x_i$
należy do tego przekroju i $V$ jest jego otwartym
otoczeniem, to $V\in\mathcal U_i$: w przeciwnym razie
$X_i\setminus V\in\mathcal U_i$, a jego domknięcie
nie zawierałoby $x_i$. W przestrzeni Hausdorffa dwa
różne punkty nie mogą być granicami jednego filtru,
więc $x_i$ jest jedyny.

Dla każdego $i$ otrzymaliśmy tym samym \emph{jednoznacznie}
określony $x_i$; funkcja $i\mapsto x_i$ istnieje
przez zastępowanie, bez aksjomatu wyboru. Wektor
$x=(x_i)_{i\in I}$ należy do $X$.
Każde podstawowe otoczenie $x$ jest skończonym
przekrojem przeciwobrazów otoczeń $x_i$, a zatem
należy do $\mathcal U$. Ultrafilter zbiega do $x$.
Aby wyprowadzić zwartość, przypuśćmy, że istnieje
pokrycie otwarte $X$ bez skończonego podpokrycia.
Dopełnienia jego elementów tworzą rodzinę domkniętych
zbiorów z własnością przecięcia skończonego; z UF
rozszerzamy generowany przez nie filtr do ultrafiltru.
Otrzymany ultrafiltr zbiega do $x$, lecz jakiś element
pokrycia zawiera $x$ i tym samym należy do ultrafiltru,
razem ze swoim dopełnieniem. Sprzeczność.

\emph{Zwartość produktów $\Rightarrow$ BPI.}
Załóżmy, że każdy produkt zwartych przestrzeni
Hausdorffa jest zwarty. W szczególności dla każdego
zbioru $B$ kostka $2^B$ jest zwarta, gdyż każdy czynnik
$\{0,1\}$ z topologią dyskretną jest skończony,
a więc zwarty i Hausdorffa. Stosujemy dokładnie
konstrukcję warunków zamkniętych z etapów 2--4
dowodu Twierdzenia~\ref{thm:bpi-ultrafilter-equivalence}:
zwartość kostki daje homomorfizm $h:B\to\{0,1\}$
przyjmujący $1$ na zadanym filtrze. Jego przeciwobraz
$\{1\}$ jest żądanym ultrafiltrem.
\end{proof}

\begin{twierdzenie}[Pełne twierdzenie Tychonoffa i AC]
\label{thm:ac-full-tychonoff}
Nad ZF następujące dwa zdania są równoważne:
(i) aksjomat wyboru AC;
(ii) produkt dowolnej rodziny przestrzeni zwartych
(niekoniecznie Hausdorffa) jest zwarty.
\end{twierdzenie}
\begin{proof}
\emph{AC $\Rightarrow$ (ii).}
AC daje lemat Kuratowskiego--Zorna
(Twierdzenie~\ref{thm:choice-equivalence}), a z niego
BPI: zbiór właściwych filtrów zawierających zadany filtr,
uporządkowany przez zawieranie, ma niepusty łańcuchowy
warunek górnego ograniczenia, bo suma niepustego łańcucha
właściwych filtrów jest właściwym filtrem, a pusty
łańcuch jest ograniczony przez zadany filtr. Zatem każdy
filtr podzbiorów można przedłużyć do ultrafiltru.
Niech $X=\prod_i X_i$ będzie produktem zwartych
przestrzeni. Dla ultrafiltru $\mathcal U$ na $X$ jego
obrazy $\mathcal U_i$ mają w każdym czynniku co najmniej
jeden punkt graniczny. AC pozwala \emph{jednocześnie}
wybrać punkty $x_i$ graniczne dla wszystkich $i$.
Wtedy $x=(x_i)_i$ jest granicą $\mathcal U$ w produkcie,
bo każde jego podstawowe otoczenie zależy od skończenie
wielu współrzędnych. Argument z pokryciem i UF z
poprzedniego dowodu zapewnia zwartość $X$.

\emph{(ii) $\Rightarrow$ AC.}
Niech $(A_i)_{i\in I}$ będzie dowolną rodziną
niepustych zbiorów. Dla każdego $i$ dołączmy jeden
nowy, nazwany punkt $\infty_i\notin A_i$ i połóżmy
$X_i=A_i\sqcup\{\infty_i\}$. Nadajmy $X_i$
topologię $\tau_i=\{\varnothing,A_i,\{\infty_i\},X_i\}$.
Każdy $X_i$ jest zwarty: ma tylko cztery zbiory otwarte;
a ponadto $A_i$ jest domknięto-otwarty. Produkt
$X=\prod_iX_i$ jest niepusty dzięki funkcji
$i\mapsto\infty_i$, nawet jeśli nie mamy wyboru
z rodziny $(A_i)$. Założenie (ii) daje zwartość $X$.
Dla każdego $i$ zbiór
$C_i=\{x\in X:x(i)\in A_i\}$ jest domknięty.
Każdy skończony przekrój zbiorów $C_i$ jest niepusty:
dla skończenie wielu indeksów wybieramy po jednym elemencie
z odpowiednich $A_i$, a na pozostałych współrzędnych
stawiamy $\infty_i$. Zwartość daje
$x\in\bigcap_{i\in I}C_i$. Wtedy $x(i)\in A_i$
dla każdego $i$, czyli $x$ jest funkcją wyboru.
\end{proof}

\begin{uwaga}[Dlaczego Hausdorffowość ma znaczenie?]
Dowód przy BPI wykorzystał \emph{jedyność} granicy
ultrafiltru w przestrzeni Hausdorffa, dzięki czemu
zebraliśmy wszystkie współrzędne w jedną funkcję
bez AC. W dowodzie dla ogólnych przestrzeni zwartych
granic może być wiele; wybranie jednej dla każdego
czynnika może wymagać pełnego AC.
Przestrzenie $X_i$ z ostatniego dowodu są dla
$|A_i|>1$ niehausdorffowskie: dowolne dwie różne
wartości z $A_i$ mają te same nietrywialne otoczenia.
\end{uwaga}

\section{Lemat Teichmüllera--Tukeya i dokładna siła wyboru}
\label{sec:weak-choice-tukey}
\begin{definicja}[Rodzina skończonego charakteru]
\label{def:tukey-finite-character}
Rodzina $\mathcal F\subseteq\mathcal P(E)$ ma
\emph{skończony charakter}, jeżeli dla każdego
$X\subseteq E$ zachodzi równoważność
\[
X\in\mathcal F\quad\Longleftrightarrow\quad
\text{każdy skończony podzbiór }Y\subseteq X\text{ należy do }\mathcal F.
\]
\emph{Lemat Teichmüllera--Tukeya} orzeka, że każda
niepusta rodzina skończonego charakteru ma element
maksymalny względem zawierania.
\end{definicja}

\begin{twierdzenie}[Równoważność lematu Tukeya i AC]
\label{thm:tukey-equivalent-ac}
W teorii ZF lemat Teichmüllera--Tukeya jest
równoważny pełnemu aksjomatowi wyboru.
\end{twierdzenie}
\begin{proof}
\emph{AC $\Rightarrow$ Tukey.}
Niech $\mathcal F\subseteq\mathcal P(E)$ będzie niepustą
rodziną skończonego charakteru, uporządkowaną przez
zawieranie. Dla każdego łańcucha $\mathcal C\subseteq\mathcal F$
jego suma $M=\bigcup\mathcal C$ należy do $\mathcal F$:
każdy skończony $Y\subseteq M$ jest zawarty w pewnym
$X\in\mathcal C$, bo finitarnie wiele jego elementów
pochodzi z finitarnie wielu elementów łańcucha;
wystarczy wybrać największy z tych skończenie wielu
zbiorów łańcucha. Stąd $Y\in\mathcal F$, a skończony
charakter daje $M\in\mathcal F$. Dla pustego łańcucha
suma jest pusta i należy do $\mathcal F$ z jej niepustości
i skończonego charakteru. Z lematu Kuratowskiego--Zorna
(Twierdzenie~\ref{thm:choice-zorn}) istnieje element
maksymalny rodziny.

\emph{Tukey $\Rightarrow$ AC.}
Niech $(A_i)_{i\in I}$ będzie rodziną niepustych zbiorów.
Niech $E=\bigcup_{i\in I}(\{i\}\times A_i)$ i niech
$\mathcal F$ składa się ze wszystkich podzbiorów $F\subseteq E$,
które są grafami funkcji częściowych $f:J\to\bigcup_iA_i$
o dziedzinie $J\subseteq I$, z $f(i)\in A_i$ dla $i\in J$.
Ta rodzina ma skończony charakter: jeżeli relacja
$F\subseteq E$ nie jest grafem funkcji, zawiera dwie
różne pary o tym samym pierwszym współrzędnym;
naruszenie własności jest więc widoczne już w
dwuelementowym podzbiorze. Jeśli każdy skończony
fragment jest grafem funkcji, to jest nim całe $F$.
Z Tukeya istnieje graf maksymalny $F_\ast\in\mathcal F$.
Gdyby jakiś indeks $i\in I$ nie należał do
$\operatorname{dom}F_\ast$, z niepustości $A_i$ moglibyśmy
wybrać jeden element $a\in A_i$ i dołączyć $(i,a)$,
otrzymując większą funkcję częściową. Sprzeczność.
Zatem $F_\ast$ jest grafem funkcji wyboru o dziedzinie $I$.
\end{proof}

\begin{uwaga}[BPI nie jest pełnym wyborem]
\label{rem:bpi-vs-ac}
Mamy w ZF łańcuch implikacji
\[
\mathrm{AC}\Longrightarrow \mathrm{BPI}
\Longleftrightarrow\mathrm{UF}
\Longleftrightarrow\text{Tychonoff dla zwartych Hausdorffa}.
\]
Lemat Tukeya jest \emph{równoważny AC}, a więc
znacznie silniejszy od BPI. Twierdzenie, że BPI nie
implikuje AC nad ZF, jest wynikiem modelowym
(teoria ZF+BPI+$\neg$AC jest względnie niesprzeczna
przy odpowiednich założeniach) i nie wynika z samego
wyżej dowiedzionego łańcucha implikacji.
\end{uwaga}

\begin{uwaga}[Zakres zasady wyboru w dowodzie Tychonoffa]
\label{rem:tychonoff-choice-distinction}
Zwartość pojedynczej przestrzeni nie wymaga wyboru punktu
z każdego czynnika produktu. W kierunku
\emph{BPI$\Rightarrow$Tychonoff dla Hausdorffa} punkt graniczny
ultrafiltru jest \emph{jednoznaczny}, co pozwala utworzyć
cały wektor przez zastępowanie. W ogólnym przypadku
zwartym, ale nie-Hausdorffa, zbiór punktów granicznych
może mieć wiele elementów: do jednoczesnego wyboru
potrzeba AC. To jest matematyczna przyczyna różnicy
między dwiema wersjami twierdzenia, a nie tylko
subtelność użytej techniki dowodu.
\end{uwaga}
\chapter{Gry nieskończone i aksjomat determinacji}\label{ch:determinacy}
\chaptermark{Gry nieskończone i determinacja}
\section{Gra o doskonałej informacji}
W skończonej grze można często wyznaczyć zwycięzcę przez indukcję wsteczną. Gry o nieskończenie wielu ruchach pokazują, jak bardzo zmienia się sytuacja po przejściu do ciągów nieskończonych.
\begin{definicja}[Gra $G(A)$]\label{def:game-a}
Niech $A\subseteq\omega^\omega$, gdzie $\omega^\omega$ jest zbiorem wszystkich funkcji z $\omega$ w $\omega$. Dwaj gracze, I i II, na przemian wybierają liczby naturalne $a_0,a_1,a_2,\ldots$; gracz I wykonuje ruchy o indeksach parzystych, a II o nieparzystych. Gracz I wygrywa dokładnie wtedy, gdy otrzymany ciąg $a\in\omega^\omega$ należy do $A$. W przeciwnym razie wygrywa II.
\end{definicja}
\begin{definicja}[Strategia i determinacja]\label{def:strategy-ad}
Strategią gracza I jest funkcja przypisująca każdemu skończonemu ciągowi ruchów o parzystej długości jego następny ruch. Strategią gracza II jest funkcja przypisująca każdemu skończonemu ciągowi o nieparzystej długości jego następny ruch. Strategia jest wygrywająca, jeżeli zapewnia zwycięstwo przeciw każdemu możliwemu ciągowi odpowiedzi przeciwnika. Gra jest \emph{zdeterminowana}, jeśli któryś z graczy ma strategię wygrywającą. Aksjomat determinacji $\mathrm{AD}$ głosi, że każda gra $G(A)$ dla $A\subseteq\omega^\omega$ jest zdeterminowana.
\end{definicja}
\begin{przyklad}[Wymuszenie pierwszego ruchu]
Niech $A=\{a\in\omega^\omega:a(0)=0\}$. Gracz I wygrywa, wybierając w pierwszym ruchu zero; dalsze ruchy nie mają znaczenia. Jeśli $A=\{a:a(1)=0\}$, analogicznie wygrywa gracz II.
\end{przyklad}
\section{Gry otwarte i zamknięte}
\begin{definicja}[Podstawowe otoczenia w przestrzeni Baire'a]\label{def:baire-cylinders}
Dla skończonego ciągu $s\in\omega^{<\omega}$ definiujemy cylinder $[s]=\{a\in\omega^\omega:s\subseteq a\}$. Zbiór $U\subseteq\omega^\omega$ jest otwarty, jeżeli dla każdego $a\in U$ istnieje skończony początkowy fragment $s\subseteq a$ taki, że $[s]\subseteq U$. Zbiór jest zamknięty, jeżeli jego dopełnienie jest otwarte.
\end{definicja}
\begin{twierdzenie}[Determinacja gier otwartych; twierdzenie Gale'a--Stewarta]\label{thm:open-determinacy}
Już w ZF, bez aksjomatu wyboru, jeśli
$A\subseteq\omega^\omega$ jest otwarty,
to gra $G(A)$ jest zdeterminowana.
\end{twierdzenie}
\begin{proof}
Niech $T:=\omega^{<\omega}$ będzie zbiorem
wszystkich skończonych pozycji gry.
Gracz I ma ruch po pozycjach parzystej
długości, gracz II po nieparzystej.
Dla $s\in T$ piszemy
$s^\frown n$ na oznaczenie pozycji
przedłużonej o jeden ruch $n$.

\emph{Krok 1: rekurencyjny zbiór pozycji wygrywanych.}
Niech
$W_0:=\{s\in T:[s]\subseteq A\}$.
Dla dowolnego $W_\alpha\subseteq T$
definiujemy
\[
\begin{split}
W_{\alpha+1}:=W_\alpha
 &\cup\{s\in T:|s|\text{ parzyste},
       \ \exists n\in\omega\
                    (s^\frown n\in W_\alpha)\}\\
 &\cup\{s\in T:|s|\text{ nieparzyste},
       \ \forall n\in\omega\
                    (s^\frown n\in W_\alpha)\}.
\end{split}
\]
Dla niezerowych ordynałów granicznych
$\lambda$ kładziemy
$W_\lambda:=\bigcup_{\alpha<\lambda}W_\alpha$.
Wszystkie wyrazy są podzbiorami $T$,
a operacja przejścia jest jednoznaczna.
Dlatego Twierdzenie~\ref{thm:ordinal-recursion}
zapewnia istnienie rodziny na dowolnym
ustalonym odcinku ordynałów, bez AC.

\emph{Krok 2: stabilizacja.}
Z Twierdzenia Hartogsa~\ref{thm:choice-hartogs}
dla zbioru $\mathcal P(T)$ mamy
dobrze uporządkowany zbiór, którego nie
można różnowartościowo odwzorować
w $\mathcal P(T)$. Biorąc jego
typ porządkowy (konstrukcja w dowodzie
Twierdzenia~\ref{thm:initial-exists},
niewymagająca wyboru dla już
dobrze uporządkowanego zbioru),
otrzymujemy ordynał $\theta$
bez iniekcji $\theta\to\mathcal P(T)$.
Skonstruujmy $W_\alpha$ dla
$\alpha\le\theta$.
Gdyby $W_\alpha\ne W_{\alpha+1}$
dla wszystkich $\alpha<\theta$,
monotoniczność rodziny dawałaby
iniekcję
$\alpha\mapsto W_\alpha$
ze zbioru $\theta$ do $\mathcal P(T)$.
To sprzeczność.
Istnieje więc $\eta<\theta$
z $W_\eta=W_{\eta+1}$,
a rekurencja daje tę samą wartość
na wszystkich dalszych etapach.
Niech $W:=W_\eta$.
Dla każdego $s\in W$ niech
$\rho(s)$ będzie najmniejszym
$\alpha\le\eta$ takim, że
$s\in W_\alpha$.
Funkcja rang istnieje przez
jednoznaczność i zastępowanie.

\emph{Krok 3: jeśli $\varnothing\in W$,
wygrywa I.}
Jeśli $s\in W_0$, to
$[s]\subseteq A$, więc każdy
wynik z przedłużeniem $s$ jest w $A$.
Dla $s\in W\setminus W_0$
ranga $\rho(s)$ jest następnikiem
$\alpha+1$, bo na etapach granicznych
nie pojawiają się nowe elementy.
Jeśli ruch ma gracz I,
istnieje $n\in\omega$
z $s^\frown n\in W_\alpha$.
I wybiera \emph{najmniejszego}
takiego $n$.
Jeżeli ruch ma II, z definicji
$W_{\alpha+1}$ każdy następnik
$s^\frown n$ ma rangę najwyżej $\alpha$.
Na wszystkich ruchach, dopóki
nie osiągnięto $W_0$, ranga ściśle maleje.
Nie istnieje nieskończony ściśle malejący
ciąg ordynałów: zbiór jego wartości
miałby element najmniejszy,
a następny wyraz byłby mniejszy.
Zatem po skończenie wielu ruchach
osiągamy $W_0$ i I wygrywa.
Dalsze ruchy I można ustalić jako $0$.

\emph{Krok 4: jeśli $\varnothing\notin W$,
wygrywa II.}
Dla pozycji $s\notin W$, z której
ruch ma I, każdy jej następnik
leży poza $W$: w przeciwnym razie
jego pewna ranga wprowadziłaby
$s$ do $W$.
Jeżeli ruch ma II, istnieje
przynajmniej jeden następnik
$s^\frown n\notin W$.
Gdyby wszystkie następniki należały
do $W$, ich rangi utworzyłyby
zbiór ordynałów. Kładąc
\[
\beta:=\sup_{n\in\omega}
              (\rho(s^\frown n)+1),
\]
otrzymalibyśmy
$s^\frown n\in W_\beta$
dla wszystkich $n$, a więc
$s\in W_{\beta+1}\subseteq W$,
sprzeczność (stabilność rodziny).
II wybiera najmniejszego $n$
z $s^\frown n\notin W$.
Cała rozgrywka pozostaje poza $W_0$.
Gdyby jej wynik należał do
otwartego $A$, pewien jego
skończony prefiks $s$ spełniałby
$[s]\subseteq A$,
a więc $s\in W_0$.
To niemożliwe, zatem II wygrywa.

Obie strategie są funkcjami
zadanymi jednoznacznie przez
najmniejszy dopuszczalny ruch
w $\omega$; w żadnym miejscu
nie zastosowano AC.
\end{proof}

\begin{przyklad}[Dwa kroki i kolejność kwantyfikatorów]
\label{ex:ad-two-step-quantifiers}
Niech
$A=\{a\in\omega^\omega:a(0)=0
\text{ oraz }a(1)\ne0\}$.
Jeżeli I w pierwszym ruchu wybierze
$a(0)\ne0$, przegrywa niezależnie od
odpowiedzi przeciwnika. Gdy wybierze
$a(0)=0$, gracz II odpowiada $a(1)=0$
i również zapewnia wypłatę spoza $A$.
Dlatego II ma strategię wygrywającą:
w pierwszym swoim ruchu gra $0$,
a później także może grać $0$.

Warunek dwóch pierwszych ruchów ma postać
\[
\forall u\in\omega\;\exists v\in\omega\;
   (u\ne0\;\lor\;v=0).
\]
Nie wystarcza więc, że istnieje para
ruchów dająca wypłatę należącą do $A$:
gracz I musi mieć odpowiedź zwycięską
na \emph{każdy} możliwy ruch przeciwnika.
\end{przyklad}
\begin{figure}[htbp]
\centering
\begin{tikzpicture}[
  level distance=13mm,
  every node/.style={font=\small,align=center},
  edge from parent/.style={draw,-latex},
  level 1/.style={sibling distance=70mm},
  level 2/.style={sibling distance=30mm}
]
\node {$\varnothing$}
  child {node {$a(0)=0$}
    child {node {$a(1)=0$\\II wygrywa}
      edge from parent node[left] {\scriptsize II}}
    edge from parent node[left] {\scriptsize I}}
  child {node {$a(0)\ne0$\\II wygrywa}
    edge from parent node[right] {\scriptsize I}};
\end{tikzpicture}
\caption{Drzewko dwóch ruchów: na każdą decyzję
gracza I istnieje odpowiedź zapewniająca
zwycięstwo graczowi II. Późniejsze ruchy
nie zmieniają wyniku.}
\label{fig:ad-two-moves}
\end{figure}

\begin{uwaga}[Dlaczego AD nie jest aksjomatem ZFC?]\label{rem:ad-ac-incompat}
Pełny aksjomat determinacji jest sprzeczny z AC: w ZFC istnieją zbiory wypłat, dla których żaden gracz nie ma strategii wygrywającej. Ścisły dowód wymaga konstrukcji z wykorzystaniem dobrego uporządkowania przestrzeni strategii; przedstawiamy ją poniżej. AD rozważa się zatem przede wszystkim jako dodatkowy aksjomat nad ZF, nie nad ZFC.
\end{uwaga}
\begin{twierdzenie}[AC implikuje istnienie gry niezdeterminowanej]\label{thm:ac-not-ad}
W ZFC istnieje zbiór $A\subseteq\omega^\omega$ taki, że gra $G(A)$ nie jest zdeterminowana.
\end{twierdzenie}
\begin{proof}
Niech $\mathfrak c=|\omega^\omega|$. Z AC dobrze porządkujemy wszystkie strategie obu graczy, łącznie jako rodzinę długości $\mathfrak c$. Dla każdej strategii gracza I wybierzemy rozgrywkę zgodną z tą strategią, która zakończy się przegraną I; dla każdej strategii gracza II wybierzemy rozgrywkę zgodną z nią, która zakończy się wygraną I. Wykonujemy te wybory rekurencyjnie po początkowym ordynale mocy $\mathfrak c$, zachowując dwie rozłączne rodziny rozgrywek $L$ i $H$.

Na etapie $\alpha<\mathfrak c$ ustalona strategia ma $\mathfrak c$ zgodnych z nią rozgrywek: przeciwnik może dowolnie wybierać nieskończony ciąg swoich ruchów. Wcześniej oznaczono mniej niż $\mathfrak c$ rozgrywek (dla nieskończonej liczby kardynalnej $\mathfrak c$ i $\alpha<\mathfrak c$ zachodzi $|\alpha|\cdot\aleph_0<\mathfrak c$), więc można wybrać świeżą rozgrywkę. Dla strategii I umieszczamy ją w $L$, dla strategii II w $H$. Wybór świeżego elementu ustalamy przez globalny dobry porządek $\omega^\omega$. Na końcu niech $A=H$. Każda strategia I ma zgodną rozgrywkę w $L\subseteq\omega^\omega\setminus A$, więc nie wygrywa zawsze. Każda strategia II ma zgodną rozgrywkę w $H=A$, więc również nie wygrywa zawsze. Gra $G(A)$ jest niezdeterminowana.
\end{proof}
\begin{uwaga}[Granice przedstawionej metody]
Determinacja gier otwartych nie oznacza determinacji wszystkich gier. Dowód powyżej wykorzystuje pełny AC; nie wolno go przenosić do ZF+AD. W kolejnych rozdziałach omówimy konstrukcję modeli wewnętrznych i wymuszanie jako metody precyzyjnego badania niezależności aksjomatów.
\end{uwaga}

\section{Przykład gry zamkniętej}
\begin{przyklad}[Warunek bezpieczeństwa]
Niech $A=\{a\in\omega^\omega:\forall n\in\omega\ a(2n)=0\}$. Zbiór $A$ jest zamknięty: jeśli $a\notin A$, to pewien parzysty wyraz $a(2n)$ jest niezerowy, a każdy ciąg o tym samym prefiksie długości $2n+1$ również leży poza $A$. Gracz I wygrywa, zawsze grając zero. Ten przykład pokazuje, że zamkniętość warunku wypłaty nie przesądza o tym, który gracz wygrywa.
\end{przyklad}
\begin{wniosek}[Determinacja gier zamkniętych]\label{cor:closed-determinacy}
W ZF każda gra $G(A)$ z zamkniętym
$A\subseteq\omega^\omega$ jest zdeterminowana.
\end{wniosek}
\begin{proof}
Niech $B:=\omega^\omega\setminus A$,
który jest otwarty.
Konstruujemy zbiory $V_\alpha$
tak samo jak w dowodzie
Twierdzenia~\ref{thm:open-determinacy},
lecz ze \emph{zamienionymi rolami}:
\[
V_0:=\{s:[s]\subseteq B\}.
\]
Na pozycji, w której ruch ma II,
włączamy $s$ do $V_{\alpha+1}$
po istnieniu choć jednego następnika
w $V_\alpha$;
na ruchu I wymagamy, aby
\emph{wszyscy} następnicy
należeli do $V_\alpha$.
Zachowujemy już uzyskane pozycje,
a na etapach granicznych bierzemy
ich sumy.
Z Twierdzenia Hartogsa i rekurencji
wynika stabilizacja dokładnie
jak dla gier otwartych.

Jeśli pozycja pusta należy
do stabilnego zbioru, II
na swoich ruchach wybiera
najmniejszego następnika
o mniejszej randze.
Na ruchach I wszystkie następniki
mają rangę mniejszą.
Dobrze ufundowane rangi wymuszają
osiągnięcie pozycji $s\in V_0$,
a wtedy każdy wynik gry leży w $B$
i wygrywa II.

W przeciwnym przypadku I
na swoich ruchach wybiera
najmniejszego następnika
poza stabilnym zbiorem,
a na ruchach II wszystkie
następniki pozostają poza nim.
Wynik nie może należeć do
otwartego $B$, gdyż miałby
prefiks w $V_0$.
Zatem wynik należy do $A$
i wygrywa I.
Wszystkie wybory są kanoniczne
i wykonane w ZF.
\end{proof}


\section{Własność Baire'a, determinacja projektywna i ich następstwa}
\label{sec:ad-regularity}

\subsection{Z jakich założeń wolno korzystać?}
W rozdziale~\ref{ch:determinacy} zdefiniowaliśmy grę $G(A)$,
strategię wygrywającą i aksjomat determinacji $\mathrm{AD}$.
Dotychczas pracowaliśmy zwykle w ZFC. Teraz wyraźnie zmieniamy
teorię bazową: przy rozważaniu pełnego $\mathrm{AD}$ pracujemy
w $\mathrm{ZF}+\mathrm{AD}$, \emph{bez} aksjomatu wyboru.
Twierdzenie~\ref{thm:ac-not-ad} pokazało, że $\mathrm{AD}$ i
$\mathrm{AC}$ nie mogą obowiązywać jednocześnie.

Nie wolno utożsamiać trzech zdań:
\begin{enumerate}
\item wszystkie gry $G(A)$ na liczbach naturalnych są zdeterminowane ($\mathrm{AD}$);
\item gry $G(A)$ są zdeterminowane tylko dla \emph{projektywnych} zbiorów wypłat ($\mathrm{PD}$);
\item każda gra z wypłatą otwartą jest zdeterminowana
      (twierdzenie Gale'a--Stewarta, udowodnione wcześniej).
\end{enumerate}
Pierwsze zdanie jest niezgodne z $\mathrm{AC}$; drugie może być
rozważane wraz z ZFC i wynika z odpowiednio silnych założeń
o dużych liczbach kardynalnych; trzecie jest zwykłym twierdzeniem ZF
(Twierdzenie~\ref{thm:open-determinacy}
i Wniosek~\ref{cor:closed-determinacy},
wykazane bez AC).

\begin{lemat}[AD pozwala wybierać z przeliczalnej rodziny zbiorów rzeczywistych]
\label{lem:ad-countable-choice-reals}
W $\mathrm{ZF}+\mathrm{AD}$ dla każdej funkcji
$n\mapsto X_n$ ($n\in\N$), której wartości są
niepustymi podzbiorami $\N^\N$, istnieje funkcja
$n\mapsto x_n$ taka, że $x_n\in X_n$ dla każdego $n$.
To jest \emph{wybór przeliczalny dla zbiorów liczb rzeczywistych},
a nie pełny aksjomat wyboru.
\end{lemat}
\begin{proof}
Zdefiniujmy grę o ruchach będących liczbami naturalnymi.
W swoim pierwszym ruchu gracz I wybiera indeks $n$.
Pozostałe ruchy gracza I nie mają znaczenia dla wypłaty.
Kolejne ruchy gracza II tworzą nieskończony ciąg $x$.
Gracz II wygrywa wtedy i tylko wtedy, gdy $x\in X_n$.

Nie istnieje wygrywająca strategia gracza I.
Istotnie, każda jego ustalona strategia na początku
wyznacza konkretny indeks $n$. Z niepustości $X_n$
możemy wybrać pojedynczy ciąg $x\in X_n$ i polecić
II rozegrać go współrzędna po współrzędnej, niezależnie
od dalszych ruchów I. Ta odpowiedź pokonuje rozważaną
strategię I. Wybór jednego świadka po ustaleniu $n$
nie jest zastosowaniem nieskończonej zasady wyboru.
Z AD wynika, że II ma strategię wygrywającą $\tau$.
Dla każdego $n$ ustalmy rozgrywkę, w której I gra
najpierw $n$, a potem same zera. Kolejne odpowiedzi
$\tau$ tworzą ciąg $x_n\in X_n$. Przyporządkowanie
$n\mapsto x_n$ jest funkcją wyboru; zdefiniowaliśmy
ją jednym wzorem z użyciem tej samej strategii $\tau$.
\end{proof}

\subsection{Zbiory pierwszej kategorii i własność Baire'a}
\begin{definicja}[Zbiory nigdziegęste i pierwszej kategorii]
\label{def:ad-meager}
Niech $X$ będzie przestrzenią topologiczną.
Zbiór $F\subseteq X$ jest \emph{nigdziegęsty}, jeżeli
wnętrze jego domknięcia jest puste:
$\operatorname{int}_X\overline F=\varnothing$.
Zbiór $M\subseteq X$ jest \emph{pierwszej kategorii},
inaczej \emph{mizerny}, jeśli istnieją zbiory nigdziegęste
$F_n\subseteq X$ dla $n\in\N$ takie, że
$M\subseteq\bigcup_{n\in\N}F_n$.
Zbiór jest \emph{rezydualny},
jeżeli jego dopełnienie jest pierwszej kategorii.
\end{definicja}
Przez \emph{rezydualny} rozumiemy więc ``dopełnienie małe
w sensie kategorii'', a nie ``dopełnienie jest przeliczalne''.
Dla przykładu $\Q$ jest pierwszej kategorii w $\R$,
gdyż jest przeliczalną sumą singletonów, z których każdy
jest nigdziegęsty; $\R\setminus\Q$ jest rezydualny.

\begin{definicja}[Własność Baire'a]
\label{def:ad-baire-property}
Zbiór $A\subseteq X$ ma \emph{własność Baire'a}, gdy
istnieje zbiór otwarty $U\subseteq X$ taki, że
$A\mathbin{\triangle}U=(A\setminus U)\cup(U\setminus A)$
jest pierwszej kategorii w $X$.
\end{definicja}
W teorii ZF przeliczalna suma zbiorów pierwszej kategorii
nie musi automatycznie być pierwszej kategorii: może
zabraknąć jednoczesnego wyboru ich przeliczalnych pokryć.
W naszych dowodach pod AD uzasadnia ją
Lemat~\ref{lem:ad-meager-countable-union}:
wykorzystujemy wcześniej dowiedziony
wybór przeliczalny dla zbiorów
liczb rzeczywistych, a nie pełny AC.


\begin{lemat}[Przeliczalna suma zbiorów pierwszej kategorii w AD]
\label{lem:ad-meager-countable-union}
W teorii ZF+$\mathrm{AD}$, jeśli
$(M_n)_{n\in\omega}$ jest ciągiem
podzbiorów przestrzeni Baire'a
$\omega^\omega$ pierwszej kategorii,
to $\bigcup_{n\in\omega}M_n$
jest pierwszej kategorii.
Nie zakładamy pełnego AC ani DC.
\end{lemat}
\begin{proof}
\emph{Krok 1: kody domkniętych zbiorów.}
Niech $T:=\omega^{<\omega}$.
Ponieważ $T$ ma ustaloną enumerację
(Twierdzenie~\ref{thm:card-n2}),
każdy jego podzbiór można jednoznacznie
zakodować elementem $2^\omega$.
Dla domkniętego $F\subseteq\omega^\omega$
połóżmy
\[
T_F:=\{s\in T:[s]\cap F\ne\varnothing\}.
\]
Jest to zbiór istniejący przez separację.
Dla $x\in\omega^\omega$ mamy
\[
x\in F\quad\Longleftrightarrow\quad
\forall k\in\omega\;(x|_k\in T_F).
\]
Kierunek wprost jest natychmiastowy;
w przeciwnym razie otwartość
$\omega^\omega\setminus F$ daje
cylinder $[x|_k]$ rozłączny z $F$.
Zatem kod $T_F$ jednoznacznie wyznacza $F$.
Ponadto $F$ jest nigdziegęsty dokładnie
wtedy, gdy żaden cylinder nie zawiera się
w $F$, co przy domknięciu $F$ daje
\[
\forall s\in T\;\exists t\sqsupseteq s
\quad t\notin T_F.
\]
Rzeczywiście z pustego wnętrza
domknięcia otrzymujemy dla każdego $[s]$
mniejszy cylinder $[t]$ rozłączny z $F$.

\emph{Krok 2: kod świadectwa małości.}
Dla każdego $M_n$ istnieje ciąg
$(F_{n,j})_{j\in\omega}$ zbiorów
nigdziegęstych pokrywających $M_n$.
Każdy z nich można zastąpić
domknięciem, które nadal jest nigdziegęste.
Zakodujmy ciąg kodów ich drzew
jednym elementem $2^\omega$,
korzystając ze zdefiniowanej wcześniej
bijekcji $\omega\times\omega\to\omega$
(Twierdzenie~\ref{thm:card-n2}).
Niech $C_n\subseteq2^\omega$ będzie
zbiorem wszystkich poprawnych kodów
takich świadectw pokrycia $M_n$.
Jest on niepusty, a przez separację
w $2^\omega$ istnieje w ZF.

\emph{Krok 3: zastosowanie AD.}
Lemat~\ref{lem:ad-countable-choice-reals}
daje jedną funkcję $n\mapsto c_n$
z $c_n\in C_n$ dla wszystkich $n$.
Dekodując $c_n$ otrzymujemy podwójnie
indeksowaną rodzinę
$(F_{n,j})_{(n,j)\in\omega^2}$
domkniętych zbiorów nigdziegęstych
z
\[
\bigcup_{n\in\omega}M_n
 \subseteq\bigcup_{(n,j)\in\omega^2}F_{n,j}.
\]
Prawa strona jest przeliczalną sumą
zbiorów nigdziegęstych, ponieważ
$\omega^2$ jest jawnie przeliczalny.
Zatem $\bigcup_nM_n$ jest pierwszej kategorii.
Nie użyto AC poza konsekwencją AD
z wcześniejszego lematu.
\end{proof}

Nie należy mylić tej własności \emph{zbioru} z twierdzeniem
Baire'a o \emph{przestrzeni}: w przestrzeni zupełnej
metrycznie żaden niepusty zbiór otwarty nie jest
pierwszej kategorii w tej przestrzeni.

\begin{lemat}[Postać konstrukcyjna twierdzenia Baire'a]
\label{lem:ad-baire-category}
Niech $X=\N^\N$ z topologią, której bazą są cylindry
$[s]=\{x\in\N^\N:s\text{ jest początkowym fragmentem }x\}$,
$s\in\N^{<\omega}$. Jeżeli $D_n$ ($n\in\N$) są otwarte
i gęste w $X$, to $\bigcap_nD_n$ przecina każdy niepusty
cylinder $[s]$. W szczególności żaden cylinder $[s]$
nie jest pierwszej kategorii.
\end{lemat}
\begin{proof}
Ustalmy $[s]$. Wybieramy najpierw najmniejsze w ustalonej
enumeracji skończonych ciągów przedłużenie $s_0\supseteq s$
takie, że $[s_0]\subseteq D_0$; istnieje ono, bo $D_0$
jest otwarty i gęsty. Następnie indukcyjnie wybieramy
najmniejsze przedłużenie $s_{n+1}\supseteq s_n$ długości
większej niż $n+1$, dla którego $[s_{n+1}]\subseteq D_{n+1}$.
Otrzymujemy funkcję $x=\bigcup_ns_n\in\N^\N$.
Dla każdego $n$ zachodzi $x\in[s_n]\subseteq D_n$,
a zarazem $x\in[s]$. Nie użyliśmy tu dowolnego wyboru
nieskończenie wielu świadków: w każdym kroku wskazaliśmy
pierwszego świadka z ustalonego przeliczalnego porządku.

Gdyby $[s]$ był zawarty w $\bigcup_nF_n$ dla zbiorów
nigdziegęstych, zbiory $D_n=X\setminus\overline F_n$
byłyby otwarte i gęste. Skonstruowany punkt leżałby
w $[s]$, ale poza każdym $F_n$, co daje sprzeczność.
\end{proof}

\subsection{Gra Banacha--Mazura}
Zapis $s\sqsubseteq t$ oznacza, że skończony ciąg $s$
jest początkowym fragmentem ciągu $t$.
\begin{definicja}[Gra Banacha--Mazura nad $A$]
\label{def:ad-banach-mazur}
Niech $A\subseteq\N^\N$. Gracze I i II na przemian
wybierają niepuste skończone ciągi $s_0,s_1,\ldots$
tak, że $s_n$ jest właściwym początkowym fragmentem
$s_{n+1}$. Powstaje jednoznacznie $x=\bigcup_ns_n\in\N^\N$.
Gracz I wygrywa wtedy i tylko wtedy, gdy $x\in A$.
Grę tę oznaczamy przez $\operatorname{BM}(A)$.
\end{definicja}
Każdy ruch jest zakodowany liczbą naturalną, ponieważ zbiór
skończonych ciągów liczb naturalnych jest przeliczalny.
Nielegalny ruch oznacza przegraną gracza, który go wykonał.
W konsekwencji z $\mathrm{AD}$ wynika determinacja tej gry:
dopuszczalne pozycje i reguła przegranej dają grę na
liczbach naturalnych, do której stosuje się definicja AD.

\begin{lemat}[Charakteryzacja strategii przez kategorię]
\label{lem:ad-banach-mazur-category}
Niech $A\subseteq\N^\N$.
\begin{enumerate}
\item Jeżeli gracz II ma strategię wygrywającą w
$\operatorname{BM}(A)$, to $A$ jest pierwszej kategorii.
\item Jeżeli gracz I ma strategię wygrywającą, to istnieje
niepusty cylinder $[s]$ taki, że $[s]\setminus A$
jest pierwszej kategorii w $\N^\N$.
\end{enumerate}
\end{lemat}
\begin{proof}
\emph{Etap 1: strategia gracza II.}
Niech $\tau$ będzie strategią wygrywającą II. Rozpatrzmy
wszystkie skończone legalne historie $h$ zgodne z $\tau$
i kończące się ruchem gracza II; dołączmy też historię pustą.
Jest ich co najwyżej przeliczalnie wiele. Niech $t_h$
będzie ostatnim ciągiem historii $h$ (dla historii pustej
przyjmujemy $[t_h]=\N^\N$).
Niech $V_h$ będzie sumą cylindrów $[u]$, gdzie $u$
jest odpowiedzią $\tau$ na jakiś dozwolony ruch gracza I
zaczynający się po $h$. Wówczas $V_h$ jest zbiorem
otwartym i gęstym w $[t_h]$: w dowolnym mniejszym
cylindrze $[v]\subseteq[t_h]$ gracz I może zagrać
odpowiednio długie przedłużenie $v$, a $\tau$ odpowie
jeszcze dłuższym ciągiem $u\supseteq v$.
Zatem $D_h=(\N^\N\setminus[t_h])\cup V_h$ jest
otwarty i gęsty w $\N^\N$.

Niech $D=\bigcap_hD_h$. Dopełnienie $D$ jest pierwszej
kategorii, gdyż historie indeksowane są przeliczalnie.
Dla $x\in D$ odtwarzamy zgodną z $\tau$ rozgrywkę,
w której wszystkie wybrane ciągi są początkowymi fragmentami
$x$: przynależność $x\in V_h$ pozwala wybrać odpowiedni
ruch I, a $\tau$ daje odpowiedź; pierwszą taką możliwość
bierzemy z ustalonej enumeracji skończonych ciągów.
Wynikiem gry jest $x$. Ponieważ $\tau$ wygrywa,
$x\notin A$, czyli $A\subseteq\N^\N\setminus D$.
Stąd $A$ jest pierwszej kategorii.

\emph{Etap 2: strategia gracza I.}
Niech $\sigma$ będzie strategią wygrywającą gracza I
i niech $s=\sigma(\varnothing)$ będzie jej pierwszym ruchem.
Ograniczmy uwagę do cylindra $[s]$. Po każdym skończonym
fragmencie gry zgodnym z $\sigma$ i zakończonym ruchem I
gracz II może wykonać dowolnie długie przedłużenie,
a $\sigma$ odpowie jeszcze dłuższym ciągiem.
Powtarzając konstrukcję otwartych i gęstych zbiorów
z etapu 1, lecz rozpoczynając od $[s]$ i zamieniając
role graczy, otrzymujemy zbiór rezydualny $D\subseteq[s]$
taki, że każdy punkt $x\in D$ jest wynikiem pewnej
rozgrywki zgodnej z $\sigma$. Wszystkie te punkty leżą
w $A$, bo $\sigma$ jest wygrywająca. Zatem $[s]\setminus A$
jest pierwszej kategorii w $[s]$, a ponieważ $[s]$ jest
otwarto-domknięty, także w całym $\N^\N$.
\end{proof}

\begin{twierdzenie}[AD implikuje własność Baire'a]
\label{thm:ad-baire-property}
W teorii $\mathrm{ZF}+\mathrm{AD}$ każdy podzbiór
$\N^\N$ ma własność Baire'a. To samo zachodzi dla
każdego podzbioru $\R$ w jej zwykłej topologii.
\end{twierdzenie}
\begin{proof}
\emph{Krok 1: część otwarta dla przestrzeni Baire'a.}
Ustalmy $A\subseteq\N^\N$. Niech
\[
U=\bigcup\{[s]:[s]\setminus A\text{ jest pierwszej kategorii}\}.
\]
Jest to zbiór otwarty. Ponieważ baza cylindrów jest
przeliczalna, $U\setminus A$ jest pierwszej kategorii:
stanowi podzbiór przeliczalnej sumy zbiorów $[s]\setminus A$.

\emph{Krok 2: część poza $U$.}
Zagrajmy w $\operatorname{BM}(A\setminus U)$.
Gdyby gracz I miał strategię wygrywającą, Lemat
\ref{lem:ad-banach-mazur-category} dałby niepusty
cylinder $[s]$ taki, że $[s]\setminus(A\setminus U)$
jest pierwszej kategorii. Ponieważ
$[s]\setminus A\subseteq[s]\setminus(A\setminus U)$,
mielibyśmy $[s]\subseteq U$. Wtedy
$[s]\setminus(A\setminus U)=[s]$ byłby pierwszej
kategorii, wbrew Lematowi~\ref{lem:ad-baire-category}.
Zatem I nie ma strategii wygrywającej. Z $\mathrm{AD}$
wynika, że wygrywa II, a Lemat
\ref{lem:ad-banach-mazur-category} pokazuje, że
$A\setminus U$ jest pierwszej kategorii.
Obie strony różnicy symetrycznej $A\triangle U$ są zatem
pierwszej kategorii.

\emph{Krok 3: przejście do prostej.}
Przestrzeń niewymiernych liczb rzeczywistych
$\R\setminus\Q$ jest homeomorficzna z $\N^\N$.
Aby zrozumieć tę identyfikację, najpierw ograniczmy się
do liczb niewymiernych z $(0,1)$. Każda taka liczba
ma jednoznaczne nieskończone rozwinięcie na ułamek
łańcuchowy $[0;a_0,a_1,\ldots]$, gdzie $a_n\geq1$.
Ucięcie rozwinięcia po skończenie wielu współczynnikach
wyznacza przedział możliwych dalszych wartości;
przedziały te tworzą bazę topologii niewymiernych
w $(0,1)$. Odwzorowanie przypisujące liczbie ciąg
$(a_0-1,a_1-1,\ldots)$ i odwzorowanie odwrotne
zachowują więc zbieżność ciągów: dostatecznie długi
wspólny początek współczynników oznacza małą różnicę
liczb, a dla liczby niewymiernej dostatecznie mała
zmiana wartości nie zmienia ustalonego początkowego
fragmentu jej rozwinięcia. Są zatem ciągłe w obie strony.
Całą $\R\setminus\Q$ rozbijamy na przeliczalnie wiele
otwarto-domkniętych podprzestrzeni $(n,n+1)\setminus\Q$,
$n\in\Z$, i stosujemy powyższą konstrukcję w każdej
z nich. Ponieważ $\N\times\N^\N$ jest homeomorficzne
z $\N^\N$, otrzymujemy żądany homeomorfizm.
Homeomorfizm zachowuje zbiory pierwszej kategorii,
więc każda część $A\cap(\R\setminus\Q)$ ma własność
Baire'a w przestrzeni liczb niewymiernych. Ponieważ
$\R\setminus\Q$ jest gęstym zbiorem $G_\delta$
w $\R$, różnica zbioru $A$ od odpowiedniego zbioru
otwartego w $\R$ jest pierwszej kategorii po dołączeniu
zbioru $\Q$. Otrzymujemy żądany wynik dla $\R$.
\end{proof}


\begin{twierdzenie}[AD i własność zbioru doskonałego]
\label{thm:ad-perfect-set}
W teorii ZF+$\mathrm{AD}$ każdy nieprzeliczalny
podzbiór przestrzeni Cantora $2^\omega$,
przestrzeni Baire'a $\omega^\omega$
lub prostej rzeczywistej $\mathbb R$
zawiera niepusty zbiór doskonały
(czyli domknięty, bez punktów izolowanych).
\end{twierdzenie}
\begin{proof}
\emph{Krok 1: gra o podzbiór doskonały.}
Najpierw ustalmy $A\subseteq2^\omega$.
Rozważmy grę, w której gracz I
w swoim $n$-tym ruchu podaje
dowolny skończony ciąg binarny $s_n$,
a gracz II odpowiada jednym bitem
$b_n\in\{0,1\}$.
Wynik rozgrywki jest ciągiem
\[
x=s_0^\frown(b_0)^\frown s_1^\frown(b_1)
          ^\frown s_2^\frown(b_2)^\frown\cdots.
\]
Jest to nieskończony ciąg, ponieważ
każdy ruch gracza II dodaje jeden bit.
Gracz I wygrywa wtedy i tylko wtedy,
gdy $x\in A$.
Zbiór skończonych ciągów binarnych
można jawnie ponumerować liczbami naturalnymi;
nielegalne numery lub ruchy
gracza II spoza $\{0,1\}$ oznaczają
przegraną gracza, który je wykonał.
Wobec tego gra jest kodowana jako gra
o ruchach w $\omega$ i podlega AD.

\emph{Krok 2: strategia II wymusza przeliczalność.}
Załóżmy, że $\tau$ jest wygrywającą strategią
gracza II.
Niech $H$ będzie zbiorem wszystkich
skończonych legalnych historii
zgodnych z $\tau$, kończących się
ruchem gracza II, wraz z historią pustą.
Dla $h\in H$ niech $p_h$
oznacza wypisany dotąd ciąg binarny,
a $\tau(h,s)\in\{0,1\}$ odpowiedź II
na kolejny ruch $s$ gracza I.
Zbiór $H$ jest przeliczalny,
gdyż należy do zbioru skończonych
ciągów liczb naturalnych.

Dla $h\in H$ połóżmy
\[
E_h:=\{x\in2^\omega:
 p_h\sqsubseteq x\ \text{oraz}\
 \forall s\in2^{<\omega}\quad
 p_h^\frown s^\frown(\tau(h,s))
              \not\sqsubseteq x\}.
\]
Każdy $E_h$ ma \emph{co najwyżej jeden}
element. Rzeczywiście, jeśli $x\in E_h$,
to następny bit po $p_h$ jest równy
$1-\tau(h,\varnothing)$.
Gdy znamy pierwszych $k$ bitów $u$
ciągu $x$ po $p_h$, jego następny bit
musi być równy $1-\tau(h,u)$.
Indukcja po $k$ wyznacza więc
wszystkie bity $x$ jednoznacznie.

Pokażemy, że
$A\subseteq\bigcup_{h\in H}E_h$.
Gdyby dla pewnego $x\in A$
nie zachodziło $x\in E_h$ dla żadnego
$h$, to zaczynając od historii pustej
moglibyśmy rekurencyjnie
wybrać pierwszy w ustalonej enumeracji
ruch $s_n$ gracza I, dla którego
\[
p_h^\frown s_n^\frown(\tau(h,s_n))
\sqsubseteq x.
\]
Taki ruch istnieje, bo $x$ przedłuża
bieżące $p_h$ i nie należy do $E_h$.
Odpowiedź $\tau$ przedłuża prefiks
o co najmniej jeden bit;
otrzymalibyśmy zatem pełną rozgrywkę
zgodną z $\tau$ o wyniku $x\in A$.
Przeczy to wygrywającemu
charakterowi $\tau$.
Zatem wskazane zawieranie zachodzi.
Ponieważ $H$ jest przeliczalny
i każdy $E_h$ zawiera najwyżej
jeden element, zbiór $A$ jest
co najwyżej przeliczalny:
przypiszmy każdemu $x\in A$
pierwszą w ustalonej enumeracji
historię $h$ z $x\in E_h$.
Jest to iniekcja do $\omega$,
bez użycia AC.

\emph{Krok 3: strategia I daje zbiór doskonały.}
Jeżeli $A$ nie jest przeliczalny,
gracz II nie ma strategii wygrywającej.
Z AD gracz I ma zatem
strategię wygrywającą $\sigma$.
Dla $b=(b_n)_{n\in\omega}\in2^\omega$
niech $f(b)$ będzie wynikiem gry,
w której II podaje kolejne bity $b_n$,
a I gra zgodnie ze strategią $\sigma$.
Funkcja $f:2^\omega\to2^\omega$
jest ciągła: aby poznać skończony
prefiks jej wartości, wystarcza
poznać skończenie wiele bitów
gracza II, ponieważ każdy z nich
wydłuża wynik o co najmniej jeden bit.
Jest też iniekcją: jeśli
$b\ne b'$, to w pierwszej
współrzędnej, w której się różnią,
poprzednie historie I są jednakowe,
ale do wyniku zostają wpisane
dwa różne bity.

Przestrzeń $2^\omega$ jest zwarta
już w ZF. Istotnie, gdyby pokrycie
otwarte nie miało skończonego
podpokrycia, drzewo cylindrów,
których nie można skończenie pokryć,
byłoby nieskończone. W każdym węźle
co najmniej jedno z dwóch dzieci
również miałoby tę własność.
Wybierając zawsze lewe takie
dziecko, gdy istnieje,
otrzymalibyśmy ciąg $b\in2^\omega$,
którego żadne podstawowe otoczenie
nie można skończenie pokryć,
sprzecznie z tym, że $b$ należy
do pewnego otwartego elementu pokrycia.

Obraz $P:=f[2^\omega]$ jest więc
zwarty, a w przestrzeni Hausdorffa
$2^\omega$ także domknięty.
Nie ma punktów izolowanych:
dla $f(b)\in P$ i jego
dowolnego otoczenia $O$ ciągłość $f$
daje cylinder $[u]\ni b$
odwzorowywany w $O$.
Istnieje $b'\ne b$ o tym samym
prefiksie $u$, a iniektywność
daje $f(b')\ne f(b)$.
Zatem $P$ jest niepusty i doskonały.
Ponieważ $\sigma$ jest wygrywająca,
$P\subseteq A$.

\emph{Krok 4: przeniesienie do
przestrzeni Baire'a i na prostą.}
Definiujemy homeomorfizm
$\theta:\omega^\omega\to D\subseteq2^\omega$,
gdzie $D$ jest zbiorem ciągów binarnych
z nieskończenie wieloma jedynkami,
przez kodowanie
\[
\theta(n_0,n_1,\ldots)
=0^{n_0}1\,0^{n_1}1\,0^{n_2}1\cdots.
\]
Z długości kolejnych bloków zer
można odzyskać wszystkie $n_k$;
obie funkcje są ciągłe,
ponieważ skończony fragment wyniku
zależy od skończenie wielu cyfr wejścia.
Dla nieprzeliczalnego
$B\subseteq\omega^\omega$
zbiór $\theta[B]\subseteq2^\omega$
jest nieprzeliczalny. Z kroków 1--3
zawiera niepusty doskonały,
a w szczególności zwarty zbiór $P$.
Zbiór $\theta^{-1}[P]$
jest zwarty i bez punktów
izolowanych, gdyż $\theta$
jest homeomorfizmem na $D$.
Jako zwarty podzbiór przestrzeni
Hausdorffa jest również domknięty,
a więc doskonały.

Na koniec zauważmy, że
$\mathbb R\setminus\mathbb Q$
jest homeomorficzny z $\omega^\omega$,
co opisaliśmy przy ułamkach łańcuchowych
w dowodzie Twierdzenia~\ref{thm:ad-baire-property}.
Jeśli $C\subseteq\mathbb R$ jest
nieprzeliczalny, to
$C\setminus\mathbb Q$ także
jest nieprzeliczalny:
w przeciwnym razie $C$ byłby
sumą dwóch przeliczalnych zbiorów,
co w ZF nie wymaga nieskończonego
wyboru.
Przenosząc przez homeomorfizm
otrzymany zwarty zbiór doskonały,
uzyskujemy niepusty doskonały
podzbiór $C$.
\end{proof}

\begin{przyklad}[Zbiór Vitaliego nie ma własności Baire'a]
\label{ex:ad-vitali}
Przypomnijmy relację $x\sim y$ wtedy i tylko wtedy, gdy
$x-y\in\Q$. Gdyby zbiór $V$ zawierał dokładnie jednego
reprezentanta każdej klasy, to przesunięcia $V+q$
($q\in\Q$) byłyby parami rozłączne i ich suma dawałaby
całą prostą. Jeśli $V$ byłby pierwszej kategorii, także
$\R=\bigcup_{q\in\Q}(V+q)$ byłaby pierwszej kategorii,
co przeczy twierdzeniu Baire'a. Jeśli $V$ miałby własność
Baire'a i nie był pierwszej kategorii, byłby rezydualny
na pewnym niepustym przedziale $I$. Dla dostatecznie
małego niezerowego $q\in\Q$ przedziały $I$ i $I-q$
przecinają się. W ich przecięciu zbiory $V$ i $V+q$
byłyby jednocześnie rezydualne, więc miałyby wspólny
punkt, sprzecznie z jednoznacznością reprezentantów.
Dlatego w $\mathrm{ZF}+\mathrm{AD}$ taki $V$ nie istnieje.
Planowana dalsza część skryptu poświęcona
ilorazowi $\R/\Q$ omawia silniejsze
ograniczenia wyboru. Nie traktujemy tu
takiej konsekwencji jako udowodnionej
bez osobnego argumentu.
\end{przyklad}

\subsection{Determinacja projektywna a pełne AD}
\begin{definicja}[Zbiory projektywne i PD]
\label{def:ad-projective}
Zbiorami borelowskimi w $\R^n$ są elementy najmniejszej
$\sigma$-algebry zawierającej zbiory otwarte.
Zbiór $A\subseteq\R^n$ jest \emph{analityczny}, jeśli
jest obrazem zbioru borelowskiego w $\R^{n+1}$ przez
rzut na pierwszych $n$ współrzędnych; po przekształceniach
standardowych równoważnie można brać rzuty zbiorów
domkniętych w $\R^n\times\N^\N$.
\emph{Zbiory projektywne} otrzymujemy ze zbiorów
borelowskich przez skończenie wiele operacji dopełnienia
i rzutowania, a więc również przez skończenie wiele
kwantyfikacji po liczbach rzeczywistych.
Aksjomat \emph{determinacji projektywnej} $\mathrm{PD}$
głosi, że każda gra $G(A)$, której zbiór wypłat
$A\subseteq\N^\N$ jest projektywny, jest zdeterminowana.
\end{definicja}
\begin{uwaga}[Granica udowodnionych rezultatów]
\label{rem:ad-deep-regularity}
Klasyczne twierdzenia teorii determinacji głoszą, że
w $\mathrm{ZF}+\mathrm{AD}$ każdy zbiór liczb rzeczywistych
jest mierzalny w sensie Lebesgue'a i ma własność zbioru
doskonałego: jest przeliczalny albo zawiera niepusty
zbiór doskonały. Pod $\mathrm{PD}$ analogiczne własności
przysługują wszystkim \emph{projektywnym} podzbiorom prostej rzeczywistej.
Dowód własności Baire'a przy AD podaliśmy powyżej.
Dowód własności zbioru doskonałego przy AD
podaliśmy w Twierdzeniu~\ref{thm:ad-perfect-set}
w bieżącym rozdziale.
Wyprowadzenie mierzalności z AD znajdziemy dalej w
Twierdzeniu~\ref{thm:ad-measurability}; używa ono
jawnie wydzielonego twierdzenia o uniwersalnej
mierzalności zbiorów analitycznych. Dowodu tego
ostatniego twierdzenia nie powielamy.
Nie wolno też przenosić następstw AD na wszystkie
podzbiory zbioru liczb rzeczywistych przy założeniach
$\mathrm{ZFC}+\mathrm{PD}$:
PD obejmuje tylko zbiory projektywne.
\end{uwaga}


\section{Determinacja a niemożliwość dobrego uporządkowania zbioru liczb rzeczywistych}
\label{sec:ad-no-wellorder-reals}
\begin{twierdzenie}[ZF+AD]
\label{thm:ad-no-wellorder-reals}
Zakładając ZF i aksjomat determinacji AD, zbiór liczb
rzeczywistych $\R$ nie może być dobrze uporządkowany.
\end{twierdzenie}
\begin{proof}
Z wcześniejszego twierdzenia o własności Baire'a pod AD
wynika, że każdy podzbiór $\R$ ma własność Baire'a.
Przypuśćmy, że istnieje dobre uporządkowanie $\R$.
Wówczas istnieje funkcja wybierająca po jednym elemencie
z każdej klasy relacji $x\sim y\iff x-y\in\Q$:
wybieramy element najmniejszy w danym dobrym porządku.
Niech $V\subseteq\R$ będzie zbiorem wybranych reprezentantów.
Pokażemy, że $V$ nie ma własności Baire'a.

Jeśli $V$ jest zbiorem pierwszej kategorii, to również
każdy jego wymierny przekład $V+q$ jest pierwszej kategorii.
Ponieważ $\R=\bigcup_{q\in\Q}(V+q)$, otrzymujemy,
że $\R$ jest zbiorem pierwszej kategorii, wbrew twierdzeniu
Baire'a dla zupełnej przestrzeni metrycznej $\R$.

Jeśli $V$ nie jest zbiorem pierwszej kategorii i ma własność
Baire'a, to istnieje niepusty przedział otwarty $I$,
w którym $V$ jest współpierwszej kategorii: $I\setminus V$
jest pierwszej kategorii. Wybierzmy niezerowe $q\in\Q$
tak małe, aby $I\cap(I+q)$ było niepuste. Zarówno
$V$, jak i $V+q$ są współpierwszej kategorii w tym
przecięciu, więc ich przecięcie jest niepuste, gdyż
niepusty przedział nie jest zbiorem pierwszej kategorii.
Jest to sprzeczne z konstrukcją $V$: dwa różne elementy
jednej klasy modulo $\Q$ nie mogą należeć do $V$.
Obie możliwości prowadzą do sprzeczności.
\end{proof}
\begin{uwaga}
Twierdzenie nie mówi, że nie można dobrze uporządkować
żadnego podzbioru $\R$. Na przykład każdy skończony
podzbiór $\R$ ma dobre uporządkowanie w ZF.
Wynik dotyczy całego zbioru liczb rzeczywistych.
\end{uwaga}


\section{Mierzalność w sensie Lebesgue'a z aksjomatu determinacji}
\label{sec:ad-lebesgue-measurability}
Wcześniej udowodniliśmy, że aksjomat determinacji
zapewnia własność Baire'a wszystkich podzbiorów
zbioru liczb rzeczywistych. To nie jest jednak
synonim mierzalności w sensie Lebesgue'a: zbiór
pierwszej kategorii może mieć dodatnią miarę,
a zbiór miary zero może być gęsty.
Poniżej przeprowadzimy niezależny argument
z użyciem \emph{gry pokryciowej}. Konstrukcję miary Lebesgue\'a,
regularność miary i kryterium Carathéodory\'ego przywołujemy
z odrębnego skryptu \emph{Teoria miary}; istotnym nowym
argumentem w tym rozdziale jest zastosowanie determinacji
do odpowiedniej gry.

\begin{definicja}[Miara zewnętrzna Lebesgue'a i mierzalność]
\label{def:ad-lebesgue}
Dla $A\subseteq\R$ przez $\lambda^*(A)$ oznaczamy
infimum liczb $\sum_{n\in\omega}(b_n-a_n)$
po wszystkich przeliczalnych pokryciach
$A\subseteq\bigcup_{n\in\omega}(a_n,b_n)$
przedziałami otwartymi o skończonych końcach.
Zbiór $E\subseteq\R$ nazywamy \emph{mierzalnym
w sensie Lebesgue'a}, jeżeli dla każdego
$T\subseteq\R$ zachodzi warunek Carathéodory'ego
\[
 \lambda^*(T)=\lambda^*(T\cap E)+\lambda^*(T\setminus E).
\]
Miara Lebesgue'a $\lambda(E)$ zbioru mierzalnego
jest równa $\lambda^*(E)$.
Zbiór $E$ nazywamy \emph{zbiorem miary zero},
jeżeli $\lambda^*(E)=0$. Definicja nie wymaga
uprzedniego założenia, że $E$ jest mierzalny:
każdy zbiór miary zewnętrznej zero jest mierzalny.
\end{definicja}

\begin{uwaga}[Dokładne odsyłacze do skryptu \emph{Teoria miary}]
\label{rem:ad-measure-theory-links}
W niniejszym rozdziale nie rekonstruujemy od podstaw
konstrukcji miary Lebesgue'a ani jej teorii regularności.
Dowody odpowiednich faktów znajdują się w skrypcie
\emph{Teoria miary} z Biblioteki matematycznej, w
następujących miejscach:
\begin{itemize}
\item rozdział 3, \emph{Miary zewnętrzne i konstrukcja
Carathéodory'ego}: kryterium mierzalności Carathéodory'ego
(etykieta źródłowa \nolinkurl{thm:caratheodory});
\item rozdział 4, \emph{Miara Lebesgue'a na przestrzeni
euklidesowej}: następujące wyniki, ze wskazaniem etykiet źródłowych:
\begin{itemize}
\item istnienie miary — \nolinkurl{thm:lebesgue-exists};
\item regularność zewnętrzna — \nolinkurl{thm:lebesgue-outer-regular};
\item regularność wewnętrzna — \nolinkurl{thm:lebesgue-inner-regular};
\item zupełność i aproksymacja borelowska — \nolinkurl{thm:lebesgue-completion}.
\end{itemize}
\end{itemize}
Odsyłacze identyfikują rozdziały i etykiety źródłowe,
lecz nie są poleceniami \texttt{\string\ref}:
kompilacja tego skryptu nie wymaga plików pomocniczych
z odrębnej książki. W szczególności regularność
zewnętrzna jest przywoływanym wynikiem teorii miary,
a nie twierdzeniem udowadnianym ponownie w tym rozdziale.
\end{uwaga}

\begin{przyklad}[Zbiór wymiernych liczb jest miary zero]
Zbiór $\Q$ jest przeliczalny. Ustalmy enumerację
$\Q=\{q_n:n\in\omega\}$. Dla $\varepsilon>0$
przedział
$J_n=(q_n-\varepsilon 2^{-n-3},
q_n+\varepsilon 2^{-n-3})$
ma długość $\varepsilon 2^{-n-2}$.
Suma długości wszystkich $J_n$ wynosi
$\varepsilon/2$, a przedziały pokrywają $\Q$.
Dlatego $\lambda^*(\Q)\le\varepsilon/2$
dla każdego $\varepsilon>0$, więc
$\lambda^*(\Q)=0$. Mimo to $\Q$ jest gęsty w $\R$.
Gęstość i dodatnia miara są więc różnymi własnościami.
\end{przyklad}

\begin{lemat}[Kodowanie pokryć zbioru miary zero]
\label{lem:ad-null-cover-coding}
Pracujemy w $\mathrm{ZF}+\mathrm{DC}$.
Niech $Z\subseteq[0,1]$ będzie zbiorem miary zero.
Dla $\varepsilon>0$ i $n\in\omega$ niech
$\mathcal K_n$ oznacza zbiór wszystkich skończonych
sum przedziałów otwartych o końcach wymiernych,
których łączna długość jest mniejsza niż
$\varepsilon\,4^{-(n+1)}$ (dopuszczamy sumę pustą).
Wtedy istnieją zbiory $H_n\in\mathcal K_n$ takie, że
$Z\subseteq\bigcup_{n\in\omega}H_n$.
Każdy $\mathcal K_n$ jest przeliczalny i posiada
ustaloną enumerację bez potrzeby stosowania
aksjomatu wyboru do nieprzeliczalnej rodziny.
\end{lemat}
\begin{proof}
Z definicji miary zewnętrznej zero możemy wybrać
przeliczalną rodzinę przedziałów otwartych
$(I_j)_{j\in\omega}$ o wymiernych końcach,
pokrywającą $Z$ i spełniającą
$\sum_j |I_j|<\varepsilon/6$.
W tym miejscu $|I_j|$ oznacza \emph{długość przedziału},
a nie moc zbioru; użycie tego samego symbolu
jest lokalnie objaśnione. Powiększenie końców do
liczb wymiernych jest możliwe przy zachowaniu
nierówności, gdyż $\Q$ jest gęsty w $\R$.

Podzielimy ciąg przedziałów na kolejne skończone
bloki. Na początku suma długości wszystkich
przedziałów jest mniejsza od $\varepsilon/4$.
Wybierzmy $N_0$ tak duże, żeby suma długości
ogona o indeksach $j>N_0$ była mniejsza niż
$\varepsilon/16$. Niech $H_0$ będzie sumą
$I_0,\dots,I_{N_0}$. Ponieważ długość sumy
jest nie większa niż suma długości jej składników,
mamy $H_0\in\mathcal K_0$.

Przypuśćmy, że wybraliśmy $N_{n-1}$ tak, aby
suma długości przedziałów $I_j$ z indeksami
$j>N_{n-1}$ była mniejsza niż
$\varepsilon 4^{-(n+1)}$.
Wybieramy $N_n>N_{n-1}$ tak duże, aby
suma długości ogona po $N_n$ była mniejsza niż
$\varepsilon 4^{-(n+2)}$.
Zbiór $H_n=\bigcup_{N_{n-1}<j\le N_n}I_j$
należy do $\mathcal K_n$. Indeksy $N_n$ mogą być
wybierane jako \emph{najmniejsze} indeksy
spełniające wymaganą nierówność, toteż cały
proces jest zdefiniowany przez rekurencję na
$\omega$; nie wymaga dodatkowej funkcji wyboru.
Ponieważ $N_n\to\infty$, wszystkie przedziały
$I_j$ należą do pewnego $H_n$, więc zbiory
$H_n$ pokrywają $Z$.
\end{proof}

\begin{uwaga}[Jeden klasyczny wynik z deskryptywnej teorii mnogości]
\label{rem:ad-analytic-input}
W poniższym dowodzie stosujemy \emph{twierdzenie
o uniwersalnej mierzalności zbiorów analitycznych}:
każdy ciągły obraz przestrzeni Baire'a $\omega^\omega$
w prostej rzeczywistej jest mierzalny w sensie
Lebesgue'a. Jest to samodzielne, nietrywialne
twierdzenie deskryptywnej teorii mnogości,
którego pełen dowód przez operację Suslina
wymaga osobnego opracowania. Nie uznajemy go
za udowodnione w niniejszym skrypcie, ale nie
przypisujemy też jego dowodu książce \emph{Teoria miary}:
w aktualnej wersji tej książki są dowody regularności
miary Lebesgue'a, natomiast nie ma osobnego dowodu
twierdzenia o uniwersalnej mierzalności zbiorów analitycznych.
W literaturze jest to również łączone z twierdzeniem
Łuzina o zbiorach analitycznych; nie należy mylić go
z twierdzeniem Łuzina o aproksymacji funkcji mierzalnych.
Jako szczegółowy tekst źródłowy można wykorzystać
D.~L. Cohna, \emph{Measure Theory}, rozdział 8,
część dotyczącą przestrzeni polskich i zbiorów analitycznych.
To zewnętrzne twierdzenie jest wyraźnie oddzielone od
argumentów teorii mnogości prowadzonych poniżej.
Następny dowód jest \emph{kompletnym wyprowadzeniem
następstwa AD z tego wyraźnie wskazanego twierdzenia}.
\end{uwaga}

\begin{lemat}[Lemat gry pokryciowej]
\label{lem:ad-covering-game}
Pracujemy w $\mathrm{ZF}+\mathrm{DC}+\mathrm{AD}$
i przyjmujemy twierdzenie o mierzalności zbiorów
analitycznych z Uwagi~\ref{rem:ad-analytic-input}.
Niech $S\subseteq[0,1]$ będzie zbiorem takim, że
każdy mierzalny w sensie Lebesgue'a podzbiór
$Z\subseteq S$ ma miarę zero. Wtedy
$\lambda^*(S)=0$.
\end{lemat}
\begin{proof}
\emph{Etap 1: ustalamy grę.}
Niech $\varepsilon>0$. Dla każdego $n$ ustalmy
enumerację $\mathcal K_n=\{K_{n,k}:k\in\omega\}$
z Lematu~\ref{lem:ad-null-cover-coding}.
Gracz I na swoim $n$-tym ruchu wybiera cyfrę
$a_n\in\{0,1\}$, a gracz II wybiera
liczbę naturalną $b_n$. Z cyfr gracza I
powstaje liczba rzeczywista
\[
 x(a)=\sum_{n\in\omega}\frac{a_n}{2^{n+1}}
 \in[0,1].
\]
Gracz I wygrywa dokładnie wtedy, gdy
$x(a)\in S$ oraz
$x(a)\notin\bigcup_{n\in\omega}K_{n,b_n}$.
W przeciwnym razie wygrywa gracz II.
Jest to gra z ruchami w $\omega$ i jasno
określonym zbiorem wypłat, więc AD zapewnia
strategię wygrywającą jednemu z graczy.

\emph{Etap 2: gracz I nie ma strategii wygrywającej.}
Przypuśćmy przeciwnie, że $\sigma$ jest taką
strategią. Po ustaleniu ciągu wszystkich ruchów
$b=(b_0,b_1,\ldots)\in\omega^\omega$ gracza II,
strategia $\sigma$ wyznacza jednoznacznie
wszystkie cyfry gracza I, a więc liczbę
$f_\sigma(b)=x(a)$.
Funkcja $f_\sigma:\omega^\omega\to[0,1]$
jest ciągła: pierwszych $m$ cyfr jej wartości
zależy tylko od pierwszych $m$ ruchów przeciwnika,
a przybliżenie liczby rzeczywistej przez pierwszych
$m$ cyfr binarnych ma błąd najwyżej $2^{-m}$.
Zatem jej obraz $Z=f_\sigma[\omega^\omega]$
jest zbiorem analitycznym, a więc mierzalnym
na mocy wskazanego twierdzenia o mierzalności zbiorów analitycznych.
Ponieważ $\sigma$ wygrywa przeciw każdemu $b$,
każda wartość $f_\sigma(b)$ należy do $S$.
Dlatego $Z\subseteq S$ i z hipotezy
$\lambda(Z)=0$.

Lemat~\ref{lem:ad-null-cover-coding} dostarcza
ciągu $(H_n)$ z $H_n\in\mathcal K_n$
pokrywającego cały $Z$. Gracz II może więc
ustalić ruchy $b_n$ tak, aby $K_{n,b_n}=H_n$.
Strategia $\sigma$ w odpowiedzi wyprodukuje
liczbę rzeczywistą $f_\sigma(b)\in Z$, która
jest pokryta przez pewien $H_n$. Taka rozgrywka
jest przegrana przez gracza I, sprzecznie z
wygrywającym charakterem $\sigma$.

\emph{Etap 3: gracz II wygrywa i pokrywa $S$.}
Z AD wynika, że gracz II ma strategię wygrywającą
$\tau$. Dla każdego skończonego ciągu binarnego
$s=(a_0,\ldots,a_n)$ oznaczmy przez $H_s$
zbiór $K_{n,b_n}$ wybrany przez $\tau$ po ruchach
$(a_0,\ldots,a_n)$ gracza I i odpowiadających
im wcześniejszych ruchach gracza II.
Wówczas
\[
 \lambda(H_s)<\varepsilon 4^{-(n+1)}.
\]
Weźmy dowolną liczbę rzeczywistą $x\in S$.
Wybierzmy jej rozwinięcie binarne $a$,
przy liczbach dyadycznych mniejszych od $1$ wybierając
rozwinięcie kończące się zerami, a dla $x=1$ wybierając
ciąg samych jedynek. Jeśli gracz I podaje
kolejne cyfry tego rozwinięcia, gracz II
postępujący według $\tau$ musi wygrać.
Ponieważ $x\in S$, oznacza to, że
$x\in H_{a\upharpoonright(n+1)}$ dla pewnego $n$.
W konsekwencji
\[
 S\subseteq\bigcup_{n\in\omega}\ \bigcup_{s\in\{0,1\}^{n+1}}H_s.
\]

\emph{Etap 4: obliczamy miarę pokrycia.}
Dla ustalonego $n$ istnieje dokładnie
$2^{n+1}$ ciągów binarnych długości $n+1$.
Podaddytywność miary zewnętrznej daje
\[
\lambda^*(S)
\le\sum_{n\in\omega}2^{n+1}\varepsilon 4^{-(n+1)}
=\varepsilon\sum_{n\in\omega}2^{-(n+1)}
=\varepsilon.
\]
Ponieważ $\varepsilon>0$ było dowolne,
$\lambda^*(S)=0$.
\end{proof}

\begin{twierdzenie}[AD implikuje mierzalność zbiorów liczb rzeczywistych]
\label{thm:ad-measurability}
W teorii $\mathrm{ZF}+\mathrm{DC}+\mathrm{AD}$,
korzystając z klasycznego twierdzenia o mierzalności
zbiorów analitycznych, dla każdego
$A\subseteq\R$ otrzymujemy, że $A$ jest
mierzalny w sensie Lebesgue'a.
\end{twierdzenie}
\begin{proof}
\emph{Etap 1: ograniczamy się do $[0,1]$.}
Niech $X\subseteq[0,1]$. Z regularności zewnętrznej miary Lebesgue\'a
(zob. Uwaga~\ref{rem:ad-measure-theory-links})
istnieją zbiory otwarte
$U_n\supseteq X$ o miarach
$\lambda(U_n)\le\lambda^*(X)+1/(n+1)$.
Zasada wyboru zależnego pozwala ustalić taki
ciąg pokryć. Niech $H=\bigcap_n U_n$.
Jest to zbiór borelowski zawierający $X$.
Monotoniczność miary zewnętrznej i nierówność
$\lambda(H)\le\lambda(U_n)$ dają
\[
 \lambda(H)=\lambda^*(X).
\]

\emph{Etap 2: dopełnienie w obwiedni ma miarę zero.}
Niech $S=H\setminus X$. Jeżeli $Z\subseteq S$
jest mierzalny, to $H\setminus Z$ jest mierzalnym
zbiorem zawierającym $X$. Gdyby $\lambda(Z)>0$,
to ponieważ $\lambda(H)\le1$, otrzymalibyśmy
\[
\lambda(H\setminus Z)=\lambda(H)-\lambda(Z)
<\lambda(H)=\lambda^*(X),
\]
sprzecznie z definicją miary zewnętrznej.
Zatem każdy mierzalny podzbiór $S$ ma miarę zero.
Z Lematu~\ref{lem:ad-covering-game} wynika
$\lambda^*(S)=0$.

\emph{Etap 3: wyprowadzamy mierzalność.}
Zbiory miary zewnętrznej zero są mierzalne,
a miara Lebesgue'a jest zupełna
(zob. Uwaga~\ref{rem:ad-measure-theory-links}). Dlatego
$X=H\setminus S$ jest mierzalny.
Dla dowolnego $A\subseteq\R$ stosujemy
argument do przekształconych liniowo zbiorów
$A\cap[n,n+1]$, $n\in\Z$. Każdy taki fragment
jest mierzalny, więc ich przeliczalna suma,
równa $A$, również jest mierzalna.
\end{proof}

\section{Co naprawdę wynika z determinacji projektywnej?}
\label{sec:pd-regularity}
\begin{twierdzenie}[Regularność zbiorów projektywnych z PD]
\label{thm:pd-regularity}
Przy założeniu $\mathrm{ZFC}+\mathrm{PD}$ każdy
projektywny podzbiór zbioru liczb rzeczywistych
jest mierzalny w sensie Lebesgue'a, ma własność
Baire'a i jest przeliczalny albo zawiera
niepusty zbiór doskonały.
W dowodzie mierzalności korzystamy ze wskazanego
klasycznego twierdzenia o zbiorach analitycznych.
\end{twierdzenie}
\begin{proof}
W dowodach mierzalności
(Twierdzenie~\ref{thm:ad-measurability}),
własności Baire'a
(Twierdzenie~\ref{thm:ad-baire-property})
i własności zbioru doskonałego
(Twierdzenie~\ref{thm:ad-perfect-set})
aksjomat AD był stosowany do gier zbudowanych
ze zbioru wypłat $A$ przez ustalone kodowania,
operacje borelowskie i przeciwobrazy przez
funkcje ciągłe. Jeśli $A$ jest projektywny,
tak otrzymane zbiory wypłat również są projektywne:
klasy projektywne są zamknięte na dopełnienie,
rzutowanie, operacje borelowskie i ciągłe
przeciwobrazy. Zatem do każdej z tych gier
wystarcza już PD. Pozostałe części dowodów
— konstrukcja pokryć miary zero, wyprowadzenie
sprzeczności ze strategią i budowanie
zbioru doskonałego — nie wymagają determinacji
innych gier. Otrzymujemy trzy własności
regularności dla $A$.
\end{proof}

\begin{uwaga}[Nie mylić PD z pełnym AD]
Z $\mathrm{ZFC}+\mathrm{PD}$ nie wynika, że
\emph{każdy} podzbiór zbioru liczb rzeczywistych
ma własność Baire'a lub jest mierzalny.
Pełny aksjomat wyboru nadal pozwala skonstruować
zbiory Vitaliego; ich istnienie nie przeczy PD,
ponieważ nie są zbiorami projektywnymi.
W szczególności nie wolno zastąpić w tezie
poprzedniego twierdzenia określenia
„każdy zbiór projektywny” przez
„każdy podzbiór zbioru liczb rzeczywistych”.
\end{uwaga}


\section{Determinacja projektywna a duże liczby kardynalne}
\label{sec:ad-pd-large-cardinal-strength}
W tej sekcji odróżniamy trzy rodzaje twierdzeń:
wyniki wynikające bezpośrednio z podanych
wcześniej definicji, głębokie twierdzenia o
modelach wewnętrznych oraz konsekwencje takich
twierdzeń, które możemy dowieść na podstawie
wcześniejszych rozdziałów. Pracujemy w otaczającej
teorii ZFC, chyba że wyraźnie wskazano inną teorię.

\begin{definicja}[Model $L(\mathbb R)$ i determinacja wewnętrzna]
\label{def:lr-ad}
Niech $V$ będzie otaczającym uniwersum zbiorów,
a $\R^V$ jego zbiorem liczb rzeczywistych.
Przez $L(\R)$ oznaczamy najmniejszy model
wewnętrzny ZF zawierający wszystkie liczby
porządkowe i wszystkie liczby rzeczywiste
z $V$, konstruowany przez iterację operacji
\emph{definiowalnych podzbiorów} z parametrem $\R^V$.
Zdanie $L(\R)\models\mathrm{AD}$ oznacza, że
\emph{każda gra o długości $\omega$ na liczbach
naturalnych, której zbiór wypłat jest elementem
$L(\R)$, jest zdeterminowana w tym modelu}.
Nie oznacza to, że $V\models\mathrm{AD}$.
\end{definicja}

\begin{uwaga}[Nie utożsamiać zbioru liczb rzeczywistych z jego kodowaniami]
Dla celów semantyki gry używamy przestrzeni
$\omega^\omega$ wszystkich funkcji $\omega\to\omega$.
Jej elementy są ciągami liczb naturalnych;
zbiór $\R$ ma inną definicję. W ZFC przestrzenie
te mają wzajemnie jednoznaczne kodowania borelowskie,
co pozwala przenosić klasy projektywne oraz
własności determinacji w ich standardowym ujęciu.
\end{uwaga}

\begin{definicja}[Liczba kardynalna Woodina — sformułowanie]
\label{def:woodin-cardinal}
Niech $\delta$ będzie nieprzeliczalną liczbą
kardynalną. Mówimy, że $\delta$ jest
\emph{liczbą kardynalną Woodina}, jeżeli dla
każdej funkcji $f:\delta\to\delta$ istnieje
$\kappa<\delta$ i elementarne zanurzenie
$j:V\to M$ do tranzytywnego modelu klasowego $M$
takie, że $\kappa$ jest punktem krytycznym $j$,
$f[\kappa]\subseteq\kappa$ oraz
$V_{j(f)(\kappa)}\subseteq M$.
Punktem krytycznym elementarnego zanurzenia jest
najmniejsza liczba porządkowa $\alpha$ spełniająca
$j(\alpha)\ne\alpha$. Zapis $f[\kappa]$ oznacza
obraz zbioru $\kappa$ przy funkcji $f$.
Sformułowanie używa semantyki klas i wymaga
precyzyjnej interpretacji w stosownej teorii klas;
jego równoważne warianty za pomocą extenderów
należą do teorii dużych liczb kardynalnych.
\end{definicja}

\begin{uwaga}[Granica samodzielności definicji]
Samo podanie definicji liczby kardynalnej Woodina
nie stanowi dowodu istnienia takiej liczby
kardynalnej ani nie daje determinacji projektywnej.
Nie twierdzimy też, że wystarczy istnienie jednej
liczby kardynalnej Woodina do udowodnienia pełnej PD.
\end{uwaga}

\begin{twierdzenie}[Martin--Steel: determinacja projektywna]
\label{thm:martin-steel-projective}
W teorii ZFC z założenia istnienia nieskończenie
wielu liczb kardynalnych Woodina wynika
aksjomat determinacji projektywnej $\mathrm{PD}$.
Jest to głębokie twierdzenie teorii modeli
wewnętrznych, przywołane tutaj ze wskazaniem
źródła, a nie wyprowadzone z definicji.
\end{twierdzenie}
\begin{uwaga}[Źródło głębokiego twierdzenia]
Opis zależności, ograniczeń i literatury znajduje się
w Stanford Encyclopedia of Philosophy,
\emph{Large Cardinals and Determinacy}:
\url{https://plato.stanford.edu/entries/large-cardinals-determinacy/}.
W książce rozwijamy szczegółowo podstawową teorię
mierzalnych liczb kardynalnych, natomiast pełny
dowód powyższego twierdzenia wymaga maszynerii
iterowalnych modeli wewnętrznych oraz extenderów.
\end{uwaga}

\begin{twierdzenie}[Martin--Steel--Woodin: determinacja w $L(\R)$]
\label{thm:martin-steel-woodin-lr}
Pracujemy w ZFC. Jeżeli istnieje rosnący ciąg
$(\delta_n)_{n\in\omega}$ liczb kardynalnych
Woodina oraz mierzalna liczba kardynalna
$\kappa$ spełniająca $\delta_n<\kappa$ dla
każdego $n\in\omega$, to
\[
 L(\R)\models\mathrm{AD}.
\]
Jest to przywołane twierdzenie głębokie;
nie zastępujemy jego dowodu stwierdzeniem,
że wynika \emph{bezpośrednio} z istnienia
mierzalnej liczby kardynalnej.
\end{twierdzenie}
\begin{uwaga}[Źródło i znaczenie hipotez]
Zob. również rozdział 33 podręcznika Jecha
\emph{Set Theory}, \emph{Determinacy},
\url{https://fa.ewi.tudelft.nl/~hart/onderwijs/set_theory/Jech/33-determinacy.pdf}.
Hipoteza mówi, że \emph{jedna} mierzalna liczba
kardynalna leży powyżej całej wskazanej rodziny
liczb kardynalnych Woodina. Nie wynika to z samej
informacji, że istnieje nieskończenie wiele
liczb kardynalnych Woodina.
\end{uwaga}

\begin{lemat}[Prawdziwość formuł projektywnych w modelach o tych samych liczbach rzeczywistych]
\label{lem:projective-same-reals}
Niech $M\subseteq V$ będzie tranzytywnym modelem
wewnętrznym ZF o tych samych liczbach naturalnych
i rzeczywistych co $V$, czyli
$\omega^M=\omega^V$ i $\R^M=\R^V$.
Niech $A\subseteq\omega^\omega$ będzie zbiorem
projektywnym zadanym formułą z parametrami będącymi
liczbami rzeczywistymi. Wtedy dla każdego
$a\in\omega^\omega$ formuła definiująca
przynależność $a\in A$ ma tę samą wartość logiczną
w $M$ i w $V$.
\end{lemat}
\begin{proof}
Zbiory borelowskie definiowane za pomocą kodów
borelowskich złożonych z liczb naturalnych są
absolutne między modelami o tych samych liczbach
rzeczywistych: relacja przynależności do zbioru
kodowanego jest obliczana przez dobrze ufundowaną
rekurencję po drzewie kodu, które należy do obu
modeli, a zbiory liczb naturalnych są te same.
Prowadzimy następnie indukcję po liczbie operacji
projektywnych. Dopełnienie zmienia prawdę
na fałsz po obu stronach, a projekcja jest
kwantyfikacją po elementach $\omega^\omega$.
Modele mają dokładnie te same takie elementy,
więc kwantyfikator egzystencjalny daje tę samą
wartość logiczną. W ten sposób wszystkie
skończone etapy definicji projektywnej są
absolutne. Nie używamy tutaj aksjomatu wyboru
\emph{wewnątrz} modelu $M$.
\end{proof}

\begin{wniosek}[Determinacja w $L(\R)$ implikuje PD w otaczającym uniwersum]
\label{cor:lr-ad-implies-pd}
Jeżeli $V\models\mathrm{ZFC}$ oraz
$L(\R)\models\mathrm{AD}$, to $V\models\mathrm{PD}$.
\end{wniosek}
\begin{proof}
Niech $A\subseteq\omega^\omega$ będzie zbiorem
projektywnym z parametrami będącymi liczbami
rzeczywistymi z $V$. Ponieważ $L(\R)$ zawiera
wszystkie liczby rzeczywiste i porządkowe,
zawiera także kod definiujący $A$, a dzięki
wyróżnianiu oraz projektywnej definicji zawiera
sam zbiór $A$.
Z założenia $L(\R)\models\mathrm{AD}$
w $L(\R)$ istnieje strategia wygrywająca jednego
z dwóch graczy w grze o wypłacie $A$. Strategia
jest funkcją na zbiorze wszystkich skończonych
ciągów liczb naturalnych; można ją zakodować
jako liczbę rzeczywistą, toteż ta sama strategia
należy także do $V$. Z Lematu~\ref{lem:projective-same-reals}
wynika, że $A$ ma identyczne elementy w obu
modelach. Dlatego każda rozgrywka zgodna ze
strategią kończy się zwycięstwem tego samego
gracza także w $V$. Gra jest zdeterminowana,
a ponieważ $A$ był dowolny, zachodzi PD.
\end{proof}

\begin{wniosek}[Regularność projektywnych podzbiorów liczb rzeczywistych]
\label{cor:woodin-projective-regularity}
Przy hipotezach Twierdzenia~\ref{thm:martin-steel-woodin-lr}
w otaczającym uniwersum ZFC każdy projektywny
podzbiór zbioru liczb rzeczywistych jest
mierzalny w sensie Lebesgue'a, ma własność
Baire'a i jest przeliczalny albo zawiera
niepusty zbiór doskonały.
\end{wniosek}
\begin{proof}
Twierdzenie~\ref{thm:martin-steel-woodin-lr}
daje $L(\R)\models\mathrm{AD}$. Na mocy
Wniosku~\ref{cor:lr-ad-implies-pd} zachodzi
PD w otaczającym uniwersum ZFC.
Twierdzenie~\ref{thm:pd-regularity} zastosowane
do dowolnego projektywnego podzbioru zbioru
liczb rzeczywistych daje wszystkie trzy własności.
\end{proof}

\begin{uwaga}[Dlaczego nie otrzymujemy pełnego AD w ZFC?]
Z aksjomatu wyboru w $V$ wynika istnienie
gier niezdeterminowanych
(Twierdzenie~\ref{thm:ac-not-ad}).
Dlatego w powyższym twierdzeniu celowo
piszemy $L(\R)\models\mathrm{AD}$, a nie
$V\models\mathrm{AD}$. To pierwsze zdanie jest
zgodne z aksjomatem wyboru w otaczającym
uniwersum, drugie — nie.
\end{uwaga}


\section{Redukcje Wadge'a i dalsza konsekwencja determinacji projektywnej}
\label{sec:pd-wadge-reductions}
W tej sekcji pracujemy w ZFC, a dodatkowo przyjmujemy
aksjomat determinacji projektywnej PD tylko tam, gdzie
został wyraźnie wymieniony. Nie używamy tu twierdzeń
o mierzalności zbiorów liczb rzeczywistych.

\begin{definicja}[Redukowalność Wadge'a]
\label{def:wadge-reducibility}
Niech $A,B$ będą podzbiorami przestrzeni Baire'a
$\omega^\omega$, czyli zbioru wszystkich ciągów liczb
naturalnych z topologią generowaną przez cylindry
$[s]$ z Definicji~\ref{def:baire-cylinders}.
Piszemy $A\leq_{\mathrm W}B$, jeżeli istnieje
funkcja ciągła $f:\omega^\omega\to\omega^\omega$
spełniająca
\[
\forall x\in\omega^\omega\qquad
   x\in A\ \Longleftrightarrow\ f(x)\in B.
\]
Oznacza to równość $A=f^{-1}[B]$.
Nie wymagamy, aby $f$ była iniekcją, surjekcją
ani bijekcją. Relacja $\leq_{\mathrm W}$ jest
preporządkiem: jest zwrotna (funkcja identycznościowa)
i przechodnia (złożenie funkcji ciągłych).
\end{definicja}

\begin{lemat}[Zbiory projektywne a kodowanie gier Wadge'a]
\label{lem:wadge-payoff-projective}
Niech $A,B\subseteq\omega^\omega$ będą zbiorami
projektywnymi, w sensie Definicji~\ref{def:ad-projective}.
W grze, w której gracz I podaje po jednej liczbie
naturalnej, a gracz II w każdej turze podaje liczbę
naturalną albo symbol pominięcia ruchu, zbiór
rozgrywek wygrywanych przez gracza II według
następujących reguł jest projektywny:
\begin{enumerate}
\item gracz II musi podać nieskończenie wiele liczb,
      aby utworzyć ciąg $y\in\omega^\omega$;
\item po oznaczeniu przez $x\in\omega^\omega$ ciągu
      liczb gracza I wymagamy równoważności
      $x\in A\iff y\in B$.
\end{enumerate}
\end{lemat}
\begin{proof}
Kodujemy jeden ruch parą liczb naturalnych,
przy czym drugi składnik jest albo symbolem pominięcia,
albo liczbą naturalną. Przy ustalonym bijektywnym
kodowaniu par całe rozgrywki tworzą przestrzeń
$\omega^\omega$. Ciąg $x$ zależy ciągle od rozgrywki.
Warunek, że II podał nieskończenie wiele liczb,
jest warunkiem borelowskim (dla każdego $n$ istnieje
późniejszy ruch niebędący pominięciem). Na zbiorze
rozgrywek spełniających ten warunek ciąg $y$ zależy
ciągle od całej rozgrywki: pierwsze $k$ jego wyrazów
są wyznaczone przez pewien skończony początkowy
fragment rozgrywki, którego długość może zależeć
od rozgrywki i od $k$. Przeciwobrazy zbiorów
projektywnych przez takie kodowania pozostają
projektywne; własność ta wynika przez indukcję
z definicji klas projektywnych, ponieważ przeciwobrazy
zbiorów borelowskich przez funkcje ciągłe są borelowskie,
a rzutowanie i dopełnienie zachowują hierarchię.
Równoważność dwóch zdań zapisujemy przez skończone
sumy, przecięcia i dopełnienia. Po uwzględnieniu
borelowskiego warunku legalności gry otrzymujemy
zbiór projektywny.
\end{proof}

\begin{twierdzenie}[Półliniowość porządku Wadge'a dla zbiorów projektywnych]
\label{thm:pd-projective-wadge}
Zakładając ZFC i PD, dla dowolnych dwóch
projektywnych zbiorów $A,B\subseteq\omega^\omega$
zachodzi co najmniej jedna z możliwości
\[
 A\leq_{\mathrm W}B
 \qquad\text{lub}\qquad
 B\leq_{\mathrm W}(\omega^\omega\setminus A).
\]
W szczególności nie chodzi o zwykłą porównywalność
bez dopełnienia, ani o porównanie mocy zbiorów.
\end{twierdzenie}
\begin{proof}
Rozważmy grę opisaną w
Lemacie~\ref{lem:wadge-payoff-projective}.
W każdej turze I podaje liczbę naturalną, a II
może podać liczbę naturalną albo pominąć ruch.
Jeśli II poda tylko skończenie wiele liczb,
przegrywa. Jeżeli poda nieskończenie wiele liczb,
to wygrywa dokładnie wtedy, gdy uzyskane ciągi
$x$ i $y$ spełniają $x\in A\iff y\in B$.
Na mocy lematu zbiór wypłat jest projektywny,
zatem PD gwarantuje strategię wygrywającą jednego
z graczy.

\emph{Przypadek pierwszy: zwycięską strategię ma II.}
Dla ustalonego wejścia $x\in\omega^\omega$
pozwalamy I kolejno podawać jego wyrazy,
a II odpowiada zgodnie ze swoją strategią.
Strategia II jest wygrywająca, więc w każdej takiej
rozgrywce podaje nieskończenie wiele liczb.
Oznaczmy utworzony ciąg przez $f(x)\in\omega^\omega$.
Aby wykazać ciągłość $f$ w punkcie $x$, ustalmy
liczbę naturalną $k$. Po pewnej skończonej liczbie
ruchów strategia II wytworzyła pierwszych $k$
wyrazów $f(x)$. Każde wejście o tym samym
wystarczająco długim prefiksie wywołuje te same
odpowiedzi strategii do tego momentu. Zatem
pierwszych $k$ wyrazów $f(x)$ jest lokalnie
stałych, co jest kryterium ciągłości w przestrzeni
Baire'a. Z warunku zwycięstwa II mamy
$x\in A\iff f(x)\in B$ dla każdego $x$.
Otrzymaliśmy $A\leq_{\mathrm W}B$.

\emph{Przypadek drugi: zwycięską strategię ma I.}
Pozwalamy graczowi II podawać w każdej turze
jedną liczbę naturalną i nigdy nie pomijać ruchu.
Dla dowolnego ciągu $y\in\omega^\omega$ jego
kolejne wyrazy stanowią taką legalną rozgrywkę
gracza II. Strategia gracza I wyznacza w odpowiedzi
ciąg $g(y)=x\in\omega^\omega$. Ponieważ każdy
skończony prefiks $g(y)$ zależy jedynie od
skończonej liczby wyrazów $y$, funkcja $g$ jest
ciągła. Gracz I wygrywa przeciw każdej rozgrywce
II, w tym przeciw wybranej, dlatego
\[
 \neg\bigl(g(y)\in A\iff y\in B\bigr).
\]
Powyższy warunek jest równoważny
$y\in B\iff g(y)\notin A$. W konsekwencji
$B=g^{-1}[\omega^\omega\setminus A]$,
czyli $B\leq_{\mathrm W}(\omega^\omega\setminus A)$.
Oba przypadki wyczerpują możliwości.
\end{proof}

\begin{przyklad}[Dlaczego potrzebne jest dopełnienie?]
Niech $z\in\omega^\omega$ będzie ciągiem stale
równym zeru, $A=\{z\}$ oraz
$B=\{y\in\omega^\omega:y(0)=0\}$. Zbiór $A$
jest domknięty, ale nieotwarty, natomiast $B$
jest jednocześnie otwarty i domknięty. Gdyby
$A=f^{-1}[B]$ dla funkcji ciągłej $f$, to
$A$ byłby otwarty, ponieważ przeciwobraz zbioru
otwartego przez funkcję ciągłą jest otwarty.
Tak więc $A\not\leq_{\mathrm W}B$.

Mimo to $B\leq_{\mathrm W}(\omega^\omega\setminus A)$:
niech $u\in\omega^\omega$ będzie ciągiem stale
równym jedności i połóżmy $g(y)=u$, gdy
$y(0)=0$, oraz $g(y)=z$ w przeciwnym razie.
Funkcja $g$ jest ciągła, bo ma stałe wartości
na dwóch rozłącznych zbiorach otwarto-domkniętych,
a dokładnie na $B$ przyjmuje wartość spoza $A$.
Ten przykład ilustruje matematyczne znaczenie
dołączenia dopełnienia w tezie twierdzenia.
\end{przyklad}

\begin{uwaga}[Zakres twierdzenia i literatura]
Twierdzenie~\ref{thm:pd-projective-wadge}
ma samodzielny dowód oparty na PD. Głębsze
rezultaty, takie jak determinacja gier o wyższej
długości, struktura rang Wadge'a i twierdzenia
o skalach projektywnych, wymagają oddzielnych
dowodów. Por. D. Marker, \emph{Descriptive Set Theory},
\url{https://www.math.uic.edu/~marker/math512/dst.pdf},
oraz Stanford Encyclopedia of Philosophy,
\emph{Large Cardinals and Determinacy},
\url{https://plato.stanford.edu/entries/large-cardinals-determinacy/}.
\end{uwaga}


\section{Pierwsze stopnie hierarchii Wadge'a i pełne znaczenie gry redukcyjnej}
\label{sec:wadge-first-levels}
Z poprzedniej sekcji wiemy, że determinacja projektywna
PD daje półliniowość redukowalności Wadge'a dla zbiorów
projektywnych (Twierdzenie~\ref{thm:pd-projective-wadge}).
Teraz zbierzemy konstrukcje, które można w całości
udowodnić w ZFC, bez zakładania PD. Odróżnimy trzy
pojęcia: \emph{redukowalność ciągłą}, \emph{równoważność
Wadge'a} i \emph{przynależność do klasy borelowskiej}.

\begin{definicja}[Równoważność Wadge'a i zupełność dla klasy zbiorów]
\label{def:wadge-equivalence-completeness}
Dla dowolnych zbiorów $A,B\subseteq\omega^\omega$
piszemy $A\equiv_{\mathrm W}B$, jeżeli
$A\leq_{\mathrm W}B$ i $B\leq_{\mathrm W}A$.
Klasa równoważności tej relacji jest nazywana
\emph{stopniem Wadge'a}. Jeśli $\Gamma$ jest
rodziną podzbiorów $\omega^\omega$, mówimy, że
$B\subseteq\omega^\omega$ jest \emph{zupełny
względem redukcji Wadge'a dla $\Gamma$}, jeżeli
$B\in\Gamma$ i każdy $A\in\Gamma$ spełnia
$A\leq_{\mathrm W}B$. Nie chodzi tutaj o
zupełność metryczną ani o moc zbioru.
\end{definicja}

\begin{lemat}[Jakie operacje zachowuje redukcja ciągła?]
\label{lem:wadge-continuous-preimages}
Niech $A,B\subseteq\omega^\omega$ i
$A\leq_{\mathrm W}B$. Jeśli $B$ jest otwarty,
to $A$ jest otwarty; jeśli $B$ jest domknięty,
to $A$ jest domknięty. Jeśli $B$ jest przeliczalną
sumą zbiorów domkniętych, to $A$ również ma tę
postać. Analogicznie przeciwobrazy przeliczalnych
przecięć zbiorów otwartych są przeliczalnymi
przecięciami zbiorów otwartych.
\end{lemat}
\begin{proof}
Istnieje funkcja ciągła $f:\omega^\omega\to\omega^\omega$
taka, że $A=f^{-1}[B]$. Z definicji ciągłości
przeciwobraz zbioru otwartego jest otwarty;
przeciwobraz zbioru domkniętego jest domknięty,
ponieważ $f^{-1}[\omega^\omega\setminus B]
=\omega^\omega\setminus f^{-1}[B]$.
Dla dowolnej rodziny $(B_n)_{n\in\omega}$ mamy
\[
 f^{-1}\!\left[\bigcup_{n\in\omega}B_n\right]
 =\bigcup_{n\in\omega}f^{-1}[B_n],
 \quad
 f^{-1}\!\left[\bigcap_{n\in\omega}B_n\right]
 =\bigcap_{n\in\omega}f^{-1}[B_n].
\]
Te równości wynikają bezpośrednio z definicji
przeciwobrazu i uzasadniają obie pozostałe tezy.
\end{proof}

\begin{twierdzenie}[Zbiór otwarty zupełny dla zbiorów otwartych]
\label{thm:wadge-open-complete}
Niech
\[
 U=\{x\in\omega^\omega:\exists n\in\omega\ (x(n)\ne0)\}.
\]
Zbiór $U$ jest otwarty i nie jest domknięty.
Dla dowolnego zbioru otwartego $A\subseteq\omega^\omega$
istnieje funkcja ciągła $f_A:\omega^\omega\to\omega^\omega$
taka, że $A=f_A^{-1}[U]$.
\end{twierdzenie}
\begin{proof}
Otwartość $U$ wynika z równości
$U=\bigcup_{n\in\omega}\{x:x(n)\ne0\}$;
każdy zbiór w tej sumie jest otwarty i domknięty.
Niech $z$ będzie ciągiem stale równym zeru, a
$x^{(m)}$ ciągiem przyjmującym wartość $1$
dokładnie na pozycji $m$ i wartość $0$
na pozostałych pozycjach. Każdy $x^{(m)}\in U$,
a $x^{(m)}\to z\notin U$, więc $U$ nie jest
domknięty.

Ustalmy otwarty zbiór $A$. Dla każdego ciągu
$x\in\omega^\omega$ i każdego $n\in\omega$ połóżmy
\[
 f_A(x)(n)=
 \begin{cases}
  1,&\text{gdy }[x\upharpoonright(n+1)]\subseteq A,\\
  0,&\text{w przeciwnym przypadku}.
 \end{cases}
\]
Symbol $x\upharpoonright(n+1)$ oznacza skończony
ciąg pierwszych $n+1$ wyrazów $x$, a cylindry
$[s]$ zdefiniowaliśmy w Definicji~\ref{def:baire-cylinders}.
Dla ustalonego $n$ wartość $f_A(x)(n)$ zależy
wyłącznie od tego skończonego ciągu, więc każda
współrzędna funkcji $f_A$ jest ciągła. Oznacza to
ciągłość $f_A$ w topologii produktowej.
Jeśli $x\in A$, otwartość daje pewien prefiks
$s=x\upharpoonright k$ z $[s]\subseteq A$.
Dla wszystkich $n+1\ge k$ mamy
$[x\upharpoonright(n+1)]\subseteq[s]\subseteq A$,
więc $f_A(x)\in U$.
Jeśli $x\notin A$, żaden cylinder zawierający
$x$ nie może być podzbiorem $A$; wówczas
$f_A(x)$ jest ciągiem stale równym zeru,
a więc $f_A(x)\notin U$. Dowiedliśmy
$A=f_A^{-1}[U]$.
\end{proof}

\begin{twierdzenie}[Zbiór domknięty zupełny dla zbiorów domkniętych]
\label{thm:wadge-closed-complete}
Niech $K=\{z\}$, gdzie $z\in\omega^\omega$
jest ciągiem stale równym zeru.
Zbiór $K$ jest domknięty i nie jest otwarty.
Dla każdego zbioru domkniętego $F\subseteq\omega^\omega$
istnieje ciągła funkcja $g_F:\omega^\omega\to\omega^\omega$
spełniająca $F=g_F^{-1}[K]$.
\end{twierdzenie}
\begin{proof}
Zbiór $K$ jest domknięty jako przecięcie
$\bigcap_{n\in\omega}\{x:x(n)=0\}$.
Nie jest otwarty, ponieważ każdy cylinder
zawierający $z$ zawiera także ciągi różne od $z$.
Dla ustalonego zbioru domkniętego $F$ definiujemy
\[
 g_F(x)(n)=
 \begin{cases}
  0,&\text{jeśli }[x\upharpoonright(n+1)]\cap F\ne\varnothing,\\
  1,&\text{jeśli }[x\upharpoonright(n+1)]\cap F=\varnothing.
 \end{cases}
\]
Także tutaj wartość każdej współrzędnej zależy
od skończonego początkowego fragmentu $x$,
zatem $g_F$ jest ciągła. Dla $x\in F$
każdy cylinder zawierający $x$ przecina $F$,
więc $g_F(x)=z\in K$.
Jeśli $x\notin F$, otwartość dopełnienia $F$
gwarantuje istnienie cylindra
$[x\upharpoonright k]$ rozłącznego z $F$.
Wtedy $g_F(x)(n)=1$ dla każdego $n+1\ge k$,
a więc $g_F(x)\notin K$.
Otrzymujemy $F=g_F^{-1}[K]$.
\end{proof}

\begin{wniosek}[Pierwsze rozróżnialne stopnie Wadge'a]
\label{cor:wadge-first-degrees}
W ZFC zachodzą następujące fakty:
\begin{enumerate}
\item wszystkie niepuste, właściwe zbiory \emph{jednocześnie otwarte i domknięte} w $\omega^\omega$ są wzajemnie równoważne Wadge'a;
\item jeśli $U$ i $K$ są zbiorami z poprzednich dwóch twierdzeń, to ani $U\leq_{\mathrm W}K$, ani $K\leq_{\mathrm W}U$;
\item mimo to $U\equiv_{\mathrm W}(\omega^\omega\setminus K)$.
\end{enumerate}
\end{wniosek}
\begin{proof}
Dla pierwszego punktu ustalmy niepuste właściwe zbiory jednocześnie otwarte i domknięte $A,B$.
Wybierzmy $u\in B$ oraz $v\notin B$ i określmy $f(x)=u$ dla $x\in A$, zaś $f(x)=v$ dla $x\notin A$.
Dla dowolnego otwartego $O\subseteq\omega^\omega$ przeciwobraz $f^{-1}[O]$ jest jednym ze zbiorów $\varnothing,A,\omega^\omega\setminus A,\omega^\omega$; każdy z nich jest otwarty. Zatem $f$ jest ciągła i $f^{-1}[B]=A$.
Zamiana ról $A,B$ daje równoważność.

Zbiór $U$ jest otwarty, ale nie jest domknięty; zbiór $K$ jest domknięty, ale nie jest otwarty (co wykazano w poprzednich twierdzeniach).
Gdyby $U\leq_{\mathrm W}K$, ciągłość przeciwobrazu wymuszałaby domkniętość $U$; analogicznie $K\leq_{\mathrm W}U$ wymuszałoby otwartość $K$.
Wreszcie $\omega^\omega\setminus K$ jest niepustym właściwym zbiorem otwartym. Z zupełności $U$ dla zbiorów otwartych wynika $\omega^\omega\setminus K\leq_{\mathrm W}U$.
Przeciwną redukcję daje zupełność $\omega^\omega\setminus K$: jest to zbiór otwarty zupełny, ponieważ $K$ jest domknięty zupełny, a dopełnianie zachowuje równanie $f^{-1}[B]=A$.
\end{proof}

\begin{twierdzenie}[Strategie gracza II i funkcje ciągłe]
\label{thm:wadge-strategy-continuity}
Niech $A,B\subseteq\omega^\omega$. W grze Wadge'a
z Lematu~\ref{lem:wadge-payoff-projective},
w której gracz I podaje kolejne wyrazy ciągu $x$,
a gracz II podaje kolejne wyrazy ciągu $y$
lub pomija ruch, przy czym musi ostatecznie podać
nieskończenie wiele wyrazów, zachodzi równoważność:
\[
 A\leq_{\mathrm W}B
 \quad\Longleftrightarrow\quad
 \text{gracz II ma strategię wygrywającą w tej grze}.
\]
Ten wynik nie wymaga założenia determinacji
wszystkich gier.
\end{twierdzenie}
\begin{proof}
\emph{Strategia daje funkcję ciągłą.}
To przejście zostało wykazane w pierwszym przypadku
dowodu Twierdzenia~\ref{thm:pd-projective-wadge}:
dla każdego wejścia $x$ rozgrywamy strategię
gracza II przeciwko kolejnym wyrazom $x$,
a ciąg wyrazów podanych przez II jest wartością
ciągłej funkcji $f(x)$. Warunek zwycięstwa
zapewnia $A=f^{-1}[B]$. Dowód ciągłości
nie wykorzystuje PD, lecz tylko określoną
strategię wygrywającą.

\emph{Funkcja ciągła daje strategię.}
Załóżmy, że $A=f^{-1}[B]$ dla pewnej funkcji
ciągłej $f:\omega^\omega\to\omega^\omega$.
Po każdej turze gracza I znamy jego dotychczasowy
skończony ciąg ruchów $s\in\omega^{<\omega}$.
Możemy wówczas określić największą liczbę
naturalną $m(s)\leq|s|$ taką, że wszystkie
ciągi $f(x)$ przy $x\in[s]$ mają jednakowe
pierwsze $m(s)$ wyrazów. Taka liczba istnieje,
bo $m=0$ zawsze spełnia ten warunek, a rozpatrujemy
skończony zbiór możliwych wartości $m$.

Gracz II przechowuje skończony ciąg $t$ już
podanych wyrazów. Jeśli $|t|<m(s)$, podaje
następny wyraz wspólnego prefiksu długości $m(s)$;
jeśli $|t|\ge m(s)$, pomija ruch.
W każdym kroku jego poprzednie odpowiedzi są
prefiksem $f(x)$ dla każdego $x\in[s]$.
Ustalmy dowolną kompletną rozgrywkę gracza I,
czyli ciąg $x\in\omega^\omega$.
Ciągłość $f$ w punkcie $x$ daje dla każdego
$k\in\omega$ pewną liczbę naturalną $N$ taką,
że wszystkie $x'\in[x\upharpoonright N]$
spełniają $f(x')\upharpoonright k=f(x)\upharpoonright k$.
Dlatego dla wszystkich $n\ge\max(N,k)$
mamy $m(x\upharpoonright(n+1))\ge k$.
Jeśli do tego czasu gracz II nie podał jeszcze
$k$ wyrazów, będzie podawał po jednym wyrazie
w kolejnych turach, aż uzyska długość $k$.
Ponieważ $k$ było dowolne, gracz II podaje
nieskończenie wiele wyrazów, tworząc dokładnie
$f(x)$. Równość $A=f^{-1}[B]$ zapewnia
wygraną gracza II w każdej takiej rozgrywce.
\end{proof}

\begin{uwaga}[Zależności od determinacji projektywnej]
Trzy pierwsze poziomy redukowalności pokazano
konstrukcyjnie w ZFC. Dopiero półliniowość
porządku Wadge'a dla wszystkich zbiorów
projektywnych używa PD
(Twierdzenie~\ref{thm:pd-projective-wadge}).
Nie twierdzimy tutaj, że pełna hierarchia Wadge'a
jest dobrze uporządkowana: jej analiza, rangi
Wadge'a i twierdzenia o skalach wymagają
oddzielnych wyników deskryptywnej teorii mnogości.
\end{uwaga}


\section{Samodualność i pierwsze stopnie redukcji Wadge'a}
\label{sec:wadge-selfduality-clopen}
Zbiór $A\subseteq\omega^\omega$ nazywamy \emph{samodualnym względem redukcji Wadge'a}, jeżeli $A\equiv_{\mathrm W}(\omega^\omega\setminus A)$.
To własność stopnia redukcji, nie równość zbiorów.

\begin{twierdzenie}[Samodualność niebanalnych zbiorów otwarto-domkniętych]
\label{thm:wadge-clopen-selfdual}
Jeśli $A\subseteq\omega^\omega$ jest niepustym, właściwym zbiorem otwartym i domkniętym, to $A\equiv_{\mathrm W}(\omega^\omega\setminus A)$.
Natomiast $\varnothing$ i $\omega^\omega$ nie są samodualne.
\end{twierdzenie}
\begin{proof}
Dopełnienie $A$ również jest niepustym, właściwym zbiorem otwartym i domkniętym. Z pierwszego punktu Wniosku~\ref{cor:wadge-first-degrees} otrzymujemy obie redukcje.
Dla dowolnej funkcji $f:\omega^\omega\to\omega^\omega$ mamy $f^{-1}[\varnothing]=\varnothing$ oraz $f^{-1}[\omega^\omega]=\omega^\omega$. Dlatego $\omega^\omega\not\leq_{\mathrm W}\varnothing$ i $\varnothing\not\leq_{\mathrm W}\omega^\omega$.
\end{proof}

\begin{twierdzenie}[Niesamodualność właściwie otwartego stopnia zupełnego]
\label{thm:wadge-open-nonselfdual}
Jeśli $U\subseteq\omega^\omega$ jest zbiorem otwartym zupełnym z Twierdzenia~\ref{thm:wadge-open-complete}, to $U\not\equiv_{\mathrm W}(\omega^\omega\setminus U)$.
\end{twierdzenie}
\begin{proof}
Dopełnienie $U$ jest domknięte. Gdyby $U\leq_{\mathrm W}(\omega^\omega\setminus U)$, ciągły przeciwobraz zbioru domkniętego byłby domknięty. To przeczy temu, że $U$ nie jest domknięty.
\end{proof}

\begin{przyklad}[Dwa różne typy stopni]
Zbiór $A=\{x\in\omega^\omega:x(0)=0\}$ jest niepusty, właściwy, otwarty i domknięty, a więc samodualny. Zbiór $U=\{x:\exists n\ x(n)=0\}$ jest otwarty, ale nie domknięty: ciągi o coraz późniejszym pierwszym zerze zbiegają do ciągu stale równego $1$, który nie należy do $U$. Ponieważ $U$ jest zupełny dla zbiorów otwartych, jego stopień nie jest samodualny.
\end{przyklad}


\begin{uwaga}[Audyt warunku samodualności]
\label{rem:wadge-selfdual-audit-32}
Samodualność dotyczy równoważności przez \emph{ciągłe przeciwobrazy}, a nie istnienia homeomorfizmu zamieniającego zbiór z dopełnieniem. W szczególności dowód Twierdzenia~\ref{thm:wadge-clopen-selfdual} wymaga, aby obie części przestrzeni były otwarte: tylko wtedy dwuwartościowa funkcja stała na obu częściach jest ciągła. Dla zbioru otwartego, który nie jest domknięty, taki argument nie działa.
\end{uwaga}
\chapter{Uniwersum konstruowalne i modele wewnętrzne}\label{ch:constructible}
\chaptermark{Uniwersum konstruowalne}
\section{Dlaczego ograniczać zbiór potęgowy?}
Aksjomat zbioru potęgowego pozwala przejść od $X$ do wszystkich jego podzbiorów. Konstrukcja Gödla $L$ zachowuje jedynie podzbiory definiowalne nad dotychczas zbudowaną strukturą. Prowadzi to do szczególnej klasy tranzytywnej, której elementy mają kanoniczne opisy.
\begin{definicja}[Definiowalne podzbiory]\label{def:constructible-def}
Niech $M$ będzie niepustym zbiorem. Zbiór
$\operatorname{Def}(M)$ składa się ze wszystkich podzbiorów
$A\subseteq M$, dla których istnieją formuła pierwszego rzędu
$\varphi(x,z_1,\ldots,z_n)$ języka $\{\in\}$ (dopuszczamy $n=0$)
oraz parametry $p_1,\ldots,p_n\in M$ takie, że
\[
 A=\{x\in M:(M,\in\cap(M\times M))\models
       \varphi(x,p_1,\ldots,p_n)\}.
\]
Definiowalność odnosimy do \emph{struktury o uniwersum $M$},
a nie do całego uniwersum teorii mnogości.
Dla $M=\varnothing$ przyjmujemy osobno
$\operatorname{Def}(\varnothing):=\{\varnothing\}$.
Ta konwencja pozwala zdefiniować pierwszy poziom $L_1$,
bez uznawania pustego zbioru za strukturę pierwszego rzędu
(o niepustej dziedzinie w przyjętej konwencji logicznej).
\end{definicja}
\begin{lemat}[Istnienie operatora definiowalnych podzbiorów]
\label{lem:def-set}
Dla każdego zbioru $M$ istnieje w ZF zbiór
$\operatorname{Def}(M)$ i zachodzi
$\operatorname{Def}(M)\subseteq\mathcal P(M)$.
Jeśli $M\ne\varnothing$, rodzina ta zawiera $\varnothing$
oraz $M$.
\end{lemat}
\begin{proof}
Dla zbioru pustego stwierdzenie wynika bezpośrednio
z konwencji w Definicji~\ref{def:constructible-def}.
Załóżmy więc $M\ne\varnothing$. Kody formuł tworzą zbiór
liczb naturalnych, a skończone krotki parametrów z $M$
tworzą zbiór $M^{<\omega}$. W ZF można dla \emph{ustalonej}
struktury zbiorowej $(M,\in)$ zdefiniować relację spełniania
$\operatorname{Sat}_M(e,\vec p,a)$, gdzie $e$ jest kodem
formuły o wyróżnionej zmiennej $x$, $\vec p$ jest krotką
jej pozostałych parametrów i $a\in M$. Relację tę
konstruujemy przez rekurencję po skończonym drzewie
podformuł: przypadki atomowe używają przynależności i
równości w $M$, spójniki odpowiednich operacji logicznych,
a kwantyfikatory przebiegają zbiór $M$.
Dla każdego kodu $(e,\vec p)$ określamy, przez wyróżnianie,
\[
 A_{e,\vec p}=\{a\in M:\operatorname{Sat}_M(e,\vec p,a)\}.
\]
Schemat zastępowania, stosowany do zbioru dopuszczalnych
kodów i krotek, daje zbiór wszystkich takich wartości
$A_{e,\vec p}$. Jest to dokładnie $\operatorname{Def}(M)$.
Każda wartość jest podzbiorem $M$. Formuły $x\ne x$
i $x=x$ definiują odpowiednio $\varnothing$ i $M$
bez dodatkowych parametrów.
\end{proof}
\begin{definicja}[Hierarchia $L$]\label{def:l-hierarchy}
Przez rekurencję pozaskończoną określamy $L_0=\varnothing$, $L_{\alpha+1}=\operatorname{Def}(L_\alpha)$ oraz $L_\lambda=\bigcup_{\beta<\lambda}L_\beta$ dla ordynałów granicznych $\lambda$. Klasę $L=\bigcup_{\alpha\in\operatorname{Ord}}L_\alpha$ rozumiemy jako klasę wszystkich zbiorów należących do pewnego $L_\alpha$; nie jest to zbiór.
\end{definicja}
\begin{lemat}[Tranzytywność i monotoniczność]\label{lem:l-transitive}
Dla każdego ordynału $\alpha$ zbiór $L_\alpha$ jest tranzytywny i $L_\alpha\subseteq L_{\alpha+1}$.
\end{lemat}
\begin{proof}
Indukcja pozaskończona. Jeśli $L_\alpha$ jest tranzytywny, to dla każdego $a\in L_\alpha$ mamy $a\subseteq L_\alpha$; jest on definiowalny w $L_\alpha$ wzorem $x\in a$ z parametrem $a$, więc $a\in\operatorname{Def}(L_\alpha)$. Każdy element $L_{\alpha+1}$ jest podzbiorem $L_\alpha$, co daje tranzytywność następnego poziomu. Na etapie granicznym $L_\lambda=\bigcup_{\beta<\lambda}L_\beta$ jest tranzytywny. Aby otrzymać $L_\lambda\subseteq L_{\lambda+1}$, zauważmy, że dla $a\in L_\lambda$ mamy $a\subseteq L_\lambda$, a zatem $a$ jest definiowalny nad $L_\lambda$ wzorem $x\in a$ z parametrem $a$.
\end{proof}
\begin{przyklad}[Cztery pierwsze poziomy hierarchii]
\label{ex:l-small-levels-34}
Zaczynamy od $L_0=\varnothing$. Z konwencji dla pustego
uniwersum otrzymujemy $L_1=\{\varnothing\}$.
Struktura o uniwersum $L_1$ ma jeden element, zatem
formuły $x\ne x$ i $x=x$ opisują jej dwa podzbiory:
$\varnothing$ i $L_1$. Stąd
$L_2=\{\varnothing,\{\varnothing\}\}$.
Zbiór $L_2$ ma dwa elementy: zbiór pusty i zbiór
jednoelementowy. W strukturze $(L_2,\in)$ pierwszy
odróżnia formuła $\forall y\ (y\notin x)$,
a drugi jej negacja. Wskutek tego dowolną kombinację
tych dwóch elementów można wskazać formułą, więc
\[
 L_3=\mathcal P(L_2)
 =\{\varnothing,\{\varnothing\},
    \{\{\varnothing\}\},L_2\}.
\]
W szczególności $|L_0|=0$, $|L_1|=1$,
$|L_2|=2$ i $|L_3|=4$.
W istocie dla każdego $n<\omega$ zachodzi $L_n=V_n$.
Uzasadnijmy indukcję. Zbiór $V_n$ jest skończony i tranzytywny.
Każdy jego element $a$ ma jednoznaczny opis bez parametrów:
dla $a=\varnothing$ stosujemy $\varphi_a(x):=\forall y\,(y\notin x)$,
a dla niepustego $a$ (który jest skończony) indukcyjnie
budujemy formułę
\[
 \varphi_a(x):=\forall y\bigl(y\in x\leftrightarrow
        \bigvee_{b\in a}\varphi_b(y)\bigr).
\]
Jest to formuła skończona, ponieważ każdy element $a$
ma skończenie wiele elementów, a relacja przynależności
na zbiorach dziedzicznie skończonych jest dobrze ufundowana.
W strukturze $(V_n,\in)$ formuła $\varphi_a$ rozpoznaje
jednoznacznie $a$ na mocy ekstensjonalności.
Każdy podzbiór $B\subseteq V_n$ jest skończony, więc
definiuje go alternatywa formuł $\varphi_b(x)$ po $b\in B$;
dla pustego $B$ wybieramy $x\ne x$. Stąd
$\operatorname{Def}(V_n)=\mathcal P(V_n)=V_{n+1}$.
Krok $n=0$ daje odrębna konwencja
$\operatorname{Def}(\varnothing)=\{\varnothing\}$.
Ta równość na skończonych poziomach nie oznacza, że
operator $\operatorname{Def}$ ogólnie jest operatorem
zbioru potęgowego; dla nieskończonych struktur
istnieją podzbiory pozbawione definicji z parametrami.
\end{przyklad}

\begin{lemat}[Jawne kodowanie zbiorów dziedzicznie skończonych]
\label{lem:lvi-ackermann-code}
Niech $V_\omega=\bigcup_{n<\omega}V_n$ oznacza zbiór
wszystkich zbiorów dziedzicznie skończonych. Istnieje
bijekcja $a\colon V_\omega\to\omega$ określona
rekurencyjnie wzorem
\[
 a(x)=\sum_{y\in x}2^{a(y)}.
\]
Suma jest skończona, gdyż każdy $x\in V_\omega$
jest zbiorem skończonym. Do konstrukcji i dowodu nie
potrzeba aksjomatu wyboru.
\end{lemat}
\begin{proof}
\emph{Krok 1: istnienie wartości.}
Dla $x\in V_n$ każdy $y\in x$ należy do $V_{n-1}$,
gdy $n>0$. Rekursja po randze pozwala więc obliczyć
wszystkie $a(y)$ przed $a(x)$. Wartość $a(x)$ jest
liczbą naturalną, ponieważ jest sumą skończenie
wielu potęg liczby $2$. Wzór wyznacza wartość jednoznacznie.

\emph{Krok 2: iniekcyjność.}
Przeprowadzamy indukcję po największej z dwóch
skończonych rang zbiorów $x,z$. Jeżeli $a(x)=a(z)$,
to jednoznaczność zapisu liczby naturalnej w systemie
binarnym pokazuje, że
$\{a(y):y\in x\}=\{a(w):w\in z\}$.
Wszystkie elementy $x$ i $z$ mają mniejsze rangi;
z założenia indukcyjnego wynika, że elementy o tym
samym kodzie są równe. Otrzymujemy $x=z$.

\emph{Krok 3: surjekcyjność.}
Stosujemy silną indukcję po $m\in\omega$. Dla $m=0$
mamy $a(\varnothing)=0$. Dla $m>0$ zapis binarny
$m=\sum_{j\in F}2^j$ ma pewien skończony zbiór
pozycji $F\subseteq m$: z $2^j\le m$ wynika $j<m$.
Dla każdego $j\in F$ założenie indukcyjne daje
jedyny $x_j\in V_\omega$ o kodzie $j$.
Zbiór $x=\{x_j:j\in F\}$ należy do $V_\omega$
(skończona rodzina zbiorów dziedzicznie skończonych
ma skończoną sumę rang), a z definicji $a(x)=m$.
To kończy dowód bijektywności.
\end{proof}
\begin{przyklad}[Ręczne obliczenie kodu i pierwszy skok w hierarchii]
\label{ex:lvi-first-infinite-level}
Kodujemy kolejno
$a(\varnothing)=0$, $a(\{\varnothing\})=2^0=1$,
$a(\{\{\varnothing\}\})=2^1=2$ oraz
$a(\{\varnothing,\{\varnothing\}\})=2^0+2^1=3$.
Zatem $L_\omega=V_\omega$ ma jawną numerację
przez liczby naturalne. Każdy element
$L_{\omega+1}=\operatorname{Def}(L_\omega)$
posiada opis składający się z kodu formuły i
skończonej krotki parametrów z $L_\omega$.
Zbiór wszystkich takich opisów jest przeliczalny:
wykorzystujemy kodowanie skończonych krotek
liczb naturalnych. Wobec tego $L_{\omega+1}$
jest przeliczalny. Tymczasem
$\mathcal P(L_\omega)$ nie jest przeliczalny
na mocy twierdzenia Cantora
(Twierdzenie~\ref{thm:zf-cantor}). Otrzymujemy ścisłą
inkluzję
\[
 L_{\omega+1}\subsetneq\mathcal P(L_\omega).
\]
Jest to pierwszy etap, na którym operacja
$\operatorname{Def}$ nie daje całego zbioru
potęgowego. Nie wybieraliśmy żadnych reprezentantów
z nieskończonej rodziny niepustych zbiorów.
\end{przyklad}

\begin{przyklad}[Dlaczego potrzebna jest absolutność]
\label{ex:l-unbounded-quantifier-34}
Ustalmy parametr $x=\varnothing$ i formułę
$\psi(x):=\exists y\ (x\in y)$. W strukturze
$(L_1,\in)$ formuła $\psi(\varnothing)$ jest fałszywa,
gdyż jedynym elementem $L_1$ jest $\varnothing$.
W strukturze $(L_2,\in)$ ta sama formuła jest
prawdziwa: świadkiem jest $y=\{\varnothing\}\in L_2$.
Nieograniczony kwantyfikator egzystencjalny nie jest
więc absolutny nawet między dwoma \emph{tranzytywnymi}
poziomami hierarchii $L$. Natomiast w formule
$\exists y\in a$ elementy $a$ nie zmieniają się
przy przejściu do większego zbioru tranzytywnego.
\end{przyklad}
\section{Absolutność i ostrożność przy kwantyfikatorach}
\begin{definicja}[Formuły ograniczone]\label{def:delta0}
Formułą $\Delta_0$ nazywamy formułę zbudowaną z formuł atomowych przez spójniki i kwantyfikatory ograniczone $\forall x\in y$ i $\exists x\in y$. Formuła jest absolutna między klasami tranzytywnymi, jeśli jej prawdziwość dla parametrów należących do obu klas jest w nich taka sama.
\end{definicja}
\begin{twierdzenie}[Absolutność formuł ograniczonych]\label{thm:delta0-absolute}
Niech $M\subseteq N$ będą tranzytywnymi klasami, w których relacja przynależności jest rzeczywistą relacją $\in$. Dla dowolnej formuły $\Delta_0$ $\varphi$ i dowolnych parametrów $a_1,\ldots,a_n\in M$ zachodzi $M\models\varphi(\vec a)$ wtedy i tylko wtedy, gdy $N\models\varphi(\vec a)$ (przy założeniu, że obie relacje spełniania są rozpatrywane w odpowiedniej metateorii).
\end{twierdzenie}
\begin{proof}
Indukcja po budowie formuły. Równość i przynależność nie zmieniają się przy przejściu z $M$ do $N$. Spójniki zachowują równoważność. Dla $\exists x\in a\,\psi(x,\vec a)$ tranzytywność $M$ zapewnia, że każdy element $a\in M$ również należy do $M$. Zatem kwantyfikator w obu klasach przebiega dokładnie ten sam zbiór elementów $a$. Korzystamy z założenia indukcyjnego dla $\psi$. Analogicznie postępujemy z kwantyfikatorem ogólnym.
\end{proof}
\begin{uwaga}[Granica argumentu]
Zwykły kwantyfikator $\exists x$ może mieć świadków w $N\setminus M$, dlatego powyższe twierdzenie nie dotyczy dowolnych formuł. Pełne twierdzenie $L\models\mathrm{ZFC}+\mathrm{GCH}$
nie wynika jeszcze z samej definicji poziomów hierarchii.
Potrzebne są zasada kondensacji, definiowalny dobry porządek
na klasie $L$ i sprawdzenie schematu zastępowania.
Dowody tych rezultatów znajdują się w dalszych częściach
niniejszego rozdziału (w szczególności w
Lemacie~\ref{lem:l-condensation} i
Twierdzeniu~\ref{thm:l-gch}).
\end{uwaga}

\section{Dlaczego definicja $L$ nie daje jeszcze ZFC?}
\begin{uwaga}[Zbiór a klasa właściwa]
Symbol $L$ oznacza klasę opisaną formułą „istnieje ordynał $\alpha$ taki,
że $x\in L_\alpha$”. Nie wolno stosować do $L$ aksjomatu zbioru
potęgowego tak, jakby był zbiorem. Zdanie $L\models\varphi$
wymaga relatywizacji kwantyfikatorów formuły $\varphi$ do klasy $L$;
nie jest to relacja spełniania dla struktury o zbiorowym uniwersum.
\end{uwaga}
\begin{przyklad}[Definiowalność zależy od struktury]
Niech $M=\{\varnothing\}$. Wtedy $\operatorname{Def}(M)$
składa się z $\varnothing$ i $M$: pierwszy podzbiór
definiuje $x\ne x$, drugi $x=x$. Dla $M=\varnothing$
opisuje go osobna konwencja z
Definicji~\ref{def:constructible-def}, nie relacja spełniania
w strukturze pierwszego rzędu o pustej dziedzinie.
Zgodnie z Przykładem~\ref{ex:l-unbounded-quantifier-34}
różne struktury mogą nadawać różne wartości logiczne
tej samej formule. Dlatego przy operatorze
$\operatorname{Def}$ zawsze wskazujemy uniwersum struktury.
\end{przyklad}

\section{Odbicie formuł do poziomu hierarchii $L$}
\label{sec:l-reflection}
Nie wystarcza powiedzieć, że $L$ jest sumą poziomów $L_\alpha$.
Trzeba jeszcze pokazać, że kwantyfikatory w aksjomatach ZF
można kontrolować wewnątrz tej klasy. Poniższy lemat
jest narzędziem do tego celu.

\begin{lemat}[Odbicie skończonej rodziny formuł]
\label{lem:l-finite-reflection}
Niech $\Phi$ będzie skończonym zbiorem formuł pierwszego
rzędu języka $\{\in\}$, domkniętym na podformuły.
Dla każdego ordynału $\eta$ istnieje ordynał graniczny
$\theta>\eta$ taki, że dla każdej $\varphi\in\Phi$
i wszystkich parametrów $\vec a\in L_\theta$ zachodzi
\[
(L_\theta,\in)\models\varphi(\vec a)
\quad\Longleftrightarrow\quad L\models\varphi(\vec a).
\]
Przez $L\models\varphi$ rozumiemy formułę otrzymaną
przez ograniczenie wszystkich kwantyfikatorów do
\emph{definiowalnej klasy} $L$; nie zakładamy istnienia
zbiorowej struktury o uniwersum $L$.
\end{lemat}
\begin{proof}
\emph{Etap 1: wybieranie etapów, nie świadków.}
Ustalmy formułę $\exists y\,\psi(y,\vec x)$ będącą
podformułą którejś formuły z $\Phi$. Jeżeli
$\vec a\in L_\alpha$ i $L\models\exists y\,\psi(y,\vec a)$,
to istnieje najmniejszy ordynał $\beta$ taki, że
$\exists y\in L_\beta\,\psi^L(y,\vec a)$.
Ta liczba porządkowa jest wyznaczona jednoznacznie
przez $\vec a$ i kod formuły. Ponieważ krotek
parametrów z $L_\alpha$ jest tylko zbiór,
aksjomat zastępowania w metateorii ZF zapewnia
wspólne ograniczenie dla ich najmniejszych etapów
oraz dla skończenie wielu formuł egzystencjalnych.
Nie wybieramy tutaj naraz elementów $y$ z całej rodziny:
wystarcza zebrać \emph{ordynały} ich pierwszego wystąpienia.

\emph{Etap 2: jednoznacznie określone domknięcie.}
Aby nie używać wyboru przeliczalnego,
ustalamy $\alpha_0:=\eta+1$.
Dla każdej $\alpha$ niech $E_\alpha$
będzie zbiorem par $(e,\vec a)$,
gdzie $e$ jest kodem podformuły
egzystencjalnej z $\Phi$,
$\vec a$ jest odpowiednią krotką
parametrów z $L_\alpha$ i w $L$
istnieje jej świadek.
Zbiór $E_\alpha$ istnieje przez
separację w skończonej sumie
zbiorów krotek parametrów.
Każdej parze z $E_\alpha$
przyporządkowujemy najmniejszy
ordynał $\beta_{e,\vec a}$ taki,
że świadek należy do
$L_{\beta_{e,\vec a}}$.
Zbiór wszystkich tych ordynałów
istnieje przez zastępowanie.
Kładziemy jednoznacznie
\[
F(\alpha):=
\sup\bigl(\{\alpha+1\}
 \cup\{\beta_{e,\vec a}+1:
                  (e,\vec a)\in E_\alpha\}\bigr).
\]
Definiujemy przez rekurencję na $\omega$
$\alpha_{n+1}:=F(\alpha_n)$.
Mamy $\alpha_{n+1}>\alpha_n$
i każdy wskazany świadek należy do
$L_{\alpha_{n+1}}$.
Niech $\theta:=\sup_{n\in\omega}\alpha_n$;
jest to ordynał graniczny większy od $\eta$.
Dla każdej skończonej krotki
$\vec a\in L_\theta^{<\omega}$
istnieje $n$ takie, że
wszystkie jej współrzędne należą
do $L_{\alpha_n}$.
Jeżeli odpowiednia podformuła
egzystencjalna ma świadka w $L$,
to ma go już w
$L_{\alpha_{n+1}}\subseteq L_\theta$.
Wszystkie etapy rekurencji
wyznaczono jednoznacznymi formułami ZF;
nie zastosowano AC ani DC.

\emph{Etap 3: indukcja po formułach.}
Dla formuł atomowych prawdziwość jest taka sama,
bo $L_\theta$ jest tranzytywny. Spójniki zachowują
równoważność. Gdy $\exists y\,\psi(y,\vec a)$ jest
prawdziwe w $L$, Etap 2 dostarcza świadka
$y\in L_\theta$, a założenie indukcyjne daje
prawdziwość $\psi$ w $L_\theta$.
W przeciwną stronę świadek z $L_\theta$ jest
również elementem $L$, więc stosujemy
założenie indukcyjne. Kwantyfikator ogólny
wyrażamy przez negację egzystencjalnego.
Otrzymujemy jednocześnie równoważność dla wszystkich
formuł z $\Phi$.
\end{proof}

\section{Dowód, że $L$ spełnia aksjomaty ZF}
\label{sec:l-model-zf}
\begin{twierdzenie}[Uniwersum konstruowalne jest modelem wewnętrznym ZF]
\label{thm:l-model-zf}
W metateorii ZF klasa $L$ jest tranzytywna, zawiera
wszystkie ordynały oraz spełnia każdy aksjomat ZF
po relatywizacji kwantyfikatorów do $L$.
\end{twierdzenie}
\begin{proof}
\emph{Krok 1: tranzytywność, ordynały i ekstensjonalność.}
Tranzytywność wynika z Lematu~\ref{lem:l-transitive}.
Uzasadnijmy dokładniej równość
$L_\alpha\cap\operatorname{Ord}=\alpha$,
która jest istotna dla wszystkich dalszych etapów.
W ZF zbiór $x$ jest ordynałem wtedy i tylko wtedy,
gdy $x$ jest przechodni i każdy element $x$
jest przechodni; dzięki ufundowaniu
jest to warunek równoważny dobremu
uporządkowaniu przez $\in$.
Warunek przechodniości można zapisać
\[
\forall u\in x\;\forall v\in u\;(v\in x),
\]
a warunek przechodniości każdego elementu
$x$ ma analogiczną postać
z kwantyfikatorami ograniczonymi.
Zatem predykat $\operatorname{Ord}(x)$
ma definicję $\Delta_0$ i jest
absolutny dla zbiorów przechodnich
(Twierdzenie~\ref{thm:delta0-absolute}).
Przez indukcję po $\alpha$, gdy
$L_\alpha\cap\operatorname{Ord}=\alpha$,
zbiór
\[
\alpha=\{x\in L_\alpha:
             (L_\alpha,\in)\models\operatorname{Ord}(x)\}
\]
jest definiowalnym podzbiorem $L_\alpha$.
Stąd $\alpha\in L_{\alpha+1}$,
a ponieważ każdy element $L_{\alpha+1}$
jest podzbiorem przechodniego $L_\alpha$,
nie może pojawić się na tym etapie
ordynał większy od $\alpha$:
jeśli $\gamma\in L_{\alpha+1}$
jest ordynałem, to
$\gamma\subseteq L_\alpha$ i
$\gamma\subseteq\alpha$, zatem
$\gamma\le\alpha$.
Na etapach granicznych korzystamy
z monotoniczności i sumy poziomów,
a dla $\alpha=0$ obie strony są puste.
To kończy indukcję i pokazuje,
że wszystkie ordynały należą do $L$.
Równość i przynależność są rzeczywistymi relacjami
na zbiorach, więc ekstensjonalność zachodzi w $L$.

\emph{Krok 2: pustka, pary, suma i nieskończoność.}
Dla $a,b\in L$ wybieramy $\alpha$ z $a,b\in L_\alpha$.
Zbiór $\{a,b\}$ jest podzbiorem $L_\alpha$
definiowalnym w $L_\alpha$ formułą $x=a\lor x=b$;
stąd $\{a,b\}\in L_{\alpha+1}$. Ponieważ
$L_\alpha$ jest tranzytywny, suma $\bigcup a$
jest podzbiorem $L_\alpha$ definiowalnym przez
$\exists y\in a\,x\in y$. Zatem $\bigcup a\in L$.
Zbiór pusty należy do $L_1$.
Wykażmy indukcyjnie $L_n=V_n$ dla $n\in\omega$.
Dla $n=0$ oba zbiory są puste.
Jeśli $L_n=V_n$ jest skończony,
to każdy jego podzbiór $B$ jest
definiowalny nad $L_n$: dla
$B=\{a_1,\ldots,a_k\}$ wystarczy
formuła
$x=a_1\lor\cdots\lor x=a_k$
z parametrami $a_i\in L_n$,
a dla $B=\varnothing$ formuła $x\ne x$.
Zatem
$\operatorname{Def}(L_n)=\mathcal P(L_n)$
(również dla $n=0$ na mocy konwencji),
czyli $L_{n+1}=V_{n+1}$.
Stąd $L_\omega=V_\omega$ jest
zbiorem wszystkich dziedzicznie
skończonych zbiorów.
Z wcześniej uzasadnionej
absolutności predykatu
$\operatorname{Ord}$ wynika,
że zbiór jego elementów będących
skończonymi ordynałami jest
definiowalny nad $L_\omega$
i równy $\omega$.
Stąd $\omega\in L_{\omega+1}$,
co daje aksjomat nieskończoności.

\emph{Krok 3: wyróżnianie.}
Niech $a,\vec p\in L$ i niech $\varphi(x,\vec p)$
będzie ustaloną formułą. Wybieramy $\theta$ z
Lematu~\ref{lem:l-finite-reflection} dla tej formuły
oraz wszystkich jej podformuł, tak duży, aby
$a,\vec p\in L_\theta$. Wtedy
\[
\{x\in a:L\models\varphi(x,\vec p)\}
=\{x\in L_\theta:x\in a\land
L_\theta\models\varphi(x,\vec p)\}.
\]
Prawa strona jest definiowalnym podzbiorem $L_\theta$,
a więc należy do $L_{\theta+1}$.

\emph{Krok 4: zbiór potęgowy.}
Ustalmy $a\in L$. Zewnętrzna rodzina
$B=\{b\subseteq a:b\in L\}$ jest zbiorem w ZF:
wydzielamy ją z $\mathcal P(a)$ za pomocą formuły
„$b\in L$”. Dla każdego $b\in B$ istnieje najmniejszy
etap $\beta_b$ z $b\in L_{\beta_b}$; przez
zastępowanie istnieje wspólne ograniczenie tych
etapów. Wybieramy $\theta$ tak duże, że
$B\subseteq L_\theta$ oraz $a\in L_\theta$.
Wtedy
\[
B=\{b\in L_\theta:\forall x\in b\ (x\in a)\}.
\]
Zbiór $B$ jest więc definiowalny nad $L_\theta$,
a zatem $B\in L_{\theta+1}$. Wewnątrz $L$
pełni dokładnie rolę zbioru potęgowego $a$:
zawiera wszystkie i tylko podzbiory $a$, które
należą do $L$. Nie należy go utożsamiać z całym
zewnętrznym $\mathcal P(a)$.

\emph{Krok 5: zastępowanie.}
Załóżmy, że $L\models\forall x\in a\ \exists!y\,
\varphi(x,y,\vec p)$, gdzie $a,\vec p\in L$.
W metateorii ZF formuła $\varphi^L$ jest zwykłą
formułą pierwszego rzędu z klasą $L$ określoną
definiowalnie. Dlatego przez zewnętrzny schemat
zastępowania istnieje zbiór $B$ wszystkich jej
jednoznacznych wartości $y$ dla $x\in a$.
Każda z tych wartości jest w $L$. Wybieramy
$\theta$ tak duże, aby $a,B,\vec p\subseteq L_\theta$
(w przypadku $a$ i parametrów chodzi także o ich
przynależność do $L_\theta$), a ponadto aby
$L_\theta$ poprawnie rozstrzygał $\varphi$
na swoich parametrach, zgodnie z Lematem
\ref{lem:l-finite-reflection}. Wówczas
\[
B=\{y\in L_\theta:\exists x\in a\;
L_\theta\models\varphi(x,y,\vec p)\}.
\]
Dzięki odbiciu nie pojawiają się tu fałszywe wartości.
Prawa strona jest definiowalna nad $L_\theta$;
wobec tego $B\in L_{\theta+1}$. Dowiedliśmy
zastępowania \emph{wewnątrz $L$}.

\emph{Krok 6: ufundowanie.}
Jeśli $a\in L$ jest niepusty w $L$, to jest niepusty
w zewnętrznym ZF. Aksjomat ufundowania daje element
$\in$-minimalny $x\in a$. Tranzytywność $L$ oznacza,
że $x\in L$, a więc jest on również świadkiem
ufundowania w $L$. Wszystkie aksjomaty ZF zostały
sprawdzone; nie wykorzystaliśmy AC w metateorii.
\end{proof}


\begin{lemat}[Absolutność hierarchii konstruowalnej wewnątrz $L$]
\label{lem:l-hierarchy-internal}
Niech $\alpha$ będzie dowolną liczbą porządkową. Poziom hierarchii
konstruowalnej obliczony wewnątrz klasy $L$ jest tym samym zbiorem,
co poziom $L_\alpha$ obliczony w otaczającej teorii ZF:
\[
 (L_\alpha)^L=L_\alpha.
\]
W szczególności, jeżeli $M=L_\alpha$, to dla każdej \emph{ustalonej}
formuły pierwszego rzędu $\varphi$ i parametrów z $M$ relacja
spełniania $M\models\varphi$ jest jednakowo obliczana w $L$ i w ZF.
\end{lemat}
\begin{proof}
Najpierw ustalmy $M\in L$. Prawdziwość formuły atomowej
w zbiorowej strukturze $(M,\in)$ zależy tylko od równości
i przynależności jej elementów. Dla spójników prawdziwość
wyznacza się z wartości logicznych podformuł; dla kwantyfikatora
$\exists x$ przebiega się dokładnie zbiór $M$, który jest ten sam
w klasie tranzytywnej $L$ i w otaczającej teorii. Indukcja
po budowie \emph{ustalonej} formuły daje więc absolutność
relacji spełniania w $(M,\in)$.

Ponadto, jeżeli $B\subseteq M$ jest podzbiorem definiowalnym
w $(M,\in)$ za pomocą formuły i parametrów z $M$, to
$B\in L$: można go wydzielić w $L$ ze zbioru $M$ formułą
wyrażającą relację spełniania zbiorowej struktury.
Odwrotnie, każda definicja obliczona w $L$ opisuje dokładnie
ten sam podzbiór $M$ co w ZF. Zatem
$\operatorname{Def}^L(M)=\operatorname{Def}(M)$.

Przeprowadzamy teraz indukcję pozaskończoną po $\alpha$.
Dla zera oba poziomy są puste. Jeżeli mają tę samą wartość
na etapie $\alpha$, to w kroku następnym wyznacza je ten sam
operator $\operatorname{Def}$. Na etapie granicznym obie
hierarchie biorą sumę tej samej rodziny wcześniejszych
poziomów. Każda liczba porządkowa należy do $L$, więc
indukcja obejmuje wszystkie liczby porządkowe zewnętrznego
uniwersum.
\end{proof}

\section{Kanoniczny dobry porządek $L$ i aksjomat wyboru}
\label{sec:l-choice}
\begin{twierdzenie}[Definiowalny dobry porządek uniwersum konstruowalnego]
\label{thm:l-definable-wellorder}
Istnieje formuła bez parametrów definiująca w $L$
relację $<_{L}$, której ograniczenie do każdego
zbioru należącego do $L$ jest dobrym porządkiem.
W szczególności $L\models\mathrm{AC}$.
\end{twierdzenie}
\begin{proof}
\emph{Etap 1: porządkujemy zbiór definicji.}
Jeśli $L_\alpha$ jest dobrze uporządkowany,
każdy element $L_{\alpha+1}=\operatorname{Def}(L_\alpha)$
ma co najmniej jeden opis: numer formuły pierwszego
rzędu i skończoną krotkę parametrów z $L_\alpha$.
Porządkujemy kody, porównując najpierw długość
krotki i numer formuły, a następnie leksykograficznie
parametry. Dla ustalonej długości iloczyn skończonej
liczby dobrych porządków z porządkiem leksykograficznym
jest dobrym porządkiem; suma po długościach i numerach
też jest dobrze uporządkowana. Wybieramy dla każdego
zbioru najwcześniejszy spośród jego kodów.

\emph{Etap 2: porządkujemy całe poziomy spójnie.}
Przez rekurencję pozaskończoną tworzymy porządek
$<_{L,\alpha}$ na $L_\alpha$. Na etapie następnym
zachowujemy dotychczasowy porządek na $L_\alpha$,
a \emph{nowe} elementy $L_{\alpha+1}\setminus L_\alpha$
umieszczamy za starymi zgodnie z Etapem 1.
Na etapie granicznym łączymy porządki z etapów
wcześniejszych. Każdy niepusty zbiór konstruktybilny
ma element pojawiający się na najwcześniejszym etapie;
w tej warstwie istnieje najmniejszy element według
wybranego porządku. Otrzymujemy $<_{L}$.
Cała rekurencja jest jednolicie definiowalna z
formuły kodującej $\alpha\mapsto L_\alpha$,
bez pozaskończonego ciągu arbitralnych wyborów.

\emph{Etap 3: wyprowadzamy AC.}
Dla każdego $a\in L$ relacja $<_{L}$ ograniczona do
$a$ jest dobrym porządkiem, a z aksjomatów ZF
spełnionych w $L$ jego wykres jest zbiorem w $L$
(otrzymujemy go przez wyróżnianie w $a\times a$).
Z twierdzenia o dobrym uporządkowaniu Zermela,
Twierdzenie~\ref{thm:choice-equivalence},
wynika aksjomat wyboru w $L$.
\end{proof}

\section{Lemat kondensacji i hipoteza continuum uogólniona}
\label{sec:l-condensation-gch}
\begin{definicja}[Kolaps Mostowskiego]
\label{def:l-mostowski-collapse}
Dla ekstensjonalnej i dobrze ufundowanej relacji
$\in$ na zbiorze $X$ kolapsem Mostowskiego nazywamy
funkcję $\pi:X\to\bar X$ zadaną rekurencją
\[
\pi(x)=\{\pi(y):y\in X\land y\in x\}.
\]
Jej obraz $\bar X$ jest tranzytywny, a $\pi$ jest
izomorfizmem $(X,\in)$ na $(\bar X,\in)$.
Gdy $X$ jest podstrukturą elementarną pewnego
zbioru tranzytywnego, relacja jest ekstensjonalna
na $X$ dzięki elementarności.
\end{definicja}

\begin{lemat}[Elementarność przekroju z elementem podstruktury]
\label{lem:l-elementary-section-75}
Niech $D$ będzie niepustym zbiorem tranzytywnym,
$X\preccurlyeq(D,\in)$,
a $A\in X$ będzie niepustym elementem $D$.
Wówczas
\[
(X\cap A,\in)\preccurlyeq(A,\in).
\]
Jeżeli $\pi:X\to N$ jest kolapsem Mostowskiego, to
$\pi\restriction(X\cap A)$ jest izomorfizmem
$(X\cap A,\in)$ na $(\pi(A),\in)$.
\end{lemat}
\begin{proof}
Ustalmy formułę $\varphi(u,\vec v)$ języka teorii mnogości
i parametry $\vec a\in(X\cap A)^n$.
Prawdziwość $\varphi$ w strukturze $(A,\in)$ jest
równoważna prawdziwości relatywizacji $\varphi^A$
w $(D,\in)$, ponieważ $A\in D$ i tranzytywność $D$ zapewnia,
że wszystkie elementy $A$ należą do uniwersum $D$.
Jeżeli $(A,\in)\models\exists u\,\varphi(u,\vec a)$,
to $(D,\in)$ spełnia formułę
$\exists u\in A\,\varphi^A(u,\vec a)$.
Wszystkie parametry tej formuły należą do $X$.
Elementarność $X\preccurlyeq D$ daje świadka
$u\in X$ takiego, że $u\in A$ i
$(A,\in)\models\varphi(u,\vec a)$.
Kryterium Tarskiego--Vaughta dowodzi pierwszej tezy.
W drugiej części definicja kolapsu daje
$\pi(A)=\{\pi(u):u\in X\cap A\}$, a zachowanie
relacji $\in$ i różnowartościowość $\pi$ zapewniają
wskazany izomorfizm.
\end{proof}

\begin{uwaga}[Co wynika, a co nie wynika z elementarności]
\label{rem:l-sections-corrected-75}
Z założeń $X\preccurlyeq L_\theta$ i $L_\gamma\in X$
\emph{rzeczywiście wynika}
$(X\cap L_\gamma,\in)\preccurlyeq(L_\gamma,\in)$,
bez dodatkowej hipotezy. Nie wynika natomiast od razu
$\pi(L_\gamma)=L_{\pi(\gamma)}$:
prawa strona wymaga przeprowadzenia indukcji po poziomach.
Tych dwóch stwierdzeń nie należy utożsamiać.
\end{uwaga}

\begin{przyklad}[Dlaczego tranzytywność struktury większej jest potrzebna]
\label{ex:l-elementary-section-transitivity-75}
Niech $A=\{\varnothing\}$, $D=\{A\}$ i $X=D$.
Wtedy $X\preccurlyeq(D,\in)$ jest tautologią
(struktura jest identyczna z własną podstrukturą),
$A\in X$ i $A$ jest niepusty. Jednak
$X\cap A=\varnothing$, ponieważ
$\varnothing\notin D$. Nie otrzymujemy zatem
podstruktury elementarnej niepustej struktury
$(A,\in)$. Przyczyną jest nietranzytywność $D$;
Lemat~\ref{lem:l-elementary-section-75}
celowo zakłada tranzytywność większego uniwersum.
\end{przyklad}

\begin{lemat}[Kolaps Mostowskiego a zbiory definiowalne]
\label{lem:l-collapse-def-compatibility}
Niech $\theta>0$ będzie graniczną liczbą porządkową,
$X\preccurlyeq(L_\theta,\in)$,
$\pi:X\to N$ kolapsem Mostowskiego, a $\gamma\in X$
liczbą porządkową taką, że $L_\gamma\in X$.
Jeżeli dodatkowo
$\pi(L_\gamma)=L_{\pi(\gamma)}$, to
\[
\pi[X\cap\operatorname{Def}(L_\gamma)]
  =\operatorname{Def}(L_{\pi(\gamma)}).
\]
Nie zakłada się osobno elementarności przekroju
$X\cap L_\gamma$: wynika ona z
Lematu~\ref{lem:l-elementary-section-75} dla $\gamma>0$.
\end{lemat}
\begin{proof}
Jeżeli $\gamma=0$, to
$L_0=\varnothing$, $\operatorname{Def}(L_0)=\{\varnothing\}$,
a $0\in X$ i $\pi(0)=0$; teza zachodzi.
Załóżmy więc $\gamma>0$ i oznaczmy
$A=L_\gamma$, $A'=L_{\pi(\gamma)}$.
Lemat~\ref{lem:l-elementary-section-75} daje
$(X\cap A,\in)\preccurlyeq(A,\in)$ oraz izomorfizm
$\pi\restriction(X\cap A)$ na $(A',\in)$.
Dla każdej \emph{ustalonej} formuły $\varphi(u,\vec v)$
jej prawdziwość z parametrami w $X\cap A$ jest więc
taka sama w $A$ i w $A'$ po przeniesieniu parametrów
przez $\pi$. Nie stosujemy predykatu prawdy dla klasy $L$.

\emph{Pierwsza inkluzja.}
Niech $z\in X\cap\operatorname{Def}(A)$.
W metateorii istnieje formuła $\varphi$ i krotka
$\vec p\in A^{<\omega}$ takie, że
$z=\{u\in A:(A,\in)\models\varphi(u,\vec p)\}$.
Po \emph{ustaleniu} tej formuły zdanie określające $z$
jest pierwszego rzędu w $L_\theta$:
kwantyfikatory formuły $\varphi$ relatywizujemy do $A$.
Ponieważ $A,z\in X$, elementarność zapewnia
krotkę $\vec p'\in(X\cap A)^{<\omega}$ definiującą
\emph{ten sam} zbiór $z$.
Dla $v\in A'$ wybierzmy $u\in X\cap A$ z
$\pi(u)=v$. Wówczas
\begin{align*}
 v\in\pi(z)
 &\iff u\in z
 \iff (A,\in)\models\varphi(u,\vec p')\\
 &\iff (A',\in)\models
            \varphi(v,\pi(\vec p')).
\end{align*}
Stąd $\pi(z)\in\operatorname{Def}(A')$.

\emph{Druga inkluzja.}
Niech $b=\{v\in A':(A',\in)\models
\varphi(v,\vec q)\}$, gdzie $\vec q\in(A')^{<\omega}$.
Wybieramy preobrazy $\vec p\in(X\cap A)^{<\omega}$
parametrów $\vec q$ i definiujemy
$z=\{u\in A:(A,\in)\models\varphi(u,\vec p)\}$.
Ponieważ $z\in\operatorname{Def}(A)=L_{\gamma+1}$,
a $\gamma+1<\theta$, mamy $z\in L_\theta$.
Zbiór $z$ jest \emph{jednoznacznie} określony w
$L_\theta$ przez $A,\vec p$ i ustaloną relatywizację
$\varphi^A$; elementarność daje $z\in X$.
Poprzednie równoważności pokazują $\pi(z)=b$.
To kończy dowód obu inkluzji.
\end{proof}


\begin{lemat}[Tablica spełniania pojedynczej formuły i jej absolutność]
\label{lem:l-satisfaction-table-transitive-73}
Niech $A$ będzie niepustym zbiorem i niech $e\in\omega$
będzie kodem ustalonej formuły pierwszego rzędu $\varphi_e$.
Istnieje jednoznacznie wyznaczony zbiór $T_{A,e}$,
który koduje, dla każdej podformuły $\psi$ formuły
$\varphi_e$ i każdego wartościowania jej zmiennych
w $A$, prawdziwość $\psi$ w $(A,\in)$.
Jeżeli $N$ jest tranzytywnym \emph{modelem ZF}
zawierającym $A$ i $e$, to tablica skonstruowana w $N$
jest dokładnie $T_{A,e}$ obliczoną w otaczającym uniwersum.
W szczególności prawdziwość każdej ustalonej formuły
w zbiorowej strukturze $(A,\in)$ jest absolutna między
takimi modelami a otaczającym uniwersum.
\end{lemat}
\begin{proof}
Niech $W$ będzie skończonym zbiorem zmiennych występujących
w $\varphi_e$, również tych związanych, a $\operatorname{Sub}(e)$
-- skończonym zbiorem kodów podformuł. Zbiór
$A^W$ wszystkich funkcji $s:W\to A$ istnieje w ZF.
Poszukujemy podzbioru
$T\subseteq\operatorname{Sub}(e)\times A^W$,
którego elementy $(\psi,s)$ oznaczają, że
$(A,\in)\models\psi[s]$.
Dla formuły atomowej $v\in w$ lub $v=w$ prawdziwość
zależy wyłącznie od tego, czy odpowiednie wartości
$s(v),s(w)\in A$ pozostają w relacji przynależności
lub są równe. Dla negacji bierzemy dopełnienie
odpowiedniego przekroju $A^W$; dla koniunkcji
-- przecięcie dwóch takich przekrojów. Jeżeli
$\psi$ ma postać $\exists v\,\chi$, przyjmujemy
$(\psi,s)\in T$ wtedy i tylko wtedy, gdy istnieje
$a\in A$ takie, że $(\chi,s[v\mapsto a])\in T$.
Jest to kwantyfikator po \emph{ustalonym zbiorze} $A$.
Indukcja po wysokości drzewa składniowego konstruuje
wszystkie przekroje $T$ za pomocą wyróżniania,
operacji zbiorowych i skończenie wielu etapów.
Ta sama indukcja daje jednoznaczność: wartości dla
formuł atomowych są ustalone, a każdy konstruktor
logiczny jednoznacznie określa wartość na podstawie
ściśle mniejszych podformuł.

Jeżeli $A,e\in N$ i $N$ jest tranzytywnym modelem ZF,
to $W$, $\operatorname{Sub}(e)$ i $A^W$ są w $N$
i mają te same elementy co w uniwersum zewnętrznym:
funkcje o skończonej dziedzinie można zakodować przez
skończone zbiory par, a tranzytywność zapewnia
zgodność ich składników. Wszystkie wymienione
operacje logiczne i kwantyfikacja po $A$ są absolutne.
Indukcja po tym samym skończonym drzewie pokazuje,
że obie konstrukcje przypisują każdemu $(\psi,s)$
tę samą wartość. Wobec jednoznaczności tablice są równe.
\end{proof}
\begin{uwaga}[Granica zastosowania lematu o tablicy]
\label{rem:l-table-vs-stage-73}
Lemat~\ref{lem:l-satisfaction-table-transitive-73}
\emph{nie} pozwala zastąpić w jego założeniach
tranzytywnego modelu ZF dowolnym poziomem $L_\theta$
ani dowolną podstrukturą elementarną $(L_\theta,\in)$.
Taki poziom nie musi spełniać całej ZF, a istnienie
\emph{jednocześnie wszystkich} tablic dla nieograniczonej
rodziny kodów formuł nie wynika ze skończoności
\emph{jednej} tablicy. Dlatego w dowodzie ogólnego
lematu kondensacji nadal potrzebny jest odrębny
argument o jednolitych historiach poziomów $L$.
\end{uwaga}

\begin{uwaga}[Dlaczego kod formuły nie jest predykatem prawdy dla $L$]
\label{rem:l-finite-satisfaction-in-limit}
W Lemacie~\ref{lem:l-collapse-def-compatibility} symbol
$\operatorname{Sat}(L_\gamma,e,\vec p,z)$ opisuje
prawdziwość \emph{jednej formuły o kodzie $e$ w
zbiorowej strukturze $L_\gamma$}. Nie jest to
jednolity predykat prawdy dla klasy właściwej $L$.
Można formalnie obliczać prawdziwość tej jednej formuły
przez rekurencję po jej \emph{skończonym} drzewie
podformuł: w kroku atomowym używamy relacji $\in$,
a każdy kolejny krok polega na operacji logicznej
albo na kwantyfikowaniu po zbiorze $L_\gamma$.
Jeżeli kod $e$ jest ustalony, tablice wyników tej
skończonej rekurencji można zakodować w pewnym poziomie
$L_{\gamma+n(e)}$, gdzie $n(e)<\omega$ zależy od
budowy formuły. Dla granicznego $\theta>\gamma$
zachodzi $\gamma+n(e)<\theta$.
Tę \emph{skończoną poprawność} wykorzystuje argument
o kolapsie; nie zakładamy, że cała nieskończona
relacja spełniania dowolnych formuł jest elementem
$L_\theta$.
\end{uwaga}



\begin{lemat}[Lokalna tablica spełniania w granicznym poziomie $L$]
\label{lem:l-local-satisfaction-bound-76}
Niech $\varphi$ będzie \emph{ustaloną} formułą pierwszego
rzędu języka $\{\in\}$, a $W$ skończonym zbiorem
wszystkich jej zmiennych. Istnieje liczba naturalna
$n_\varphi>0$ zależna wyłącznie od składni $\varphi$
taka, że dla każdej liczby porządkowej $\gamma>0$
zbiór
\[
T_{\gamma,\varphi}=\bigl\{(\psi,s):
  \psi\text{ jest podformułą }\varphi,\ 
  s:W\to L_\gamma,\ (L_\gamma,\in)\models\psi[s]\bigr\}
\]
należy do $L_{\gamma+n_\varphi}$.
W szczególności, jeżeli $\theta>0$ jest graniczny
i $0<\gamma<\theta$, to $T_{\gamma,\varphi}\in L_\theta$.
Nie twierdzimy, że istnieje jedna liczba $n$
wspólna dla \emph{wszystkich} formuł.
\end{lemat}
\begin{proof}
Niech $A=L_\gamma$. Ponieważ $\gamma>0$, zbiór
$A$ jest niepusty i $A\in L_{\gamma+1}$.
Zbiór funkcji $A^W$ otrzymujemy ze skończonej
liczby iloczynów kartezjańskich i ograniczeń do
krotek reprezentujących funkcje o dziedzinie $W$.
Liczba potrzebnych operacji zależy od $|W|$, nie
od liczby elementów $A$. Kody zmiennych, podformuł
i ich drzewa składniowego są dziedzicznie skończone.

Budujemy osobno dla każdej podformuły $\psi$ relację
$S_\psi=\{s\in A^W:(A,\in)\models\psi[s]\}$.
Dla równości i przynależności relacja taka powstaje
przez przecięcie zbioru wartościowań z relacją
atomową. Dla negacji używamy różnicy
$A^W\setminus S_\chi$, a dla koniunkcji — przecięcia.
Jeśli $\psi=\exists v\,\chi$, to
\[
S_\psi=\{s\in A^W:\exists a\in A\ 
                   s[v\mapsto a]\in S_\chi\}.
\]
Zbiór par $(s,a)$ z taką własnością otrzymujemy
przez skończone operacje na produktach $A^W\times A$;
następnie rzutujemy na współrzędną $s$. Jest to
skończone złożenie operacji Gödla z
Lematu~\ref{lem:l-godel-normal-closure}.
W każdym kroku korzystamy wyłącznie z wyników
odpowiadających \emph{ściśle mniejszym} podformułom.

Każda z użytych operacji ma stałą realizację za pomocą
formuł ograniczonych, iloczynów, różnic i rzutów.
Jeżeli jej argumenty należą do $L_\delta$, to wynik
należy do pewnego $L_{\delta+c}$, gdzie stała skończona
$c$ zależy od \emph{rodzaju} operacji, a nie od
$\delta$ ani liczności jej argumentów. Istotnie,
$L_\delta$ jest tranzytywny, zaś odpowiednie pośrednie
zbiory otrzymuje się z $L_\delta$ przez ustaloną liczbę
operacji $\operatorname{Def}$; równoważnie można
wykorzystać normalną postać z przywołanego lematu.
W drzewie $\varphi$ jest tylko skończenie wiele węzłów.
Sumując liczbę operacji i koszt kodowania par
$(\psi,s)$, wybieramy jedno $n_\varphi$ tak, żeby
wszystkie potrzebne relacje i ich skończona suma
$T_{\gamma,\varphi}$ należały do $L_{\gamma+n_\varphi}$.

Dla granicznego $\theta$ i $\gamma<\theta$ każdy
skończony następnik $\gamma+n_\varphi$ jest
mniejszy od $\theta$, więc tablica należy do
$L_\theta$. Argument nie stosuje zastępowania
wewnątrz $L_\theta$ i nie zbiera w $L_\theta$
tablic wszystkich formuł naraz.
\end{proof}

\begin{przyklad}[Dwa kwantyfikatory i granica jednolitości]
\label{ex:l-local-satisfaction-76}
Niech $\varphi(u)$ będzie formułą
$\exists v\,(u\in v\land\forall w\,(w\in v\Rightarrow w=u))$.
Dla $A=L_\gamma$ i $\gamma>0$ budujemy najpierw
relację atomową $u\in v$ na $A^2$, następnie
relację $w\in v\Rightarrow w=u$ na $A^3$.
Jej dopełnienie i rzut względem $w$ opisują
kwantyfikator $\forall w$; przecięcie z relacją
$u\in v$ i rzut względem $v$ opisują
kwantyfikator $\exists v$. Każdy krok jest skończoną
operacją na zbiorach wartościowań. Poziom konstruktywny
wystarczający do uzyskania całej tablicy zależy
od składni tej jednej formuły, a nie od $\gamma$.
Zmiana $\varphi$ na formuły o rosnącej liczbie
zagnieżdżonych kwantyfikatorów zmienia ten koszt;
nie wolno więc wyciągać wniosku o jednej wspólnej
tablicy wszystkich formuł w dowolnym $L_\theta$.
\end{przyklad}


\begin{lemat}[Operacje Gödla i postać normalna definicji]
\label{lem:l-godel-normal-closure}
Dla zbiorów $A,B$ definiujemy następujące operacje
($G_6,G_7,G_8,G_9,G_{10}$ są jednoargumentowe):
\begin{align*}
G_1(A,B)&=\{A,B\}, &G_2(A,B)&=A\times B,\\
G_3(A,B)&=\{(u,v)\in A\times B:u\in v\},
 &G_4(A,B)&=A\setminus B,\\
G_5(A,B)&=A\cap B, &G_6(A)&=\bigcup A,\\
G_7(A)&=\operatorname{dom}(A),
 &G_8(A)&=\{(v,u):(u,v)\in A\},\\
G_9(A)&=\{(u,w,v):(u,v,w)\in A\},
 &G_{10}(A)&=\{(w,u,v):(u,v,w)\in A\}.
\end{align*}
Krotki w ostatnich dwóch wzorach są kodowane jednolicie.
Każda z relacji $B=G_i(\vec A)$ jest wyrażalna
formułą $\Delta_0$; w szczególności jej wartość jest
absolutna dla zbiorów tranzytywnych, w których istnieje.
Jeżeli $A$ jest tranzytywny, to
\[
\operatorname{Def}(A)
=\operatorname{cl}_{G_1,\ldots,G_{10}}(A\cup\{A\})
\cap\mathcal P(A),
\]
gdzie $\operatorname{cl}_{G_1,\ldots,G_{10}}$ oznacza
zbiór wszystkich wyników \emph{skończonych} złożeń tych
operacji z argumentami z podanego zbioru.
\end{lemat}
\begin{proof}
\emph{Absolutność.} Na przykład relację $z=A\setminus B$
wyraża koniunkcja
$\forall u\in z\,(u\in A\land u\notin B)$ oraz
$\forall u\in A\,(u\notin B\Rightarrow u\in z)$.
Dla iloczynu kartezjańskiego sprawdzamy, że każdy jego
element jest parą złożoną z elementów odpowiednich
argumentów i że wszystkie takie pary należą do wyniku.
Kwantyfikatory po parach ograniczamy do argumentów,
a po ich składnikach do sum i zbiorów par: takie
operacje, podobnie jak przestawienia współrzędnych,
mają opisy $\Delta_0$. Dla sumy używamy
$u\in\bigcup A\iff\exists a\in A\;(u\in a)$,
a dla dziedziny
$u\in\operatorname{dom}(A)\iff
\exists p\in A\;\exists v\in\bigcup\bigcup\{p\}\;(p=(u,v))$.
Pozostałe operacje są połączeniami tych samych wzorców.
Indukcja po długości złożenia zachowuje absolutność
każdej operacji otrzymanej w ten sposób.

\emph{Postać normalna.} Dowodzimy przez indukcję po
budowie formuły ograniczonej $\psi$,
że relację
$\{\vec u\in A_1\times\cdots\times A_n:\psi(\vec u)\}$
można otrzymać przez skończone złożenie operacji
$G_i$. Dla formuły atomowej $u_i\in u_j$ używamy
$G_3$ i przestawień współrzędnych; równość sprowadzamy
do dwóch inkluzji z ograniczonymi kwantyfikatorami.
Negacja odpowiada różnicy od całego iloczynu, koniunkcja
-- przecięciu, a kwantyfikator ograniczony
$\exists v\in u_i$ -- rzutowi odpowiedniej relacji
po uprzednim przecięciu z relacją $v\in u_i$.
Rzut realizuje $G_7$ po przestawieniu współrzędnych.
Wszystkie te operacje wymagają skończenie wielu kroków;
nie wprowadzamy kwantyfikatora po nieokreślonej klasie.

Dla dowolnej formuły pierwszego rzędu $\varphi$ jej
relatywizacja do $(A,\in)$ zastępuje wszystkie
kwantyfikatory przez kwantyfikatory ograniczone do $A$.
Otrzymujemy zatem, po dodaniu $A$ oraz parametrów,
wynik skończonego złożenia operacji $G_i$.
To dowodzi inkluzji $\subseteq$.
Dla inkluzji odwrotnej indukcja po długości złożenia
pokazuje, że każda operacja $G_i$ ma definicję
$\Delta_0$ z użytymi argumentami jako parametrami.
Jeśli jej końcowy wynik jest podzbiorem tranzytywnego
$A$, ograniczenie kwantyfikatorów do $A$ daje formułę
definiującą ten podzbiór nad $(A,\in)$.
Dla $A=\varnothing$ obie strony są równe
$\{\varnothing\}$ zgodnie z przyjętą konwencją.
\end{proof}


\begin{lemat}[Jednolity kod $\operatorname{Def}$ klasy $\Sigma_1$]
\label{lem:l-def-sigma1-certificate-75}
W ZF istnieje jedna formuła $\mathsf D(A,B)$ klasy
$\Sigma_1$ taka, że dla każdego \emph{tranzytywnego}
zbioru $A$ i dowolnego $B$ zachodzi
\[
\mathsf D(A,B)\quad\Longleftrightarrow\quad
B=\operatorname{Def}(A).
\]
Dla tranzytywnych argumentów graf funkcji
$A\mapsto\operatorname{Def}(A)$ jest też definiowalny
formułą $\Pi_1$, a zatem jest $\Delta_1$ \emph{nad ZF}.
\end{lemat}
\begin{proof}
Oznaczmy $U_0=A\cup\{A\}$ i zdefiniujmy
jednoznacznie, dla $n\in\omega$,
\[
U_{n+1}=U_n\cup
\{G_i(x,y):x,y\in U_n,\ 1\le i\le10\},
\]
przy czym operacje jednoargumentowe ignorują drugi
argument. Zbiór wszystkich użytych argumentów i wartości
jest objęty jednym $U_{n+1}$. W ZF funkcja
$W:n\mapsto U_n$ istnieje przez rekurencję po $\omega$.
Definiujemy $C=\bigcup_{n\in\omega}U_n$.
Z Lematu~\ref{lem:l-godel-normal-closure} wynika
\[
B=\{z\in C:z\subseteq A\}
  =\operatorname{Def}(A).
\]
Relację ``$W$ jest funkcją o dziedzinie $\omega$,
$W(0)=U_0$ i dla każdego $n\in\omega$ spełnia powyższe
równanie'' można zapisać formułą $\Delta_0$
\emph{z parametrami $W,A,\omega$}. W istocie
$z\in U_{n+1}$ oznacza: $z\in U_n$ albo
$z=G_i(x,y)$ dla pewnych $x,y\in U_n$ i skończenie
wielu indeksów $i$; wszystkie kwantyfikatory są
ograniczone do $U_n$, $U_{n+1}$ lub zbioru indeksów.
Własności $C=\bigcup\operatorname{ran}(W)$ oraz
$B=\{z\in C:z\subseteq A\}$ mają również zapisy
$\Delta_0$ (oba kierunki inkluzji sprawdza się
kwantyfikatorami po $C$ i $B$).
Predykat ``$w=\omega$'' można zapisać formułą
$\Delta_0$, charakteryzując $w$ jako pierwszy
niezerowy graniczny ordynał. Wobec tego zdanie
``istnieją $w,W,C$ spełniające powyższe warunki''
ma postać $\exists w\,\exists W\,\exists C\,\delta$
z $\delta$ klasy $\Delta_0$, czyli jest $\Sigma_1$.
Jednoznaczność $B$ wynika z jednoznaczności
$U_n$ i równości $B=C\cap\mathcal P(A)$.

Aby otrzymać równoważny opis $\Pi_1$,
zakładamy tranzytywność $A$ (warunek $\Delta_0$)
i zapisujemy, że \emph{każdy} zestaw świadków
$w,W,C,B'$ spełniający ograniczone warunki
musi mieć $B'=B$. Jest to formuła $\Pi_1$.
W ZF świadkowie zawsze istnieją, zatem warunek
powszechny jest równoważny poprzedniemu \emph{egzystencjalnemu}. Nie wynika z tego, że ich świadkowie
należą do dowolnego tranzytywnego podzbioru uniwersum.
\end{proof}

\begin{lemat}[Definiowalność $\Delta_1$ całej hierarchii]
\label{lem:l-hierarchy-delta1-derived-75}
W ZF istnieją ustalone formuły
$\Phi_{\Sigma}(\alpha,y)$ klasy $\Sigma_1$ i
$\Phi_{\Pi}(\alpha,y)$ klasy $\Pi_1$ takie, że
ZF dowodzi dla każdej liczby porządkowej $\alpha$:
\[
\Phi_{\Sigma}(\alpha,y)
\ \Longleftrightarrow\
\Phi_{\Pi}(\alpha,y)
\ \Longleftrightarrow\ y=L_\alpha.
\]
Ponadto dla każdego $\alpha$ istnieje dokładnie
jedna funkcja $h_\alpha$ o dziedzinie $\alpha+1$,
która koduje wszystkie etapy konstrukcji do $\alpha$.
\end{lemat}
\begin{proof}
Wybieramy \emph{jeden} predykat $\mathsf D$ z
Lematu~\ref{lem:l-def-sigma1-certificate-75}.
Dla funkcji $h$ o dziedzinie $\alpha+1$
sprawdzamy ograniczonymi kwantyfikatorami warunki:
$h(0)=\varnothing$, na każdym etapie granicznym
$h(\lambda)=\bigcup_{\xi<\lambda}h(\xi)$,
a na etapach następnikowych
$\mathsf D(h(\xi),h(\xi+1))$.
Każde użycie $\mathsf D$ ma świadków $w_\xi,W_\xi,C_\xi$
z poprzedniego lematu.
Ponieważ $\xi$ przebiega \emph{zbiór} $\alpha$,
a świadkowie są wyznaczeni przez kanoniczną rekurencję
po $\omega$, schemat zastępowania w ZF pozwala
zebrać ich rodziny w jedną funkcję $Q$ o dziedzinie
zbioru indeksów następnikowych $\xi<\alpha$.
Jest to zastosowanie zastępowania do
jednoznacznie określonych danych, nie aksjomatu wyboru.

Po włączeniu $Q$ jako dodatkowego świadka wszystkie
warunki dotyczące $h,Q$ można wypisać jako jedną
formułę $\Delta_0$: kwantyfikatory przebiegają
$\alpha$, $h$, $Q$, ich wartości lub zbiory kodujące
operacje, a równości etapów sprawdza się oboma
kierunkami inkluzji. Stąd zdanie
``istnieją $h,Q$ będące poprawnym kodem etapów
z $h(\alpha)=y$'' jest $\Sigma_1$.
Istnienie takich $h,Q$ dla wszystkich ordynałów
wynika z rekurencji pozaskończonej w ZF; jednoznaczność
$h$ wynika z indukcji po $\alpha$ i jednoznaczności
$\operatorname{Def}$. Dlatego równoważna w ZF jest
formuła powszechna: dla \emph{każdych} świadków
$h,Q$ spełniających te same warunki wartość
$h(\alpha)$ wynosi $y$. Ma ona postać $\Pi_1$.
Obie formuły definiują więc graf funkcji
$\alpha\mapsto L_\alpha$ nad ZF.
\end{proof}

\begin{przyklad}[Dlaczego $\Delta_1$ nad ZF nie wystarcza]
\label{ex:l-delta1-not-arbitrary-level-75}
Zbiór $L_\omega=V_\omega$ jest tranzytywny, lecz
nie zawiera $\omega$. Świadek $W$ w formule
$\mathsf D$ powyżej musi być funkcją o dziedzinie
$\omega$; zatem $W\notin L_\omega$, nawet jeśli
$A=\varnothing$ i $\operatorname{Def}(A)$ jest
skończony oraz należy do $L_\omega$.
W szczególności egzystencjalne kodowanie
$\Sigma_1$ \emph{nie jest automatycznie prawdziwe}
w każdym granicznym $L_\theta$.
Równoważność opisów $\Sigma_1$ i $\Pi_1$ jest
twierdzeniem \emph{ZF}, nie twierdzeniem o każdej
tranzytywnej strukturze niespełniającej ZF.
To jest konkretna przeszkoda, którą musi rozwiązać
skończony pakiet kondensacji.
\end{przyklad}

\begin{przyklad}[Absolutność operacji nie daje domknięcia]
\label{ex:l-delta0-not-closure-2026}
Niech $N=\{\varnothing\}$. Jest to niepusty zbiór
tranzytywny. Zdanie $z=G_1(a,b)$ jest $\Delta_0$
— przy argumentach i wartości należących do $N$
jego prawdziwość jest absolutna. Mimo to dla
$a=b=\varnothing\in N$ zachodzi
$G_1(a,b)=\{\varnothing\}\notin N$.
Absolutność \emph{grafu} operacji nie oznacza więc
istnienia jej \emph{wartości} w modelu.
Analogicznie sprawdzenie jednej skończonej tablicy
operacji nie dowodzi, że zbiór wszystkich potrzebnych
tablic jest elementem rozpatrywanego modelu.
\end{przyklad}

\begin{lemat}[Identyfikacja poziomu na podstawie poprawnych kodów]
\label{lem:l-external-codes-identify-2026}
Niech $N$ będzie niepustym zbiorem tranzytywnym,
$\beta=N\cap\operatorname{Ord}$ niezerową graniczną
liczbą porządkową i niech zachodzą dwa warunki:
\begin{enumerate}[label=(\roman*)]
\item dla każdego $\alpha<\beta$ rzeczywista,
  \emph{zewnętrznie obliczona} funkcja
  $h_\alpha=\langle L_\xi:\xi\leq\alpha\rangle$
  należy do $N$;
\item dla każdego $x\in N$ istnieje
  $\alpha<\beta$ takie, że $x\in L_\alpha$.
\end{enumerate}
Wówczas $N=L_\beta$.
\end{lemat}
\begin{proof}
Ustalmy $\alpha<\beta$. Na mocy (i) mamy
$h_\alpha\in N$, a ponieważ $N$ jest tranzytywny,
każda para uporządkowana należąca do $h_\alpha$
i każdy jej składnik należy do $N$. W szczególności
$L_\alpha\in N$, a z tranzytywności otrzymujemy
$L_\alpha\subseteq N$; po zsumowaniu tych
inkluzji po $\alpha<\beta$ mamy
$L_\beta\subseteq N$.
Z kolei (ii) daje każdemu $x\in N$ poziom
$L_\alpha$, do którego należy. Stąd
$N\subseteq\bigcup_{\alpha<\beta}L_\alpha=L_\beta$.
Obie inkluzje kończą dowód. W szczególności nie
utożsamialiśmy tu obiektów \emph{wewnętrznie
zadeklarowanych} jako kody z prawdziwymi kodami:
poprawność została jawnie założona w (i).
\end{proof}

\begin{definicja}[Adekwatność dla konstrukcji $L$]
\label{def:l-adequacy-2026}
Niech $N$ będzie zbiorem tranzytywnym.
Nazywamy go \emph{adekwatnym}, jeżeli:
\begin{enumerate}[label=(A\arabic*)]
\item $N$ jest domknięty na dziesięć operacji
  Gödla z Lematu~\ref{lem:l-godel-normal-closure};
\item dla każdego $U\in N$ zbiór wszystkich
  wartości $G_i(a,b)$ dla $a,b\in U$ i
  $1\leq i\leq10$ (z powtórzeniem argumentu
  dla operacji jednoargumentowych) należy do $N$;
\item dla każdej liczby porządkowej $\alpha\in N$
  \emph{zewnętrznie poprawny} początkowy odcinek
  $\langle L_\xi:\xi<\alpha\rangle$ należy do $N$.
\end{enumerate}
Warunek (A3) jest tu najpierw sformułowany
zewnętrznie, aby nie przemycić jego absolutności
przez samą definicję. Jednolitą formułę
$\mathsf{Hist}$ i pełny dowód (A3) podajemy w
Lemacie~\ref{lem:l-all-stage-histories-78}.
\end{definicja}

\begin{lemat}[Jednolity koszt operacji Gödla]
\label{lem:l-godel-finite-stage-bound-77}
Istnieje \emph{jedna} liczba naturalna $c>0$ taka,
że dla dowolnej liczby porządkowej $\eta$ i każdego
$i\in\{1,\ldots,10\}$, jeżeli wszystkie argumenty
operacji $G_i$ należą do $L_\eta$, to
\[
G_i(\vec a)\in L_{\eta+c}.
\]
Dla operacji jednoargumentowych nie wprowadza się
fikcyjnego dodatkowego argumentu. Dowód nie zakłada
aksjomatu wyboru ani istnienia kodu wszystkich poziomów
$L_\xi$ w jednym ustalonym modelu $L_\theta$.
\end{lemat}
\begin{proof}
Używamy własności udowodnionych przy definiowaniu
hierarchii: każdy $L_\eta$ jest tranzytywny,
$L_\eta\subseteq L_{\eta+1}$, a każdy podzbiór
$L_\rho$ definiowalny nad $(L_\rho,\in)$
z parametrami z $L_\rho$ należy do $L_{\rho+1}$.

Niech wszystkie argumenty leżą w $L_\eta$.
Każdy element argumentu, jak również każdy element
jego elementu, należy do $L_\eta$, dzięki
tranzytywności. Przy kodowaniu Kuratowskiego pary
$\langle u,v\rangle=\{\{u\},\{u,v\}\}$
dla $u,v\in L_\eta$ mamy
$\{u\},\{u,v\}\in L_{\eta+1}$, a więc
$\langle u,v\rangle\in L_{\eta+2}$.
Dla ustalonego kodowania trójek przez zagnieżdżone
pary otrzymujemy analogicznie
$\langle u,v,w\rangle\in L_{\eta+4}$.

Sprawdźmy teraz dziesięć typów operacji.
Elementami $G_1(a,b)$ są $a,b$;
elementami $G_2(a,b)$ i $G_3(a,b)$ są odpowiednie
pary o współrzędnych z argumentów;
elementy $G_4(a,b)$, $G_5(a,b)$ i $G_6(a)$
leżą w $L_\eta$; elementy $G_7(a)$ są
pierwszymi współrzędnymi par należących do $a$.
Elementy $G_8(a)$ są odwróconymi parami, a
$G_9(a)$ i $G_{10}(a)$ są przestawionymi
trójkami o współrzędnych z $L_\eta$.
W każdym z tych przypadków wszystkie elementy
wyniku należą do $L_{\eta+4}$.

Ponadto każda relacja $z\in G_i(\vec a)$ jest
wyrażalna \emph{jedną ustaloną} formułą $\Delta_0$
(Lemat~\ref{lem:l-godel-normal-closure}); jej
prawdziwość z parametrami należącymi do
$L_{\eta+4}$ jest absolutna między tym
tranzytywnym zbiorem a uniwersum.
Zatem $G_i(\vec a)$ jest dokładnie podzbiorem
$L_{\eta+4}$ definiowanym nad
$(L_{\eta+4},\in)$ przez wspomnianą formułę.
Wynika stąd
$G_i(\vec a)\in L_{\eta+5}$.
Można więc przyjąć wspólną stałą $c=6$,
niezależną od $i$, $\eta$ i argumentów.
\end{proof}

\begin{lemat}[Dwa autonomiczne warunki adekwatności]
\label{lem:l-limit-godel-adequacy-77}
Dla każdej niezerowej granicznej liczby porządkowej
$\theta$ zbiór $L_\theta$ spełnia warunki
\emph{(A1)} i \emph{(A2)}
Definicji~\ref{def:l-adequacy-2026}.
Dokładniej, dla każdego $U\in L_\theta$ zbiór
\[
T(U)=\{G_i(a,b):a,b\in U,\ 1\le i\le10\}
\]
(z umową dotyczącą operacji jednoargumentowych)
należy do $L_\theta$.
Sama ta konkluzja nie stwierdza warunku (A3).
\end{lemat}
\begin{proof}
Najpierw weźmy dowolne argumenty $a,b\in L_\theta$.
Istnieje $\eta<\theta$ takie, że $a,b\in L_\eta$:
wybieramy wspólny etap dla skończonej rodziny
argumentów, a następnie, jeżeli trzeba,
zwiększamy go o jeden. Lemat~
\ref{lem:l-godel-finite-stage-bound-77}
daje $G_i(\vec a)\in L_{\eta+6}\subseteq L_\theta$,
gdyż dla granicznej $\theta$ i $\eta<\theta$
zachodzi $\eta+6<\theta$. To dowodzi (A1).

Dla (A2) ustalmy $U\in L_\theta$ i wybierzmy
$\eta<\theta$ z $U\in L_\eta$. Tranzytywność
$L_\eta$ implikuje $U\subseteq L_\eta$.
Stąd, na mocy poprzedniego lematu,
każda wartość $G_i(a,b)$ dla $a,b\in U$
należy do $L_{\eta+6}$, czyli
$T(U)\subseteq L_{\eta+6}$.

Zbiór $T(U)$ jest definiowalny nad
$(L_{\eta+6},\in)$ z parametrem $U$ za pomocą
jednej formuły:
\[
 z\in T(U)\ \longleftrightarrow\
 \bigvee_{i=1}^{10}\ \exists a,b\in U\
                 \bigl(z=G_i(a,b)\bigr).
\]
Każdy graf $G_i$ zastępujemy przy tym jego
formułą $\Delta_0$, a wszystkie kwantyfikatory
po argumentach są ograniczone do $U$.
Ze względu na wcześniejszą inkluzję oraz
absolutność tych grafów formuła definiuje
\emph{dokładnie} $T(U)$, nie tylko jego fragment.
Wobec tego
$T(U)\in L_{\eta+7}\subseteq L_\theta$.
Nie korzystaliśmy z zastępowania wewnątrz
$L_\theta$: zbiór wartości został otrzymany
jako definiowalny podzbiór wcześniejszego poziomu.
\end{proof}

\begin{lemat}[Warunek historii: indeksy skończone i $\omega$]
\label{lem:l-initial-history-omega-77}
Niech $\theta>0$ będzie graniczną liczbą porządkową.
Dla każdej $n<\omega$ rzeczywisty zbiór
$H_n=\langle L_k:k<n\rangle$ należy do $L_\theta$.
Jeżeli $\omega\in L_\theta$, to dodatkowo
$H_\omega=\langle L_n:n<\omega\rangle\in L_\theta$.
\end{lemat}
\begin{proof}
Dla $n<\omega$ wszystkie $L_k=V_k$ są zbiorami
dziedzicznie skończonymi; stąd skończona funkcja
$H_n$ należy do $L_\omega=V_\omega\subseteq L_\theta$.

Dla drugiej części zauważmy najpierw, że
$H_\omega\subseteq L_\omega$:
każda para $\langle n,V_n\rangle$ jest
zbiorem dziedzicznie skończonym.
Pokażemy, że jest to podzbiór
\emph{definiowalny} nad $(L_\omega,\in)$.
Dla $n\in\omega$ formuła $\rho(n,y)$ stwierdza
istnienie skończonej funkcji $f$ o dziedzinie
$n+1$, dla której $f(0)=\varnothing$,
$f(k+1)=\mathcal P(f(k))$ dla każdego $k<n$
oraz $f(n)=y$. Wystąpienie $\mathcal P$ jest
tu skrótem jednej formuły pierwszego rzędu
\[
\forall z\,\bigl(z\in f(k+1)\longleftrightarrow
                 z\subseteq f(k)\bigr).
\]
Ponieważ $f(k)$ jest skończony, każdy jego
podzbiór należy do $L_\omega$; wszystkie
funkcje $f$ potrzebne dla skończonego $n$
również należą do $L_\omega$.
Z tego wynika, że $\rho(n,y)$ jest prawdziwe
w $(L_\omega,\in)$ dokładnie dla $y=V_n$.
Zatem nad $L_\omega$ formuła
\[
\exists n,y\,\bigl(\rho(n,y)\land
  z=\langle n,y\rangle\bigr)
\]
definiuje podzbiór $H_\omega$ zbioru $L_\omega$.
W rezultacie
$H_\omega\in\operatorname{Def}(L_\omega)
=L_{\omega+1}$.
Jeśli $\omega\in L_\theta$, to $\omega<\theta$
i z graniczności $\theta$ wynika
$\omega+1<\theta$, a więc $H_\omega\in L_\theta$.
\end{proof}

\begin{uwaga}[Co brakowało w audycie LXXVII]
\label{rem:l-adequacy-remaining-history-77}
Dowód warunków (A1) i (A2) sam nie pozwala zebrać
historii na etapie granicznym: z przynależności
$H_\xi\in L_\theta$ dla każdego $\xi<\alpha$ nie wynika,
że $\{H_\xi:\xi<\alpha\}\in L_\theta$.
Następne dwa lematy usuwają tę lukę bez zakładania
schematu zastępowania w $L_\theta$. Zamiast zbierać
wartości funkcji na całej $\alpha$, definiujemy
$H_\alpha$ jako \emph{jeden} podzbiór $L_\alpha$;
kluczowa jest absolutność skończonych kodów operacji.
\end{uwaga}

\begin{lemat}[Absolutność skończonych kodów operacji Gödla]
\label{lem:l-finite-godel-term-absoluteness-78}
Istnieją ustalone formuły pierwszego rzędu
$\mathsf{Eval}_G(t,A,\vec p,z)$ oraz
$\mathsf{Def}_G(A,B)$ o następujących własnościach.
Niech $N$ będzie \emph{niepustym tranzytywnym} zbiorem
zamkniętym na operacje $G_1,\ldots,G_{10}$.
Dla każdego tranzytywnego $A\in N$ i $B\in N$ mamy
\[
 (N,\in)\models\mathsf{Def}_G(A,B)
 \quad\Longleftrightarrow\quad B=\operatorname{Def}(A).
\]
Formuła $\mathsf{Eval}_G$ opisuje wartość dowolnego
\emph{skończonego} termu zbudowanego z $G_i$, stałej
$A$ i skończonej krotki parametrów z $A$; jest ona
absolutna między $(N,\in)$ i uniwersum dla tych
argumentów. Nie zakładamy, że $N$ spełnia ZF.
\end{lemat}
\begin{proof}
Z niepustości i tranzytywności $N$ wynika
$\varnothing\in N$: wybieramy element o najmniejszej
randze i stosujemy tranzytywność. Każdy zbiór
dziedzicznie skończony należy do $N$. Rzeczywiście,
jeżeli wszystkie elementy skończonego zbioru $u$
należą do $N$, to wielokrotne zastosowanie operacji
$G_1(a,b)=\{a,b\}$ daje $u\in N$ (dla $u=\varnothing$
korzystamy z poprzedniego zdania). W szczególności
$N$ zawiera wszystkie \emph{standardowe} kody formuł,
termów i skończonych krotek.

Kod termu wybieramy jako skończony \emph{program acykliczny}
$t=(n,c)$: $n$ jest dodatnią skończoną liczbą
porządkową, a funkcja $c$ o dziedzinie $n$
w każdej pozycji $j<n$ wskazuje jedną z instrukcji:
``wstaw $A$'', ``wstaw parametr numer $k$''
albo ``zastosuj $G_i$ do wyników instrukcji
o numerach mniejszych od $j$''.
Dla $G_6,\ldots,G_{10}$ instrukcja ma jeden
argument, a dla pozostałych dwa. Skończoność
$n$ można wyrazić \emph{bez używania $\omega$ jako
parametru}: $n$ jest liczbą porządkową i ani $n$,
ani żaden niezerowy $\xi\le n$ nie jest graniczny.
Jest to warunek absolutny w rozważanych zbiorach
tranzytywnych zawierających wszystkie zbiory
dziedzicznie skończone.

Formuła $\mathsf{Eval}_G(t,A,\vec p,z)$ mówi, że
istnieje funkcja $v$ o dziedzinie $n$, której wartość
$v(j)$ jest wskazaną stałą, parametrem albo
$G_i(v(j_1),v(j_2))$ zgodnie z instrukcją $c(j)$,
przy czym $j_1,j_2<j$, oraz $v(n-1)=z$.
Warunki dla $v$ są jedną formułą, ponieważ indeks
$i$ przebiega ustalony skończony zbiór i każdy
graf $G_i$ ma ustalony opis $\Delta_0$ na mocy
Lematu~\ref{lem:l-godel-normal-closure}.
Dla skończonej krotki parametrów $\vec p$
kod ten ma tylko skończenie wiele instrukcji;
nie kwantyfikujemy po dowolnej klasie operacji.
Ponieważ $N$ jest domknięty na wszystkie
$G_i$, rzeczywiste wartości wszystkich węzłów oraz
cała ich skończona tablica należą do $N$. Odwrotnie,
każda tablica uznana w $N$ jest poprawna zewnętrznie:
wykorzystujemy absolutność grafów $G_i$ i indukcję
po skończonym drzewie. Otrzymujemy żądaną
absolutność $\mathsf{Eval}_G$.

Formuła $\mathsf{Def}_G(A,B)$ stwierdza, że $A$ jest
tranzytywny, każdy element $B$ jest podzbiorem $A$
i jest wynikiem pewnego takiego termu oraz że
\emph{każda} wartość termu z argumentami w
$A\cup\{A\}$, będąca podzbiorem $A$, należy do $B$.
Kwantyfikacja po kodach termów jest zwykłą
kwantyfikacją pierwszego rzędu; warunek, że kod jest
skończonym poprawnym termem, wyrażamy przez relację
jego skończonej tablicy. Dla dowolnego termu i
krotki parametrów z $A$ ich kody należą do $N$,
a wartość termu również należy do $N$. Wobec tego
oba kierunki warunku definiującego $B$ są absolutne.
Z postaci normalnej Gödla
(Lemat~\ref{lem:l-godel-normal-closure}) wynika, że
zewnętrznie $B$ jest dokładnie $\operatorname{Def}(A)$.
Nie twierdzimy, że $\operatorname{Def}(A)\in N$
\emph{dla każdego} $A\in N$: lemat mówi o poprawności
relacji, jeżeli rozpatrywana wartość $B$ już należy
do $N$. To rozróżnienie będzie ważne w dowodzie
następnikowym.
\end{proof}

\begin{lemat}[Wszystkie historie hierarchii mają skończony koszt]
\label{lem:l-all-stage-histories-78}
Dla każdej liczby porządkowej $\alpha$ istnieje
$n_\alpha\in\omega\setminus\{0\}$ takie, że rzeczywista
historia $H_\alpha=\langle L_\xi:\xi<\alpha\rangle$
należy do $L_{\alpha+n_\alpha}$. Co więcej, jeżeli
$\alpha>0$ jest graniczna, to
\[
 H_\alpha\in L_{\alpha+1}.
\]
W konsekwencji dla \emph{każdej} niezerowej granicznej
liczby porządkowej $\theta$ i każdego $\alpha<\theta$
zachodzi $H_\alpha\in L_\theta$, czyli warunek (A3).
\end{lemat}
\begin{proof}
Prowadzimy indukcję pozaskończoną po $\alpha$.
Dla $\alpha=0$ historia jest funkcją pustą, więc
$H_0=\varnothing\in L_1$.

Jeżeli $\alpha=\beta+1$, to
$H_\alpha=H_\beta\cup\{\langle\beta,L_\beta\rangle\}$.
Na mocy indukcji $H_\beta$ należy do pewnego
$L_{\beta+n_\beta}$, natomiast $\beta$ i $L_\beta$
należą do $L_{\beta+1}$. Utworzenie pary
uporządkowanej, singletonu i sumy jest skończonym
złożeniem operacji Gödla. Lemat~
\ref{lem:l-godel-finite-stage-bound-77}
i jego dowód dają zatem pewne $m<\omega$ takie, że
$H_{\beta+1}\in L_{\beta+n_\beta+m}$.
Wybierając dostatecznie duże skończone
$n_{\beta+1}>0$, otrzymujemy tezę dla następnika.

Niech teraz $\alpha>0$ będzie graniczna i
załóżmy tezę dla wszystkich $\xi<\alpha$.
Dla $\xi<\alpha$ także $\xi+1<\alpha$, a więc
$H_{\xi+1}\in L_{\xi+1+n_{\xi+1}}\subseteq L_\alpha$.
Ta ostatnia inkluzja zachodzi, ponieważ
$\alpha$ jest graniczna i dodanie skończonej liczby
do liczby porządkowej mniejszej niż $\alpha$
pozostawia wynik poniżej $\alpha$.
Podobnie $\langle\xi,L_\xi\rangle\in L_\alpha$:
$\xi,L_\xi\in L_{\xi+1}$ i kodowanie pary wymaga
tylko skończenie wielu dalszych poziomów.
Dlatego $H_\alpha\subseteq L_\alpha$.

Zdefiniujmy jedną formułę $\mathsf{Hist}(h,\delta)$:
``$h$ jest funkcją o dziedzinie $\delta$,
$h(0)=\varnothing$ gdy $\delta>0$,
$h(\eta+1)$ spełnia $\mathsf{Def}_G(h(\eta),h(\eta+1))$
a na etapach granicznych $h(\lambda)$ jest sumą
wcześniejszych wartości''. W definicji wymagamy
tranzytywności wszystkich poziomów $h(\eta)$.
Jest to jedna formuła pierwszego rzędu: na każdym
etapie kwantyfikujemy tylko po indeksach, wartościach
$h$ i sprawdzamy ustaloną formułę $\mathsf{Def}_G$.
Dla dowolnej tranzytywnej struktury domkniętej na
operacje Gödla, w szczególności dla
$M=L_\alpha$, każdy kod $h\in M$ spełniający
tę formułę jest \emph{zewnętrznie poprawny}.
Dowodzimy tego indukcją po $\eta<\delta$:
krok zerowy i graniczny są absolutne, a krok
następnikowy wynika z poprzedniego lematu.
To dowodzi także jednoznaczności kodu.

Ponieważ $L_\alpha$ jest domknięty na operacje Gödla,
formuła
\[
 \exists\xi\,\exists h\;\bigl(
  \operatorname{Ord}(\xi)\land
  \mathsf{Hist}(h,\xi+1)\land
  z=\langle\xi,h(\xi)\rangle\bigr)
\]
rozpatrywana w $(L_\alpha,\in)$ definiuje
\emph{dokładnie} $H_\alpha$.
W istocie liczby porządkowe należące do $L_\alpha$
to dokładnie $\xi<\alpha$. Dla każdego z nich
indukcja dostarcza prawdziwy kod $H_{\xi+1}\in L_\alpha$,
więc żaden wymagany element $H_\alpha$ nie jest
pominięty. Odwrotnie, każdy kod $h$ uznany w
$L_\alpha$ jest poprawny zewnętrznie na mocy
absolutności $\mathsf{Def}_G$, więc nie daje
żadnego fałszywego elementu. Ponieważ
$H_\alpha\subseteq L_\alpha$ i jest definiowalny
nad $(L_\alpha,\in)$, mamy
$H_\alpha\in\operatorname{Def}(L_\alpha)=L_{\alpha+1}$.
To kończy krok graniczny.

Na koniec, gdy $\alpha<\theta$ i $\theta$ jest
niezerową graniczną liczbą porządkową,
$\alpha+n_\alpha<\theta$. Zatem
$H_\alpha\in L_{\alpha+n_\alpha}\subseteq L_\theta$.
Warunek (A3) jest udowodniony bez schematu
zastępowania wewnątrz $L_\theta$.
\end{proof}

\begin{lemat}[Redukcja (A3) do poprawności kodu]
\label{lem:l-adequacy-to-external-codes-2026}
Niech $N$ będzie tranzytywnym zbiorem adekwatnym,
$\beta=N\cap\operatorname{Ord}$ niezerową graniczną
liczbą porządkową i załóżmy, że
\begin{enumerate}[label=(\roman*)]
\item dla każdego $x\in N$ istnieje liczba
  porządkowa $\alpha\in N$, dla której $N$
  stwierdza $x\in L_\alpha$;
\item za pomocą warunków (A1)--(A3) poprawność
  wewnętrznego kodu $\langle L_\xi:\xi<\alpha\rangle$
  implikuje jego zgodność z kodem zewnętrznym,
  a definicja $x\in L_\alpha$ jest absolutna
  dla $N$ i metateorii.
\end{enumerate}
Wtedy $N=L_\beta$.
\end{lemat}
\begin{proof}
Dla każdego $\alpha<\beta$, ponieważ $\beta$ jest
graniczny, mamy $\alpha+1<\beta$. Warunek (A3)
zastosowany do $\alpha+1$ dostarcza w $N$ kodu
$\langle L_\xi:\xi\leq\alpha\rangle$.
Warunek (ii) identyfikuje ten kod z zewnętrznym
obiektem; stąd warunek (i) Lematu
\ref{lem:l-external-codes-identify-2026}.
Dla dowolnego $x\in N$ wybieramy $\alpha$ z (i).
Na mocy (ii) jego wewnętrzne świadectwo
$x\in L_\alpha$ jest prawdziwe zewnętrznie.
Otrzymujemy również drugą hipotezę tego lematu,
a więc $N=L_\beta$.
\end{proof}

\begin{lemat}[Skończony pakiet absolutności dla poziomów $L$]
\label{lem:l-finite-condensation-package}
Istnieje jedna \emph{ustalona skończona} koniunkcja
zdań $\sigma_*$ języka teorii mnogości taka, że:
\begin{enumerate}[label=(\roman*)]
\item dla każdej niezerowej granicznej liczby
  porządkowej $\theta$ zachodzi
  $(L_\theta,\in)\models\sigma_*$;
\item jeżeli $N$ jest niepustym zbiorem tranzytywnym,
  $(N,\in)\models\sigma_*$ i
  $\beta=N\cap\operatorname{Ord}$, to $\beta$ jest
  graniczny i $N=L_\beta$.
\end{enumerate}
\end{lemat}
\begin{proof}
Wykorzystujemy ustalone formuły
$\mathsf{Def}_G$ i $\mathsf{Hist}$ z Lematów
\ref{lem:l-finite-godel-term-absoluteness-78} i
\ref{lem:l-all-stage-histories-78}.
Wszystkie poniższe warunki są \emph{pojedynczymi}
zdaniami pierwszego rzędu, nie schematami:
\begin{enumerate}[label=(\arabic*)]
\item domknięcie na dziesięć operacji Gödla:
 dla wszystkich argumentów istnieją ich wartości,
 opisane dziesięcioma ustalonymi grafami $\Delta_0$;
\item każda liczba porządkowa ma w strukturze
 większą liczbę porządkową;
\item dla każdej liczby porządkowej $\alpha$
 istnieje $h$ spełniające
 $\mathsf{Hist}(h,\alpha+1)$;
\item dla każdego $x$ istnieje liczba porządkowa
 $\alpha$ i $h$ spełniające
 $\mathsf{Hist}(h,\alpha+1)$ oraz $x\in h(\alpha)$.
\end{enumerate}
Koniunkcję tych zdań oznaczamy przez $\sigma_*$.
Można do niej dodać warunek (A2), ponieważ
Lemat~\ref{lem:l-limit-godel-adequacy-77}
dowodzi go dla wszystkich rozpatrywanych poziomów;
nie jest on jednak potrzebny do samej identyfikacji
kolapsu.

Niech $\theta>0$ będzie graniczna. Warunek (1)
wynika z Lematu~\ref{lem:l-limit-godel-adequacy-77}.
Warunek (2) wynika z graniczności $\theta$ i
równości $L_\theta\cap\operatorname{Ord}=\theta$.
Dla $\alpha<\theta$ mamy
$H_{\alpha+1}\in L_\theta$ na mocy
Lematu~\ref{lem:l-all-stage-histories-78}, a
absolutność kodu z tego samego lematu zapewnia
$(L_\theta,\in)\models\mathsf{Hist}(H_{\alpha+1},\alpha+1)$.
Stąd (3). Dla $x\in L_\theta$ wybieramy
$\alpha<\theta$ takie, że $x\in L_\alpha$,
co jest możliwe, ponieważ
$L_\theta=\bigcup_{\xi<\theta}L_\xi$.
Ten sam kod $H_{\alpha+1}$ jest świadkiem (4).
Wykazaliśmy (i).

Niech teraz $N$ spełnia założenia (ii).
Ponieważ $N$ jest niepusty i tranzytywny, zawiera
$\varnothing$, a warunek (2) pokazuje, że
$\beta=N\cap\operatorname{Ord}$ jest niezerową
graniczną liczbą porządkową. Na mocy (1) i
Lematu~\ref{lem:l-finite-godel-term-absoluteness-78}
każdy kod $h\in N$ spełniający w $N$ formułę
$\mathsf{Hist}$ jest \emph{zewnętrznie} prawdziwym
odcinkiem hierarchii $L$; wynika to z indukcji
po dziedzinie kodu. Z (3) otrzymujemy dla każdego
$\alpha<\beta$ rzeczywistą historię
$H_{\alpha+1}\in N$. Ponieważ $N$ jest tranzytywny,
$L_\alpha\in N$, a zatem $L_\alpha\subseteq N$.
Po zsumowaniu po $\alpha<\beta$ mamy
$L_\beta\subseteq N$.
Z kolei dla dowolnego $x\in N$ warunek (4)
dostarcza prawdziwego kodu historii i
$\alpha<\beta$ z $x\in L_\alpha$. Stąd
$N\subseteq L_\beta$. Otrzymujemy $N=L_\beta$.
Cały dowód jest wewnętrzny: opiera się na
skończonym kodowaniu operacji, indukcji
pozaskończonej w metateorii ZF i definicji $L$.
\end{proof}

\begin{uwaga}[Zakres uzyskanego dowodu]
\label{rem:l-finite-package-source}
Skończoność $\sigma_*$ wynika z tego, że
$\mathsf{Eval}_G$, $\mathsf{Def}_G$ i $\mathsf{Hist}$
są \emph{ustalonymi} formułami języka teorii
mnogości. Kwantyfikowanie po wszystkich
skończonych kodach termów nie oznacza użycia
predykatu prawdy dla całego uniwersum.
W przeciwieństwie do starszej wersji rozdziału
warunek historii (A3) jest tu udowodniony
(Lemat~\ref{lem:l-all-stage-histories-78}), a nie
przyjmowany jako wynik zewnętrzny.
Zasadnicza zależność tego dowodu to już
udowodniona w rozdziale postać normalna Gödla
(Lemat~\ref{lem:l-godel-normal-closure}).
\end{uwaga}

\begin{lemat}[Dziedzicznie skończony początek kadłuba]
\label{lem:l-hf-fixed-collapse-76}
Niech $\theta>0$ będzie graniczną liczbą porządkową
i $X\preccurlyeq(L_\theta,\in)$.
Wówczas $L_\omega\subseteq X$, a kolaps Mostowskiego
$\pi:X\to N$ spełnia $\pi(a)=a$ dla każdego
$a\in L_\omega$. W szczególności dla $\theta=\omega$
mamy $X=L_\omega$ i $N=L_\omega$.
\end{lemat}
\begin{proof}
Dla dowolnego dziedzicznie skończonego zbioru $a$
konstruujemy przez indukcję po randze formułę
$\delta_a(x)$ bez parametrów, która jednoznacznie
określa $a$ w każdej tranzytywnej strukturze
zawierającej $a$ i jego elementy. Dla $a=\varnothing$
można wziąć $\forall z\,(z\notin x)$.
Jeżeli $a=\{a_0,\ldots,a_{k-1}\}$, to przyjmujemy
\[
\delta_a(x)\equiv\forall z\,
 \bigl(z\in x\longleftrightarrow
  \bigvee_{i<k}\delta_{a_i}(z)\bigr),
\]
a dla $k=0$ zastępujemy alternatywę formułą
fałszywą. Indukcja uzasadnia jednoznaczność:
w tranzytywnym $L_\theta$ każde $a_i$ jest
jedynym rozwiązaniem $\delta_{a_i}$, więc
ekstensjonalność wyznacza jedyny możliwy zbiór $x=a$.
Ponieważ $L_\omega\subseteq L_\theta$ i
$X\preccurlyeq L_\theta$, jedyny świadek zdania
$\exists!x\,\delta_a(x)$ należy do $X$.
Stąd $L_\omega\subseteq X$.

Przez indukcję po randze $a\in L_\omega$
obliczamy kolaps:
$\pi(a)=\{\pi(b):b\in a\cap X\}$.
Wszystkie elementy $b\in a$ są dziedzicznie
skończone, zatem należą do $X$ i przez indukcję
$\pi(b)=b$. Wobec tego $\pi(a)=a$.
Dla $\theta=\omega$ inkluzja
$L_\omega\subseteq X\subseteq L_\omega$
daje równość; w tym przypadku nie potrzebujemy
jednolitego kodowania poziomów konstruowalnych.
\end{proof}

\begin{lemat}[Klasyczny lemat kondensacji Gödla]
\label{lem:l-condensation}
Niech $\theta>0$ będzie graniczną liczbą porządkową,
a $X\preccurlyeq(L_\theta,\in)$ podstrukturą elementarną.
Wówczas kolaps Mostowskiego $(X,\in)$ jest izomorficzny
z $(L_\beta,\in)$ dla pewnej granicznej liczby
porządkowej $\beta\le\theta$.
\end{lemat}
\begin{proof}
Rozważmy $(X,\in)$ z tezy i jego kolaps Mostowskiego
$\pi:X\to N$. Jest on dobrze ufundowany jako
podrelacja rzeczywistej przynależności na zbiorze
$L_\theta$, a ekstensjonalność wynika z
$X\preccurlyeq L_\theta$: formuła odróżniająca dwa
różne zbiory przez element ma świadka w $X$.
Zatem $N$ jest tranzytywnym zbiorem, a $\pi$ jest
izomorfizmem struktur.

Na mocy Lematu~\ref{lem:l-finite-condensation-package}(i)
struktura $L_\theta$ spełnia jedno ustalone zdanie
$\sigma_*$. Pełna elementarność daje
$X\models\sigma_*$, a zachowanie prawdziwości zdań
przez izomorfizm daje $N\models\sigma_*$.
Lemat~\ref{lem:l-finite-condensation-package}(ii)
dostarcza granicznej liczby porządkowej $\beta$
z $N=L_\beta$.

Aby wykazać $\beta\leq\theta$, przypomnijmy,
że dla liczby porządkowej $\gamma\in X$ wartość
$\pi(\gamma)$ jest typem porządkowym zbioru
$X\cap\gamma$, więc ordynały obrazu kolapsu
tworzą podzbiór $\theta$ w sensie porządkowym.
Dokładniej, funkcja $\pi$ ograniczona do
$X\cap\operatorname{Ord}$ jest ściśle rosnąca,
a jej obraz to $\beta$. Ponieważ
$X\cap\operatorname{Ord}\subseteq\theta$,
niemożliwe jest $\beta>\theta$.
Wynika stąd teza.
\end{proof}
\begin{uwaga}[Niezależna kontrola kroku kolapsu]
\label{rem:lvi-condensation-dependency}
Ogólny Lemat~\ref{lem:l-condensation} jest teraz
konsekwencją \emph{udowodnionego wewnętrznie}
skończonego pakietu z
Lematu~\ref{lem:l-finite-condensation-package}.
Nie wykorzystujemy nieuzasadnionej równości
$\pi[X\cap L_\gamma]=L_{\pi(\gamma)}$;
jednolitość kodu jest zapewniona przez
Lematy~\ref{lem:l-finite-godel-term-absoluteness-78}
i~\ref{lem:l-all-stage-histories-78}.
Alternatywny dowód dla kadłuba w języku
rozszerzonym pozostaje w
Lemacie~\ref{lem:l-expanded-hull-condensation};
nie jest już potrzebny do zamknięcia samej kondensacji.
\end{uwaga}

\begin{twierdzenie}[Kondensacja z jednolitego rozpoznawania poziomów]
\label{thm:l-condensation-from-uniform-recognition-75}
Niech $\theta>0$ będzie graniczną liczbą porządkową.
Załóżmy, że istnieje \emph{jedna ustalona formuła}
$\Lambda(\gamma,y)$ języka teorii mnogości taka, że
\begin{enumerate}[label=(\roman*)]
\item dla wszystkich $\gamma<\theta$ mamy
$L_\gamma\in L_\theta$;
\item dla wszystkich liczb porządkowych
$\gamma<\theta$ i $y\in L_\theta$ zachodzi
\[
(L_\theta,\in)\models\Lambda(\gamma,y)
\quad\Longleftrightarrow\quad y=L_\gamma.
\]
\end{enumerate}
Wówczas dla \emph{każdego}
$X\preccurlyeq(L_\theta,\in)$ jego kolaps Mostowskiego
jest dokładnie $(L_\beta,\in)$ dla pewnej
\emph{granicznej} liczby porządkowej $\beta\le\theta$.
\end{twierdzenie}
\begin{proof}
Niech $\pi:X\to N$ będzie kolapsem i
$\beta=N\cap\operatorname{Ord}$.
Elementarność zapewnia $0\in X$ oraz
$\gamma+1\in X$ dla każdej liczby porządkowej
$\gamma\in X$, ponieważ $0$ i operacja następnika
są jednoznacznie definiowalne w $L_\theta$.
Kolaps zachowuje następnik ordynalny, więc $\beta$
jest niezerową liczbą porządkową bez elementu
największego, a zatem jest graniczny.
Dla każdej liczby porządkowej $\gamma\in X$
prawdziwe jest w $L_\theta$ zdanie
$\exists!y\,\Lambda(\gamma,y)$, którego jedynym
świadkiem jest $L_\gamma$. Z elementarności otrzymujemy
$L_\gamma\in X$.

\emph{Indukcja po liczbach porządkowych z $X$.}
Pokażemy, że
\[
\pi(L_\gamma)=L_{\pi(\gamma)}
\qquad(\gamma\in X\cap\operatorname{Ord}).
\]
Dla $\gamma=0$ mamy $L_0=\varnothing$ i $\pi(0)=0$.
Jeżeli $\gamma=\eta+1$, to $\eta\in X$, bo jest
jednoznacznie definiowalnym poprzednikiem
$\gamma$ w $L_\theta$.
Ponieważ $\pi(\gamma)=\pi(\eta)+1$, hipoteza
indukcyjna i Lemat~\ref{lem:l-collapse-def-compatibility}
dają
\begin{align*}
\pi(L_\gamma)
 &=\pi[X\cap\operatorname{Def}(L_\eta)]\\
 &=\operatorname{Def}(L_{\pi(\eta)})
  =L_{\pi(\gamma)}.
\end{align*}
Pierwsza równość wynika z definicji kolapsu
zastosowanej do $L_\gamma\in X$.

Niech teraz $\lambda\in X$ będzie niezerową
graniczną liczbą porządkową.
Dla każdego $z\in X\cap L_\lambda$ w $L_\theta$
zachodzi formuła
``$\exists\xi<\lambda\,\exists y\,
(\Lambda(\xi,y)\land z\in y)$''.
Elementarność daje świadków $\xi,y\in X$;
na mocy (ii) koniecznie $y=L_\xi$.
Otrzymujemy
\[
\pi(L_\lambda)=
 \bigcup_{\xi\in X\cap\lambda}\pi(L_\xi).
\]
Równość odwrotnej inkluzji wynika z
$L_\xi\subseteq L_\lambda$ dla $\xi<\lambda$.
Ponieważ $\pi(\lambda)=\{\pi(\xi):\xi\in X\cap\lambda\}$
i stosujemy hipotezę indukcyjną, prawa strona
jest $\bigcup_{\zeta<\pi(\lambda)}L_\zeta
=L_{\pi(\lambda)}$.

\emph{Identyfikacja całego kolapsu.}
Każde $z\in L_\theta$ należy do pewnego $L_\gamma$
z $\gamma<\theta$. Zatem w $L_\theta$ jest prawdziwe
zdanie $\forall z\,\exists\gamma\,\exists y
(\operatorname{Ord}(\gamma)\land\Lambda(\gamma,y)
\land z\in y)$.
Dla $z\in X$ świadkowie $\gamma,y$ mogą być
wybrane w $X$. Po kolapsie $\pi(z)\in L_{\pi(\gamma)}$,
a więc $N\subseteq L_\beta$.
Z drugiej strony dla każdego $\zeta<\beta$ istnieje
$\gamma\in X\cap\operatorname{Ord}$ z
$\pi(\gamma)=\zeta$. Wówczas
$L_\zeta=\pi(L_\gamma)\in N$, a tranzytywność
$N$ daje $L_\zeta\subseteq N$.
Sumując po $\zeta<\beta$, uzyskujemy
$L_\beta\subseteq N$. Ostatecznie $N=L_\beta$;
nierówność $\beta\le\theta$ wynika z faktu,
że $\pi$ porządkowo odwzorowuje $X\cap\theta$
na $\beta$.
\end{proof}
\begin{uwaga}[Precyzyjna granica redukcji]
\label{rem:l-recognition-gap-75}
Twierdzenie~\ref{thm:l-condensation-from-uniform-recognition-75}
ma kompletny dowód \emph{warunkowy}; nie używa
pakietu $\sigma_*$ ani tablicy spełniania dla
klasy $L$. Lemat~\ref{lem:l-hierarchy-delta1-derived-75}
udowadnia, że nad \emph{pełną ZF} istnieje jedna
formuła definiująca $\gamma\mapsto L_\gamma$.
Nadal jednak nie wykazaliśmy, że dla \emph{każdego}
granicznego $L_\theta$ ta sama formuła rozpoznaje
własne poziomy poprawnie: brak może dotyczyć
wewnętrznych świadków egzystencjalnej definicji.
Przykład~\ref{ex:l-delta1-not-arbitrary-level-75}
pokazuje, dlaczego samo zaklasyfikowanie formuły
do $\Delta_1^{\mathrm{ZF}}$ nie wystarcza.
To jest teraz \emph{pozostający istotny krok wymagany}
w tym alternatywnym dowodzie kondensacji; nie należy
mylić go z elementarnością przekrojów, która jest
już dowiedziona w Lemacie~\ref{lem:l-elementary-section-75}.
\end{uwaga}

\begin{wniosek}[Samodzielna kondensacja na współkońcowo wielu poziomach]
\label{cor:l-condensation-reflecting-levels-75}
Dla każdej liczby porządkowej $\eta$ istnieje
graniczna liczba porządkowa $\theta>\eta$ taka, że
\emph{dla każdej} podstruktury
$X\preccurlyeq(L_\theta,\in)$ jej kolaps Mostowskiego
jest strukturą $(L_\beta,\in)$ dla pewnej
$\beta\le\theta$. Wniosek jest dowiedziony w ZF
bez używania zewnętrznego pakietu adekwatności
z Lematu~\ref{lem:l-finite-condensation-package}.
\end{wniosek}
\begin{proof}
Niech $\Phi_{\Sigma}(\gamma,y)$ będzie
\emph{jedną ustaloną} formułą z
Lematu~\ref{lem:l-hierarchy-delta1-derived-75}.
Ponieważ klasa $L$ spełnia ZF
(Twierdzenie~\ref{thm:l-model-zf}), wewnątrz $L$
ta formuła definiuje graf funkcji
$\gamma\mapsto(L_\gamma)^L$.
Lemat~\ref{lem:l-hierarchy-internal} dowodzi, że
$(L_\gamma)^L=L_\gamma$ także w metateorii ZF.
Zatem dla każdych $\gamma,y\in L$ z $\gamma$
liczbą porządkową mamy
\[
L\models\Phi_{\Sigma}(\gamma,y)
\quad\Longleftrightarrow\quad y=L_\gamma.
\]
Stosujemy Lemat~\ref{lem:l-finite-reflection} do
\emph{skończonej} rodziny składającej się ze
wszystkich podformuł $\Phi_{\Sigma}$, otrzymując
graniczne $\theta>\eta$ z odbiciem tej formuły.
Dla $\gamma<\theta$ zbiór $L_\gamma$ należy do
$L_{\gamma+1}=\operatorname{Def}(L_\gamma)$,
gdyż jest nad sobą definiowany formułą $u=u$.
Ponieważ $\theta$ jest graniczny, mamy
$\gamma+1<\theta$, zatem $L_\gamma\in L_\theta$.
Dla każdego $y\in L_\theta$ odbicie daje
\[
(L_\theta,\in)\models\Phi_{\Sigma}(\gamma,y)
\iff L\models\Phi_{\Sigma}(\gamma,y)
\iff y=L_\gamma.
\]
Znaleźliśmy więc jedną formułę $\Lambda=\Phi_{\Sigma}$
spełniającą obie hipotezy
Twierdzenia~\ref{thm:l-condensation-from-uniform-recognition-75}.
Jego zastosowanie dowodzi tezy \emph{jednocześnie
we wszystkich} podstrukturach elementarnych
wybranego $L_\theta$.
\end{proof}
\begin{uwaga}[Dlaczego wniosek nie zastępuje lematu ogólnego]
\label{rem:l-reflecting-levels-not-all-75}
Odbicie zapewnia \emph{współkońcowo wiele}
odpowiednich poziomów $L_\theta$, nie wszystkie
poziomy graniczne. Wniosek
\ref{cor:l-condensation-reflecting-levels-75}
nie dowodzi zatem bezpośrednio ogólnego
Lematu~\ref{lem:l-condensation}, którego teza
obejmuje \emph{każde} graniczne $\theta$.
Jest natomiast odrębnym, autonomicznym wynikiem
o kondensacji, użytecznym wszędzie tam, gdzie
można swobodnie wybrać dostatecznie duży poziom
za pomocą odbicia.
\end{uwaga}

\begin{lemat}[Liczność poziomów hierarchii $L$]
\label{lem:l-level-cardinality}
W teorii $\mathrm{ZFC}+V=L$, dla każdej nieskończonej liczby
porządkowej $\xi$ zachodzi $|L_\xi|=|\xi|$.
\end{lemat}
\begin{proof}
W każdym kroku hierarchii mamy $L_\xi\cap\operatorname{Ord}=\xi$,
a więc $\xi\subseteq L_\xi$ i $|\xi|\le |L_\xi|$.
Wystarczy dowieść nierówności odwrotnej. Dla $\xi=\omega$
wszystkie poziomy $L_n$ są skończone, zatem $L_\omega$
jest zbiorem przeliczalnym.

Jeśli $\xi=\eta+1$, gdzie $\eta\ge\omega$, każdy zbiór
$z\in L_{\eta+1}=\operatorname{Def}(L_\eta)$ ma kod postaci
$(e,\vec p)$: $e\in\omega$ koduje formułę, a
$\vec p\in L_\eta^{<\omega}$ jest skończoną krotką parametrów.
Wobec tego istnieje surjekcja z pewnego podzbioru
$\omega\times L_\eta^{<\omega}$ na $L_{\eta+1}$.
W ZFC moc zbioru kodów nie przekracza
$\max(\aleph_0,|L_\eta|)=|\eta|$.
Z założenia indukcyjnego wynika
$|L_{\eta+1}|\le|\eta|=|\eta+1|$.

Jeśli $\xi>\omega$ jest graniczny, to
$L_\xi=\bigcup_{\omega\le\eta<\xi}L_\eta$.
Rozważamy zbiór par $(\eta,z)$, gdzie
$\omega\le\eta<\xi$ i $z\in L_\eta$.
Odwzorowanie $(\eta,z)\mapsto z$ jest surjekcją na $L_\xi$.
Na mocy indukcji zbiór par ma moc co najwyżej
$\sum_{\eta<\xi}\max(\aleph_0,|\eta|)\le
|\xi|\cdot|\xi|=|\xi|$, gdzie użyliśmy aksjomatu
wyboru, dostępnego w $V=L$. Otrzymujemy
$|L_\xi|\le|\xi|$, co kończy indukcję.
\end{proof}

\begin{lemat}[Kadłub elementarny o zadanej mocy]
\label{lem:l-skolem-hull-kappa}
Załóżmy $\mathrm{ZFC}+V=L$. Niech $\kappa$ będzie nieskończoną
liczbą kardynalną, $\theta>0$ liczbą porządkową, a
$A\subseteq L_\theta$ zbiorem mocy nie większej niż $\kappa$.
Wówczas istnieje $X\subseteq L_\theta$ takie, że
\[
 A\subseteq X,\qquad X\preccurlyeq (L_\theta,\in),
 \qquad |X|\le\kappa.
\]
Jeżeli $|A|=\kappa$, to $|X|=\kappa$.
\end{lemat}
\begin{proof}
Dla każdej formuły $\exists y\,\psi(y,\vec x)$
ustalamy funkcję Skolema na krotkach $\vec a\in L_\theta$:
jeżeli istnieje świadek $y\in L_\theta$, wybieramy
$<_{L}$-najmniejszego takiego świadka; w przeciwnym
wypadku wybieramy ustalony element $L_\theta$.
Ponieważ $\theta>0$, struktura ma niepusty zbiór
obiektów; funkcje mają więc wartości w $L_\theta$.
Ponieważ język teorii mnogości jest przeliczalny,
funkcji Skolema jest przeliczalnie wiele. Porządek
$<_{L}$ pozwala wyznaczyć świadków bez dodatkowych
arbitralnych wyborów.

Kładziemy $X_0=A$ (lub dodajemy jeden ustalony element,
aby pracować z niepustą strukturą) i indukcyjnie definiujemy
$X_{n+1}$ jako sumę $X_n$ i zbiorów wartości
wszystkich funkcji Skolema na skończonych krotkach
z $X_n$. Wykorzystując arytmetykę nieskończonych liczb
kardynalnych w ZFC, otrzymujemy $|X_n|\le\kappa$
indukcyjnie dla każdego $n\in\omega$. Stąd
$X=\bigcup_{n\in\omega}X_n$ ma moc co najwyżej $\kappa$.
Każdy prawdziwy w $L_\theta$ warunek egzystencjalny
z parametrami w $X$ ma świadka w pewnym $X_{n+1}$.
Z kryterium Tarskiego--Vaughta (konsekwencji
Twierdzenia~\ref{thm:downward-ls}) wynika elementarność
$X\preccurlyeq L_\theta$. Jeżeli $|A|=\kappa$,
inkluzja $A\subseteq X$ daje nierówność odwrotną.
\end{proof}



\begin{lemat}[Niezmienniczość definiowalności przy izomorfizmie]
\label{lem:l-def-isomorphism-74}
Niech $A,B$ będą zbiorami niepustymi, $E\subseteq A\times A$,
$F\subseteq B\times B$, a $f:(A,E)\to(B,F)$ izomorfizmem
struktur dla jednego symbolu relacyjnego. Dla każdej formuły
pierwszego rzędu $\varphi(v_0,\ldots,v_n)$ i każdego
$\vec a\in A^{n+1}$ zachodzi
\[
 (A,E)\models\varphi(\vec a)
 \quad\Longleftrightarrow\quad
 (B,F)\models\varphi(f(\vec a)).
\]
Jeżeli $\operatorname{Def}_E(A)$ oznacza rodzinę podzbiorów
$A$ definiowalnych w $(A,E)$ z parametrami z $A$, to
\[
 \{f[Z]:Z\in\operatorname{Def}_E(A)\}
       =\operatorname{Def}_F(B).
\]
\end{lemat}
\begin{proof}
Dowodzimy pierwszej równoważności indukcyjnie po budowie
$\varphi$. Dla $v_i=v_j$ wynika ona z iniekcyjności $f$.
Dla $v_i\,E\,v_j$ wynika bezpośrednio z definicji
izomorfizmu. Krok negacji polega na zanegowaniu obu stron
równoważności indukcyjnej, a krok koniunkcji na połączeniu
dwóch takich równoważności. W kroku egzystencjalnym
$(A,E)\models\exists v\,\psi(v,\vec a)$ wtedy i tylko
wtedy, gdy istnieje $c\in A$ spełniające $\psi(c,\vec a)$.
Z założenia indukcyjnego jest to równoważne istnieniu
$f(c)\in B$ spełniającego $\psi(f(c),f(\vec a))$;
surjektywność $f$ daje również implikację odwrotną.
Pozostałe spójniki i kwantyfikator ogólny są skrótami
poprzednich konstrukcji.

Jeżeli $Z=\{a\in A:(A,E)\models\psi(a,\vec p)\}$,
to pierwsza część dowodu daje
$f[Z]=\{b\in B:(B,F)\models\psi(b,f(\vec p))\}$,
więc $f[Z]\in\operatorname{Def}_F(B)$. Stosując ten sam
argument do izomorfizmu odwrotnego $f^{-1}$ otrzymujemy
inkluzję przeciwną. W szczególności nie zastępujemy tu
zbioru $Z$ jego przecięciem z innym uniwersum bez
uzasadnienia; przy przejściu do kolapsu użyjemy jawnie
$\pi(z)=\{\pi(u):u\in z\cap X\}$ dla $z\in X$.
\end{proof}
\begin{uwaga}[Granica absolutności w powyższym lemacie]
\label{rem:l-isomorphism-fixed-vs-uniform-74}
Indukcja dotyczy \emph{każdej ustalonej} formuły
$\varphi$ oraz dwóch już określonych struktur zbiorowych.
Nie wynika z niej bez dodatkowej pracy, że poziom
$L_\theta$ posiada wewnętrzny zbiór kodów wszystkich
prawdziwych formuł $(L_\delta,\in)$ dla \emph{każdego}
$\delta<\theta$; taki wniosek wymaga także jednolitej
konstrukcji i dostępności odpowiednich zbiorów kodów
na danym poziomie. Dlatego Lemat~\ref{lem:l-def-isomorphism-74}
usuwa przeskok w kroku izomorfizmu kadłuba
rozszerzonego, ale nie zastępuje technicznego
pakietu z Lematu~\ref{lem:l-finite-condensation-package}.
\end{uwaga}

\begin{lemat}[Kondensacja dla kadłuba w języku rozszerzonym]
\label{lem:l-expanded-hull-condensation}
Niech $\theta$ będzie niezerową graniczną liczbą porządkową.
W metateorii ZF rozważmy strukturę $\mathfrak A_\theta$
o uniwersum $L_\theta$, której język oprócz przynależności
zawiera następujące trzy relacje (ich interpretacje są
\emph{zewnętrznymi} zbiorami krotek):
\begin{itemize}
\item $R(\delta,z)$ wtedy i tylko wtedy, gdy $\delta<\theta$
  jest liczbą porządkową i $z\in L_\delta$;
\item $S(\delta,e,\vec p)$ wtedy i tylko wtedy, gdy
  $0<\delta<\theta$, $e$ jest kodem formuły
  $\varphi_e$ o liczbie zmiennych wolnych zgodnej z
  długością $\vec p\in L_\delta^{<\omega}$ oraz
  $(L_\delta,\in)\models\varphi_e(\vec p)$;
\item dla $0<\delta<\theta$ relacja
  $D(\delta,z,e,\vec p)$ zachodzi wtedy i tylko wtedy,
  gdy $e$ koduje formułę $\varphi_e(x,\vec y)$,
  $\vec p\in L_\delta^{<\omega}$ i
  $z=\{u\in L_\delta:(L_\delta,\in)\models
  \varphi_e(u,\vec p)\}$. Dla $\delta=0$
  przyjmujemy $D(0,z,e,\vec p)$ wtedy i tylko wtedy,
  gdy $z=\varnothing$, $\vec p$ jest krotką pustą
  i $e$ jest kodem formuły z wyróżnioną zmienną $x$;
  w ten sposób $D$ koduje
  $\operatorname{Def}(L_0)=\{\varnothing\}$
  bez relacji spełniania dla pustej struktury.
\end{itemize}
Skończone krotki kodujemy standardowo jako zbiory.
Ponieważ $\theta$ jest graniczny, każdy taki kod
złożony z elementów $L_\delta$ należy do $L_\theta$.
Relacje $R,S,D$ istnieją jako zbiory w metateorii:
ich argumenty pochodzą z ustalonego zbioru $L_\theta$,
a spełnianie dla struktury o zbiorowym uniwersum
$L_\delta$ jest definiowalne w ZF.

Jeżeli $X\preccurlyeq\mathfrak A_\theta$ i
$\pi:X\to\bar X$ jest kolapsem Mostowskiego
\emph{reduktu} $(X,\in)$, to
$\bar X=L_\beta$ dla pewnej granicznej liczby
porządkowej $\beta\le\theta$. Ponadto dla każdej
liczby porządkowej $\delta\in X$ zachodzi
\[
\pi[X\cap L_\delta]=L_{\pi(\delta)}.
\]
Nie zakładamy, że relacje $R,S,D$ są definiowalne
\emph{wewnątrz} $(L_\theta,\in)$.
\end{lemat}
\begin{proof}
\emph{Krok 1: elementarność przekrojów.}
Niech $\delta\in X$ będzie dodatnią liczbą porządkową.
Z kryterium Tarskiego--Vaughta wykażemy
$(X\cap L_\delta,\in)\preccurlyeq(L_\delta,\in)$.
Dla $\delta=0$ obie strony są puste; nie używamy
pojęcia elementarności struktur o pustym uniwersum.
W przypadku $\delta>0$ obie strony mają niepuste
uniwersa: $\varnothing\in X\cap L_\delta$ dzięki
elementarności $X$ i definowalności zbioru pustego.
Ustalmy formułę $\exists y\,\psi(y,\vec p)$
i parametry $\vec p\in(X\cap L_\delta)^{<\omega}$.
Jeżeli $(L_\delta,\in)\models\exists y\,\psi(y,\vec p)$,
to w rozszerzonej strukturze $\mathfrak A_\theta$
prawdziwa jest formuła
\[
\exists y\bigl(R(\delta,y)\land
 S(\delta,\ulcorner\psi\urcorner,(y,\vec p))\bigr).
\]
Kod $\ulcorner\psi\urcorner$ jest liczbą naturalną,
a każda liczba naturalna należy do $X$:
struktura $L_\theta$ zawiera wszystkie skończone
liczby porządkowe, które można kolejno jednoznacznie
definiować. Skończona krotka parametrów i jej kod
również należą do $X$, gdyż $L_\theta$ jest domknięty
na skończone kodowanie, a $X$ jest elementarny.
Elementarność $X$ \emph{w rozszerzonym języku}
daje świadka $y\in X$ spełniającego $R(\delta,y)$,
a więc $y\in X\cap L_\delta$. Jest to dokładnie
warunek Tarskiego--Vaughta. Nie użyliśmy tutaj
wewnętrznej definicji relacji $S$ w $L_\theta$.

\emph{Krok 2: indukcja po liczbach porządkowych z $X$.}
Niech $\delta\in X\cap\operatorname{Ord}$. Wykażemy
\[
(\ast)_\delta:\quad
\pi[X\cap L_\delta]=L_{\pi(\delta)}.
\]
Dla $\delta=0$ obie strony są puste.
Jeżeli $\delta=1$, to $0=\varnothing\in X$
i $L_1=\{0\}$, więc
$\pi[X\cap L_1]=\{0\}=L_1=L_{\pi(1)}$.
Na etapach następnikowych $\delta=\eta+1$ z
$\eta\geq1$ stosujemy niżej etap indukcyjny
oraz elementarność przekrojów niepustych.
Wówczas $\eta\in X$, ponieważ operacja
poprzednika liczby porządkowej następnikowej jest
jednoznacznie definiowalna w $(L_\theta,\in)$.
Załóżmy $(\ast)_\eta$. Weźmy $z\in X\cap L_{\eta+1}$.
Z definicji $L_{\eta+1}$ istnieją kod formuły $e$
i skończona krotka parametrów $\vec p\in L_\eta^{<\omega}$
takie, że $D(\eta,z,e,\vec p)$. Dzięki elementarności
\emph{rozszerzonej} struktury wybieramy
$e,\vec p\in X$. Dla każdego $u\in X\cap L_\eta$
Krok 1 daje równoważność
\[
(L_\eta,\in)\models\varphi_e(u,\vec p)
\iff
(X\cap L_\eta,\in)\models\varphi_e(u,\vec p).
\]
Na mocy Lematu~\ref{lem:l-def-isomorphism-74}
relacja definiowalności przechodzi przez izomorfizm.
Kolaps jest izomorfizmem $X\cap L_\eta$ na
$L_{\pi(\eta)}$ na mocy $(\ast)_\eta$, zatem
\[
\pi(z)=\{v\in L_{\pi(\eta)}:
(L_{\pi(\eta)},\in)\models
\varphi_e(v,\pi(\vec p))\}\in
L_{\pi(\eta)+1}=L_{\pi(\delta)}.
\]
I odwrotnie: każdy $b\in L_{\pi(\eta)+1}$ ma taki
kod $e$ i parametry $\vec q\in L_{\pi(\eta)}$.
Z $(\ast)_\eta$ otrzymujemy parametry
$\vec p\in X\cap L_\eta$ o obrazie $\vec q$.
Relacja $D$ określa jednoznacznie $z\in L_{\eta+1}$;
ponieważ $\eta,e,\vec p\in X$, elementarność daje
$z\in X$. Poprzedni rachunek pokazuje $\pi(z)=b$.
To dowodzi $(\ast)_\delta$ w kroku następnikowym.

Jeżeli $\delta$ jest graniczny i $z\in X\cap L_\delta$,
to $z\in L_\eta$ dla pewnej liczby porządkowej
$\eta<\delta$. Relacja $R$ należy do języka
$\mathfrak A_\theta$, dlatego elementarność daje
takie $\eta\in X$. Z założeń indukcyjnych wynika
\[
\pi[X\cap L_\delta]
=\bigcup_{\eta\in X\cap\delta}L_{\pi(\eta)}
=L_{\pi(\delta)}.
\]
Ostatnia równość korzysta z faktu, że obraz
$X\cap\delta$ przez $\pi$ jest dokładnie
liczbą porządkową $\pi(\delta)$.

\emph{Krok 3: identyfikacja całego kolapsu.}
Dla dowolnego $z\in X$ istnieje $\delta<\theta$
z $z\in L_\delta$, ponieważ $\theta$ jest graniczny.
Elementarność dla relacji $R$ daje takie $\delta\in X$.
Krok 2 pokazuje, że $\pi(z)\in L_{\pi(\delta)}$.
Niech $\beta=\bar X\cap\operatorname{Ord}$.
Z definicji kolapsu dla liczb porządkowych
$\beta$ jest liczbą porządkową, a jego zbiór
poprzedników ma postać
$\{\pi(\delta):\delta\in X\cap\operatorname{Ord}\}$.
Ponieważ $X$ jest domknięty na następniki liczb
porządkowych, $\beta$ jest graniczny.
Stąd $\bar X\subseteq L_\beta$.
Dla każdego $\varepsilon<\beta$ istnieje
$\delta\in X\cap\operatorname{Ord}$
z $\pi(\delta)=\varepsilon$. Z Kroku 2 wynika
$L_\varepsilon\subseteq\bar X$.
Biorąc sumę po $\varepsilon<\beta$,
otrzymujemy $L_\beta\subseteq\bar X$.
Zatem $\bar X=L_\beta$. Ponieważ liczby porządkowe
w $X$ są mniejsze od $\theta$, otrzymujemy
$\beta\le\theta$.
\end{proof}

\begin{uwaga}[Porównanie dwóch samodzielnych dowodów]
\label{rem:l-expanded-vs-full-condensation}
Lemat~\ref{lem:l-expanded-hull-condensation}
konstruuje szczególny kadłub elementarny
w języku rozszerzonym o $R,S,D$ i daje osobną
ścieżkę dowodu GCH. Ogólny
Lemat~\ref{lem:l-condensation} został natomiast
udowodniony przez skończony pakiet rozpoznający
wszystkie graniczne poziomy $L_\theta$.
Obie drogi nie powinny być ze sobą utożsamiane:
pierwsza wymaga wyboru rozszerzonego języka,
a druga dowodu absolutności kodów historii.
\end{uwaga}

\begin{twierdzenie}[GCH w uniwersum konstruowalnym]
\label{thm:l-gch}
W teorii ZF dowodzimy, że $L\models\mathrm{ZFC}+\mathrm{GCH}$.
Dokładniej, dla każdej nieskończonej liczby kardynalnej
$\kappa$ \emph{obliczonej wewnątrz klasy $L$} zachodzi
\[
 L\models 2^\kappa=\kappa^+,
\]
gdzie $\kappa^+$ oznacza następną liczbę kardynalną
\emph{obliczoną wewnątrz $L$}.
\end{twierdzenie}
\begin{proof}
Rozdzielimy dwie płaszczyzny argumentu. Najpierw
przeprowadzimy rachunek w teorii $\mathrm{ZFC}+V=L$.
Następnie zastosujemy go wewnątrz klasy $L$.
Możliwość wykonania tego drugiego kroku wynika
z Twierdzenia~\ref{thm:l-model-zf},
Twierdzenia~\ref{thm:l-definable-wellorder} oraz
Lematu~\ref{lem:l-hierarchy-internal}.

\emph{Etap 1: ustalenie liczby kardynalnej i podzbioru.}
Pracujemy chwilowo przy założeniu $V=L$. Niech $\kappa$
będzie dowolną nieskończoną liczbą kardynalną,
a $A$ dowolnym podzbiorem $\kappa$. Ponieważ $V=L$,
zbiór $A$ należy do pewnego poziomu $L_\eta$.
Wybieramy graniczną liczbę porządkową $\theta>\eta$
tak dużą, że $A,\kappa\in L_\theta$.
Możemy to uczynić, ponieważ hierarchia $L$ jest
ciągła na etapach granicznych i jej suma obejmuje $V$.

\emph{Etap 2: kadłub obejmujący wszystkie liczby
porządkowe poniżej $\kappa$.}
Bierzemy zbiór parametrów
$D=\kappa\cup\{\kappa,A\}$, gdzie $\kappa$ po prawej
stronie rozumiemy jako \emph{zbiór wszystkich liczb
porządkowych mniejszych niż $\kappa$}.
Ponieważ $\kappa$ jest nieskończoną liczbą kardynalną,
$|D|=\kappa$. Stosujemy konstrukcję kadłuba z
Lematu~\ref{lem:l-skolem-hull-kappa} do
\emph{rozszerzonej} struktury $\mathfrak A_\theta$
z Lematu~\ref{lem:l-expanded-hull-condensation}.
Jej język jest skończony, a więc przeliczalny;
zawarte w nim relacje są zbiorami w metateorii.
Otrzymujemy $X\preccurlyeq\mathfrak A_\theta$ takie, że
\[
 \kappa\subseteq X,\qquad \kappa,A\in X,
 \qquad |X|=\kappa.
\]
Uwaga: warunek $\kappa\in X$ nie zastępuje
silniejszego warunku $\kappa\subseteq X$.
Gdybyśmy wiedzieli tylko, że zbiór $\kappa$ należy
do $X$, kolaps mógłby zmienić poszczególne jego
elementy i argument dotyczący zbioru $A$ zawiódłby.

\emph{Etap 3: kolaps i kontrola jego wysokości.}
Niech $\pi:X\to\bar X$ będzie kolapsem Mostowskiego.
Lemat~\ref{lem:l-expanded-hull-condensation} daje
$\bar X=L_\beta$ dla pewnej granicznej liczby porządkowej
$\beta\le\theta$. Nie korzystamy tu z niezweryfikowanego
kroku o wewnętrznym kodowaniu wszystkich poziomów $L$
w dowolnej podstrukturze $(L_\theta,\in)$. Ponieważ $\pi$ jest bijekcją,
$|L_\beta|=|X|=\kappa$. Zbiór $L_\beta$ zawiera
każdą liczbę porządkową mniejszą niż $\beta$;
stąd $|\beta|\le\kappa$. Z definicji następnej
liczby kardynalnej wynika $\beta<\kappa^+$.

Dla każdej liczby porządkowej $\xi\le\kappa$
dowodzimy indukcyjnie $\pi(\xi)=\xi$.
Istotnie, gdy $\xi\in X$ i wszystkie liczby porządkowe
$\zeta<\xi$ należą do $X$, definicja kolapsu daje
\[
 \pi(\xi)=\{\pi(\zeta):\zeta\in X\cap\xi\}
 =\{\zeta:\zeta<\xi\}=\xi.
\]
Warunki te zachodzą, bo $\kappa\subseteq X$
i $\kappa\in X$. W szczególności
$\pi(\xi)=\xi$ dla każdego $\xi<\kappa$.

\emph{Etap 4: kolaps nie zmienia zbioru $A$.}
Ponieważ $A\subseteq\kappa$, wszystkie elementy $A$
są liczbami porządkowymi mniejszymi niż $\kappa$,
a więc należą do $X$. Mamy
\[
 \pi(A)=\{\pi(\xi):\xi\in A\cap X\}
       =\{\xi:\xi\in A\}=A.
\]
Jednocześnie $A\in X$, zatem $\pi(A)\in L_\beta$.
Otrzymujemy $A\in L_\beta\subseteq L_{\kappa^+}$.
Z dowolności podzbioru $A$ wynika
\[
 \mathcal P(\kappa)\subseteq L_{\kappa^+}
 \qquad\text{przy założeniu }V=L.
\]

\emph{Etap 5: obliczenie mocy.}
Lemat~\ref{lem:l-level-cardinality} daje
$|L_{\kappa^+}|=\kappa^+$. Stąd
$2^\kappa=|\mathcal P(\kappa)|\le\kappa^+$.
Twierdzenie Cantora (Twierdzenie~\ref{thm:zf-cantor})
daje nierówność $2^\kappa>\kappa$. Ponieważ
porządek liczb kardynalnych jest liniowy w ZFC,
zachodzi $2^\kappa\ge\kappa^+$. Wobec tego
$2^\kappa=\kappa^+$.

\emph{Etap 6: przeniesienie rachunku do klasy $L$.}
Dowód Etapów 1--5 był dowodem formalnym z aksjomatów
$\mathrm{ZFC}+V=L$, a nie argumentem wymagającym
istnienia przeliczalnego tranzytywnego modelu.
W metateorii ZF klasa $L$ spełnia ZFC, a jej
własna hierarchia konstruowalna jest tożsama
z hierarchią zewnętrzną na mocy
Lematu~\ref{lem:l-hierarchy-internal}.
Ponadto $L\models V=L$. Relatywizując do $L$
każdy krok dowodu z Etapów 1--5, otrzymujemy
GCH dla wszystkich nieskończonych liczb
kardynalnych obliczonych wewnątrz $L$.
Nie stwierdzamy, że liczby kardynalne następniki
w $L$ i w otaczającym uniwersum ZF muszą być równe.
\end{proof}

\begin{przyklad}[Znaczenie warunku $\kappa\subseteq X$]
\label{ex:l-hull-kappa-subset}
Jeśli podstruktura elementarna $X$ zawiera liczbę
porządkową $\kappa$ tylko jako \emph{element}, nie wynika
z tego, że zawiera wszystkie liczby porządkowe
mniejsze niż $\kappa$. Na przykład przeliczalna
podstruktura elementarna dużego poziomu $L_\theta$
może zawierać $\omega_1$ jako element, lecz nie może
zawierać wszystkich elementów $\omega_1$.
Jej kolaps wysyła wówczas $\omega_1$ na przeliczalną
liczbę porządkową. Dlatego w dowodzie GCH
wprowadziliśmy do kadłuba \emph{każdy} element $\kappa$,
a nie tylko sam zbiór $\kappa$.
\end{przyklad}

\begin{wniosek}[Względna niesprzeczność AC, GCH i CH]
\label{cor:l-relative-consistency}
Jeżeli ZF jest niesprzeczna, to niesprzeczna jest
również ZFC+GCH, a więc także ZFC+CH.
\end{wniosek}
\begin{proof}
Wszystkie aksjomaty ZFC+GCH po relatywizacji
kwantyfikatorów do definiowalnej klasy $L$
są twierdzeniami ZF na mocy powyższych wyników.
Gdyby z ZFC+GCH dało się wyprowadzić sprzeczność,
pewien skończony dowód używałby tylko skończenie
wielu instancji jej schematów aksjomatów.
Relatywizując każdy wiersz tego dowodu do $L$
i zastępując aksjomaty udowodnionymi w ZF
odpowiednikami, otrzymalibyśmy dowód sprzeczności
w ZF. Zatem niesprzeczność ZF implikuje
niesprzeczność ZFC+GCH.
\end{proof}


\begin{uwaga}[Status dowodu niezależności CH po audycie]
\label{rem:ch-gch-independent-of-full-condensation}
Dowód Twierdzenia~\ref{thm:l-gch} korzysta bezpośrednio
z Lematu~\ref{lem:l-expanded-hull-condensation}.
Wariant dla dowolnego $X\preccurlyeq L_\theta$
zredukowano osobno do skończonego pakietu
z Lematu~\ref{lem:l-finite-condensation-package}.
Pakiet został już udowodniony w całości w
Lemacie~\ref{lem:l-finite-condensation-package},
a dowód GCH ma dodatkowo niezależną wersję dla
kadłuba w języku rozszerzonym.
\end{uwaga}

\begin{uwaga}[Audyt: status lematu kondensacji]
\label{rem:l-condensation-proof-scope}
Bezpośredni dowód Lematu~\ref{lem:l-condensation} można oprzeć
na zgodności operacji $\operatorname{Def}$ z kolapsem, lecz
nie wynika ona z samej elementarności podstruktury. Jeżeli $\pi:X\to\bar X$ jest kolapsem, to
$\pi$ zachowuje relację przynależności między elementami $X$;
nie wolno jednak bez dodatkowego uzasadnienia utożsamiać
$\pi[X\cap L_\gamma]$ z $L_{\pi(\gamma)}$.
Przekrój $X\cap L_\gamma$ nie musi być całym $L_\gamma$,
ale przy $L_\gamma\in X$ jest jego podstrukturą elementarną
na mocy Lematu~
\ref{lem:l-elementary-section-75}.
Prawidłowe uzasadnienie kroku następnikowego wykorzystuje
jednolity kod formuły i krotki parametrów, elementarność
$X$ oraz absolutność relacji spełniania dla zbiorowej struktury.
To rozróżnienie jest istotne przy stosowaniu lematu do GCH.
\end{uwaga}



\begin{uwaga}[Od lokalnych tablic do jednolitego kodu]
\label{rem:l-condensation-scope-76}
Lemat~\ref{lem:l-local-satisfaction-bound-76}
kontroluje tablice ustalonych formuł, a
Lemat~\ref{lem:l-hf-fixed-collapse-76}
przypadek $L_\omega$.
Dopiero Lemat~\ref{lem:l-all-stage-histories-78}
zapewnia wszystkie historie jednocześnie w sensie
\emph{jednej formuły} pierwszego rzędu.
Konkretnie można przyjąć
\[
\Lambda(\gamma,y)\;\equiv\;
 \exists h\bigl(\mathsf{Hist}(h,\gamma+1)
                     \land h(\gamma)=y\bigr).
\]
Na mocy absolutności z
Lematu~\ref{lem:l-finite-godel-term-absoluteness-78}
dla każdego granicznego $L_\theta$ i
$\gamma<\theta$ struktura $(L_\theta,\in)$
rozpoznaje przez $\Lambda(\gamma,y)$ dokładnie
$y=L_\gamma$. Świadek $H_{\gamma+1}$ należy
do $L_\theta$ dzięki
Lematowi~\ref{lem:l-all-stage-histories-78}.
Spełnione są więc także hipotezy
Twierdzenia~\ref{thm:l-condensation-from-uniform-recognition-75}.
\end{uwaga}

\section{Jednolite kody poziomów konstruowalnych — zakres formalizacji}
\label{sec:l-uniform-codes-32}
\begin{definicja}[Kod zbiorowego etapu konstrukcji]
\label{def:l-stage-code-32}
Niech $\alpha$ będzie liczbą porządkową. Przyjmujemy
$\operatorname{Def}(\varnothing)=\{\varnothing\}$
zgodnie z Definicją~\ref{def:constructible-def}. \emph{Kodem etapu} nazywamy funkcję $h$ o dziedzinie $\alpha+1$, dla której $h(0)=\varnothing$, $h(\xi+1)=\operatorname{Def}(h(\xi))$ dla $\xi<\alpha$ i $h(\lambda)=\bigcup_{\xi<\lambda}h(\xi)$ dla każdej niezerowej granicznej liczby porządkowej $\lambda\leq\alpha$. Dla niepustego $A$ relacja $\operatorname{Def}(A)=B$
oznacza, że $B$ jest zbiorem wszystkich podzbiorów $A$
definiowalnych w strukturze $(A,\in)$ przez formuły
pierwszego rzędu z parametrami z $A$; przypadek pusty
reguluje wskazana konwencja.
\end{definicja}
\begin{lemat}[Jednoznaczność kodu etapu]
\label{lem:l-stage-code-unique-32}
W ZF dla każdej liczby porządkowej $\alpha$ istnieje dokładnie jeden kod etapu $h$ o dziedzinie $\alpha+1$; jego ostatnia wartość wynosi $L_\alpha$.
\end{lemat}
\begin{proof}
Dla zbioru $A$ relacja spełniania formuł w zbiorowej strukturze $(A,\in)$ jest definiowalna za pomocą rekurencji po skończonej budowie formuł. Zatem $\operatorname{Def}(A)$ istnieje na mocy aksjomatów zbioru potęgowego i wyodrębniania: wybieramy spośród podzbiorów $A$ te, dla których istnieje kod formuły i skończona krotka parametrów z $A$ definiujące dokładnie dany podzbiór. Jest to operacja jednoznaczna. Twierdzenie o rekurencji pozaskończonej, stosowane do tej operacji i do sum na etapach granicznych, daje funkcję $h$ o wskazanej dziedzinie. Jeśli $h$ i $k$ są dwoma takimi funkcjami, to indukcja pozaskończona po $\xi\leq\alpha$ dowodzi $h(\xi)=k(\xi)$: dla zera wynika to z definicji, dla następnika z jednoznaczności $\operatorname{Def}$, a dla etapu granicznego z równości wcześniejszych wartości. Wartość końcowa jest więc z definicji $L_\alpha$.
\end{proof}
\begin{uwaga}[Formalizacja i jej niezależna kontrola]
\label{rem:l-stage-internality-open-32}
Istnienie kodów z poprzedniego lematu wynika
z rekurencji w metateorii ZF. Dodatkowa treść,
niezbędna do kondensacji, została udowodniona
w Lematach~\ref{lem:l-finite-godel-term-absoluteness-78}
i~\ref{lem:l-all-stage-histories-78}:
każdy graniczny $L_\theta$ zawiera właściwe
kody etapów oraz prawidłowo interpretuje
jedną formułę $\mathsf{Hist}$. Pozwala to używać
skończonego pakietu z
Lematu~\ref{lem:l-finite-condensation-package}
bez założenia ZF wewnątrz $L_\theta$.
\end{uwaga}
\chapter{Podstawy wymuszania i modele generyczne}\label{ch:forcing-foundations}
\chaptermark{Podstawy wymuszania}
\section{Motywacja i założenia metateoretyczne}
W Rozdziale~\ref{ch:constructible} zbudowaliśmy uniwersum
konstruowalne $L$, ograniczając zbiory do tych, które dają się
zdefiniować na kolejnych etapach hierarchii. Wymuszanie służy
odwrotnemu zadaniu: chcemy rozszerzyć model teorii mnogości
o nowy obiekt, zachowując możliwie wiele aksjomatów.
Najprostszy przykład polega na dodaniu ciągu binarnego
przez stopniowe ustalanie jego skończonych fragmentów.

\begin{uwaga}[Model wyjściowy]\label{rem:forcing-background}
W rozdziale pracujemy w zewnętrznym uniwersum ZFC i zakładamy,
że istnieje przeliczalny tranzytywny model zbiorowy
$M\models\mathrm{ZFC}$. Przeliczalność jest własnością
zewnętrzną, nie orzeczeniem modelu $M$ o sobie samym.
Tranzytywność oznacza
$\forall x\in M\ \forall y\in x\ (y\in M)$.
Istnienie takiego modelu nie wynika z samej syntaktycznej
niesprzeczności ZFC. W szczególności poniższych konstrukcji
nie wolno bez dalszego argumentu traktować jako dowodu
względnej niesprzeczności ZFC i niezależności CH.
\end{uwaga}

\section{Warunki i zgodność informacji}
\begin{definicja}[Pojęcie wymuszania]\label{def:forcing-poset}
\emph{Pojęciem wymuszania} nazywamy niepusty zbiór
$\mathbb P$ z relacją $\leq$, która jest zwrotna
i przechodnia. Napis $q\leq p$ oznacza, że $q$ jest
\emph{silniejszym} warunkiem od $p$.
Nie wymagamy antysymetrii. Dla uproszczenia dalszej
konstrukcji zakładamy istnienie największego warunku
$1_{\mathbb P}$, dla którego
$p\leq1_{\mathbb P}$ dla każdego $p\in\mathbb P$.
\end{definicja}

\begin{definicja}[Zgodność i antyłańcuch]\label{def:forcing-compatible}
Warunki $p,q\in\mathbb P$ są zgodne, co zapisujemy
$p\parallel q$, jeżeli istnieje $r\in\mathbb P$
takie, że $r\leq p$ i $r\leq q$.
W przeciwnym przypadku piszemy $p\perp q$.
Zbiór $A\subseteq\mathbb P$ jest antyłańcuchem,
gdy jego dwa różne elementy są niezgodne.
\end{definicja}

\begin{przyklad}[Skończone przybliżenia ciągu binarnego]
\label{ex:cohen-poset}
Niech $\operatorname{Fn}(\omega,2)$ oznacza zbiór
wszystkich funkcji częściowych
$p:F\to\{0,1\}$ o skończonej dziedzinie
$F\subseteq\omega$. Definiujemy
\[
 \mathbb P_{\mathrm C}:=\operatorname{Fn}(\omega,2),
 \qquad q\leq p\iff q\supseteq p.
\]
Największym elementem jest funkcja pusta.
Warunki $p,q$ są zgodne wtedy i tylko wtedy,
gdy zgadzają się na części wspólnej dziedzin:
w takim przypadku ich suma jako funkcji
$p\cup q$ jest wspólnym wzmocnieniem.
Na przykład
$\{(0,0),(3,1)\}$ i $\{(1,1),(3,1)\}$
są zgodne, natomiast
$\{(3,0)\}\perp\{(3,1)\}$.
\end{przyklad}

\begin{figure}[htbp]
\centering
\begin{tikzpicture}[
 level distance=14mm,
 sibling distance=39mm,
 every node/.style={draw,rounded corners,inner sep=3pt,font=\small},
 edge from parent/.style={draw,-{Stealth[length=1.8mm]}}
]
\node {$\varnothing$}
 child {node {$\{(0,0)\}$}
   child {node {$\{(0,0),(1,0)\}$}}
   child {node {$\{(0,0),(1,1)\}$}}}
 child {node {$\{(0,1)\}$}
   child {node {$\{(0,1),(1,0)\}$}}
   child {node {$\{(0,1),(1,1)\}$}}};
\end{tikzpicture}
\caption{Sześć przykładowych wzmocnień warunków Cohena.
Kierunek w dół oznacza silniejsze warunki.}
\label{fig:forcing-cohen-tree}
\end{figure}

\begin{lemat}[Monotoniczność niezgodności]
\label{lem:forcing-compatibility}
Dla dowolnych $p,q,r\in\mathbb P$:
jeśli $p\perp q$ i $r\leq p$, to $r\perp q$.
Jeżeli $p\parallel q$ i $p\leq p'$, $q\leq q'$,
to $p'\parallel q'$.
\end{lemat}
\begin{proof}
Gdyby $r\parallel q$, wspólne wzmocnienie
$s\leq r,q$ spełniałoby też $s\leq p,q$,
wbrew niezgodności. Jeśli natomiast $s\leq p,q$
oraz $p\leq p'$ i $q\leq q'$, to z przechodniości
$s\leq p',q'$, co dowodzi drugiego stwierdzenia.
\end{proof}

\section{Filtry i zbiory gęste}
\begin{definicja}[Filtr]\label{def:forcing-filter}
Filtr w pojęciu wymuszania $(\mathbb P,\leq)$ to niepusty
zbiór $G\subseteq\mathbb P$, który jest domknięty
ku słabszym warunkom i skierowany ku silniejszym:
\begin{enumerate}
\item jeśli $p\in G$ i $p\leq q$, to $q\in G$;
\item jeśli $p,q\in G$, istnieje $r\in G$ z $r\leq p,q$.
\end{enumerate}
\end{definicja}

\begin{definicja}[Gęstość]\label{def:forcing-dense}
Zbiór $D\subseteq\mathbb P$ jest \emph{gęsty},
jeśli $\forall p\in\mathbb P\;\exists q\in D\ (q\leq p)$.
Nazywamy go \emph{otwartym w porządku},
jeżeli $p\in D$ i $q\leq p$ implikują $q\in D$.
\end{definicja}

\begin{lemat}[Otwarte domknięcie zbioru gęstego]
\label{lem:forcing-open-dense}
Jeżeli $D$ jest gęsty, to
$D^\downarrow=\{p:\exists q\in D\ (p\leq q)\}$
jest gęsty i otwarty w porządku. Dla dowolnego
filtru $G$ zachodzi
$G\cap D\ne\varnothing$ wtedy i tylko wtedy,
gdy $G\cap D^\downarrow\ne\varnothing$.
\end{lemat}
\begin{proof}
Jeśli $p\in\mathbb P$, gęstość daje
$q\in D$ z $q\leq p$. Zatem $q\in D^\downarrow$,
więc jest gęsty. Gdy $r\leq p\in D^\downarrow$
oraz $p\leq q\in D$, to $r\leq q$, więc
$r\in D^\downarrow$. Jedna implikacja dotycząca
filtrów wynika z $D\subseteq D^\downarrow$.
Dla drugiej, z $p\in G\cap D^\downarrow$
wybieramy $q\in D$ z $p\leq q$; wtedy
$q\in G\cap D$, ponieważ filtr jest
domknięty ku słabszym warunkom.
\end{proof}

\begin{definicja}[Filtr generyczny]\label{def:forcing-generic}
Niech $M$ będzie tranzytywnym modelem ZFC,
$(\mathbb P,\leq)\in M$ pojęciem wymuszania.
Filtr $G\subseteq\mathbb P$ nazywamy
\emph{$M$-generycznym}, jeśli dla każdego
gęstego $D\subseteq\mathbb P$ należącego do
$M$ zachodzi $G\cap D\ne\varnothing$.
Nie wymagamy $G\in M$. Gęstość $D\in M$ jest
absolutna, gdyż wszystkie kwantyfikatory
biegną po $\mathbb P\subseteq M$.
\end{definicja}

\begin{lemat}[Generyczność rozstrzyga zgodność]
\label{lem:forcing-generic-decision}
Jeśli $G$ jest filtrem $M$-generycznym, to
dla każdego $p\in\mathbb P$ albo $p\in G$,
albo istnieje $q\in G$ niezgodny z $p$.
Dokładnie jedna z tych możliwości zachodzi.
\end{lemat}
\begin{proof}
Z tranzytywności $p\in M$. Zbiór
$D_p=\{q\in\mathbb P:q\leq p\text{ lub }q\perp p\}$
należy do $M$ i jest gęsty. Dla każdego $r$
albo $r\perp p$, albo istnieje wspólne
wzmocnienie $s\leq r,p$. Generyczność
daje $q\in G\cap D_p$. Jeśli $q\leq p$,
to $p\in G$ przez domknięcie ku słabszym
warunkom. W drugim przypadku $q\perp p$.
Oba przypadki nie mogą zachodzić razem,
gdyż dwa elementy filtru są zgodne.
\end{proof}

\section{Lemat Rasiowej--Sikorskiego}
\begin{twierdzenie}[Lemat Rasiowej--Sikorskiego]
\label{thm:rasiowa-sikorski}
Niech $\mathbb P$ będzie niepustym
przeliczalnym pojęciem wymuszania,
$(D_n)_{n\in\omega}$ ciągiem gęstych
podzbiorów $\mathbb P$, a $p_*\in\mathbb P$.
W ZFC istnieje filtr $G\subseteq\mathbb P$
zawierający $p_*$ i przecinający
każdy zbiór $D_n$.
\end{twierdzenie}
\begin{proof}
Ustalmy enumerację $\mathbb P$. Definiujemy
$p_0=p_*$, a $p_{n+1}$ jako pierwszy
w tej enumeracji element $D_n$ spełniający
$p_{n+1}\leq p_n$. Gęstość zapewnia istnienie.
Niech
\[
 G=\{q\in\mathbb P:\exists n\in\omega\ (p_n\leq q)\}.
\]
Zbiór zawiera $p_0$, więc jest niepusty.
Jeżeli $q\in G$ i $q\leq r$,
to $p_n\leq q\leq r$ dla pewnego $n$,
czyli $r\in G$. Jeśli $q,r\in G$
i $p_n\leq q$, $p_m\leq r$, to
$p_{\max(m,n)}\in G$ jest wspólnym
wzmocnieniem obu warunków. Zatem $G$
jest filtrem. Ponieważ $p_{n+1}\in G\cap D_n$,
przecina wszystkie wskazane zbiory.
\end{proof}

\begin{wniosek}[Filtr generyczny nad modelem przeliczalnym]
\label{cor:countable-generic}
Jeżeli $M$ jest przeliczalnym tranzytywnym
modelem ZFC, $(\mathbb P,\leq)\in M$
i $p_*\in\mathbb P$, to istnieje zewnętrzny
filtr $M$-generyczny zawierający $p_*$.
\end{wniosek}
\begin{proof}
Tranzytywność zapewnia $\mathbb P\subseteq M$,
zatem $\mathbb P$ jest przeliczalny zewnętrznie.
Rodzina gęstych podzbiorów $\mathbb P$
należących do $M$ jest także przeliczalna,
gdyż jest podzbiorem $M$.
Wyliczmy ją jako $(D_n)$, powtarzając
elementy w przypadku rodziny skończonej.
Twierdzenie~\ref{thm:rasiowa-sikorski}
daje odpowiedni filtr.
\end{proof}

\begin{uwaga}[Zewnętrzna perspektywa]
\label{rem:forcing-external-generic}
Wyliczenie gęstych zbiorów należących do $M$
jest dokonane na zewnątrz i nie musi należeć
do samego $M$. Nie dowodzi więc istnienia
generycznego filtru jako elementu $M$.
Nie wymaga też przecięcia wszystkich
gęstych zbiorów zewnętrznego uniwersum.
\end{uwaga}

\chapter{Relacja wymuszania i twierdzenie o wymuszaniu}
\label{ch:forcing-theorem}
\chaptermark{Twierdzenie o wymuszaniu}

\section{Język wymuszania i gęstość lokalna}
W Rozdziale~\ref{ch:forcing-foundations} zdefiniowaliśmy
nazwy, ich wartości i tranzytywny zbiór $M[G]$.
Nie wykazaliśmy jeszcze, że $M[G]$ spełnia ZFC.
Zrobimy to poprzez relację wymuszania, definiowaną
w modelu wyjściowym, a następnie porównaną
z prawdziwością zdań w rozszerzeniu.

\begin{uwaga}[Założenia]
\label{rem:forcing33-background}
Przyjmujemy, że $M$ jest przeliczalnym
tranzytywnym modelem zbiorowym ZFC,
a $(\mathbb P,\leq,1_{\mathbb P})\in M$
jest pojęciem wymuszania z największym warunkiem.
Nie zakładamy antysymetrii ani separatywności.
Zawsze gdy piszemy o nazwach $\sigma,\tau,\rho$,
są to nazwy $\mathbb P$ należące do $M$;
$G$ jest filtrem $M$-generycznym.
Przeliczalność $M$ jest zewnętrzna.
Nie wywodzimy istnienia takiego modelu
z samej $\operatorname{Con}(\mathrm{ZFC})$.
\end{uwaga}

\begin{definicja}[Gęstość poniżej warunku]
\label{def:forcing33-local-dense}
Niech $p\in\mathbb P$, $D\subseteq\mathbb P$.
Zbiór $D$ jest \emph{gęsty poniżej $p$},
gdy dla każdego $q\leq p$ istnieje
$r\in D$ z $r\leq q$.
\end{definicja}

\begin{lemat}[Lokalna generyczność]
\label{lem:forcing33-local-generic}
Jeśli $p\in G$ i $D\in M$ jest gęsty
poniżej $p$, to $D\cap G\neq\varnothing$.
\end{lemat}
\begin{proof}
Niech $I_p=\{q\in\mathbb P:q\perp p\}$.
Zbiory $I_p$ i $D\cup I_p$ należą do $M$.
Jeżeli $q\perp p$, to $q\in I_p$.
W przeciwnym przypadku istnieje
$s\leq q,p$, a gęstość $D$ poniżej $p$
dostarcza $r\in D$ z $r\leq s\leq q$.
Zatem $D\cup I_p$ jest gęsty w $\mathbb P$.
Generyczność daje $t\in G\cap(D\cup I_p)$.
Ponieważ $p\in G$, każdy warunek $t\in G$
jest zgodny z $p$, więc $t\notin I_p$.
Stąd $t\in G\cap D$.
\end{proof}

\begin{przyklad}[Gęstość lokalna nie jest globalna]
\label{ex:forcing33-local}
W $\operatorname{Fn}(\omega,2)$ ustalmy
$p=\{(0,0)\}$ oraz
$D=\{q:q\leq p,\ 1\in\operatorname{dom}q\}$.
Każde wzmocnienie $p$ można przedłużyć
do elementu $D$. Jednak warunek
$\{(0,1)\}$ jest niezgodny z każdym
elementem $D$, więc $D$ nie jest
gęsty w całym wymuszaniu.
\end{przyklad}

\section{Atomowa relacja wymuszania}
W języku z parametrami będącymi nazwami
formuły atomowe mają postać $\sigma\in\tau$
albo $\sigma=\tau$. Nie definiujemy
wymuszania przez kwantyfikację po
rozszerzeniach generycznych: potrzebna
jest relacja \emph{wewnętrznie definiowalna}
w modelu $M$.

\begin{definicja}[Atomowe klauzule wymuszania]
\label{def:forcing33-atomic}
Niech $\sigma,\tau\in M$ będą nazwami,
a $p\in\mathbb P$. Przez jednoczesną
rekurencję ustalamy:
\[
\begin{aligned}
p\Vdash\sigma\in\tau
&\iff\forall q\leq p\ \exists r\leq q\
 \exists(\pi,s)\in\tau\
 (r\leq s\land r\Vdash\sigma=\pi),\\
p\Vdash\sigma=\tau
&\iff
 \bigl[\forall(\pi,s)\in\sigma\
  \forall q\leq p,s\
  \exists r\leq q\ (r\Vdash\pi\in\tau)\bigr]\\
&\quad{}\land
 \bigl[\forall(\pi,s)\in\tau\
  \forall q\leq p,s\
  \exists r\leq q\ (r\Vdash\pi\in\sigma)\bigr].
\end{aligned}
\]
W skrócie $\forall q\leq p,s$ wymagamy
jednocześnie $q\leq p$ i $q\leq s$.
\end{definicja}

\begin{uwaga}[Dobrze ufundowana rekurencja]
\label{rem:forcing33-ranks}
Dla nazwy $\tau$ niech $\rho(\tau)$ będzie
najmniejszym $\alpha$ z
$\tau\in N_{\alpha+1}$; hierarchia $N_\alpha$
pochodzi z Definicji~\ref{def:forcing-names}.
Każda pierwsza współrzędna $\pi$
pary $(\pi,s)\in\tau$ ma rangę
$\rho(\pi)<\rho(\tau)$.
Do uszeregowania dwóch argumentów
stosujemy sumę naturalną liczb
porządkowych $\alpha\mathbin{\#}\beta$,
zdefiniowaną przez dodawanie współczynników
przy jednakowych potęgach $\omega$
w postaci normalnej Cantora.
W odróżnieniu od zwykłego dodawania
jest ona ściśle rosnąca względem
każdego argumentu. W obu klauzulach
definicji zastąpienie jednej z nazw
jej pierwszą współrzędną zmniejsza
$\rho(\sigma)\mathbin{\#}\rho(\tau)$.
To dowodzi poprawności rekurencji
pozaskończonej dla \emph{par nazw},
również gdy nazwy są nieskończone.
\end{uwaga}

\begin{lemat}[Monotoniczność atomowego wymuszania]
\label{lem:forcing33-atomic-monotone}
Niech $\theta$ oznacza
$\sigma\in\tau$ albo $\sigma=\tau$.
Jeżeli $p\Vdash\theta$ i $q\leq p$,
to $q\Vdash\theta$.
\end{lemat}
\begin{proof}
Indukujemy jednocześnie po obu
klauzulach Definicji~\ref{def:forcing33-atomic}.
Każde wzmocnienie warunku $q$ jest
także wzmocnieniem $p$.
Wymagane przez klauzule świadectwa
istnieją więc już z założenia
$p\Vdash\theta$. Rekurencja po
rangach zapewnia, że przywołane
relacje pomocnicze są określone.
\end{proof}

\begin{lemat}[Domknięcie atomowe na gęstość]
\label{lem:forcing33-atomic-dense}
Niech $\theta$ będzie atomową formułą
z parametrami nazwami z $M$,
a $p\in\mathbb P$. Jeśli zbiór
$\{q\leq p:q\Vdash\theta\}$ jest
gęsty poniżej $p$, to $p\Vdash\theta$.
\end{lemat}
\begin{proof}
W przypadku $\theta=\sigma\in\tau$
dla każdego $q\leq p$ wybieramy
$r\leq q$ wymuszające przynależność.
Pierwsza klauzula dla $r$ dostarcza
dalszego wzmocnienia i pary z $\tau$,
a więc spełniona jest pierwsza
klauzula dla $p$.
Gdy $\theta=\sigma=\tau$, ustalmy
$(\pi,s)\in\sigma$ i $q\leq p,s$.
Gęstość daje $t\leq q$ wymuszające
$\sigma=\tau$. Z klauzuli równości
dla $t$ i warunku $t\leq s$
otrzymujemy $r\leq t$ wymuszające
$\pi\in\tau$. Analogicznie otrzymujemy
drugą połowę klauzuli równości.
\end{proof}

\section{Klauzule logiczne i definiowalność}
Każdą formułę pierwszego rzędu języka
teorii mnogości przedstawiamy przy użyciu
$=$, $\in$, $\neg$, $\land$ oraz $\exists$.
Pozostałe spójniki i $\forall$
są skrótami składniowymi.

\begin{definicja}[Wymuszanie formuł złożonych]
\label{def:forcing33-formulas}
Dla ustalonych formuł $\varphi,\psi$
z parametrami nazwami z $M$ określamy:
\[
\begin{aligned}
p\Vdash\neg\varphi
&\iff \neg\exists q\leq p\ (q\Vdash\varphi),\\
p\Vdash(\varphi\land\psi)
&\iff (p\Vdash\varphi)\land(p\Vdash\psi),\\
p\Vdash\exists x\,\varphi(x)
&\iff \forall q\leq p\ \exists r\leq q\
 \exists\tau\in M\
 \bigl(\tau\text{ jest nazwą}\land
                 r\Vdash\varphi(\tau)\bigr).
\end{aligned}
\]
Ostatni kwantyfikator po nazwach
oznacza kwantyfikację \emph{wewnątrz $M$}
po definiowalnej klasie nazw,
nie po utworzonym zbiorze wszystkich nazw.
\end{definicja}

\begin{lemat}[Monotoniczność i regularność wymuszania]
\label{lem:forcing33-regular}
Dla dowolnej ustalonej formuły
$\varphi(\vec\tau)$ z nazwami z $M$
oraz $p,q\in\mathbb P$:
\begin{enumerate}
\item z $q\leq p$ i $p\Vdash\varphi(\vec\tau)$
wynika $q\Vdash\varphi(\vec\tau)$;
\item jeśli
$\{q\leq p:q\Vdash\varphi(\vec\tau)\}$
jest gęsty poniżej $p$, to
$p\Vdash\varphi(\vec\tau)$.
\end{enumerate}
\end{lemat}
\begin{proof}
Pierwszy punkt dowodzimy indukcyjnie
po budowie formuły, zaczynając od
Lematu~\ref{lem:forcing33-atomic-monotone}.
Dla negacji zmniejszamy zakres
kwantyfikacji nad warunkami,
dla koniunkcji stosujemy indukcję
do obu składników, a w przypadku
egzystencji posługujemy się tą samą
gęstością poniżej każdego wzmocnienia.

Drugi punkt dla formuł atomowych jest
Lematem~\ref{lem:forcing33-atomic-dense}.
Dla koniunkcji jej wymuszanie
na zbiorze gęstym oznacza wymuszanie
każdego składnika na zbiorze gęstym;
hipoteza indukcyjna kończy dowód.
Jeżeli gęsto wiele warunków poniżej
$p$ wymusza $\neg\psi$, a istnieje
$r\leq p$ z $r\Vdash\psi$,
to jakieś $s\leq r$ wymusza negację.
Monotoniczność daje równocześnie
$s\Vdash\psi$, co przeczy definicji
negacji. Zatem $p\Vdash\neg\psi$.
Dla egzystencji: jeśli gęsto
wiele $r\leq p$ wymusza
$\exists x\,\psi(x)$, to dla każdego
$q\leq p$ możemy wybrać takie
$r\leq q$, a następnie $s\leq r$
i nazwę $\tau$ z
$s\Vdash\psi(\tau)$. Jest to
klauzula egzystencjalna dla $p$.
\end{proof}

\begin{twierdzenie}[Lemat definiowalności]
\label{thm:forcing33-definability}
Niech $M$ będzie tranzytywnym
modelem ZFC, a
$(\mathbb P,\leq,1_{\mathbb P})\in M$
pojęciem wymuszania.
Dla każdej \emph{ustalonej}
formuły pierwszego rzędu
$\varphi(x_1,\ldots,x_n)$ istnieje
formuła języka teorii mnogości
$\operatorname{Force}_\varphi
(p,\tau_1,\ldots,\tau_n,\mathbb P)$,
która w $M$ wyraża relację
$p\Vdash\varphi(\tau_1,\ldots,\tau_n)$
dla $p\in\mathbb P$ i nazw
$\tau_1,\ldots,\tau_n\in M$.
W konsekwencji zbiory warunków
opisane za pomocą ustalonej relacji
wymuszania i parametrów z $M$
należą do $M$.
\end{twierdzenie}
\begin{proof}
Rekurencja atomowa
z Definicji~\ref{def:forcing33-atomic}
jest dobrze ufundowana dzięki
Uwadze~\ref{rem:forcing33-ranks}.
Dla ustalonej pary nazw można
zamknąć zbiór obu nazw na
występujące dziedzicznie pierwsze
współrzędne par. To jest zbiór
należący do $M$ (wykorzystujemy
rekurencję, sumę i zastępowanie).
Rekurencja po odpowiednich parach
nazw i wszystkich warunkach
$p\in\mathbb P$ tworzy w $M$
jednoznaczny wynik. Dla danej pary
zdanie „istnieje poprawny wynik
tej rekurencji, w którym wskazana
relacja zachodzi” jest formułą
pierwszego rzędu. Jednoznaczność
gwarantuje zgodność wyników
uzyskanych z różnych ograniczonych
zbiorów nazw. Tak otrzymujemy
definiowalność relacji atomowych.

Następnie przeprowadzamy indukcję
po \emph{ustalonej formule}
$\varphi$. Negacja i koniunkcja
używają jedynie znanych już
formuł definiujących wymuszanie
składników. W przypadku
$\exists x\,\psi(x)$ kwantyfikator
po nazwach jest kwantyfikatorem
pierwszego rzędu wewnątrz $M$
z definiowalnym warunkiem
bycia nazwą. Otrzymujemy zatem
formułę $\operatorname{Force}_\varphi$.
Wyróżnianie wewnątrz $M$
tworzy wymagane podzbiory
$\mathbb P$.
\end{proof}

\begin{uwaga}[Schemat, nie predykat prawdy]
\label{rem:forcing33-tarski}
Lemat definiowalności dotyczy
\emph{każdej ustalonej formuły}
oddzielnie; nie jest wewnętrzną
definicją jednego predykatu prawdy
dla wszystkich kodów zdań.
Dla wymuszania jednoelementowego
ten ostatni postulat przeczyłby
twierdzeniu Tarskiego o
niedefiniowalności prawdy.
\end{uwaga}

\section{Lemat prawdy i twierdzenie o wymuszaniu}

\begin{lemat}[Gęstość rozstrzygnięć]
\label{lem:forcing33-decision}
Dla ustalonej formuły $\varphi(\vec\tau)$
z nazwami z $M$ zbiór
\[
D_\varphi=\{p\in\mathbb P:
p\Vdash\varphi(\vec\tau)
\text{ lub }p\Vdash\neg\varphi(\vec\tau)\}
\]
jest gęsty w $\mathbb P$ i należy do $M$.
\end{lemat}
\begin{proof}
Przynależność do $M$ wynika z
Twierdzenia~\ref{thm:forcing33-definability}.
Dla danego $p\in\mathbb P$ albo
istnieje $q\leq p$ wymuszające
$\varphi$, albo żadne wzmocnienie
$p$ nie wymusza $\varphi$.
W pierwszym przypadku $q\in D_\varphi$,
w drugim $p\Vdash\neg\varphi$,
więc $p\in D_\varphi$.
\end{proof}

\begin{twierdzenie}[Atomowy lemat prawdy]
\label{thm:forcing33-atomic-truth}
Niech $\sigma,\tau\in M$ będą
nazwami, a $G$ filtrem $M$-generycznym.
Dla $\theta$ postaci $\sigma\in\tau$
lub $\sigma=\tau$ zachodzi
\[
M[G]\models\theta(\sigma^G,\tau^G)
\iff\exists p\in G\
(p\Vdash\theta(\sigma,\tau)).
\]
\end{twierdzenie}
\begin{proof}
Przeprowadzamy jednoczesną indukcję
po sumie naturalnej rang par nazw
z Uwagi~\ref{rem:forcing33-ranks}.

\emph{Przynależność, od prawej do lewej.}
Jeśli $p\in G$ wymusza $\sigma\in\tau$,
to zbiór $q\leq p$, dla których istnieje
$(\pi,s)\in\tau$ z
$q\leq s$ i $q\Vdash\sigma=\pi$,
jest gęsty poniżej $p$ i należy do $M$.
Z Lematu~\ref{lem:forcing33-local-generic}
wybieramy taki $q\in G$.
Ponieważ $q\leq s$, mamy $s\in G$.
Indukcja dla równości daje
$\sigma^G=\pi^G\in\tau^G$.

\emph{Przynależność, od lewej do prawej.}
Jeśli $\sigma^G\in\tau^G$, to
istnieje $(\pi,s)\in\tau$
z $s\in G$ i $\pi^G=\sigma^G$.
Indukcja dla równości daje
$p\in G$ wymuszające $\sigma=\pi$.
Wspólne wzmocnienie $q\in G$
warunków $p,s$ wymusza $\sigma=\pi$.
Każde $r\leq q$ pozostaje
silniejsze od $s$ i wymusza
tę samą równość; zatem
$q\Vdash\sigma\in\tau$.

\emph{Równość, od prawej do lewej.}
Jeśli $p\in G$ wymusza
$\sigma=\tau$ i
$x=\pi^G\in\sigma^G$,
wybieramy $(\pi,s)\in\sigma$,
$s\in G$, i wspólne
wzmocnienie $q\in G$ poniżej $p,s$.
Klauzula równości daje gęstość
poniżej $q$ zbioru warunków
wymuszających $\pi\in\tau$.
Z lokalnej generyczności
otrzymujemy taki warunek w $G$;
indukcja daje $x\in\tau^G$.
Symetrycznie $\tau^G\subseteq\sigma^G$.

\emph{Równość, od lewej do prawej.}
Załóżmy $\sigma^G=\tau^G$.
Definiujemy $E\in M$ jako zbiór
warunków $p$ spełniających jeden
z warunków:
\begin{enumerate}
\item $p\Vdash\sigma=\tau$;
\item istnieje $(\pi,s)\in\sigma$
z $p\leq s$ i bez żadnego
$q\leq p$ wymuszającego $\pi\in\tau$;
\item istnieje $(\pi,s)\in\tau$
z $p\leq s$ i bez żadnego
$q\leq p$ wymuszającego $\pi\in\sigma$.
\end{enumerate}
Jest to zbiór definiowalny,
a jego gęstość wynika z definicji:
gdy $p$ nie wymusza równości,
zanegowanie jednej z dwóch klauzul
daje jej świadectwo i wzmocnienie
$q\leq p$ należące do drugiego
lub trzeciego członu $E$.
Wybieramy $p\in G\cap E$.
W przypadku drugim z $p\leq s$
wynika $s\in G$, zatem
$\pi^G\in\sigma^G=\tau^G$.
Indukcja dla przynależności
daje $q\in G$ wymuszające
$\pi\in\tau$. Wspólne wzmocnienie
$p,q$ przeczy drugiemu warunkowi
definiującemu $E$. Trzeci
przypadek odrzucamy analogicznie.
Pozostaje $p\Vdash\sigma=\tau$.
\end{proof}

\begin{twierdzenie}[Lemat prawdy]
\label{thm:forcing33-truth}
Dla ustalonej formuły
$\varphi(x_1,\ldots,x_n)$, nazw
$\tau_1,\ldots,\tau_n\in M$
oraz filtru $M$-generycznego $G$
zachodzi równoważność
\[
M[G]\models\varphi(\tau_1^G,\ldots,\tau_n^G)
\iff\exists p\in G\
(p\Vdash\varphi(\tau_1,\ldots,\tau_n)).
\]
\end{twierdzenie}
\begin{proof}
Indukujemy po budowie formuły.
Podstawą jest
Twierdzenie~\ref{thm:forcing33-atomic-truth}.
Dla koniunkcji z prawdziwości
obu składników otrzymujemy warunki
$p,q\in G$ wymuszające składniki,
a ich wspólne wzmocnienie wymusza
koniunkcję. Implikacja odwrotna
jest bezpośrednia.

Gdy $p\in G$ wymusza $\neg\psi$,
prawdziwość $\psi$ w $M[G]$
prowadziłaby do warunku $q\in G$
wymuszającego $\psi$.
Wspólne wzmocnienie $p,q$
przeczy klauzuli negacji.
Odwrotnie, jeśli $M[G]\models\neg\psi$,
generyczność daje $p\in G\cap D_\psi$
z Lematu~\ref{lem:forcing33-decision}.
Nie może on wymuszać $\psi$,
więc wymusza $\neg\psi$.

Jeśli $M[G]\models\exists x\,\psi(x,\vec\tau)$,
świadek ma postać $\rho^G$ dla nazwy
$\rho\in M$. Indukcja daje
$p\in G$ wymuszające $\psi(\rho,\vec\tau)$.
Każde wzmocnienie $p$ wymusza tę
formułę, więc $p$ wymusza egzystencję.
Odwrotnie, gdy $p\in G$ wymusza
$\exists x\,\psi(x,\vec\tau)$, zbiór
\[
D=\{q\leq p:\exists\rho\text{ nazwa w }M\
(q\Vdash\psi(\rho,\vec\tau))\}
\]
należy do $M$ i jest gęsty poniżej $p$.
Lemat~\ref{lem:forcing33-local-generic}
daje $q\in G\cap D$ oraz nazwę $\rho$
wymuszającą odpowiednią formułę.
Indukcja daje
$M[G]\models\psi(\rho^G,\vec\tau^G)$,
a więc prawdziwość egzystencji.
\end{proof}

\begin{twierdzenie}[Pełne twierdzenie o wymuszaniu]
\label{thm:forcing33-main}
Niech $M$ będzie przeliczalnym tranzytywnym
modelem ZFC i
$(\mathbb P,\leq,1_{\mathbb P})\in M$
pojęciem wymuszania. Dla każdej
ustalonej formuły $\varphi(\vec x)$,
nazw $\vec\tau\in M$ oraz $p\in\mathbb P$:
\begin{enumerate}
\item relacja $p\Vdash\varphi(\vec\tau)$
jest definiowalna w $M$;
\item dla każdego filtru $M$-generycznego $G$
zachodzi lemat prawdy:
$M[G]\models\varphi(\vec\tau^G)$
dokładnie wtedy, gdy pewien $q\in G$
wymusza $\varphi(\vec\tau)$;
\item zachodzi równoważność semantyczna
\[
 p\Vdash\varphi(\vec\tau)
 \iff \forall G\,
 \bigl(G\text{ jest $M$-generyczny},\
 p\in G\ \Longrightarrow\
 M[G]\models\varphi(\vec\tau^G)\bigr).
\]
\end{enumerate}
\end{twierdzenie}
\begin{proof}
Punkty pierwszy i drugi to
Twierdzenia~\ref{thm:forcing33-definability}
i~\ref{thm:forcing33-truth}.
Implikacja w prawo w punkcie trzecim
wynika z lematu prawdy.
Jeżeli $p$ nie wymusza formuły,
to na mocy Lematu~\ref{lem:forcing33-regular}
zbiór warunków ją wymuszających
nie jest gęsty poniżej $p$.
Istnieje więc $q\leq p$, poniżej
którego nikt nie wymusza tej formuły,
czyli $q\Vdash\neg\varphi(\vec\tau)$.
Wniosek~\ref{cor:countable-generic}
zapewnia filtr $M$-generyczny $G$
zawierający $q$, a więc i $p$.
Z lematu prawdy wynika
$M[G]\models\neg\varphi(\vec\tau^G)$.
To przeczy prawej stronie równoważności.
\end{proof}

\begin{uwaga}[Brak błędnego koła]
\label{rem:forcing33-no-circularity}
W dowodzie lematu prawdy i twierdzenia
o wymuszaniu nie korzystaliśmy z
$M[G]\models\mathrm{ZFC}$.
Znaliśmy tylko tranzytywny
zbiór $M[G]$ z Rozdziału~\ref{ch:forcing-foundations}
i zwykłą semantykę pierwszego rzędu
dla struktury $(M[G],\in)$.
Dopiero poniżej zastosujemy otrzymane
twierdzenie do sprawdzenia aksjomatów ZFC.
\end{uwaga}

\begin{przyklad}[Rozstrzygnięcie bitu Cohena]
\label{ex:forcing33-cohen-bit}
W $\operatorname{Fn}(\omega,2)$
niech $\dot C$ będzie nazwą
generycznego podzbioru $\omega$
z Przykładu~\ref{ex:forcing-cohen-name}.
Jeżeli $n\in\operatorname{dom}p$,
to $p$ wymusza $\check n\in\dot C$
wtedy i tylko wtedy, gdy $p(n)=1$;
gdy $p(n)=0$, wymusza negację.
Każdy filtr generyczny zawierający $p$
daje tę samą wartość $c_G(n)=p(n)$,
a więc wniosek jest konsekwencją
Twierdzenia~\ref{thm:forcing33-main}.
Jeżeli $n\notin\operatorname{dom}p$,
istnieją niezgodne wzmocnienia
$p\cup\{(n,0)\}$ i $p\cup\{(n,1)\}$.
Mają one przedłużenia do filtrów
generycznych, które nadają
różne wartości współrzędnej $n$;
sam $p$ nie rozstrzyga tej formuły.
\end{przyklad}


\chapter{Hipoteza continuum i zastosowania wymuszania}
\label{ch:forcing-continuum}
\chaptermark{Wymuszanie i hipoteza continuum}

\section{Cel i zakres argumentu}
Hipoteza continuum (CH) jest zdaniem
$2^{\aleph_0}=\aleph_1$. Równoważnie:
nie istnieje moc ściśle pomiędzy
$\aleph_0$ a mocą $\mathcal P(\omega)$.
W poprzednim rozdziale uzasadniliśmy,
że dla przeliczalnego tranzytywnego
modelu zbiorowego $M$ teorii ZFC,
pojęcia wymuszania $\mathbb P\in M$
i filtru $M$-generycznego $G$,
rozszerzenie $M[G]$ również spełnia ZFC.

Pokażemy dwie odmienne konstrukcje.
Pierwsza dodaje wiele ciągów binarnych
i zachowuje liczby kardynalne modelu
wyjściowego, uzyskując $\neg\mathrm{CH}$.
Druga zmienia moce większych zbiorów
bez dodawania nowych podzbiorów $\omega$,
co pozwala uzyskać CH.

\begin{uwaga}[Perspektywa wewnętrzna i zewnętrzna]
\label{rem:forcing34-perspective}
W całym rozdziale $M$ jest
\emph{przeliczalnym tranzytywnym modelem
zbiorowym ZFC}. Własności
takie jak „nieprzeliczalny”,
„$\omega_1$”, „$\omega_2$” i
„warunek przeliczalny”,
jeśli dotyczą obiektów $M$,
są rozumiane \emph{wewnątrz $M$}.
Przykładowo $\omega_2^M$ jest
liczbą porządkową należącą do $M$,
choć zewnętrznie jest przeliczalna.
Warunki przeliczalności modelu $M$
używamy wyłącznie zewnętrznie, aby
otrzymać filtr generyczny.
Istnienia takiego $M$ nie wyprowadzamy
z samej $\operatorname{Con}(\mathrm{ZFC})$.
\end{uwaga}

\section{Liczby porządkowe w rozszerzeniu}

\begin{lemat}[Wymuszanie nie dodaje liczb porządkowych]
\label{lem:forcing34-no-ordinals}
Niech $(\mathbb P,\leq,1_{\mathbb P})\in M$
będzie pojęciem wymuszania, a
$G\subseteq\mathbb P$ filtrem
$M$-generycznym. Wtedy
\[
 \operatorname{Ord}^{M[G]}
   =\operatorname{Ord}^M,
\]
przy czym równość oznacza, że
każda liczba porządkowa należąca
do $M[G]$ należy do $M$ i odwrotnie.
\end{lemat}
\begin{proof}
Włączenie od prawej do lewej jest
natychmiastowe, bo $M\subseteq M[G]$
i oba modele są tranzytywne.
Dla drugiego włączenia wybierzmy
nazwę $\tau\in M$ o wartości
$\tau^G=\alpha$, gdzie $\alpha$
jest liczbą porządkową w $M[G]$.
Z hierarchii nazw poprzedniego
rozdziału otrzymujemy w $M$
liczbę porządkową $\gamma$
ograniczającą rangi nazw
występujących dziedzicznie w $\tau$.
Indukcja po tej hierarchii pokazuje
w otaczającym uniwersum, że
$\operatorname{rank}(\tau^G)\leq\gamma$:
każdy element $\tau^G$ jest wartością
nazwy o mniejszej randze.
Ponieważ $\alpha$ jest liczbą
porządkową, $\operatorname{rank}(\alpha)
=\alpha$, a więc $\alpha\leq\gamma$.
Tranzytywność $M$ i fakt
$\gamma\in M$ dają $\alpha\in M$
(również gdy $\alpha=\gamma$).
Nie otrzymujemy nowych liczb
porządkowych, choć ich \emph{moce}
mogą się zmieniać.
\end{proof}

\section{Warunek przeliczalnych antyłańcuchów}

\begin{definicja}[Antyłańcuch i warunek ccc]
\label{def:forcing34-ccc}
Niech $(\mathbb P,\leq)$ będzie
niepustym pojęciem wymuszania.
Zbiór $A\subseteq\mathbb P$ nazywamy
\emph{antyłańcuchem}, gdy każde dwa
różne elementy $A$ są niezgodne,
czyli nie mają wspólnego wzmocnienia.
Powiemy, że $\mathbb P$ spełnia
\emph{warunek przeliczalnych
antyłańcuchów} (ccc), gdy każdy
jego antyłańcuch jest co najwyżej
przeliczalny. W zastosowaniu do
$\mathbb P\in M$ wymagamy
$M\models\text{„$\mathbb P$ spełnia ccc”}$.
\end{definicja}

\begin{definicja}[System $\Delta$]
\label{def:forcing34-delta-system}
Niech $I$ będzie zbiorem, a
$\mathcal F$ rodziną jego skończonych
podzbiorów. Rodzina $\mathcal F$
jest \emph{systemem $\Delta$ o rdzeniu $R$},
gdy $R\subseteq I$ oraz dla
wszystkich różnych $A,B\in\mathcal F$
mamy $A\cap B=R$.
\end{definicja}

\begin{lemat}[Lemat o systemie $\Delta$]
\label{lem:forcing34-delta}
W ZFC każda nieprzeliczalna rodzina
różnych skończonych zbiorów
zawiera nieprzeliczalną podrodzinę,
która jest systemem $\Delta$.
\end{lemat}
\begin{proof}
Z dowolnej nieprzeliczalnej rodziny
wybieramy podrodzinę mocy $\aleph_1$.
Ponieważ wszystkie jej elementy
są skończone, z regularności
$\aleph_1$ istnieje $n\in\omega$
oraz podrodzina mocy $\aleph_1$
złożona z $n$-elementowych zbiorów.
Indukujemy po $n$.

Dla $n=0$ twierdzenie jest oczywiste;
rodzina różnych zbiorów pustych nie
ma zresztą mocy $\aleph_1$.
Załóżmy, że $n>0$. Jeżeli pewien
element $x$ należy do $\aleph_1$
wielu zbiorów rodziny, usuwamy
$x$ z każdego z nich. Dostajemy
$\aleph_1$ różnych zbiorów o
$n-1$ elementach. Hipoteza
indukcyjna daje system $\Delta$
o rdzeniu $R$, a po ponownym
dołączeniu $x$ otrzymujemy
system o rdzeniu $R\cup\{x\}$.

W przeciwnym przypadku każdy element
należy do co najwyżej przeliczalnie
wielu zbiorów rozpatrywanej rodziny.
Rekurencyjnie przez wszystkie
$\xi<\omega_1$ wybierzemy parami
rozłączne zbiory $F_\xi$.
Przed krokiem $\xi$ wybrano
przeliczalnie wiele zbiorów,
gdyż $\xi<\omega_1$ jest
liczbą porządkową przeliczalną.
Ich suma jest przeliczalna;
każdy jej element występuje
w co najwyżej przeliczalnie
wielu zbiorach rodziny.
Zatem tylko przeliczalnie wiele
zbiorów jest wykluczonych i
można wybrać następny rozłączny.
Otrzymana rodzina mocy $\aleph_1$
jest systemem $\Delta$ o pustym
rdzeniu. To kończy indukcję.
\end{proof}

\begin{definicja}[Dodawanie wielu ciągów Cohena]
\label{def:forcing34-many-cohen}
Niech $\kappa\in M$ będzie niezerową
liczbą kardynalną modelu $M$.
Definiujemy w $M$ pojęcie wymuszania
\[
 \operatorname{Add}(\omega,\kappa)^M
 =\{p:p\text{ jest skończoną funkcją }
       \kappa\times\omega\longrightarrow2\}.
\]
Porządek jest odwrotną inkluzją:
$q\leq p$ wtedy i tylko wtedy,
gdy $q\supseteq p$. Największym
warunkiem jest funkcja pusta.
\end{definicja}

\begin{lemat}[Warunek ccc dla wielu ciągów Cohena]
\label{lem:forcing34-add-ccc}
Niech $M$ będzie tranzytywnym modelem
ZFC, a $\kappa$ niezerową liczbą
kardynalną należącą do $M$.
Wówczas
\[
 M\models
 \text{„$\operatorname{Add}(\omega,\kappa)$
       spełnia ccc”}.
\]
\end{lemat}
\begin{proof}
Prowadzimy cały argument wewnątrz $M$.
Gdyby istniał nieprzeliczalny
antyłańcuch $A$, wybralibyśmy
$\aleph_1$ różnych warunków $p$
należących do $A$.
Ich dziedziny są skończonymi
podzbiorami $\kappa\times\omega$.
Ponieważ dla jednej ustalonej
dziedziny istnieje tylko skończenie
wiele funkcji do $\{0,1\}$,
zbiór wybranych dziedzin jest
nieprzeliczalny. Lemat
\ref{lem:forcing34-delta} daje
nieprzeliczalną podrodzinę
o parami jednakowych przecięciach
dziedzin równych skończonemu
rdzeniowi $R$.
Na $R$ istnieje tylko $2^{|R|}$
możliwych przypisań wartości.
Wybieramy dwie różne funkcje
$p,q$ z tej podrodziny,
które zgadzają się na $R$.
Ponieważ ich dziedziny poza $R$
są rozłączne, $p\cup q$ jest
funkcją i wspólnym wzmocnieniem
$p$ i $q$. Przeczy to temu,
że $A$ jest antyłańcuchem.
\end{proof}

\begin{przyklad}[Zgodność pomimo różnych dziedzin]
\label{ex:forcing34-compatible}
Niech $\kappa$ ma co najmniej
trzy elementy i niech
\[
\begin{aligned}
p&=\{((0,0),1),((1,2),0)\},\\
q&=\{((1,2),0),((2,1),1)\}.
\end{aligned}
\]
Wspólny rdzeń dziedzin to
$\{(1,2)\}$, gdzie oba warunki
przypisują $0$. Dlatego
$p\cup q$ jest ich wspólnym
wzmocnieniem. Gdyby natomiast
$q(1,2)=1$, warunki byłyby
niezgodne. Nie chodzi więc o
samą wielkość dziedzin, lecz
o zgodność funkcji na przecięciu.
\end{przyklad}

\section{Dlaczego ccc zachowuje moce}

\begin{lemat}[Gęste rozstrzyganie wartości porządkowej]
\label{lem:forcing34-decide-ordinal}
Niech $\mathbb P\in M$ będzie pojęciem
wymuszania, $\dot f\in M$ nazwą,
$\beta,\lambda\in M$ liczbami
porządkowymi oraz $p\in\mathbb P$
warunkiem wymuszającym, że
$\dot f$ jest funkcją $\check\beta$
w $\check\lambda$. Dla każdego
$\xi<\beta$ zbiór
\[
 D_\xi=\{q\leq p:\exists\eta<\lambda\
      q\Vdash\dot f(\check\xi)=\check\eta\}
\]
jest gęsty poniżej $p$ i należy
do $M$.
\end{lemat}
\begin{proof}
Jest to podzbiór $\mathbb P$
definiowalny w $M$ przy użyciu
Twierdzenia~\ref{thm:forcing33-definability},
a więc należy do $M$.
Gdyby nie był gęsty poniżej $p$,
istniałby $q\leq p$, poniżej
którego żaden warunek nie wymusza
równości z żadną $\check\eta$.
Z konstrukcji filtru generycznego
w Rozdziale~\ref{ch:forcing-foundations}
wybierzmy filtr $H$ zawierający $q$.
Skoro $p\in H$, w $M[H]$
wartość $\dot f^H(\xi)$ jest
elementem $\lambda$, a zatem
równym pewnej $\eta<\lambda$
należącej do $M$. Lemat prawdy
daje $r\in H$ wymuszające
$\dot f(\check\xi)=\check\eta$.
Wspólne wzmocnienie $r$ i $q$
należałoby do $D_\xi$,
sprzecznie z wyborem $q$.
\end{proof}

\begin{twierdzenie}[Zachowanie liczb kardynalnych przez ccc]
\label{thm:forcing34-ccc-preserves}
Niech $M$ będzie przeliczalnym
tranzytywnym modelem ZFC,
a $\mathbb P\in M$ pojęciem
wymuszania spełniającym ccc
w sensie $M$. Dla każdego filtru
$M$-generycznego $G$ każda
liczba kardynalna modelu $M$
pozostaje liczbą kardynalną
modelu $M[G]$. W szczególności
$\omega_1^{M[G]}=\omega_1^M$
i $\omega_2^{M[G]}=\omega_2^M$.
\end{twierdzenie}
\begin{proof}
Z Lematu~\ref{lem:forcing34-no-ordinals}
wiemy, że $M[G]$ i $M$
mają te same liczby porządkowe.
Załóżmy, że pewna nieprzeliczalna
liczba kardynalna $\lambda$ z $M$
przestała być kardynalna w $M[G]$.
Wtedy dla pewnego $\beta<\lambda$
istnieje w $M[G]$ surjekcja
$f:\beta\to\lambda$.
Wybierzmy nazwę $\dot f\in M$
tej funkcji i warunek $p\in G$
wymuszający, że $\dot f$ jest
taką surjekcją.

Dla każdego $\xi<\beta$
rozpatrujemy w $M$ zbiór
$D_\xi$ z Lematu~\ref{lem:forcing34-decide-ordinal}.
Jest gęsty poniżej $p$,
więc na mocy aksjomatu wyboru
można wybrać maksymalny
antyłańcuch $A_\xi\subseteq D_\xi$,
\emph{maksymalny poniżej $p$}.
Jego istnienie uzasadnia się
lematem Kuratowskiego--Zorna:
porządek na antyłańcuchach
przez inkluzję ma górne
ograniczenia dla łańcuchów.
Wobec ccc każdy $A_\xi$
jest co najwyżej przeliczalny
\emph{w $M$}. Niech
\[
 S_\xi=\{\eta<\lambda:
  \exists q\in A_\xi\
  q\Vdash\dot f(\check\xi)=\check\eta\}.
\]
Zatem $S_\xi\in M$ jest co najwyżej
przeliczalny. Zastępowanie w $M$
pozwala utworzyć wszystkie
zbiory $S_\xi$ jednocześnie.
W konsekwencji
$S=\bigcup_{\xi<\beta}S_\xi$
spełnia w $M$ oszacowanie
\[
 |S|\leq\max(\aleph_0,|\beta|)<\lambda.
\]

Dlaczego w każdym rozszerzeniu
generycznym zawierającym $p$
mamy $f(\xi)\in S_\xi$?
Maksymalność $A_\xi$ poniżej
$p$ zapewnia, że zbiór
\[
 E_\xi=\{q:\exists a\in A_\xi\
              q\leq a\}
       \cup\{q:q\perp p\}
\]
jest gęsty w całym $\mathbb P$
i należy do $M$.
Jeżeli $p\in H$ dla filtru
$M$-generycznego $H$,
to $H$ przecina $E_\xi$,
a ponieważ żaden warunek
$H$ nie jest niezgodny z $p$,
filtr $H$ zawiera pewien
$a\in A_\xi$. Warunek $a$
ustala wartość $f^H(\xi)$
w zbiorze $S_\xi$.
Zatem
$\operatorname{ran}f^H\subseteq S$,
wbrew temu, że $f^H$ ma być
surjekcją na $\lambda$.
Sprzeczność kończy dowód.

Liczba $\omega$ pozostaje
kardynalna także w każdym
tranzytywnym modelu ZFC
o tych samych liczbach porządkowych.
Ponieważ wszystkie kardynały
modelu $M$ są zachowane,
jego kolejne kardynały
$\omega_1$ i $\omega_2$
mają tę samą wartość
w $M[G]$.
\end{proof}

\section{Wymuszanie zaprzeczenia CH}

\begin{twierdzenie}[Wiele różnych nowych ciągów]
\label{thm:forcing34-many-reals}
Niech $\kappa\in M$ będzie
nieskończoną liczbą kardynalną,
$\mathbb P=\operatorname{Add}(\omega,\kappa)^M$,
a $G$ filtrem $M$-generycznym.
Wówczas funkcja
$F=\bigcup G$ ma dziedzinę
$\kappa\times\omega$ i wartości
w $\{0,1\}$. Dla każdego
$\alpha<\kappa$ funkcja
$c_\alpha(n)=F(\alpha,n)$
jest elementem $2^\omega\cap M[G]$.
Ponadto
$c_\alpha\neq c_\beta$
dla różnych $\alpha,\beta<\kappa$.
\end{twierdzenie}
\begin{proof}
Skierowanie $G$ zapewnia, że
każde dwa warunki należące
do $G$ są zgodnymi funkcjami.
Dlatego ich suma $F=\bigcup G$
jest funkcją i należy do $M[G]$.
Dla każdej pary $(\alpha,n)$
zbiór warunków określonych
na tej współrzędnej jest
gęsty w $\mathbb P$ i należy
do $M$. Filtr $G$ go przecina,
więc $F$ jest całkowitą funkcją
na $\kappa\times\omega$.

Ustalmy różne $\alpha,\beta<\kappa$.
W modelu $M$ rozpatrzmy
\[
\begin{aligned}
 D_{\alpha,\beta}
  =\{p\in\mathbb P:\exists n<\omega\
  (&(\alpha,n),(\beta,n)
            \in\operatorname{dom}p\\
   &{}\land p(\alpha,n)\ne p(\beta,n))\}.
\end{aligned}
\]
Zbiór jest gęsty: dla dowolnego
skończonego $p$ wybieramy
$n$, dla którego obie współrzędne
nie należą do $\operatorname{dom}p$,
i przedłużamy $p$ o wartości
$0$ i $1$ na tych współrzędnych.
Ponieważ $D_{\alpha,\beta}\in M$,
mamy $G\cap D_{\alpha,\beta}\ne\varnothing$.
Wtedy dla wskazanego $n$
wartości $c_\alpha(n)$ i
$c_\beta(n)$ są różne.
Każde $c_\alpha$ jest podzbiorem
$\omega$ kodowanym przez funkcję
binarną, zatem należy do $M[G]$.
\end{proof}

\begin{twierdzenie}[Generyczny model zaprzeczający CH]
\label{thm:forcing34-not-ch}
Niech $M$ będzie przeliczalnym
tranzytywnym modelem zbiorowym ZFC.
Połóżmy
$\kappa=\omega_2^M$ i
$\mathbb P=\operatorname{Add}(\omega,\kappa)^M$.
Dla każdego filtru $M$-generycznego
$G\subseteq\mathbb P$ zachodzi
\[
 M[G]\models\mathrm{ZFC}+\neg\mathrm{CH}.
\]
Nie trzeba zakładać $M\models\mathrm{CH}$.
\end{twierdzenie}
\begin{proof}
Twierdzenie~\ref{thm:forcing33-zfc}
daje $M[G]\models\mathrm{ZFC}$.
Z Lematu~\ref{lem:forcing34-add-ccc}
i Twierdzenia~\ref{thm:forcing34-ccc-preserves}
wynika, że $\kappa$ pozostaje
$\omega_2$ w modelu $M[G]$.
Natomiast Twierdzenie
\ref{thm:forcing34-many-reals}
daje w $M[G]$ iniekcję
$\alpha\mapsto c_\alpha$
ze zbioru $\kappa$ do $2^\omega$.
Stąd wewnątrz $M[G]$ mamy
\[
 2^{\aleph_0}\geq\aleph_2>\aleph_1.
\]
To jest zaprzeczenie CH.
\end{proof}

\section{Przeliczalna domkniętość i brak nowych ciągów}

\begin{definicja}[Przeliczalna domkniętość]
\label{def:forcing34-countably-closed}
Pojęcie wymuszania $(\mathbb P,\leq)$
jest \emph{przeliczalnie domknięte},
jeżeli dla każdego ciągu
$p_0\geq p_1\geq\cdots$
warunków z $\mathbb P$
istnieje $q\in\mathbb P$
takie, że $q\leq p_n$
dla każdego $n\in\omega$.
Jeśli $\mathbb P\in M$, to
wymagamy tej własności dla
ciągów należących do $M$,
czyli w sensie
$M\models\text{„$\mathbb P$
jest przeliczalnie domknięte”}$.
\end{definicja}

\begin{lemat}[Domkniętość nie dodaje ciągów liczb porządkowych]
\label{lem:forcing34-closed-no-sequences}
Niech $\mathbb P\in M$ będzie
przeliczalnie domkniętym pojęciem
wymuszania, a $G$ filtrem
$M$-generycznym. Jeśli
$f\in M[G]$ jest funkcją
$f:\omega\to\operatorname{Ord}^{M[G]}$,
to $f\in M$. W szczególności
$M[G]$ i $M$ mają dokładnie
te same podzbiory $\omega$
i tę samą liczbę porządkową
$\omega_1$.
\end{lemat}
\begin{proof}
Wybierzmy nazwę $\dot f\in M$
oraz $p\in G$ wymuszające, że
$\dot f$ jest funkcją z $\omega$
w klasę liczb porządkowych.
Z Lematu~\ref{lem:forcing34-no-ordinals}
wiemy, że jej wartości są
liczbami porządkowymi należącymi do $M$.

Dla każdego $n\in\omega$ zbiór
warunków poniżej $p$ ustalających
wartość $\dot f(\check n)$
jako pewną liczbę porządkową
z $M$ jest gęsty poniżej $p$.
Dowód jest identyczny z argumentem
w Lemacie~\ref{lem:forcing34-decide-ordinal},
z tą różnicą, że kwantyfikator
po możliwych liczbach porządkowych
nie ma ustalonego ograniczenia,
lecz jest zwykłym kwantyfikatorem
pierwszego rzędu w $M$.

Wewnątrz $M$ ustalamy
dobre uporządkowanie zbioru
$\mathbb P$ i przez rekurencję
po $n\in\omega$ wybieramy
$p_{n+1}\leq p_n$,
$p_0=p$, tak by $p_{n+1}$
ustalał jedną wartość
$\dot f(\check n)$.
Cały ciąg $(p_n)_{n<\omega}$
należy do $M$ dzięki rekurencji
i zastępowaniu w $M$.
Przeliczalna domkniętość w $M$
daje $q\leq p_n$ dla wszystkich $n$.
Każdy $p_{n+1}$ wymusza
$\dot f(\check n)=\check\gamma_n$
dla jednoznacznej $\gamma_n\in M$.
Monotoniczność zapewnia, że
$q$ wymusza wszystkie te równości.
Ciąg $(\gamma_n)_{n<\omega}$
należy do $M$ przez zastępowanie,
gdyż zbiory warunków
$p_{n+1}$ oraz wymuszane
wartości są definowalne w $M$.
Dla każdego filtru generycznego
zawierającego $q$ interpretacja
$\dot f$ jest równa temu ciągowi.

Zbiór warunków $q\leq p$,
które rozstrzygają cały ciąg
jako ciąg należący do $M$,
jest gęsty poniżej $p$ i należy
do $M$ przez definiowalność
wymuszania. Lemat o lokalnej
generyczności z Rozdziału
\ref{ch:forcing-theorem} pokazuje,
że $G$ zawiera taki warunek.
Zatem $f\in M$.

Każdy podzbiór $\omega$
jest kodowany ciągiem binarnym,
więc nie powstają nowe podzbiory
$\omega$. Gdyby $\omega_1^M$
stało się przeliczalne w $M[G]$,
istniałaby nowa funkcja
z $\omega$ na $\omega_1^M$.
Wykazaliśmy, że taka funkcja
musiałaby należeć do $M$,
co przeczy definicji $\omega_1^M$.
Ponieważ nie ma nowych liczb
porządkowych, otrzymujemy
$\omega_1^{M[G]}=\omega_1^M$.
\end{proof}

\begin{uwaga}[Różnica między domkniętością a ccc]
\label{rem:forcing34-two-properties}
Przeliczalna domkniętość i ccc
wyrażają odmienne mechanizmy.
Pierwsza pozwala uzyskać wspólne
wzmocnienie całego przeliczalnego
ciągu warunków i tym samym
rozstrzygnąć jednocześnie
przeliczalnie wiele zdań.
Druga kontroluje rozmiar zbiorów
parami niezgodnych warunków
i zapobiega zmniejszaniu liczb
kardynalnych. Nie wolno
wnioskować z samej przeliczalnej
domkniętości o zachowaniu
wszystkich liczb kardynalnych:
w następnym przykładzie
celowo zmienimy moc dużego zbioru.
\end{uwaga}

\section{Wymuszanie hipotezy continuum}

\begin{definicja}[Pojęcie wymuszania kolapsu]
\label{def:forcing34-collapse}
Niech $\lambda=(2^{\aleph_0})^M$.
Wewnątrz $M$ definiujemy
$\mathbb Q=\operatorname{Coll}(\omega_1,\lambda)$
jako zbiór wszystkich funkcji częściowych
$p:\omega_1^M\longrightarrow\lambda$,
których dziedzina jest
\emph{co najwyżej przeliczalna w $M$}.
Porządkiem jest odwrotna inkluzja,
a największym warunkiem funkcja pusta.
\end{definicja}

\begin{lemat}[Kolaps jest przeliczalnie domknięty]
\label{lem:forcing34-collapse-closed}
Pojęcie wymuszania
$\operatorname{Coll}(\omega_1,\lambda)$
z Definicji~\ref{def:forcing34-collapse}
jest przeliczalnie domknięte w $M$.
\end{lemat}
\begin{proof}
Niech $(p_n)_{n<\omega}\in M$
będzie malejącym ciągiem
warunków, czyli $p_{n+1}\supseteq p_n$.
Funkcja $q=\bigcup_{n<\omega}p_n$
jest dobrze określona, ponieważ
warunki w ciągu są zgodne.
Z aksjomatu wyboru w $M$
wynika, że suma przeliczalnie
wielu przeliczalnych dziedzin
jest przeliczalna. Stąd
$q\in\mathbb Q$, a $q\leq p_n$
dla każdego $n$.
\end{proof}

\begin{lemat}[Generyczna surjekcja kolapsu]
\label{lem:forcing34-collapse-surjection}
Jeśli $H\subseteq\mathbb Q$
jest filtrem $M$-generycznym,
to $F=\bigcup H\in M[H]$
jest surjekcją
$F:\omega_1^M\longrightarrow\lambda$.
\end{lemat}
\begin{proof}
Dla każdego $\xi<\omega_1^M$
zbiór warunków, których dziedzina
zawiera $\xi$, jest gęsty i należy
do $M$: jeśli $\xi$ nie występuje,
przedłużamy funkcję dowolną
wartością, na przykład $0$.
Dlatego $F$ jest funkcją całkowitą.

Dla każdego $\eta<\lambda$
zbiór warunków, których obraz
zawiera $\eta$, jest także gęsty.
Istotnie, dziedzina dowolnego
warunku jest przeliczalna w $M$,
a $\omega_1^M$ jest w $M$
nieprzeliczalna, więc można
wybrać nową współrzędną $\xi$
i przypisać jej $\eta$.
Generyczność zapewnia
$\eta\in\operatorname{ran}F$.
Ponieważ $\eta$ była dowolna,
$F$ jest surjekcją.
\end{proof}

\begin{twierdzenie}[Generyczny model spełniający CH]
\label{thm:forcing34-ch}
Niech $M$ będzie przeliczalnym
tranzytywnym modelem zbiorowym ZFC,
$\lambda=(2^{\aleph_0})^M$
oraz $\mathbb Q$ pojęciem wymuszania
z Definicji~\ref{def:forcing34-collapse}.
Wtedy dla każdego filtru
$M$-generycznego $H\subseteq\mathbb Q$
zachodzi
\[
 M[H]\models\mathrm{ZFC}+\mathrm{CH}.
\]
\end{twierdzenie}
\begin{proof}
Twierdzenie~\ref{thm:forcing33-zfc}
daje $M[H]\models\mathrm{ZFC}$.
Z Lematu~\ref{lem:forcing34-collapse-closed}
oraz Lematu~\ref{lem:forcing34-closed-no-sequences}
otrzymujemy
\[
 \mathcal P(\omega)^{M[H]}
 =\mathcal P(\omega)^M,
 \qquad
 \omega_1^{M[H]}=\omega_1^M.
\]
Ponieważ w $M$ istnieje bijekcja
między $\lambda$ a $\mathcal P(\omega)^M$,
taka bijekcja należy również
do $M[H]$.
Z Lematu~\ref{lem:forcing34-collapse-surjection}
istnieje tam surjekcja
$\omega_1^{M[H]}\to\lambda$.
Jej złożenie z powyższą bijekcją
jest surjekcją z $\omega_1^{M[H]}$
na wszystkie podzbiory $\omega$
w $M[H]$. Stąd
$2^{\aleph_0}\leq\aleph_1$
wewnątrz $M[H]$.
Z twierdzenia Cantora i ZFC
mamy zawsze
$2^{\aleph_0}\geq\aleph_1$.
Zatem $M[H]\models\mathrm{CH}$.
\end{proof}

\begin{przyklad}[Co zmienia kolaps]
\label{ex:forcing34-collapse-change}
Jeżeli $M\models
2^{\aleph_0}=\aleph_3$, to
$\lambda=\aleph_3^M$.
Po wymuszeniu kolapsu nadal
obowiązuje
$\mathcal P(\omega)^{M[H]}
=\mathcal P(\omega)^M$:
żaden stary ani nowy podzbiór
$\omega$ nie znika i nie
powstaje nowy. Zmienia się
jednak relacja równoliczności:
istnieje nowa surjekcja
$\omega_1^M\to\lambda$.
Liczba porządkowa $\lambda$
pozostaje taka sama, ale nie
jest już mocą $\aleph_3$
w modelu rozszerzonym.
Nie ma tu sprzeczności:
wymuszanie może zmienić
moc starego zbioru bez
zmiany jego elementów.
\end{przyklad}

\section{Wynik modelowy i granice niezależności}

\begin{twierdzenie}[Dwa modele o różnych wartościach CH]
\label{thm:forcing34-two-models}
Jeżeli istnieje przeliczalny
tranzytywny model zbiorowy
$M\models\mathrm{ZFC}$, to
istnieją przeliczalne tranzytywne
modele zbiorowe $M_0$ i $M_1$,
dla których
\[
 M_0\models\mathrm{ZFC}+\mathrm{CH},
 \qquad
 M_1\models\mathrm{ZFC}+\neg\mathrm{CH}.
\]
\end{twierdzenie}
\begin{proof}
Niech $M_0=M[H]$, gdzie
$H$ jest filtrem generycznym
dla kolapsu z
Twierdzenia~\ref{thm:forcing34-ch}.
Niech $M_1=M[G]$, gdzie
$G$ jest filtrem generycznym
dla $\operatorname{Add}(\omega,\omega_2^M)$
z Twierdzenia~\ref{thm:forcing34-not-ch}.
Istnienie filtrów zapewnia
twierdzenie Rasiowej--Sikorskiego
z Rozdziału~\ref{ch:forcing-foundations}.
Oba rozszerzenia są zbiorami
tranzytywnymi spełniającymi ZFC.
Oba rozszerzenia są także zewnętrznie
przeliczalne: każdy ich element jest
wartością pewnej nazwy $\tau\in M$,
a $M$ jest zewnętrznie przeliczalnym
zbiorem. Nie twierdzimy przy tym,
że \emph{wewnątrz} $M$ istnieje
zbiór wszystkich nazw:
nazwy tworzą tam klasę właściwą.
Korzystamy tylko z zewnętrznego
faktu, że zbiór wszystkich nazw
należących do $M$ jest podzbiorem $M$.
Twierdzenia o CH i $\neg\mathrm{CH}$
kończą dowód.
\end{proof}

\begin{uwaga}[Czego wymaga pełna niezależność]
\label{rem:forcing34-consistency}
Właśnie udowodniliśmy implikację
\emph{z istnienia przeliczalnego
tranzytywnego modelu ZFC} do istnienia
dwóch modeli o przeciwnych wartościach
CH. Nie jest to jeszcze dowód obu
implikacji metamatematycznych
\[
\begin{aligned}
 \operatorname{Con}(\mathrm{ZFC})
  &\Longrightarrow
       \operatorname{Con}(\mathrm{ZFC}+\mathrm{CH}),\\
 \operatorname{Con}(\mathrm{ZFC})
  &\Longrightarrow
       \operatorname{Con}(\mathrm{ZFC}+\neg\mathrm{CH}).
\end{aligned}
\]
Istnienie przeliczalnego
tranzytywnego modelu zbiorowego
jest mocniejszym założeniem niż
sama niesprzeczność ZFC.
Pełny wynik o względnej
niesprzeczności można uzyskać
przez formalizację argumentu
wymuszania, modele o wartościach
boolowskich lub odpowiedni
argument metamatematyczny.
Nie przedstawiamy żadnej z tych
metod jako udowodnionej bez
wykonania jej poszczególnych kroków.
Konstrukcja modeli wewnętrznych
z Rozdziału~\ref{ch:constructible}
stanowi drugi istotny kierunek
badań nad CH i GCH.
\end{uwaga}

\chapter{Modele o wartościach boolowskich i względna niesprzeczność CH}
\label{ch:boolean-valued-independence}
\chaptermark{Modele boolowskie i niesprzeczność}

\section{Dlaczego potrzebujemy innej konstrukcji?}

W Rozdziale~\ref{ch:forcing-continuum} wykazaliśmy, że z
\emph{istnienia przeliczalnego tranzytywnego modelu zbiorowego
ZFC} wynikają modele spełniające CH oraz modele spełniające
jej zaprzeczenie. Samo założenie
$\operatorname{Con}(\mathrm{ZFC})$ nie daje jednak takiego
modelu tranzytywnego. Poniżej uzasadnimy silniejszy wynik:
\[
\begin{aligned}
\operatorname{Con}(\mathrm{ZFC})
 &\Longrightarrow \operatorname{Con}(\mathrm{ZFC}+\mathrm{CH}),\\
\operatorname{Con}(\mathrm{ZFC})
 &\Longrightarrow \operatorname{Con}(\mathrm{ZFC}+\neg\mathrm{CH}).
\end{aligned}
\]
Argument będzie przebiegał przez modele o wartościach w
\emph{zupełnych algebrach Boole'a}. W przeciwieństwie do
poprzedniego rozdziału nie wymaga on, by model wyjściowy
był tranzytywny, dobrze ufundowany ani przeliczalny.

\begin{uwaga}[Metateoria i oznaczenia]
\label{rem:boolean35-metatheory}
Rozumujemy w zwykłej metateorii matematycznej opartej
na ZFC. Przez $M\models\mathrm{ZFC}$ rozumiemy dowolną
\emph{strukturę zbiorową} języka $\{\in\}$ spełniającą
wszystkie standardowe aksjomaty ZFC. Zbiór jej elementów
oznaczamy przez $|M|$; relację przynależności w $M$
przez $\in^M$. Nie zakładamy, że $\in^M$ pokrywa się z
zewnętrzną relacją $\in$. Jeśli $b\in|M|$ jest według $M$
algebrą Boole'a, przez $\mathbb B^M$ oznaczamy
\emph{zewnętrzny} zbiór jej elementów z działaniami
interpretowanymi przez $M$. Kresy w $\mathbb B$ są
zupełne \emph{w sensie $M$}; nie twierdzimy, że
istnieją kresy dowolnych zewnętrznych podzbiorów
$\mathbb B^M$.

Formuła $\varphi$ jest zawsze \emph{ustaloną standardową}
formułą pierwszego rzędu. Nie postulujemy pojedynczego
predykatu prawdy dla wszystkich wewnętrznych kodów
formuł w modelu $M$.
\end{uwaga}

\section{Od pojęcia wymuszania do algebry zupełnej}

\begin{definicja}[Algebra Boole'a i zupełność]
\label{def:boolean35-algebra}
Niech $\mathbb B$ będzie zbiorem z działaniami
$\wedge,\vee$, dopełnieniem $b\mapsto\neg b$
i elementami $0\ne1$, spełniającym aksjomaty
algebry Boole'a. Porządek określamy przez
$b\le c$ wtedy i tylko wtedy, gdy $b\wedge c=b$.
Algebra jest \emph{zupełna}, gdy każdy
podzbiór $A\subseteq\mathbb B$ ma kres górny
$\bigvee A$ i dolny $\bigwedge A$ w $\mathbb B$.
Implikację boolowską oznaczamy
$b\Rightarrow c:=\neg b\vee c$.
\end{definicja}

\begin{definicja}[Algebra zbiorów regularnie otwartych]
\label{def:boolean35-regular-open}
Niech $(X,\mathcal T)$ będzie przestrzenią
topologiczną. Zbiór otwarty $U\subseteq X$
nazywamy \emph{regularnie otwartym}, gdy
$U=\operatorname{int}(\overline U)$.
Rodzinę takich zbiorów oznaczamy
$\operatorname{RO}(X)$.
\end{definicja}

\begin{lemat}[Zupełność algebry zbiorów regularnie otwartych]
\label{lem:boolean35-ro-complete}
Dla każdej przestrzeni topologicznej
$(X,\mathcal T)$ zbiór $\operatorname{RO}(X)$
jest zupełną algebrą Boole'a z działaniami
\[
 U\wedge V=U\cap V,\qquad
 \neg U=\operatorname{int}(X\setminus\overline U),
 \qquad
 \bigvee_{i\in I}U_i=
 \operatorname{int}\overline{\bigcup_{i\in I}U_i}.
\]
Zbiór pusty i $X$ są odpowiednio $0$ i $1$.
\end{lemat}
\begin{proof}
Operator $r(A)=\operatorname{int}(\overline A)$
jest idempotentny na zbiorach otwartych:
dla otwartego $U$ zachodzi
$\overline{r(U)}=\overline U$, gdyż
$U\subseteq r(U)\subseteq\overline U$.
Stąd $r(r(U))=r(U)$.
Jeżeli $U,V$ są regularnie otwarte,
to $U\cap V$ jest otwarty i
\[
 U\cap V\subseteq r(U\cap V)
 \subseteq r(U)\cap r(V)=U\cap V.
\]
Wskazany wzór na dopełnienie daje
zbiór regularnie otwarty rozłączny z $U$;
suma boolowska obu jest równa $X$,
ponieważ ich suma jest gęsta w $X$.
Wzór na dowolny kres górny jest
najmniejszym regularnie otwartym
zbiorem zawierającym wszystkie $U_i$.
Kres dolny dowolnej rodziny definiujemy
z praw de Morgana; dla pustej rodziny
są to odpowiednio $0$ i $1$.
Pozostałe tożsamości boolowskie
wynikają z powyższych działań,
w szczególności z rozdzielności
przecięcia względem regularnego
domknięcia sumy zbiorów otwartych.
\end{proof}

\begin{definicja}[Uzupełnienie boolowskie wymuszania]
\label{def:boolean35-completion}
Niech $(\mathbb P,\leq)$ będzie niepustym
pojęciem wymuszania. Na $\mathbb P$
rozpatrujemy topologię, której bazą są
stożki $\mathbb P_{\le p}=\{q:q\le p\}$.
Algebrę
$\mathbb B_{\mathbb P}=\operatorname{RO}(\mathbb P)$
nazywamy jego \emph{uzupełnieniem boolowskim}.
Warunkowi $p$ odpowiada niezerowy element
$i(p)=\operatorname{int}\overline{\mathbb P_{\le p}}$.
Przy porównywaniu wymuszań korzystamy z
dodatnich elementów $\mathbb B_{\mathbb P}$
i ich porządku; równoważne warunki
preorderu mogą dawać ten sam element $i(p)$.
\end{definicja}

\begin{lemat}[Gęstość obrazów warunków]
\label{lem:boolean35-dense-completion}
Dla niepustego pojęcia wymuszania
$(\mathbb P,\leq)$ oraz algebry
$\mathbb B_{\mathbb P}$ z
Definicji~\ref{def:boolean35-completion}
zbiór $\{i(p):p\in\mathbb P\}$
jest gęsty w $\mathbb B_{\mathbb P}^+
=\mathbb B_{\mathbb P}\setminus\{0\}$:
dla każdego $b>0$ istnieje
$p\in\mathbb P$ z $0<i(p)\le b$.
Jeśli $p,q$ mają wspólne wzmocnienie,
to $i(p)\wedge i(q)>0$; jeżeli są niezgodne,
to $i(p)\wedge i(q)=0$.
\end{lemat}
\begin{proof}
Niech $b$ będzie niepustym regularnie
otwartym zbiorem warunków. Wybieramy
$p\in b$. Ponieważ $b$ jest otwarty,
istnieje $q\le p$ takie, że
$\mathbb P_{\le q}\subseteq b$.
Z regularnej otwartości $b$ wynika
$i(q)\subseteq b$. Stożek nie jest pusty,
więc $i(q)>0$. Jeżeli istnieje
$r\le p,q$, to $r$ należy do przecięcia
stożków i jego regularne otwarte
otoczenie świadczy o niezerowości
$i(p)\cap i(q)$. Gdy wspólnego
wzmocnienia nie ma, stożki są rozłączne,
a przecięcie ich regularnych otwartych
domknięć jest puste: gdyby było
niepuste, zawierałoby stożek i na nim
oba pierwotne stożki byłyby gęste,
co dostarczyłoby wspólnego wzmocnienia.
\end{proof}

\section{Nazwy boolowskie i wartości formuł}

\begin{definicja}[Hierarchia nazw boolowskich]
\label{def:boolean35-names}
Niech $\mathbb B$ będzie niezerową
zupełną algebrą Boole'a.
Rekurencyjnie po liczbach porządkowych
definiujemy
\[
 V_0^{\mathbb B}=\varnothing,\quad
 V_{\alpha+1}^{\mathbb B}
 =\{f:\operatorname{dom}f\subseteq
   V_\alpha^{\mathbb B},\
   f:\operatorname{dom}f\to\mathbb B\},
 \quad
 V_\lambda^{\mathbb B}
 =\bigcup_{\alpha<\lambda}V_\alpha^{\mathbb B}.
\]
Elementy klasy
$V^{\mathbb B}=\bigcup_{\alpha\in\operatorname{Ord}}
V_\alpha^{\mathbb B}$ nazywamy
\emph{nazwami boolowskimi}.
Dla $\tau\in V^{\mathbb B}$ zapis
$\tau(\sigma)$ oznacza współczynnik
$\sigma\in\operatorname{dom}\tau$.
Dla $x$ z modelu wyjściowego
nazwa kanoniczna $\check x$ ma dziedzinę
$\{\check y:y\in x\}$ i przypisuje każdemu
$\check y$ wartość $1$.
\end{definicja}

\begin{definicja}[Wartości boolowskie formuł]
\label{def:boolean35-values}
Dla nazw $\sigma,\tau$ definiujemy
jednoczesną rekurencję atomową
\[
\begin{aligned}
\llbracket\sigma\in\tau\rrbracket
 &=\bigvee_{\pi\in\operatorname{dom}\tau}
       \bigl(\tau(\pi)\wedge
             \llbracket\sigma=\pi\rrbracket\bigr),\\
\llbracket\sigma=\tau\rrbracket
 &=\bigwedge_{\pi\in\operatorname{dom}\sigma}
       \bigl(\sigma(\pi)\Rightarrow
             \llbracket\pi\in\tau\rrbracket\bigr)
   \wedge
   \bigwedge_{\pi\in\operatorname{dom}\tau}
       \bigl(\tau(\pi)\Rightarrow
             \llbracket\pi\in\sigma\rrbracket\bigr).
\end{aligned}
\]
Kolejne klauzule dla ustalonych formuł są:
\[
\begin{aligned}
\llbracket\neg\varphi\rrbracket
   &=\neg\llbracket\varphi\rrbracket,\qquad
\llbracket\varphi\land\psi\rrbracket
   =\llbracket\varphi\rrbracket
       \wedge\llbracket\psi\rrbracket,\\
\llbracket\exists x\,\varphi(x)\rrbracket
   &=\bigvee\{\llbracket\varphi(\tau)\rrbracket:
               \tau\in V^{\mathbb B}\}.
\end{aligned}
\]
Pozostałe spójniki i kwantyfikator
ogólny definiujemy zwyczajowo.
\end{definicja}

\begin{lemat}[Poprawność definicji wartości]
\label{lem:boolean35-well-defined}
Niech $\mathbb B$ będzie niezerową
zupełną algebrą Boole'a.
Dla każdej ustalonej formuły
$\varphi(x_1,\ldots,x_n)$ i wszystkich
nazw $\tau_1,\ldots,\tau_n$
wartość $\llbracket
\varphi(\tau_1,\ldots,\tau_n)\rrbracket$
jest jednoznacznie określonym
elementem $\mathbb B$. Operacja ta
jest definiowalna w ZFC przy
ustalonej formule $\varphi$.
\end{lemat}
\begin{proof}
Klauzule atomowe są poprawnie
ufundowane przez rekurencję
po \emph{sumie naturalnej rang}
dwóch nazw, podobnie jak w
Uwadze~\ref{rem:forcing33-ranks}.
Przy każdej rekurencyjnej zależności
jedna z nazw jest zastępowana nazwą
ze zbioru wartości współrzędnych
pierwszej nazwy, o ściśle mniejszej
randze. Kresy w obu klauzulach
dotyczą zbiorów, a ich istnienie
zapewnia zupełność $\mathbb B$.

Pozostaje uzasadnić kwantyfikator
egzystencjalny. Choć
$V^{\mathbb B}$ jest klasą właściwą,
\emph{zbiór różnych wartości}
\[
 S_\varphi=\{b\in\mathbb B:
      \exists\tau\ (\tau\text{ jest nazwą}
        \land b=\llbracket\varphi(\tau)\rrbracket)\}
\]
jest podzbiorem \emph{zbioru}
$\mathbb B$. Powstaje przez
wyróżnianie, gdy klauzula dla
$\varphi$ jest już ustalona.
Zupełność $\mathbb B$ daje
$\bigvee S_\varphi$.
Indukcja po budowie każdej
ustalonej formuły kończy dowód.
Nie twierdzimy, że zbiór
wszystkich nazw jest elementem
uniwersum.
\end{proof}

\begin{lemat}[Tożsamości równości i podstawiania]
\label{lem:boolean35-equality}
Niech $\mathbb B$ będzie niezerową
zupełną algebrą Boole'a i
$\sigma,\tau,\rho$ nazwami.
Wtedy
\[
\begin{aligned}
\llbracket\sigma=\sigma\rrbracket&=1,\\
\llbracket\sigma=\tau\rrbracket
 &=\llbracket\tau=\sigma\rrbracket,\\
\llbracket\sigma=\tau\rrbracket\wedge
 \llbracket\tau=\rho\rrbracket
 &\le\llbracket\sigma=\rho\rrbracket.
\end{aligned}
\]
Dla każdej ustalonej formuły $\varphi(x)$
z dodatkowymi parametrami nazwami mamy
\[
 \llbracket\sigma=\tau\rrbracket\wedge
 \llbracket\varphi(\sigma)\rrbracket
 \le\llbracket\varphi(\tau)\rrbracket.
\]
\end{lemat}
\begin{proof}
Refleksyjność, symetrię i przechodniość
weryfikujemy równoczesną indukcją
po rangach nazw, korzystając z
rozdzielności zupełnej algebry Boole'a:
$b\wedge\bigvee_i c_i
=\bigvee_i(b\wedge c_i)$.
Dla przechodniości najpierw
przechodzimy od współczynnika
$\sigma(\pi)$ do przynależności
$\pi\in\tau$, a następnie do
$\pi\in\rho$; nierówności
$\llbracket\sigma=\tau\rrbracket
 \wedge\llbracket\tau=\rho\rrbracket
 \le\llbracket\sigma=\rho\rrbracket$
wynikają z dwóch klauzul
ograniczających współczynniki
równości. Analogicznie dla
$\llbracket\sigma\in\tau\rrbracket$.

Podstawianie dowodzimy indukcją
po budowie $\varphi$. W przypadku
atomowym stosujemy poprzednie
nierówności i definicję wartości
przynależności. Koniunkcja
zachowuje nierówność przez
monotoniczność $\wedge$, negacja
przez symetrię równości i
odwrócenie porządku przy
dopełnieniu. Dla egzystencji
stosujemy rozdzielność
$\wedge$ względem dowolnych
kresów górnych.
\end{proof}

\section{Mieszanie nazw i lemat o świadkach}

\begin{definicja}[Antyłańcuch w algebrze Boole'a]
\label{def:boolean35-antichain}
Niech $\mathbb B$ będzie algebrą Boole'a.
Zbiór $A\subseteq\mathbb B\setminus\{0\}$
jest \emph{antyłańcuchem boolowskim},
jeżeli $a\wedge a'=0$ dla wszystkich
różnych $a,a'\in A$. Jest maksymalny
poniżej $b\in\mathbb B$, gdy
$a\le b$ dla każdego $a\in A$
oraz $\bigvee A=b$.
\end{definicja}

\begin{lemat}[Mieszanie nazw]
\label{lem:boolean35-mixing}
Niech $\mathbb B$ będzie zupełną
algebrą Boole'a, $A\subseteq\mathbb B^+$
antyłańcuchem i
$(\tau_a)_{a\in A}$ rodziną nazw
boolowskich \emph{indeksowaną zbiorem} $A$.
Wtedy istnieje nazwa $\tau$ taka, że
\[
 a\le\llbracket\tau=\tau_a\rrbracket
 \qquad(a\in A).
\]
\end{lemat}
\begin{proof}
Niech $D=\bigcup_{a\in A}
\operatorname{dom}\tau_a$. Zastępowanie
i suma zapewniają, że $D$ jest zbiorem
nazw o wspólnie ograniczonych rangach.
Definiujemy na $D$ funkcję
\[
 \tau(\pi)=
  \bigvee\{a\wedge\tau_a(\pi):
       a\in A,\ \pi\in\operatorname{dom}\tau_a\}.
\]
Jest to nazwa, ponieważ wszystkie
jej współczynniki należą do $\mathbb B$.
Dla ustalonego $a\in A$ i $\pi\in D$
rozdzielność i rozłączność
antyłańcucha dają
$a\wedge\tau(\pi)=
 a\wedge\tau_a(\pi)$, gdy
$\pi\in\operatorname{dom}\tau_a$,
oraz $a\wedge\tau(\pi)=0$ w przeciwnym
przypadku. Wstawiając te równości
do atomowych klauzul
Definicji~\ref{def:boolean35-values}
i korzystając z refleksyjności
równości z Lematu~\ref{lem:boolean35-equality},
otrzymujemy
$a\le\llbracket\tau=\tau_a\rrbracket$.
\end{proof}

\begin{twierdzenie}[Lemat o pełności wartości]
\label{thm:boolean35-fullness}
Niech $\mathbb B$ będzie niezerową
zupełną algebrą Boole'a.
Dla każdej ustalonej formuły
$\varphi(x,\tau_1,\ldots,\tau_n)$
oraz nazw $\tau_1,\ldots,\tau_n$
istnieje nazwa $\rho$ spełniająca
\[
 \llbracket\exists x\,\varphi(x,\vec\tau)\rrbracket
   =\llbracket\varphi(\rho,\vec\tau)\rrbracket.
\]
\end{twierdzenie}
\begin{proof}
Niech $u=\bigvee S_\varphi$, gdzie
$S_\varphi\subseteq\mathbb B$
jest zbiorem wartości
z Lematu~\ref{lem:boolean35-well-defined}.
Jeżeli $u=0$, dowolna nazwa
jest świadkiem, ponieważ jej wartość
nie przekracza $u$. Załóżmy $u>0$.
Zdefiniujmy
\[
 D=\{a\in\mathbb B^+:
    a\le u,\ \exists b\in S_\varphi\ (a\le b)\}.
\]
Zbiór $D$ jest gęsty poniżej $u$:
jeżeli $0<c\le u$, to z
$c\wedge\bigvee S_\varphi=c>0$
i rozdzielności istnieje
$b\in S_\varphi$ z $c\wedge b>0$;
wówczas $c\wedge b\in D$.
Wybierzmy maksymalny antyłańcuch
$A\subseteq D$, którego kres górny
jest równy $u$.

Dla każdego $a\in A$ istnieje
nazwa $\rho_a$ taka, że
$a\le\llbracket\varphi(\rho_a,\vec\tau)
\rrbracket$. Nie stosujemy wyboru
z klasy właściwej: dla każdego $a$
wybieramy \emph{najmniejszą rangę}
nazwy spełniającej nierówność.
Istnienie minimum wynika ze
znalezienia jednej nazwy i
minimalizacji w zbiorze liczb
porządkowych poniżej jej rangi.
Zastępowanie na zbiorze $A$
ogranicza wszystkie minimalne
rangi wspólną liczbą porządkową
$\gamma$. Teraz wszystkie
kandydatury leżą w jednym
\emph{zbiorze} $V_{\gamma+1}^{\mathbb B}$;
zastosowanie aksjomatu wyboru
daje rodzinę $(\rho_a)_{a\in A}$.

Lemat~\ref{lem:boolean35-mixing}
daje nazwę $\rho$ równą $\rho_a$
z wartością co najmniej $a$.
Z Lematu~\ref{lem:boolean35-equality}
i podstawiania wynika
\[
 a\le\llbracket\varphi(\rho,\vec\tau)\rrbracket
 \quad(a\in A).
\]
Ponieważ $\bigvee A=u$,
mamy $u\le\llbracket\varphi(\rho,\vec\tau)
\rrbracket$. Odwrotna nierówność
wynika z definicji $u$ jako kresu
wszystkich wartości formuły.
\end{proof}

\section{Aksjomaty ZFC w uniwersum boolowskim}

\begin{twierdzenie}[Boolowska poprawność ZFC]
\label{thm:boolean35-zfc}
Niech $\mathbb B$ będzie
niezerową zupełną algebrą Boole'a.
Wtedy dla każdego \emph{ustalonego}
aksjomatu ZFC oraz każdego
ustalonego wystąpienia schematów
wyróżniania i zastępowania,
oznaczonego przez $\theta$,
zachodzi
$\llbracket\theta\rrbracket=1$.
Dotyczy to także wszystkich
aksjomatów logicznych i
reguł poprawnego wnioskowania.
\end{twierdzenie}
\begin{proof}
Poniżej wskazujemy konstrukcje nazw
i uzasadnienie wartości $1$ dla
każdej grupy aksjomatów. W całym
dowodzie korzystamy z
Lematu~\ref{lem:boolean35-equality},
w szczególności z podstawiania.

\emph{Ekstensjonalność.}
Klauzula równości stwierdza, że
każdy element nazwy $\sigma$,
z odpowiednim współczynnikiem,
należy do $\tau$ i odwrotnie.
Użycie obu implikacji boolowskich
oraz podstawiania daje
\[
 \llbracket\forall x\,
 (x\in\sigma\leftrightarrow x\in\tau)
 \rrbracket
 \le\llbracket\sigma=\tau\rrbracket.
\]
W drugą stronę stosujemy podstawianie.
Zatem ekstensjonalność ma wartość $1$.

\emph{Pusty zbiór, para i nieskończoność.}
Pusta funkcja jest nazwą zbioru
pustego. Nazwa o współczynnikach
$1$ przy $\sigma,\tau$ reprezentuje
$\{\sigma,\tau\}$, z uwzględnieniem
możliwego zlania współrzędnych,
gdy nazwy są identyczne.
Nazwa $\check\omega$ reprezentuje
wewnętrznie zbiór induktywny;
definicja $\check x$ zachowuje
przynależność i podstawowe
konstrukcje zbiorowe z wartością $1$.

\emph{Aksjomat sumy.}
Dla nazwy $\sigma$ ustalamy
$D=\bigcup_{\pi\in\operatorname{dom}\sigma}
\operatorname{dom}\pi$. To jest zbiór.
Tworzymy na $D$ nazwę
\[
 U_\sigma(\rho)=
 \bigvee_{\pi\in\operatorname{dom}\sigma}
   \bigl(\sigma(\pi)\wedge
     \pi(\rho)\bigr),
\]
przyjmując $\pi(\rho)=0$ dla
$\rho\notin\operatorname{dom}\pi$.
Klauzule atomowe, rozdzielność
kresów i podstawianie pokazują,
że dla każdej nazwy $\eta$
\[
 \llbracket\eta\in U_\sigma\rrbracket
 =\llbracket
    \exists x\in\sigma\ (\eta\in x)\rrbracket.
\]
To jest aksjomat sumy z wartością $1$.

\emph{Schemat wyróżniania.}
Niech $\varphi(x,\vec\tau)$ będzie
ustaloną formułą, a $\sigma$ nazwą.
Na zbiorze $\operatorname{dom}\sigma$
definiujemy nazwę
\[
 S(\rho)=\sigma(\rho)\wedge
      \llbracket\varphi(\rho,\vec\tau)\rrbracket.
\]
Podstawianie oraz rozdzielność
dają dla każdej nazwy $\eta$
\[
 \llbracket\eta\in S\rrbracket
 =\llbracket\eta\in\sigma\land
                  \varphi(\eta,\vec\tau)\rrbracket.
\]
To potwierdza każdy ustalony
przypadek schematu wyróżniania.

\emph{Aksjomat zbioru potęgowego.}
Niech $D=\operatorname{dom}\sigma$.
W otaczającym modelu wyjściowym
rozpatrujemy \emph{zbiór} wszystkich
funkcji $f:D\to\mathbb B$ takich,
że $f(\rho)\le\sigma(\rho)$
dla każdej $\rho\in D$.
Każda taka funkcja jest nazwą.
Tworzymy nazwę $\Pi_\sigma$,
która każdej z tych funkcji
przypisuje wartość $1$.
Jeżeli $f$ spełnia wskazane
ograniczenia, podstawianie daje
$\llbracket f\subseteq\sigma\rrbracket=1$.

Odwrotnie, niech $\tau$ będzie
dowolną nazwą i niech
$b=\llbracket\tau\subseteq\sigma\rrbracket$.
Definiujemy funkcję na $D$
przez
\[
 f_\tau(\rho)=
   \sigma(\rho)\wedge
   \llbracket\rho\in\tau\rrbracket.
\]
Z aksjomatu ekstensjonalności
w postaci boolowskiej wynika
$b\le\llbracket\tau=f_\tau\rrbracket$:
pod wartością $b$ każdy element
$\tau$ jest elementem $\sigma$,
a jego przynależność można więc
odczytać przez współczynniki
na ustalonej dziedzinie $D$.
Wartość $b$ nie przekracza zatem
$\llbracket\tau\in\Pi_\sigma\rrbracket$.
Obie strony definicji zbioru
potęgowego mają wartość $1$.

\emph{Schemat zastępowania.}
Niech $\varphi(x,y,\vec\tau)$
będzie ustaloną formułą, a
$\sigma$ nazwą. Jeżeli
\[
 b=\llbracket
  \forall x\in\sigma\ \exists!y\
            \varphi(x,y,\vec\tau)\rrbracket,
\]
to dla każdego
$\pi\in\operatorname{dom}\sigma$
Lemat~\ref{thm:boolean35-fullness}
dostarcza nazwy $\rho_\pi$
uzyskującej wartość
$\llbracket\exists y\,
\varphi(\pi,y,\vec\tau)\rrbracket$.
Rodzinę świadków dla \emph{zbioru}
$\operatorname{dom}\sigma$ wybieramy
za pomocą minimalnych rang,
zastępowania i wyboru,
tak samo jak w dowodzie
Lematu~\ref{thm:boolean35-fullness}.
Z nazw $\rho_\pi$ tworzymy
nazwę obrazu o współczynnikach
$\sigma(\pi)$, łącząc powtarzające
się nazwy przez kres górny.
Pod wartością $b$ każde
$x\in\sigma$ ma jedyny obraz
spełniający $\varphi$, a wskazana
nazwa zawiera dokładnie te obrazy.
Dla każdej ustalonej formuły
jest to wymagany przypadek
schematu zastępowania.

\emph{Aksjomat ufundowania.}
Ustalmy nazwę $\sigma$.
Jeżeli $b=\llbracket
\sigma\ne\varnothing\rrbracket>0$,
to pod każdym niezerowym
$c\le b$ istnieje nazwa $\pi$
i niezerowe $d\le c$ takie,
że $d\le\llbracket\pi\in\sigma\rrbracket$.
Wybieramy na odpowiednim
niezerowym fragmencie
\emph{najmniejszą możliwą rangę}
takiego reprezentanta.
Ta minimalizacja jest poprawna,
bo każdą przynależność do $\sigma$
można reprezentować nazwą z
\emph{zbioru} $\operatorname{dom}\sigma$;
zbiór możliwych rang jest więc
ograniczony. Jeśli na niezerowym
fragmencie istniałby jednocześnie
element $\rho\in\pi\cap\sigma$,
atomowa klauzula przynależności
do $\pi$ dawałaby, po ewentualnym
zastąpieniu $\rho$ równą jej nazwą
ze współrzędnych $\pi$, nazwę
o \emph{mniejszej randze}
reprezentującą element $\sigma$
na tym fragmencie. To przeczy
wyborowi najmniejszej rangi.
Wybór maksymalnego antyłańcucha
takich fragmentów i
Lemat~\ref{lem:boolean35-mixing}
sklejają lokalne świadectwa
w jedną nazwę elementu
$\pi\in\sigma$ rozłącznego
z $\sigma$ z wartością co
najmniej $b$. Stąd
\[
 \llbracket\forall x\,
   (x\ne\varnothing\Longrightarrow
    \exists y\in x\ (x\cap y=\varnothing))
 \rrbracket=1.
\]

\emph{Aksjomat wyboru.}
Dla nazwy $\sigma$ zbiór
$D=\operatorname{dom}\sigma$
ma w modelu wyjściowym dobre
uporządkowanie, a więc surjekcję
$e:\kappa\to D$ dla pewnej
liczby porządkowej $\kappa$
(gdy $D$ jest pusty, wybieramy
$\kappa=0$). Nazwa funkcji
$\xi\mapsto e(\xi)$ ma z wartością
$1$ obraz zawierający wszystkie
elementy $\sigma$. Wartości
$\llbracket e(\xi)\in\sigma\rrbracket$
wyznaczają nazwę podzbioru
$\kappa$ indeksów faktycznie
należących do $\sigma$.
Ograniczenie tej funkcji do
owej nazwanej dziedziny jest
surjekcją na $\sigma$ z
wartością $1$. Ponieważ
dziedzinę można dobrze
uporządkować, każdy element
$\sigma$ otrzymuje najmniejszy
indeks swojego przeciwobrazu.
Wcześniej zweryfikowane
wyróżnianie i zastępowanie
pozwalają przenieść to dobre
uporządkowanie na $\sigma$.
Aksjomat wyboru ma więc
wartość $1$.

Pozostają logiczne reguły
rachunku pierwszego rzędu.
Wszystkie aksjomaty logiczne
mają wartość $1$ na mocy
tożsamości algebry Boole'a,
a modus ponens zachowuje
wartość $1$. Reguła uogólniania
zachowuje wartość $1$, ponieważ
kres dolny rodziny wartości
równych $1$ jest równy $1$.
Jest to poprawność boolowska
całego rachunku.
\end{proof}

\chapter{Modele symetryczne i niezależność aksjomatu wyboru}
\label{ch:symmetric-choice}
\chaptermark{Modele symetryczne i aksjomat wyboru}

\section{Dlaczego nie wystarcza zwykłe rozszerzenie generyczne?}

W Rozdziałach~\ref{ch:forcing-foundations}--\ref{ch:forcing-theorem}
zbudowaliśmy rozszerzenia generyczne, a w
Rozdziale~\ref{ch:boolean-valued-independence}
wyprowadziliśmy względną niesprzeczność CH i jej zaprzeczenia
z samej niesprzeczności ZFC.
W zwykłym rozszerzeniu modelu ZFC przez zbiorowe pojęcie
wymuszania aksjomat wyboru pozostaje prawdziwy.
Aby skonstruować model ZF, w którym wybór zawodzi,
trzeba zatem ograniczyć zbiór dopuszczalnych nazw.

Celem niniejszego rozdziału jest dowód implikacji
\[
 \operatorname{Con}(\mathrm{ZFC})\Longrightarrow
 \operatorname{Con}(\mathrm{ZF}+\neg\mathrm{AC})
\]
bez zakładania istnienia przeliczalnego tranzytywnego modelu
ZFC. Po połączeniu z uniwersum konstruowalnym
otrzymamy niezależność AC od ZF w zwykłym
znaczeniu względnej niesprzeczności.
Zasadniczy przykład będzie miał \emph{nieskończony zbiór
bez przeliczalnie nieskończonego podzbioru}.

\begin{uwaga}[Dwa poziomy konstrukcji]
\label{rem:sym36-two-levels}
Najpierw opisujemy budowę w przeliczalnym
tranzytywnym modelu $M$ dla przejrzystości notacji
$\tau^G$. Na końcu formułujemy ten sam argument
wewnątrz dowolnego modelu zbiorowego ZFC,
posługując się ilorazem nazw przez filtr generyczny.
Samo istnienie modelu tranzytywnego jest
założeniem silniejszym niż $\operatorname{Con}(\mathrm{ZFC})$;
nie będziemy go używać w końcowym wniosku.
\end{uwaga}

\section{Rodzina Cohenów i grupa permutacji współrzędnych}

\begin{definicja}[Wymuszanie przeliczalnej rodziny ciągów]
\label{def:sym36-forcing}
W modelu ZFC niech
\[
 \mathbb P=\operatorname{Fn}
       (\omega\times\omega,\{0,1\},<\omega).
\]
Warunkami są funkcje częściowe o \emph{skończonych}
dziedzinach, a $q\le p$ oznacza $q\supseteq p$.
Elementem największym jest funkcja pusta $1_{\mathbb P}$.
Jeżeli $G\subseteq\mathbb P$ jest filtrem generycznym,
to $c=\bigcup G$ jest funkcją
$\omega\times\omega\to\{0,1\}$.
Jej $n$-ty wiersz oznaczamy
$c_n(k)=c(n,k)$.
\end{definicja}

\begin{proof}[Uzasadnienie całkowitości funkcji]
Dla każdej pary $(n,k)$ zbiór
$D_{n,k}=\{p:(n,k)\in\operatorname{dom}p\}$
jest gęsty: każdy warunek można przedłużyć,
przypisując brakującej parze dowolny bit.
Generyczność daje przecięcie $G\cap D_{n,k}$.
Zgodność warunków w filtrze pokazuje,
że suma funkcji z $G$ jest jednoznaczna.
\end{proof}

\begin{definicja}[Działanie grupy i skończony nośnik]
\label{def:sym36-support}
Niech $\Gamma$ będzie grupą wszystkich permutacji
zbioru $\omega$ należących do modelu wyjściowego.
Dla $\pi\in\Gamma$ i $p\in\mathbb P$ określamy
\[
 (\pi p)(\pi(n),k)=p(n,k).
\]
Działanie na nazwach wymuszania definiujemy
rekurencyjnie przez
$\pi\tau=\{(\pi\sigma,\pi p):(\sigma,p)\in\tau\}$.
Dla skończonego $E\subseteq\omega$ oznaczamy
\[
 \operatorname{Fix}(E)
 =\{\pi\in\Gamma:\forall n\in E\ \pi(n)=n\}.
\]
Mówimy, że $E$ jest \emph{nośnikiem} nazwy $\tau$,
gdy $\pi\tau=\tau$ dla wszystkich
$\pi\in\operatorname{Fix}(E)$.
Równość tutaj jest \emph{dosłowną równością nazw},
nie tylko równością ich wartości po interpretacji.
\end{definicja}

\begin{lemat}[Podstawowe własności nośników]
\label{lem:sym36-supports}
Dla dowolnych nazw $\sigma,\tau$, skończonych
$E,F\subseteq\omega$ oraz $\pi\in\Gamma$ zachodzą:
\begin{enumerate}
\item jeśli $E$ jest nośnikiem $\tau$, to każdy
skończony nadzbiór $E$ również jest jej nośnikiem;
\item jeśli $E$ jest nośnikiem $\tau$, a $F$ nośnikiem
$\sigma$, to $E\cup F$ jest wspólnym nośnikiem obu nazw;
\item $\pi\operatorname{Fix}(E)\pi^{-1}
=\operatorname{Fix}(\pi[E])$;
\item jeżeli $\pi\tau=\tau$, to
$\pi[\operatorname{dom}\tau]=\operatorname{dom}\tau$.
\end{enumerate}
\end{lemat}
\begin{proof}
Dla pierwszego punktu wystarczy zauważyć,
że $\operatorname{Fix}(F)\subseteq\operatorname{Fix}(E)$,
gdy $E\subseteq F$. Drugi punkt wynika z zastosowania
tej obserwacji do obu nazw. W trzecim punkcie
permutacja $\rho$ ustala $\pi[E]$ dokładnie wtedy,
gdy $\pi^{-1}\rho\pi$ ustala $E$.
Ostatnia własność wynika z porównania
pierwszych współrzędnych par tworzących
nazwę $\tau$ oraz $\pi\tau$.
\end{proof}

\begin{definicja}[Nazwy hereditarne symetryczne]
\label{def:sym36-hs}
Nazwa $\tau$ jest \emph{symetryczna}, gdy ma skończony
nośnik. Jest \emph{heredytarnie symetryczna},
gdy sama ma skończony nośnik i każda
$\sigma\in\operatorname{dom}\tau$ jest
heredytarnie symetryczna.
Klasę wszystkich takich nazw oznaczamy
$\mathrm{HS}$. Dla filtru generycznego $G$
określamy
\[
 N=\mathrm{HS}^G
   =\{\tau^G:\tau\in\mathrm{HS}\}.
\]
\end{definicja}

\begin{lemat}[Definiowalność i niezmienniczość klasy HS]
\label{lem:sym36-hs-invariant}
Klasa $\mathrm{HS}$ jest definiowalna w ZFC.
Dla każdej $\pi\in\Gamma$ zachodzi
$\pi[\mathrm{HS}]=\mathrm{HS}$.
Jeżeli $\tau\in\mathrm{HS}$ ma nośnik $E$,
to $\pi\tau$ ma nośnik $\pi[E]$.
Każda nazwa kanoniczna $\check x$
dla zbioru $x$ z modelu wyjściowego
należy do $\mathrm{HS}$ i ma pusty nośnik.
\end{lemat}
\begin{proof}
Predykat „$E$ jest skończonym nośnikiem $\tau$”
jest formułą teorii mnogości, ponieważ
działanie na nazwach jest zadane
rekurencją po ich rangach.
Rekurencja po randze definiuje także predykat
„$\tau$ oraz wszystkie jej podnazwy
mają skończony nośnik”.
W szczególności nie trzeba kwantyfikować
po zbiorze wszystkich nazw.

Dla $\pi\in\Gamma$ i $\rho\in\operatorname{Fix}(\pi[E])$
mamy $\pi^{-1}\rho\pi\in\operatorname{Fix}(E)$.
Zatem $\rho(\pi\tau)=\pi\tau$.
Indukcja po randze pokazuje równocześnie,
że wszystkie podnazwy $\pi\tau$ należą do HS.
Działając przez $\pi^{-1}$ dostajemy
odwrotną inkluzję klas.

Dla nazwy kanonicznej
$\check x=\{(\check y,1_{\mathbb P}):y\in x\}$
indukcja po randze $x$ daje
$\pi\check x=\check x$, ponieważ
$\pi1_{\mathbb P}=1_{\mathbb P}$.
Pusty nośnik działa również na wszystkie podnazwy.
\end{proof}

\begin{lemat}[Symetria wymuszania]
\label{lem:sym36-symmetry}
Niech $\pi\in\Gamma$, $p\in\mathbb P$,
$\varphi(x_1,\ldots,x_m)$ będzie ustaloną formułą
języka $\{\in\}$, a $\tau_1,\ldots,\tau_m$
nazwami wymuszania. Wówczas
\[
 p\Vdash_{\mathbb P}\varphi(\vec\tau)
 \quad\Longleftrightarrow\quad
 \pi p\Vdash_{\mathbb P}\varphi(\pi\vec\tau).
\]
\end{lemat}
\begin{proof}
Odwzorowanie $p\mapsto\pi p$
jest automorfizmem porządku $\mathbb P$:
zachowuje porządek, zgodność i wszystkie
zbiory gęste przenosi na zbiory gęste.
W atomowych klauzulach wymuszania
z Rozdziału~\ref{ch:forcing-theorem}
zastępujemy każdą parę $(\sigma,p)$
parą $(\pi\sigma,\pi p)$.
Równoczesna indukcja po rangach nazw
daje równoważność dla $\in$ i $=$.
Indukcja po budowie ustalonej formuły
zachowuje negację i koniunkcję,
ponieważ automorfizm zachowuje wszystkie
wzmocnienia warunku, a egzystencję,
ponieważ przenosi nazwy i gęstość
świadków na odpowiadające im nazwy i gęstość.
Zastosowanie $\pi^{-1}$ daje oba kierunki.
\end{proof}

\section{Wymuszanie ograniczone do nazw symetrycznych}

\begin{definicja}[Relacja wymuszania dla HS]
\label{def:sym36-relative-forcing}
Dla \emph{ustalonej standardowej formuły}
$\varphi$ oraz nazw z $\mathrm{HS}$
definiujemy $p\Vdash_{\mathrm{HS}}\varphi$
rekurencyjnie:
\begin{enumerate}
\item dla formuły atomowej stosujemy zwykłą
relację $p\Vdash_{\mathbb P}\varphi$;
\item $p\Vdash_{\mathrm{HS}}(\varphi\land\psi)$,
gdy wymusza osobno oba człony;
\item $p\Vdash_{\mathrm{HS}}\neg\varphi$
wtedy i tylko wtedy, gdy
\emph{żaden} $q\le p$ nie wymusza
$\varphi$ w sensie HS;
\item $p\Vdash_{\mathrm{HS}}\exists x\,\varphi(x,\vec\tau)$
wtedy i tylko wtedy, gdy
\[
 \forall q\le p\ \exists r\le q\
 \exists\sigma\in\mathrm{HS}\
 r\Vdash_{\mathrm{HS}}\varphi(\sigma,\vec\tau).
\]
\end{enumerate}
Pozostałe spójniki i kwantyfikator ogólny
są skrótami zbudowanymi z podanych operacji.
\end{definicja}

\begin{lemat}[Definiowalność, monotoniczność i symetria]
\label{lem:sym36-relative-properties}
Niech $\varphi(\vec x)$ będzie ustaloną
standardową formułą, a $\vec\tau\in\mathrm{HS}$.
W ZFC predykat
$p\Vdash_{\mathrm{HS}}\varphi(\vec\tau)$
jest definiowalny. Jeżeli $q\le p$ i
$p\Vdash_{\mathrm{HS}}\varphi(\vec\tau)$,
to $q\Vdash_{\mathrm{HS}}\varphi(\vec\tau)$.
Ponadto dla $\pi\in\Gamma$ zachodzi
\[
 p\Vdash_{\mathrm{HS}}\varphi(\vec\tau)
 \iff
 \pi p\Vdash_{\mathrm{HS}}\varphi(\pi\vec\tau).
\]
\end{lemat}
\begin{proof}
Dla atomów korzystamy z
Twierdzenia~\ref{thm:forcing33-definability}
i Lematu~\ref{lem:sym36-symmetry}.
Dla formuły złożonej indukcja po jej budowie
zamienia kwantyfikator „$\sigma\in\mathrm{HS}$”
na zwykły kwantyfikator po zbiorach,
z dodatkiem definiowalnego predykatu HS.
Każdy krok jest zatem formułą ZFC,
choć klasa HS nie jest zbiorem.
Monotoniczność wynika ze zmniejszania
zbioru dopuszczalnych wzmocnień.
Przy egzystencji, jeśli warunek $p$
spełnia kryterium gęstości,
to spełnia je każdy $q\le p$.
Niezmienniczość dla negacji i egzystencji
wynika z bijektywnego przenoszenia
wzmocnień i z Lematu~\ref{lem:sym36-hs-invariant}.
\end{proof}

\begin{twierdzenie}[Lemat prawdy dla modeli symetrycznych]
\label{thm:sym36-truth}
Niech $M$ będzie przeliczalnym tranzytywnym
modelem ZFC, $\mathbb P,\mathrm{HS}$ konstrukcjami
z Definicji~\ref{def:sym36-forcing} oraz
Definicji~\ref{def:sym36-hs} wykonanymi w $M$,
a $G\subseteq\mathbb P^M$ filtrem $M$-generycznym.
Niech $N=\mathrm{HS}^{M,G}$.
Dla każdej ustalonej formuły
$\varphi(x_1,\ldots,x_m)$ i
$\tau_1,\ldots,\tau_m\in(\mathrm{HS})^M$ mamy
\[
 N\models\varphi(\tau_1^G,\ldots,\tau_m^G)
 \iff
 \exists p\in G\
 \bigl(M\models
   p\Vdash_{\mathrm{HS}}\varphi(\vec\tau)\bigr).
\]
\end{twierdzenie}
\begin{proof}
Prowadzimy indukcję po budowie formuły.
Dla formuł atomowych zwykły lemat prawdy
z Rozdziału~\ref{ch:forcing-theorem}
obowiązuje bez zmian: prawdziwość
$\sigma^G\in\tau^G$ lub $\sigma^G=\tau^G$
nie zależy od tego, czy nazwę rozpatrujemy
w $N$, czy w pełnym $M[G]$.

Dla koniunkcji wybieramy wspólne
wzmocnienie dwóch warunków z $G$.
Dla negacji zakładamy najpierw,
że $N\models\neg\psi(\vec\tau)$.
Zbiór warunków, które wymuszają
$\neg\psi$ albo mają wzmocnienie
wymuszające $\psi$, jest gęsty i
należy do $M$ z definiowalności
Lematu~\ref{lem:sym36-relative-properties}.
Generyczność daje warunek z $G$
rozstrzygający jedną z dwóch możliwości;
druga przeczy hipotezie indukcyjnej.
Odwrotnie, gdy $p\in G$ wymusza negację,
żaden warunek z $G$ nie może wymuszać $\psi$:
istniałoby ich wspólne wzmocnienie w $G$.

Jeżeli $N\models\exists x\,\psi(x,\vec\tau)$,
to świadkiem jest $\sigma^G$ dla
pewnej $\sigma\in(\mathrm{HS})^M$.
Indukcja daje $p\in G$ wymuszający
$\psi(\sigma,\vec\tau)$.
Warunki wymuszające tę konkretną
formułę są gęste poniżej $p$
w trywialny sposób (sam $p$
jest odpowiednim wzmocnieniem);
zatem $p$ wymusza egzystencję.
W drugą stronę, jeśli $p\in G$
wymusza egzystencję, to zbiór
wzmocnień $r$ poniżej $p$
mających świadka $\sigma\in\mathrm{HS}$
jest gęsty poniżej $p$ i należy do $M$.
Generyczność $G$ daje taki $r\in G$.
Hipoteza indukcyjna świadczy,
że $N\models\psi(\sigma^G,\vec\tau^G)$.
\end{proof}

\section{Zachowanie aksjomatów ZF}

\begin{twierdzenie}[Twierdzenie o modelu symetrycznym]
\label{thm:sym36-zf}
Przy założeniach
Twierdzenia~\ref{thm:sym36-truth}
klasa $N=\mathrm{HS}^{M,G}$ jest
tranzytywnym modelem wewnętrznym ZF
zawartym w $M[G]$, zawierającym $M$
i mającym te same liczby porządkowe co $M[G]$.
W szczególności $N\models\mathrm{ZF}$.
\end{twierdzenie}
\begin{proof}
Z Lematu~\ref{lem:sym36-hs-invariant}
wynika, że każda nazwa kanoniczna jest w HS,
więc $M\subseteq N$. Jeżeli $x\in y=\tau^G$
dla $\tau\in\mathrm{HS}$, to
$x=\sigma^G$ dla pewnej pary
$(\sigma,p)\in\tau$ z $p\in G$.
Heredytarna symetria daje $\sigma\in\mathrm{HS}$,
czyli $x\in N$; stąd tranzytywność.
Żadne wymuszanie zbiorowe nie dodaje
liczb porządkowych, a wszystkie
liczby porządkowe modelu wyjściowego
pozostają w $N$. Zatem klasa $N$ ma
te same liczby porządkowe co $M[G]$.

Sprawdzimy teraz aksjomaty.
W całym dowodzie „$p$ wymusza”
oznacza relację $\Vdash_{\mathrm{HS}}$.
W szczególności wszystkie nazwy parametrów
należą do HS, a używane
formuły są ustalone i standardowe.

\emph{Ekstensjonalność i ufundowanie.}
Ponieważ $N$ jest klasą tranzytywną
w modelu $M[G]$, oba aksjomaty dziedziczy
po $M[G]$. Dla ufundowania trzeba jedynie
zauważyć, że jeśli $a\in N$ jest niepusty,
to dowolny element $\in$-minimalny
zbioru $a$, znaleziony w $M[G]$,
należy do $N$ dzięki tranzytywności.

\emph{Pusty zbiór, para, suma i nieskończoność.}
Pustą nazwę i $\check\omega$ mamy w HS.
Jeśli $\sigma,\tau\in\mathrm{HS}$
mają nośniki $E,F$, to nazwa pary
$\{(\sigma,1_{\mathbb P}),(\tau,1_{\mathbb P})\}$
ma nośnik $E\cup F$ i wszystkie podnazwy
w HS. Dla sumy nazwy $\sigma$
używamy zwykłej konstrukcji
$\bigcup\sigma$ w postaci nazwy
\[
 \mathrm U(\sigma)=
 \{(\rho,r):\exists(\pi,p)\in\sigma\
       \exists(\rho,q)\in\pi\
       (r\le p\ \land\ r\le q)\}.
\]
Wszystkie jej podnazwy leżą w HS,
a każdy automorfizm ustalający
$\sigma$ ustala także $\mathrm U(\sigma)$.
Jest to nazwa sumy zbiorów
zakodowanych przez elementy $\sigma^G$,
ponieważ filtr generyczny przecina
gęste zbiory wzmocnień rozstrzygających
obie przynależności. To daje aksjomat sumy.

\emph{Wyróżnianie.}
Dla $\tau,\vec\eta\in\mathrm{HS}$
i ustalonej formuły
$\varphi(x,\vec\eta)$ definiujemy
\[
 S_{\varphi,\tau}=
 \{(\rho,p):\rho\in\operatorname{dom}\tau,\
   p\Vdash_{\mathrm{HS}}
    (\rho\in\tau\land\varphi(\rho,\vec\eta))\}.
\]
Zbiór ten istnieje w $M$,
ponieważ $\operatorname{dom}\tau\times\mathbb P$
jest zbiorem i relacja wymuszania
jest definiowalna. Wszystkie jego podnazwy
są w HS, a wspólny nośnik nazw parametrów
ustala $S_{\varphi,\tau}$ na mocy symetrii
relacji wymuszania. Lemat prawdy pokazuje
dla każdego filtru generycznego, że
$S_{\varphi,\tau}^G=
 \{x\in\tau^G:N\models\varphi(x,\vec\eta^G)\}$.
To dowodzi wszystkich instancji
schematu wyróżniania.

\emph{Zbiór potęgowy.}
Niech $D=\operatorname{dom}\tau$ i
\[
 T_\tau=
 \{\sigma\in\mathrm{HS}:
       \sigma\subseteq D\times\mathbb P\}.
\]
To jest \emph{zbiór} w $M$,
gdyż jest podzbiorem zbioru
potęgowego $D\times\mathbb P$.
Definiujemy nazwę
\[
 \dot{\mathcal P}_N(\tau)=
 \{(\sigma,p):\sigma\in T_\tau,\
              p\Vdash_{\mathrm{HS}}
                  \sigma\subseteq\tau\}.
\]
Jeśli $E$ jest nośnikiem $\tau$,
to $\pi[D]=D$ dla $\pi\in\operatorname{Fix}(E)$,
więc $T_\tau$ jest niezmienniczy
i również $\dot{\mathcal P}_N(\tau)$
ma nośnik $E$. Każda nazwa w jej
dziedzinie jest w HS, dlatego
$\dot{\mathcal P}_N(\tau)\in\mathrm{HS}$.
Każdy jej element jest podzbiorem $\tau^G$.

Odwrotnie, niech $x=\sigma^G\in N$
i $x\subseteq\tau^G$.
Definiujemy \emph{ograniczoną nazwę}
\[
 \sigma_*=
 \{(\rho,p):\rho\in D,\
   p\Vdash_{\mathrm{HS}}
    (\rho\in\sigma\land\rho\in\tau)\}.
\]
Nazwę tę ustala każdy automorfizm
ustalający obie nazwy $\sigma,\tau$,
a jej podnazwy należą do HS.
Stąd $\sigma_*\in T_\tau$.
Z lematu prawdy i z faktu,
że wszystkie elementy $\tau^G$
są interpretacjami nazw z $D$,
otrzymujemy
$\sigma_*^G=\sigma^G\cap\tau^G=x$.
Co więcej, już
$1_{\mathbb P}\Vdash_{\mathrm{HS}}
 \sigma_*\subseteq\tau$,
bo każdy warunek występujący w
$\sigma_*$ wymusza przynależność do $\tau$.
Zatem para
$(\sigma_*,1_{\mathbb P})$ należy do
$\dot{\mathcal P}_N(\tau)$.
Każdy podzbiór $\tau^G$ należący
do $N$ występuje więc w tej samej
nazwie zbiorowej. Jest to wymagany
zbiór potęgowy \emph{obliczony wewnątrz $N$};
nie należy go mylić z całym
$\mathcal P(\tau^G)$ w $M[G]$.

\emph{Schemat zbierania (Collection).}
Niech $\varphi(x,y,\vec\eta)$ będzie
ustaloną formułą, $\tau,\vec\eta\in\mathrm{HS}$.
W $M$ rozpatrujemy \emph{zbiór} par
$(p,\rho)$, gdzie $p\in\mathbb P$,
$\rho\in\operatorname{dom}\tau$ i
\[
 p\Vdash_{\mathrm{HS}}
   \bigl(\rho\in\tau\land
          \exists y\,\varphi(\rho,y,\vec\eta)\bigr).
\]
Dla każdej takiej pary relacja
wymuszania egzystencji zapewnia,
że poniżej $p$ istnieje warunek $q$
i nazwa $\sigma\in\mathrm{HS}$ taka, że
$q\Vdash_{\mathrm{HS}}\varphi(\rho,\sigma,\vec\eta)$.
Wybieramy \emph{najmniejszą rangę}
możliwej nazwy $\sigma$.
Dla ustalonej pary minimum można
znaleźć, wybierając jednego świadka
i minimalizując w zbiorze liczb
porządkowych nie większych od jego rangi.
Zastępowanie w $M$ ogranicza te rangi
jedną liczbą porządkową $\alpha$,
tak że każdy rozpatrywany $(p,\rho)$
ma jakieś $q\le p$ i świadka
$\sigma\in\mathrm{HS}$ o randze $<\alpha$.

Niech $V_\alpha^{\mathbb P}$ oznacza
zbiór nazw wymuszania o randze $<\alpha$
(zgodnie z rekurencją nazw
z Rozdziału~\ref{ch:forcing-foundations}).
Definiujemy
\[
 \dot C=
 \{(\sigma,p):
   \sigma\in \mathrm{HS}\cap V_\alpha^{\mathbb P},\
   p\Vdash_{\mathrm{HS}}
      \exists x\in\tau\
        \varphi(x,\sigma,\vec\eta)\}.
\]
Jest to nazwa zbiorowa w $M$.
Wspólny nośnik $\tau,\vec\eta$ ustala
jej dziedzinę i warunek przynależności,
bo automorfizmy zachowują rangi;
podnazwy należą do HS.
Zatem $\dot C\in\mathrm{HS}$.

Załóżmy teraz w $N$, że
$\forall x\in\tau^G\ \exists y\,\varphi(x,y,\vec\eta^G)$.
Dla ustalonego $x=\rho^G\in\tau^G$
lemat prawdy dostarcza $p\in G$
wymuszający jednocześnie
$\rho\in\tau$ i istnienie obrazu.
Dla \emph{każdego} $r\le p$
istnieje $q\le r$ i świadek
$\sigma$ rangi $<\alpha$, ponieważ
$r$ wymusza te same dwa warunki
i jest jedną z par uwzględnionych
przy ograniczaniu rang.
Zbiór takich $q$ jest gęsty poniżej $p$
i należy do $M$. Filtr $G$ przecina
ten zbiór w pewnym $q$.
Wtedy $q$ wymusza równocześnie
$\rho\in\tau$ i $\varphi(\rho,\sigma,\vec\eta)$,
zatem $(\sigma,q)\in\dot C$.
W konsekwencji $\sigma^G\in\dot C^G$
jest wymaganym świadkiem dla $x$.
Ponieważ każdy element $\dot C^G$
z definicji jest wartością obrazu
jakiegoś $x\in\tau^G$, otrzymujemy
schemat zbierania. Po zastosowaniu
wyróżniania daje on także
schemat zastępowania.

Sprawdziliśmy zatem ekstensjonalność,
pusty zbiór, parę, sumę, nieskończoność,
zbiór potęgowy, wyróżnianie,
zastępowanie i ufundowanie.
Nie używaliśmy aksjomatu wyboru
\emph{wewnątrz} $N$.
\end{proof}

\section{Nieskończony zbiór Dedekinda skończony}

\begin{definicja}[Zbiór Dedekinda skończony]
\label{def:sym36-dedekind}
Zbiór $A$ jest \emph{Dedekinda skończony},
jeśli nie istnieje iniekcja $f:\omega\to A$.
W ZF zbiór może być nieskończony
(nie być równoliczny z żadną liczbą naturalną)
i jednocześnie Dedekinda skończony.
\end{definicja}

\begin{lemat}[Nazwy wierszy i całej rodziny]
\label{lem:sym36-cohen-family}
Dla każdego $n\in\omega$ niech
\[
 \dot c_n=
 \{(\check k,p):
     (n,k)\in\operatorname{dom}p,\
     p(n,k)=1\},\qquad
 \dot A=\{(\dot c_n,1_{\mathbb P}):n\in\omega\}.
\]
Wtedy $\dot c_n\in\mathrm{HS}$
ma nośnik $\{n\}$, a
$\dot A\in\mathrm{HS}$ ma nośnik pusty.
Dla każdego filtru generycznego $G$
nazwy te reprezentują odpowiednio
podzbiory $\omega$ oraz zbiór
$A=\{c_n:n<\omega\}$;
elementy $c_n$ są parami różne.
\end{lemat}
\begin{proof}
Permutacja ustalająca $n$ nie zmienia
pierwszej współrzędnej $(n,k)$, więc
ustala $\dot c_n$. Wszystkie jej podnazwy
$\check k$ mają pusty nośnik.
Ponadto $\pi\dot c_n=\dot c_{\pi(n)}$,
wobec czego działanie dowolnej $\pi$
jedynie przestawia elementy nazwy $\dot A$.
Zatem $\dot A$ ma pusty nośnik,
a jej podnazwy są w HS.

Dla ustalonych $n\ne m$ zbiór
\[
 D_{n,m}=\{p:\exists k\
 [(n,k),(m,k)\in\operatorname{dom}p
   \land p(n,k)\ne p(m,k)]\}
\]
jest gęsty. Dla danego warunku $p$
wybieramy $k$ spoza obu skończonych
zbiorów użytych współrzędnych
i przedłużamy $p$, dopisując
odpowiednio wartości $0$ i $1$.
Generyczność zapewnia, że wiersze
$c_n,c_m$ różnią się na pewnym $k$.
Definicja nazwy $\dot c_n$ i gęstość
warunków rozstrzygających każdy bit
pokazują, że $\dot c_n^G
=\{k:c(n,k)=1\}$. Stąd teza.
\end{proof}

\begin{lemat}[Skończony nośnik ogranicza wartości funkcji]
\label{lem:sym36-finite-range}
Niech $\dot f\in\mathrm{HS}$ ma skończony
nośnik $E\subseteq\omega$ i niech
$p\in\mathbb P$ wymusza w sensie HS,
że $\dot f$ jest funkcją o dziedzinie
$\check\omega$ i wartościach w $\dot A$.
Wówczas $p$ wymusza, że zbiór wartości
$\dot f$ jest podzbiorem
$\{\dot c_k:k\in E\}$.
\end{lemat}
\begin{proof}
Ustalmy $m\in\omega$.
Z definicji \emph{funkcji o wartościach w $\dot A$},
z przynależności do nazwy
$\dot A$ i z gęstości świadków wynika,
że poniżej każdego $q\le p$ istnieje
$r\le q$ i $k\in\omega$ takie, że
\[
 r\Vdash_{\mathrm{HS}}
 \dot f(\check m)=\dot c_k.
\]
Pokażemy, że koniecznie $k\in E$.

Załóżmy przeciwnie, że $k\notin E$.
Ponieważ $r$ jest skończoną funkcją,
istnieje $\ell\notin E\cup\{k\}$
niewystępujące jako pierwsza
współrzędna w $\operatorname{dom}r$.
Niech $\pi$ zamienia $k$ z $\ell$
i ustala pozostałe liczby naturalne.
Permutacja ta ustala $E$, a więc
$\pi\dot f=\dot f$ oraz
$\pi\check m=\check m$.
Lemat~\ref{lem:sym36-relative-properties}
daje
\[
 \pi r\Vdash_{\mathrm{HS}}
  \dot f(\check m)=\dot c_\ell.
\]
Warunki $r$ i $\pi r$ są zgodne:
wszystkie współrzędne różne od
wierszy $k,\ell$ są identyczne,
a na wierszu $\ell$ warunek $r$
nie ma żadnego przypisania.
Mają zatem wspólne wzmocnienie $s$.

Warunek $s$ wymusza jednocześnie
$\dot f(\check m)=\dot c_k$
i $\dot f(\check m)=\dot c_\ell$,
czyli $\dot c_k=\dot c_\ell$.
Przeczy to Lematu~\ref{lem:sym36-cohen-family},
gdyż $k\ne\ell$ i już
$1_{\mathbb P}$ wymusza ich nierówność.
Stąd $k\in E$.

Dla każdego $m$ warunki, które
rozstrzygają $\dot f(\check m)$
jako jeden z wierszy o indeksach w $E$,
są gęste poniżej $p$.
Definicja wymuszania kwantyfikatora
ogólnego i dowiedziony lemat prawdy
dają zapowiedzianą inkluzję
zbioru wartości dla całej funkcji.
\end{proof}

\begin{twierdzenie}[Nieskończony zbiór Dedekinda skończony]
\label{thm:sym36-dedekind}
Przy założeniach
Twierdzenia~\ref{thm:sym36-zf}
zbiór $A=\dot A^G\in N$ jest
nieskończony, lecz w $N$
nie istnieje iniekcja $\omega\to A$.
W szczególności $N\models\neg\mathrm{AC}$.
\end{twierdzenie}
\begin{proof}
Lemat~\ref{lem:sym36-cohen-family}
pokazuje, że $A$ ma przynajmniej
$n$ elementów dla każdego
$n\in\omega$: wystarczy
$c_0,\ldots,c_{n-1}$.
Są one elementami $N$,
bo ich nazwy należą do HS.
Zatem $A$ nie jest zbiorem skończonym.

Przypuśćmy, że w $N$ istnieje
iniekcja $f:\omega\to A$.
Z definicji $N$ wybieramy
$\dot f\in\mathrm{HS}$ z
$\dot f^G=f$ i skończony
nośnik $E$ nazwy $\dot f$.
Lemat prawdy dostarcza $p\in G$
wymuszający, że $\dot f$
jest funkcją $\omega\to A$
i jest iniekcją.
Lemat~\ref{lem:sym36-finite-range}
ogranicza jej zbiór wartości
do skończonego zbioru
$\{c_k:k\in E\}$.
Nie istnieje iniekcja zbioru
$\omega$ w zbiór skończony;
sprzeczność.

W ZF z aksjomatu wyboru wynika,
że każdy nieskończony zbiór
ma podzbiór przeliczalnie nieskończony:
dobrze porządkujemy zbiór i wybieramy
rekurencyjnie pierwszy element,
który nie znalazł się w skończonej
liście wcześniej wybranych elementów.
Otrzymana funkcja z $\omega$ jest
iniekcją. Ponieważ nasz zbiór $A$
jest kontrprzykładem, AC zawodzi w $N$.
\end{proof}

\begin{przyklad}[Dlaczego nie można „po prostu ponumerować” wierszy?]
\label{ex:sym36-indexing}
W pełnym rozszerzeniu $M[G]$
istnieje funkcja $n\mapsto c_n$,
zadana generyczną tablicą bitów.
W modelu $N$ każdy pojedynczy $c_n$
oraz zbiór $A$ istnieją, lecz
funkcja wyliczająca wszystkie $c_n$
nie jest elementem $N$.
Jej zwykła nazwa nie ma skończonego
nośnika: permutacja zmieniająca
indeks spoza skończonego $E$
zmienia odpowiednią parę w grafie.
Ścisłym uzasadnieniem
\emph{nieistnienia jakiejkolwiek}
iniekcji, a nie tylko tej konkretnej
funkcji, jest
Lemat~\ref{lem:sym36-finite-range}.
\end{przyklad}

\section{Od niesprzeczności ZFC do niezależności AC od ZF}

\begin{lemat}[Interpretacja bez tranzytywności modelu wyjściowego]
\label{lem:sym36-arbitrary-model}
Niech $M\models\mathrm{ZFC}$ będzie dowolnym
\emph{przeliczalnym modelem zbiorowym}.
Nie zakładamy zewnętrznego dobrego ufundowania
relacji $\in^M$.
W metateorii ZFC istnieje model zbiorowy
$N\models\mathrm{ZF}+\neg\mathrm{AC}$
zbudowany z wewnętrznych nazw
$(\mathrm{HS})^M$ dla
$(\mathbb P,\Gamma)^M$.
\end{lemat}
\begin{proof}
Ponieważ $|M|$ jest zewnętrznie
przeliczalny, możemy wyliczyć
wszystkie zbiory warunków
$D\subseteq\mathbb P^M$, które
$M$ uważa za gęste, i kolejno
wybierać ich warunki. Używając
skierowania porządku przez wzmocnienia,
konstruujemy \emph{zewnętrzny} filtr
$M$-generyczny $G$ na $\mathbb P^M$.
Nie potrzeba w tym celu zewnętrznej
rekurencji po nieufundowanych rangach nazw.

Dla $\sigma,\tau\in(\mathrm{HS})^M$
definiujemy
\[
 \sigma\equiv_G\tau
 \iff\exists p\in G\
     M\models p\Vdash_{\mathrm{HS}}\sigma=\tau,
\]
oraz
\[
 [\sigma]_G\in_G[\tau]_G
 \iff\exists p\in G\
     M\models p\Vdash_{\mathrm{HS}}\sigma\in\tau.
\]
Zwykła poprawność wymuszania dla
atomów, podstawianie równych nazw
i własności filtru dowodzą,
że pierwsza relacja jest równoważnością,
a druga nie zależy od reprezentantów.
Klasy i relacja są \emph{zbiorami zewnętrznymi},
ponieważ $(\mathrm{HS})^M\subseteq|M|$
jest zewnętrznie zbiorem.
Otrzymujemy więc strukturę zbiorową.

Wszystkie dowody definicji
$\Vdash_{\mathrm{HS}}$, symetrii,
zbierania i potęgowania z
Twierdzenia~\ref{thm:sym36-zf}
są dowodami pierwszego rzędu
\emph{w ZFC modelu wyjściowego}:
kwantyfikują po zbiorowych nazwach,
skończonych nośnikach i porządkach,
a przy ustalonej formule mają
definiowalny predykat wymuszania.
Stąd każdy ich krok pozostaje
prawdziwy, gdy interpretujemy go
w nietranzytywnym $M$.
Indukcja po \emph{standardowej}
formule, analogiczna do
Twierdzenia~\ref{thm:sym36-truth},
daje dla klas nazw lemat prawdy:
\[
 N\models\varphi([\vec\tau]_G)
 \iff\exists p\in G\
       (M\models p\Vdash_{\mathrm{HS}}
          \varphi(\vec\tau)).
\]
Argument o skończonym nośniku
z Lematu~\ref{lem:sym36-finite-range}
również jest dowodem wewnętrznym
ZFC. Otrzymany model spełnia
każdą \emph{standardową}
instancję ZF oraz zaprzeczenie AC.
Jego zewnętrzna relacja przynależności
nie musi być dobrze ufundowana;
do niesprzeczności teorii nie
jest to potrzebne.
\end{proof}

\begin{twierdzenie}[Względna niesprzeczność i niezależność AC]
\label{thm:sym36-relative-consistency}
Dla teorii ZF i ZFC zachodzą implikacje
metamatematyczne
\[
 \operatorname{Con}(\mathrm{ZFC})
 \Longrightarrow
 \operatorname{Con}(\mathrm{ZF}+\neg\mathrm{AC}),
\]
oraz
\[
 \operatorname{Con}(\mathrm{ZF})
 \Longrightarrow
 \begin{cases}
 \operatorname{Con}(\mathrm{ZF}+\mathrm{AC}),\\
 \operatorname{Con}(\mathrm{ZF}+\neg\mathrm{AC}).
 \end{cases}
\]
W szczególności, jeśli ZF jest niesprzeczna,
to ZF nie dowodzi ani AC, ani $\neg$AC.
\end{twierdzenie}
\begin{proof}
Z założenia $\operatorname{Con}(\mathrm{ZFC})$
i twierdzenia Gödla o pełności otrzymujemy
model zbiorowy ZFC. Dolne twierdzenie
Löwenheima--Skolema pozwala wybrać model
przeliczalny. Lemat~\ref{lem:sym36-arbitrary-model}
daje model ZF+$\neg$AC, a więc
niesprzeczność tej teorii.

Z drugiej strony ZF dowodzi, że
uniwersum konstruowalne $L$
spełnia ZFC, co wykazaliśmy w
Rozdziale~\ref{ch:constructible}.
Tym samym
$\operatorname{Con}(\mathrm{ZF})$
implikuje
$\operatorname{Con}(\mathrm{ZFC})
=\operatorname{Con}(\mathrm{ZF}+\mathrm{AC})$.
Składając tę implikację z
pierwszą częścią dowodu, otrzymujemy
niesprzeczność także ZF+$\neg$AC.
Gdyby ZF dowodziła jednego z dwóch
sprzecznych zdań AC, $\neg$AC,
odpowiednia druga teoria byłaby sprzeczna.
\end{proof}

\begin{uwaga}[AC, AD i zakres wniosku]
\label{rem:sym36-ac-ad}
Niezależność AC \emph{od ZF} jest innym
zagadnieniem niż relacja między AC i
aksjomatem determinacji AD.
W Rozdziale~\ref{ch:determinacy}
udowodniliśmy, że ZF+AC wyklucza
pełne AD. Obecna konstrukcja daje
model ZF+$\neg$AC, lecz \emph{nie}
daje automatycznie modelu ZF+AD:
sam brak wyboru nie implikuje determinacji.
Nie utożsamiamy także zbioru Dedekinda
skończonego ze zbiorem skończonym
w zwykłym znaczeniu.
\end{uwaga}

\begin{uwaga}[Granice formalnej weryfikacji]
\label{rem:sym36-formal-status}
Dowody w tym rozdziale są tekstowymi
dowodami matematycznymi, nie wynikiem
kompilacji ich formalizacji w Lean.
Kontrola jednostkowa semantyki nazw,
odnośników i kompilacji LaTeX nie zastępuje
weryfikacji poprawności dowodów przez
jądro asystenta dowodzenia.
\end{uwaga}

\chapter{Duże liczby kardynalne: nieosiągalność i mierzalność}
\label{ch:large-cardinals}
\chaptermark{Kardynały nieosiągalne i mierzalne}

\section{Od niezależności zdań do silniejszych aksjomatów}

W Rozdziałach~\ref{ch:constructible}, \ref{ch:boolean-valued-independence}
i~\ref{ch:symmetric-choice} badaliśmy, jak różne modele teorii
mnogości rozstrzygają CH i aksjomat wyboru. Duże liczby kardynalne
dostarczają innego rodzaju informacji: postulują istnienie obiektów,
których własności wymuszają bogatą strukturę całego uniwersum zbiorów.

Pracujemy w ZFC jako metateorii, chyba że wyraźnie zaznaczymy
inaczej. Wszystkie występujące tutaj kardynały utożsamiamy
z początkowymi liczbami porządkowymi. Symbol $V_\alpha$
oznacza $\alpha$-ty etap hierarchii kumulatywnej, wprowadzony
w Rozdziale~\ref{ch:zf-axioms}.
Nie zakładamy istnienia żadnego nieosiągalnego ani mierzalnego
kardynału: twierdzenia postaci „jeżeli istnieje $\kappa$, to ...”
są implikacjami udowadnianymi w ZFC.

\section{Regularność, silna graniczność i nieosiągalność}

\begin{definicja}[Współkońcowość kardynału]
\label{def:lc37-cofinality}
Niech $\kappa$ będzie nieskończoną liczbą porządkową.
Jej \emph{współkońcowością} $\operatorname{cf}(\kappa)$ nazywamy
najmniejszą liczbę porządkową $\delta$, dla której istnieje
ściśle rosnąca funkcja $f:\delta\to\kappa$ o obrazie
nieograniczonym w $\kappa$.
Kardynał $\kappa$ jest \emph{regularny},
gdy $\operatorname{cf}(\kappa)=\kappa$,
a \emph{singularny}, gdy $\operatorname{cf}(\kappa)<\kappa$.
\end{definicja}

\begin{definicja}[Silna graniczność i kardynały nieosiągalne]
\label{def:lc37-inaccessible}
Kardynał nieskończony $\kappa$ jest \emph{silnie graniczny},
gdy $2^\lambda<\kappa$ dla każdego kardynału $\lambda<\kappa$.
Kardynał $\kappa$ jest \emph{słabo nieosiągalny},
gdy jest nieprzeliczalny, regularny i jest granicznym
kardynałem. Jest \emph{silnie nieosiągalny},
gdy jest nieprzeliczalny, regularny i silnie graniczny.
Bez dodatkowego przymiotnika przez
\emph{kardynał nieosiągalny} rozumiemy w tym rozdziale
kardynał silnie nieosiągalny.
\end{definicja}

\begin{uwaga}[Dlaczego z definicji wyłączamy $\omega$?]
\label{rem:lc37-omega}
Kardynał $\omega$ jest regularny i spełnia
$2^n<\omega$ dla każdej skończonej liczby $n$.
Dlatego bez warunku nieprzeliczalności sam $\omega$
byłby silnie graniczny i regularny, choć nie
jest dużym kardynałem. Również warunek „graniczny”
nie implikuje silnej graniczności:
zapisujemy oba pojęcia osobno.
\end{uwaga}

\begin{lemat}[Kardynał nieosiągalny jest graniczny]
\label{lem:lc37-limit}
Jeżeli $\kappa$ jest silnie nieosiągalny,
to nie jest następnikiem żadnego kardynału.
W szczególności $\kappa$ jest graniczną liczbą porządkową.
\end{lemat}
\begin{proof}
Gdyby $\kappa=\lambda^+$, to twierdzenie Cantora
dawałoby $\lambda^+\le 2^\lambda$, podczas gdy
silna graniczność wymagałaby $2^\lambda<\kappa=\lambda^+$.
Otrzymujemy sprzeczność.
Nieprzeliczalny kardynał, który nie jest następnikiem
kardynału, jest granicznym kardynałem, a więc
graniczną liczbą porządkową.
\end{proof}

\begin{lemat}[Rozmiary etapów hierarchii poniżej $\kappa$]
\label{lem:lc37-small-ranks}
Jeżeli $\kappa$ jest silnie nieosiągalny,
to $|V_\alpha|<\kappa$ dla każdej liczby porządkowej
$\alpha<\kappa$.
W konsekwencji $|x|<\kappa$ dla każdego $x\in V_\kappa$.
\end{lemat}
\begin{proof}
Indukujemy pozaskończenie po $\alpha<\kappa$.
Dla $\alpha=0$ teza jest oczywista.
Jeżeli $|V_\alpha|<\kappa$, to
$V_{\alpha+1}=\mathcal P(V_\alpha)$ i
$|V_{\alpha+1}|=2^{|V_\alpha|}<\kappa$.

Niech $0<\lambda<\kappa$ będzie graniczną liczbą
porządkową. Mamy
$V_\lambda=\bigcup_{\alpha<\lambda}V_\alpha$.
Rodzina ma moc $|\lambda|<\kappa$, a każdy
jej wyraz ma moc mniejszą od $\kappa$.
Ponieważ $\kappa$ jest regularny, supremum
tych mocy jest pewnym $\mu<\kappa$.
Zatem
$|V_\lambda|\le|\lambda|\cdot\mu<\kappa$,
przy czym dla mocy skończonych nierówność
wynika tym bardziej bezpośrednio.
Jeśli $x\in V_\kappa$, to $x\in V_\alpha$
dla pewnego $\alpha<\kappa$ i
$|x|\le|V_\alpha|<\kappa$.
\end{proof}

\begin{twierdzenie}[Nieosiągalny etap jest modelem ZFC]
\label{thm:lc37-rank-model}
Niech $\kappa$ będzie kardynałem silnie nieosiągalnym
w uniwersum spełniającym ZFC. Wówczas zbiorowa,
tranzytywna struktura $(V_\kappa,\in)$ spełnia ZFC.
\end{twierdzenie}
\begin{proof}
Wszystkie elementy elementów $V_\kappa$ nadal
należą do $V_\kappa$, stąd tranzytywność.
Ekstensjonalność i ufundowanie są dziedziczone
z uniwersum. Pusty zbiór i $\omega$ należą do
$V_\kappa$, gdyż $\kappa>\omega$ jest graniczny.

Jeżeli $x,y\in V_\kappa$, to ich para i suma
mają rangę nie większą niż pewna skończona
liczba następczych kroków ponad
$\max\{\operatorname{rank}(x),\operatorname{rank}(y)\}$.
Z graniczności $\kappa$ wynika, że para i suma
należą do $V_\kappa$.
Tak samo $\mathcal P(x)\in V_\kappa$ dla
$x\in V_\kappa$.

Jeżeli formuła $\varphi(x,\vec p)$ z parametrami
$\vec p\in V_\kappa$ definiuje podzbiór
$a\in V_\kappa$, to relacja prawdziwości
$\varphi$ w \emph{zbiorowej strukturze} $V_\kappa$
jest definiowalna zewnętrznie.
Wyróżnianie w metateorii daje taki podzbiór $b$;
ponieważ $b\subseteq a$, jego ranga
jest nie większa od rangi $a$, a więc
$b\in V_\kappa$.

Sprawdzimy schemat zastępowania.
Niech $a,\vec p\in V_\kappa$ i załóżmy, że
$(V_\kappa,\in)$ spełnia
$\forall x\in a\ \exists!y\,\varphi(x,y,\vec p)$.
Zastępowanie w otaczającym uniwersum,
z wykorzystaniem predykatu spełniania
struktury zbiorowej, daje zbiór
$B$ wszystkich odpowiednich wartości $y\in V_\kappa$.
Z Lematu~\ref{lem:lc37-small-ranks} mamy
$|a|<\kappa$, a zatem $|B|<\kappa$.
Dla każdego $y\in B$ mamy
$\operatorname{rank}(y)<\kappa$.
Regularność daje
$\sup_{y\in B}(\operatorname{rank}(y)+1)<\kappa$;
stąd $B\in V_\kappa$.

Wreszcie dla dowolnego $a\in V_\kappa$
wybieramy w otaczającym uniwersum
relację dobrze porządkującą $a$,
korzystając z AC. Relacja ta jest
podzbiorem $a\times a$, więc jej ranga
jest ograniczona przez skończoną liczbę
następców rangi $a$. Należy zatem do
$V_\kappa$, co dowodzi AC w $V_\kappa$.
\end{proof}

\begin{wniosek}[Względna niesprzeczność z nieosiągalności]
\label{cor:lc37-inaccessible-con}
ZFC dowodzi implikacji arytmetycznej
\[
 \bigl(\exists\kappa\
    \text{$\kappa$ jest silnie nieosiągalny}\bigr)
 \ \Longrightarrow\ \operatorname{Con}(\mathrm{ZFC}).
\]
Jeżeli ZFC jest niesprzeczna, to
ZFC nie dowodzi istnienia kardynału silnie
nieosiągalnego.
\end{wniosek}
\begin{proof}
Z Twierdzenia~\ref{thm:lc37-rank-model} otrzymujemy
model zbiorowy ZFC. Formalizacja zasady poprawności
rachunku dowodowego daje wtedy
$\operatorname{Con}(\mathrm{ZFC})$;
argument jest przeprowadzalny w samej ZFC.
Jeśli ZFC dowodziłaby istnienia takiego kardynału,
dowodziłaby własnej niesprzeczności.
Przeczy to drugiemu twierdzeniu Gödla
omówionemu w Rozdziale~\ref{ch:second-incompleteness},
przy założeniu niesprzeczności ZFC.
\end{proof}

\section{Ultrafiltry zupełne i mierzalność}

\begin{definicja}[Ultrafiltr i zupełność]
\label{def:lc37-complete-ultrafilter}
Niech $I$ będzie niepustym zbiorem.
\emph{Filtrem właściwym} na $I$ nazywamy
$\mathcal U\subseteq\mathcal P(I)$ spełniający:
$I\in\mathcal U$, $\varnothing\notin\mathcal U$;
z $A,B\in\mathcal U$ wynika
$A\cap B\in\mathcal U$; z $A\in\mathcal U$
i $A\subseteq B\subseteq I$ wynika $B\in\mathcal U$.
Filtr jest \emph{ultrafiltrem}, gdy
dla każdego $A\subseteq I$ dokładnie jeden
ze zbiorów $A$, $I\setminus A$ należy do
$\mathcal U$. Jest \emph{niegłówny}, gdy
nie zawiera żadnego singletonu.

Dla nieskończonego kardynału $\kappa$ mówimy,
że filtr jest \emph{$\kappa$-zupełny}, gdy
dla każdego $\mu<\kappa$ i rodziny
$(A_\xi)_{\xi<\mu}$ jego elementów
przecięcie $\bigcap_{\xi<\mu}A_\xi$ również
należy do filtru.
\end{definicja}

\begin{definicja}[Kardynał mierzalny]
\label{def:lc37-measurable}
Nieprzeliczalny kardynał $\kappa$ jest
\emph{mierzalny}, jeśli na zbiorze $\kappa$
istnieje niegłówny $\kappa$-zupełny
ultrafiltr $\mathcal U$.
Taki ultrafiltr będziemy również nazywać
\emph{miarą zero-jedynkową} na $\kappa$;
nie chodzi tutaj o miarę liczbową o dowolnych
wartościach z przedziału $[0,1]$.
\end{definicja}

\begin{lemat}[Małe zbiory nie należą do miary]
\label{lem:lc37-uniform}
Niech $\kappa>\omega$ będzie kardynałem,
a $\mathcal U$ niegłównym $\kappa$-zupełnym
ultrafiltrem na $\kappa$. Wówczas
żaden zbiór $A\subseteq\kappa$ o mocy
$|A|<\kappa$ nie należy do $\mathcal U$.
W szczególności każdy podzbiór $\kappa$,
którego dopełnienie ma moc mniejszą od
$\kappa$, należy do $\mathcal U$.
\end{lemat}
\begin{proof}
Dla każdego $\alpha\in A$ mamy
$\kappa\setminus\{\alpha\}\in\mathcal U$.
Zupełność pozwala przeciąć $|A|<\kappa$
tych zbiorów, więc
$\kappa\setminus A\in\mathcal U$.
Właściwy filtr nie może zawierać
jednocześnie $A$ i jego dopełnienia.
Drugi wniosek jest już udowodniony.
\end{proof}

\begin{twierdzenie}[Kardynał mierzalny jest nieosiągalny]
\label{thm:lc37-measurable-inaccessible}
Jeżeli nieprzeliczalny kardynał $\kappa$
jest mierzalny, to jest regularny i
silnie graniczny, a więc silnie nieosiągalny.
\end{twierdzenie}
\begin{proof}
Niech $\mathcal U$ będzie niegłównym
$\kappa$-zupełnym ultrafiltrem na $\kappa$.

\emph{Regularność.} Gdyby
$\operatorname{cf}(\kappa)=\mu<\kappa$,
to pewien rosnący, współkońcowy ciąg
$(\gamma_\xi)_{\xi<\mu}$ dzieliłby
$\kappa$ na $\mu$ przedziałów
o mocy mniejszej od $\kappa$.
Żaden z tych przedziałów nie należałby
do $\mathcal U$ na mocy
Lematu~\ref{lem:lc37-uniform}.
Ich dopełnienia należałyby więc do
$\mathcal U$, a przecięcie tych
$\mu<\kappa$ dopełnień byłoby puste.
To przeczy $\kappa$-zupełności
i właściwości filtru.

\emph{Silna graniczność.} Załóżmy,
że istnieje kardynał $\lambda<\kappa$
taki, że $2^\lambda\ge\kappa$.
W ZFC istnieje wtedy iniekcja
$f:\kappa\to\mathcal P(\lambda)$.
Dla każdego $\xi<\lambda$ ultrafiltr
wybiera dokładnie jeden z dwóch zbiorów
\[
 A_\xi=\{\alpha<\kappa:\xi\in f(\alpha)\},
 \qquad \kappa\setminus A_\xi.
\]
Niech $B_\xi$ oznacza wybrany element
$\mathcal U$. Ponieważ $\lambda<\kappa$,
zbiór $B=\bigcap_{\xi<\lambda}B_\xi$
należy do $\mathcal U$.
Każde dwie liczby $\alpha,\beta\in B$
dają jednak te same odpowiedzi
na pytania o przynależność
$\xi\in f(\alpha)$, dla wszystkich
$\xi<\lambda$. Zatem
$f(\alpha)=f(\beta)$, a iniektywność
$f$ pokazuje, że $|B|\le1$.
Jest to sprzeczne z
Lematem~\ref{lem:lc37-uniform}.
\end{proof}

\begin{przyklad}[Dlaczego nie wystarcza zwykły ultrafiltr?]
\label{ex:lc37-frechet}
Niegłówny ultrafiltr na $\omega$
można otrzymać, stosując lemat o
ultrafiltrze do filtru zbiorów
o skończonych dopełnieniach.
Jest on $\omega$-zupełny, ponieważ
przecięcia \emph{skończenie} wielu
elementów zawsze należą do filtru.
Nie jest jednak $\omega_1$-zupełny:
dla $n<\omega$ zbiory
$A_n=\{m\in\omega:m\ge n\}$
należą do takiego ultrafiltru, lecz
$\bigcap_{n<\omega}A_n=\varnothing$.
Wymóg nieprzeliczalności w definicji
mierzalności nie jest więc kosmetyczny.
\end{przyklad}

\section{Ultraprodukt uniwersum i twierdzenie Łosia}

\begin{definicja}[Klasa ultraproduktu uniwersum]
\label{def:lc37-ultrapower}
Niech $\kappa$ będzie mierzalny, a
$\mathcal U$ niegłównym $\kappa$-zupełnym
ultrafiltrem na $\kappa$. Rozważamy
klasę wszystkich funkcji zbiorowych
$f:\kappa\to V$. Określamy
\[
 f\equiv_{\mathcal U}g
 \iff\{\alpha<\kappa:f(\alpha)=g(\alpha)\}
           \in\mathcal U.
\]
Przez $[f]_{\mathcal U}$ oznaczamy
klasę równoważności (jako obiekt
klasowy), a relację przynależności
ustalamy przez
\[
 [f]_{\mathcal U}\in_{\mathcal U}[g]_{\mathcal U}
 \iff\{\alpha<\kappa:f(\alpha)\in g(\alpha)\}
           \in\mathcal U.
\]
Tak powstałą strukturę klasową
oznaczamy $\operatorname{Ult}(V,\mathcal U)$.
\end{definicja}

\begin{uwaga}[Klasy i ich reprezentanci]
\label{rem:lc37-class-ultrapower}
W ZFC nie traktujemy zbioru wszystkich
funkcji $\kappa\to V$ jako zbioru.
Zamiast ilorazu będącego zbiorem można
wybrać \emph{reprezentantów Scotta}:
spośród równoważnych funkcji wybieramy
te o najmniejszej randze i bierzemy
zbiór wszystkich takich reprezentantów.
Pozwala to kodować każdy element
ultraproduktu jednym zbiorem
i rozważać definiowalną klasę
tych kodów.
\end{uwaga}

\begin{twierdzenie}[Twierdzenie Łosia dla ultraproduktu $V$]
\label{thm:lc37-los}
Niech $\kappa$ będzie mierzalny i
$\mathcal U$ jak w
Definicji~\ref{def:lc37-ultrapower}.
Dla każdej ustalonej formuły
$\varphi(x_1,\ldots,x_n)$ języka
teorii mnogości oraz funkcji
$f_1,\ldots,f_n:\kappa\to V$ zachodzi
\[
 \operatorname{Ult}(V,\mathcal U)
 \models\varphi([f_1]_{\mathcal U},
                \ldots,[f_n]_{\mathcal U})
 \quad\Longleftrightarrow\quad
 \{\alpha<\kappa:
 \varphi(f_1(\alpha),\ldots,f_n(\alpha))\}
 \in\mathcal U.
\]
Prawa strona jest predykatem ZFC dla
każdej ustalonej formuły.
\end{twierdzenie}
\begin{proof}
Dla atomów $x=y$ i $x\in y$
równoważność jest definicją obu
relacji. Przy koniunkcji używamy
zamknięcia $\mathcal U$ na przecięcia,
a przy negacji własności ultrafiltru:
dla dowolnego podzbioru $\kappa$
należy doń dokładnie jeden spośród
tego zbioru i jego dopełnienia.

Rozważmy kwantyfikator egzystencjalny.
Jeżeli klasa $[g]_{\mathcal U}$
jest świadkiem formuły, to hipoteza
indukcyjna daje zbiór dużych indeksów,
na których $g(\alpha)$ jest świadkiem
dla argumentów $f_i(\alpha)$.
W drugą stronę niech
\[
 D=\{\alpha<\kappa:
  \exists y\ \psi(y,f_1(\alpha),\ldots,
                         f_n(\alpha))\}\in\mathcal U.
\]
Dla każdego $\alpha\in D$ wybieramy
jeden zbiór $g(\alpha)$ spełniający
formułę; dla $\alpha\notin D$
przyjmujemy $g(\alpha)=\varnothing$.
Wybór jest poprawny w ZFC: najpierw
wybieramy najmniejszą możliwą rangę
świadka, następnie jednego ze świadków
w odpowiednim zbiorze tej rangi,
stosując AC do rodziny indeksowanej
\emph{zbiorem} $D$.
Otrzymujemy funkcję zbiorową
$g:\kappa\to V$. Zbiór indeksów,
dla których $\psi(g(\alpha),\vec f(\alpha))$
jest prawdziwa, zawiera $D$,
a więc należy do $\mathcal U$.
Hipoteza indukcyjna daje świadka $[g]$
po lewej stronie. To zamyka indukcję
po budowie ustalonej formuły.
\end{proof}

\begin{lemat}[Dobre ufundowanie ultraproduktu]
\label{lem:lc37-wellfounded}
Przy założeniach
Twierdzenia~\ref{thm:lc37-los}
relacja $\in_{\mathcal U}$ jest
dobrze ufundowana w sensie
zewnętrznym: nie istnieje ciąg
$([f_m]_{\mathcal U})_{m<\omega}$
spełniający
$[f_{m+1}]_{\mathcal U}\in_{\mathcal U}
 [f_m]_{\mathcal U}$ dla wszystkich $m$.
\end{lemat}
\begin{proof}
Gdyby taki ciąg istniał,
dla każdego $m$ mielibyśmy
\[
 A_m=\{\alpha<\kappa:
      f_{m+1}(\alpha)\in f_m(\alpha)\}
      \in\mathcal U.
\]
Ponieważ $\omega<\kappa$, z zupełności
$\mathcal U$ wynikałoby
$\bigcap_{m<\omega}A_m\in\mathcal U$.
Wybierając $\alpha$ z tego
niepustego przecięcia,
dostalibyśmy nieskończony ciąg
$f_0(\alpha)\ni f_1(\alpha)\ni
f_2(\alpha)\ni\cdots$, sprzeczny
z aksjomatem ufundowania w V.
\end{proof}

\begin{twierdzenie}[Zanurzenie elementarne z miary]
\label{thm:lc37-embedding}
Jeżeli $\kappa$ jest mierzalny,
to istnieje definiowalny z parametru
$\mathcal U$ tranzytywny model klasowy
$M\models\mathrm{ZFC}$ oraz
zanurzenie elementarne $j:V\to M$
takie, że
\[
 j(\alpha)=\alpha\quad(\alpha<\kappa),
 \qquad j(\kappa)>\kappa.
\]
W szczególności najmniejsza liczba
porządkowa przesuwana przez $j$,
zwana \emph{punktem krytycznym},
jest równa $\kappa$.
\end{twierdzenie}
\begin{proof}
Niech $\mathcal U$ będzie miarą na
$\kappa$. Dzięki Twierdzeniu~\ref{thm:lc37-los}
ultraprodukt jest elementarnym obrazem
uniwersum przy odwzorowaniu
$x\mapsto[c_x]_{\mathcal U}$, gdzie
$c_x(\alpha)=x$.
Z Lematu~\ref{lem:lc37-wellfounded}
relacja członkostwa jest dobrze ufundowana
dla ciągów długości $\omega$.
Jest również ekstensjonalna
z elementarności, a ponadto
\emph{zbiorolokalna}: dla ustalonej
$f$ członkowie $[f]$ mają reprezentantów
$g$ takich, że
$g(\alpha)\in f(\alpha)$ na odpowiednim
zbiorze indeksów. Można więc ograniczyć
zakres wszystkich takich $g$ do zbioru
$\bigl(\bigcup_{\alpha<\kappa}f(\alpha)\bigr)
\cup\{\varnothing\}$,
zmieniając wartości poza zbiorem
miary jeden. Zbiór odpowiednich funkcji
jest zbiorem. Istnieje zatem
kolaps Mostowskiego tej dobrze
ufundowanej, ekstensjonalnej,
zbiorolokalnej klasy definiowalnej
do tranzytywnej klasy $M$.
Złożenie odwzorowania stałych funkcji
z kolapsem oznaczamy przez $j$.
Twierdzenie Łosia daje elementarność,
a więc $M\models\mathrm{ZFC}$.

Pokażemy indukcyjnie, że
$j(\alpha)=\alpha$ dla $\alpha<\kappa$.
Dla $\alpha=0$ teza jest oczywista.
Jeśli $j(\beta)=\beta$ dla $\beta<\alpha$,
to dowolny element $j(\alpha)$
jest reprezentowany funkcją
$f:\kappa\to\alpha$ na zbiorze
miary jeden. Podzielenie tego zbioru
na mniej niż $\kappa$ włókien
$f^{-1}(\{\beta\})$, $\beta<\alpha$,
i $\kappa$-zupełność ultrafiltru
dają jedno włókno należące do
$\mathcal U$. Zatem $[f]=[c_\beta]$
dla pewnej $\beta<\alpha$,
co po kolapsie oznacza, że każdy
element $j(\alpha)$ jest mniejszy
od $\alpha$. Ponieważ
$j(\beta)\in j(\alpha)$ dla
wszystkich $\beta<\alpha$,
otrzymujemy $j(\alpha)=\alpha$.

Niech $i:\kappa\to\kappa$
będzie funkcją identycznościową.
Dla każdego $\alpha<\kappa$
zbiór $\{\xi:\alpha<\xi\}$ ma
dopełnienie mocy mniejszej od
$\kappa$, więc należy do $\mathcal U$
na mocy Lematu~\ref{lem:lc37-uniform}.
Stąd w uporządkowaniu ultraproduktu
\[
 [c_\alpha]_{\mathcal U}
 <[i]_{\mathcal U}
 <[c_\kappa]_{\mathcal U}.
\]
Po kolapsie dostajemy
$\alpha<[i]$ dla wszystkich
$\alpha<\kappa$ oraz
$[i]<j(\kappa)$. Zatem
$\kappa\le[i]<j(\kappa)$,
czyli $j(\kappa)>\kappa$.
\end{proof}

\section{Miary normalne i diagonalne przecięcia}

\begin{definicja}[Normalny ultrafiltr]
\label{def:lc37-normal}
Niech $\kappa$ będzie nieprzeliczalnym
kardynałem regularnym, a $\mathcal U$
niegłównym $\kappa$-zupełnym
ultrafiltrem na $\kappa$.
Mówimy, że $\mathcal U$ jest
\emph{normalny}, gdy dla każdego
$A\in\mathcal U$ i funkcji
$f:A\to\kappa$ spełniającej
$f(\alpha)<\alpha$ dla
wszystkich $0<\alpha\in A$,
istnieją $B\subseteq A$,
$B\in\mathcal U$ oraz $\xi<\kappa$
takie, że $f|_B$ jest stale
równa $\xi$.
Warunek w punkcie $0$ nie ma
znaczenia, ponieważ singleton
$\{0\}$ nie należy do $\mathcal U$.
\end{definicja}

\begin{definicja}[Przecięcie diagonalne]
\label{def:lc37-diagonal}
Dla ciągu $(A_\xi)_{\xi<\kappa}$
podzbiorów $\kappa$ definiujemy
\emph{przecięcie diagonalne}
\[
 \mathop{\triangle}_{\xi<\kappa}A_\xi
 =\{\alpha<\kappa:
    \forall\xi<\alpha\ (\alpha\in A_\xi)\}.
\]
\end{definicja}

\begin{twierdzenie}[Charakterystyka normalności]
\label{thm:lc37-normal-diagonal}
Niech $\kappa$ będzie nieprzeliczalnym
regularnym kardynałem, a $\mathcal U$
niegłównym $\kappa$-zupełnym
ultrafiltrem na $\kappa$.
Wówczas następujące warunki są równoważne:
\begin{enumerate}
\item $\mathcal U$ jest normalny;
\item dla każdej rodziny
$(A_\xi)_{\xi<\kappa}\subseteq\mathcal U$
jej przecięcie diagonalne należy
do $\mathcal U$.
\end{enumerate}
\end{twierdzenie}
\begin{proof}
Przyjmijmy normalność i niech
$D=\mathop{\triangle}_{\xi<\kappa}A_\xi$.
Jeżeli $D\notin\mathcal U$, to
$C=\kappa\setminus D\in\mathcal U$.
Dla $\alpha\in C\setminus\{0\}$
niech $f(\alpha)$ będzie
najmniejszym $\xi<\alpha$
takim, że $\alpha\notin A_\xi$.
Funkcja jest regresywna.
Z normalności dostajemy
$B\in\mathcal U$, $B\subseteq C$,
na którym $f$ jest stale równa
pewnej $\xi<\kappa$.
Wtedy $B\cap A_\xi=\varnothing$,
co jest sprzeczne z tym, że
oba zbiory należą do ultrafiltru.

Odwrotnie, załóżmy własność
przecięć diagonalnych.
Niech $A\in\mathcal U$,
a $f:A\setminus\{0\}\to\kappa$
będzie regresywna.
Przypuśćmy, że żadne włókno
$f^{-1}(\{\xi\})$ nie należy
do $\mathcal U$.
Dla każdej $\xi<\kappa$
zbiór
$A_\xi=\kappa\setminus f^{-1}(\{\xi\})$
należy do $\mathcal U$.
Wtedy $A\cap
 \mathop{\triangle}_{\xi<\kappa}A_\xi$
należy do $\mathcal U$,
a więc jest niepusty.
Dla $\alpha$ z tego przecięcia,
$\alpha>0$, mamy $f(\alpha)<\alpha$
oraz $\alpha\notin
 f^{-1}(\{f(\alpha)\})$, sprzeczność.
Jedno z włókien należy więc
do $\mathcal U$, co daje normalność.
\end{proof}

\begin{twierdzenie}[Kardynał mierzalny ma miarę normalną]
\label{thm:lc37-normal-exists}
Dla każdego kardynału mierzalnego
$\kappa$ istnieje niegłówny
$\kappa$-zupełny normalny ultrafiltr
na $\kappa$.
\end{twierdzenie}
\begin{proof}
Niech $\mathcal U$ będzie miarą,
a $j:V\to M$ zanurzeniem z
Twierdzenia~\ref{thm:lc37-embedding}.
Definiujemy
\[
 \mathcal W=\{X\subseteq\kappa:
                  \kappa\in j(X)\}.
\]
Jest to zbiór, a nie klasa
właściwa: funkcja $j$ jest
definiowalna z ustalonego parametru
$\mathcal U$, więc działa
wyodrębnianie na $\mathcal P(\kappa)$.

Z elementarności $j$ mamy
$j(\kappa\setminus X)
 =j(\kappa)\setminus j(X)$
i $j(X\cap Y)=j(X)\cap j(Y)$.
Ponieważ $\kappa<j(\kappa)$,
dokładnie jeden spośród $X$,
$\kappa\setminus X$ należy do
$\mathcal W$, a przecięcia i
nadzbiory zachowują przynależność.
Dla $\alpha<\kappa$ mamy
$j(\{\alpha\})=\{\alpha\}$,
zatem $\{\alpha\}\notin\mathcal W$.

Jeżeli $\mu<\kappa$ i
$(X_\xi)_{\xi<\mu}$ jest rodziną
elementów $\mathcal W$, to
$j(\mu)=\mu$ i elementarność
daje
\[
 j\left(\bigcap_{\xi<\mu}X_\xi\right)
  =\bigcap_{\xi<\mu}j(X_\xi).
\]
Wszystkie prawe zbiory zawierają
$\kappa$, więc przecięcie
należy do $\mathcal W$.
To dowodzi $\kappa$-zupełności.

Na koniec niech $A\in\mathcal W$,
a $f:A\to\kappa$ będzie regresywna
poza punktem $0$.
Ponieważ $\kappa\in j(A)$ i
$\kappa>0$, elementarność daje
$j(f)(\kappa)<\kappa$.
Jest to więc pewna
$\xi<\kappa$, a $j(\xi)=\xi$.
Zatem
\[
 \kappa\in j(\{\alpha\in A:
                          f(\alpha)=\xi\}),
\]
co znaczy, że włókno
$f^{-1}(\{\xi\})$ należy do
$\mathcal W$. Ultrafiltr
$\mathcal W$ jest normalny.
\end{proof}

\section{Hierarchia założeń i granice otrzymanych wyników}

\begin{twierdzenie}[Wynikające implikacje]
\label{thm:lc37-hierarchy}
W ZFC prawdziwe są następujące twierdzenia:
\begin{enumerate}
\item jeśli istnieje mierzalny kardynał,
to istnieje silnie nieosiągalny kardynał;
\item jeśli istnieje silnie nieosiągalny
kardynał, to istnieje tranzytywny
model zbiorowy ZFC;
\item jeśli istnieje mierzalny kardynał,
to istnieje nietrywialne zanurzenie
elementarne $j:V\to M$, dla pewnego
tranzytywnego, definiowalnego
modelu klasowego $M$ ZFC.
\end{enumerate}
Żadnej z odwrotnych implikacji nie
stwierdzamy w tym rozdziale.
\end{twierdzenie}
\begin{proof}
Pierwsza część jest
Twierdzeniem~\ref{thm:lc37-measurable-inaccessible},
druga wynika z
Twierdzenia~\ref{thm:lc37-rank-model},
a trzecia z
Twierdzenia~\ref{thm:lc37-embedding}.
\end{proof}

\begin{uwaga}[Co jest nowym aksjomatem?]
\label{rem:lc37-new-axioms}
ZFC nie zapewnia istnienia
silnie nieosiągalnego ani mierzalnego
kardynału. W odróżnieniu od CH i AC,
których niezależność omawialiśmy
wcześniej, tutaj rozważamy
\emph{wzmocnienia} teorii aksjomatami
istnienia dużych kardynałów.
Drugie twierdzenie Gödla wyklucza
dowód istnienia nieosiągalnego w
niesprzecznej ZFC, lecz samo
to rozumowanie nie daje dowodu
względnej niesprzeczności
takiego dodatkowego aksjomatu.
\end{uwaga}

\begin{uwaga}[Dalsze zagadnienia]
\label{rem:lc37-next}
Słaba zwartość kardynałów,
liczby kardynalne Mahla,
silna zwartość, superzwartość
oraz rozbudowane twierdzenia o
zanurzeniach elementarnych
wymagają odrębnych definicji
i dowodów. W tym rozdziale
nie zakładamy ich własności.
Nie utożsamiamy także
ultrafiltrów mierzalnych z
miarami Lebesgue'a.
\end{uwaga}

\begin{uwaga}[Zakres kontroli formalnej]
\label{rem:lc37-formal}
Sprawdzenie skończonych przykładów
ultrafiltrów i kompilacja LaTeX
nie są formalizacją twierdzeń
o dużych kardynałach w Lean.
Taką weryfikację należy
raportować osobno na podstawie
rzeczywistego przebiegu kompilatora.
\end{uwaga}

\chapter{Liczby kardynalne Mahla i słaba zwartość}
\label{ch:mahlos-weak-compactness}
\chaptermark{Liczby kardynalne Mahla i słaba zwartość}

\section{Plan i założenia}
W Rozdziale~\ref{ch:large-cardinals} wykazaliśmy,
że mierzalność implikuje silną nieosiągalność,
oraz skonstruowaliśmy zanurzenie elementarne
$j:V\to M$ o punkcie krytycznym $\kappa$.
Tutaj wprowadzimy zbiory klubowe i stacjonarne,
liczby kardynalne Mahla, własność drzewa, słabą zwartość
oraz relacje podziałowe. Udowodnimy w ZFC, że
\[
 \kappa\ \text{jest mierzalną liczbą kardynalną}\quad\Longrightarrow\quad
 \begin{cases}
 \kappa\ \text{jest liczbą kardynalną Mahla},\\
 \kappa\ \text{jest słabo zwartą liczbą kardynalną},\\
 \kappa\longrightarrow(\kappa)^2_2.
 \end{cases}
\]
Nie postulujemy istnienia mierzalnych liczb kardynalnych.
Jeżeli nie wskazano inaczej, $\kappa$ oznacza
nieprzeliczalną regularną liczbę kardynalną, a liczby kardynalne
utożsamiamy z początkowymi liczbami porządkowymi.

\section{Zbiory klubowe i stacjonarne}

\begin{definicja}[Zbiór klubowy]\label{def:lc38-club}
Niech $\kappa$ będzie nieprzeliczalną regularną
liczbą kardynalną. Zbiór $C\subseteq\kappa$
jest \emph{nieograniczony}, jeżeli
$\forall\alpha<\kappa\ \exists\beta\in C$
takie, że $\alpha<\beta$.
Jest \emph{domknięty}, jeżeli dla każdej
niezerowej granicznej liczby porządkowej
$\delta<\kappa$ z nieograniczoności
$C\cap\delta$ w $\delta$ wynika $\delta\in C$.
Zbiór spełniający oba warunki jest
\emph{klubowy} (ang.\ \emph{club}).
\end{definicja}

\begin{definicja}[Zbiór stacjonarny]\label{def:lc38-stationary}
Dla nieprzeliczalnego regularnego $\kappa$
zbiór $S\subseteq\kappa$ jest
\emph{stacjonarny}, gdy
$S\cap C\ne\varnothing$ dla każdego
zbioru klubowego $C\subseteq\kappa$.
W przeciwnym razie jest
\emph{niestacjonarny}.
\end{definicja}

\begin{lemat}[Przecięcie mniej niż $\kappa$ zbiorów klubowych]
\label{lem:lc38-club-intersections}
Niech $\mu<\kappa$ i $(C_\xi)_{\xi<\mu}$
będzie rodziną zbiorów klubowych w
nieprzeliczalną regularną liczbą kardynalną $\kappa$.
Wtedy $\bigcap_{\xi<\mu}C_\xi$ jest klubowy.
\end{lemat}
\begin{proof}
Przecięcie jest domknięte, ponieważ
każdy jego punkt skupienia jest
punktem skupienia wszystkich $C_\xi$.
Dla $\mu=0$ teza jest oczywista.
Załóżmy $\mu>0$.
Ustalmy $\beta<\kappa$ oraz
$\beta_0>\beta$. Dla $\xi<\mu$
i $\gamma<\kappa$ określamy
$d_\xi(\gamma)=
 \min\{\delta\in C_\xi:\delta\ge\gamma\}$,
co jest możliwe dzięki
nieograniczoności $C_\xi$.
Rekurencyjnie przyjmujemy
\[
 \beta_{n+1}=
 \sup\bigl(\{\beta_n+1\}\cup
       \{d_\xi(\beta_n)+1:\xi<\mu\}\bigr).
\]
Regularność $\kappa$ zapewnia
$\beta_n<\kappa$ dla każdego $n<\omega$
oraz $\delta=\sup_{n<\omega}\beta_n<\kappa$.
Dla każdego $\xi<\mu$ mamy
$\beta_n\le d_\xi(\beta_n)<\beta_{n+1}$,
więc ciąg punktów $C_\xi$ o tych
wartościach jest współkońcowy w $\delta$.
Domkniętość $C_\xi$ daje
$\delta\in C_\xi$.
Zatem wspólne przecięcie
zawiera punkt $\delta>\beta$.
\end{proof}

\begin{definicja}[Przecięcie diagonalne]
\label{def:lc38-diagonal}
Dla rodziny $(C_\xi)_{\xi<\kappa}$
zbiorów w $\kappa$ określamy
\[
 \mathop{\triangle}_{\xi<\kappa}C_\xi
 =\{\alpha<\kappa:
       \forall\xi<\alpha\ (\alpha\in C_\xi)\}.
\]
\end{definicja}

\begin{lemat}[Diagonalne przecięcie zbiorów klubowych]
\label{lem:lc38-diagonal-club}
Jeżeli $(C_\xi)_{\xi<\kappa}$
jest rodziną zbiorów klubowych
w nieprzeliczalnej regularnej liczbie kardynalnej $\kappa$,
to $D=\mathop{\triangle}_{\xi<\kappa}C_\xi$
jest klubowy.
\end{lemat}
\begin{proof}
Niech graniczna $\delta<\kappa$
będzie punktem skupienia $D$.
Dla ustalonego $\xi<\delta$
wszystkie dostatecznie duże
$\alpha\in D\cap\delta$ należą do
$C_\xi$. Domkniętość $C_\xi$
daje $\delta\in C_\xi$.
Ponieważ $\xi$ było dowolne,
$\delta\in D$, co dowodzi domkniętości.

Ustalmy $\beta<\kappa$.
Konstruujemy rosnący ciąg
$\gamma_0>\beta$ oraz
\[
 \gamma_{n+1}\in
 \bigcap_{\xi\le\gamma_n}C_\xi,\qquad
 \gamma_{n+1}>\gamma_n.
\]
Przecięcie po prawej jest klubowe
na mocy Lematu~\ref{lem:lc38-club-intersections},
bo $\gamma_n+1<\kappa$.
Niech $\delta=\sup_n\gamma_n<\kappa$.
Dla każdego $\xi<\delta$
wszystkie dostatecznie późne
$\gamma_{n+1}$ należą do $C_\xi$;
ich supremum $\delta$ także należy
do $C_\xi$. Zatem $\delta\in D$
i $\delta>\beta$.
\end{proof}

\begin{definicja}[Funkcja regresywna]
\label{def:lc38-regressive}
Niech $S\subseteq\kappa$.
Funkcja $f:S\to\kappa$ jest
\emph{regresywna}, jeżeli
$f(\alpha)<\alpha$
dla wszystkich $\alpha\in S\setminus\{0\}$.
\end{definicja}

\begin{twierdzenie}[Lemat Fodora]\label{thm:lc38-fodor}
Niech $\kappa$ będzie nieprzeliczalną
regularną liczbą kardynalną, $S\subseteq\kappa$
zbiorem stacjonarnym, a $f:S\to\kappa$
funkcją regresywną.
Wtedy istnieje $\xi<\kappa$
takie, że $f^{-1}(\{\xi\})$
jest stacjonarny.
\end{twierdzenie}
\begin{proof}
Przecięcie zbioru stacjonarnego
z ogonem ponad $0$ pozostaje
stacjonarne, więc pomijamy argument $0$.
Przypuśćmy, że każde włókno
$S_\xi=f^{-1}(\{\xi\})$
jest niestacjonarne.
Dla każdego $\xi<\kappa$
wybieramy klubowy $C_\xi$
rozłączny z $S_\xi$.
ZFC pozwala wybrać je jednocześnie,
ponieważ indeksy tworzą zbiór.
Na mocy Lematu~\ref{lem:lc38-diagonal-club}
zbiór $D=\mathop{\triangle}_{\xi<\kappa}C_\xi$
jest klubowy.
Wybierzmy $\alpha\in S\cap D$,
$\alpha>0$, i połóżmy
$\xi=f(\alpha)<\alpha$.
Z definicji $D$ otrzymujemy
$\alpha\in C_\xi$, podczas gdy
$\alpha\in S_\xi$ wymaga
$\alpha\notin C_\xi$.
Jest to sprzeczność.
\end{proof}

\begin{przyklad}[Klub i zbiór niestacjonarny]
\label{ex:lc38-omega-one}
Dla $\kappa=\omega_1$ zbiór
niezerowych granicznych liczb
porządkowych poniżej $\omega_1$
jest klubowy: nad każdą
$\alpha<\omega_1$ leży
$\alpha+\omega<\omega_1$, a
granica granicznych liczb porządkowych
jest graniczna. Zbiór następników
$\{\alpha+1:\alpha<\omega_1\}$
jest z nim rozłączny,
a więc niestacjonarny.
W szczególności samo bycie
nieograniczonym nie oznacza
stacjonarności.
\end{przyklad}

\begin{uwaga}[Stacjonarność a normalność]
\label{rem:lc38-fodor-normal}
Lemat Fodora przypomina własność
normalnego ultrafiltru z
Rozdziału~\ref{ch:large-cardinals}.
Zbiory stacjonarne nie tworzą
jednak filtru: dwa takie zbiory
mogą być rozłączne.
Nie utożsamiamy stacjonarności
z przynależnością do jednej miary.
\end{uwaga}

\section{Liczby kardynalne Mahla}

\begin{definicja}[Liczba kardynalna Mahla]\label{def:lc38-mahlo}
Nieprzeliczalną liczbę kardynalną $\kappa$
nazywamy \emph{liczbą kardynalną Mahla},
gdy jest silnie nieosiągalna i zbiór
\[
 I_\kappa=\{\lambda<\kappa:
      \lambda\text{ jest silnie nieosiągalną liczbą kardynalną}\}
\]
jest stacjonarny w $\kappa$.
Nie wystarcza sama nieograniczoność
zbioru $I_\kappa$.
\end{definicja}

\begin{twierdzenie}[Liczba kardynalna Mahla jest granicą nieosiągalnych liczb kardynalnych]
\label{thm:lc38-mahlo-limit}
Jeżeli $\kappa$ jest liczbą kardynalną Mahla,
to dla każdego $\alpha<\kappa$
istnieje silnie nieosiągalna liczba kardynalna
$\lambda$ taki, że
$\alpha<\lambda<\kappa$.
Ponadto zbiór takich $\lambda$
przecina każdy zbiór klubowy w $\kappa$.
\end{twierdzenie}
\begin{proof}
Ogon $C_\alpha=\{\beta<\kappa:
\beta>\alpha\}$ jest klubowy:
jest nieograniczony i domknięty.
Z definicji stacjonarności
$I_\kappa\cap C_\alpha\ne\varnothing$.
Każdy element tego przecięcia
spełnia pierwszą część tezy.
Druga część jest dokładnie
warunkiem stacjonarności $I_\kappa$.
\end{proof}

\begin{uwaga}[Granica a stacjonarność]
\label{rem:lc38-mahlo-warning}
Własność „$\kappa$ jest granicą
nieosiągalnych liczb kardynalnych”
jest słabszym warunkiem niż
własność Mahla. Z faktu
nieograniczoności $I_\kappa$
nie wolno bez dowodu wyprowadzać
przecinania każdego klubu.
\end{uwaga}

\section{Własność drzewa i słaba zwartość}

\begin{definicja}[Drzewo, poziomy i gałąź]
\label{def:lc38-tree}
\emph{Drzewo} to zbiór częściowo
uporządkowany $(T,<_{T})$,
w którym ścisłe poprzedniki
każdego węzła $t$ są dobrze
uporządkowane przez $<_{T}$.
Ich typ porządkowy oznaczamy
$\operatorname{ht}_T(t)$.
Poziom $\alpha$ to
$T_\alpha=\{t\in T:
                 \operatorname{ht}_T(t)=\alpha\}$,
a wysokość drzewa to
$\sup\{\operatorname{ht}_T(t)+1:t\in T\}$.

Dla nieskończonej liczby kardynalnej
$\kappa$ \emph{drzewo $\kappa$-owe}
ma wysokość $\kappa$ i
$|T_\alpha|<\kappa$ dla
wszystkich $\alpha<\kappa$.
\emph{Gałąź współkońcowa} jest
łańcuchem przecinającym
każdy poziom $T_\alpha$.
\end{definicja}

\begin{definicja}[Słaba zwartość]
\label{def:lc38-weak-compact}
Nieprzeliczalna regularna
liczba kardynalna $\kappa$ ma
\emph{własność drzewa}, gdy każde
drzewo $\kappa$-owe ma
gałąź współkońcową.
Jest \emph{słabo zwarta},
gdy jest silnie nieosiągalna
i spełnia własność drzewa.
\end{definicja}

\begin{przyklad}[Drzewo binarne]
\label{ex:lc38-binary}
Niech $\kappa$ będzie silnie
nieosiągalny i
$T=2^{<\kappa}
 =\bigcup_{\alpha<\kappa}2^\alpha$,
uporządkowany przez właściwe
przedłużanie funkcji.
Każdy poziom ma moc
$2^{|\alpha|}<\kappa$,
więc jest to drzewo $\kappa$-owe.
Funkcje stale równe zeru
tworzą gałąź współkońcową.
Ten przykład nie rozstrzyga
własności drzewa: trzeba
uwzględniać wszystkie drzewa
$\kappa$-owe, nie tylko jedno.
\end{przyklad}

\begin{lemat}[Obrazy małych zbiorów]
\label{lem:lc38-small-image}
Niech $\kappa$ będzie mierzalną liczbą kardynalną,
a $j:V\to M$ zanurzeniem elementarnym
z Rozdziału~\ref{ch:large-cardinals}
o punkcie krytycznym $\kappa$.
Jeżeli $A$ jest zbiorem mocy
mniejszej od $\kappa$, to
$j(A)=\{j(a):a\in A\}$.
\end{lemat}
\begin{proof}
Dla $A=\varnothing$ jest to oczywiste.
W przeciwnym razie istnieje
$\mu<\kappa$ oraz surjekcja
$f:\mu\to A$. W $M$ funkcja
$j(f):j(\mu)\to j(A)$ jest
surjekcją. Ponieważ $\mu<\kappa$,
mamy $j(\mu)=\mu$.
Z elementarności
$j(f)(\xi)=j(f(\xi))$
dla każdego $\xi<\mu$.
Wszystkie elementy $j(A)$
są więc obrazami elementów $A$.
Odwrotna inkluzja jest
konsekwencją elementarności
relacji przynależności.
\end{proof}

\begin{twierdzenie}[Mierzalność implikuje własność drzewa]
\label{thm:lc38-measurable-tree}
Jeżeli liczba kardynalna $\kappa$ jest mierzalna,
to każde drzewo $\kappa$-owe
ma gałąź współkońcową.
W szczególności $\kappa$
jest słabo zwarty.
\end{twierdzenie}
\begin{proof}
Z Twierdzenia~\ref{thm:lc37-measurable-inaccessible}
wiemy, że $\kappa$ jest
silnie nieosiągalną liczbą kardynalną.
Niech $T$ będzie dowolnym
drzewem $\kappa$-owym i
$j:V\to M$ zanurzeniem
z punktem krytycznym $\kappa$
z Rozdziału~\ref{ch:large-cardinals}.
W $M$ drzewo $j(T)$ ma
wysokość $j(\kappa)>\kappa$.
Zatem na jego poziomie
$\kappa$ istnieje węzeł $t$.

Dla każdego $\alpha<\kappa$
węzeł $t$ ma dokładnie
jednego poprzednika $u_\alpha$
na poziomie $\alpha$.
Ponieważ $j(\alpha)=\alpha$,
elementarność daje
$(j(T))_\alpha=j(T_\alpha)$.
Z $|T_\alpha|<\kappa$
i Lematu~\ref{lem:lc38-small-image}
wynika
$j(T_\alpha)=\{j(s):s\in T_\alpha\}$.
Możemy więc jednoznacznie
wybrać $s_\alpha\in T_\alpha$
takie, że $j(s_\alpha)=u_\alpha$.

Dla $\alpha<\beta<\kappa$
mamy $u_\alpha<_{j(T)}u_\beta$,
a elementarność $j$ daje
$s_\alpha<_{T}s_\beta$.
Ponieważ $j$ jest definiowalny
z ustalonej miary, zastępowanie
w ZFC tworzy zbiór
$\{s_\alpha:\alpha<\kappa\}$.
Jest to gałąź przechodząca
przez wszystkie poziomy drzewa.
\end{proof}

\section{Miara normalna i stacjonarność nieosiągalnych}

\begin{lemat}[Kluby należą do miary normalnej]
\label{lem:lc38-clubs-in-measure}
Niech $\kappa$ będzie mierzalną liczbą kardynalną
i $\mathcal W$ niegłównym,
$\kappa$-zupełnym normalnym
ultrafiltrem na $\kappa$.
Każdy zbiór klubowy
$C\subseteq\kappa$ należy do
$\mathcal W$.
\end{lemat}
\begin{proof}
Gdyby $A=\kappa\setminus C$
należał do $\mathcal W$,
wybierzmy $\gamma_0\in C$
oraz $A'=A\cap\{\alpha:
\gamma_0<\alpha<\kappa\}$.
Z Lematu~\ref{lem:lc37-uniform}
ogon należy do $\mathcal W$,
więc $A'\in\mathcal W$.
Określmy
$f(\alpha)=\sup(C\cap\alpha)$,
$\alpha\in A'$.
Ponieważ $C$ jest domknięty,
a $\alpha\notin C$,
zachodzi $f(\alpha)<\alpha$.
Normalność daje
$B\in\mathcal W$, $B\subseteq A'$
i $\xi<\kappa$, dla których
$f|_B$ jest stale równa $\xi$.
Wybieramy $\gamma\in C$
z $\gamma>\xi$ i
$\alpha\in B$ z $\alpha>\gamma$:
taki $\alpha$ istnieje,
ponieważ $B$ ma moc $\kappa$
na mocy Lematu~\ref{lem:lc37-uniform}.
Wtedy $\gamma\in C\cap\alpha$,
czyli $f(\alpha)\ge\gamma>\xi$,
co jest sprzecznością.
\end{proof}

\begin{twierdzenie}[Mierzalna liczba kardynalna jest liczbą kardynalną Mahla]
\label{thm:lc38-measurable-mahlo}
Jeżeli liczba kardynalna $\kappa$ jest mierzalna,
to zbiór silnie nieosiągalnych
liczb kardynalnych mniejszych od $\kappa$
jest stacjonarny w $\kappa$.
\end{twierdzenie}
\begin{proof}
Niech $j:V\to M$ będzie
zanurzeniem elementarnym
o punkcie krytycznym $\kappa$
z Rozdziału~\ref{ch:large-cardinals}.
Przez
$I_\kappa=\{\alpha<\kappa:
\alpha\text{ jest silnie nieosiągalną liczbą kardynalną}\}$
oznaczamy odpowiedni zbiór.
W Rozdziale~\ref{ch:large-cardinals}
skonstruowano normalny ultrafiltr
$\mathcal W=\{X\subseteq\kappa:
                        \kappa\in j(X)\}$.

Liczba kardynalna $\kappa$ jest silnie
nieosiągalną liczbą kardynalną w $V$ na mocy
Twierdzenia~\ref{thm:lc37-measurable-inaccessible}.
Jest także silnie nieosiągalną liczbą kardynalną
w tranzytywnym $M$.
Istotnie, funkcja świadcząca
w $M$ o singularności $\kappa$
byłaby taką samą funkcją w $V$,
co przeczyłoby regularności.
Gdyby $M$ uważał, że $\kappa$
nie jest liczbą kardynalną, odpowiednia
bijekcja z mniejszą liczbą
porządkową także należałaby do $V$.
Jeśli $M$ uważałby, że
$2^\lambda\ge\kappa$ dla
pewnego $\lambda<\kappa$,
to wybór w $M$ dawałby iniekcję
$\kappa\to\mathcal P^M(\lambda)$.
W $V$ byłaby to iniekcja
$\kappa\to\mathcal P^V(\lambda)$,
sprzeczna z silną granicznością.

Z elementarności wynika, że
$j(I_\kappa)$ jest w $M$
zbiorem silnie nieosiągalnych
liczb kardynalnych poniżej $j(\kappa)$.
Ponieważ $\kappa<j(\kappa)$
i $M$ uważa $\kappa$ za
silnie nieosiągalną liczbę kardynalną,
mamy $\kappa\in j(I_\kappa)$.
Stąd $I_\kappa\in\mathcal W$.
Każdy klub także należy do
$\mathcal W$ na mocy
Lematu~\ref{lem:lc38-clubs-in-measure}.
Właściwy filtr nie zawiera
dwóch rozłącznych zbiorów.
Dlatego $I_\kappa$ przecina
każdy klub i jest stacjonarny.
\end{proof}

\section{Kolorowania par}

\begin{definicja}[Relacja podziałowa]
\label{def:lc38-partition}
Niech $\kappa,\lambda,\theta$
będą liczbami kardynalnymi, a $r\ge1$
liczbą naturalną.
Przez $[\kappa]^r$ rozumiemy
zbiór $r$-elementowych
podzbiorów $\kappa$.
Zapis
\[
 \kappa\longrightarrow(\lambda)^r_\theta
\]
oznacza, że dla każdego
$c:[\kappa]^r\to\theta$
istnieje $H\subseteq\kappa$
o mocy $\lambda$, na którym
$c$ jest stałe dla wszystkich
argumentów z $[H]^r$.
Zbiór $H$ nazywamy
\emph{jednorodnym}.
\end{definicja}

\begin{twierdzenie}[Mierzalność i jednorodność par]
\label{thm:lc38-measurable-partition}
Jeśli liczba kardynalna $\kappa$ jest mierzalna,
to $\kappa\longrightarrow(\kappa)^2_2$.
\end{twierdzenie}
\begin{proof}
Niech $c:[\kappa]^2\to\{0,1\}$
i niech $\mathcal U$ będzie
niegłównym $\kappa$-zupełnym
ultrafiltrem na $\kappa$.
Dla $\alpha<\kappa$ i $i<2$
definiujemy
\[
 B_{\alpha,i}
 =\{\beta<\kappa:\beta>\alpha,\
                   c(\{\alpha,\beta\})=i\}.
\]
Ogon ponad $\alpha$ należy
do $\mathcal U$, więc dokładnie
jeden z dwóch rozłącznych
zbiorów $B_{\alpha,0}$,
$B_{\alpha,1}$ należy do
$\mathcal U$. Oznaczmy
jego kolor przez $i_\alpha$.
Jeden ze zbiorów
$X_i=\{\alpha<\kappa:i_\alpha=i\}$
należy do $\mathcal U$.
Ustalmy taki $i$.

Rekurencyjnie zbudujemy rosnący
ciąg $(\alpha_\xi)_{\xi<\kappa}$
o wyrazach w $X_i$.
Po wybraniu $\alpha_\eta$
dla $\eta<\xi$, $\xi<\kappa$,
rozważamy przecięcie
\[
 X_i\cap\bigcap_{\eta<\xi}
                      B_{\alpha_\eta,i}.
\]
Jest ono w $\mathcal U$,
ponieważ $\xi<\kappa$.
Nie może być ograniczone w
$\kappa$, bo wszystkie zbiory
ograniczone mają moc mniejszą
od $\kappa$, a nie należą do
$\mathcal U$ na mocy
Lematu~\ref{lem:lc37-uniform}.
Wybieramy najmniejszy element
przecięcia większy od wszystkich
wcześniejszych $\alpha_\eta$.
Zbiór
$H=\{\alpha_\xi:\xi<\kappa\}$
ma moc $\kappa$ i dla
$\eta<\xi$ spełnia
$\alpha_\xi\in B_{\alpha_\eta,i}$.
Zatem każde dwie różne
liczby z $H$ otrzymują
ten sam kolor $i$.
\end{proof}

\begin{uwaga}[Właściwa hierarchia implikacji]
\label{rem:lc38-implications}
Wykazaliśmy bezpośrednio, że
mierzalność implikuje słabą
zwartość, własność Mahla i
powyższą relację podziałową.
Nie dowodziliśmy odwrotnych
implikacji ani pełnych
charakterystyk słabej zwartości
przez logikę infinitarną.
Tego rodzaju równoważności
wymagają osobnego dowodu,
a nie samego rysunku hierarchii.
\end{uwaga}

\begin{wniosek}[Dwa następstwa mierzalności]
\label{cor:lc38-hierarchy}
Jeśli istnieje mierzalna liczba kardynalna,
to istnieje liczba kardynalna Mahla
i słabo zwarta liczba kardynalna.
Dokładniej, każda mierzalna liczba kardynalna ma obie te własności.
\end{wniosek}
\begin{proof}
Pierwsza własność wynika
z Twierdzenia~\ref{thm:lc38-measurable-mahlo},
druga z
Twierdzenia~\ref{thm:lc38-measurable-tree}
oraz z nieosiągalności
mierzalnych liczb kardynalnych
w Rozdziale~\ref{ch:large-cardinals}.
\end{proof}

\begin{uwaga}[Nie mylić implikacji z niesprzecznością]
\label{rem:lc38-consistency}
Nieosiągalność implikuje
istnienie zbiorowego
modelu ZFC, jak pokazano
w Twierdzeniu~\ref{thm:lc37-rank-model}.
Zatem niesprzeczna ZFC nie
dowodzi istnienia liczb kardynalnych
Mahla, słabo zwartych ani
mierzalnych: przeczyłoby to
drugiemu twierdzeniu Gödla.
Sam ten argument nie dowodzi
względnej niesprzeczności
żadnego z tych aksjomatów.
\end{uwaga}

\begin{uwaga}[Zakres formalnej kontroli]
\label{rem:lc38-formal}
Kompilacja LaTeX i skończone
testy drzew i relacji podziałowych
nie są weryfikacją przez jądro Lean.
Status formalizacji wymaga
rzeczywistej kompilacji kodu Lean.
\end{uwaga}

\chapter{Słaba zwartość: charakterystyki kombinatoryczne i logiczne}
\label{ch:weak-compactness-characterizations}
\chaptermark{Słaba zwartość i logika infinitarna}

\section{Trzy postacie jednej własności}

W Rozdziale~\ref{ch:mahlos-weak-compactness} wprowadziliśmy
własność drzewa oraz relacje podziałowe. Wykazaliśmy tam,
że mierzalna liczba kardynalna jest słabo zwarta
i spełnia $\kappa\longrightarrow(\kappa)^2_2$.
W tym rozdziale udowodnimy odwrotną zależność
\emph{na poziomie silnie nieosiągalnych liczb kardynalnych}
oraz powiążemy obie własności ze zwartością
infinitarnej logiki zdaniowej.

W całym rozdziale pracujemy w ZFC.
Niech $\kappa$ będzie ustaloną nieprzeliczalną
silnie nieosiągalną liczbą kardynalną.
Oznacza to, że $\kappa$ jest regularna
i $2^\lambda<\kappa$ dla każdej liczby
kardynalnej $\lambda<\kappa$.
Nie przyjmujemy, że liczba kardynalna
$\kappa$ ma własność drzewa, chyba że
występuje ona w hipotezie konkretnego twierdzenia.

\begin{twierdzenie}[Charakterystyki słabej zwartości]
\label{thm:lc39-main-equivalence}
Niech $\kappa$ będzie nieprzeliczalną
silnie nieosiągalną liczbą kardynalną.
Następujące warunki są równoważne:
\begin{enumerate}
\item $\kappa$ ma własność drzewa:
każde drzewo $\kappa$-owe ma gałąź współkońcową;
\item $\kappa\longrightarrow(\kappa)^2_2$;
\item dla każdej niezerowej liczby kardynalnej
$\mu<\kappa$ zachodzi
$\kappa\longrightarrow(\kappa)^2_\mu$;
\item każda teoria infinitarnej logiki
zdaniowej o co najwyżej $\kappa$ formułach,
z których każda zależy od mniej niż
$\kappa$ zmiennych, mająca model dla
każdej podteorii mocy mniejszej niż
$\kappa$, ma model całej teorii.
\end{enumerate}
Pojęcia z punktu czwartego zostaną dokładnie
zdefiniowane w
Definicji~\ref{def:lc39-propositional-infinitary}
oraz Definicji~\ref{def:lc39-prop-compactness}.
Dowody wszystkich implikacji podajemy poniżej.
\end{twierdzenie}

\section{Oszacowania mocy i drzewa profili kolorowania}

\begin{lemat}[Małość potęg poniżej nieosiągalności]
\label{lem:lc39-small-exponent}
Niech $\kappa$ będzie silnie nieosiągalną
liczbą kardynalną. Dla każdej liczby
kardynalnej $0<\mu<\kappa$ i każdej
liczby porządkowej $\alpha<\kappa$
zachodzi $\mu^{|\alpha|}<\kappa$.
\end{lemat}
\begin{proof}
Jeżeli $\mu=1$ albo $\alpha=0$,
lewa strona jest równa $1$.
W przeciwnym razie istnieje iniekcja
$\mu\hookrightarrow\mathcal P(\mu)$,
więc funkcję $\alpha\to\mu$ można
kodować funkcją
$\alpha\to\mathcal P(\mu)$,
a następnie podzbiorem $\alpha\times\mu$.
Otrzymujemy
\[
 \mu^{|\alpha|}
 \le 2^{|\alpha|\cdot\mu}<\kappa.
\]
Druga nierówność wynika z silnej
graniczności $\kappa$, ponieważ
$|\alpha|\cdot\mu<\kappa$:
iloczyn dwóch nieskończonych liczb
kardynalnych jest ich maksimum,
a przypadki skończone są prostsze.
\end{proof}

\begin{definicja}[Drzewo profili]
\label{def:lc39-profile-tree}
Niech $0<\mu<\kappa$ będzie liczbą
kardynalną i niech
$c:[\kappa]^2\to\mu$ będzie kolorowaniem.
Dla każdej liczby porządkowej
$\alpha<\kappa$ definiujemy
$f_\alpha:\alpha\to\mu$ przez
$f_\alpha(\xi)=c(\{\xi,\alpha\})$
dla $\xi<\alpha$.
Niech
\[
 T_c=\{f_\alpha\restriction\delta:
                   \delta\le\alpha<\kappa\},
\]
uporządkowany przez właściwe
przedłużanie funkcji.
\emph{Poziomem profili} o długości
$\delta$ jest zbiór
$T_{c,\delta}=T_c\cap\mu^\delta$.
\end{definicja}

\begin{lemat}[Drzewo profili jest drzewem $\kappa$-owym]
\label{lem:lc39-profile-kappa-tree}
Niech $\kappa$ będzie silnie nieosiągalną
liczbą kardynalną, $0<\mu<\kappa$
i $c:[\kappa]^2\to\mu$.
Drzewo $T_c$ z
Definicji~\ref{def:lc39-profile-tree}
ma wysokość $\kappa$, a
$|T_{c,\delta}|<\kappa$
dla wszystkich $\delta<\kappa$.
\end{lemat}
\begin{proof}
Zbiór $T_c$ jest zamknięty na
ograniczanie funkcji: jeśli
$t=f_\alpha\restriction\delta$,
to dla każdego $\gamma<\delta$
mamy
$t\restriction\gamma
 =f_\alpha\restriction\gamma\in T_c$.
Wysokość węzła wynosi długość
jego dziedziny. Dla każdego
$\delta<\kappa$ istnieje
$f_\alpha\restriction\delta$
z $\alpha\ge\delta$, więc
poziom $\delta$ jest niepusty.
Wysokość drzewa to $\kappa$.

Ponadto
$T_{c,\delta}\subseteq\mu^\delta$,
a z Lematu~\ref{lem:lc39-small-exponent}
wynika
$|T_{c,\delta}|
 \le\mu^{|\delta|}<\kappa$.
\end{proof}

\begin{twierdzenie}[Własność drzewa daje wszystkie małe kolorowania par]
\label{thm:lc39-tree-to-partition}
Niech $\kappa$ będzie nieprzeliczalną
silnie nieosiągalną liczbą kardynalną
mającą własność drzewa.
Dla każdej niezerowej liczby
kardynalnej $\mu<\kappa$ zachodzi
\[
 \kappa\longrightarrow(\kappa)^2_\mu.
\]
\end{twierdzenie}
\begin{proof}
Ustalmy $0<\mu<\kappa$ i
$c:[\kappa]^2\to\mu$.
Z Lematu~\ref{lem:lc39-profile-kappa-tree}
drzewo profili $T_c$ jest
drzewem $\kappa$-owym.
Na mocy własności drzewa
ma gałąź współkońcową.
Jej suma jest funkcją
$b:\kappa\to\mu$;
dla każdego $\delta<\kappa$
jej ograniczenie
$b\restriction\delta$
należy do $T_{c,\delta}$.

Dla każdego $\delta<\kappa$
istnieje $\alpha\ge\delta$
taka, że
$f_\alpha\restriction\delta
 =b\restriction\delta$.
Wybieramy \emph{najmniejszą}
taką liczbę porządkową
i oznaczamy ją $a(\delta)$.
Zatem
\[
 a(\delta)\ge\delta,\qquad
 f_{a(\delta)}\restriction\delta
       =b\restriction\delta.
\]
Definiujemy $h:\kappa\to\mu$,
$h(\delta)=b(a(\delta))$.

Ponieważ $\kappa$ jest regularna,
pewien kolor $i<\mu$ ma
\emph{nieograniczony} przeciwobraz
$D=\{\delta<\kappa:h(\delta)=i\}$.
W istocie, gdyby każdy przeciwobraz
był ograniczony, supremum
mniej niż $\kappa$ ich ograniczeń
pozostałoby mniejszy od $\kappa$,
co przeczyłoby temu, że
ich suma zbiorów jest równa $\kappa$.

Rekurencyjnie, dla $\xi<\kappa$,
wybieramy $\delta_\xi\in D$
tak, by
\[
 \delta_\xi>
    \sup_{\eta<\xi}(a(\delta_\eta)+1).
\]
Istnieje odpowiedni punkt,
gdyż $\xi<\kappa$,
regularność zapewnia ograniczoność
wyrażenia po prawej, a $D$
jest nieograniczony.
Możemy wybierać najmniejszy
punkt spełniający wymaganie.
Niech $\alpha_\xi=a(\delta_\xi)$.
Z konstrukcji mamy
$\alpha_\eta<\delta_\xi\le\alpha_\xi$
dla każdej pary $\eta<\xi$,
więc $(\alpha_\xi)_{\xi<\kappa}$
jest ściśle rosnący.

Dla $\eta<\xi$:
\[
 \begin{aligned}
 c(\{\alpha_\eta,\alpha_\xi\})
 &=f_{\alpha_\xi}(\alpha_\eta)\\
 &=b(\alpha_\eta)
   =h(\delta_\eta)=i.
 \end{aligned}
\]
Pierwsza równość jest definicją
$f_{\alpha_\xi}$, druga wynika z
$\alpha_\eta<\delta_\xi$ i
$f_{\alpha_\xi}\restriction\delta_\xi
 =b\restriction\delta_\xi$,
a ostatnia z $\delta_\eta\in D$.
Zatem $H=\{\alpha_\xi:\xi<\kappa\}$
ma moc $\kappa$ i wszystkie
jego pary otrzymują kolor $i$.
\end{proof}

\section{Własność podziałowa wymusza gałąź}

\begin{definicja}[Porządek leksykograficzny drzewa]
\label{def:lc39-tree-lex}
Niech $(T,<_{T})$ będzie drzewem
o wysokości $\kappa$.
Dla każdego $\delta<\kappa$
wybieramy porządek liniowy
$\prec_\delta$ na poziomie $T_\delta$.

Każdemu węzłowi $t$ przyporządkowujemy
\emph{ścieżkę poprzedników}
$p_t$, której dziedziną jest
$\operatorname{ht}_{T}(t)+1$ i
$p_t(\delta)$ jest jedynym
węzłem na poziomie $\delta$
nie większym od $t$.
Na $T$ definiujemy $\prec$
jako porządek leksykograficzny
tych ścieżek: ścieżka będąca
właściwym początkiem drugiej
jest wcześniejsza; w przeciwnym
razie porównujemy węzły
w pierwszym poziomie $\delta$,
na którym ścieżki się różnią,
korzystając z $\prec_\delta$.
\end{definicja}

\begin{lemat}[Rozdzielenie stożków w porządku leksykograficznym]
\label{lem:lc39-lex-cones}
Niech $\prec$ będzie porządkiem
z Definicji~\ref{def:lc39-tree-lex}.
Jest to porządek liniowy
przedłużający porządek $<_{T}$.
Jeśli węzły $s,t\in T$ są
nieporównywalne w $T$ i
$s\prec t$, to dla wszystkich
$s'\ge_T s$ i $t'\ge_T t$
mamy $s'\prec t'$.
\end{lemat}
\begin{proof}
Dwie różne ścieżki poprzedników
są porównywalne leksykograficznie:
jeżeli żadna nie jest początkiem
drugiej, istnieje najmniejszy
poziom, na którym się różnią.
Porządki $\prec_\delta$ są liniowe,
więc opisane porównanie jest
jednoznaczne. Przechodniość
wynika z rozważenia pierwszego
poziomu różnicy trzech ścieżek;
daje to standardowy dowód
przechodniości porządku
leksykograficznego.
Jeżeli $s<_{T}t$, to
$p_s$ jest właściwym początkiem
$p_t$, więc $s\prec t$.

Załóżmy, że $s,t$ są
nieporównywalne.
Ich ścieżki różnią się
na pewnym najmniejszym
poziomie $\delta$.
Wszystkie przedłużenia $s'$
i $t'$ mają na poziomach
$\le\delta$ te same węzły,
co odpowiednio $s$ i $t$.
Zatem pierwsza różnica
między ścieżkami $s',t'$
występuje na tym samym
poziomie $\delta$ i ma
ten sam znak w
porządku $\prec_\delta$.
Stąd $s'\prec t'$.
\end{proof}

\begin{twierdzenie}[Relacja podziałowa implikuje własność drzewa]
\label{thm:lc39-partition-to-tree}
Niech $\kappa$ będzie nieprzeliczalną
regularną liczbą kardynalną
spełniającą
$\kappa\longrightarrow(\kappa)^2_2$.
Wtedy każde drzewo wysokości
$\kappa$, którego poziomy mają
moc mniejszą od $\kappa$,
ma gałąź współkońcową.
\end{twierdzenie}
\begin{proof}
Niech $T$ będzie dowolnym
takim drzewem. Ponieważ
ma wysokość $\kappa$,
każdy poziom $\delta<\kappa$
jest niepusty. Z jednej strony
$|T|\ge\kappa$, z drugiej
$|T|\le\kappa\cdot\kappa=\kappa$,
więc $|T|=\kappa$.
W ZFC wybieramy wyliczenie
$T=\{t_\alpha:\alpha<\kappa\}$
bez powtórzeń i porządek
leksykograficzny $\prec$
z Definicji~\ref{def:lc39-tree-lex}.

Definiujemy kolorowanie
$c:[\kappa]^2\to\{0,1\}$:
dla $\alpha<\beta$ kładziemy
\[
 c(\{\alpha,\beta\})=
 \begin{cases}
 0,&t_\alpha\prec t_\beta,\\
 1,&t_\beta\prec t_\alpha.
 \end{cases}
\]
Z założenia istnieje
$H\subseteq\kappa$ o mocy
$\kappa$, na którym $c$
jest stale równe $i<2$.
Dla $s\in T$ określamy
\[
 U_s=\{\alpha\in H:s\le_T t_\alpha\},
 \qquad
 A=\{s\in T:|U_s|=\kappa\}.
\]

Pokażemy, że dla każdego
$\delta<\kappa$ zbiór
$A\cap T_\delta$ jest niepusty.
Ponieważ $\kappa$ jest regularna,
zbiór wszystkich węzłów z
poziomów mniejszych od $\delta$
ma moc mniejszą od $\kappa$:
jest sumą zbiorów mniej niż
$\kappa$ poziomów, każdy
o mocy mniejszej od $\kappa$.
Zatem dla $\kappa$ wielu
indeksów $\alpha\in H$
węzeł $t_\alpha$ ma
wysokość co najmniej $\delta$.
Każdy taki węzeł ma jednego
poprzednika na poziomie $\delta$.
Poziom ten ma moc mniejszą
od $\kappa$; regularność
gwarantuje, że dla co najmniej
$\kappa$ wielu rozważanych
indeksów poprzednik jest
tym samym węzłem $s$.
Wtedy $s\in A\cap T_\delta$.

Pozostaje pokazać, że
węzły $A$ są parami
porównywalne w $T$.
Przypuśćmy przeciwnie,
że $s,t\in A$ są
nieporównywalne i,
po zamianie oznaczeń,
$s\prec t$.
Z Lematu~\ref{lem:lc39-lex-cones}
dla każdych
$\alpha\in U_s$ i
$\beta\in U_t$ mamy
$t_\alpha\prec t_\beta$.
Zbiory $U_s$ i $U_t$
są rozłączne: jeden węzeł
nie może przedłużać dwóch
nieporównywalnych węzłów.

Jeśli $i=0$, to jednorodność
$H$ oznacza, że
$\alpha<\beta$ zawsze, gdy
$\alpha\in U_s$ i
$\beta\in U_t$.
Jednak $|U_s|=\kappa$
oznacza, że $U_s$ jest
nieograniczony w $\kappa$;
wybierając dowolne
$\beta\in U_t$ i większe
$\alpha\in U_s$, otrzymujemy
sprzeczność.
Jeśli $i=1$, jednorodność
oznacza $\beta<\alpha$
dla wszystkich takich par.
Wybierając dowolne
$\alpha\in U_s$ i większe
$\beta\in U_t$ ponownie
otrzymujemy sprzeczność.
Zatem $A$ jest łańcuchem.
Przecina wszystkie poziomy
$T_\delta$, więc stanowi
gałąź współkońcową.
\end{proof}

\begin{wniosek}[Równoważność dwóch charakterystyk]
\label{cor:lc39-tree-partition}
Dla każdej nieprzeliczalnej
silnie nieosiągalnej liczby
kardynalnej $\kappa$
następujące warunki są
równoważne:
\[
 \kappa\text{ ma własność drzewa}
 \quad\Longleftrightarrow\quad
 \kappa\longrightarrow(\kappa)^2_2.
\]
\end{wniosek}
\begin{proof}
Implikacja z lewej na prawą
jest przypadkiem $\mu=2$
Twierdzenia~\ref{thm:lc39-tree-to-partition}.
Implikacja odwrotna wynika
z Twierdzenia~\ref{thm:lc39-partition-to-tree}.
\end{proof}

\begin{uwaga}[Co zawiera hipoteza nieosiągalności?]
\label{rem:lc39-scope}
W tym rozdziale dowodzimy
równoważności \emph{przy założeniu}
silnej nieosiągalności.
W pełnej charakterystyce
nieprzeliczalnych słabo zwartych
liczb kardynalnych sam warunek
$\kappa\to(\kappa)^2_2$
implikuje również silną
nieosiągalność. Dowód tego
dodatkowego faktu wymaga
osobnego argumentu o porządku
leksykograficznym na $2^\lambda$;
nie używamy go tutaj bez dowodu.
\end{uwaga}

\section{Zwartość infinitarnej logiki zdaniowej}

\begin{definicja}[Infinitarna logika zdaniowa o małych nośnikach]
\label{def:lc39-propositional-infinitary}
Niech $\kappa$ będzie nieprzeliczalną
regularną liczbą kardynalną.
Dla każdego $\xi<\kappa$
niech $p_\xi$ będzie odrębną
zmienną zdaniową.
Rozważamy formuły utworzone
z tych zmiennych przez
negację oraz skończone
lub nieskończone koniunkcje
i alternatywy indeksowane
zbiorami o mocy mniejszej
niż $\kappa$.
Zbiór zmiennych występujących
w danej formule $\varphi$
nazywamy jej \emph{nośnikiem}
i oznaczamy $\operatorname{supp}(\varphi)$.
Dopuszczamy w naszej teorii
wyłącznie formuły o nośniku
mocy mniejszej niż $\kappa$.

\emph{Wartościowanie} to funkcja
$v:\kappa\to\{0,1\}$.
Relację $v\models\varphi$
określamy rekurencyjnie:
$v\models p_\xi$ wtedy i tylko
wtedy, gdy $v(\xi)=1$;
negacja zmienia wartość
logiczną, koniunkcja
wymaga prawdziwości
wszystkich składników,
a alternatywa co najmniej
jednego ze składników.
Z definicji
$v\models\varphi$ zależy
wyłącznie od ograniczenia
$v\restriction\operatorname{supp}(\varphi)$.
\end{definicja}

\begin{definicja}[Lokalna spełnialność i zwartość]
\label{def:lc39-prop-compactness}
Niech $\kappa$ będzie nieprzeliczalną
regularną liczbą kardynalną,
a $\Sigma$ zbiorem co najwyżej
$\kappa$ formuł z
Definicji~\ref{def:lc39-propositional-infinitary}.
Mówimy, że $\Sigma$ jest
\emph{$\kappa$-lokalnie spełnialna},
gdy dla każdego
$\Sigma_0\subseteq\Sigma$
o mocy $|\Sigma_0|<\kappa$
istnieje wartościowanie
$v:\kappa\to\{0,1\}$
spełniające wszystkie
formuły z $\Sigma_0$.

Mówimy, że zachodzi
\emph{zasada $\kappa$-zwartości
infinitarnej logiki zdaniowej},
gdy każda taka
$\kappa$-lokalnie spełnialna
$\Sigma$ ma wartościowanie
spełniające wszystkie
jej formuły.
\end{definicja}

\begin{lemat}[Drzewo częściowych wartościowań]
\label{lem:lc39-valuations-tree}
Niech $\kappa$ będzie nieprzeliczalną
silnie nieosiągalną liczbą
kardynalną, a
$\Sigma=\{\varphi_\xi:\xi<\kappa\}$
niech będzie
$\kappa$-lokalnie spełnialną
teorią z Definicji~\ref{def:lc39-prop-compactness}.
Dla $\alpha<\kappa$ oznaczmy
\[
 \begin{aligned}
 T_\alpha=\{v\in2^\alpha:\ 
 &v\text{ spełnia }\varphi_\xi\
   \text{ zawsze, gdy }\xi<\alpha\\
 &\text{oraz }
   \operatorname{supp}(\varphi_\xi)
                           \subseteq\alpha\}.
 \end{aligned}
\]
Przez spełnianie formuły
przez częściowe $v$
rozumiemy jej obliczenie
na wszystkich zmiennych
w jej nośniku.
Wtedy
$T=\bigcup_{\alpha<\kappa}T_\alpha$
uporządkowany przez
przedłużanie funkcji
jest drzewem $\kappa$-owym.
\end{lemat}
\begin{proof}
Dla ustalonego $\alpha$
formuł o indeksach $\xi<\alpha$,
których nośniki zawierają się
w $\alpha$, jest mniej niż $\kappa$.
Z lokalnej spełnialności
istnieje pełne wartościowanie
spełniające te formuły.
Jego ograniczenie do $\alpha$
należy do $T_\alpha$,
więc każdy poziom jest niepusty.

Jeśli $v\in T_\alpha$
oraz $\beta<\alpha$,
to $v\restriction\beta$
spełnia wszystkie
$\varphi_\xi$ z $\xi<\beta$
o nośniku zawartym w $\beta$:
były one również warunkami
dla $v$ na poziomie $\alpha$.
Zatem $v\restriction\beta
\in T_\beta$.
Zbiór $T$ jest więc drzewem
wysokości $\kappa$ z
poziomami $T_\alpha$.
Ponadto
$|T_\alpha|
 \le 2^{|\alpha|}<\kappa$
na mocy silnej graniczności.
\end{proof}

\begin{twierdzenie}[Własność drzewa daje zwartość zdaniową]
\label{thm:lc39-tree-to-compactness}
Niech $\kappa$ będzie nieprzeliczalną
silnie nieosiągalną liczbą
kardynalną mającą własność drzewa.
Wtedy zachodzi zasada
$\kappa$-zwartości infinitarnej
logiki zdaniowej z
Definicji~\ref{def:lc39-prop-compactness}.
\end{twierdzenie}
\begin{proof}
Ustalmy $\kappa$-lokalnie
spełnialną teorię $\Sigma$
o mocy co najwyżej $\kappa$.
Jeżeli $|\Sigma|<\kappa$,
jej spełnialność jest
bezpośrednio częścią hipotezy.
Załóżmy więc $|\Sigma|=\kappa$
i wybierzmy wyliczenie
$\Sigma=\{\varphi_\xi:\xi<\kappa\}$.

Na mocy
Lematu~\ref{lem:lc39-valuations-tree}
drzewo częściowych wartościowań
jest drzewem $\kappa$-owym.
Z własności drzewa ma gałąź
współkońcową. Jej suma
jest pełnym wartościowaniem
$v:\kappa\to\{0,1\}$,
bo gałąź przechodzi przez
każdy poziom $\alpha<\kappa$.

Ustalmy dowolne $\xi<\kappa$.
Z założenia
$|\operatorname{supp}(\varphi_\xi)|
 <\kappa$ i z regularności
$\kappa$ wynika, że
$\operatorname{supp}(\varphi_\xi)$
jest ograniczony w $\kappa$.
Możemy więc wybrać
$\alpha<\kappa$ większe od
$\xi$ i wszystkich
indeksów z tego nośnika.
Węzeł $v\restriction\alpha$
należy do $T_\alpha$,
a zatem spełnia
$\varphi_\xi$.
Ponieważ prawdziwość
zależy tylko od nośnika,
pełne $v$ również spełnia
$\varphi_\xi$.
Formuła była dowolna,
więc $v\models\Sigma$.
\end{proof}

\begin{twierdzenie}[Zwartość zdaniowa daje własność drzewa]
\label{thm:lc39-compactness-to-tree}
Niech $\kappa$ będzie
nieprzeliczalną regularną
liczbą kardynalną.
Załóżmy zasadę
$\kappa$-zwartości
infinitarnej logiki zdaniowej
z Definicji~\ref{def:lc39-prop-compactness}.
Wtedy każde drzewo
$\kappa$-owe ma gałąź
współkońcową.
\end{twierdzenie}
\begin{proof}
Niech $T$ będzie drzewem
$\kappa$-owym.
Ma ono dokładnie
$\kappa$ węzłów,
ponieważ ma wysokość $\kappa$
i mniej niż $\kappa$ węzłów
na każdym z $\kappa$ poziomów.
Wybieramy dla każdego
$t\in T$ odrębną zmienną
zdaniową $p_t$.
Po wyliczeniu węzłów
możemy utożsamić te zmienne
z rodziną $(p_\xi)_{\xi<\kappa}$.

Tworzymy teorię $\Sigma_T$
z dwóch rodzajów formuł.
Dla każdego poziomu
$\alpha<\kappa$ dodajemy
\[
 P_\alpha=
 \bigvee_{t\in T_\alpha}p_t.
\]
Jest to dopuszczalna formuła
infinitarna, ponieważ
$|T_\alpha|<\kappa$.
Dla każdej pary
nieporównywalnych węzłów
$s,t\in T$ dodajemy
\[
 Q_{s,t}=\neg(p_s\land p_t).
\]
Cała teoria ma moc
co najwyżej $\kappa$,
ponieważ liczba poziomów
wynosi $\kappa$,
a liczba par węzłów
jest co najwyżej
$\kappa^2=\kappa$.
Każda formuła zależy od
mniej niż $\kappa$ zmiennych.

Wykażemy lokalną spełnialność.
Niech $\Sigma_0\subseteq\Sigma_T$
ma moc mniejszą od $\kappa$.
Zbiór poziomów $\alpha$,
dla których
$P_\alpha\in\Sigma_0$,
ma moc mniejszą od $\kappa$.
Ponieważ $\kappa$ jest regularna,
istnieje $\delta<\kappa$
większa od wszystkich
tych poziomów.
Wybieramy dowolny
$t\in T_\delta$.
Definiujemy wartościowanie
\[
 v_t(p_s)=1
   \quad\Longleftrightarrow\quad s\le_T t.
\]
Wszystkie węzły na ścieżce
poprzedników $t$ są
porównywalne, więc
$v_t$ spełnia wszystkie
formuły typu $Q_{s,u}$.
Dla każdego poziomu
$\alpha$ uwzględnionego
w $\Sigma_0$ węzeł $t$
ma poprzednika na poziomie
$\alpha$, a więc
$v_t$ spełnia także
$P_\alpha$.
W konsekwencji
$v_t\models\Sigma_0$.

Z założonej zwartości
otrzymujemy wartościowanie
$v\models\Sigma_T$.
Dla każdego $\alpha<\kappa$
formuła $P_\alpha$ wymusza,
że istnieje węzeł
$t_\alpha\in T_\alpha$
z $v(p_{t_\alpha})=1$.
Wszystkie wybrane węzły
muszą być parami
porównywalne, ponieważ
dla nieporównywalnej pary
odpowiednia formuła
$Q_{s,t}$ zakazuje
jednoczesnej wartości $1$.
Zbiór węzłów o wartości $1$
jest zatem łańcuchem
przecinającym każdy poziom,
a więc gałęzią
współkońcową.
\end{proof}

\begin{wniosek}[Kombinatoryczna i zdaniowa zwartość]
\label{cor:lc39-compactness-equivalence}
Niech $\kappa$ będzie nieprzeliczalną
silnie nieosiągalną liczbą
kardynalną. Wtedy
własność drzewa, relacja
$\kappa\longrightarrow(\kappa)^2_2$
oraz zasada $\kappa$-zwartości
infinitarnej logiki zdaniowej
są równoważne.
\end{wniosek}
\begin{proof}
Twierdzenia~\ref{thm:lc39-tree-to-compactness}
i~\ref{thm:lc39-compactness-to-tree}
dają równoważność własności
drzewa i zwartości zdaniowej.
Wniosek~\ref{cor:lc39-tree-partition}
daje równoważność
własności drzewa i relacji
podziałowej.
\end{proof}

\begin{przyklad}[Teoria opisująca gałąź w skończonym drzewie]
\label{ex:lc39-finite-branch}
Rozważmy drzewo z trzema poziomami
\[
 \{\varnothing\},\quad
 \{(0),(1)\},\quad
 \{(0,0),(0,1),(1,0),(1,1)\},
\]
uporządkowane przez
przedłużanie skończonych ciągów.
Wprowadzamy zmienną $p_s$
dla każdego węzła $s$.
Formuła pierwszego poziomu
to $p_\varnothing$.
Drugi poziom daje
$p_{(0)}\lor p_{(1)}$,
a trzeci alternatywę
czterech odpowiednich zmiennych.
Dodatkowo wymagamy
$\neg(p_{(0)}\land p_{(1)})$
oraz analogicznych
ograniczeń dla
wszystkich nieporównywalnych
par węzłów.

Wartościowanie przyjmujące
wartość $1$ dokładnie
na węzłach
$\varnothing,(0),(0,1)$
spełnia wszystkie formuły:
na każdym poziomie
wybiera co najmniej jeden
węzeł, a wybrane
węzły są porównywalne.
Tak działa kodowanie
gałęzi w dowodzie
Twierdzenia~\ref{thm:lc39-compactness-to-tree}.
\end{przyklad}

\section{Dodatkowa konieczna własność regularności}

\begin{twierdzenie}[Relacja podziałowa wyklucza singularność]
\label{thm:lc39-partition-regular}
Niech $\kappa$ będzie
nieprzeliczalną liczbą kardynalną.
Jeśli
$\kappa\longrightarrow(\kappa)^2_2$,
to $\kappa$ jest regularna.
\end{twierdzenie}
\begin{proof}
Przypuśćmy, że
$\operatorname{cf}(\kappa)=\mu<\kappa$.
Wybieramy rosnący
ciąg $(\gamma_\xi)_{\xi<\mu}$
współkońcowy w $\kappa$
i dzielimy $\kappa$
na $\mu$ rozłącznych
przedziałów $I_\xi$,
z których każdy ma
moc mniejszą od $\kappa$.
Można je otrzymać,
biorąc przedziały między
kolejnymi punktami ciągu;
występujące ewentualnie
puste przedziały nie
zmieniają argumentu.

Kolorujemy parę
$\{\alpha,\beta\}$ kolorem
$0$, gdy oba elementy
należą do jednego $I_\xi$,
oraz kolorem $1$
w przeciwnym wypadku.
Jeżeli $H\subseteq\kappa$
jest jednorodny
w kolorze $0$, to
wszystkie jego elementy
należą do jednego
przedziału, zatem
$|H|<\kappa$.
Jeśli $H$ jest
jednorodny w kolorze $1$,
to zawiera co najwyżej
jeden punkt każdego
przedziału, a więc
$|H|\le\mu<\kappa$.
Nie istnieje zatem
jednorodny zbiór
mocy $\kappa$;
sprzeczność z hipotezą.
\end{proof}

\section{Podsumowanie i zakres dalszego rozwoju}

\begin{proof}[Dowód Twierdzenia~\ref{thm:lc39-main-equivalence}]
Implikacja (1)$\Rightarrow$(3)
została udowodniona w
Twierdzeniu~\ref{thm:lc39-tree-to-partition}.
Implikacja (3)$\Rightarrow$(2)
wynika przez podstawienie
$\mu=2$.
Implikacja (2)$\Rightarrow$(1)
jest Twierdzeniem~\ref{thm:lc39-partition-to-tree}.
Równoważność punktów (1) i (4)
zawierają
Twierdzenia~\ref{thm:lc39-tree-to-compactness}
oraz~\ref{thm:lc39-compactness-to-tree}.
Wszystkie warunki są więc
równoważne.
\end{proof}

\begin{uwaga}[Zwartość zdaniowa i logika pierwszego rzędu]
\label{rem:lc39-first-order}
Zasada zwartości użyta w tym
rozdziale dotyczy
\emph{infinitarnej logiki zdaniowej}
z formułami zależnymi od
mniej niż $\kappa$ zmiennych.
Nie utożsamiamy jej bez dowodu
z pełnym twierdzeniem
o zwartości logiki pierwszego
rzędu $L_{\kappa,\kappa}$,
dopuszczającej również bloki
kwantyfikatorów długości
mniejszej niż $\kappa$.
Pełne przejście między tymi
językami wymaga osobnego
dowodu konstrukcji modelu
i odpowiedniej wersji
rachunku Henkina.
\end{uwaga}

\begin{uwaga}[Zależność od silnej nieosiągalności]
\label{rem:lc39-inaccessible-usage}
W dowodzie z własności
drzewa do relacji podziałowej
kluczowe było oszacowanie
$\mu^{|\delta|}<\kappa$,
które zapewnia małe poziomy
drzewa profili.
Z kolei przy zwartości
zdaniowej małe poziomy
drzewa częściowych
wartościowań daje
$2^{|\delta|}<\kappa$.
Nie można pominąć
tych oszacowań i twierdzić,
że samo istnienie
drzewa o wysokości
$\kappa$ wystarcza.
\end{uwaga}

\begin{uwaga}[Status formalnej weryfikacji]
\label{rem:lc39-formal}
Przedstawione dowody są
dowodami matematycznymi
w tekście LaTeX.
Automatyczna kontrola
odnośników, środowisk
i skończonych testów
kombinatorycznych nie
zastępuje formalizacji
w Lean. Status formalizacji
wymaga odrębnego,
rzeczywistego przebiegu
kompilatora Lean.
\end{uwaga}

\chapter{Logika rozmyta i zbiory rozmyte}\label{ch:lst-fuzzy-logic}

\section{Wielowartościowość i zbiory rozmyte}
Stopień prawdziwości zdania nie jest prawdopodobieństwem zdarzenia.
W logice rozmytej interpretacje spójników zależą od wybranego systemu,
a nie tylko od zbioru wartości $[0,1]$.
\begin{definicja}\label{def:lst-fuzzy-set}
Zbiór rozmyty $A$ w $X$ jest funkcją przynależności $\mu_A:X\to[0,1]$.
Klasyczne podzbiory odpowiadają funkcjom charakterystycznym $\chi_E$.
Porządek zawierania określa nierówność punktowa funkcji przynależności.
Jądro to $\{x:\mu_A(x)=1\}$, nośnik to $\{x:\mu_A(x)>0\}$, a
przekrój na poziomie $\alpha\in(0,1]$ to $A_\alpha=\{x:\mu_A(x)\geq\alpha\}$.
\end{definicja}
\begin{twierdzenie}\label{thm:lst-fuzzy-cuts}
Przekroje wyznaczają zbiór rozmyty:
\[
\mu_A(x)=\sup(\{0\}\cup\{\alpha\in(0,1]:x\in A_\alpha\}).
\]
Ponadto dla $0<\alpha\leq\beta\leq1$ mamy $A_\beta\subseteq A_\alpha$,
a $A_\alpha=\bigcap_{0<\beta<\alpha}A_\beta$.
\end{twierdzenie}
\begin{proof}
Dodatnie poziomy zawierające $x$ tworzą przedział $(0,\mu_A(x)]$,
pusty gdy $\mu_A(x)=0$. Stąd pierwszy wzór i monotoniczność.
Jeśli $\mu_A(x)<\alpha$, istnieje $\beta$ między $\mu_A(x)$ a $\alpha$,
więc $x$ nie należy do przecięcia. Druga inkluzja wynika z monotoniczności.
\end{proof}
\begin{przyklad}[Trójkątna funkcja przynależności]
Dla $a<b<c$ połóżmy
\[
\mu(x)=\max\left(0,\min\left(\frac{x-a}{b-a},\frac{c-x}{c-b}\right)\right).
\]
Rozwiązanie obu nierówności liniowych daje
$A_\alpha=[a+\alpha(b-a),c-\alpha(c-b)]$.
\begin{center}
\begin{tikzpicture}[x=1.1cm,y=2.6cm,>=Stealth]
\draw[->] (0,0)--(6,0) node[right]{$x$};
\draw[->] (0,0)--(0,1.15) node[above]{$\mu(x)$};
\fill[blue!12] (1,0)--(3,1)--(5,0)--cycle;
\draw[thick,blue!70!black] (0,0)--(1,0)--(3,1)--(5,0)--(6,0);
\draw[dashed] (0,0.5) node[left]{$\alpha$}--(4,0.5)--(4,0);
\draw[dashed] (2,0.5)--(2,0);
\draw[very thick,teal!70!black] (2,0)--(4,0) node[midway,below=5pt]{$A_\alpha$};
\draw (0,1) node[left]{$1$};
\end{tikzpicture}
\end{center}
\end{przyklad}
\section{Normy trójkątne i konormy}
Kres dolny i górny w porządku punktowym to minimum i maksimum.
Ich przekroje są odpowiednio przecięciem i sumą przekrojów.
\begin{definicja}\label{def:lst-fuzzy-tnorm}
Norma trójkątna $T:[0,1]^2\to[0,1]$ jest przemienna,
łączna, niemalejąca w obu argumentach i spełnia $T(a,1)=a$.
\end{definicja}
\begin{lemat}\label{lem:lst-fuzzy-tnorm-bound}
$T(a,0)=0$ oraz $T(a,b)\leq\min(a,b)$.
\end{lemat}
\begin{proof}
$0\leq T(a,0)\leq T(1,0)=0$ i $T(a,b)\leq T(a,1)=a$;
przemienność daje ograniczenie przez $b$.
\end{proof}
\begin{przyklad}\label{ex:lst-fuzzy-three-tnorms}
Trzy ciągłe normy trójkątne to
\[
T_G(a,b)=\min(a,b),\quad T_P(a,b)=ab,\quad
T_L(a,b)=\max(0,a+b-1).
\]
Dla minimum i iloczynu aksjomaty wynikają z własności porządku
i mnożenia. Dla $T_L$ łączność wynika z równości obu sposobów
złożenia do $\max(0,a+b+c-2)$; pozostałe aksjomaty wynikają wprost
ze wzoru. Dla $a=b=0{,}7$ wartości wynoszą $0{,}7$, $0{,}49$, $0{,}4$.
\end{przyklad}
Konorma dualna względem negacji $N(a)=1-a$ ma postać
$S(a,b)=1-T(1-a,1-b)$. Przemienność i łączność wynikają z
przemienności i łączności $T$, monotoniczność z dwukrotnego
odwrócenia porządku, a $S(a,0)=a$ z elementu neutralnego $T$.
Definicja daje $N(S(a,b))=T(N(a),N(b))$, a inwolucja $N$
daje drugie prawo de Morgana.

\section{Residuum i implikacja}
\begin{twierdzenie}\label{thm:lst-fuzzy-residuation}
Niech $T$ będzie normą trójkątną lewostronnie ciągłą. Wówczas
\[
a\Rightarrow_T b=\sup\{c:T(a,c)\leq b\}
\]
spełnia $T(a,c)\leq b$ wtedy i tylko wtedy, gdy $c\leq a\Rightarrow_T b$.
\end{twierdzenie}
\begin{proof}
Zbiór pod supremum jest niepustym przedziałem początkowym na mocy
Lematu~\ref{lem:lst-fuzzy-tnorm-bound} i monotoniczności.
Jeśli jego supremum $r>0$, ciąg elementów tego zbioru rosnący do $r$
i lewostronna ciągłość dają $T(a,r)\leq b$; dla $r=0$ wynika to z lematu.
Zbiór jest więc równy $[0,r]$.
\end{proof}
W szczególności $a\Rightarrow_T b=1$ wtedy i tylko wtedy, gdy $a\leq b$,
przez podstawienie $c=1$, a $T(a,a\Rightarrow_T b)\leq b$ przez
podstawienie $c=a\Rightarrow_T b$.
\begin{przyklad}\label{ex:lst-fuzzy-implications}
\[
a\Rightarrow_G b=\begin{cases}1&a\leq b,\\b&a>b,\end{cases}
\quad
a\Rightarrow_P b=\begin{cases}1&a\leq b,\\b/a&a>b,\end{cases}
\quad
a\Rightarrow_L b=\min(1,1-a+b).
\]
Dla $a>b$ nierówności $\min(a,c)\leq b$ i $ac\leq b$ są
równoważne odpowiednio $c\leq b$ i $c\leq b/a$.
Dla $T_L$ nierówność jest równoważna $c\leq1-a+b$.
Negacja rezydualna $a\Rightarrow_T0$ jest dla $G,P$ równa $1$
w zerze i $0$ poza zerem, a dla $L$ równa $1-a$.
\end{przyklad}

\section{Semantyka i algebra}
Wartościowanie przypisuje zmiennym zdaniowym liczby w $[0,1]$.
Spójniki $\bot,\wedge,\vee,\mathbin{\&},\to$ interpretujemy przez
$0,\min,\max,T,\Rightarrow_T$. Negacja to $\varphi\to\bot$.
Tautologia ma wartość jeden przy każdym wartościowaniu.
Konsekwencja $\Gamma\models_T\varphi$ oznacza, że wartość jeden
wszystkich przesłanek pociąga wartość jeden wniosku.
\begin{przyklad}
Dla $v(p)=1/2$ i normy Łukasiewicza mamy $v(p\vee\neg p)=1/2$,
a $v(p\mathbin{\&}p)=0\ne v(p\wedge p)=1/2$.
Koniunkcja kratowa i monoidalna nie są więc tym samym spójnikiem.
\end{przyklad}
\begin{definicja}\label{def:lst-fuzzy-residuated-lattice}
Ograniczona przemienna całkowita krata rezyduowana jest kratą
$(L,\wedge,\vee,0,1)$ z przemiennym monoidem $(L,\mathbin{\&},1)$
oraz operacją $\to$ spełniającą
$a\mathbin{\&}c\leq b\iff c\leq a\to b$. Element $1$ jest największy.
Algebra MTL spełnia dodatkowo prelinearność
$(a\to b)\vee(b\to a)=1$; algebra BL ponadto podzielność
$a\wedge b=a\mathbin{\&}(a\to b)$, a algebra MV może być
scharakteryzowana jako algebra BL z negacją inwolucyjną.
\end{definicja}
\begin{stwierdzenie}
Trzy struktury z Przykładu~\ref{ex:lst-fuzzy-three-tnorms}
są algebrami BL; struktura Łukasiewicza jest algebrą MV.
\end{stwierdzenie}
\begin{proof}
Residuacja wynika z Twierdzenia~\ref{thm:lst-fuzzy-residuation},
a prelinearność z liniowości porządku. Dla $a\leq b$
podzielność wynika z $a\to b=1$. Dla $a>b$ podstawienie wzorów
implikacji daje odpowiednio $\min(a,b)=b$, $a(b/a)=b$ i
$\max(0,a+1-a+b-1)=b$. Dla Łukasiewicza $\neg a=1-a$.
\end{proof}

\section{Logika pierwszego rzędu}
W niepustym uniwersum $M$ interpretujemy symbole funkcyjne klasycznie,
a predykaty $n$-argumentowe jako funkcje $M^n\to[0,1]$.
Dla wartościowania $s$ zmiennych
\[
\|\forall x\,\varphi\|_s=\inf_{a\in M}\|\varphi\|_{s[x\mapsto a]},
\qquad
\|\exists x\,\varphi\|_s=\sup_{a\in M}\|\varphi\|_{s[x\mapsto a]}.
\]
Kresy istnieją przez zupełność porządkową $[0,1]$.
Jeżeli $M=\N$ oraz $P(n)=1-1/(n+1)$, to $\|\exists xP(x)\|=1$,
ale żadna wartość $P(n)$ nie jest równa jeden.

\section{Relacje rozmyte}
Dla $R:X\times Y\to[0,1]$ i $S:Y\times Z\to[0,1]$ definiujemy
$(S\circ_T R)(x,z)=\sup_{y\in Y}T(R(x,y),S(y,z))$,
gdzie $\sup\varnothing=0$.
\begin{stwierdzenie}
Dla lewostronnie ciągłej $T$ składanie jest łączne.
\end{stwierdzenie}
\begin{proof}
Residuacja pokazuje, że $T(a,\cdot)$ zachowuje dowolne suprema:
$T(a,\sup C)\leq b$ wtedy i tylko wtedy, gdy
$\sup C\leq a\Rightarrow_T b$, równoważnie
$T(a,c)\leq b$ dla wszystkich $c\in C$.
Oba złożenia trzech relacji dają więc supremum po parach
elementów pośrednich z tej samej wartości na mocy łączności $T$.
\end{proof}

\section{Liczby rozmyte i zasada rozszerzania}
Liczba rozmyta to zbiór rozmyty w $\R$ z górnie półciągłą funkcją
przynależności osiągającą jeden, ograniczonym nośnikiem i przedziałowymi
dodatnimi przekrojami. Te przekroje są niepuste i zwarte.
Dla ciągłej $f:\R^2\to\R$ definiujemy
$\mu_{f(A,B)}(z)=\sup_{f(x,y)=z}\min(\mu_A(x),\mu_B(y))$.
\begin{twierdzenie}\label{thm:lst-fuzzy-extension-cuts}
Dla liczb rozmytych i $\alpha>0$, $\alpha\leq1$,
$(f(A,B))_\alpha=f(A_\alpha\times B_\alpha)$.
\end{twierdzenie}
\begin{proof}
Inkluzja od prawej do lewej wynika z definicji. Dla punktu
lewej strony wybieramy pary o obrazie $z$ i stopniach co najmniej
$\alpha-1/n$. Dla dużych $n$ należą do zwartego
$A_{\alpha/2}\times B_{\alpha/2}$. Podciąg zbiega do pary
o obrazie $z$ przez ciągłość $f$ i o stopniach co najmniej $\alpha$
przez górną półciągłość funkcji przynależności.
\end{proof}
Przy dodawaniu przedziałów przekroje sumują się końcami, ponieważ
dodawanie jest ciągłe i monotoniczne na zwartym spójnym prostokącie.
Przykładowo trójkątne liczby $(1,2,4)$ i $(0,3,5)$ dają $(1,5,9)$.

\section{Wnioskowanie rozmyte}
Dwie reguły Mamdaniego z aktywacjami $0{,}25$, $0{,}75$
i konsekwencjami $\nu_1,\nu_2$ dają funkcję
$\nu(y)=\max(\min(0{,}25,\nu_1(y)),\min(0{,}75,\nu_2(y)))$.
Dla całkowalnej $\nu$ o dodatniej całce i skończonym pierwszym momencie
defuzyfikacja środkiem ciężkości daje
\[
y_*=\frac{\int_{\R} y\nu(y)\,\mathrm d\lambda_1(y)}
{\int_{\R}\nu(y)\,\mathrm d\lambda_1(y)}.
\]
W systemie Sugeno zerowego rzędu, dla konsekwencji liczbowych $2,6$,
znormalizowana średnia wynosi
$(0{,}25\cdot2+0{,}75\cdot6)/(0{,}25+0{,}75)=5$.
Nie tworzy się tu agregowanej funkcji przynależności wyjścia.

\end{document}
